\DocumentMetadata{
  tagging=on,
  lang=en-US,
  pdfstandard=UA-2,
  tagging-setup={
%    math/alt/use,
    math/setup=mathml-SE
  }
}
\documentclass[openany]{book}

\usepackage{../wtrench}
% !TEX root = TRENCH_DIFFEQ.tex

\usepackage{mathtools}
\usepackage{lua-unicode-math}
\usepackage[usenames, dvipsnames, svgnames, table]{xcolor}
%\usepackage[T1]{fontenc}
%\usepackage{textcomp}
%\renewcommand{\rmdefault}{ptm}
%\usepackage[scaled=0.92]{helvet}
%\usepackage[psamsfonts]{amsfonts}
\usepackage{amsmath, amsbsy}
\usepackage{calc}
\usepackage{booktabs}
\usepackage[per-mode = symbol]{siunitx}
\usepackage{capt-of}
\usepackage{zref-clever}
\usepackage{etoolbox}
\usepackage{pgfplots}
\pgfplotsset{compat=1.18}

\usepackage[%dvips,
 bookmarks = true,
 colorlinks=true,
 plainpages = false,
 citecolor = blue,
 urlcolor = blue,
 filecolor = blue,
 pageanchor = true]
{hyperref}

\zcRefTypeSetup{chapter}{cap}
\zcRefTypeSetup{definition}{cap}
\zcRefTypeSetup{equation}{cap}
\zcRefTypeSetup{example}{cap}
\zcRefTypeSetup{figure}{cap}
\zcRefTypeSetup{section}{cap}
\zcRefTypeSetup{table}{cap}
\zcRefTypeSetup{theorem}{cap}
\zcRefTypeSetup{enumXi}{
  Name-sg=Exercise,
  Name-pl=Exercises,
  name-sg=exercise,
  name-pl=exercises,
  rangesep={\textendash},
  cap
}
\zcRefTypeSetup{enumXvii}{
  Name-sg=Exercise,
  Name-pl=Exercises,
  name-sg=exercise,
  name-pl=exercises,
  rangesep={\textendash},
  cap
}

\newtoggle{boundaryvalues}

\EditInstance{heading}{section}{title-decls=\raggedright}

\DeclareMathOperator{\tr}{tr}

\DeclareSIUnit\dyne{\text{dyn}}
\DeclareSIUnit[quantity-product = {}]\degreeFahrenheit{\text{\textdegree F}}
\DeclareSIUnit\feet{\text{ft}}
\DeclareSIUnit\foot{\text{ft}}
\DeclareSIUnit\gallon{\text{gal}}
\DeclareSIUnit\pound{\text{lb}}

\newcommand{\vect}[1]{\mathbf{#1}}
\newcommand{\ivect}{\vect{i}}
\newcommand{\jvect}{\vect{j}}

\newcommand{\laplace}{{\mathcal L}}

\DeclareTextCommandDefault{\textRightarrow}{\ensuremath{\Rightarrow}}

\newcommand{\amatrix}[2]{
 \left[
  \begin{matrix}#1\end{matrix}
 \hspace{\tabcolsep}\middle|\hspace{\tabcolsep}
  \begin{matrix}#2\end{matrix}
 \right]
}

\ExplSyntaxOn
% TODO?
\keys_define:nn { zref-clever/reference }
  {
    CAP .bool_set:N = \l_tmpa_bool,
    CAP .default:n = false
  }
\ExplSyntaxOff

\usepackage{xstring}

\newcommand*{\personHref}[2]{%
 \IfStrEqCase{#1}{%
  {Libby}{\href{https://www.nobelprize.org/prizes/chemistry/1960/libby/facts/}{\color{blue}\slshape#2}}%
  {Malthus}{\href{http://en.wikipedia.org/wiki/Thomas_Robert_Malthus}{\color{blue}\slshape#2}}%
 }[\href{https://mathshistory.st-andrews.ac.uk/Biographies/#1/}{\color{blue}\slshape#2}]%
}



%\colorlet{blue}{DarkBlue}
\colorlet{red}{DarkRed}

\setcounter{secnumdepth}{2}
\usepackage[a4paper]{geometry}
\usepackage{float}
\usepackage{upref}

\usepackage{array}
\newcolumntype{H}{>{\setbox0=\hbox\bgroup}c<{\egroup}@{}}

\usepackage{graphicx}
\graphicspath{%
    {converted_graphics/}% inserted by PCTeX
    {EPS/}% inserted by PCTeX
}
%\usepackage{epstopdf}
%\epstopdfsetup{update}

\numberwithin{example}{section}
\numberwithin{equation}{section}
\numberwithin{theorem}{section}
\numberwithin{table}{section}
\numberwithin{figure}{section}

%\newcommand{\place}{\bigskip\hrule\bigskip\noindent}
%\newcommand{\threecol}[3]{\left[\begin{array}{r}#1\\#2\\#3\end{array}\right]}
\newcommand{\threecol}[3]{\begin{bmatrix}#1\\#2\\#3\end{bmatrix}}
%\newcommand{\threecolj}[3]{\left[\begin{array}{r}#1\\[1\jot]#2\\[1\jot]#3\end{array}\right]}
\newcommand{\threecolj}{\threecol}
\newcommand{\lims}[2]{\,\bigg|_{#1}^{#2}}
%\newcommand{\twocol}[2]{\left[\begin{array}{l}#1\\#2\end{array}\right]}
\newcommand{\twocol}[2]{\begin{bmatrix}#1\\#2\end{bmatrix}}
%\newcommand{\ctwocol}[2]{\left[\begin{array}{c}#1\\#2\end{array}\right]}
\newcommand{\ctwocol}{\twocol}
%\newcommand{\cthreecol}[3]{\left[\begin{array}{c}#1\\#2\\#3\end{array}\right]}
\newcommand{\cthreecol}{\threecol}
\newcommand{\eqline}[1]{\centerline{\hfill$\displaystyle#1$\hfill}}
%\newcommand{\twochar}[4]{\left|\begin{array}{cc}
%#1-\lambda&#2\\#3&#4-\lambda\end{array}\right|}
\newcommand{\twochar}[4]{\begin{vmatrix}#1-\lambda&#2\\#3&#4-\lambda\end{vmatrix}}
%\newcommand{\twobytwo}[4]{\left[\begin{array}{rr}
%#1&#2\\#3&#4\end{array}\right]}
\newcommand{\twobytwo}[4]{\begin{bmatrix}#1&#2\\#3&#4\end{bmatrix}}
%\newcommand{\threechar}[9]{\left[\begin{array}{ccc}
%#1-\lambda&#2&#3\\#4&#5-\lambda&#6\\#7&#8&#9-\lambda\end{array}\right]}
\newcommand{\threechar}[9]{\begin{bmatrix}#1-\lambda&#2&#3\\#4&#5-\lambda&#6\\#7&#8&#9-\lambda\end{bmatrix}}
%\newcommand{\threebythree}[9]{\left[\begin{array}{rrr}
%#1&#2&#3\\#4&#5&#6\\#7&#8&#9\end{array}\right]}
\newcommand{\threebythree}[9]{\begin{bmatrix}#1&#2&#3\\#4&#5&#6\\#7&#8&#9\end{bmatrix}}
% TODO
%\newcommand{\solutionpart}[1]{\bigskip\noindent{%\underbar
%\color{blue}\scshape Solution({\bfseries #1})\ }
%}
\newcommand{\Cex}{\fbox{\textcolor{red}{C}}\, }
\newcommand{\CGex}{\fbox{\textcolor{red}{C/G}}\, }
\newcommand{\Lex}{\fbox{\textcolor{red}{L}}\, }
%\newcommand{\matfunc}[3]{\left[\begin{array}{cccc}#1_{11}(t)&#1_{12}(t)&\cdots
%&#1_{1#3}(t)\\#1_{21}(t)&#1_{22}(t)&\cdots&#1_{2#3}(t)\\\vdots&
%\vdots&\ddots&\vdots\\#1_{#21}(t)&#1_{#22}(t)&\cdots&#1_{#2#3}(t)
%\end{array}\right]}
\newcommand{\matfunc}[3]{\begin{bmatrix}#1_{11}(t)&#1_{12}(t)&\cdots
&#1_{1#3}(t)\\#1_{21}(t)&#1_{22}(t)&\cdots&#1_{2#3}(t)\\\vdots&
\vdots&\ddots&\vdots\\#1_{#21}(t)&#1_{#22}(t)&\cdots&#1_{#2#3}(t)
\end{bmatrix}}

%\newcommand{\col}[2]{\left[\begin{array}{c}#1_1\\#1_2\\\vdots\\#1_#2\end{array}\right]}
\newcommand{\col}[2]{\begin{bmatrix}#1_1\\#1_2\\\vdots\\#1_#2\end{bmatrix}}
%\newcommand{\colfunc}[2]{\left[\begin{array}{c}#1_1(t)\\#1_2(t)\\\vdots\\#1_#2(t)\end{array}\right]}
\newcommand{\colfunc}[2]{\begin{bmatrix}#1_1(t)\\#1_2(t)\\\vdots\\#1_#2(t)\end{bmatrix}}

%\newcommand{\cthreebythree}[9]{\left[\begin{array}{ccc}
%#1&#2&#3\\#4&#5&#6\\#7&#8&#9\end{array}\right]}
\newcommand{\cthreebythree}[9]{\begin{bmatrix}#1&#2&#3\\#4&#5&#6\\#7&#8&#9\end{bmatrix}}

\newcommand{\technology}{\bigskip\noindent\fbox{\color{red}USING
TECHNOLOGY}\par}

\newcommand{\bigbox}[1]{\vskip10pt\vbox{\hrule\hbox{\vrule\kern5pt
\vbox{\noindent
\strut\kern5pt\mbox{}\newline#1\mbox{}\newline
\kern5pt\strut}\kern5pt\vrule}\hrule}}

\newcommand{\blankbox}[1]{\vskip10pt\vbox{\hbox{\kern5pt
\vbox{\noindent\strut\kern5pt\vspace*{-15pt}#1\vspace*{-20pt}
\kern5pt\strut}\kern5pt}}}

%\newenvironment{newtable}[1]{ \begin{quote}
%Table\space\refstepcounter{table}\thetable. #1}{\end{quote}}

\newcommand{\sectionanswers}[1]{%
 \ParseLaTeXeHeading{section}{\BooleanTrue}{
  heading-decls = \normalfont\normalsize\bfseries,
  after-vspace = \medskipamount,
  before-vspace = \medskipamount,
  bookmark = {Section #1 Answers},
%  toc = {Section #1 Answers},
  running = {Section #1 Answers}
 }{Section #1 Answers, pp.~\pageref{exer:#1.1}\thinspace ff}%
}

\newcommand{\indexterm}[1]{\emph{\color{blue}#1}\index{#1}}

\setlength{\oddsidemargin}{15.5pt}
\setlength{\evensidemargin}{15.5pt}

\usepackage{makeidx}
\makeindex

\title{Elementary Differential Equations}
\author{William F. Trench}

\makeatletter
\newcommand{\maybeStepEnum}{
 \ifnum\@listdepth<2
  \stepcounter{enumi}
 \fi
}
\makeatother
\newenvironment{sectionSolutions}{
 \begin{enumerate}[item-label={\thesection.\arabic{enumi}.}]
}{
 \end{enumerate}
}

%\toggletrue{boundaryvalues}
\toggletrue{solutionManual}

% TODO see https://github.com/josephwright/siunitx/issues/877
\setlength\hfuzz{8pt}

\begin{document}

\frontmatter

% !TEX root = ../TRENCH_DIFFEQ.tex

\begin{center}


\thispagestyle{empty}
{
\makeatletter
\sloppy
\tagstructbegin{tag=Title}
\tagmcbegin{tag=Title}
\bfseries\huge\MakeUppercase\@title
\tagmcend
\tagstructend
\par
}

 \vspace*{.15in}

\begin{center}
 \includegraphics[width=5.67in,height=3.3in,keepaspectratio,
  alt={Arrows moving through the plain up and left, but circling clockwise around the origin.  Lines follow the arrows, creating a denser line that starts and ends along y=-x, but also circles clockwise around the origin.}
 ]{\iftoggle{solutionManual}{../EPS/cover}{cover}}
\end{center}
\vspace*{.15in}
\bfseries\LARGE
\makeatletter
\href{http://ramanujan.math.trinity.edu/wtrench/index.shtml}
{\@author}
\makeatother
%{William F. Trench}
\medskip

\large
Andrew G. Cowles Distinguished Professor Emeritus\\
Department of Mathematics\\
Trinity University \\
San Antonio, Texas, USA
%\\\href{mailto:wtrench@trinity.edu}{wtrench@trinity.edu}
% He passed in 2016
\large

\begin{quote}
This book has been judged to meet the evaluation criteria set by the
Editorial  Board
of the American Institute of Mathematics in connection with the Institute's
\href{http://www.aimath.org/textbooks/}
{Open
Textbook Initiative}.
It may be copied, modified, redistributed, translated,  and
built upon  subject to the  Creative Commons
\href{http://creativecommons.org/licenses/by-nc-sa/3.0/deed.en_G}
{Attribution-NonCommercial-ShareAlike 3.0 Unported
License}.
\end{quote}

\iftoggle{solutionManual}{}{
\vspace{.25in}
\href{http://ramanujan.math.trinity.edu/wtrench/texts/TRENCH_DE_STUDENT_MANUAL.PDF}
{\bfseries FREE DOWNLOAD: STUDENT SOLUTIONS MANUAL}
}

\end{center}
\newpage

\vspace*{2in}

\noindent
\medskip

\noindent
%Free Edition 1.01 (December 2013)
Tagged Free Edition 1.1 (\makeatletter\@date\makeatother)%Summer 2026)

\medskip
\thispagestyle{empty}

\noindent
This book was  published previously by Brooks/Cole Thomson Learning, 2001.
This free edition is made available in the hope that it will be useful
as a textbook or reference.
 Reproduction  is
permitted
for any valid  noncommercial educational, mathematical, or scientific
purpose.
However,
charges for
profit beyond reasonable printing costs
are  prohibited.

%\vspace*{.51in}

\newpage

\noindent
\thispagestyle{empty}
\vspace*{3in}

\begin{center}\color{blue}\bfseries\Huge TO BEVERLY\end{center}

\newpage

\tableofcontents

\rmfamily


\mainmatter

\raggedright
\allowdisplaybreaks

\AddToHook{cmd/item/before}{\maybeStepEnum}

\chapter{Introduction}

\stepcounter{section}
\section{Basic Concepts}

\begin{sectionSolutions}
\item
\begin{alist}
\item If $y=ce^{2x}$, then $y'=2ce^{2x}=2y$.

\item  If $y=\dst \frac{x^2}{3}+\frac{c}{ x}$, then
$y'=\dst\frac{2x}{3}-\frac{c}{ x^2}$, so $xy'+y=\dst\frac{2x^2}{3}-\frac{c}{ x
}+\frac{x^2}{3}+\frac{c}{ x}=x^2$.

\item If
\[
y=\dst\frac{1}{2}+ce^{-x^2}, \quad\text{then}\quad
y'=-2xce^{-x^2}
\]
and
\[
y'+2xy=-2xce^{-x^2}+2x\dst\left(\frac{1}{2}+ce^{-x^2}\right)
=-2xce^{-x^2}+x+2cxe^{-x^2}=x.
\]

\item
If
\[
y=\dfrac{1+ce^{-x^2/2}}{1-ce^{-x^2/2}}
\]
then
\begin{align*}
y'&=\frac{(1-ce^{-x^2/2})(-cxe^{-x^2/2})-(1+ce^{-x^2/2})cxe^{-x^2/2}}{
(1-cxe^{-x^2/2})^2}\\[2\jot]
&=\frac{-2cxe^{-x^2/2}}{(1-ce^{-x^2/2})^2}
\end{align*}
and
\begin{align*}
y^2-1&=\left(\frac{1+ce^{-x^2/2}}{1-ce^{-x^2/2}}\right)^2-1\\[2\jot]
&=\frac{(1+ce^{-x^2/2})^2-(1-ce^{-x^2/2})^2}{(1-ce^{-x^2/2})^2}\\[2\jot]
&=\frac{4ce^{-x^2/2}}{(1-ce^{-x^2/2})^2},
\end{align*}
so
\[
2y'+x(y^2-1)=\frac{-4cx+4cx}{(1-ce^{-x^2/2})^2}=0.
\]


\item
If $y=\dst\tan\left( \frac{x^3}{3}+c\right)$, then
$y'=x^2\dst\sec^2\left(\frac{x^3}{3}+c\right)=
x^2\left(1+\tan^2\dst\left(\frac{x^3}{3+c}\right)\right)=
x^2(1+y^2)$.

\item
If
\begin{align*}
y&=(c_1+c_2x)e^x+\sin x+x^2,\quad\text{then}\\
y'&=(c_1+2c_2x)e^x+\cos x+2x,\\
y'&=(c_1+3c_2x)e^x-\sin x+2,\\
\intertext{and}
y''-2y'+y&=c_1e^x(1-2+1)+c_2xe^x(3-4+1)\\
&-\sin x-2\cos x+\sin x+2-4x+x^2\\
&=-2 \cos x+x^2-4x+2.
\end{align*}

\item If $y=c_1e^x+c_2x+\dst\frac{2}{ x}$, then
$y'=c_1e^x+c_2-\dst\frac{2}{ x^2}$ and $y''=c_1e^x+\dst\frac{4}{ x^3}$, so
$(1-x)y''+xy'- y=c_1(1-x+x-1)+c_2(x-x)+\dfrac{4(1-x)}{ x^3}-\dfrac{2}{
x}-\dfrac{2}{ x}= \dfrac{4(1-x-x^2)}{ x^3}$

\item
If $y=\dst\frac{c_1\sin x+c_2 \cos x}{ x^{1/2}}+4x+8$
then
 $y'=\dst\frac{c_1\cos x-c_2 \sin x}{ x^{1/2}}
-\frac{c_1\sin x+c_2 \cos x}{ 2x^{3/2}}+4$ and
$y''=-\dst\frac{c_1\sin x+c_2 \cos x}{ x^{1/2}}
-\frac{c_1\sin x-c_2 \cos x}{ x^{3/2}}+\frac{3}{4}
\frac{c_1\sin x+c_2 \cos x}{ x^{5/2}}$, so
$\dst x^2y''+xy'+\left(x^2-\frac{1}{4}\right)y=
c_1\left(-x^{-3/2}\sin x-x^{1/2}\cos x+
\frac{3}{4}x^{-1/2}\sin x +x^{1/2}\cos x-\right.$\\
$\dst\left.
\frac{1}{2}x^{-1/2}\sin x+x^{3/2}\sin
x-\frac{1}{4}x^{-1/2}
\sin x\right)+
\dst c_2\left(-x^{-3/2}\cos x+x^{1/2}\sin x
+\frac{3}{4}x^{-1/2}\cos x \right.$\\
$\dst\left.-x^{1/2}\sin x-\frac{1}{2}x^{-1/2}\cos x+
 x^{3/2}\cos
x-\frac{1}{4}x^{-1/2}\cos x\right)+4x+\left(x^2-\frac{1}{4}\right)
(4x+8)=
4x^3+8x^2+3x-2$.
\end{alist}

\item
\begin{alist}
\item If $y'=-xe^x$, then
$y=-xe^x+\int e^x\,dx+c=(1-x)e^x+c$,
and $y(0)=1\Rightarrow 1=1+c$, so $c=0$ and $y=(1-x)e^x$.

\item
If
$y'=x\sin x^2$, then
$y=-\dst\frac{1}{2}\cos x^2+c$;
$\dst y\left({\sqrt\frac{\pi}{2}}\right)=1 \Rightarrow 1=0+c$,
so $c=1$ and $y=1-\dst\frac{1}{2}\cos x^2$.

\item Write $y'=\tan x=\dst\frac{\sin x}{\cos x}=-\dst\frac{1}{\cos
x}\frac{d}{ dx}(\cos x)$. Integrating this yields $y=-\ln|\cos x|+c$;
 $y(\pi/4)=3\Rightarrow 3=-\ln\left(\cos(\pi/4)\right)+c$, or
$3=\ln\sqrt2+c$, so $c=3-\ln\sqrt2$, so
$y=-\ln(|\cos x|)+3-\ln\sqrt2=3-\ln(\sqrt2|\cos x|)$.

\item If $y''=x^4$, then $y'=\dst\frac{x^5}{5}+c_1$; $y'(2)=-1
\Rightarrow\dst\frac{32}{5}+c_1=-1\Rightarrow c_1=-\dst\frac{37}{15}$,
so $y'=\dst\frac{x^5}{5}-\frac{37}{15}$. Therefore, $y=\dst\frac{x^6}{30}
-\frac{37}{15}(x-2)+c_2$; $y(2)=-1\Rightarrow\dst\frac{64}{30}+c_2=-1
\Rightarrow c_2=-\dst\frac{47}{15}$, so
$y=-\dst\frac{47}{15}-\frac{37}{5}(x-2)+\frac{x^6}{30}$.

\item (A) $\int xe^{2x}\,dx=\dfrac{xe^{2x}}{2}-\dfrac{1}{2}
\int e^{2x}\,dx=\dfrac{xe^{2x}}{2}-\dfrac{e^{2x}}{4}$. Therefore,
$y'=\dfrac{xe^{2x}}{2}-\dfrac{e^{2x}}{4}+c_1$; $y'(0)=1\Rightarrow
-\dst\frac{1}{4}+c_1=\dfrac{5}{4}\Rightarrow c_1=\dst\frac{5}{4}$,
so $y'=\dfrac{xe^{2x}}{2}-\dfrac{e^{2x}}{4}+\dfrac{5}{4}$;
Using (A) again,
$y=\dfrac{xe^{2x}}{4}-\dfrac{e^{2x}}{8}-\dfrac{e^{2x}}{8}+\dfrac{5}{4}x+c_2
=\dfrac{xe^{2x}}{4}-\dfrac{e^{2x}}{4}+\dfrac{5}{4}x+c_2$; $y(0)=7
\Rightarrow-\dst\frac{1}{4}+c_2=7\Rightarrow c_2=\dst\frac{29}{4}$, so
$y=\dfrac{xe^{2x}}{4}-\dfrac{e^{2x}}{4}+\frac{5}{4}x+\frac{29}{4}$.

\item (A) $\int x\sin x\,dx=-x\cos x+\int \cos x\,dx=-x\cos x+\sin
x$ and (B) $\int x\cos x\,dx=x\sin x-\int\sin x\,dx=x\sin x+\cos x$.
If $y''=-x\sin x$, then (A) implies that $y'=x\cos x-\sin x+c_1$;
$y'(0)=-3\Rightarrow c=-3$, so $y'=x\cos x-\sin x-3$. Now (B) implies
that $y=x\sin x+\cos x+\cos x-3x+c_2=x\sin x+2\cos x-3x+c_2$;
$y(0)=1\Rightarrow 2+c_2=1\Rightarrow c_2=-1$, so $y=x \sin x+2 \cos
x-3x-1$.


\item If $y'''=x^2e^x$, then $y''=\int x^2e^x\,dx=x^2e^x-2\int
xe^x\,dx=x^2e^x-2xe^x+2e^x+c_1$; $y''(0)=3\Rightarrow
2+c_1=3\Rightarrow c_1=1$, so (A) $y''=(x^2-2x+2)e^x+1$. Since $\int
(x^2-2x+2)e^x\,dx=(x^2-2x+2)e^x-\int (2x-2)e^x\,dx
=(x^2-2x+2)e^x-(2x-2)e^x+2e^x=(x^2-4x+6)e^x$, (A) implies that
$y'=(x^2-4x+6)e^x+x+c_2$; $y'(0)=-2\Rightarrow 6+c_2=-2\Rightarrow
c_2=-8$, so (B) $y'=(x^2-4x+6)e^x+x-8$; Since $\int
(x^2-4x+6)e^x\,dx=(x^2-4x+6)e^x-\int (2x-4)e^x\,dx
=(x^2-4x+6)e^x-(2x-4)e^x+2e^x=(x^2-6x+12)e^x$, (B) implies that
$y=(x^2-6x+12)e^x+\dst\frac{x^2}{2}-8x+c_3$; $y(0)=1\Rightarrow
12+c_3=1\Rightarrow c_3=-11$, so
$y=(x^2-6x+12)e^x+\dst\frac{x^2}{2}-8x-11$.

\item If $y'''=2+\sin2x$, then $y''=\dst2x-\frac{\cos 2x}{ 2}+c_1$;
$y''(0)=3\Rightarrow -\dst\frac{1}{ 2}+c_1=3\Rightarrow c_1=\dst\frac{7}{
2}$, so $y''=\dst2x-\frac{\cos 2x}{ 2}+\frac{7}{ 2}$. Then
$y'=x^2-\dst\frac{\sin 2x}{ 4}+\frac{7}{ 2}x+c_2$; $y'(0)=-6\Rightarrow
c_2=-6$, so $y'=x^2-\dst\frac{\sin 2x}{ 4}+\frac{7}{ 2}x-6$. Then
$y=\dst\frac{x^3}{ 3}+\frac{\cos 2x}{ 8}+\frac{7}{ 4}x^2-6x+c_3$;
$y(0)=1\Rightarrow \dst\frac{1}{ 8}+c_3=1\Rightarrow c_3=\dst\frac{7}{
8}$, so $\dst y=\frac{x^3}{3}+\frac{\cos 2x}{8}+
\frac{7}{4}x^2-6x+\frac{7}{8}$.

\item If $y'''=2x+1$, then $y''=x^2+x+c_1$; $y''(2)=7\Rightarrow
6+c_1=7\Rightarrow c_1=1$; so $y''=x^2+x+1$. Then $y'=\dst\frac{x^3}{
3}+\frac{x^2}{ 2}+(x-2)+c_2$; $y'(2)=-4\Rightarrow \dst\frac{14}{
3}+c_2=-4\Rightarrow c_2=-\dst\frac{26}{ 3}$, so $y'=\dst\frac{x^3}{
3}+\frac{x^2}{ 2}+(x-2)-\frac{26}{ 3}$. Then $y=\dst\frac{x^4}{
12}+\frac{x^3}{ 6}+\frac{1}{ 2}(x-2)^2-\frac{26}{ 3}(x-2)+c_3$;
$y(2)=1\Rightarrow \dst\frac{8}{ 3}+c_3=1\Rightarrow c_3=-\dst\frac{5}{
3}$, so $\dst y=\frac{x^4}{ 12}+\frac{x^3}{6}+\frac{1}{2} (x-2)^2-\frac{26}{3}
(x-2)-\frac{5}{3}$.
\end{alist}

\item
\begin{alist}
\item If $y=x^2(1+\ln x)$, then $y(e)=e^2(1+\ln e)=2e^2$;
$y'=2x(1+\ln x)+x=3x+2x\ln x$, so $y'(e)=3e+2e\ln e=5e$; (A)
$y''=3+2+2\ln x=5+2\ln x$. Now, $3xy'-4y=3x(3x+2x\ln x)-4x^2(1+\ln
x)=5x^2+2x^2\ln x=x^2y''$, from (A).

\item If $y=\dst\frac{x^2}{3}+x-1$, then $y(1)=\dst\frac{1}{
3}+1-1=\frac{1}{ 3}$; $y'=\dst\frac{2}{ 3}x+1$, so $y'(1)=\dst\frac{2}{
3}+1=\frac{5}{ 3}$; (A) $y''=\dst\frac{2}{ 3}$. Now
$x^2-xy'+y+1=x^2-\dst x\left(\frac{2}{ 3}x+1\right)+\frac{x^2}{
3}+x-1+1=\frac{2}{ 3}x^2=x^2y''$, from (A).

\item If $y=(1+x^2)^{-1/2}$, then $y(0)=(1+0^2)^{-1/2}=1$;
$y'=-x(1+x^2)^{-3/2}$, so $y'(0)=0$; (A) $y''=(2x^2-1)(1+x^2)^{-5/2}$.
Now,
$(x^2-1)y-x(x^2+1)y'=(x^2-1)(1+x^2)^{-1/2}-x(x^2+1)(-x)(1+x^2)^{-3/2}
=(2x^2-1)(1+x^2)^{-1/2}=y''(1+x^2)^2$ from (A), so
$y''=\dfrac{(x^2-1)y-x(x^2+1)y'}{ (x^2+1)^2}$.

\item If $y=\dst\frac{x^2}{ 1-x}$, then $y(1/2)=\dfrac{1/4}{
1-1/2}=\frac{1}{ 2}$; $y'=-\dfrac{x(x-2)}{(1-x)^2}$, so
$y'(1/2)=\dfrac{(-1/2)(-3/2)}{(1-1/2)^2}=3$; (A) $y''=\dst\frac{2}{
(1-x)^3}$. Now, (B) $x+y=x+\dst\frac{x^2}{1-x}=\frac{x}{1-x}$ and (C)
$xy'-y=\dst-\frac{x^2(x-2)}{(1-x)^2}-\frac{x^2}{1-x}=\frac{x^2}{(1-x)^2}$.
From (B) and (C), $(x+y)(xy'-y)=\dst\frac{x^3}{(1-x)^3}=\frac{x^3}{2}y''$,
so $y''=\dfrac{2(x+y)(xy'-y)}{ x^3}$.
\end{alist}


\item
\begin{alist}
\item
$y=(x-c)^a$ is defined and  $x-c=y^{1/a}$ on $(c,\infty)$;
moreover, $y'=a(x-c)^{a-1}=a\left(y^{1/a}\right)^{a-1}=ay^{(a-1)/a}$.

\item  if $a>1$ or $a<0$, then $y\equiv0$  is a solution of
(B) on $(-\infty,\infty)$.
\end{alist}

\item
\begin{alist}
\item Since $y'=c$ we must show that the right side of
(B)  reduces to $c$ for all values of $x$
in some interval. If $y=c^2+cx+2c+1$,
\begin{align*}
x^2+4x+4y&=x^2+4x+4c^2+4cx+8c+4\\
&=x^2+4(1+c)x+4(c^2+2c+1)\\
&=x^2+4(1+c)+2(c+1)^2=(x+2c+2)^2.
\end{align*}

Therefore, $\sqrt{x^2+4x+4y}=x+2c+2$ and the right side of
(B) reduces to $c$ if
$x>-2c-2$.

\item If $y_1=-\dfrac{x(x+4)}{4}$, then $y_1'=-\dst\frac{x+2}{2}$
and $x^2+4x+4y=0$ for all $x$. Therefore, $y_1$ satisfies
(A) on $(-\infty,\infty)$.
\end{alist}
\end{sectionSolutions}

\chapter{First Order Equations}

\section{Linear First Order Equations}

\begin{sectionSolutions}
\item
$\dst\frac{y'}{ y}=-3x^2$;\quad $|\ln|y|=-x^3+k$;\quad $y=ce^{-x^3}$.
$y=ce^{-(\ln x)^2/2}$.

\item
$\dst\frac{y'}{ y}=-\dst\frac{3}{ x}$;\quad $\ln|y|=-3\ln|x|+k=-\ln|x|^3+k$;\quad
$y=\dst\frac{c}{ x^3}$.

\item
$\dst\frac{y'}{ y}=-\dst\frac{1+x}{ x}=-\dst\frac{1}{ x}-1$;\quad
$|\ln|y|=-\ln|x|-x+k$;\quad
$y=\dfrac{ce^{-x}}{ x}$;\quad
$y(1)=1\Rightarrow
 c=e$;\quad
$y=\dfrac{e^{-(x-1)}}{ x}$.


\item
$\dst\frac{y'}{ y}=-\dst\frac{1}{ x}-\cot x$;
\quad $|\ln|y|=-\ln|x|-\ln|\sin x|+k=-\ln|x\sin x|+k$;\quad
$y=\dst\frac{c}{ x\sin x}$;\quad $y(\pi/2)=2\Rightarrow
c=\pi$;\quad $y=\dst\frac{\pi}{ x\sin x}$.


\item
$\dst\frac{y'}{ y}=-\dst\frac{k}{ x}$;\quad
$|\ln|y|=-k\ln|x|+k_1=\ln|x^{-k}|+k_1$;\quad
$y=c|x|^{-k}$;\quad
$y(1)=3\Rightarrow c=3$;\quad $y=3x^{-k}$.


\item
 $\dst\frac{y_1'}{ y_1}=-3$;\quad $\ln|y_1|=-3x$;\quad
$y_1=e^{-3x}$;\quad $y=ue^{-3x}$;\quad
$u'e^{-3x}=1$;\quad $u'=e^{3x}$;\quad $u=\dfrac{e^{3x}}{3}+c$;\quad
$y=\dst\frac{1}{3}+ce^{-3x}$.



\item
$\dst\frac{y_1'}{ y_1}=-2x$;\quad
$\ln|y_1|=-x^2$;\quad
$y_1=e^{-x^2}$;\quad
$y=ue^{-x^2}$;\quad
$u'e^{-x^2}=xe^{-x^2}$;\quad
$u'=x$;\quad
$u=\dst\frac{x^2}{2}+c$;\quad
$y=e^{-x^2}\dst\left(\frac{x^2}{2}+c\right)$.

\item
$\dst\frac{y_1'}{ y_1}=-\frac{1}{ x}$;\quad
$\ln|y_1|=-\ln|x|$;\quad $y_1=\dst\frac{1}{ x}$;\quad $y=\dst\frac{u}{
x}$;\quad
$\dst\frac{u'}{ x}=\dst\frac{7}{ x^2}+3$;\quad $u'=\dst\frac{7}{
x}+3x$;\quad
$u=7\ln|x|+\dst\frac{3x^2}{2}+c$;\quad $y=\dfrac{7\ln|x|}{
x}+\dst\frac{3x}{2}+\dst\frac{c}{ x}$.

\item
$\dst\frac{y_1'}{ y_1}=-\frac{1}{ x}-2x$; \quad
$\ln|y_1|=-\ln|x|-x^2$;\quad $y_1=\dfrac{e^{-x^2}}{ x}$;\quad
$y=\dfrac{ue^{-x^2}}{ x}$;\quad
$\dfrac{u'e^{-x^2}}{ x}=x^2e^{-x^2}$;\quad
 $u'=x^3$;\quad
$u=\dst\frac{x^4}{4}+c$;\quad
$y=e^{-x^2}\dst{\left(\frac{x^3}{4}+\frac{c}{ x}\right)}$.

\item
 $\dst\frac{y_1'}{ y_1}=-\tan x$;\quad
$\ln|y_1|=\ln|\cos x|$;\quad
$y_1=\cos x$;\quad $y=u\cos x$;\quad $u'\cos x=\cos x$;\quad
$u'=1$;\quad $u=x+c$;\quad $y=(x+c)\cos x$.


\item
$\dst\frac{y_1'}{ y_1}=\dst\frac{4x-3}{(x-2)(x-1)}=\dst\frac{5}{
x-2}-\dst\frac{1}{ x-1}$;\quad
$\ln|y_1|=5\ln|x-2|-\ln|x-1|=\dst{\ln\left|\frac{(x-2)^5}{
x-1}\right|}$;\quad
$y_1=\dfrac{(x-2)^5}{ x-1}$;\quad
 $y=\dfrac{u(x-2)^5}{ x-1}$;\quad
$\dfrac{u'(x-2)^5}{ x-1}=\dfrac{(x-2)^2}{ x-1}$;\quad
$u'=\dst\frac{1}{ (x-2)^3}$;\quad
$u=-\dst\frac{1}{ 2}\dst\frac{1}{(x-2)^2} +c$;\quad $y=-\dst\frac{1}{
2}\frac{(x-2)^3}{ (x-1)}+c\dfrac{(x-2)^5}{ (x-1)}$.


\item
$\dst\frac{y_1'}{ y_1}=-\dst\frac{3}{ x}$;\quad
$\ln|y_1|=-3\ln|x|=\ln|x|^{-3}$;\quad
$y_1=\dst\frac{1}{ x^3}$;\quad
$y=\dst\frac{u}{ x^3}$;\quad
$\dst\frac{u'}{ x^3}=\dst\frac{e^x}{ x^2}$;\quad
$u'=xe^x$;\quad
$u=xe^x-e^x+c$;\quad
$y=\dst\frac{e^x}{ x^2}-\dst\frac{e^x}{ x^3}+\dst\frac{c}{ x^3}$.



\item
$\dst\frac{y_1'}{ y_1}=-\dst\frac{4x}{1+x^2}$;\quad
$\ln|y_1|=-2\ln(1+x^2)=\ln(1+x^2)^{-2}$;\quad
$y_1=\dst\frac{1}{(1+x^2)^2}$;\quad
$y=\dst\frac{u}{(1+x^2)^2}$;\quad
$\dst\frac{u'}{(1+x^2)^2}=\frac{2}{(1+x^2)^2}$;\quad
$u'=2$;\quad
$u=2x+c$;\quad
$y=\dst\frac{2x+c}{(1+x^2)^2}$;\quad
$y(0)=1\Rightarrow c=1$;\quad
$y=\dst\frac{2x+1}{(1+x^2)^2}$.



\item
$\dst\frac{y_1'}{ y_1}=-\cot x$;\quad
$\ln|y_1|=-\ln|\sin x|$;\quad
$y_1=\dst\frac{1}{\sin x}$;\quad
$y=\dst\frac{u}{\sin x}$;\quad
$\dst\frac{u'}{\sin x}=\cos x$;\quad
$u'=\sin x\cos x$;\quad
$u=\dst\frac{\sin^2x}{2}+c$;\quad
$y=\dst\frac{\sin x}{2}+c\csc x$;\quad
$y(\pi/2)=1\Rightarrow c=\frac{1}{2}$;\quad
$y=\dst\frac{1}{2}(\sin x+\csc x)$.

\item
$\dst\frac{y_1'}{ y_1}=-\dst\frac{3}{ x-1}$;\quad
$\ln|y_1|=-3\ln|x-1|
=\ln|x-1|^{-3}$;\quad
$y_1=\dst\frac{1}{(x-1)^3}$;\quad
$y=\dst\frac{u}{(x-1)^3}$;\quad
$\dst\frac{u'}{(x-1)^3}=
\dst\frac{1}{(x-1)^4}+\dfrac{\sin
x}{(x-1)^3}$;\quad
$u'=\dst\frac{1}{ x-1}+\sin x$;\quad
$u=\ln|x-1|-\cos x+c$;\quad
$y=\dfrac{\ln|x-1|-\cos x+c}{(x-1)^3}$;\quad
$y(0)=1\Rightarrow c=0$;\quad
$y=\dfrac{\ln|x-1|-\cos x}{(x-1)^3}$.


\item
$\dst\frac{y_1'}{ y_1}=-\dst\frac{2}{ x}$;\quad
$\ln|y_1|=2\ln|x|=\ln(x^2)$;\quad
$y_1=x^2$;\quad
$y=ux^2$;\quad
$u'x^2=-x$;\quad
$u'=-\dst\frac{1}{ x}$;\quad
$u=-\ln|x|+c$;\quad
$y=x^2(c-\ln|x|)$;\quad
$y(1)=1\Rightarrow c=1$;\quad
$y=x^2(1-\ln x)$.


\item
$\dst\frac{y_1'}{ y_1}=-\dst\frac{3}{ x-1}$;\quad
$\ln|y_1|=-3\ln|x-1|=\ln|x-1|^{-3}$;\quad
$y_1=\dst\frac{1}{(x-1)^3}$;\quad
$y=\dst\frac{u}{(x-1)^3}$;\quad
$\dst\frac{u'}{(x-1)^3}=\dfrac{1+(x-1)\sec^2x}{(x-1)^4}$;\quad
$u'=\dst\frac{1}{ x-1}+\sec^2x$;\quad
$u=\ln|x-1|+\tan x+c$;\quad
$y=\dfrac{\ln|x-1|+\tan x+c}{(x-1)^3}$;\quad
$y(0)=-1\Rightarrow c=1$;\quad
$y=\dfrac{\ln|x-1|+\tan
x+1}{(x-1)^3}$.



\item
$\dst\frac{y_1'}{ y_1}=\dst\frac{2x}{ x^2-1}$;\quad
$\ln|y_1|=\ln|x^2-1|$;\quad
$y_1=x^2-1$;\quad
$y=u(x^2-1)$;\quad
$u'(x^2-1)=x$;\quad
$u'=\dst\frac{x}{ x^2-1}$;\quad
$u=\dst\frac{1}{2}\ln|x^2-1|+c$;\quad
$y=(x^2-1)\left(\dst\frac{1}{2}\ln|x^2-1|+c\right)$;\quad
$y(0)=4\Rightarrow c=-4$;\quad
$y=(x^2-1)\left(\dst\frac{1}{2}\ln|x^2-1|-4\right)$.


\item
$\dst\frac{y_1'}{ y_1}=-2x$;\quad
$\ln|y_1|=-x^2$;\quad
$y_1=e^{-x^2}$;\quad
$y=ue^{-x^2}$;\quad
$u'e^{-x^2}=x^2$;\quad
$u'=x^2e^{x^2}$;\quad
$u=c+\dst{\int_0^xt^2e^{t^2}\,dt}$;\quad
$y=e^{-x^2}\left(c+\dst{\int_0^xt^2e^{t^2}\,dt}\right)$;\quad
$y(0)=3\Rightarrow c=3$;\quad
$y=e^{-x^2}\left(3+\dst{\int_0^xt^2e^{t^2}\,dt}\right)$.


\item
$\dst\frac{y_1'}{ y_1}=-1$;\quad
$\ln|y_1|=-x$;\quad
$y_1=e^{-x}$;\quad
$y=ue^{-x}$;\quad
$u'e^{-x}=\dfrac{e^{-x}\tan x}{ x}$;\quad
$u'=\dst\frac{\tan x}{ x}$;\quad
$u=c+\dst{\int_1^x\dst\frac{\tan t}{ t}\,dt}$;\quad
$y=e^{-x}\left(c+\dst{\int_1^x\dst\frac{\tan t}{ t}\,dt}\right)$;\quad
$y(1)=0\Rightarrow c=0$;\quad
$y=e^{-x}\dst{\int_1^x\dst\frac{\tan t}{ t}\,dt}$.


\item
$\dst\frac{y_1'}{ y_1}=-1-\frac{1}{ x}$;\quad
$\ln|y_1|=-x-\ln|x|$;\quad
$y_1=\dfrac{e^{-x}}{ x}$;\quad
$y=\dfrac{ue^{-x}}{ x}$;\quad
$\dfrac{u'e^{-x}}{ x}=\dfrac{e^{x^2}}{ x}$;\quad
$u'=e^xe^{x^2}$;\quad
$u=c+\dst{\int_1^xe^te^{t^2}\,dt}$;\quad
$y=\dst{\frac{e^{-x}}{ x}\left(c+\int_1^xe^te^{t^2}\,dt\right)}$;\quad
$y(1)=2\Rightarrow c=2e$;\quad
$y=\dst{\frac{1}{
x}\left(2e^{-(x-1)}+e^{-x}\int_1^xe^te^{t^2}\,dt\right)}$.


\item
\begin{alist}[start=2]
\item Eqn. (A) is equivalent to
\[
y'-\frac{2}{ x}=-\frac{1}{ x}
\tag{B}
\]
on $(-\infty,0)$ and $(0,\infty)$. Here
$\dst\frac{y_1'}{ y_1}=\dst\frac{2}{ x}$;\quad $\ln|y_1|=2\ln|x|$;\quad
$y_1=x^2$;\quad
$y=ux^2$;\quad $u'x^2=-\dst\frac{1}{ x}$;\quad $u'=-\dst\frac{1}{
x^3}$;\quad
$u=\dst\frac{1}{2x^2}+c$, so $y=\dst\frac{1}{2}+cx^2$ is the general
solution of (A) on $(-\infty,0)$ and $(0,\infty)$.

\item
From the proof of \part{b}, any solution of (A) must be of the form
\[
y=\begin{cases}\dst{\frac{1}{2}+c_1x^2}, &x
\ge 0,\\
\dst{\frac{1}{2}+c_2x^2}, &x < 0,\end{cases}
\tag{C}
\]
for $x\ne0$, and any function of the form (C)
satisfies (A) for $x\ne0$. To complete the proof we
must show that any function of the form (C) is
differentiable and satisfies (A) at $x=0$. By
definition,
\[
y'(0)=\lim_{x\to0}\frac{y(x)-y(0)}{ x-0}=\lim_{x\to0}\frac{y(x)-1/2}{
x}
\]
if the limit exists. But
\[
\frac{y(x)-1/2}{ x}=\begin{cases} c_1x,&x>0\\
c_2x,&x<0,\end{cases}
\]
so $y'(0)=0$. Since $0y'(0)-2y(0)=0\cdot0-2(1/2)=-1$, any function of
the form (C) satisfies (A) at
$x=0$.

\item From \part{b} any solution $y$ of (A) on
$(-\infty,\infty)$ is of the form (C), so
$y(0)=1/2$.

\item If $x_0>0$, then every function of the form
(C) with $c_1=\dfrac{y_0-1/2}{ x_0^2}$ and $c_2$
arbitrary is a solution of the initial value problem on
$(-\infty,\infty)$. Since these functions are all identical on
$(0,\infty)$, this does not contradict Theorem~2.1.1, which
implies that (B) (so
(A)) has exactly one solution on $(0,\infty)$ such
that $y(x_0)=y_0$. A similar argument applies if $x_0<0$.
\end{alist}

\item
\begin{alist}
\item Let $y=c_1y_1+c_2y_2$. Then
\begin{align*}
y'+p(x)y&=(c_1y_1+c_2y_2)'+p(x)(c_1y_1+c_2y_2)\\[2\jot]
&=c_1y_1'+c_2y_2'+c_1p(x)y_1+c_2p(x)y_2\\[2\jot]
&=c_1(y_1'+p(x)y_1)+c_2(y_2+p(x)y_2)=c_1f_1(x)+c_2f_2(x).
\end{align*}

\item Let $f_1=f_2=f$ and $c_1=-c_2=1$.

\item Let $f_1=f$, $f_2=0$, and $c_1=c_2=1$.
\end{alist}


\item
\begin{alist}
\item
If $z=\tan y$, then $z'=(\sec^2y)y'$, so $z'-3z=-1$;\;
$z_1=e^{3x}$;\;
$z=ue^{3x}$;\;
$u'e^{3x}=-1$;\;
$u'=-e^{-3x}$;\;
$u=\dfrac{e^{-3x}}{3}+c$;\;
$z=\dst\frac{1}{3}+ce^{3x}=\tan y$;\;
$y=\tan^{-1}\dst{\left(\frac{1}{3}+ce^{3x}\right)}$.

\item
If $z=e^{y^2}$, then $z'=2yy' e^{y^2}$, so $z'+\dst\frac{2}{
x}z=\dst\frac{1}{ x^2}$;\;
$z_1=\dst\frac{1}{ x^2}$;\;
$z=\dst\frac{u}{ x^2}$;\;
$\dst\frac{u'}{ x^2}=\dst\frac{1}{ x^2}$;\;
$u'=1$;\;
$u=x+c$;\;
$z=\dst\frac{1}{ x}+\dst\frac{c}{ x^2}=e^{y^2}$;\;
$y=\pm\dst{\left[\ln\left(\dst\frac{1}{ x}+\dst\frac{c}{
x^2}\right)\right]^{1/2}}$.

\item Rewrite the equation as $\dst\frac{y'}{ y}+\dst\frac{2}{ x}\ln y=4x$.
If $z=\ln y$, then $z'=\dst\frac{y'}{ y}$, so $z'+\dst\frac{2}{ x}z=4x$;
\
$z_1=\dst\frac{1}{ x^2}$;\;
$z=\dst\frac{u}{ x^2}$;\;
$\dst\frac{u'}{ x^2}=4x$;\;
$u'=4x^3$;\;
$u=x^4+c$;\;
$z=x^2+\dst\frac{c}{ x^2}=\ln y$;\;
$y=\exp\dst\left(x^2+\frac{c}{ x^2}\right)$.

\item
If $z=-\dst\frac{1}{ 1+y}$, then $z'=\dst\frac{y'}{(1+y)^2}$, so
$z'+\dst\frac{1}{ x}z=-\dst\frac{3}{ x^2}$;\;
$z_1=\dst\frac{1}{ x}$;\;
$z=\dst\frac{u}{ x}$;\;
$\dst\frac{u'}{ x}=-\dst\frac{3}{ x^2}$;\;
$u'=-\dst\frac{3}{ x}$;\;
$u=-3\ln|x|-c$;\;
$z=-\dfrac{3\ln|x|+c}{ x}=-\dst\frac{1}{1+y}$;\;
$y=-1+\dst\frac{x}{3\ln|x|+c}$.
\end{alist}
\end{sectionSolutions}

\section{Separable Equations}

\begin{sectionSolutions}
\item
By inspection, $y\equiv k\pi$ ($k=$integer) is a constant solution.
Separate variables to find others: $\dst{\left(\frac{\cos y}{\sin
y}\right)}y'=-\sin x$;\ $\ln (|\sin y|)=\cos x+c$.



\item
 $y\equiv0$ is a  constant solution.
Separate variables to find others: $\dst{\left(\frac{\ln y}{
y}\right)y'}=-x^2$;\ $\dfrac{(\ln y)^2}{2}=-\frac{x^3}{3}+c$.


\item
 $y\equiv1$ and $y\equiv-1$ are constant solutions.
For others, separate variables:
$(y^2-1)^{-3/2}yy'=\dst\frac{1}{ x^2}$;\;
$-(y^2-1)^{-1/2}=-\dst\frac{1}{ x}-c=-\dst{\left(\frac{1+cx}{ x}\right)}$;\;
$(y^2-1)^{1/2}=\dst\left(\frac{x}{1+cx}\right)$;\;
$(y^2-1)=\dst{\left(\frac{x}{1+cx}\right)^2}$;\;
$y^2=1+\dst{\left(\frac{x}{1+cx}\right)^2}$;\;
$y=\pm\dst{\left(1+\left(\frac{x}{1+cx}\right)^2\right)^{1/2}}$.


\item
By inspection, $y\equiv0$ is a constant solution. Separate variables
to find others: $\dst\frac{y'}{ y}=-\dst\frac{x}{1+x^2}$;\;
$\ln|y|=-\dst\frac{1}{2}\ln(1+x^2)+k$;\;
 $y=\dst\frac{c}{\sqrt{1+x^2}}$,
which includes the constant solution $y\equiv0$.


\item
$(y-1)^2y'=2x+3$;\ $\dfrac{(y-1)^3}{3}=x^2+3x+c$;\;
$(y-1)^3=3x^2+9x+c$;\;
$\dst{y=1+\big(3x^2+9x+c)^{1/3}}$.



\item
$\dst\frac{y'}{ y(y+1)}=-x$;\ $\dst{\left[\frac{1}{ y}-\frac{1}{
y+1}\right]y'}=-x$;\ $\dst{\ln\left|\frac{y}{
y+1}\right|}=-\dst\frac{x^2}{2}+k$;\ $\dst\frac{y}{
y+1}=ce^{-x^2/2}$;\ $y(2)=1\Rightarrow c=\dst\frac{e^2}{2}$;\;
$y=(y+1)ce^{-x^2/2}$;\ $y(1-ce^{-x^2/2})=ce^{-x^2/2}$;\;
$y=\dfrac{ce^{-x^2/2}}{1-ce^{-x^2}/2}$;\ setting
$c=\dst\frac{e^2}{2}$ yields
 $\dst{y=\frac{e^{-(x^2-4)/2}
}{2-e^{-(x^2-4)/2}}}$.


\item
$\dst\frac{y'}{(y+1)(y-1)(y-2)}=-\dst\frac{1}{ x+1}$;\;
$\dst{\left[\frac{1}{6}\frac{1}{ y+1}-\frac{1}{2}\frac{1}{
y-1}+\frac{1}{3}\frac{1}{ y-2}\right]}y'=-\frac{1}{ x+1}$;\;
$\dst{\left[\frac{1}{ y+1}-\frac{3}{ y-1}+\frac{2}{
y-2}\right]y'}=-\dst\frac{6}{ x+1}$;\;
$\ln|y+1|-3\ln|y-1|+2\ln|y-2|=-6\ln|x+1|+k$;\;
$\dst\frac{(y+1)(y-2)^2}{(y-1)^3}=\dst\frac{c}{(x+1)^6}$;\;
$y(1)=0\Rightarrow c=-256$;\;
$\dfrac{(y+1)(y-2)^2}{(y-1)^3}=-\dst\frac{256}{(x+1)^6}$.


\item
$\dst\frac{y'}{ y(1+y^2)}=2x$;\ $\dst{\left[\frac{1}{ y}-\frac{y}{
y^2+1}\right]y'}=2x$;\ $\dst{\ln\left(\frac{|y|}{
\sqrt{y^2+1}}\right)}=x^2+k$;\ $\dst\frac{y}{
\sqrt{y^2+1}}=ce^{x^2}$;\ $y(0)=1\Rightarrow c=\dst\frac{1}{\sqrt2}$;\;
$\dst\frac{y}{\sqrt{y^2+1}}=\dfrac{e^{x^2}}{\sqrt2}$;\;
$2y^2=(y^2+1)e^{x^2}$;\ $y^2(2-e^{x^2})=e^{2x^2}$;
$y=\dst\frac{1}{\sqrt{2e^{-2x^2}-1}}$.


\item
$\dst\frac{y'}{(y-1)(y-2)}=-2x$;\ $\dst{\left[\frac{1}{
y-2}-\frac{1}{ y-1}\right]}y'=-2x$;\ $\dst{\ln\left|\frac{y-2}{
y-1}\right|}=-x^2+k$;\ $\dst\frac{y-2}{ y-1}=ce^{-x^2}$;\;
$y(0)=3\Rightarrow c=\dst\frac{1}{2}$;\ $\dst\frac{y-2}{
y-1}=\dfrac{e^{-x^2}}{2}$;\ $y-2=\dfrac{e^{-x^2}}{2}(y-1)$;\;
$\dst{y\left(1-\frac{e^{-x^2}}{2}\right)}=2-\dfrac{e^{-x^2}}{2}$;\;
$y=\dfrac{4-e^{-x^2}}{2-e^{-x^2}}$. \
The interval of validity is $(-\infty,\infty)$.



\item
$\dst\frac{y'}{ y(y-2)}=-1$;\ $\dst{\frac{1}{2}\left[\frac{1}{
y-2}-\frac{1}{ y}\right]y'}=-1$;\;
 $\dst{\left[\frac{1}{
y-2}-\frac{1}{ y}\right]y'}=-2$;\ $\dst{\ln\left|\frac{y-2}{
y}\right|}=-2x+k$;\ $\dst\frac{y-2}{ y}=ce^{-2x}$;\;
$y(0)=1\Rightarrow c=-1$;\ $\dst\frac{y-2}{ y}=-e^{-2x}$;\;
$y-2=-ye^{-2x}$;\ $y(1+e^{-2x})=2$;\;
   $y=\dst\frac{2}{
1+e^{-2x}}$. The interval of validity is $(-\infty,\infty)$.


\item
$y\equiv2$ is a constant solution of the differential equation, and
it satisfies the initial condition. Therefore, $y\equiv2$ is a
solution of the initial value problem. The interval of validity is
$(-\infty,\infty)$.



\item
$\dst\frac{y'}{1+y^2}=\dst\frac{1}{1+x^2}$;\ $\tan^{-1}y=
\tan^{-1}x+k$;\ $y=\tan(\tan^{-1}x+k)$. Now use
the identity $\tan(A+B)=\dst\frac{\tan A+\tan B}{1-\tan A\tan B}$ with
$A=\tan^{-1}x$ and $B=\tan^{-1}c$ to rewrite $y$ as
$y=\dst\frac{x+c}{1-cx}$, where $c=\tan k$.

\item
$(\sin y)y'=\cos x$;\ $-\cos
y=\sin x+c$;\  $y(\pi)=\dst\frac{\pi}{2}\Rightarrow
c=0$, so (A)  $\cos y=-\sin
x$. To obtain $y$ explicitly we  note that $-\sin x=\cos(x+\pi/2)$, so
(A) can be rewritten as  $\cos y=\cos(x+\pi/2)$. This equation
holds if an only if one of the following conditions holds for some
integer $k$:
\[
\text{(B)}\qquad y=x+\frac{\pi}{2}+2k\pi;\qquad
\text{(C)}\qquad y=-x-\frac{\pi}{2}+2k\pi.
\]
Among these choices the only way to satisfy the initial condition
is to let $k=1$ in (C), so $y=-x+\dst\frac{3\pi}{2}.$



\item
Rewrite the equation as $P'=-a\alpha P(P-1/\alpha)$.
By inspection, $P\equiv0$ and $P\equiv1/\alpha$ are constant
solutions.
Separate variables to find others: $\dst\frac{P'}{
P(P-1/\alpha)}=-a\alpha$;\;
 $\dst{\left[\frac{1}{ P-1/\alpha}-\frac{1}{ P}\right]P'}=-a
$;\ $\dst{\ln\left|\frac{P-1/\alpha}{ P}\right|}=-at+k$;\;
(A) $\dfrac{P-1/\alpha}{
P}=ce^{-\alpha t}$; $P(1-ce^{-\alpha t})=1/\alpha$;\;
(B) $P=\dst\frac{1}{\alpha(1-ce^{-\alpha t})}$. From (A),
 $P(0)=P_0\Rightarrow c=\dfrac{P_0-1/\alpha}{ P_0}$.
Substituting this into (B) yields
$P=\dst\frac{P_0}{\alpha
P_0+(1-\alpha P_0)e^{-at}}$. From this
$\lim_{t\to\infty}P(t)=1/\alpha$.


\item
If $q=rS$ the equation for $I$ reduces to $I'=-rI^2$, so
$\dst\frac{I'}{ I^2}=-r$; $-\dst\frac{1}{ I}=-rt-\dst\frac{1}{ I_0}$;
so $I=\dst\frac{I_0}{1+rI_0t}$ and $\lim_{t\to\infty}I(t)=0$. If $q\ne
rS$, then rewrite the equation for $I$ as $I'=-rI(I-\alpha)$ with
$\alpha=S-\dst\frac{q}{ r}$. Separating variables yields
$\dst\frac{I'}{ I(I-\alpha)}=-r$;\;
 $\dst{\left[\frac{1}{ I-\alpha}-\frac{1}{ I}\right]I'}=-r\alpha
$;\ $\dst{\ln\left|\frac{I-\alpha}{ I}\right|}=-r\alpha t+k$;\;
(A) $\dst\frac{I-\alpha}{
I}=ce^{-r\alpha t}$; $I(1-ce^{-r\alpha t})=\alpha$;\;
(B) $I=\dst\frac{\alpha}{ 1-ce^{-r\alpha t}}$. From (A),
 $I(0)=I_0\Rightarrow c=\dst\frac{I_0-\alpha}{ I_0}$.
Substituting this into (B) yields
$I=\dst\frac{\alpha I_0}{ I_0+(\alpha-I_0)e^{-r\alpha t}}$.
If $q<rS$, then $\alpha>0$ and
$\lim_{t\to\infty}I(t)=\alpha=S-\dst\frac{q}{ r}$.
If $q>rS$, then $\alpha<0$ and $\lim_{t\to\infty}I(t)=0$.

\addtocounter{enumi}{2}
\item%2.2.34
The given equation is separable if $f=ap$, where $a$ is a constant.
In this case the equation is
\[
y'+p(x)y=ap(x).
\tag{A}
\]
Let $P$ be an antiderivative  of $p$; that is, $P'=p$.

{\sc Solution by Separation of Variables.} $y'=-p(x)(y-a)$;\;
$\dst\frac{y'}{ y-a}=-p(x)$;\ $\ln|y-a|=-P(x)+k$;\;
$y-a=ce^{-P(x)}$;\ $y=a+ce^{-P(x)}$.

{\sc Solution by Variation of Parameters.} $y_1=e^{-P(x)}$
is a solution of the complementary equation, so solutions of
(A) are of the form $y=ue^{-P(x)}$ where
$u'e^{-P(x)}=ap(x)$. Hence, $u'=ap(x)e^{P(x)}$;\;
$u=ae^{P(x)}+c$;\ $y=a+ce^{-P(x)}$.


\item Rewrite the given equation as
(A) $y'-\dst\frac{2}{ x}y=\dst\frac{x^5}{ y+x^2}$.
 $y_1=x^2$ is a solution of $y'-\dst\frac{2}{ x}y=0$.
Look for solutions of (A)  of the form $y=ux^2$.
Then $u'x^2=\dst\frac{x^5}{(u+1)x^2}=\dst\frac{x^3}{ u+1}$;\;
$u'=\dst\frac{x}{ u+1}$;\;
$(u+1)u'=x$;\;
$\dfrac{(1+u)^2}{2}=\dst\frac{x^2}{2}+\dst\frac{c}{2}$;\;
$u=-1\pm\sqrt{x^2+c}$;\ $y=x^2\left(-1\pm\sqrt{x^2+c}\right)$.


\item
$y_1=e^{2x}$ is a solution of $y'-2y=0$.
Look for solutions of the nonlinear equation of the form $y=ue^{2x}$.
Then $u'e^{2x}=\dfrac{xe^{2x}}{1-u}$;\ $u'=\dst\frac{x}{1-u}$;\;
$(1-u)u'=x$;\;
$-\dfrac{(1-u)^2}{2}=\dst\frac{1}{2}(x^2-c)$;\;
$u=1\pm\sqrt{c-x^2}$;\;
$y=e^{2x}\left(1\pm\sqrt{c-x^2}\right)$.

\end{sectionSolutions}

\section{Existence and Uniqueness of Solutions of Nonlinear Equations}

\begin{sectionSolutions}
\item
$f(x,y)=\dst\frac{e^x+y}{ x^2+y^2}$ and $f_y(x,y)=\dst\frac{1}{
x^2+y^2}-\dfrac{2y(e^x+y)}{(x^2+y^2)^2}$ are both continuous at all
$(x,y)\ne(0,0)$. Hence,  Theorem~2.3.1 implies that if
$(x_0,y_0)\ne(0,0)$, then  the initial
value problem has a a unique solution on some open interval containing
$x_0$.  Theorem~2.3.1 does not apply
if $(x_0,y_0)=(0,0)$.

\item
$f(x,y)=\dst\frac{x^2+y^2}{\ln xy}$ and $f_y(x,y)=\dst\frac{2y}{\ln
xy}-\dst\frac{x^2+y^2}{ x(\ln xy)^2}$ are both continuous at all $(x,y)$
such that $xy>0$ and $xy\ne1$. Hence,  Theorem~2.3.1 implies
that if
$x_0y_0>0$ and $x_0y_0\ne1$, then the
initial value problem has unique solution on an open interval
containing $x_0$.   Theorem~2.3.1 does not apply
if $x_0y_0\le0$ or $x_0y_0=1$.


\item
$f(x,y)=2xy$ and $f_y(x,y)=2x$ are both continuous at all $(x,y)$.
Hence,  Theorem~2.3.1 implies that if $(x_0,y_0)$
is arbitrary, then the initial value problem
has a  unique solution on some open interval containing  $x_0$.


\item
$f(x,y)=\dst\frac{2x+3y}{ x-4y}$ and $f_y(x,y)=\dst\frac{3}{
x-4y}+4\dst\frac{2x+3y}{(x-4y)^2}$ are both continuous at all $(x,y)$
such that $x\ne4y$. Hence,  Theorem~2.3.1 implies that if
$x_0\ne4y_0$, then the
initial value problem has a  unique solution on some open interval
containing  $x_0$.
  Theorem~2.3.1 does not apply if $x_0=4y_0$.

\item
$f(x,y)=x(y^2-1)^{2/3}$ is continuous at all $(x,y)$, but
$f_y(x,y)=\dst\frac{4}{3}xy(y^2-1)^{1/3}$ is continuous at $(x,y)$ if
and only if $y\ne\pm1$. Hence,  Theorem~2.3.1 implies that
if $y_0\ne\pm1$, then the
initial value problem has a  unique solution on some open interval
containing  $x_0$, while  if $y_0=\pm1$, then the initial value problem
 has at least one solution (possibly not unique on any
open interval containing  $x_0$).


\item
$f(x,y)=(x+y)^{1/2}$ and $f_y(x,y)=\dst\frac{1}{2(x+y)^{1/2}}$ are both
continuous at all $(x,y)$ such that $x+y>0$ Hence,
 Theorem~2.3.1 implies that if $x_0+y_0>0$, then the initial
value problem has a
 unique solution on some open interval containing  $x_0$.
 Theorem~2.3.1 does not apply if $x_0+y_0\le0$.

\item
To apply  Theorem~2.3.1, rewrite the given initial value problem
as (A) $y'=f(x,y),\ y(x_0)=y_0$, where $f(x,y)=-p(x)y+q(x)$ and
$f_y(x,y)=-p(x)$. If $p$ and $f$ are continuous on some open interval
$(a,b)$ containing $x_0$, then $f$ and $f_y$ are continuous on some
open
rectangle containing $(x_0,y_0)$, so  Theorem~2.3.1 implies
that (A) has a unique solution  on {\it some\/} open
interval
containing $x_0$. The conclusion of Theorem~2.1.2 is more
specific:
the solution of (A) exists and is unique on $(a,b)$. For example, in
the extreme case where $(a,b)=(-\infty,\infty)$,  Theorem~2.3.1
still implies only existence and uniqueness on {\it some\/} open
interval containing $x_0$, while  Theorem~2.1.2 implies that the
solution exists and is unique on $(-\infty,\infty)$.


\item
First find solutions of (A) $y'=y^{2/5}$. Obviously $y\equiv0$ is a
solution. If $y\not\equiv0$, then we can separate variables on any
open interval where $y$ has no zeros: $y^{-2/5}y'=1$;\;
$\dst\frac{5}{3}y^{3/5}=x+c$;\;
$y=\dst{\left(\frac{3}{5}(x+c)^{5/3}\right)}$. (Note that this solution
is also defined  at $x=-c$, even though  $y(-c)=0$. To satisfy the
initial condition, let $c=1$. Thus,
$y=\dst{\left(\frac{3}{5}(x+1)^{5/3}\right)}$ is a solution of the
initial value problem on $(-\infty,\infty)$; moreover,
since $f(x,y)=y^{2/5}$ and $f_y(x,y)=\dst\frac{2}{5}y^{-3/5}$
are both continuous  at all $(x,y)$ such that $y\ne0$,
this is the only
solution on $(-5/3,\infty)$, by an argument similar to that given in
Example~2.3.7, the function
\[
y =
\begin{cases}
0, &  -\infty< x \le -\frac{5}{ 3} \\
\left(\frac{3}{ 5}x + 1\right)^{5/3}, & -\frac{5}{3} < x < \infty
\end{cases}
\]
(To see that $y$ satisfies $y'=y^{2/5}$  at $x=-\dst\frac{5}{3}$
use an argument similar to that of
Discussion~2.3.15-2)
For every $\dst{a \ge \frac{5}{ 3}}$, the following function is also a
solution:
\[
y =
\begin{cases}
\left( \frac{3}{ 5} (x+a) \right)^{5/3}, & - \infty < x < -a, \\[3pt]
0, &  -a \le x \le -\frac{5}{ 3} \\[3pt]
\left(\frac{3}{ 5}x + 1\right)^{5/3}, & -\frac{5}{3} < x < \infty.
\end{cases}
\]


\item
Obviously, $y_1\equiv1$ is a solution.
From Discussion~2.3.18 (taking $c=0$
in the two families of solutions) yields
$y_2=1+|x|^3$ and $y_3=1-|x|^3$. Other solutions are
$y_4=1+x^3$, $y_5=1-x^3$,
\begin{gather*}
y_6=\dst{\begin{cases}1+x^3,&x\ge 0,\\
1,&x<0\end{cases}};\quad
y_7=\dst{\begin{cases}1-x^3,&x\ge 0,\\
1,&x<0\end{cases}};
\\
y_8=\dst{\begin{cases}1,&x\ge 0,\\
1+x^3,&x<0\end{cases}};\quad
y_9=\dst{\begin{cases}1,&x\ge 0,\\
1-x^3,&x<0\end{cases}}
\end{gather*}


It is straightforward to verify that all these functions satisfy
$y'=3x(y-1)^{1/3}$  for all $x\ne0$. Moreover,
 $y_i'(0)=\dst{\lim_{x\to0}\frac{y_i(x)-1}{ x}}=0$ for
$1\le i\le 9$, which implies that they also satisfy the equation at
$x=0$.



\item
Let $y$ be any solution of (A) $y'=3x(y-1)^{1/3},\ y(3)=-7$.
By continuity, there is some open interval $I$ containing  $x_0=3$
on which $y(x)<1$. From
Discussion~2.3.18,
 $y=1+(x^2+c)^{3/2}$ on $I$;\ $y(3)=-7\Rightarrow c=-5$;\;
(B) $y=1-(x^2-5)^{3/2}$.  It now follows that every solution of (A)
satisfies $y(x)<1$
and is given by (B) on $(\sqrt5,\infty)$; that is,
(B) is the unique solution of (A) on
$(\sqrt5,\infty)$.
 This solution can be extended uniquely to
$(0,\infty)$ as
\[
y=
\begin{cases}
1, & 0< x\le\sqrt5,\\
1-(x^2-5)^{3/2}, & \sqrt5<x<\infty
\end{cases}
\]
It can be extended to $(-\infty,\infty)$ in infinitely many ways.
Thus,
\[
\dst y= \begin{cases}
 1, &-\infty< x\le\sqrt5,\\
1-(x^2-5)^{3/2},& \sqrt5<x<\infty
\end{cases}
\]
is a solution of the initial value problem on $(-\infty,\infty)$.
Moreover, if
$\alpha\ge0$, then
\[
y= \begin{cases}
 1+(x^2-\alpha^2)^{3/2},& -\infty<x<-\alpha, \\
1, & -\alpha\le x\le\sqrt5,\\
1-(x^2-5)^{3/2},& \sqrt5<x<\infty,
\end{cases}
\]
and
\[
y= \begin{cases}
 1-(x^2-\alpha^2)^{3/2},& -\infty<x<-\alpha, \\
1, & -\alpha\le x\le\sqrt5,\\
1-(x^2-5)^{3/2},& \sqrt5<x<\infty,
\end{cases}
\]
are also solutions of the initial value problem  on
$(-\infty,\infty)$.

\end{sectionSolutions}

\section{Transformation of Nonlinear Equations into Separable Equations}

\begin{sectionSolutions}
\item
 Rewrite as $y'-\dst\frac{2}{7x}y=-\dst\frac{x}{7y^6}$. Then
$\dst\frac{y_1'}{ y_1}=\dst\frac{2}{7x}$;\;
$\ln|y_1|=\dst\frac{2}{7}\ln|x|=\ln|x|^{2/7}$;\;
$y_1=x^{2/7}$;\;
$y=ux^{2/7}$;\;
$u'x^{2/7}=-\dst\frac{1}{7u^6x^{5/7}}$;\;
$u^6u'=-\dst\frac{1}{7x}$;\;
$\dst\frac{u^7}{7}=-\dst\frac{1}{7}\ln|x|+\dst\frac{c}{7}$;\;
$u=(c-\ln|x|)^{1/7}$;\;
$y=x^{2/7}(c-\ln|x|)^{1/7}$.


\item
 Rewrite as $y'+\dst\frac{2x}{1+x^2}y=\dst\frac{1}{(1+x^2)^2y}$.
Then
$\dst\frac{y_1'}{ y_1}=-\dst\frac{2x}{1+x^2}$;\;
$\ln|y_1|=-\ln(1+x^2)$;\;
$y_1=\dst\frac{1}{1+x^2}$;\;
$y=\dst\frac{u}{1+x^2}$;\;
$\dst\frac{u'}{1+x^2}=\dst\frac{1}{ u(1+x^2)}$;\;
$u'u=1$;\;
$\dst\frac{u^2}{2}=x+\dst\frac{c}{2}$;\;
$u=\pm\sqrt{2x+c}$;\;
$y=\pm\dst{\frac{\sqrt{2x+c}}{1+x^2}}$.


\item
$\dst\frac{y_1'}{ y_1}=\dst{\frac{1}{3}\left(\frac{1}{ x}+1\right)}$;\;
$\ln|y_1|=\dst\frac{1}{3}(\ln|x|+x)$;\;
$y_1=x^{1/3}e^{x/3}$;\;
$y=ux^{1/3}e^{x/3}$;\;
$u'x^{1/3}e^{x/3}=x^{4/3}e^{4x/3}u^4$;\;
$\dst\frac{u'}{ u^4}=xe^x$;\;
$-\dst\frac{1}{3u^3}=(x-1)e^x-\dst\frac{c}{3}$;\;
$u=\dst\frac{1}{[3(1-x)e^x+c]^{1/3}}$;\;
$y=\dst{\left[\frac{x}{3(1-x)+ce^{-x}}\right]^{1/3}}$.


\item
$\dst\frac{y_1'}{ y_1}=x$;\;
$\ln|y_1|=\dst\frac{x^2}{2}$;\;
$y_1=e^{x^2/2}$;\;
$y=ue^{x^2/2}$;\;
$u'e^{x^2/2}=xu^{3/2}e^{3x^2/4}$;\;
$\dst\frac{u'}{ u^{3/2}}=xe^{x^2/4}$;\;
(A) $-\dst\frac{2}{ u^{1/2}}=2e^{x^2/4}+2c$;\;
$u^{1/2}=-\dst\frac{1}{ c+e^{x^2/4}}$;\;
$u^=\dst\frac{1}{(c+e^{x^2/4})^2}$;\;
$y=\dst\frac{1}{(1+ce^{-x^2/4})^2}$. Because of (A) we must choose $c$
so that $y(1)=4$ and $1+ce^{-1/4}<0$. This implies that
$c=-3e^{1/4}$;\;
$y=\dst{\left[1-\frac{3}{2}e^{-(x^2-1)/4}\right]^{-2}}$.


\item
$\dst\frac{y_1'}{ y_1}=2$;\;
$\ln|y_1|=2x$;\;
$y_1=e^{2x}$;\;
$y=ue^{2x}$;\;
$u'e^{2x}=2u^{1/2}e^x$;\;
$u^{-1/2}u'=2e^{-x}$;\;
$2u^{1/2}=-2e^{-x}+2c$;\;
$u^{1/2}=c-e^{-x}>0$;\;
$y(0)=1\Rightarrow u(0)=1\Rightarrow
c=2$;\;
$u=(2-e^{-x})^2$;\;
$y=(2e^x-1)^2$.


\item
Rewrite as $y'+\dst\frac{2}{ x}y=\dst\frac{y^3}{ x^2}$. Then
$\dst\frac{y_1'}{ y_1}=-\dst\frac{2}{ x}$;\;
$\ln|y_1|=-2\ln|x|=\ln x^{-2}$;\;
$y_1=\dst\frac{1}{ x^2}$;\;
$y=\dst\frac{u}{ x^2}$;\;
$\dst\frac{u'}{ x^2}=\dst\frac{u^3}{ x^8}$;\;
$\dst\frac{u'}{ u^3}=\dst\frac{1}{ x^6}$;\;
$-\dst\frac{1}{2u^2}=-\dst\frac{1}{5x^5}+c$;\;
$y(1)=\dst\frac{1}{\sqrt2}\Rightarrow u(1)=\dst\frac{1}{\sqrt2}\Rightarrow
c=-\dst\frac{4}{5}$;\;
$u=\dst{\left[\frac{5x^5}{2(1+4x^5)}\right]^{1/2}}$;\;
$y=\dst{\left[\frac{5x}{2(1+4x^5)}\right]^{1/2}}$.


\item
$P=ue^{at}$;\ $u'e^{at}=-a\alpha u^2e^{2at}$;\ $\dst\frac{u'}{
u^2}=-a\alpha e^{at}$;\ $-\dst\frac{1}{
u}=-a\int_0^t\alpha(\tau)e^{a\tau}\,d\tau-\dst\frac{1}{ P_0}$;\;
$P=\dfrac{P_0e^{at}}{1+aP_0\int_0^t\alpha(\tau)e^{a\tau}\,d\tau}$,
which can also be written as
$P=\dst\frac{P_0}{
e^{-at}+aP_0e^{-at}\int_0^t\alpha(\tau)e^{a\tau}\,d\tau}$.
Therefore, $\lim_{t\to\infty}P(t)=
\dst{\begin{cases}\infty&\text{ if }L=0,\\0&\text{ if
}L=\infty,\\ 1/aL&\text{ if } 0<L<\infty.\end{cases}}$

\item
$y=ux$;\ $u'x+u=u^2+2u$;\;
(A) $u'x=u(u+1)$. Since $u\equiv0$ and $u\equiv-1$ are constant
solutions of (A), $y\equiv0$ and $y=-x$ are
solutions of the given equation. The nonconstant solutions of (A)
satisfy
$=\dst\frac{u'}{ u(u+1)}=\frac{1}{ x}$;\;
$\dst{\left[\frac{1}{ u}-\frac{1}{ u+1}\right]u'}=\dst\frac{1}{ x}$;\;
$\dst{\ln\left|\frac{u}{ u+1}\right|}=\ln|x|+k$;\;
$\dst\frac{u}{ u+1}=cx$;\ $u=(u+1)cx$;\ $u(1-cx)=cx$;\;
$u=\dst\frac{cx}{1-cx}$;\;
$y=\dst\frac{cx^2}{1-cx}$.

\item
$y=ux$;\ $u'x+u=u+\sec$;\;
$u'x=\sec u$;\;
$(\cos u)u'=\dst\frac{1}{ x}$;\;
$\sin u=\ln|x|+c$;\;
$u=\sin^{-1}(\ln|x|+c)$;\;
$y=x\sin^{-1}(\ln|x|+c)$.



\item
Rewrite the given equation as $y'=\dst\frac{x^2+2y^2}{ xy}$;\;
$y=ux$;\ $u'x+u=\dst\frac{1}{ u}+2u$;\;
$u'x=\dst\frac{1+u^2}{ u}$;\;
$\dst\frac{uu'}{1+u^2}=\dst\frac{1}{ x}$;\;
$\dst\frac{1}{2}\ln(1+u^2)=\ln|x|+k$;\;
$\dst{\ln\left(1+\frac{y^2}{ x^2}\right)}=\ln x^2+2k$;\;
$\dst{1+\frac{y^2}{ x^2}}=cx^2$;\;
$x^2+y^2=cx^4$;\ $y=\pm x\sqrt{cx^2-1}$.



\item
$y=ux$;\ $u'x+u=u+u^2$;\;
$u'x=u^2$;\;
$\dst\frac{u'}{ u^2}=\dst\frac{1}{ x}$;\;
$-\dst\frac{1}{ u}=\ln|x|+c$;\;
$y(-1)=2\Rightarrow u(-1)=-2\Rightarrow c=\dst\frac{1}{2}$;\;
$u=-\dst\frac{2}{2\ln|x|+1}$;\;
$y=-\dst\frac{2x}{2\ln|x|+1}$.


\item
Rewrite the given equation as $y'=-\dst\frac{x^2+y^2}{ xy}$;\;
$y=ux$;\ $u'x+u=-\dst\frac{1}{ u}-u$;\;
$u'x=-\dst\frac{1+2u^2}{ u}$;\;
$-\dst\frac{uu'}{1+2u^2}=\dst\frac{1}{ x}$;\;
$-\dst\frac{1}{4}\ln(1+2u^2)=\ln|x|+k$;\;
$x^4(1+2u^2)=c$;\;
$y(1)=2\Rightarrow u(1)=2\Rightarrow c=9$;\;
$x^4(1+2u^2)=9$;\;
$u^2=\dst\frac{9-x^4}{2x^4}$;\;
$u=\dst\frac{1}{ x^2}\dst{\left(\frac{9-x^4}{2}\right)^{1/2}}$;\;
$y=\dst\frac{1}{ x}\dst{\left(\frac{9-x^4}{2}\right)^{1/2}}$.

\item
Rewrite the given equation as $y'=2+\dst\frac{y^2}{
x^2}+4\dst\frac{y}{ x}$;\;
$y=ux$;\ $u'x+u=2+u^2+4u$;\;
$u'x=u^2+3u+2=(u+1)(u+2)$;\;
$\dst\frac{u'}{(u+1)(u+2)}=\dst\frac{1}{ x}$;\;
$\dst{\left[\frac{1}{ u+1}-\frac{1}{ u+2}\right]u'}=\dst\frac{1}{ x}$;\;
$\dst{\ln\left|\frac{u+1}{ u+2}\right|}=\ln|x|+k$;\;
$\dst\frac{u+1}{ u+2}=cx$;\;
$y(1)=1\Rightarrow u(1)=1\Rightarrow c=\dst\frac{2}{3}$;\;
$\dst\frac{u+1}{ u+2}=\dst\frac{2}{3}x$;\;
$u+1=\dst\frac{2}{3}x(u+2)$;\;
$u\dst{\left(1-\frac{2}{3}x\right)}=-1+\dst\frac{4}{3}x$;\;
$u=-\dst\frac{4x-3}{2x-3}$;\;
$y=-\dfrac{x(4x-3)}{2x-3}$.

\item
$y=ux$;\ $u'x+u=\dst\frac{1+u}{1-u}$;\;
$u'x=\dst\frac{1+u^2}{1-u}$;\;
$\dfrac{(1-u)u'}{1+u^2}=\dst\frac{1}{ x}$;\;
$\tan^{-1}u-\dst\frac{1}{2}\ln(1+u^2)=\ln|x|+c$;\;
$\tan^{-1}\dst\frac{y}{
x}-\dst\frac{1}{2}\ln\dst{\left(1+\frac{y^2}{ x^2}\right)}=\ln|x|+c$;\;
$\tan^{-1}\dst\frac{y}{
x}-\dst\frac{1}{2}\ln(x^2+y^2)=c$.

\item
$y=ux$;
$u'x+u=\dst\frac{u^3+2u^2+u+1}{(u+1)^2}=\dfrac{u(u+1)^2+1}{(u+1)^2}
=u+\dst\frac{1}{(u+1)^2}$;
$u'x=\dst\frac{1}{(u+1)^2}$;\;
$(u+1)^2u'=\dst\frac{1}{ x}$;\;
$\dfrac{(u+1)^3}{3}=\ln|x|+c$;\;
$(u+1)^3=3(\ln|x|+c)$;\;
$\dst\left(\frac{y}{ x}+1\right)^3=3(\ln|x|+c)$;\;
$(y+x)^3=3x^3(\ln|x|+c)$.



\item
$y=ux$;\ $u'x+u=\dst\frac{u}{ u-2}$;\;
(A)  $u'x=\dfrac{u(u-3)}{2-u}$;\;
Since $u\equiv0$ and $u\equiv3$
are constant solutions of (A), $y\equiv0$ and
$y=3x$ are solutions of the given equation.
The nonconstant solutions of (A) satisfy
$\dfrac{(2-u)u'}{ u(u-3)}=\dst\frac{1}{ x}$;\;
$\dst{\left[\frac{1}{ u-3}+\frac{2}{
u}\right]u'}=-\dst\frac{3}{ x}$;\;
$\ln|u-3|+2\ln|u|=-3\ln|x|+k$;\;
$u^2(u-3)=\dst\frac{c}{ x^3}$;\;
$y^2(y-3x)=c$.


\item
$y=ux$;\;
$u'x+u=
\dst \frac{1+u+3u^3 }{ 1+3u^2}=u+\frac{1 }{ 1+3u^2}$;\;
$(1+3u^2)u'=\dst\frac{1}{ x}$;\;
$u+u^3=\ln |x|+c$;\;
$\dst{\frac{y}{ x}+\frac{y^3}{ x^3}}=\ln |x|+c$.


\item
Rewrite the given equation as $y'=\dst\frac{x^2-xy+y^2}{ xy}$;\;
$y=ux$;\ $u'x+u=\dst\frac{1}{ u}-1+u$;\;
$u'x=\dst\frac{1-u}{ u}$;\;
$\dst\frac{uu'}{ u-1}=-\dst\frac{1}{ x}$;\;
$\dst{\left[1+\frac{1}{ u-1}\right]u'}=-\dst\frac{1}{ x}$;\;
$u+\ln|u-1|=-\ln|x|+k$;\;
$e^u(u-1)=\dst\frac{c}{ x}$;\;
$e^{y/x}(y-x)=c$.



\item
$y=ux$;\;
$u'x+u=1+\dst\frac{1}{ u}+u$;\;
(A) $u'x=\dst\frac{u+1}{ u}$.
Since (A)  has the constant solution $u=-1$;
$y=-x$ is
a solution of the given equation. The nonconstant solutions of (A)
satisfy
$\dst\frac{uu'}{ u+1}=\dst\frac{1}{ x}$;\;
$\dst{\left[1-\frac{1 }{ u+1}\right]u'}=\frac{1}{ x}$;\;
$u-\ln |u+1|=\ln |x|+c$;\;
$\dst\frac{y}{ x}-\dst{\ln\left|\frac{y}{ x}-1\right|}=\ln|x|+c$;\;
$y-x\ln|y-x|=cy$.

\item
If $x=X-X_0$ and $y=Y-Y_0$, then $\dst\frac{dy}{ dx}=\dst\frac{dY}{
dx}=\dst\frac{dY}{ dX}\dst\frac{dX}{ dx}=\dst\frac{dY}{ dX}$, so $y=y(x)$
satisfies the given equation if and only if $Y=Y(X)$ satisfies
\[
\frac{dY}{ dX}=\dfrac{a(X-X_0)+b(Y-Y_0)+\alpha}{
c(X-X_0)+d(Y-Y_0)+\beta},
\]
 which reduces to the nonlinear homogeneous equation
\[
\frac{dY}{ dX}=\dst\frac{aX+bY}{ cX+dY}
\]
 if and only if
\[
\begin{aligned}
aX_0+bY_0&=\alpha\\
cX_0+dY_0&=\beta.
\end{aligned}
\tag{B}
\]
We will now show that
if $ad-bc\ne0$, then it is possible (for any choice of $\alpha$
and $\beta$) to solve (B).
Multiplying the first equation in (B) by $d$ and
 the second by  $b$ yields
\begin{align*}
daX_0+dbY_0&=d\alpha  \\
bcX_0+bdY_0&=b\beta.
\end{align*}
Subtracting the second of these equations from the first yields
 $(ad-bc)X_0=\alpha d-\beta b$.
Since $ad-bc\ne0$, this implies that $X_0=\dst\frac{\alpha d-\beta b}{
ad-bc}$. Multiplying the first equation in (B) by
$c$ and the second by $a$ yields
\begin{align*}
caX_0+cbY_0&=c\alpha  \\
acX_0+adY_0&=a\beta.
\end{align*}
Subtracting the first of these equation from the second yields
 $(ad-bc)Y_0=\alpha c-\beta a$. Since $ad-bc\ne0$ this implies that
$Y_0=\dst\frac{\alpha c-\beta a}{ ad-bc}$.


\item
For the given equation, (B)
of Exercise~2.4.40 is
\begin{align*}
2X_0+Y_0&=\phantom{-}1\\
X_0+2Y_0&=-4.
\end{align*}
Solving this pair of equations yields   $X_0=2$ and $Y_0=-3$.
The transformed differential equation is
\[
\frac{dY}{ dX}=\dst\frac{2X+Y}{ X+2Y}.
\tag{A}
\]
Let $Y=uX$;\;
$u'X+U=\dst\frac{2+u}{1+2u}$;\;
(B) $u'X=-\dfrac{2(u-1)(u+1)}{ 2u+1}$. Since $u\equiv1$ and
$u\equiv-1$
satisfy (B), $Y=X$ and $Y=-X$ are solutions of (A).
 Since $X=x+2$ and $Y=y-3$, it follows that $y=x+5$
and $y=-x+1$ are solutions of the given equation. The nonconstant
solutions of (B) satisfy
$\dfrac{(2u+1)u'}{(u-1)(u+1)}=-\dst\frac{2}{ X}$;\;
$\dst{\left[\frac{1}{ u+1}+\frac{3}{ u-1}\right]u'}=-\dst\frac{4}{
X}$;\;
$\ln|u+1|+3\ln|u-1|=-4\ln|X|+k$;\;
$\dst{(u+1)(u-1)^3}=\dst\frac{c}{ X^4}$;\;
$(Y+X)(Y-X)^3=c$;\;
Setting $X=x+2$ and $Y=y-3$ yields
$(y+x-1)(y-x-5)^3=c$.


\item
Rewrite the given equation as $y'=\dst\frac{y^3+x}{3xy^2}$;\;
$y=ux^{1/3}$;\;
$u'x^{1/3}+\dst\frac{1}{3x^{2/3}}u=\dst\frac{u^3+1}{3u^2x^{2/3}}$;\;
$u'x^{1/3}=\dst\frac{1}{3x^{2/3}u^2}$;\;
$u^2u'=\dst\frac{1}{3x}$;\;
$\dst\frac{u^3}{3}=\dst\frac{1}{3}(\ln|x|+c)$;\;
$u=(\ln|x|+c)^{1/3}$;\;
$y=x^{1/3}(\ln|x|+c)^{1/3}$.


\item
Rewrite the given equation as $y'=\dfrac{2(y^2+x^2y-x^4)}{
x^3}$;\;
$y=ux^2$;\;
$u'x^2+2xu=2x(u^2+u-1)$;\;
(A) $u'x^2=2x(u^2-1)$.
Since $u\equiv1$ and $u\equiv-1$ are constant solutions of (A),
$y=x^2$ and $y=-x^2$ are solutions of the given equation. The
nonconstant solutions of (A) satisfy
$\dst\frac{u'}{ u^2-1}=\dst\frac{2}{ x}$;\;
$\dst{\left[\frac{1}{ u-1}-\frac{1}{ u+1}\right]u'}=\dst\frac{4}{ x}$;\;
$\dst{\ln\left|\frac{u-1}{ u+1}\right|}=4\ln|x|+k$;\;
$\dst\frac{u-1}{ u+1}=cx^4$;\;
$(u-1)=(u+1)cx^4$;\;
$u(1-cx^4)=1+cx^4$;\;
$u=\dst\frac{1+cx^4}{1-cx^4}$;\;
$y=\dfrac{x^2(1+cx^4)}{1-cx^4}$.


\item
$y=u\tan x$;\;
$u'\tan x+u\sec^2x=(u^2+u+1)\sec^2x$;\;
$u'\tan x=(u^2+1)\sec^2x$;\;
$\dst\frac{u'}{ u^2+1}=\sec^2x\cot x=\cot x+\tan x$;\;
$\tan^{-1}u=\ln|\sin x|-\ln|\cos x|+c=\ln|\tan x|+c$;\;
$u=\tan(\ln|\tan x|+c)$;\;
$y=\tan x\ \tan(\ln|\tan x|+c)$.


\item
Rewrite the given equation as $y'=\dfrac{(y+\sqrt
x)^2}{2x(y+2\sqrt x)}$;\;
$y=ux^{1/2}$;\;
$u'x^{1/2}+\dst\frac{1}{2\sqrt x}u=\dfrac{(u+1)^2}{2\sqrt x\
(u+2)}$;\;
$u'x^{1/2}=\dst\frac{1}{2\sqrt x\ (u+2)}$;\;
$(u+2)u'=\dst\frac{1}{2x}$;\;
$\dfrac{(u+2)^2}{2}=\dst\frac{1}{2}(\ln|x|+c)$;\;
$(u+2)^2=\ln|x|+c$;\;
$u=-2\pm\sqrt{\ln|x|+c}$;\;
$y=x^{1/2}(-2\pm\sqrt{\ln|x|+c})$.


\item
$y_1=\dst\frac{1}{ x^2}$ is a solution of $y'+\dst\frac{2}{ x}y=0$.
Let $y=\dst\frac{u}{ x^2}$; then
\[
\dst\frac{u'}{ x^2}=\dfrac{3x^2(u^2/x^4)+6x(u/x^2)+2}{
x^2\left(2x(u/x^2)+3\right)}=
\dfrac{3(u/x)^2+6(u/x)+2}{ x^2\left(2(u/x)+3\right)},
\]
so (A) $u'=\dfrac{3(u/x)^2+6(u/x)+2}{2(u/x)+3}$. Since (A)
is a homogeneous nonlinear equation, we now substitute $u=vx$
into (A). This yields
$v'x+v=\dst\frac{3v^2+6v+2}{2v+3}$;\;
$v'x=\dfrac{(v+1)(v+2)}{2v+3}$;\;
$\dfrac{(2v+3)v'}{(v+1)(v+2)}=\dst\frac{1}{ x}$;\;
$\dst{\left[\frac{1}{ v+1}+\frac{1}{ v+2}\right]v'}=\dst\frac{1}{ x}$;\;
$\ln|(v+1)(v+2)|=\ln|x|+k$;\;
(B) $(v+1)(v+2)=cx$.
Since $y(2)=2\Rightarrow u(2)=8\Rightarrow v(2)=4$, (B) implies that
$c=15$.\ $(v+1)(v+2)=15x$;\;
$v^2+3v+2-15x=0$. From the quadratic formula,
$v=\dfrac{-3+\sqrt{1+60x}}{2}$;\;
$u=vx=\dfrac{x(-3+\sqrt{1+60x})}{2}$;\;
$y=\dst\frac{u}{ x^2}=\dfrac{-3+\sqrt{1+60x}}{2x}$.


\item
Differentiating  (A) $y_1(x)=\dfrac{y(ax)}{ a}$   yields
(B) $y_1'(x)=\dst\frac{1}{ a}y'(ax)\cdot a=y'(ax)$. Since
$y'(x)=q(y(x)/x)$ on some interval $I$, (C) $y'(ax)=q(y(ax)/ax)$
on some interval $J$. Substituting (A) and (B) into (C)
yields $y_1'(x)=q(y_1(x)/x)$ on $J$.


\item
  If $y=z+1$, then
$z'+z=xz^2$;\;
$z=ue^{-x}$;\;
$u'e^{-x}=xu^2e^{-2x}$;\;
$\dst\frac{u'}{ u^2}=xe^{-x}$;\;
$-\dst\frac{1}{ u}=-e^{-x}(x+1)-c$;\;
$u=\dst\frac{1}{ e^{-x}(x+1)+c}$;\;
$z=\dst\frac{1}{ x+1+ce^x}$;\;
$y=1+\dst\frac{1}{ x+1+ce^x}$.


\item
If $y=z+1$, then
$z'+\dst\frac{2}{ x}z=z^2$;\;
$z_1=\dst\frac{1}{ x^2}$;\;
$z=\dst\frac{u}{ x^2}$;\;
$\dst\frac{u'}{ x^2}=\dst\frac{u^2}{ x^4}$;\;
$\dst\frac{u'}{ u^2}=\dst\frac{1}{ x^2}$;\;
$-\dst\frac{1}{ u}=-\dst\frac{1}{ x}+c=-\dst\frac{1-cx}{ x}$;\;
$u=-\dst\frac{x}{ 1-cx}$;\;
$z=-\dst\frac{1}{ x(1-cx)}$;\;
$y=1-\dst\frac{1}{ x(1-cx)}$.

\end{sectionSolutions}

\section{Exact Equations}

\begin{sectionSolutions}
\item
$M(x,y)=3y\cos x+4xe^x+2x^2e^x$;\;
$N(x,y)=3\sin x+3$;\;
$M_y(x,y)=3\cos x=N_x(x,y)$,
so the  equation is exact.
We must find $F$ such that
(A) $F_x(x,y)=3y\cos x+4xe^x+2x^2e^x$ and
(B) $F_y(x,y)=3\sin x+3$.
Integrating (B) with respect to $y$ yields
(C) $F(x,y)=3y\sin x+3y+\psi(x)$.
Differentiating (C) with respect to $x$  yields
(D) $F_x(x,y)=3y\cos x+\psi'(x)$.
Comparing (D) with (A)  shows that
(E) $\psi'(x)=4xe^x+2x^2e^x$.
Integration by parts yields
$\dst{\int}xe^x\,dx=xe^x-e^x$ and
$\dst{\int}x^2e^x\,dx=x^2e^x-2xe^x+2e^x$.
Substituting from the last two equations into (E) yields
$\psi(x)=2x^2e^x$.
Substituting this into (C) yields
$F(x,y)=3y\sin x+3y+2x^2e^x$.
Therefore, $3y\sin x+3y+2x^2e^x=c$.


\item
$M(x,y)=2x-2y^2$;\;
$N(x,y)=12y^2-4xy$;\;
$M_y(x,y)=-4y=N_x(x,y)$,
so the  equation is exact.
We must find $F$ such that
(A) $F_x(x,y)=2x-2y^2$ and
(B) $F_y(x,y)=12y^2-4xy$.
Integrating (A) with respect to $x$ yields
(C) $F(x,y)=x^2-2xy^2+\phi(y)$.
Differentiating (C) with respect to $y$  yields
(D) $F_y(x,y)=-4xy+\phi'(y)$.
Comparing (D) with (B)  shows that
$\phi'(y)=12y^2$, so we take
$\phi(y)=4y^3$.
Substituting this into (C) yields
$F(x,y)=x^2-2xy^2+4y^3$.
Therefore, $x^2-2xy^2+4y^3=c$.


\item
$M(x,y)=4x+7y$;\;
$N(x,y)=3x+4y$;\;
$M_y(x,y)=7\ne3=N_x(x,y)$,
so the  equation is not exact.

\item
$M(x,y)=2x+y$;\;
$N(x,y)=2y+2x$;\;
$M_y(x,y)=1\ne2=N_x(x,y)$,
so the  equation is not exact.


\item
$M(x,y)=2x^2+8xy+y^2$;\;
$N(x,y)=2x^2+\dst\frac{xy^3}{3}$;\;
$M_y(x,y)=8x+2y\ne4x+\dst\frac{y^3}{3}=N_x(x,y)$,
so the  equation is not exact.


\item
$M(x,y)=y\sin xy+xy^2\cos xy$;\;
$N(x,y)=x\sin xy+xy^2\cos xy$;\;
$M_y(x,y)=3xy\cos xy+(1-x^2y^2)\sin xy\ne(xy+y^2)\cos
xy+(1-xy^3)\sin xy=N_x(x,y)$, so the equation is not exact.


\item
$M(x,y)=e^x(x^2y^2+2xy^2)+6x$;\;
$N(x,y)=2x^2ye^x+2$;\;
$M_y(x,y)=2xye^x(x+2)=N_x(x,y)$,
so the  equation is  exact.
We must find $F$ such that
(A) $F_x(x,y)=e^x(x^2y^2+2xy^2)+6x$ and
(B) $F_y(x,y)=2x^2ye^x+2$.
Integrating (B) with respect to $y$ yields
(C) $F(x,y)=x^2y^2e^x+2y+\psi(x)$.
Differentiating (C) with respect to $x$  yields
(D) $F_x(x,y)=e^x(x^2y^2+2xy^2)+\psi'(x)$.
Comparing (D) with (A)  shows that
$\psi'(x)=6x$, so we take
$\psi(x)=3x^2$.
Substituting this into (C) yields
$F(x,y)=x^2y^2e^x+2y+3x^2$.
Therefore, $x^2y^2e^x+2y+3x^2=c$.


\item
$M(x,y)=e^{xy}(x^4y+4x^3)+3y$;\;
$N(x,y)=x^5e^{xy}+3x$;\;
$M_y(x,y)=x^4e^{xy}(xy+5)+3=N_x(x,y)$,
so the  equation is  exact.
We must find $F$ such that
(A) $F_x(x,y)=e^{xy}(x^4y+4x^3)+3y$ and
(B) $F_y(x,y)=x^5e^{xy}+3x$.
Integrating (B) with respect to $y$ yields
(C) $F(x,y)=x^4e^{xy}+3xy+\psi(x)$.
Differentiating (C) with respect to $x$  yields
(D) $F_x(x,y)=e^{xy}(x^4y+4x^3)+3y+\psi'(x)$.
Comparing (D) with (A)  shows that
$\psi'(x)=0$, so we take
$\psi(x)=0$.
Substituting this into (C) yields
$F(x,y)=x^4e^{xy}+3xy$.
Therefore,  $x^4e^{xy}+3xy=c$.


\item
$M(x,y)=4x^3y^2-6x^2y-2x-3$;\;
$N(x,y)=2x^4y-2x^3$;\;
$M_y(x,y)=8x^3y-6x^2=N_x(x,y)$,
so the  equation is exact.
We must find $F$ such that
(A) $F_x(x,y)=4x^3y^2-6x^2y-2x-3$ and
(B) $F_y(x,y)=2x^4y-2x^3$.
Integrating (A) with respect to $x$ yields
(C) $F(x,y)=x^4y^2-2x^3y-x^2-3x+\phi(y)$.
Differentiating (C) with respect to $y$  yields
(D) $F_y(x,y)=2x^4y-2x^3+\phi'(y)$.
Comparing (D) with (B)  shows that
$\phi'(y)=0$, so we take
$\phi(y)=0$.
Substituting this into (C) yields
$F(x,y)=x^4y^2-2x^3y-x^2-3x$.
Therefore, $x^4y^2-2x^3y-x^2-3x=c$.
Since $y(1)=3\Rightarrow c=-1$,
$x^4y^2-2x^3y-x^2-3x+1=0$ is an implicit solution of the initial
value problem. Solving this for $y$ by means of the quadratic formula
yields $y=\dfrac{x+\sqrt{2x^2+3x-1}}{ x^2}$.


\item
$M(x,y)=(y^3-1)e^x$;\;
$N(x,y)=3y^2(e^x+1)$;\;
$M_y(x,y)=3y^2e^x=N_x(x,y)$,
so the  equation is exact.
We must find $F$ such that
(A) $F_x(x,y)=(y^3-1)e^x$ and
(B) $F_y(x,y)=3y^2(e^x+1)$.
Integrating (A) with respect to $x$ yields
(C) $F(x,y)=(y^3-1)e^x+\phi(y)$.
Differentiating (C) with respect to $y$  yields
(D) $F_y(x,y)=3y^2e^x+\phi'(y)$.
Comparing (D) with (B)  shows that
$\phi'(y)=3y^2$, so we take
$\phi(y)=y^3$.
Substituting this into (C) yields
$F(x,y)=(y^3-1)e^x+y^3$.
Therefore, $(y^3-1)e^x+y^3=c$.
Since $y(0)=0\Rightarrow c=-1$,
$(y^3-1)e^x+y^3=-1$ is an implicit solution of  the initial value
problem. Therefore, $y^3(e^x+1)=e^x-1$, so
$y=\dst{\left(\frac{e^x-1}{ e^x+1}\right)^{1/3}}$.


\item
$M(x,y)=(2x-1)(y-1)$;\;
$N(x,y)=(x+2)(x-3)$;\;
$M_y(x,y)=2x-1=N_x(x,y)$,
so the  equation is exact.
We must find $F$ such that
(A) $F_x(x,y)=(2x-1)(y-1)$ and
(B) $F_y(x,y)=(x+2)(x-3)$.
Integrating (A) with respect to $x$ yields
(C) $F(x,y)=(x^2-x)(y-1)+\phi(y)$.
Differentiating (C) with respect to $y$  yields
(D) $F_y(x,y)=x^2-x+\phi'(y)$.
Comparing (D) with (B)  shows that
$\phi'(y)=-6$, so we take
$\phi(y)=-6y$.
Substituting this into (C) yields
$F(x,y)=(x^2-x)(y-1)-6y$.
Therefore, $(x^2-x)(y-1)-6y=c$.
Since $y(1)=-1\Rightarrow c=6$,
$(x^2-x)(y-1)-6y=6$ is an implicit solution of  the initial value
problem. Therefore, $(x^2-x-6)y=x^2-x+6$, so
$y=\dst\frac{x^2-x+6}{(x-3)(x+2)}$.


\item
$M(x,y)=e^x(x^4y^2+4x^3y^2+1)$;\;
$N(x,y)=2x^4ye^x+2y$;\;
$M_y(x,y)=2x^3ye^x(x+4)=N_x(x,y)$,
so the  equation is  exact.
We must find $F$ such that
(A) $F_x(x,y)=e^x(x^4y^2+4x^3y^2+1)$ and
(B) $F_y(x,y)=2x^4ye^x+2y$.
Integrating (B) with respect to $y$ yields
(C) $F(x,y)=x^4y^2e^x+y^2+\psi(x)$.
Differentiating (C) with respect to $x$  yields
(D) $F_x(x,y)=e^xy^2(x^4+4x^3)+\psi'(x)$.
Comparing (D) with (A)  shows that
$\psi'(x)=e^x$, so we take
$\psi(x)=e^x$.
Substituting this into (C) yields
$F(x,y)=(x^4y^2+1)e^x+y^2$.
Therefore, $(x^4y^2+1)e^x+y^2=c$.

\addtocounter{enumi}{2}
\item%2.5.28
$M(x,y)=x^2+y^2$;\;
$N(x,y)=2xy$;\;
$M_y(x,y)=2y=N_x(x,y)$,
so the  equation is exact.
We must find $F$ such that
(A) $F_x(x,y)=x^2+y^2$ and
(B) $F_y(x,y)=2xy$.
Integrating (A) with respect to $x$ yields
(C) $F(x,y)=\dst\frac{x^3}{3}+xy^2+\phi(y)$.
Differentiating (C) with respect to $y$  yields
(D) $F_y(x,y)=2xy+\phi'(y)$.
Comparing (D) with (B)  shows that
$\phi'(y)=0$, so we take
$\phi(y)=0$.
Substituting this into (C) yields
$F(x,y)=\dst\frac{x^3}{3}+xy^2$.
Therefore, $\dst\frac{x^3}{3}+xy^2=c$.


\item
\begin{alist}
\item
Exactness requires that
$N_x(x,y)=M_y(x,y)=\dst\frac{\partial}{\partial
y}(x^3y^2+2xy+3y^2)=2x^3y+2x+6y$.
 Hence, $N(x,y)=\dst\frac{x^4y}{4}+x^2+6xy+g(x)$,
where $g$ is  differentiable.

\item
Exactness requires that
$N_x(x,y)=M_y(x,y)=\dst\frac{\partial}{\partial y}(\ln xy+2y\sin
x)=\dst\frac{1}{ y}+2\sin x$.
 Hence, $N(x,y)=\dst\frac{x}{ y}-2\cos x+g(x)$,
where $g$ is  differentiable.

\item
Exactness requires that
$N_x(x,y)=M_y(x,y)=\dst\frac{\partial}{\partial y}(x\sin x+y\sin
y)=y\cos y+\sin y$.
 Hence, $N(x,y)=x(y\cos y+\sin y)+g(x)$,
where $g$ is  differentiable.
\end{alist}

\item
 The assumptions imply that
$\dst\frac{\partial M_1}{\partial y}=\dst\frac{\partial N_1}{\partial x}$\
and\ $\dst\frac{\partial M_2}{\partial y}=\dst\frac{\partial N_2}{\partial
x}$. Therefore, $\dst\frac{\partial }{\partial
y}(M_1+M_2)=\dst\frac{\partial M_1}{\partial y}+\dfrac{\partial
M_2}{\partial y}=\dst\frac{\partial N_1}{\partial x}+\dfrac{\partial
N_2}{\partial x}=\dst\frac{\partial }{\partial x}(N_1+N_2)$,
which implies that $(M_1+M_2)\,dx+(N_1+N_2)\,dy=0$ is exact on $R$.


\item
Here $M(x,y)=Ax^2+Bxy+Cy^2$ and
$N(x,y)=Dx^2+Exy+Fy^2$. Since $M_y=Bx+2Cy$ and $N_x=2Dx+Ey$, the
equation is exact if and only if $B=2D$ and $E=2C$.


\item
Differentiating (A)
$F(x,y)=\dst{\int^y_{y_0}N(x_0,s)\,ds}+\dst{\int^x_{x_0}M(t,y)\,dt}$
with respect to $x$ yields $F_x(x,y)=M(x,y)$, since the first
integral
in (A) is independent of $x$ and $M(t,y)$ is a continuous function
of
$t$ for each fixed $y$. Differentiating  (A)  with respect to $y$
and using the assumption that $M_y=N_x$ yields
$F_y(x,y)=N(x_0,y)+\dst{\int^x_{x_0}\frac{\partial M
}{\partial y}(t,y)\,dt}
=N(x_0,y)+\dst{\int^x_{x_0}\frac{\partial N
}{\partial x} (t,y)\,dt}
=N(x_0,y)+N(x,y)-N(x_0,y)=N(x,y)$.


\item
$y_1=\dst\frac{1}{ x^2}$ is a solution of $y'+\dst\frac{2}{ x}y=0$.
Let $y=\dst\frac{u}{ x^2}$; then
\[
\dst\frac{u'}{ x^2}=-\dfrac{2x(u/x^2)}{
\left(x^2+2x^2(u/x^2)+1\right)}=
-\dst\frac{2xu}{ x^2(x^2+2u+1)},
\]
so  $u'=-\dst\frac{2xu}{ x^2+2u+1}$, which can be rewritten as (A)
$2xu\,dx+(x^2+2u+1)\,du=0$. Since $\dst{\dst\frac{\partial }{\partial
u}}(2xu)=\dst{\dst\frac{\partial }{\partial x}}(x^2+2u+1)=2x$, (A)
is exact. To solve (A)
we must find $F$ such that
(A) $F_x(x,u)=2xu$ and
(B) $F_u(x,u)=x^2+2u+1$.
Integrating (A) with respect to $x$ yields
(C) $F(x,u)=x^2u+\phi(u)$.
Differentiating (C) with respect to $u$  yields
(D) $F_u(x,u)=x^2+\phi'(u)$.
Comparing (D) with (B)  shows that
$\phi'(u)=2u+1$, so we take
$\phi(u)=u^2+u$.
Substituting this into (C) yields
$F(x,u)=x^2u+u^2+u=u(x^2+u+1)$.
Therefore, $u(x^2+u+1)=c$.
Since $y(1)=-2\Rightarrow u(1)=-2,
c=0$. Therefore,  $u(x^2+u+1)=0$. Since $u\equiv0$ does not satisfy
$u(1)=-2$, it follows that $u=-x^2-1$ and $y=-1-\dst\frac{1}{ x^2}$.


\item
$y_1=e^{-x^2}$ is a solution of $y'+2xy=0$.
Let $y=ue^{-x^2}$; then
$u'e^{-x^2}=-e^{-x^2}\dst{\left(\frac{3x+2u}{ 2x+3u}\right)}$,
so  $u'=-\dst\frac{3x+2u}{2x+3u}$, which can be rewritten as (A)
$(3x+2u)\,dx+(2x+3u)\,du=0$. Since $\dst{\dfrac{\partial
}{\partial
u}}(3x+2u)=\dst{\dst\frac{\partial }{\partial x}}(2x+3u)=2$, (A)
is exact. To solve (A)
we must find $F$ such that
(A) $F_x(x,u)=3x+2u$ and
(B) $F_u(x,u)=2x+3u$.
Integrating (A) with respect to $x$ yields
(C) $F(x,u)=\dst\frac{3x^2}{2}+2xu+\phi(u)$.
Differentiating (C) with respect to $u$  yields
(D) $F_u(x,u)=2x+\phi'(u)$.
Comparing (D) with (B)  shows that
$\phi'(u)=3u$, so we take
$\phi(u)=\dst\frac{3u^2}{2}$.
Substituting this into (C) yields
$F(x,u)=\dst\frac{3x^2}{2}+2xu+\dst\frac{3u^2}{2}$.
Therefore, $\dst\frac{3x^2}{2}+2xu+\dst\frac{3u^2}{2}=c$.
Since $y(0)=-1\Rightarrow u(0)=-1,
c=\dst\frac{3}{2}$. Therefore,  $3x^2+4xu+3u^2=3$
is an implicit solution of the initial value problem.
Rewriting this as $3u^2+4xu+(3x^2-3)=0$ and
 solving  for $u$ by means of the quadratic formula
yields $u=-\dst{\left(\frac{2x+\sqrt{9-5x^2}}{3}\right)}$, so
 $y=-e^{-x^2}\dst{\left(\frac{2x+\sqrt{9-5x^2}}{3}\right)}$.


\item
Since $M\,dx+N\,dy=0$ is exact,  (A) $M_y=N_x$.
Since $-N\,dx+M\,dy=0$ is exact,  (B) $M_x=-N_y$.
Differentiating (A) with respect to $y$ and (B) with respect to
$x$ yields (C) $M_{yy}=N_{xy}$ and (D) $M_{xx}=-N_{yx}$.
Since $N_{xy}=N_{yx}$, adding (C) and (D) yields $M_{xx}+M_{yy}=0$.
Differentiating (A) with respect to $x$ and (B) with respect to
$y$ yields (E) $M_{yx}=N_{xx}$ and (F) $M_{xy}=-N_{yy}$.
Since $M_{xy}=M_{yx}$, subtracting (F) from (E) yields
$N_{xx}+N_{yy}=0$.


\item
\begin{alist}
\item
If $F(x,y)=x^2-y^2$, then $F_x(x,y)=2x$, $F_y(x,y)=-2y$,
$F_{xx}(x,y)=2$, and $F_{yy}(x,y)=-2$.
Therefore, $F_{xx}+F_{yy}=0$, and $G$ must satisfy
(A) $G_x(x,y)=2y$ and (B) $G_y(x,y)=2x$.
Integrating (A) with respect to $x$ yields
(C) $G(x,y)=2xy+\phi(y)$.
Differentiating
(C) with respect to
$y$ yields
(D) $G_y(x,y)=2x+\phi'(y)$.
Comparing (D) with (B)  shows that
$\phi'(y)=0$, so we take
$\phi(y)=c$.
Substituting this into (C) yields
$G(x,y)=2xy+c$.

\item
If $F(x,y)=e^x\cos y$, then $F_x(x,y)=e^x\cos y$,
 $F_y(x,y)=-e^x\sin y$,
$F_{xx}(x,y)=e^x\cos y$, and $F_{yy}(x,y)=-e^x\cos y$.
Therefore, $F_{xx}+F_{yy}=0$, and $G$ must satisfy
(A) $G_x(x,y)=e^x\sin y$ and (B) $G_y(x,y)=e^x\cos y$.
Integrating (A) with respect to $x$ yields
(C) $G(x,y)=e^x\sin y+\phi(y)$.
Differentiating (C) with respect to $y$  yields
(D) $G_y(x,y)=e^x\cos y+\phi'(y)$.
Comparing (D) with (B)  shows that
$\phi'(y)=0$, so we take
$\phi(y)=c$.
Substituting this into (C) yields
$G(x,y)=e^x\sin y+c$.

\item
If $F(x,y)=x^3-3xy^2$ , then $F_x(x,y)=3x^2-3y^2$, $F_y(x,y)=-6xy$,
$F_{xx}(x,y)=6x$, and $F_{yy}(x,y)=-6x$.
Therefore, $F_{xx}+F_{yy}=0$, and $G$ must satisfy
(A) $G_x(x,y)=6xy$ and (B) $G_y(x,y)=3x^2-3y^2$.
Integrating (A) with respect to $x$ yields
(C) $G(x,y)=3x^2y+\phi(y)$.
Differentiating (C) with respect to $y$  yields
(D) $G_y(x,y)=3x^2+\phi'(y)$.
Comparing (D) with (B)  shows that
$\phi'(y)=-3y^2$, so we take
$\phi(y)=-y^3+c$.
Substituting this into (C) yields
$G(x,y)=3x^2y-y^3+c$.

\item
If $F(x,y)=\cos x\cosh y$, then $F_x(x,y)=-\sin x\cosh y$,
$F_y(x,y)=\cos x\sinh y$,
$F_{xx}(x,y)=-\cos x\cosh y$, and $F_{yy}(x,y)=\cos x\cosh y$.
Therefore, $F_{xx}+F_{yy}=0$, and $G$ must satisfy
(A) $G_x(x,y)=-\cos x\sinh y$ and (B) $G_y(x,y)=-\sin x\cosh y$.
Integrating (A) with respect to $x$ yields
(C) $G(x,y)=-\sin x\sinh y+\phi(y)$.
Differentiating (C) with respect to $y$  yields
(D) $G_y(x,y)=-\sin x\cosh y+\phi'(y)$.
Comparing (D) with (B)  shows that
$\phi'(y)=0$, so we take
$\phi(y)=c$.
Substituting this into (C) yields
$G(x,y)=-\sin x\sinh y+c$.

\item
If $F(x,y)=\sin x\cosh y$, then $F_x(x,y)=\cos x\cosh y$,
 $F_y(x,y)=\sin x\sinh y$,
$F_{xx}(x,y)=-\sin x\cosh y$, and $F_{yy}(x,y)=\sin x\cosh y$.
Therefore, $F_{xx}+F_{yy}=0$, and $G$ must satisfy
(A) $G_x(x,y)=-\sin x\sinh y$ and (B) $G_y(x,y)=\cos x\cosh y$.
Integrating (A) with respect to $x$ yields
(C) $G(x,y)=\cos x\sinh y+\phi(y)$.
Differentiating (C) with respect to $y$  yields
(D) $G_y(x,y)=\cos x\cosh y+\phi'(y)$.
Comparing (D) with (B)  shows that
$\phi'(y)=0$, so we take
$\phi(y)=c$.
Substituting this into (C) yields
$G(x,y)=\cos x\sinh y+c$.
\end{alist}

\end{sectionSolutions}

\section{Exact Equations}

\begin{sectionSolutions}
\item
\begin{alist}
\item[a and b]\qquad\quad To show that $\mu(x,y)=\dst\frac{1}{(x-y)^2}$
is an integrating
factor for (A) and that (B) is exact, it suffices to observe that
$\dst{\frac{\partial }{\partial x}\left(\frac{xy}{ x-y}\right)}=-
\dst\frac{y^2}{(x-y)^2}$ and
$\dst{\frac{\partial }{\partial y}\left(\frac{xy}{ x-y}\right)}=
\dst\frac{x^2}{(x-y)^2}$. By Theorem~2.5.1 this also shows that
(C) is an implicit solution of (B). Since
$\mu(x,y)$ is never zero, any solution of (B)  is a solution of~(A).

\addtocounter{enumXi}{2}
\item If we interpret (A) as $-y^2+x^2y'=0$, then substituting
$y=x$ yields $-x^2+x^2\cdot1=0$.
\end{alist}

\note{In Exercises~2.6.3--2.6.23, the given
equation is multiplied by an integrating factor to produce an exact
equation, and an implicit solution is found for the latter. For a
complete analysis of the relationship between the sets of solutions of
the two equations it is necessary to check for additional solutions of
the given equation ``along which'' the integrating factor is undefined,
or for solutions of the exact equation ``along which'' the integrating
factor vanishes. In the interests of brevity we omit these tedious
details except in cases where there actually is a difference between
the sets of solutions of the two equations.}


\item
$M(x,y)=3x^2y$;\;
$N(x,y)=2x^3$;\;
$M_y(x,y)-N_x(x,y)=3x^2-6x^2=-3x^2$;\;
$p(x)=\dfrac{M_y(x,y)-N_x(x,y)}{
N(x,y)}=-\dst\frac{3x^2}{2x^3}=-\dst\frac{3}{2x}$;\;
$\int p(x)\,dx=-\dst\frac{3}{2}\ln|x|$;\;
$\mu(x)=P(x)=x^{-3/2}$;
therefore
$3x^{1/2}y\,dx+2x^{3/2}\,dy=0$
is exact.
We must find $F$ such that
(A) $F_x(x,y)=3x^{1/2}y$ and
(B) $F_y(x,y)=2x^{3/2}$.
Integrating (A) with respect to $x$ yields
(C) $F(x,y)=2x^{3/2}y+\phi(y)$.
Differentiating (C) with respect to $y$  yields
(D) $F_y(x,y)=2x^{3/2}+\phi'(y)$.
Comparing (D) with (B)  shows that
$\phi'(y)=0$, so we take
$\phi(y)=0$.
Substituting this into (C) yields
$F(x,y)=2x^{3/2}y$,
so $x^{3/2}y=c$.


\item
$M(x,y)=5xy+2y+5$;\;
$N(x,y)=2x$;\;
$M_y(x,y)-N_x(x,y)=(5x+2)-2=5x$;\;
$p(x)=\dfrac{M_y(x,y)-N_x(x,y)}{
N(x,y)}=\dst\frac{5x}{2x}=\dst\frac{5}{2}$;\;
$\int p(x)\,dx=\dst\frac{5x}{2}$;\;
$\mu(x)=P(x)=e^{5x/2}$;
therefore
$e^{5x/2}(5xy+2y+5)\,dx+2xe^{5x/2}\,dy=0$
is exact.
We must find $F$ such that
(A) $F_x(x,y)=e^{5x/2}(5xy+2y+5)$ and
(B) $F_y(x,y)=2xe^{5x/2}$.
Integrating (B) with respect to $y$ yields
(C) $F(x,y)=2xye^{5x/2}+\psi(x)$.
Differentiating (C) with respect to $x$  yields
(D) $F_x(x,y)=5xye^{5x/2}+2ye^{5x/2}+\psi'(x)$.
Comparing (D) with (A)  shows that
$\psi'(x)=5e^{5x/2}$, so we take
$\psi(x)=2e^{5x/2}$.
Substituting this into (C) yields
$F(x,y)=2e^{5x/2}(xy+1)$,
so $e^{5x/2}(xy+1)=c$.


\item
$M(x,y)=27xy^2+8y^3$;\;
$N(x,y)=18x^2y+12xy^2$;\;
$M_y(x,y)-N_x(x,y)=(54xy+24y^2)-(36xy+12y^2)=18xy+12y^2$;\;
$p(x)=\dfrac{M_y(x,y)-N_x(x,y)}{ N(x,y)}=\dst\frac{18xy+12y^2}{
18x^2y+12y^2x}=\dst\frac{1}{ x}$;\;
$\int p(x)\,dx=\ln|x|$;\;
$\mu(x)=P(x)=x$;
therefore
$(27x^2y^2+8xy^3)\,dx+(18x^3y+12x^2y^2)\,dy=0$
is exact.
We must find $F$ such that
(A) $F_x(x,y)=27x^2y^2+8xy^3$ and
(B) $F_y(x,y)=18x^3y+12x^2y^2$.
Integrating (A) with respect to $x$ yields
(C) $F(x,y)=9x^3y^2+4x^2y^3+\phi(y)$.
Differentiating (C) with respect to $y$  yields
(D) $F_y(x,y)=18x^3y+12x^2y^2+\phi'(y)$.
Comparing (D) with (B)  shows that
$\phi'(y)=0$, so we take
$\phi(y)=0$.
Substituting this into (C) yields
$F(x,y)=9x^3y^2+4x^2y^3$,
so $x^2y^2(9x+4y)=c$.

\item
$M(x,y)=y^2$;\;
$N(x,y)= \dst{\left(xy^2+3xy+\frac{1}{ y}\right)}$;\;
$M_y(x,y)-N_x(x,y)=2y-(y^2+3y)=-y(y+1)$;\;
$q(y)=\dfrac{N_x(x,y)-M_y(x,y)}{ M(x,y)}=\dfrac{y(y+1)}{ y^2}=
1+\dst\frac{1}{ y}$;\;
$\int q(y)\,dy=y\ln|y|$;\;
$\mu(y)=Q(y)=ye^y$;
therefore
$y^3e^y\,dx+e^y(xy^3+3xy^2+1)\,dy=0$
is exact.
We must find $F$ such that
(A) $F_x(x,y)=y^3e^y$ and
(B) $F_y(x,y)=e^y(xy^3+3xy^2+1)$.
Integrating (A) with respect to $x$ yields
(C) $F(x,y)=xy^3e^y+\phi(y)$.
Differentiating (C) with respect to $y$  yields
(D) $F_y(x,y)=xy^3e^y+3xy^2e^y+\phi'(y)$.
Comparing (D) with (B)  shows that
$\phi'(y)=e^y$, so we take
$\phi(y)=e^y$.
Substituting this into (C) yields
$F(x,y)=xy^3e^y+e^y$,
so $e^y(xy^3+1)=c$.




\item
$M(x,y)=x^2y+4xy+2y$;\;
$N(x,y)=x^2+x$;\;
$M_y(x,y)-N_x(x,y)=(x^2+4x+2)-(2x+1)=x^2+2x+1=(x+1)^2$;\;
$p(x)=\dfrac{M_y(x,y)-N_x(x,y)}{ N(x,y)}=\dfrac{(x+1)^2}{
x(x+1)}=1+\dst\frac{1}{ x}$;\;
$\int p(x)\,dx=x+\ln|x|$;\;
$\mu(x)=P(x)=xe^x$;
therefore
$e^x(x^3y+4x^2y+2xy)\,dx+e^x(x^3+x^2)\,dy=0$
is exact.
We must find $F$ such that
(A) $F_x(x,y)=e^x(x^3y+4x^2y+2xy)$ and
(B) $F_y(x,y)=e^x(x^3+x^2)$.
Integrating (B) with respect to $y$ yields
(C) $F(x,y)=y(x^3+x^2)e^x+\psi(x)$.
Differentiating (C) with respect to $x$  yields
(D) $F_x(x,y)=e^x(x^3y+4x^2y+2xy)+\psi'(x)$.
Comparing (D) with (A)  shows that
$\psi'(x)=0$, so we take
$\psi(x)=0$.
Substituting this into (C) yields
$F(x,y)=y(x^3+x^2)e^x=x^2y(x+1)e^x$,
so $x^2y(x+1)e^x=c$.



\item
$M(x,y)=\cos x\cos y$;\;
$N(x,y)=\sin x\cos y-\sin x\sin y+y$;\;
$M_y(x,y)-N_x(x,y)=-\cos x\sin y-(\cos x\cos y-\cos x\sin y)
=-\cos x\cos y$;\;
$q(y)=\dfrac{N_x(x,y)-M_y(x,y)}{ M(x,y)}=\dfrac{\cos x\cos y
}{\cos x\cos y}=1$;\;
$\int q(y)\,dy=1$;\;
$\mu(y)=Q(y)=e^y$;
therefore
$e^y\cos x\cos y\,dx+e^y(\sin x\cos y-\sin x\sin y+y)\,dy=0$
is exact.
We must find $F$ such that
(A) $F_x(x,y)=e^y\cos x\cos y$ and
(B) $F_y(x,y)=e^y(\sin x\cos y-\sin x\sin y+y)$.
Integrating (A) with respect to $x$ yields
(C) $F(x,y)=e^y\sin x\cos y+\phi(y)$.
Differentiating (C) with respect to $y$  yields
(D) $F_y(x,y)=e^y(\sin x\cos y-\sin x\sin y)+\phi'(y)$.
Comparing (D) with (B)  shows that
$\phi'(y)=ye^y$, so we take
$\phi(y)=e^y(y-1)$.
Substituting this into (C) yields
$F(x,y)=e^y(\sin x\cos y+y-1)$,
so $e^y(\sin x\cos y+y-1)=c$.

\item
$M(x,y)=y\sin y$;\;
$N(x,y)=x(\sin y-y\cos y)$;\;
 $M_y(x,y)-N_x(x,y)=(y\cos y+\sin y)-(\sin y-y\cos y)=2y\cos y$;
$q(y)=\dfrac{N_x(x,y)-M_y(x,y)}{ N(x,y)}=-\dst\frac{2\cos y}{\sin y}$;\;
$\int q(y)\,dy=-2\ln|\sin y|$;\;
$\mu(y)=Q(y)=\dst\frac{1}{\sin^2y}$;
therefore
$\dst{\left(\frac{y}{\sin y}\right)}\,dx+
x\dst{\left(\frac{1}{\sin y}-\frac{y\cos y}{\sin^2y}\right)}\,dy=0$
is exact. We must find $F$ such that
(A) $F_x(x,y)=\dst\frac{y}{\sin y}$ and
(B) $F_y(x,y)=x\dst{\left(\frac{1}{\sin y}-\frac{y\cos y}{\sin^2y}\right)}$.
Integrating (A) with respect to $x$ yields
(C) $F(x,y)=\dst\frac{xy}{\sin y}+\phi(y)$.
Differentiating (C) with respect to $y$  yields
(D) $F_y(x,y)=x\dst{\left(\frac{1}{\sin y}-\frac{y\cos y}{\sin^2y}\right)}
+\phi'(y)$.
Comparing (D) with (B)  shows that
$\phi'(y)=0$, so we take
$\phi(y)=0$.
Substituting this into (C) yields
$F(x,y)=\dst\frac{xy}{\sin y}$,
so $\dst\frac{xy}{\sin y}=c$. In addition, the given equation has the
constant solutions $y=k\pi$, where $k$ is an integer.


\item
$M(x,y)=\alpha y+\gamma xy$;\;
$N(x,y)=\beta x+ \delta xy$;\;
 $M_y(x,y)-N_x(x,y)=(\alpha+\gamma x)-(\beta+\delta y)$;
and  $p(x)N(x,y)-q(y)M(x,y)=p(x)x(\beta + \delta y)-
q(y)y(\alpha +\gamma x)$.
so exactness requires that
 $(\alpha+\gamma x)-(\beta+\delta y)=
p(x)x(\beta+\delta y)-
q(y)y(\alpha+\gamma x)$,
which holds if
  $p(x)x=-1$ and $q(y)y=-1$. Thus
 $p(x)=-\dst\frac{1}{ x}$;\;
 $q(y)=-\dst\frac{1}{ y}$;\;
$\int p(x)\,dx=-\ln|x|$;\;
$\int q(y)\,dy=-\ln|y|$;\;
$P(x)=\dst\frac{1}{ x}$;
$Q(y)=\dst\frac{1}{ y}$;
$\mu(x,y)=\dst\frac{1}{ xy}$.
Therefore,
$\dst{\left(\frac{\alpha}{ x}+\gamma\right)}\,dx+
\dst{\left(\frac{\beta}{ y}+\delta\right)}\,dy=0$
is exact.
We must find $F$ such that
(A) $F_x(x,y)=\dst{\frac{\alpha}{ x}+\gamma}$ and
(B) $F_y(x,y)=\dst{\frac{\beta}{ y}+\delta}$.
Integrating (A) with respect to $x$ yields
(C) $F(x,y)=\alpha\ln|x|+\gamma x+\phi(y)$.
Differentiating (C) with respect to $y$  yields
(D) $F_y(x,y)=\phi'(y)$.
Comparing (D) with (B)  shows that
$\phi'(y)=\dst{\frac{\beta}{ y}+\delta}$, so we take
$\phi(y)=\beta\ln|y|+\delta y$.
Substituting this into (C) yields
$F(x,y)=\alpha\ln|x|+\gamma x+\beta\ln|y|+\delta y$,
so $|x|^\alpha|y|^\beta e^{\gamma x}e^{\delta y}=c$.
The given equation also has the solutions $x\equiv0$
and $y\equiv0$.



\item
$M(x,y)=2y$;\;
$N(x,y)=3(x^2+x^2y^3)$;\;
 $M_y(x,y)-N_x(x,y)=2-(6x+6xy^3)$;
and  $p(x)N(x,y)-q(y)M(x,y)=3p(x)(x^2+x^2y^3)-2q(y)y$.
so exactness requires that
(A) $2-6x-6xy^3=3p(x)x(x+xy^3)-2q(y)y$.
To obtain  similar terms on the two sides of (A)
we let  $p(x)x=a$ and $q(y)y=b$  where $a$ and $b$ are constants
such that $2-6x-6xy^3=3a(x+xy^3)-2b$, which  holds if
$a=-2$ and $b=-1$. Thus,
 $p(x)=-\dst\frac{2}{ x}$;\;
 $q(y)=-\dst\frac{1}{ y}$;\;
$\int p(x)\,dx=-2\ln|x|$;\;
$\int q(y)\,dy=-\ln|y|$;\;
$P(x)=\dst\frac{1}{ x^2}$;
$Q(y)=\dst\frac{1}{ y}$;
$\mu(x,y)=\dst\frac{1}{ x^2y}$.
Therefore,
$\dst\frac{2}{ x^2}\,dx+3\dst{\left(\frac{1}{ y}+y^2\right)}\,dy=0$
is exact.
We must find $F$ such that
(B) $F_x(x,y)=\dst\frac{2}{ x^2}$ and
(C) $F_y(x,y)=3\dst{\left(\frac{1}{ y}+y^2\right)}$.
Integrating (B) with respect to $x$ yields
(D) $F(x,y)=-\dst\frac{2}{ x}+\phi(y)$.
Differentiating (D) with respect to $y$  yields
(E) $F_y(x,y)=\phi'(y)$.
Comparing (E) with (C)  shows that
$\phi'(y)=3\dst{\left(\frac{1}{ y}+y^2\right)}$, so we take
$\phi(y)=y^3+3\ln|y|$.
Substituting this into (D) yields
$F(x,y)=-\dst\frac{2}{ x}+y^3+3\ln|y|$,
so $-\dst\frac{2}{ x}+y^3+3\ln|y|=c$.
The given equation also has the solutions $x\equiv0$
and $y\equiv0$.

\item
$M(x,y)=x^4y^4$;\;
$N(x,y)=x^5y^3$;\;
 $M_y(x,y)-N_x(x,y)=4x^4y^3-5x^4y^3=-x^4y^3$;
and  $p(x)N(x,y)-q(y)M(x,y)=p(x)x^5y^3-q(y)x^4y^4$.
so exactness requires that
 $-x^4y^3=p(x)x^5y^3-q(y)x^4y^4$, which is equivalent to
 $p(x)x-q(y)y=-1$. This holds if
 $p(x)x=a$ and $q(y)y=a+1$  where $a$ is an arbitrary real number.
Thus,
 $p(x)=\dst\frac{a}{ x}$;\;
 $q(y)=\dst\frac{a+1}{ y}$;\;
$\int p(x)\,dx=a\ln|x|$;\;
$\int q(y)\,dy=(a+1)\ln|y|$;\;
$P(x)=|x|^a$;
$Q(y)=|y|^{a+1}$;
$\mu(x,y)=|x^a||y|^{a+1}$.
Therefore,
$|x|^a|y|^{a+1}\left(x^4y^4\,dx+x^5y^3\,dy\right)=0$
is exact for any choice of $a$. For simplicity we let $a=-4$,
so (A)  is equivalent to $y\,dx+x\,dy=0$.
We must find $F$ such that
(B) $F_x(x,y)=y$ and
(C) $F_y(x,y)=x$.
Integrating (B) with respect to $x$ yields
(D) $F(x,y)=xy+\phi(y)$.
Differentiating (D) with respect to $y$  yields
(E) $F_y(x,y)=x+\phi'(y)$.
Comparing (E) with (C)  shows that
$\phi'(y)=0$, so we take
$\phi(y)=0$.
Substituting this into (D) yields
$F(x,y)=xy$,
so $xy=c$.


\item
$M(x,y)=x^4y^3+y$;\;
$N(x,y)=x^5y^2-x$;\;
$M_y(x,y)-N_x(x,y)=(3x^4y^2+1)-(5x^4y^2-1)=-2x^4y^2+2$;\;
$p(x)=\dfrac{M_y(x,y)-N_x(x,y)}{ N(x,y)}=-\dfrac{2x^4y^2-2
}{ x^5y^2-x}=-\dst\frac{2}{ x}$;\;
$\int p(x)\,dx=-2\ln|x|$;\;
$\mu(x)=P(x)=\dst\frac{1}{ x^2}$;
therefore
$\dst{\left(x^2y^3+\frac{y}{ x^2}\right)}\,dx+
\dst{\left(x^3y^2-\frac{1}{ x}\right)}\,dy=0$
is exact.
We must find $F$ such that
(A) $F_x(x,y)=\dst{\left(x^2y^3+\frac{y}{ x^2}\right)}$ and
(B) $F_y(x,y)=\dst{\left(x^3y^2-\dst\frac{1}{ x}\right)}$.
Integrating (A) with respect to $x$ yields
(C) $F(x,y)=\dst\frac{x^3y^3}{3}-\dst\frac{y}{ x}+\phi(y)$.
Differentiating (C) with respect to $y$  yields
(D) $F_y(x,y)=x^3y^2-\dst\frac{1}{ x}+\phi'(y)$.
Comparing (D) with (B)  shows that
$\phi'(y)=0$, so we take
$\phi(y)=0$.
Substituting this into (C) yields
$F(x,y)=\dst\frac{x^3y^3}{3}-\dst\frac{y}{ x}$,
so $\dst\frac{x^3y^3}{3}-\dst\frac{y}{ x}=c$.


\item
$M(x,y)=12xy+6y^3$;\;
$N(x,y)=9x^2+10xy^2$;\;
 $M_y(x,y)-N_x(x,y)=(12x+18y^2)-(18x+10y^2)=-6x+8y^2$;
and  $p(x)N(x,y)-q(y)M(x,y)=p(x)x(9x+10y^2)-q(y)y(12x+6y^2)$,
so exactness requires that
(A) $-6x+8y^2=p(x)x(9x+10y^2)-q(y)y(12x+6y^2)$.
To obtain  similar terms on the two sides of (A)
we let  $p(x)x=a$ and $q(y)y=b$  where $a$ and $b$ are constants
such that $-6x+8y^2=a(9x+10y^2)-b(12x+6y^2)$, which holds if
  $9a-12b=-6$, $10a-6b=8$; that is, $a=b=2$. Thus
 $p(x)=\dst\frac{2}{ x}$;\;
 $q(y)=\dst\frac{2}{ y}$;\;
$\int p(x)\,dx=2\ln|x|$;\;
$\int q(y)\,dy=2\ln|y|$;\;
$P(x)=x^2$;
$Q(y)=y^2$;
$\mu(x,y)=x^2y^2$.
Therefore,
$(12x^3y^3+6x^2y^5)\,dx+(9x^4y^2+10x^3y^4)\,dy=0$
is exact.
We must find $F$ such that
(B) $F_x(x,y)=12x^3y^3+6x^2y^5$ and
(C) $F_y(x,y)=9x^4y^2+10x^3y^4$.
Integrating (B) with respect to $x$ yields
(D) $F(x,y)=3x^4y^3+2x^3y^5+\phi(y)$.
Differentiating (D) with respect to $y$  yields
(E) $F_y(x,y)=9x^4y^2+10x^3y^4+\phi'(y)$.
Comparing (E) with (C)  shows that
$\phi'(y)=0$, so we take
$\phi(y)=0$.
Substituting this into (D) yields
$F(x,y)=3x^4y^3+2x^3y^5$,
so $x^3y^3(3x+2y^2)=c$.



\item
$M(x,y)=ax^my+by^{n+1}$;\ $N(x,y)=cx^{m+1}+dxy^n$;\;
$M_y(x,y)-N_x(x,y)=\left[ax^{m+1}+(n+1)by^n\right]-\left[(m+1)cx^m
+dy^n\right]$;\;
$p(x)N(x,y)-q(y)M(x,y)=xp(x)(cx^m+dy^n)-yp(y)(ax^m+by^n)$.
Let (A)  $xp(x)=\alpha$ and (B) $yp(y)=\beta$, where $\alpha$
and $\beta$ are to be chosen so that
$\left[ax^{m+1}+(n+1)by^n\right]-\left[(m+1)cx^m
+dy^n\right]=\alpha(cx^m+dy^n)-\beta(ax^m+by^n)$, which will hold if
\[
\begin{aligned}
c\alpha-a\beta&=\phantom{-}a-(m+1)c=_{\mathrm{df}}&A\\
d\alpha-b\beta&=-d+\,(n+1)b=_{\mathrm{df}}&B.
\end{aligned}
\tag{C}
\]
Since $ad-bc\ne0$ it can be verified that
$\alpha=\dst\frac{aB-bA}{ ad-bc}$ and $\beta=\dst\frac{cB-dA}{ ad-bc}$
satisfy (C). From (A) and (B),
$p(x)=\dst\frac{\alpha}{ x}$ and $q(y)=\dst\frac{\beta}{ y}$, so
$\mu(x,y)=x^\alpha y^\beta$  is an integrating factor for the given
equation.


\item
\begin{alist}
\item
Since $M(x,y)=p(x)y-f(x)$  and $N(x,y)=1$,
$\dfrac{M_y(x,y)-N_x(x,y)}{ N(x,y)}=p(x)$ and Theorem~2.6.1
implies that $\mu(x)\pm e^{\int p(x)\,dx}$  is an integrating factor
for (C).

\item  Multiplying (A) through $\mu=\pm e^{\int p(x)\,dx}$
yields  (D) $\mu(x)y'+\mu'(x)y=\mu(x)f(x)$, which is equivalent to
$(\mu(x)y)'=\mu(x)f(x)$. Integrating this yields
$\dst\mu(x)y=c+\int\mu(x)f(x)\,dx$, so
$y=\dst\frac{1}{\mu(x)}\left(c+\int
\mu(x)f(x)\,dx\right)$, which  is equivalent to (B) since
$y_1=\dst\frac{1}{\mu}$ is a nontrivial solution of $y'+p(x)y=0$.
\end{alist}
\end{sectionSolutions}

\chapter{Numerical Methods}

\section{Euler's Method}

\begin{sectionSolutions}
\item
$y_1=1.200000000,\ y_2=1.440415946,\ y_3=1.729880994$

\item
$y_1=2.962500000,\ y_2=2.922635828,\ y_3=2.880205639$
\bigskip

\item
{\small
\tagpdfsetup{table/header-rows={1},table/header-columns={1}}
\begin{tabular}{S[table-format=1.1] *4{S[table-format=2.9]} }\toprule
{$x$}&{$h=0.1$}&{$h=0.05$}&{$h=0.025$}&{Exact}\\ \midrule
0.0 &  2.000000000 &   2.000000000 &  2.000000000 &   2.000000000 \\
0.1 &  2.100000000 &   2.169990965 &  2.202114518 &   2.232642918 \\
0.2 &  2.514277288 &   2.649377900 &  2.713011720 &   2.774352565 \\
0.3 &  3.317872752 &   3.527672599 &  3.628465025 &   3.726686582 \\
0.4 &  4.646592772 &   4.955798226 &  5.106379369 &   5.254226636 \\
0.5 &  6.719737638 &   7.171467977 &  7.393322991 &   7.612186259 \\
0.6 &  9.876155616 &  10.538384528 & 10.865186799 &  11.188475269 \\
0.7 & 14.629532397 &  15.605686107 & 16.088630652 &  16.567103199 \\
0.8 & 21.751925418 &  23.197328550 & 23.913328531 &  24.623248150 \\
0.9 & 32.399118931 &  34.545932627 & 35.610005377 &  36.665439956 \\
1.0 & 48.298147362 &  51.492825643 & 53.076673685 &  54.647937102 \\
\bottomrule
\end{tabular}}\bigskip

\item
{\small
\tagpdfsetup{table/header-rows={1},table/header-columns={1}}
\begin{tabular}{S[table-format=1.2] *3{S[table-format=2.9]} S[table-format=2.9] }\toprule
{$x$}&{$h=0.05$}&{$h=0.025$}&{$h=0.0125$}&{Exact}\\ \midrule
1.00 &  2.000000000 & 2.000000000 & 2.000000000 &  2.000000000 \\
1.05 &  2.250000000 & 2.259280190 & 2.264490570 &  2.270158103 \\
1.10 &  2.536734694 & 2.559724746 & 2.572794280 &  2.587150838 \\
1.15 &  2.867950854 & 2.910936426 & 2.935723355 &  2.963263785 \\
1.20 &  3.253613825 & 3.325627715 & 3.367843117 &  3.415384615 \\
1.25 &  3.706750613 & 3.820981064 & 3.889251900 &  3.967391304 \\
1.30 &  4.244700641 & 4.420781829 & 4.528471927 &  4.654198473 \\
1.35 &  4.891020001 & 5.158883503 & 5.327348558 &  5.528980892 \\
1.40 &  5.678467290 & 6.085075790 & 6.349785943 &  6.676923077 \\
1.45 &  6.653845988 & 7.275522641 & 7.698316221 &  8.243593315 \\
1.50 &  7.886170437 & 8.852463793 & 9.548039907 & 10.500000000 \\
\bottomrule
\end{tabular}}\bigskip

\item
{\small
\tagpdfsetup{table/header-rows={1,2},table/header-columns={1}}
\begin{tabular}{S[table-format=1.1] *3{S[table-format=1.9]} *3{S[table-format=-1.4]} }\toprule
&\multicolumn{3}{c}{Approximate Solutions}&
\multicolumn{3}{c}{Residuals}\\\cmidrule(r){2-4}\cmidrule(l){5-7}
{$x$}&{$h=0.1$}&{$h=0.05$}&{$h=0.025$}&{$h=0.1$}&{$h=0.05$}&{$h=0.025$}\\ \midrule
1.0 & 1.000000000 & 1.000000000 & 1.000000000 & 0.0000 & 0.0000  &  0.0000  \\
1.1 & 0.920000000 & 0.921898275 & 0.922822717 &-0.0384 & -0.0189 & -0.0094  \\
1.2 & 0.847469326 & 0.851018464 & 0.852746371 &-0.0745 & -0.0368 & -0.0183  \\
1.3 & 0.781779403 & 0.786770087 & 0.789197876 &-0.1092 & -0.0540 & -0.0268  \\
1.4 & 0.722453556 & 0.728682209 & 0.731709712 &-0.1428 & -0.0707 & -0.0351  \\
1.5 & 0.669037867 & 0.676299618 & 0.679827306 &-0.1752 & -0.0868 & -0.0432  \\
1.6 & 0.621054176 & 0.629148585 & 0.633080163 &-0.2062 & -0.1023 & -0.0509  \\
1.7 & 0.578000416 & 0.586740390 & 0.590986601 &-0.2356 & -0.1170 & -0.0583  \\
1.8 & 0.539370187 & 0.548588902 & 0.553070392 &-0.2631 & -0.1310 & -0.0653  \\
1.9 & 0.504674296 & 0.514228603 & 0.518877246 &-0.2889 & -0.1441 & -0.0719  \\
2.0 & 0.473456737 & 0.483227470 & 0.487986391 &-0.3129 & -0.1563 & -0.0781 \\
\bottomrule
\end{tabular}}\bigskip

\item
{\small
\tagpdfsetup{table/header-rows={1},table/header-columns={1}}
\begin{tabular}{S[table-format=1.1] *4{S[table-format=-1.9]}  }\toprule
{$x$}&{$h=0.1$}&{$h=0.05$}&{$h=0.025$}&{``Exact''}\\ \midrule
1.0 &  0.000000000 &  0.000000000 &  0.000000000 &  0.000000000 \\
1.1 & -0.100000000 & -0.099875000 & -0.099780455 & -0.099664000 \\
1.2 & -0.199000000 & -0.198243434 & -0.197800853 & -0.197315517 \\
1.3 & -0.294996246 & -0.293129862 & -0.292110713 & -0.291036003 \\
1.4 & -0.386095345 & -0.382748403 & -0.380986158 & -0.379168221 \\
1.5 & -0.470695388 & -0.465664569 & -0.463078857 & -0.460450590 \\
1.6 & -0.547627491 & -0.540901018 & -0.537503081 & -0.534085626 \\
1.7 & -0.616227665 & -0.607969574 & -0.603849795 & -0.599737720 \\
1.8 & -0.676329533 & -0.666833345 & -0.662136956 & -0.657473792 \\
1.9 & -0.728190908 & -0.717819639 & -0.712718751 & -0.707670533 \\
2.0 & -0.772381768 & -0.761510960 & -0.756179726 & -0.750912371 \\
\bottomrule
\end{tabular}}\bigskip



\item
{\small Euler's method:
\tagpdfsetup{table/header-rows={1},table/header-columns={1}}
\begin{tabular}{S[table-format=1.1] *4{S[table-format=2.9]}  }\toprule
{$x$}&{$h=0.1$}&{$h=0.05$}&{$h=0.025$}&{``Exact''}\\ \midrule
2.0 &  2.000000000 &  2.000000000 &  2.000000000 &  2.000000000 \\
2.1 &  2.420000000 &  2.440610764 &  2.451962006 &  2.464119569 \\
2.2 &  2.922484288 &  2.972198224 &  2.999753046 &  3.029403212 \\
2.3 &  3.524104434 &  3.614025082 &  3.664184099 &  3.718409925 \\
2.4 &  4.244823572 &  4.389380160 &  4.470531822 &  4.558673929 \\
2.5 &  5.108581185 &  5.326426396 &  5.449503467 &  5.583808754 \\
2.6 &  6.144090526 &  6.459226591 &  6.638409411 &  6.834855438 \\
2.7 &  7.385795229 & 7.828984275  & 8.082588076  & 8.361928926 \\
2.8 &  8.875017001 & 9.485544888  & 9.837137672  & 10.226228709 \\
2.9 & 10.661332618 & 11.489211987 & 11.969020902 & 12.502494409 \\
3.0 &12.804226135 & 13.912944662 & 14.559623055 & 15.282004826 \\
\bottomrule
\end{tabular}}\bigskip

{\small Euler semilinear method:
\tagpdfsetup{table/header-rows={1},table/header-columns={1}}
\begin{tabular}{S[table-format=1.1] *4{S[table-format=2.9]}  }\toprule
{$x$}&{$h=0.1$}&{$h=0.05$}&{$h=0.025$}&{``Exact''}\\ \midrule
2.0 &  2.000000000 &  2.000000000 &  2.000000000 &  2.000000000 \\
2.1 &  2.467233571 &  2.465641081 &  2.464871435 &  2.464119569 \\
2.2 &  3.036062650 &  3.032657307 &  3.031011316 &  3.029403212 \\
2.3 &  3.729169725 &  3.723668026 &  3.721008466 &  3.718409925 \\
2.4 &  4.574236356 &  4.566279470 &  4.562432696 &  4.558673929 \\
2.5 &  5.605052990 &  5.594191643 &  5.588940276 &  5.583808754 \\
2.6 &  6.862874116 &  6.848549921 &  6.841623814 &  6.834855438 \\
2.7 &  8.398073101 &  8.379595572 &  8.370660695 &  8.361928926 \\
2.8 & 10.272163096 & 10.248681420 & 10.237326199 & 10.226228709 \\
2.9 & 12.560265110 & 12.530733531 & 12.516452106 & 12.502494409 \\
3.0 & 15.354122287 & 15.317257705 & 15.299429421 & 15.282004826 \\
\hline
\end{tabular}}\bigskip


\item
{\small Euler's method:
\tagpdfsetup{table/header-rows={1},table/header-columns={1}}
\begin{tabular}{S[table-format=1.1] *4{S[table-format=1.9]}  }\toprule
{$x$}&{$h=0.2$}&{$h=0.1$}&{$h=0.05$}&{``Exact''}\\ \midrule
1.0 & 2.000000000 & 2.000000000 & 2.000000000 & 2.000000000 \\
1.2 & 1.768294197 & 1.786514499 & 1.794412375 & 1.801636774 \\
1.4 & 1.603028371 & 1.628427487 & 1.639678822 & 1.650102616 \\
1.6 & 1.474580412 & 1.502563111 & 1.515157063 & 1.526935885 \\
1.8 & 1.368349549 & 1.396853671 & 1.409839229 & 1.422074283 \\
2.0 & 1.276424761 & 1.304504818 & 1.317421794 & 1.329664953 \\
2.2 & 1.194247156 & 1.221490111 & 1.234122458 & 1.246155344 \\
2.4 & 1.119088175 & 1.145348276 & 1.157607418 & 1.169334346 \\
2.6 & 1.049284410 & 1.074553688 & 1.086419453 & 1.097812069 \\
2.8 & 0.983821745 & 1.008162993 & 1.019652023 & 1.030719114 \\
3.0 & 0.922094379 & 0.945604800 & 0.956752868 & 0.967523153 \\
\bottomrule
\end{tabular}}\bigskip

{\small Euler semilinear method:
\tagpdfsetup{table/header-rows={1},table/header-columns={1}}
\begin{tabular}{S[table-format=1.1] *4{S[table-format=1.9]}  }\toprule
{$x$}&{$h=0.2$}&{$h=0.1$}&{$h=0.05$}&{``Exact''}\\ \midrule
1.0 & 2.000000000 & 2.000000000 & 2.000000000 & 2.000000000 \\
1.2 & 1.806911831 & 1.804304958 & 1.802978526 & 1.801636774 \\
1.4 & 1.659738603 & 1.654968381 & 1.652547436 & 1.650102616 \\
1.6 & 1.540257861 & 1.533652916 & 1.530308405 & 1.526935885 \\
1.8 & 1.438532932 & 1.430361800 & 1.426232584 & 1.422074283 \\
2.0 & 1.348782285 & 1.339279577 & 1.334486249 & 1.329664953 \\
2.2 & 1.267497415 & 1.256876924 & 1.251528766 & 1.246155344 \\
2.4 & 1.192497494 & 1.180958765 & 1.175157264 & 1.169334346 \\
2.6 & 1.122416379 & 1.110147777 & 1.103988310 & 1.097812069 \\
2.8 & 1.056405906 & 1.043585743 & 1.037158237 & 1.030719114 \\
3.0 & 0.993954754 & 0.980751307 & 0.974140320 & 0.967523153 \\
\bottomrule
\end{tabular}}\bigskip


\item
{\small Euler's method:
\tagpdfsetup{table/header-rows={1},table/header-columns={1}}
\begin{tabular}{S[table-format=1.1] *4{S[table-format=1.9]}  }\toprule
{$x$}&{$h=0.2$}&{$h=0.1$}&{$h=0.05$}&{``Exact''}\\ \midrule
0.0 & 1.000000000 & 1.000000000 & 1.000000000 & 1.000000000 \\
0.2 & 1.200000000 & 1.186557290 & 1.179206574 & 1.171515153 \\
0.4 & 1.333543409 & 1.298441890 & 1.280865289 & 1.263370891 \\
0.6 & 1.371340142 & 1.319698328 & 1.295082088 & 1.271251278 \\
0.8 & 1.326367357 & 1.270160237 & 1.243958980 & 1.218901287 \\
1.0 & 1.233056306 & 1.181845667 & 1.158064902 & 1.135362070 \\
1.2 & 1.122359136 & 1.080477477 & 1.060871608 & 1.042062625 \\
1.4 & 1.013100262 & 0.981124989 & 0.965917496 & 0.951192532 \\
1.6 & 0.914000211 & 0.890759107 & 0.879460404 & 0.868381328 \\
1.8 & 0.827848558 & 0.811673612 & 0.803582000 & 0.795518627 \\
2.0 & 0.754572560 & 0.743869878 & 0.738303914 & 0.732638628 \\
\bottomrule
\end{tabular}}\bigskip

{\small Euler semilinear method:
\tagpdfsetup{table/header-rows={1},table/header-columns={1}}
\begin{tabular}{S[table-format=1.1] *4{S[table-format=1.9]}  }\toprule
{$x$}&{$h=0.2$}&{$h=0.1$}&{$h=0.05$}&{``Exact''}\\ \midrule
0.0 & 1.000000000 & 1.000000000 & 1.000000000 & 1.000000000 \\
0.2 & 1.153846154 & 1.162906599 & 1.167266650 & 1.171515153 \\
0.4 & 1.236969953 & 1.250608357 & 1.257097924 & 1.263370891 \\
0.6 & 1.244188456 & 1.258241892 & 1.264875987 & 1.271251278 \\
0.8 & 1.195155456 & 1.207524076 & 1.213335781 & 1.218901287 \\
1.0 & 1.115731189 & 1.125966437 & 1.130768614 & 1.135362070 \\
1.2 & 1.025938754 & 1.034336918 & 1.038283392 & 1.042062625 \\
1.4 & 0.937645707 & 0.944681597 & 0.948002346 & 0.951192532 \\
1.6 & 0.856581823 & 0.862684171 & 0.865583126 & 0.868381328 \\
1.8 & 0.784832910 & 0.790331183 & 0.792963532 & 0.795518627 \\
2.0 & 0.722610454 & 0.727742966 & 0.730220211 & 0.732638628 \\
\bottomrule
\end{tabular}}\bigskip

\item
{\small Euler's method:
\tagpdfsetup{table/header-rows={1},table/header-columns={1}}
\begin{tabular}{S[table-format=1.1] *4{S[table-format=1.9]}  }\toprule
{$x$}&{$h=0.1$}&{$h=0.05$}&{$h=0.025$}&{``Exact''}\\ \midrule
0.0 & 1.000000000 & 1.000000000 & 1.000000000 & 1.000000000 \\
0.1 & 0.700000000 & 0.725841563 & 0.736671690 & 0.746418339 \\
0.2 & 0.498330000 & 0.532982493 & 0.547988831 & 0.561742917 \\
0.3 & 0.356272689 & 0.392592562 & 0.408724303 & 0.423724207 \\
0.4 & 0.254555443 & 0.289040639 & 0.304708942 & 0.319467408 \\
0.5 & 0.181440541 & 0.212387189 & 0.226758594 & 0.240464879 \\
0.6 & 0.128953069 & 0.155687255 & 0.168375130 & 0.180626161 \\
0.7 & 0.091393543 & 0.113851516 & 0.124744976 & 0.135394692 \\
0.8 & 0.064613612 & 0.083076641 & 0.092230966 & 0.101293057 \\
0.9 & 0.045585102 & 0.060505907 & 0.068068776 & 0.075650324 \\
1.0 & 0.032105117 & 0.043997045 & 0.050159310 & 0.056415515 \\
\bottomrule
\end{tabular}}\bigskip

{\small Euler semilinear method:
\tagpdfsetup{table/header-rows={1},table/header-columns={1}}
\begin{tabular}{S[table-format=1.1] *4{S[table-format=1.9]}  }\toprule
{$x$}&{$h=0.1$}&{$h=0.05$}&{$h=0.025$}&{``Exact''}\\ \midrule
0.0 & 1.000000000 & 1.000000000 & 1.000000000 & 1.000000000 \\
0.1 & 0.740818221 & 0.743784320 & 0.745143557 & 0.746418339 \\
0.2 & 0.555889275 & 0.558989106 & 0.560410719 & 0.561742917 \\
0.3 & 0.418936461 & 0.421482025 & 0.422642541 & 0.423724207 \\
0.4 & 0.315890439 & 0.317804400 & 0.318668549 & 0.319467408 \\
0.5 & 0.237908421 & 0.239287095 & 0.239902094 & 0.240464879 \\
0.6 & 0.178842206 & 0.179812811 & 0.180239888 & 0.180626161 \\
0.7 & 0.134165506 & 0.134840668 & 0.135133367 & 0.135394692 \\
0.8 & 0.100450939 & 0.100918118 & 0.101117514 & 0.101293057 \\
0.9 & 0.075073968 & 0.075396974 & 0.075532643 & 0.075650324 \\
1.0 & 0.056020154 & 0.056243980 & 0.056336491 & 0.056415515 \\
\bottomrule
\end{tabular}}\bigskip

\item
{\small Euler's method:
\tagpdfsetup{table/header-rows={1},table/header-columns={1}}
\begin{tabular}{S[table-format=1.1] *4{S[table-format=1.9]}  }\toprule
{$x$}&{$h=0.1$}&{$h=0.05$}&{$h=0.025$}&{``Exact''}\\ \midrule
2.0 & 1.000000000 & 1.000000000 & 1.000000000 & 1.000000000 \\
2.1 & 1.000000000 & 1.005062500 & 1.007100815 & 1.008899988 \\
2.2 & 1.020500000 & 1.026752091 & 1.029367367 & 1.031723469 \\
2.3 & 1.053489840 & 1.059067423 & 1.061510137 & 1.063764243 \\
2.4 & 1.093521685 & 1.097780573 & 1.099748225 & 1.101614730 \\
2.5 & 1.137137554 & 1.140059654 & 1.141496651 & 1.142903776 \\
2.6 & 1.182269005 & 1.184090031 & 1.185056276 & 1.186038851 \\
2.7 & 1.227745005 & 1.228755801 & 1.229350441 & 1.229985178 \\
2.8 & 1.272940309 & 1.273399187 & 1.273721920 & 1.274092525 \\
2.9 & 1.317545833 & 1.317651554 & 1.317786528 & 1.317967533 \\
3.0 & 1.361427907 & 1.361320824 & 1.361332589 & 1.361383810 \\
\bottomrule
\end{tabular}}\bigskip

{\small Euler semilinear method:
\tagpdfsetup{table/header-rows={1},table/header-columns={1}}
\begin{tabular}{S[table-format=1.1] *4{S[table-format=1.9]}  }\toprule
{$x$}&{$h=0.1$}&{$h=0.05$}&{$h=0.025$}&{``Exact''}\\ \midrule
2.0 & 1.000000000 & 1.000000000 & 1.000000000 & 1.000000000 \\
2.1 & 0.982476904 & 0.996114142 & 1.002608435 & 1.008899988 \\
2.2 & 0.988105346 & 1.010577663 & 1.021306044 & 1.031723469 \\
2.3 & 1.009495813 & 1.037358814 & 1.050731634 & 1.063764243 \\
2.4 & 1.041012955 & 1.071994816 & 1.086964414 & 1.101614730 \\
2.5 & 1.078631301 & 1.111346365 & 1.127262285 & 1.142903776 \\
2.6 & 1.119632590 & 1.153300133 & 1.169781376 & 1.186038851 \\
2.7 & 1.162270287 & 1.196488725 & 1.213325613 & 1.229985178 \\
2.8 & 1.205472927 & 1.240060456 & 1.257146091 & 1.274092525 \\
2.9 & 1.248613584 & 1.283506001 & 1.300791772 & 1.317967533 \\
3.0 & 1.291345518 & 1.326535737 & 1.344004102 & 1.361383810 \\
\bottomrule
\end{tabular}}\bigskip
\end{sectionSolutions}

\section{The Improved Euler Method and Related Methods}

\begin{sectionSolutions}
\item
$y_1=1.220207973,\ y_2=1.489578775\ y_3=1.819337186$


\item
$y_1=2.961317914$;\ $y_2=2.920132727$;\ $y_3=2.876213748$.



\item
{\small
\tagpdfsetup{table/header-rows={1},table/header-columns={1}}
\begin{tabular}{S[table-format=1.1] *4{S[table-format=2.9]}  }\toprule
{$x$}&{$h=0.1$}&{$h=0.05$}&{$h=0.025$}&{Exact}\\ \midrule
0.0 &  2.000000000 &  2.000000000 &  2.000000000 &   2.000000000 \\
0.1 &  2.257138644 &  2.238455342 &  2.234055168 &   2.232642918 \\
0.2 &  2.826004666 &  2.786634110 &  2.777340360 &   2.774352565 \\
0.3 &  3.812671926 &  3.747167263 &  3.731674025 &   3.726686582 \\
0.4 &  5.387430580 &  5.285996803 &  5.261969043 &   5.254226636 \\
0.5 &  7.813298361 &  7.660199197 &  7.623893064 &   7.612186259 \\
0.6 & 11.489337756 & 11.260349005 & 11.206005869 &  11.188475269 \\
0.7 & 17.015861211 & 16.674352914 & 16.593267820 &  16.567103199 \\
0.8 & 25.292140630 & 24.783149862 & 24.662262731 &  24.623248150 \\
0.9 & 37.662496723 & 36.903828191 & 36.723608928 &  36.665439956 \\
1.0 & 56.134480009 & 55.003390448 & 54.734674836 &  54.647937102 \\
\bottomrule
\end{tabular}}\bigskip

\item
{\small
\tagpdfsetup{table/header-rows={1},table/header-columns={1}}
\begin{tabular}{S[table-format=1.2] *4{S[table-format=2.9]}  }\toprule
{$x$}&{$h=0.05$}&{$h=0.025$}&{$h=0.0125$}&{Exact}\\ \midrule
1.00 & 2.000000000  &  2.000000000 &  2.000000000 &  2.000000000 \\
1.05 & 2.268367347  &  2.269670336 &  2.270030868 &  2.270158103 \\
1.10 & 2.582607299  &  2.585911295 &  2.586827341 &  2.587150838 \\
1.15 & 2.954510022  &  2.960870733 &  2.962638822 &  2.963263785 \\
1.20 & 3.400161788  &  3.411212150 &  3.414293964 &  3.415384615 \\
1.25 & 3.942097142  &  3.960434900 &  3.965570792 &  3.967391304 \\
1.30 & 4.612879780  &  4.642784826 &  4.651206769 &  4.654198473 \\
1.35 & 5.461348619  &  5.510188575 &  5.524044591 &  5.528980892 \\
1.40 & 6.564150753  &  6.645334756 &  6.668600859 &  6.676923077 \\
1.45 & 8.048579617  &  8.188335998 &  8.228972215 &  8.243593315 \\
1.50 &10.141969585  & 10.396770409 & 10.472502111 & 10.500000000 \\
\bottomrule
\end{tabular}}\bigskip

\item
{\small
\tagpdfsetup{table/header-rows={1,2},table/header-columns={1}}
\begin{tabular}{S[table-format=1.1] *3{S[table-format=1.9]} *3{S[table-format=-1.6]} }\toprule
&\multicolumn{3}{c}{Approximate Solutions}&
\multicolumn{3}{c}{Residuals}\\\cmidrule(r){2-4}\cmidrule(l){5-7}
{$x$}&{$h=0.1$}&{$h=0.05$}&{$h=0.025$}&{$h=0.1$}&{$h=0.05$}&{$h=0.025$}\\ \midrule
1.0 & 1.000000000 & 1.000000000 & 1.000000000 &0.00000 & 0.000000 &  0.000000  \\
1.1 & 0.923734663 & 0.923730743 & 0.923730591 &0.00004 & 0.000001 & -0.000001  \\
1.2 & 0.854475600 & 0.854449616 & 0.854444697 &0.00035 & 0.000068 &  0.000015  \\
1.3 & 0.791650344 & 0.791596016 & 0.791584634 &0.00078 & 0.000167 &  0.000039  \\
1.4 & 0.734785779 & 0.734703826 & 0.734686010 &0.00125 & 0.000277 &  0.000065  \\
1.5 & 0.683424095 & 0.683318666 & 0.683295308 &0.00171 & 0.000384 &  0.000091  \\
1.6 & 0.637097057 & 0.636973423 & 0.636945710 &0.00213 & 0.000483 &  0.000115  \\
1.7 & 0.595330359 & 0.595193634 & 0.595162740 &0.00250 & 0.000572 &  0.000137  \\
1.8 & 0.557658422 & 0.557513000 & 0.557479947 &0.00283 & 0.000650 &  0.000156  \\
1.9 & 0.523638939 & 0.523488343 & 0.523453958 &0.00311 & 0.000718 &  0.000173  \\
2.0 & 0.492862999 & 0.492709931 & 0.492674855 &0.00335 & 0.000777 &  0.000187  \\
\bottomrule
\end{tabular}}\bigskip



\item
{\small
\tagpdfsetup{table/header-rows={1},table/header-columns={1}}
\begin{tabular}{S[table-format=1.1] *4{S[table-format=-1.9]}  }\toprule
{$x$}&{$h=0.1$}&{$h=0.05$}&{$h=0.025$}&{``Exact''}\\ \midrule
1.0 &  0.000000000 &  0.000000000 &  0.000000000 &  0.000000000 \\
1.1 & -0.099500000 & -0.099623114 & -0.099653809 & -0.099664000 \\
1.2 & -0.196990313 & -0.197235180 & -0.197295585 & -0.197315517 \\
1.3 & -0.290552949 & -0.290917718 & -0.291006784 & -0.291036003 \\
1.4 & -0.378532718 & -0.379013852 & -0.379130237 & -0.379168221 \\
1.5 & -0.459672297 & -0.460262848 & -0.460404546 & -0.460450590 \\
1.6 & -0.533180153 & -0.533868468 & -0.534032512 & -0.534085626 \\
1.7 & -0.598726853 & -0.599496413 & -0.599678824 & -0.599737720 \\
1.8 & -0.656384109 & -0.657214624 & -0.657410640 & -0.657473792 \\
1.9 & -0.706530934 & -0.707400266 & -0.707604759 & -0.707670533 \\
2.0 & -0.749751364 & -0.750637632 & -0.750845571 & -0.750912371 \\
\bottomrule
\end{tabular}}\bigskip

\item
{\small Improved Euler method:
\tagpdfsetup{table/header-rows={1},table/header-columns={1}}
\begin{tabular}{S[table-format=1.1] *4{S[table-format=2.9]} }\toprule
{$x$}&{$h=0.1$}&{$h=0.05$}&{$h=0.025$}&{``Exact''}\\ \midrule
2.0 &  2.000000000 &  2.000000000 &  2.000000000 &  2.000000000 \\
2.1 &  2.461242144 &  2.463344439 &  2.463918368 &  2.464119569 \\
2.2 &  3.022367633 &  3.027507237 &  3.028911026 &  3.029403212 \\
2.3 &  3.705511610 &  3.714932709 &  3.717507170 &  3.718409925 \\
2.4 &  4.537659565 &  4.553006531 &  4.557202414 &  4.558673929 \\
2.5 &  5.551716960 &  5.575150456 &  5.581560437 &  5.583808754 \\
2.6 &  6.787813853 &  6.822158665 &  6.831558101 &  6.834855438 \\
2.7 &  8.294896222 &  8.343829180 &  8.357227947 &  8.361928926 \\
2.8 & 10.132667135 & 10.200955596 & 10.219663917 & 10.226228709 \\
2.9 & 12.373954732 & 12.467758807 & 12.493470722 & 12.502494409 \\
3.0 & 15.107600968 & 15.234856000 & 15.269755072 & 15.282004826 \\
\bottomrule
\end{tabular}}\bigskip

{\small Improved Euler semilinear method:
\tagpdfsetup{table/header-rows={1},table/header-columns={1}}
\begin{tabular}{S[table-format=1.1] *4{S[table-format=2.9]} }\toprule
{$x$}&{$h=0.1$}&{$h=0.05$}&{$h=0.025$}&{``Exact''}\\ \midrule
2.0 &  2.000000000 &  2.000000000 &  2.000000000 &  2.000000000 \\
2.1 &  2.464261688 &  2.464155139 &  2.464128464 &  2.464119569 \\
2.2 &  3.029706047 &  3.029479005 &  3.029422165 &  3.029403212 \\
2.3 &  3.718897663 &  3.718531995 &  3.718440451 &  3.718409925 \\
2.4 &  4.559377397 &  4.558849990 &  4.558717956 &  4.558673929 \\
2.5 &  5.584766724 &  5.584048510 &  5.583868709 &  5.583808754 \\
2.6 &  6.836116246 &  6.835170986 &  6.834934347 &  6.834855438 \\
2.7 &  8.363552464 &  8.362335253 &  8.362030535 &  8.361928926 \\
2.8 & 10.228288880 & 10.226744312 & 10.226357645 & 10.226228709 \\
2.9 & 12.505082132 & 12.503142042 & 12.502656361 & 12.502494409 \\
3.0 & 15.285231726 & 15.282812424 & 15.282206780 & 15.282004826 \\
\bottomrule
\end{tabular}}\bigskip

\item
{\small Improved Euler method:
\tagpdfsetup{table/header-rows={1},table/header-columns={1}}
\begin{tabular}{S[table-format=1.1] *4{S[table-format=1.9]} }\toprule
{$x$}&{$h=0.2$}&{$h=0.1$}&{$h=0.05$}&{``Exact''}\\ \midrule
1.0 & 2.000000000 & 2.000000000 & 2.000000000 & 2.000000000 \\
1.2 & 1.801514185 & 1.801606135 & 1.801629115 & 1.801636774 \\
1.4 & 1.649911580 & 1.650054870 & 1.650090680 & 1.650102616 \\
1.6 & 1.526711768 & 1.526879870 & 1.526921882 & 1.526935885 \\
1.8 & 1.421841570 & 1.422016119 & 1.422059743 & 1.422074283 \\
2.0 & 1.329441172 & 1.329609020 & 1.329650971 & 1.329664953 \\
2.2 & 1.245953205 & 1.246104819 & 1.246142713 & 1.246155344 \\
2.4 & 1.169162994 & 1.169291515 & 1.169323639 & 1.169334346 \\
2.6 & 1.097677870 & 1.097778523 & 1.097803683 & 1.097812069 \\
2.8 & 1.030626179 & 1.030695880 & 1.030713305 & 1.030719114 \\
3.0 & 0.967473721 & 0.967510790 & 0.967520062 & 0.967523153 \\
\bottomrule
\end{tabular}}\bigskip

{\small Improved Euler semilinear method:
\tagpdfsetup{table/header-rows={1},table/header-columns={1}}
\begin{tabular}{S[table-format=1.1] *4{S[table-format=1.9]} }\toprule
{$x$}&{$h=0.2$}&{$h=0.1$}&{$h=0.05$}&{``Exact''}\\ \midrule
1.0 & 2.000000000 & 2.000000000 & 2.000000000 & 2.000000000 \\
1.2 & 1.801514185 & 1.801606135 & 1.801629115 & 1.801636774 \\
1.4 & 1.649911580 & 1.650054870 & 1.650090680 & 1.650102616 \\
1.6 & 1.526711768 & 1.526879870 & 1.526921882 & 1.526935885 \\
1.8 & 1.421841570 & 1.422016119 & 1.422059743 & 1.422074283 \\
2.0 & 1.329441172 & 1.329609020 & 1.329650971 & 1.329664953 \\
2.2 & 1.245953205 & 1.246104819 & 1.246142713 & 1.246155344 \\
2.4 & 1.169162994 & 1.169291515 & 1.169323639 & 1.169334346 \\
2.6 & 1.097677870 & 1.097778523 & 1.097803683 & 1.097812069 \\
2.8 & 1.030626179 & 1.030695880 & 1.030713305 & 1.030719114 \\
3.0 & 0.967473721 & 0.967510790 & 0.967520062 & 0.967523153 \\
\bottomrule
\end{tabular}}\bigskip

\item
{\small Improved Euler method:
\tagpdfsetup{table/header-rows={1},table/header-columns={1}}
\begin{tabular}{S[table-format=1.1] *4{S[table-format=1.9]} }\toprule
{$x$}&{$h=0.2$}&{$h=0.1$}&{$h=0.05$}&{``Exact''}\\ \midrule
0.0 & 1.000000000 & 1.000000000 & 1.000000000 & 1.000000000 \\
0.2 & 1.166771705 & 1.170394902 & 1.171244037 & 1.171515153 \\
0.4 & 1.255835116 & 1.261642355 & 1.262958788 & 1.263370891 \\
0.6 & 1.263517157 & 1.269528214 & 1.270846761 & 1.271251278 \\
0.8 & 1.212551997 & 1.217531648 & 1.218585457 & 1.218901287 \\
1.0 & 1.130812573 & 1.134420589 & 1.135150284 & 1.135362070 \\
1.2 & 1.039104333 & 1.041487727 & 1.041938536 & 1.042062625 \\
1.4 & 0.949440052 & 0.950888923 & 0.951132561 & 0.951192532 \\
1.6 & 0.867475787 & 0.868263999 & 0.868364849 & 0.868381328 \\
1.8 & 0.795183973 & 0.795523696 & 0.795530315 & 0.795518627 \\
2.0 & 0.732679223 & 0.732721613 & 0.732667905 & 0.732638628 \\
\bottomrule
\end{tabular}}\bigskip

{\small Improved Euler semilinear method:
\tagpdfsetup{table/header-rows={1},table/header-columns={1}}
\begin{tabular}{S[table-format=1.1] *4{S[table-format=1.9]} }\toprule
{$x$}&{$h=0.2$}&{$h=0.1$}&{$h=0.05$}&{``Exact''}\\ \midrule
0.0 & 1.000000000 & 1.000000000 & 1.000000000 & 1.000000000 \\
0.2 & 1.170617859 & 1.171292452 & 1.171459576 & 1.171515153 \\
0.4 & 1.261629934 & 1.262938347 & 1.263262919 & 1.263370891 \\
0.6 & 1.269173253 & 1.270734290 & 1.271122186 & 1.271251278 \\
0.8 & 1.216926014 & 1.218409355 & 1.218778420 & 1.218901287 \\
1.0 & 1.133688235 & 1.134944960 & 1.135257876 & 1.135362070 \\
1.2 & 1.040721691 & 1.041728386 & 1.041979126 & 1.042062625 \\
1.4 & 0.950145706 & 0.950931597 & 0.951127345 & 0.951192532 \\
1.6 & 0.867573431 & 0.868179975 & 0.868331028 & 0.868381328 \\
1.8 & 0.794899034 & 0.795364245 & 0.795480063 & 0.795518627 \\
2.0 & 0.732166678 & 0.732521078 & 0.732609267 & 0.732638628 \\
\bottomrule
\end{tabular}}\bigskip

\item
{\small Improved Euler method:
\tagpdfsetup{table/header-rows={1},table/header-columns={1}}
\begin{tabular}{S[table-format=1.1] *4{S[table-format=1.9]} }\toprule
{$x$}&{$h=0.1$}&{$h=0.05$}&{$h=0.025$}&{``Exact''}\\ \midrule
0.0 & 1.000000000 & 1.000000000 & 1.000000000 & 1.000000000 \\
0.1 & 0.749165000 & 0.747022742 & 0.746561141 & 0.746418339 \\
0.2 & 0.565942699 & 0.562667885 & 0.561961242 & 0.561742917 \\
0.3 & 0.428618351 & 0.424803657 & 0.423978964 & 0.423724207 \\
0.4 & 0.324556426 & 0.320590918 & 0.319732571 & 0.319467408 \\
0.5 & 0.245417735 & 0.241558658 & 0.240723019 & 0.240464879 \\
0.6 & 0.185235654 & 0.181643813 & 0.180866303 & 0.180626161 \\
0.7 & 0.139546094 & 0.136310496 & 0.135610749 & 0.135394692 \\
0.8 & 0.104938506 & 0.102096319 & 0.101482503 & 0.101293057 \\
0.9 & 0.078787731 & 0.076340645 & 0.075813072 & 0.075650324 \\
1.0 & 0.059071894 & 0.056999028 & 0.056553023 & 0.056415515 \\
\bottomrule
\end{tabular}}\bigskip

{\small Improved Euler semilinear method:
\tagpdfsetup{table/header-rows={1},table/header-columns={1}}
\begin{tabular}{S[table-format=1.1] *4{S[table-format=1.9]} }\toprule
{$x$}&{$h=0.1$}&{$h=0.05$}&{$h=0.025$}&{``Exact''}\\ \midrule
0.0 & 1.000000000 & 1.000000000 & 1.000000000 & 1.000000000 \\
0.1 & 0.745595127 & 0.746215164 & 0.746368056 & 0.746418339 \\
0.2 & 0.560827568 & 0.561515647 & 0.561686492 & 0.561742917 \\
0.3 & 0.422922083 & 0.423524585 & 0.423674586 & 0.423724207 \\
0.4 & 0.318820339 & 0.319306259 & 0.319427337 & 0.319467408 \\
0.5 & 0.239962317 & 0.240339716 & 0.240433757 & 0.240464879 \\
0.6 & 0.180243441 & 0.180530866 & 0.180602470 & 0.180626161 \\
0.7 & 0.135106416 & 0.135322934 & 0.135376855 & 0.135394692 \\
0.8 & 0.101077312 & 0.101239368 & 0.101279714 & 0.101293057 \\
0.9 & 0.075489492 & 0.075610310 & 0.075640381 & 0.075650324 \\
1.0 & 0.056295914 & 0.056385765 & 0.056408124 & 0.056415515 \\
\bottomrule
\end{tabular}}\bigskip

\item
{\small Improved Euler method:
\tagpdfsetup{table/header-rows={1},table/header-columns={1}}
\begin{tabular}{S[table-format=1.1] *4{S[table-format=1.9]} }\toprule
{$x$}&{$h=0.1$}&{$h=0.05$}&{$h=0.025$}&{``Exact''}\\ \midrule
2.0 & 1.000000000 & 1.000000000 & 1.000000000 & 1.000000000 \\
2.1 & 1.010250000 & 1.009185754 & 1.008965733 & 1.008899988 \\
2.2 & 1.033547273 & 1.032105322 & 1.031811002 & 1.031723469 \\
2.3 & 1.065562151 & 1.064135919 & 1.063849094 & 1.063764243 \\
2.4 & 1.103145347 & 1.101926450 & 1.101685553 & 1.101614730 \\
2.5 & 1.144085693 & 1.143140125 & 1.142957158 & 1.142903776 \\
2.6 & 1.186878796 & 1.186202854 & 1.186075600 & 1.186038851 \\
2.7 & 1.230530804 & 1.230088035 & 1.230007943 & 1.229985178 \\
2.8 & 1.274404357 & 1.274147657 & 1.274104430 & 1.274092525 \\
2.9 & 1.318104153 & 1.317987551 & 1.317971490 & 1.317967533 \\
3.0 & 1.361395309 & 1.361379259 & 1.361382239 & 1.361383810 \\
\bottomrule
\end{tabular}}\bigskip

{\small Improved Euler semilinear method:
\tagpdfsetup{table/header-rows={1},table/header-columns={1}}
\begin{tabular}{S[table-format=1.1] *4{S[table-format=1.9]} }\toprule
{$x$}&{$h=0.1$}&{$h=0.05$}&{$h=0.025$}&{``Exact''}\\ \midrule
2.0 & 1.000000000 & 1.000000000 & 1.000000000 & 1.000000000 \\
2.1 & 1.012802674 & 1.009822081 & 1.009124116 & 1.008899988 \\
2.2 & 1.038431870 & 1.033307426 & 1.032108359 & 1.031723469 \\
2.3 & 1.072484834 & 1.065821457 & 1.064263950 & 1.063764243 \\
2.4 & 1.111794329 & 1.104013534 & 1.102197168 & 1.101614730 \\
2.5 & 1.154168041 & 1.145554968 & 1.143547198 & 1.142903776 \\
2.6 & 1.198140189 & 1.188883373 & 1.186728849 & 1.186038851 \\
2.7 & 1.242762459 & 1.232984559 & 1.230712361 & 1.229985178 \\
2.8 & 1.287441845 & 1.277221941 & 1.274850828 & 1.274092525 \\
2.9 & 1.331821976 & 1.321210992 & 1.318753047 & 1.317967533 \\
3.0 & 1.375699933 & 1.364730937 & 1.362193997 & 1.361383810 \\
\bottomrule
\end{tabular}}\bigskip

\item
{\small
\tagpdfsetup{table/header-rows={1},table/header-columns={1}}
\begin{tabular}{S[table-format=1.1] *4{S[table-format=1.9]} }\toprule
{$x$}&{$h=0.1$}&{$h=0.05$}&{$h=0.025$}&{Exact}\\ \midrule
1.0 & 1.000000000 & 1.000000000 & 1.000000000 & 1.000000000 \\
1.1 & 1.151019287 & 1.153270661 & 1.153777957 & 1.153937085 \\
1.2 & 1.238798618 & 1.241884421 & 1.242580821 & 1.242799540 \\
1.3 & 1.289296258 & 1.292573128 & 1.293313355 & 1.293546032 \\
1.4 & 1.317686801 & 1.320866599 & 1.321585242 & 1.321811247 \\
1.5 & 1.333073855 & 1.336036248 & 1.336705820 & 1.336916440 \\
1.6 & 1.341027170 & 1.343732006 & 1.344343232 & 1.344535503 \\
1.7 & 1.345001345 & 1.347446389 & 1.347998652 & 1.348172348 \\
1.8 & 1.347155352 & 1.349355473 & 1.349852082 & 1.350008229 \\
1.9 & 1.348839325 & 1.350816158 & 1.351261995 & 1.351402121 \\
2.0 & 1.350890736 & 1.352667599 & 1.353067951 & 1.353193719\\
\bottomrule
\end{tabular}}\bigskip

\item
{\small
\tagpdfsetup{table/header-rows={1},table/header-columns={1}}
\begin{tabular}{S[table-format=1.1] *4{S[table-format=1.9]} }\toprule
{$x$}&{$h=0.05$}&{$h=0.025$}&{$h=0.0125$}&{Exact}\\ \midrule
1.00 &  2.000000000 &  2.000000000 &  2.000000000 &  2.000000000 \\
1.05 &  2.268496358 &  2.269703943 &  2.270043628 &  2.270158103 \\
1.10 &  2.582897367 &  2.585985695 &  2.586855275 &  2.587150838 \\
1.15 &  2.954995034 &  2.960992388 &  2.962683751 &  2.963263785 \\
1.20 &  3.400872342 &  3.411384294 &  3.414355862 &  3.415384615 \\
1.25 &  3.943047906 &  3.960651794 &  3.965644965 &  3.967391304 \\
1.30 &  4.614039436 &  4.643018510 &  4.651277424 &  4.654198473 \\
1.35 &  5.462568051 &  5.510357362 &  5.524069547 &  5.528980892 \\
1.40 &  6.564985580 &  6.645224236 &  6.668472955 &  6.676923077 \\
1.45 &  8.047824947 &  8.187384679 &  8.228413044 &  8.243593315 \\
1.50 & 10.136329642 & 10.393419681 & 10.470731411 & 10.500000000 \\
\bottomrule
\end{tabular}}\bigskip

\item
{\small
\tagpdfsetup{table/header-rows={1},table/header-columns={1}}
\begin{tabular}{S[table-format=1.1] *4{S[table-format=1.9]} }\toprule
{$x$}&{$h=0.1$}&{$h=0.05$}&{$h=0.025$}&{``Exact''}\\ \midrule
0.0 & 1.000000000 & 1.000000000 & 1.000000000 & 1.000000000 \\
0.1 & 0.984142840 & 0.984133302 & 0.984130961 & 0.984130189 \\
0.2 & 0.965066124 & 0.965044455 & 0.965039117 & 0.965037353 \\
0.3 & 0.942648578 & 0.942611457 & 0.942602279 & 0.942599241 \\
0.4 & 0.916705578 & 0.916648569 & 0.916634423 & 0.916629732 \\
0.5 & 0.886970525 & 0.886887464 & 0.886866778 & 0.886859904 \\
0.6 & 0.853066054 & 0.852948011 & 0.852918497 & 0.852908668 \\
0.7 & 0.814458249 & 0.814291679 & 0.814249848 & 0.814235883 \\
0.8 & 0.770380571 & 0.770143777 & 0.770083998 & 0.770063987 \\
0.9 & 0.719699643 & 0.719355385 & 0.719267905 & 0.719238519 \\
1.0 & 0.660658411 & 0.660136630 & 0.660002840 & 0.659957689 \\
\bottomrule
\end{tabular}}\bigskip

\item
{\small
\tagpdfsetup{table/header-rows={1},table/header-columns={1}}
\begin{tabular}{S[table-format=1.1] *4{S[table-format=1.9]} }\toprule
{$x$}&{$h=0.1$}&{$h=0.05$}&{$h=0.025$}&{``Exact''}\\ \midrule
1.0 & 0.000000000 & 0.000000000  &  0.000000000 &  0.000000000 \\
1.1 &-0.099666667 &-0.099665005  & -0.099664307 & -0.099674132 \\
1.2 &-0.197322275 &-0.197317894  & -0.197316222 & -0.197355914 \\
1.3 &-0.291033227 &-0.291036361  & -0.291036258 & -0.291123993 \\
1.4 &-0.379131069 &-0.379160444  & -0.379166504 & -0.379315647 \\
1.5 &-0.460350276 &-0.460427667  & -0.460445166 & -0.460662347 \\
1.6 &-0.533897316 &-0.534041581  & -0.534075026 & -0.534359685 \\
1.7 &-0.599446325 &-0.599668984  & -0.599721072 & -0.600066382 \\
1.8 &-0.657076288 &-0.657379719  & -0.657450947 & -0.657845646 \\
1.9 &-0.707175010 &-0.707553135  & -0.707641993 & -0.708072516 \\
2.0 &-0.750335016 &-0.750775571  & -0.750879100 & -0.751331499 \\
\bottomrule
\end{tabular}}\bigskip

\item
\begin{alist}
\item
Let $x_i=a+ih$, $i=0,1,\dots,n$. If $y$ is the solution of the initial
value problem $y'=f(x)$, $y(a)=0$, then $y(b)=\int_a^bf(x)\,dx$.
The improved  Euler method yields
$y_{i+1}=y_i+.5h\left(f(a+ih)+f(a+(i+1)h)\right)$,
$i=0,1,\dots,n-1$,
where $y_0=a$ and $y_n$ is an approximation to $\dst\int_a^bf(x)\,dx$.
But
\[
y_n=\dst\sum_{i=0}^{n-1}(y_{i+1}-y_i) =.5h\left(f(a)+f(b)\right)
+h\sum_{i=1}^{n-1}f(a+ih).
\]

\stepcounter{enumXi}
\item The local truncation error is a multiple of
$y'''(\tilde x_i)=f''(\tilde x_i)$, where $x_i<\tilde
x_i<x_{i+1}$. Therefore, the quadrature formula is exact
if $f$ is a polynomial of degree  $<2$.

\item
Let
$E(f)=\dst\int_a^bf(x)\,dx-y_n$.  Note that E is linear.
If $f$ is a polynomial of degree $2$, then $f(x)=f_0(x)+K(x-a)^2$
where $\deg(f_0)\le1$. Since
$E(f_0)=0$ from \part{c} and
\begin{align*}
E((x-a)^2)&=\frac{(b-a)^3}{3}-\frac{(b-a)^2h}{2}-h^3\sum_{i=1}^{n-1}i^2\\
&=h^3\left[\frac{n^3}{3}-\frac{n^2}{2}-\frac{n(n-1)(2n-1)}{6}\right]
=-\frac{nh^3}{6}=-\frac{(b-a)h^2}{6},
\end{align*}
$E(f)=-\dfrac{K(b-a)h^2}{6}$; therefore the error is proportional
to  $h^2$.
\end{alist}
\end{sectionSolutions}

\section{The Runge-Kutta Method}

\begin{sectionSolutions}
\item
$y_1=1.221551366$; $y_2=1.492920208$


\item
$ y_1=2.961316248$; $y_2=2.920128958$.

\item
{\small
\tagpdfsetup{table/header-rows={1},table/header-columns={1}}
\begin{tabular}{S[table-format=1.1] *4{S[table-format=2.9]} }\toprule
{$x$}&{$h=0.1$}&{$h=0.05$}&{$h=0.025$}&{Exact}\\ \midrule
0.0 &   2.000000000 &  2.000000000 &  2.000000000 &   2.000000000 \\
0.1 &   2.232752507 &  2.232649573 &  2.232643327 &   2.232642918 \\
0.2 &   2.774582759 &  2.774366625 &  2.774353431 &   2.774352565 \\
0.3 &   3.727068686 &  3.726710028 &  3.726688030 &   3.726686582 \\
0.4 &   5.254817388 &  5.254263005 &  5.254228886 &   5.254226636 \\
0.5 &   7.613077020 &  7.612241222 &  7.612189662 &   7.612186259 \\
0.6 &  11.189806778 & 11.188557546 & 11.188480365 &  11.188475269 \\
0.7 &  16.569088310 & 16.567225975 & 16.567110808 &  16.567103199 \\
0.8 &  24.626206255 & 24.623431201 & 24.623259496 &  24.623248150 \\
0.9 &  36.669848687 & 36.665712858 & 36.665456874 &  36.665439956 \\
1.0 &  54.654509699 & 54.648344019 & 54.647962328 &  54.647937102 \\
\bottomrule
\end{tabular}}\bigskip

\item
{\small
\tagpdfsetup{table/header-rows={1},table/header-columns={1}}
\begin{tabular}{S[table-format=1.2] *4{S[table-format=2.9]} }\toprule
{$x$}&{$h=0.05$}&{$h=0.025$}&{$h=0.0125$}&{Exact}\\ \midrule
1.00 &  2.000000000 &  2.000000000 &  2.000000000 &  2.000000000 \\
1.05 &  2.270153785 &  2.270157806 &  2.270158083 &  2.270158103 \\
1.10 &  2.587139846 &  2.587150083 &  2.587150789 &  2.587150838 \\
1.15 &  2.963242415 &  2.963262317 &  2.963263689 &  2.963263785 \\
1.20 &  3.415346864 &  3.415382020 &  3.415384445 &  3.415384615 \\
1.25 &  3.967327077 &  3.967386886 &  3.967391015 &  3.967391304 \\
1.30 &  4.654089950 &  4.654191000 &  4.654197983 &  4.654198473 \\
1.35 &  5.528794615 &  5.528968045 &  5.528980049 &  5.528980892 \\
1.40 &  6.676590929 &  6.676900116 &  6.676921569 &  6.676923077 \\
1.45 &  8.242960669 &  8.243549415 &  8.243590428 &  8.243593315 \\
1.50 & 10.498658198 & 10.499906266 & 10.499993820 & 10.500000000 \\
\bottomrule
\end{tabular}}\bigskip

\item
{\small
\tagpdfsetup{table/header-rows={1,2},table/header-columns={1}}
\begin{tabular}{ @{}
 S[table-format=1.2]
 *3{S[table-format=1.9]}
 S[table-format=-1.2e-1,exponent-mode=scientific]
 S[table-format=-1.2e-2,exponent-mode=scientific]
 S[table-format=-1.2e-2,exponent-mode=scientific] @{} }\toprule
&\multicolumn{3}{c}{Approximate Solutions}&
\multicolumn{3}{c}{Residuals}\\\cmidrule(r){2-4}\cmidrule(l){5-7}
{$x$}&{$h=0.1$}&{$h=0.05$}&{$h=0.025$}&{$h=0.1$}&{$h=0.05$}&{$h=0.025$}\\ \midrule
1.0 & 1 & 1 & 1 & 0 &  0 &   0\\
1.1 & 0.923730622 & 0.923730677 & 0.923730681 &-0.000000608 & -0.0000000389 &  -0.00000000245 \\
1.2 & 0.854443253 & 0.854443324 & 0.854443328 &-0.000000819 & -0.0000000529 &  -0.00000000335 \\
1.3 & 0.791581155 & 0.791581218 & 0.791581222 &-0.000000753 & -0.0000000495 &  -0.00000000316 \\
1.4 & 0.734680497 & 0.734680538 & 0.734680541 &-0.000000523 & -0.0000000359 &  -0.00000000233 \\
1.5 & 0.683288034 & 0.683288051 & 0.683288052 &-0.000000224 & -0.0000000178 &  -0.00000000122 \\
1.6 & 0.636937046 & 0.636937040 & 0.636937040 & 0.000000079 &  0.0000000006 &  -0.00000000009 \\
1.7 & 0.595153053 & 0.595153029 & 0.595153028 & 0.000000351 &  0.0000000171 &   0.00000000093 \\
1.8 & 0.557469558 & 0.557469522 & 0.557469520 & 0.000000578 &  0.0000000309 &   0.00000000179 \\
1.9 & 0.523443129 & 0.523443084 & 0.523443081 & 0.000000760 &  0.0000000421 &   0.00000000248 \\
2.0 & 0.492663789 & 0.492663738 & 0.492663736 & 0.000000902 &  0.0000000508 &   0.00000000302 \\
\bottomrule
\end{tabular}}\bigskip

\item
{\small
\tagpdfsetup{table/header-rows={1},table/header-columns={1}}
\begin{tabular}{S[table-format=1.1] *4{S[table-format=-1.9]} }\toprule
{$x$}&{$h=0.1$}&{$h=0.05$}&{$h=0.025$}&{``Exact''}\\ \midrule
1.0 &  0.000000000 &  0.000000000 &  0.000000000 &  0.000000000 \\
1.1 & -0.099663901 & -0.099663994 & -0.099664000 & -0.099664000 \\
1.2 & -0.197315322 & -0.197315504 & -0.197315516 & -0.197315517 \\
1.3 & -0.291035700 & -0.291035983 & -0.291036001 & -0.291036003 \\
1.4 & -0.379167790 & -0.379168194 & -0.379168220 & -0.379168221 \\
1.5 & -0.460450005 & -0.460450552 & -0.460450587 & -0.460450590 \\
1.6 & -0.534084875 & -0.534085579 & -0.534085623 & -0.534085626 \\
1.7 & -0.599736802 & -0.599737663 & -0.599737717 & -0.599737720 \\
1.8 & -0.657472724 & -0.657473726 & -0.657473788 & -0.657473792 \\
1.9 & -0.707669346 & -0.707670460 & -0.707670529 & -0.707670533 \\
2.0 & -0.750911103 & -0.750912294 & -0.750912367 & -0.750912371 \\
\bottomrule
\end{tabular}}\bigskip

\item
{\small Runge-Kutta method:
\tagpdfsetup{table/header-rows={1},table/header-columns={1}}
\begin{tabular}{S[table-format=1.1] *4{S[table-format=2.9]} }\toprule
{$x$}&{$h=0.1$}&{$h=0.05$}&{$h=0.025$}&{``Exact''}\\ \midrule
2.0 &  2.000000000 &  2.000000000 &  2.000000000 &  2.000000000 \\
2.1 &  2.464113907 &  2.464119185 &  2.464119544 &  2.464119569 \\
2.2 &  3.029389360 &  3.029402271 &  3.029403150 &  3.029403212 \\
2.3 &  3.718384519 &  3.718408199 &  3.718409812 &  3.718409925 \\
2.4 &  4.558632516 &  4.558671116 &  4.558673746 &  4.558673929 \\
2.5 &  5.583745479 &  5.583804456 &  5.583808474 &  5.583808754 \\
2.6 &  6.834762639 &  6.834849135 &  6.834855028 &  6.834855438 \\
2.7 &  8.361796619 &  8.361919939 &  8.361928340 &  8.361928926 \\
2.8 & 10.226043942 & 10.226216159 & 10.226227891 & 10.226228709 \\
2.9 & 12.502240429 & 12.502477158 & 12.502493285 & 12.502494409 \\
3.0 & 15.281660036 & 15.281981407 & 15.282003300 & 15.282004826 \\
\bottomrule
\end{tabular}}\bigskip

{\small Runge-Kutta semilinear method:
\tagpdfsetup{table/header-rows={1},table/header-columns={1}}
\begin{tabular}{S[table-format=1.1] *4{S[table-format=2.9]} }\toprule
{$x$}&{$h=0.1$}&{$h=0.05$}&{$h=0.025$}&{``Exact''}\\ \midrule
2.0 &  2.000000000 &  2.000000000 &  2.000000000 &  2.000000000 \\
2.1 &  2.464119623 &  2.464119573 &  2.464119570 &  2.464119569 \\
2.2 &  3.029403325 &  3.029403219 &  3.029403212 &  3.029403212 \\
2.3 &  3.718410105 &  3.718409936 &  3.718409925 &  3.718409925 \\
2.4 &  4.558674188 &  4.558673945 &  4.558673930 &  4.558673929 \\
2.5 &  5.583809105 &  5.583808776 &  5.583808755 &  5.583808754 \\
2.6 &  6.834855899 &  6.834855467 &  6.834855440 &  6.834855438 \\
2.7 &  8.361929516 &  8.361928963 &  8.361928928 &  8.361928926 \\
2.8 & 10.226229456 & 10.226228756 & 10.226228712 & 10.226228709 \\
2.9 & 12.502495345 & 12.502494468 & 12.502494413 & 12.502494409 \\
3.0 & 15.282005990 & 15.282004899 & 15.282004831 & 15.282004826 \\
\bottomrule
\end{tabular}}\bigskip

\item
{\small Runge-Kutta method:
\tagpdfsetup{table/header-rows={1},table/header-columns={1}}
\begin{tabular}{S[table-format=1.1] *4{S[table-format=1.9]} }\toprule
{$x$}&{$h=0.2$}&{$h=0.1$}&{$h=0.05$}&{``Exact''}\\ \midrule
1.0 & 2.000000000 & 2.000000000 & 2.000000000 & 2.000000000 \\
1.2 & 1.801636785 & 1.801636775 & 1.801636774 & 1.801636774 \\
1.4 & 1.650102633 & 1.650102617 & 1.650102616 & 1.650102616 \\
1.6 & 1.526935904 & 1.526935886 & 1.526935885 & 1.526935885 \\
1.8 & 1.422074302 & 1.422074284 & 1.422074283 & 1.422074283 \\
2.0 & 1.329664970 & 1.329664954 & 1.329664953 & 1.329664953 \\
2.2 & 1.246155357 & 1.246155345 & 1.246155344 & 1.246155344 \\
2.4 & 1.169334355 & 1.169334347 & 1.169334346 & 1.169334346 \\
2.6 & 1.097812074 & 1.097812070 & 1.097812069 & 1.097812069 \\
2.8 & 1.030719113 & 1.030719114 & 1.030719114 & 1.030719114 \\
3.0 & 0.967523147 & 0.967523152 & 0.967523153 & 0.967523153 \\
\bottomrule
\end{tabular}}\bigskip

{\small Runge-Kutta semilinear method:
\tagpdfsetup{table/header-rows={1},table/header-columns={1}}
\begin{tabular}{S[table-format=1.1] *4{S[table-format=1.9]} }\toprule
{$x$}&{$h=0.2$}&{$h=0.1$}&{$h=0.05$}&{``Exact''}\\ \midrule
1.0 & 2.000000000 & 2.000000000 & 2.000000000 & 2.000000000 \\
1.2 & 1.801636785 & 1.801636775 & 1.801636774 & 1.801636774 \\
1.4 & 1.650102633 & 1.650102617 & 1.650102616 & 1.650102616 \\
1.6 & 1.526935904 & 1.526935886 & 1.526935885 & 1.526935885 \\
1.8 & 1.422074302 & 1.422074284 & 1.422074283 & 1.422074283 \\
2.0 & 1.329664970 & 1.329664954 & 1.329664953 & 1.329664953 \\
2.2 & 1.246155357 & 1.246155345 & 1.246155344 & 1.246155344 \\
2.4 & 1.169334355 & 1.169334347 & 1.169334346 & 1.169334346 \\
2.6 & 1.097812074 & 1.097812070 & 1.097812069 & 1.097812069 \\
2.8 & 1.030719113 & 1.030719114 & 1.030719114 & 1.030719114 \\
3.0 & 0.967523147 & 0.967523152 & 0.967523153 & 0.967523153 \\
\bottomrule
\end{tabular}}\bigskip

\item
{\small Runge-Kutta method:
\tagpdfsetup{table/header-rows={1},table/header-columns={1}}
\begin{tabular}{S[table-format=1.1] *4{S[table-format=1.9]} }\toprule
{$x$}&{$h=0.2$}&{$h=0.1$}&{$h=0.05$}&{``Exact''}\\ \midrule
0.0 & 1.000000000 & 1.000000000 & 1.000000000 & 1.000000000 \\
0.2 & 1.171515610 & 1.171515156 & 1.171515152 & 1.171515153 \\
0.4 & 1.263365845 & 1.263370556 & 1.263370869 & 1.263370891 \\
0.6 & 1.271238957 & 1.271250529 & 1.271251232 & 1.271251278 \\
0.8 & 1.218885528 & 1.218900353 & 1.218901230 & 1.218901287 \\
1.0 & 1.135346772 & 1.135361174 & 1.135362016 & 1.135362070 \\
1.2 & 1.042049558 & 1.042061864 & 1.042062579 & 1.042062625 \\
1.4 & 0.951181964 & 0.951191920 & 0.951192495 & 0.951192532 \\
1.6 & 0.868372923 & 0.868380842 & 0.868381298 & 0.868381328 \\
1.8 & 0.795511927 & 0.795518241 & 0.795518603 & 0.795518627 \\
2.0 & 0.732633229 & 0.732638318 & 0.732638609 & 0.732638628 \\
\bottomrule
\end{tabular}}\bigskip

{\small Runge-Kutta semilinear method:
\tagpdfsetup{table/header-rows={1},table/header-columns={1}}
\begin{tabular}{S[table-format=1.1] *4{S[table-format=1.9]} }\toprule
{$x$}&{$h=0.2$}&{$h=0.1$}&{$h=0.05$}&{``Exact''}\\ \midrule
0.0 & 1.000000000 & 1.000000000 & 1.000000000 & 1.000000000 \\
0.2 & 1.171517316 & 1.171515284 & 1.171515161 & 1.171515153 \\
0.4 & 1.263374485 & 1.263371110 & 1.263370904 & 1.263370891 \\
0.6 & 1.271254636 & 1.271251485 & 1.271251291 & 1.271251278 \\
0.8 & 1.218903802 & 1.218901442 & 1.218901297 & 1.218901287 \\
1.0 & 1.135363869 & 1.135362181 & 1.135362077 & 1.135362070 \\
1.2 & 1.042063952 & 1.042062706 & 1.042062630 & 1.042062625 \\
1.4 & 0.951193560 & 0.951192595 & 0.951192536 & 0.951192532 \\
1.6 & 0.868382157 & 0.868381378 & 0.868381331 & 0.868381328 \\
1.8 & 0.795519315 & 0.795518669 & 0.795518629 & 0.795518627 \\
2.0 & 0.732639212 & 0.732638663 & 0.732638630 & 0.732638628 \\
\bottomrule
\end{tabular}}\bigskip

\item
{\small Runge-Kutta method:
\tagpdfsetup{table/header-rows={1},table/header-columns={1}}
\begin{tabular}{S[table-format=1.1] *4{S[table-format=1.9]} }\toprule
{$x$}&{$h=0.1$}&{$h=0.05$}&{$h=0.025$}&{``Exact''}\\ \midrule
0.0 & 1.000000000 & 1.000000000 & 1.000000000 & 1.000000000 \\
0.1 & 0.746430962 & 0.746418992 & 0.746418376 & 0.746418339 \\
0.2 & 0.561761987 & 0.561743921 & 0.561742975 & 0.561742917 \\
0.3 & 0.423746057 & 0.423725371 & 0.423724274 & 0.423724207 \\
0.4 & 0.319489811 & 0.319468612 & 0.319467478 & 0.319467408 \\
0.5 & 0.240486460 & 0.240466046 & 0.240464947 & 0.240464879 \\
0.6 & 0.180646105 & 0.180627244 & 0.180626225 & 0.180626161 \\
0.7 & 0.135412569 & 0.135395665 & 0.135394749 & 0.135394692 \\
0.8 & 0.101308709 & 0.101293911 & 0.101293107 & 0.101293057 \\
0.9 & 0.075663769 & 0.075651059 & 0.075650367 & 0.075650324 \\
1.0 & 0.056426886 & 0.056416137 & 0.056415552 & 0.056415515 \\
\bottomrule
\end{tabular}}\bigskip

{\small Runge-Kutta semilinear method:
\tagpdfsetup{table/header-rows={1},table/header-columns={1}}
\begin{tabular}{S[table-format=1.1] *4{S[table-format=1.9]} }\toprule
{$x$}&{$h=0.1$}&{$h=0.05$}&{$h=0.025$}&{``Exact''}\\ \midrule
0.0 & 1.000000000 & 1.000000000 & 1.000000000 & 1.000000000 \\
0.1 & 0.746416306 & 0.746418217 & 0.746418332 & 0.746418339 \\
0.2 & 0.561740647 & 0.561742780 & 0.561742908 & 0.561742917 \\
0.3 & 0.423722193 & 0.423724084 & 0.423724199 & 0.423724207 \\
0.4 & 0.319465760 & 0.319467308 & 0.319467402 & 0.319467408 \\
0.5 & 0.240463579 & 0.240464800 & 0.240464874 & 0.240464879 \\
0.6 & 0.180625156 & 0.180626100 & 0.180626158 & 0.180626161 \\
0.7 & 0.135393924 & 0.135394645 & 0.135394689 & 0.135394692 \\
0.8 & 0.101292474 & 0.101293021 & 0.101293055 & 0.101293057 \\
0.9 & 0.075649884 & 0.075650297 & 0.075650322 & 0.075650324 \\
1.0 & 0.056415185 & 0.056415495 & 0.056415514 & 0.056415515 \\
\bottomrule
\end{tabular}}\bigskip

\item
{\small Runge-Kutta method:
\tagpdfsetup{table/header-rows={1},table/header-columns={1}}
\begin{tabular}{S[table-format=1.1] *4{S[table-format=1.9]} }\toprule
{$x$}&{$h=0.1$}&{$h=0.05$}&{$h=0.025$}&{``Exact''}\\ \midrule
2.0 & 1.000000000 & 1.000000000 & 1.000000000 & 1.000000000 \\
2.1 & 1.008912398 & 1.008900636 & 1.008900025 & 1.008899988 \\
2.2 & 1.031740789 & 1.031724368 & 1.031723520 & 1.031723469 \\
2.3 & 1.063781819 & 1.063765150 & 1.063764295 & 1.063764243 \\
2.4 & 1.101630085 & 1.101615517 & 1.101614774 & 1.101614730 \\
2.5 & 1.142915917 & 1.142904393 & 1.142903811 & 1.142903776 \\
2.6 & 1.186047678 & 1.186039295 & 1.186038876 & 1.186038851 \\
2.7 & 1.229991054 & 1.229985469 & 1.229985194 & 1.229985178 \\
2.8 & 1.274095992 & 1.274092692 & 1.274092535 & 1.274092525 \\
2.9 & 1.317969153 & 1.317967605 & 1.317967537 & 1.317967533 \\
3.0 & 1.361384082 & 1.361383812 & 1.361383809 & 1.361383810 \\
\bottomrule
\end{tabular}}\bigskip

{\small Runge-Kutta semilinear method:
\tagpdfsetup{table/header-rows={1},table/header-columns={1}}
\begin{tabular}{S[table-format=1.1] *4{S[table-format=1.9]} }\toprule
{$x$}&{$h=0.1$}&{$h=0.05$}&{$h=0.025$}&{``Exact''}\\ \midrule
2.0 & 1.000000000 & 1.000000000 & 1.000000000 & 1.000000000 \\
2.1 & 1.008913934 & 1.008900843 & 1.008900041 & 1.008899988 \\
2.2 & 1.031748526 & 1.031725001 & 1.031723564 & 1.031723469 \\
2.3 & 1.063798300 & 1.063766321 & 1.063764371 & 1.063764243 \\
2.4 & 1.101656264 & 1.101617259 & 1.101614886 & 1.101614730 \\
2.5 & 1.142951721 & 1.142906691 & 1.142903955 & 1.142903776 \\
2.6 & 1.186092475 & 1.186042105 & 1.186039051 & 1.186038851 \\
2.7 & 1.230043983 & 1.229988742 & 1.229985397 & 1.229985178 \\
2.8 & 1.274156172 & 1.274096377 & 1.274092762 & 1.274092525 \\
2.9 & 1.318035787 & 1.317971658 & 1.317967787 & 1.317967533 \\
3.0 & 1.361456502 & 1.361388196 & 1.361384079 & 1.361383810 \\
\bottomrule
\end{tabular}}\bigskip

\item
{\small
\tagpdfsetup{table/header-rows={1},table/header-columns={1}}
\begin{tabular}{S[table-format=1.1] *4{S[table-format=1.9]} }\toprule
{$x$}&{$h=0.1$}&{$h=0.05$}&{$h=0.025$}&{Exact}\\ \midrule
1.00  &  0.142854841 &  0.142857001 &  0.142857134 &  0.142857143 \\
1.10  &  0.053340745 &  0.053341989 &  0.053342066 &  0.053342071 \\
1.20  & -0.046154629 & -0.046153895 & -0.046153849 & -0.046153846 \\
1.30  & -0.153363206 & -0.153362764 & -0.153362736 & -0.153362734 \\
1.40  & -0.266397049 & -0.266396779 & -0.266396762 & -0.266396761 \\
1.50  & -0.383721107 & -0.383720941 & -0.383720931 & -0.383720930 \\
1.60  & -0.504109696 & -0.504109596 & -0.504109589 & -0.504109589 \\
1.70  & -0.626598326 & -0.626598268 & -0.626598264 & -0.626598264 \\
1.80  & -0.750437351 & -0.750437320 & -0.750437318 & -0.750437318 \\
1.90  & -0.875050587 & -0.875050574 & -0.875050573 & -0.875050573 \\
2.00  & -1.000000000 & -1.000000000 & -1.000000000 & -1.000000000 \\
\bottomrule
\end{tabular}}\bigskip

\item
{\small
\tagpdfsetup{table/header-rows={1},table/header-columns={1}}
\begin{tabular}{S[table-format=1.1] *4{S[table-format=1.9]} }\toprule
{$x$}&{$h=0.1$}&{$h=0.05$}&{$h=0.025$}&{Exact}\\ \midrule
 0.50  &  -8.954103230 & -8.954063245 &  -8.954060698 & -8.954060528 \\
 0.60  &  -5.059648314 & -5.059633293 &  -5.059632341 & -5.059632277 \\
 0.70  &  -2.516755942 & -2.516749850 &  -2.516749465 & -2.516749439 \\
 0.80  &  -0.752508672 & -0.752506238 &  -0.752506084 & -0.752506074 \\
 0.90  &   0.530528482 &  0.530529270 &   0.530529319 &  0.530529323 \\
 1.00  &   1.500000000 &  1.500000000 &   1.500000000 &  1.500000000 \\
 1.10  &   2.256519743 &  2.256519352 &   2.256519328 &  2.256519326 \\
 1.20  &   2.863543039 &  2.863542454 &   2.863542417 &  2.863542415 \\
 1.30  &   3.362731379 &  3.362730700 &   3.362730658 &  3.362730655 \\
 1.40  &   3.782361948 &  3.782361231 &   3.782361186 &  3.782361183 \\
 1.50  &   4.142171279 &  4.142170553 &   4.142170508 &  4.142170505 \\
\bottomrule
\end{tabular}}\bigskip

\item
\begin{alist}
\item
Let $x_i=a+ih$, $i=0,1,\dots,n$. If $y$ is the solution of the initial
value problem $y'=f(x)$, $y(a)=0$, then $y(b)=\int_a^bf(x)\,dx$.
The Runge-Kutta method yields
$y_{i+1}=\dst
y_i+\frac{h}{6}\left(f(a+ih)+4f(a+(2i+1)h/2)+f(a+(i+1)h)\right)$,
$i=0,1,\dots,n-1$,
where $y_0=a$ and $y_n$ is an approximation to $\dst\int_a^bf(x)\,dx$.
But
\[
y_n=\dst\sum_{i=0}^{n-1}(y_{i+1}-y_i) =
\frac{h}{6}(f(a)+f(b))+
\frac{h}{3}\sum_{i=1}^{n-1}f(a+ih)+\frac{2h}{3}\sum_{i=1}^n
f\left(a+(2i-1)h/2\right).
\]

\stepcounter{enumXi}
\item The local truncation error is a multiple of
$y^{(5)}(\tilde x_i)=f^{(4)}(\tilde x_i)$, where $x_i<\tilde
x_i<x_{i+1}$. Therefore, the quadrature formula is exact
if $f$ is a polynomial of degree $<4$.

\item
Let
$E(f)=\dst\int_a^bf(x)\,dx-y_n$. Note that $E$ is linear.
If $f$ is a polynomial of degree $4$, then $f(x)=f_0(x)+K(x-a)^4$
where $\deg(f_0)\le3$ and $K$ is constant. Since
$E(f_0)=0$ from \part{c} and
\begin{align*}
E((x-a)^4)&=\frac{(b-a)^5}{5}-\frac{(b-a)^4h}{6}-\frac{h^5}{3}\sum_{i=1}^{n-1}
i^4-\frac{2h^5}{3}\sum_{i=1}^n(i-1/2)^4\\
&=h^5\left[\frac{n^5}{5}-\frac{n^4}{6}-\left(\frac{n^5}{15}
-\frac{n^4}{6}+\frac{n^3}{9}-\frac{n}{90}\right)-\left(
\frac{2n^5}{15}-\frac{n^3}{9}+\frac{7n}{360}\right)\right]\\
&=-\frac{nh^5}{120}=-\frac{(b-a)h^4}{120},
\end{align*}
$E(f)=-\dfrac{(b-a)h^4}{120}$; thus, the error is proportional to
$h^4$.
\end{alist}
\end{sectionSolutions}

\chapter{Applications of First Order Equations}

\section{Growth and Decay}

\begin{sectionSolutions}
\item
$k\tau=\ln2$ and $\tau=2\Rightarrow k=\dst\frac{\ln2}{2}$;
$Q(t)=Q_0e^{-t\ln2/2}$; if $Q(T)=\dst\frac{Q_0}{10}$, then
$\dst\frac{Q_0}{10}=Q_0e^{-T\ln2/2}$; $\ln10=\dst\frac{T\ln2}{2}$;
$T=\dst\frac{2\ln10}{\ln2}$ days.


\item
Let $t_1$ be the elapsed time since the tree died.
Since $p(t)=e^{-(t\ln2)\tau}$, it follows that
$p_1=p_0e^{-(t_1\ln2)/\tau}$, so $\dst{\ln\left(\frac{p_1}{ p_0}\right)}
=-\dst\frac{t_1}{\tau}\ln2$ and
 $t_1=\dst{\tau\frac{\ln(p_0/p_1)}{\ln2}}$.




\item
$Q=Q_0e^{-kt}$;
$Q_1=Q_0e^{-kt_1}$;
$Q_2=Q_0e^{-kt_2}$; $\dst\frac{Q_2}{ Q_1}=e^{-k(t_2-t_1)}$;
$\dst{\ln\left(\frac{Q_1}{ Q_2}\right)}=k(t_2-t_1)$; $k=\dst{\frac{1}{
t_2-t_1}\ln\left(\frac{Q_1}{ Q_2}\right)}$.


\item
$Q'=.06Q,\ Q(0)=Q_0$;\ $Q=Q_0e^{.06t}$. We must find $\tau$
such that $Q(\tau)=2Q_0$; that is, $Q_0e^{.06\tau}=2Q_0$, so
$.06\tau=\ln2$ and $\tau=\dst\frac{\ln2}{.06} =\dst\frac{50\ln2}{3}$ yr.


\item
\begin{alist}
\item If $T$ is the time to triple the value, then
$Q(T)=Q_0e^{.05T}=3Q_0$, so $e^{.05T}=3$.
Therefore, $.05T=\ln3$ and $T=20\ln3$.

\item If $Q(10)=100000$, then  $Q_0e^{.5}=100000$,
so $Q_0=100000e^{-.5}$
\end{alist}

\item
$Q'=-\dst\frac{Q^2}{2},\ Q(0)=50$;\, $\dst\frac{Q'}{
Q^2}=-\dst\frac{1}{2}$;\,$-\dst\frac{1}{ Q}=-\dst\frac{t}{2}+c$;\,
$Q(0)=50\Rightarrow c=-\dst\frac{1}{50}$;\, $\dst\frac{1}{
Q}=\dst\frac{t}{2}+\dst\frac{1}{50}=\dst\frac{1+25t}{50}$;\,
$Q=\dst\frac{50}{1+25t}$. Now $Q(T)=25\Rightarrow1+25T=2\Rightarrow
25T=1\Rightarrow T=\dst\frac{1}{25}$ years.


\item
Since $\tau=1500$, $k=\dst\frac{\ln2}{1500}$; hence
$Q=Q_0e^{-(t\ln2)/1500}$. If $Q(t_1)=\dst\frac{3Q_0}{4}$, then
$e^{-(t_1\ln2)/1500}=\dst\frac{3}{4}$;\,
$-t_1\dst\frac{\ln2}{1500}=\ln\dst{\left(\frac{3}{4}\right)}=
-\ln\dst{\left(\frac{4}{3}\right)}$;\,
$t_1=1500\dfrac{\ln\left(\frac{4}{3}\right)}{\ln2}$. Finally,
$Q(2000)=Q_0e^{-\frac{4}{3}\ln2}=2^{-4/3}Q_0$.

\item
(A) $S'=1-\dst\frac{S}{10},\ S(0)=20$. Rewrite the
differential equation in (A) as (B) $S'+\dst\frac{S}{10}=1$. Since
$S_1=e^{-t/10}$ is a solution of the complementary equation, the
solutions of (B) are given by $S=ue^{-t/10}$, where $u'e^{-t/10}=1$.
Therefore, $u'=e^{t/10}$;\ $u=10e^{t/10}+c$;\;
$S=10+ce^{-t/10}$. Now $S(0)=20\Rightarrow c=10$, so
$S=10+10e^{-t/10}$ and
 $\lim_{t\to\infty}S(t)=10$ g.


\item
 (A) $V'=-750+\dst\frac{V}{20},\ V(0)=25000$. Rewrite the
differential equation in (A) as (B) $V'-\dst\frac{V}{20}=-750$. Since
$V_1=e^{t/20}$ is a solution of the complementary equation, the
solutions of (B) are given by $V=ue^{t/20}$, where $u'e^{t/20}=-750$.
Therefore, $u'=-750e^{-t/20}$;\ $u=15000e^{-t/20}+c$;\;
$V=15000+ce^{t/20}$;\ $V(0)=25000\Rightarrow c=10000$. Therefore,
$V=15000+10000e^{t/20}$.


\item
$p'=\dst\frac{p}{2}-\dst\frac{p^2}{8}=-\dst\frac{1}{8}p(p-4)$;\,
$\dst\frac{p'}{ p(p-4)}=-\dst\frac{1}{8}$;\, $\dst{\frac{1}{4}\left[
\frac{1}{ p-4}-\frac{1}{ p}\right]}p'=-\dst\frac{1}{8}$;\,
 $\dst{\left[
\frac{1}{ p-4}-\frac{1}{ p}\right]}p'=-\dst\frac{1}{2}$;\,
$\dst{\left|\frac{p-4}{ p}\right]}=-\dst\frac{t}{2}+k$;\,
$\dst\frac{p-4}{ p}=ce^{-t/2}$;\, $p(0)=100\Rightarrow
c=\dst\frac{24}{25}$;\, $\dst\frac{p-4}{ p}=\dst\frac{24}{25}e^{-t/2}$;\,
$p-4=\dst\frac{24}{25}pe^{-t/2}$;\,
$p\dst{\left(1-\frac{24}{25}e^{-t/2}\right)}=4$;\,
$p=\dst\frac{4}{1-\frac{24}{25}e^{-t/2}}=\dst\frac{100}{24-24e^{-t/2}}$.


\item
\begin{alist}
\item $P'=rP-12M$.

\item
$P=ue^{rt}$; $u'e^{rt}=-12M$; $u'=-12Me^{rt}$;
$u=\dst\frac{12M}{ r}e^{-rt}+c$; $P=\dst\frac{12M}{ r}+ce^{rt}$;
$P(0)=P_0\Rightarrow c=P_0-\dst\frac{12M}{ r}$;
$\dst{P=\frac{12M}{ r}(1-e^{rt})+P_0e^{rt}}$.

\item Since $P(N)=0$, the answer to \part{b} implies that
 $\dst{M=\frac{rP_0}{ 12(1-e^{-rN})}}$
\end{alist}

\item
The researcher's salary is the solution of the initial value problem
$S'=aS,\ S(0)=S_0$. Therefore, $S=S_0e^{at}$. If $P=P(t)$ is the value
of the trust fund, then $P'=-S_0e^{at}+rP$, or $P'-rP=-S_0e^{at}$.
Therefore, (A) $P=ue^{rt}$, where $u'e^{rt}=-S_0e^{at}$,
 so (B) $u'=-S_0e^{(a-r)t}$. If $a\ne r$, then (B) implies that
$u=\dst\frac{S_0}{ r-a}e^{(a-r)t}+c$, so (A) implies that
$P=\dst\frac{S_0}{ r-a}e^{at}+ce^{rt}$. Now $P(0)=P_0\Rightarrow
c=P_0-\dst\frac{S_0}{ r-a}$; therefore
$P=\dst\frac{S_0}{ r-a}e^{at}+\left(P_0-\frac{S_0}{ r-a}\right)e^{rt}$.
We must choose $P_0$ so that $P(T)=0$; that is,
$P=\dst\frac{S_0}{ r-a}e^{aT}+\left(P_0-\frac{S_0}{ r-a}\right)e^{rT}=0$.
Solving this for $P_0$ yields
$P_0=\dfrac{S_0(1-e^{(a-r)T})}{ r-a}$.
If $a=r$, then (B) becomes
$u'=-S_0$, so $u=-S_0t+c$ and (A) implies that
$P=(-S_0t+c)e^{rt}$. Now $P(0)=P_0\Rightarrow
c=P_0$; therefore
$P=(-S_0t+P_0)e^{rt}$.
To make $P(T)=0$ we must take $P_0=S_0T$.



\item
$Q'=\dst\frac{at}{1+btQ^2}-kQ$;\  $\lim_{t\to\infty}Q(t)=(a/bk)^{1/3}$.
\end{sectionSolutions}

\section{Cooling and Mixing}

\begin{sectionSolutions}
\item
Since $T_0=100$ and $T_M=-10$, $T=-10+110e^{-kt}$. Now
$T(1)=80\Rightarrow 80=-10+110e^{-k}$, so $e^{-k}=\dst\frac{9}{11}$
and $k=\ln\dst\frac{11}{9}$. Therefore, $T=-10+110e^{-t\ln\frac{11}{9}}$.


\item
Let $T$ be the thermometer reading.
Since $T_0=212$ and $T_M=70$, $T=70+142e^{-kt}$. Now
$T(2)=125\Rightarrow 125=70+142e^{-2k}$, so
$e^{-2k}=\dst\frac{55}{142}$
and $k=\dst\frac{1}{2}\ln\dst\frac{142}{55}$. Therefore,
(A) $T=70+142e^{-\frac{t}{2}\ln\frac{142}{55}}$.

\begin{alist}
\item
$T(2)=70+142e^{-2\ln\frac{142}{55}}=70+142\dst{\left(\frac{55}{142}\right)^2}
\approx91.30^\circ$F.

\item Let $\tau$ be the time when
 $T(\tau)=72$, so $72=70+142e^{-\frac{\tau}{2}\ln\frac{142}{55}}$, or
$e^{-\frac{\tau}{2}\ln\frac{142}{55}}=\dst\frac{1}{71}$. Therefore,
$\tau=2\dst\frac{\ln71}{\ln\frac{142}{55}}\approx\qty{8.99}{\minute}$.

\item Since (A) implies that $T>70$ for all $t>0$, the
thermometer will never read $\qty{69}{\degreeFahrenheit}$.
\end{alist}

\item
Since $T_M=20$, $T=20+(T_0-20)e^{-kt}$. Now
$T_0-5=20+(T_0-20)e^{-4k}$  and $T_0-7=20+(T_0-20)e^{-8k}$. Therefore,
$\dst\frac{T_0-25}{ T_0-20}=e^{-4k}$ and $\dst\left(\frac{T_0-27}{
T_0-20}\right)=e^{-8k}$, so  $\dst\frac{T_0-27}{
T_0-20}=\dst\left(\frac{T_0-25}{ T_0-20}\right)^2$,
which implies that $(T_0-20)(T_0-27)=(T_0-25)^2$, or
$T_0^2-47T_0+540=T_0^2-50T_0+625$; hence $3T_0=85$
and  $T_0=(85/3)\unit{\degreeCelsius}$.



\item
$Q'=3-\dst\frac{3}{40}Q,\quad Q(0)=0$. Rewrite the differential equation
as (A) $Q'+\dst\frac{3}{40}Q=3$. Since $Q_1=e^{-3t/40}$ is a solution of
the complementary equation, the solutions of (A) are given by
$Q=ue^{-3t/40}$ where $u'e^{-3t/40}=3$. Therefore,$u'=3e^{3t/40}$,
$u=40e^{3t/40}+c$, and $Q=40+ce^{-3t/40}$. Now $Q(0)=0\Rightarrow
c=-40$, so $Q=40(1-e^{-3t/40})$.


\item
$Q'=\dst\frac{3}{2}-\dst\frac{Q}{20},\quad Q(0)=10$. Rewrite the
differential equation as (A) $Q'+\dst\frac{Q}{20}=\dst\frac{3}{2}$. Since
$Q_1=e^{-t/20}$ is a solution of the complementary equation, the
solutions of (A) are given by $Q=ue^{-t/20}$ where
$u'e^{-t/20}=\dst\frac{3}{2}$. Therefore,$u'=\dst\frac{3}{2}e^{t/20}$,
$u=30e^{t/20}+c$, and $Q=30+ce^{-t/20}$. Now $Q(0)=10\Rightarrow
c=-20$, so $Q=30-20e^{-t/20}$ and $K=\dst\frac{Q}{100}=.3-.2e^{-t/20}$.



\item
$Q'=10-\dst\frac{Q}{5}$, or
 (A) $Q'+\dst\frac{Q}{5}=10$. Since $Q_1=e^{-t/5}$ is a solution
of the complementary equation, the solutions of (A) are given by
$Q=ue^{-t/5}$ where $u'e^{-t/5}=10$. Therefore,$u'=10e^{t/5}$,
$u=50e^{t/10}+c$, and $Q=50+ce^{-t/5}$. Since
$\lim_{t\to\infty}Q(t)=50$, the minimum capacity is 50 gallons.



\item
 Since there are $2t+600$ gallons of mixture in the tank at
time $t$ and mixture is being drained at 4
gallons/min, $Q'=3-\dst\frac{2}{ t+300}Q,\quad Q(0)=40$.
 Rewrite the differential
equation
as (A) $Q'+\dst\frac{2}{ t+300}Q=3$. Since $Q_1=\dst\frac{1}{(t+300)^2}$
is a
solution of the complementary equation, the solutions of (A) are given
by
$Q=\dst\frac{u}{(t+300)^2}$ where $\dst\frac{u'}{(t+300)^2}=3$. Therefore,
$u'=3(t+300)^2$,
$u=(t+300)^3+c$, and $Q=t+300+\dst\frac{c}{(t+300)^2}$.
 Now
$Q(0)=40\Rightarrow c=-234\times10^5$, so
 $Q= t+300-\dst\frac{234 \times 10^5}{(t+300)^2},
\, 0\le t\le300$.



\item
\begin{alist}
\item $S'=-k_m(S-T_m),\ S(0)=0$, so (A)
$S=T_m+(S_0-T_m)e^{-k_mt}$.
$T'=-k(T-S)=-k\left(T-T_m-(S_0-T_m)e^{-k_mt}\right)$,
from
(A). Therefore,$T'+kT=kT_m+k(S_0-T_m)e^{-k_mt}$; $T=ue^{-kt}$;
(B) $u'=kT_me^{kt}+k(S_0-T_m)e^{(k-k_m)t}$;
 $u=T_me^{kt}+\dst\frac{k}{
k-k_m}(S_0-T_m)e^{(k-k_m)t}+c$;
 $T(0)=T_0\Rightarrow
 c=T_0-T_m-\dst\frac{k}{ k-k_m}(S_0-T_m)$;
 $u=T_me^{kt}+\dst\frac{k}{
k-k_m}(S_0-T_m)e^{(k-k_m)t}+T_0-T_m-\dst\frac{k}{ k-k_m}(S_0-T_m)$;
$T=T_m+(T_0-T_m)e^{-kt}+\dfrac{k(S_0-T_m)}{
(k-k_m)}\left(e^{-k_mt}-e^{-kt}\right)$.

\item If $k=k_m$ (B) becomes
(B) $u'=kT_me^{kt}+k(S_0-T_m)$; $u=T_me^{kt}+k(S_0-T_m)t+c$;
$T(0)=T_0\Rightarrow c=T_0-T_m$;
$u=T_me^{kt}+k(S_0-T_m)t+(T_0-T_m)$;
$T=T_m+k(S_0-T_m)te^{-kt}+(T_0-T_m)e^{-kt}$.

\item $\lim_{t\to\infty}T(t)=\lim_{t\to\infty}S(t)=T_m$
in either  case.
\end{alist}

\item
$V'=aV-bV^2=-bV(V-b/a)$; $\dst\frac{V'}{ V(V-a/b)}=-b$;
$\dst\left[\frac{1}{ V-a/b}-\frac{1}{ V}\right]V'=-a$;
$\ln\dst\left|\frac{V-a/b}{ V}\right|=-at+k$; (A) $\dfrac{V-a/b}{
V}=ce^{-at}$; (B) $V=\dst{\frac{a}{ b}\frac{1}{1-ce^{-at}}}$. Since
$V(0)=V_0$,  (A) $\Rightarrow c=\dfrac{V_0-a/b}{ V_0}$. Substituting
this into (B) yields
$V=\dst\frac{a}{ b}\dst\frac{V_0}{ V_0-\left(V_0-a/b
\right)e^{-at}}$ so $\lim_{t\to\infty}V(t)=a/b$

\item
If $Q_n(t)$ is the number of pounds of salt in $T_n$ at time $t$, then
$Q_{n+1}'+\dst\frac{r}{ W}Q_{n+1}=rc_n(t),\ n=0,1,\dots$, where
$c_0(t)\equiv c$. Therefore,$Q_{n+1}=u_{n+1}e^{-rt/w}$;
(A) $u_{n+1}'=re^{rt/W}c_n(t)$. In particular, with $n=0$,
$u_1=cW(e^{rt/W}-1)$, so $Q_1=cW(1-e^{-rt/W})$  and
$c_1=c(1-e^{-rt/W})$. We will shown by induction that
$c_n=c\dst{\left(1-e^{-rt/W}\sum_{j=0}^{n-1}\frac{1}{
j!}\left(\frac{rt}{ W}\right)^j\right)}$.
 This is true for $n=1$; if it is
true for a given $n$, then, from (A),
\[
u_{n+1}'=cre^{rt/W}
\left(1-e^{-rt/W}\sum_{j=0}^{n-1}\frac{1}{
j!}\left(\frac{rt}{ W}\right)^j\right)
=cre^{rt/W}-cr
\sum_{j=0}^{n-1}\frac{1}{
j!}\left(\frac{rt}{ W}\right)^j,
\]
so (since $Q_{n+1}(0)=0$),
\[
u_{n+1}=cW(e^{rt/W}-1)-c\sum_{j=0}^{n-1}\frac{1}{(j+1)!}
\frac{r^{j+1}}{ W^j}t^{j+1}.
\]
Therefore,
\[
c_{n+1}=\frac{1}{ W}u_{n+1}e^{-rt/W}=
c\dst{\left(1-e^{-rt/W}\sum_{j=0}^n\frac{1}{
j!}\left(\frac{rt}{ W}\right)^j\right)},
\]
which completes the induction. From this, $\lim_{t\to\infty}c_n(t)=c$.



\item
Since the incoming solution contains 1/2 lb of salt per gallon
and there are always 600 gallons in the tank, we conclude
intuitively that $\lim_{t\to\infty}Q(t)=300$.
To verify this rigorously,
note that $Q_1(t)=\dst\exp\left(-\frac{1}{150}\int_0^t
a(\tau)\,d\tau\right)$ is a solution of the complementary equation,
(A) $Q_1(0)=1$, and (B) $\lim_{t\to\infty}Q_1(t)=0$ (since
$\lim_{t\to\infty}a(t)=1$).  Therefore,$Q=Q_1u$; $Q_1u'=2$;
$u'=\dst\frac{2}{ Q_1}$;
$u=\dst Q_0+2\int_0^t\frac{d\tau}{ Q_1(\tau)}$ (see (A)), and
$Q(t)=\dst Q_0Q_1(t)+2Q_1(t)\int_0^t\frac{d\tau}{ Q_1(\tau)}$.
From (B),
$\dst\lim_{t\to\infty}Q(t)=2\lim_{t\to\infty}Q_1(t)
\int_0^t\frac{d\tau}{ Q_1(\tau)}$, a $0\cdot\infty$ indeterminate
form. By L'Hôpital's rule,
$\dst\lim_{t\to\infty}Q(t)=2\lim_{t\to\infty}
\dst\frac{1}{ Q_1(t)}\bigg/\left(\frac{-Q_1'(t)}{ Q_1^2(t)}\right)
=-2\lim_{t\to\infty}\frac{Q_1(t)}{ Q_1'(t)}=300$.
\end{sectionSolutions}

\section{Elementary Mechanics}

\begin{sectionSolutions}
\item
The firefighter's mass is $m=\dst\frac{192}{32}=6$ sl, so
 $6v'=-192-kv$, or (A) $v'+\dst\frac{k}{6}v=-32$.
Since $v_1=e^{-kt/6}$ is a solution of the complementary equation,
the solutions of (A) are $v=ue^{-kt/6}$ where $u'e^{-kt/6}=-32$.
Therefore,$u'=-32e^{kt/6}$;\ $u=-\dst\frac{192}{ k}e^{kt/6}+c$;\;
$v=-\dst\frac{192}{ k}+ce^{-kt/6}$. Now $v(0)=0\Rightarrow
c=\dst\frac{192}{ k}$. Therefore,$v=-\dst\frac{192}{ k}(1-e^{-kt/6})$
and $\lim_{t\to\infty}v(t)=-\dst\frac{192}{ k}=-16$ ft/s, so $k=12$
lb-s/ft and $v=-16(1-e^{-2t})$.

\addtocounter{enumi}{-1}
\item%4.3.3
 $m=\dst\frac{64000}{32}=2000$, so
 $2000v'=50000-2000v$, or (A)
$v'+v=25$.
Since $v_1=e^{-t}$ is a solution of the complementary equation,
the solutions of (A) are $v=ue^{-t}$ where $u'e^{-t}=25$.
Therefore,$u'=25e^t$;\ $u=25e^t+c$;\;
$v=25+ce^{-t}$. Now $v(0)=0\Rightarrow
c=-25$. Therefore,$v=25(1-e^{-t})$
and $\lim_{t\to\infty}v(t)=25$ ft/s.

\addtocounter{enumi}{-1}
\item%4.3.4
 $20v'=10-\dst\frac{1}{2}v$, or (A)
$v'+\dst\frac{1}{20}v=\dst\frac{1}{2}$.
Since $v_1=e^{-t/40}$ is a solution of the complementary equation,
the solutions of (A) are $v=ue^{-t/40}$ where
$u'e^{-t/40}=\dst\frac{1}{2}$. Therefore,$u'=\dfrac{e^{t/40}}{2}$;\;
$u=20e^{t/40}+c$;\;
$v=20+ce^{-t/40}$. Now $v(0)=-7\Rightarrow
c=-27$. Therefore,$v=20-27e^{-t/40}$.



\item
 $m=\dst\frac{3200}{32}=100$ sl. The
component of the gravitational force in the direction of motion is
$-3200\cos(\pi/3)=-1600$ lb. Therefore,
$100v'=-1600+v^2$. Separating variables yields
$\dst\frac{v'}{(v-40)(v+40)}=\dst\frac{1}{100}$, or $\dst{\left[\frac{1}{
v-40}-\frac{1}{ v+40}\right]}=\frac{4}{5}$. Therefore,
$\dst{\ln\left|\frac{v-40}{
v+40}\right|}=\dst\frac{4t}{5}+k$ and $\dst\frac{v-40}{ v+40}=ce^{4t/5}$.
Now $v(0)=-64\Rightarrow c=\dst\frac{13}{3}$; therefore $\dst\frac{v-40}{
v+40}=\dfrac{13e^{4t/5}}{3}$, so
$v=\dfrac{40(3+13e^{4t/5})}{3-13e^{4t/5}}$, or
$v=-\dfrac{40(13+3e^{-4t/5})}{13-3e^{-4t/5}}$.



\item
From Example~4.3.1, (A) $v=-\dst\frac{mg}{
k}+\dst\left(v_0+\frac{mg}{ k}\right)e^{-kt/m}$. Integrating this yields
(B) $y=-\dst\frac{mgt}{
k}-\dst{\frac{m}{ k}\left(v_0+\frac{mg}{ k}\right)}e^{-kt/m}+c$.
Now $y(0)=y_0\Rightarrow c=y_0+\dst{\frac{m}{ k}\left(v_0+\frac{mg}{
k}\right)}$. Substituting this into (B) yields
\begin{align*}
y&=-\dst\frac{mgt}{
k}-\dst{\frac{m}{ k}\left(v_0+\frac{mg}{ k}\right)}e^{-kt/m}+
y_0+\dst{\frac{m}{ k}\left(v_0+\frac{mg}{ k}\right)}\\
&=y_0+\frac{m}{ k}\left(v_0-gt+\frac{mg}{
k}-\left(v_0+\frac{mg}{ k}\right)e^{-kt/m}\right)\\
&=y_0+\frac{m}{ k}(v_0-v-gt)
\end{align*}
where the last equality follows from (A).


\item
 $m=\dst\frac{256}{32}=8$ sl.
Since the
resisting force is 1 lb when $|v|=\qty{4}{\feet\per\second}$,
$k=\dst\frac{1}{16}$.
 Therefore,
$8v'=-256+\dst\frac{1}{16}v^2=\dst\frac{1}{16}\left(v^2-(64)^2\right)$.
Separating variables yields
$\dst\frac{v'}{(v-64)(v+64)}=\dst\frac{1}{128}$, or $\dst{\left[\frac{1}{
v-64}-\frac{1}{ v+64}\right]}v'=1$. Therefore,$\dst{\ln\left|\frac{v-64}{
v+64}\right|}=t+k$ and $\dst\frac{v-64}{ v+64}=ce^t$.
Now $v(0)=0\Rightarrow c=-1$; therefore $\dst\frac{v-64}{
v+64}=-e^t$, so
$v=\dfrac{64(1-e^t)}{1+e^t}$, or
$v=-\dfrac{64(1-e^{-t})}{1+e^{-t}}$.  Therefore,
$\lim_{t\to\infty}v(t)=-64$.




\item
\begin{alist}
\item $mv'=-mg-kv^2=-mg(1+\gamma^2v^2)$, where
$\gamma=\dst{\sqrt\frac{k}{
mg}}$. Therefore,(A)  $\dst\frac{v'}{1+\gamma^2v^2}=-g$. With the
substitution $u=\gamma v$,
 $\dst\int\frac{dv}{1+\gamma^2 v^2}
=\dst{\frac{1}{\gamma}\int\frac{du}{1+u^2}}=\dst\frac{1}{\gamma}\tan^{-1}u=
\dst\frac{1}{\gamma}\tan^{-1}(\gamma v)$.
Therefore,
$\dst{\frac{1}{\gamma}\tan^{-1}(\gamma v)}=-gt+c$. Now
$v(0)=v_0\Rightarrow c=\dst{\frac{1}{\gamma}\tan^{-1}(\gamma v_0)}$, so
$\dst{\frac{1}{\gamma}\tan^{-1}(\gamma v)}=-gt+
\dst{\frac{1}{\gamma}\tan^{-1}(\gamma v_0)}$.
Since $v(T)=0$, it follows that $T=\dst\frac{1}{\gamma g}\tan^{-1}\gamma v_0
=\dst{\sqrt{\frac{m}{ kg}}
\tan^{-1}\left(v_0 \sqrt{\frac{k}{ mg}}\right)}$.

\item Replacing $t$ by $t-T$ and setting $v_0=0$in the answer to
the previous exercise yields
$v =-\dst{\sqrt\frac{mg}{ k} \,
\frac{1-e^{-2\sqrt\frac{gk}{ m}\, (t-T)}}{ {1+e^{-2\sqrt\frac{gk}{ m}\,
(t-T)}}}}$.
\end{alist}

\item
\begin{alist}
\item $mv'=-mg+f(|v|)$; since $s=|v|=-v$, (A) $ms'=mg-f(s)$.

\item Since $f$ is increasing and $\lim_{t\to\infty}f(s)\le mg$,
$mg-f(s)>0$ for all~$s$. This and (A) imply that $s$ is an increasing
function of $t$, so either (B) $\lim_{t\to\infty}s(t)=\infty$
or (C ) $\lim_{t\to\infty}s(t)=\overline s<\infty$. However,
(A) and (C) imply that $s'(t)>K=g-f(\overline s)/m$ for all $t>0$.
Consequently, $s(t)>s_0+Kt$ for all $t>0$, which contradicts (C)
because $K>0$.

\item There is a unique positive number $s_T$ such that
$f(s_T)=mg$, and $s\equiv s_T$ is a constant solution of (A).
Now suppose that $s(0)<s_T$. Then Theorem~2.3.1
implies that (D) $s(t)<s_T$ for all $t>0$, so (A) implies that
$s$ is strictly increasing. This and (D) imply that
$\lim_{t\to\infty}s(t)=\overline s\le s_T$. If $\overline s<s_T$
then (A) implies that $s'(t)>K=g-f(\overline s)/m$. Consequently,
$s(t)>s(0)+Kt$, which contradicts (D) because $K>0$. Therefore,
$s(0)<s_T\Rightarrow\lim_{t\to\infty}s(t)=s_T$. A similar proof with
inequalities reversed shows that
$s(0)>s_T\Rightarrow\lim_{t\to\infty}s(t)=s_T$.
\end{alist}

\item
\begin{alist}
\item (A) $mv'=-mg+k\sqrt{|v|}$; since the magnitude of the
resistance is \qty{64}{\pound} when $v=\qty{16}{\feet\per\second}$, $4k=64$,
so $k=\qty{16}{\pound\second\tothe{0.5}\per\foot\tothe{0.5}}$.
Since $m=2$ and $g=32$,
(A) becomes $2v'=-64+16\sqrt{|v|}$, or $v'=-32+8\sqrt{|v|}$.

\item From Exercise~4.3.14\part{c}, $v_T$ is the negative
number such that $-32+8\sqrt{|v_T|}=0$; thus, $v_T=\qty{-16}{\feet\per\second}$.
\end{alist}

\item
With $h=0$, $v_e=\sqrt{2gR}$, where $R$ is the radius of the moon and
$g$ is the acceleration due to gravity at the moon's surface. With
length in miles, $g=\dst\frac{5.31}{5280}$ \unit{\mile\per\second\squared}, so
$v_e=\dst{\sqrt\frac{2\cdot5.31\cdot1080}{5280}} \approx \qty{1.47}{\mile\per\second}$.



\item
Suppose that there is a number $y_m$ such that $y(t)\le y_m$ for all
$t\ge0$ and let $\alpha=\dst\frac{gR^2}{(y_m+R)^2}$. Then
$\dst\frac{d^2y}{ dt^2}\le-\alpha$ for all $t\ge 0$. Integrating this
inequality from $t=0$ to $t=T>0$ yields $v(T)-v_0\le -\alpha T$, or
$v(T)\le v_0-\alpha T$, so $v(T) < 0$ for $T>\dst\frac{v_0}{\alpha}$.
This implies that the vehicle must eventually fall back to Earth,
which contradicts the assumption that it continues to climb forever.

\end{sectionSolutions}

\section{Autonomous Second Order Equations}

\begin{sectionSolutions}

\addtocounter{enumi}{-1}
\item%4.4.1
$\overline y=0$ is a stable equilibrium.
The phase plane equivalent is $\dst{v\frac{dv}{ dy}}+y^3=0$, so the
 trajectories are $v^2+\dst\frac{y^4}{4}=c$.




\addtocounter{enumi}{-1}
\item%4.4.2
$\overline y=0$ is an unstable equilibrium.
The phase plane equivalent is $\dst{v\frac{dv}{ dy}}+y^2=0$, so the
 trajectories are $v^2+\dst\frac{2y^3}{3}=c$.



\item
$\overline y=0$ is a stable equilibrium.
The phase plane equivalent is $\dst{v\frac{dv}{ dy}}+ye^{-y}=0$, so the
 trajectories are $v^2-e^{-y}(y+1)=c$.




\item
$p(y)=y^3-4y=(y+2)y(y-2)$, so the
 equilibria are $-2,0,2$.
Since
\begin{align*}
y(y-2)(y+2)&<0\quad \text{if } y<-2\text{ or }0<y<2,\\
&>0\quad\text{ if }-2<y<0\text{ or }y>2,
\end{align*}
$0$ is unstable and $-2,2$ are stable. The phase plane equivalent is
$\dst{v\frac{dv}{ dy}}+y^3-4y=0$, so the trajectories are
$2v^2+y^4-8y^2=c$. Setting $(y,v)=(0,0)$ yields $c=0$, so the equation
of the separatrix is $2v^2-y^4+8y^2=0$.


\item
$p(y)=y(y-2)(y-1)(y+2)$, so the
 equilibria are $-2,0,1,2$.
Since
\begin{align*}
y(y-2)(y-1)(y+2)&>0\quad \text{ if } y<-2\text{ or }0<y<1\text{ or }y>2,\\
&<0\quad\text{ if }-2<y<0\text{ or }1<y<2,
\end{align*}
$0,2$ are stable and $-2,1$ are unstable. The phase plane
equivalent is
$\dst{v\frac{dv}{ dy}}+y(y-2)(y-1)(y+2)=0$, so the trajectories are
$30v^2+y^2(12y^3-15y^2-80y+120)=c$. Setting $(y,v)=(-2,0)$
and $(y,v)=(1,0)$
yields
$c=496$ and $c=37$ respectively, so the equations of the separatrices
are $30v^2+y^2(12y^3-15y^2-80y+120)=496$
and $30v^2+y^2(12y^3-15y^2-80y+120)=37$.


\item
$p(y)=y^3-ay$.
If $a\le0$, then $p(0)=0$, $p(y)>0$ if $y>0$, and $p(y)<0$ if $y<0$, so
$0$ is  stable.
 If $a>0$, then
\begin{align*}
y^3-ay=y(y-\sqrt a)(y+\sqrt a)&>0\text{ if }-\sqrt a<y<0\text{ or }
y>\sqrt a,\\
&<0\text{ if }y<-\sqrt a\text{ or }0<y<\sqrt a,
\end{align*}
so $-\sqrt a$ and $\sqrt a$ are  stable  and $0$
is unstable. We say that $a=0$ is a critical value
because  it separates the two cases.



\item
$p(y)=y-ay^3$. If $a\le0$, then $p(0)=0$, $p(y)>0$ if $y>0$, and
$p(y)<0$ if $y<0$, so $0$ is stable. If $a>0$, then
\begin{align*}
y-ay^3=-ay(y-1/\sqrt a)(y+1/\sqrt a)
&>0\text{ on }(-\infty,-1/\sqrt a)\cup(0,1/\sqrt a)\\
&<0\text{ on }(-1/\sqrt a,0)\cup(1/\sqrt a,\infty)
%\\
%y-ay^3=-ay(y-1/\sqrt a)(y+1/\sqrt a)
%&>0\text{ if }y<-1/\sqrt a\text{ or } 0<y<1/\sqrt a\\
%&<0\text{ if }-1/\sqrt a<y<0\text{ or }1/\sqrt a<y,
\end{align*}
so $-\sqrt a$ and $\sqrt a$ are unstable and $0$ is stable. We say
that $a=0$ is a critical value because it separates the two cases.


\addtocounter{enumi}{10}
\item%4.4.24
\begin{alist}
\item Since $v'=-p(y)\ge k$ and $v(0)=0$, $v\ge kt$ and therefore
$y\ge y_0+kt^2/2$ for $0\le t<T$.

\item
Let $0<\epsilon<\rho$.
Suppose that $y$ is the solution of the initial value problem
(A) $y''+p(y)=0,\quad y(0)=y_0,\quad y'(0)=0$, where $\overline
y<y_0<\overline y+\epsilon$. Now let $Y=y-\overline y$ and
$P(Y)=p(Y+\overline y)$. Then $P(0)=0$ and $P(Y)<0$ if
$0<Y\le\rho$ . Moreover, $Y$ is the solution of
 $Y''+p(Y)=0,\quad Y(0)=Y_0,\quad Y''(0)=0$, where $Y_0=y_0-\overline
y$, so $0<Y_0<\epsilon$. From \part{a}, $Y(t)\ge\epsilon$ for some
$t>0$. Therefore,$y(t)>\overline y+\epsilon$ for some $t>0$,
so $\overline y$ is an unstable equilibrium of $y''+p(y)=0$.
\end{alist}
\end{sectionSolutions}


\section{Applications to Curves}

\begin{sectionSolutions}
\item
Differentiating (A) $e^{xy}=cy$ yields (B) $(xy'+y)e^{xy}=cy'$. From
(A), $c=\dfrac{e^{xy}}{ y}$. Substituting this into (B) and
cancelling $e^{xy}$ yields $xy'+y=\dst\frac{y'}{ y}$, so
$y'=-\dst\frac{y^2}{(xy-1)}$.

\item
Solving $y=x^{1/2}+cx$ for $c$ yields $c=\dst\frac{y}{ x}-x^{-1/2}$, and
differentiating yields $0=\dst\frac{y'}{ x}-\dst\frac{y}{
x^2}+\dfrac{x^{-3/2}}{2}$, or $xy'-y=-\dfrac{x^{1/2}}{2}$.


\item
Rewriting $\dst{y=x^3+\frac{c}{ x}}$ as $xy=x^4+c$ and differentiating
yields $xy'+y=4x^3$.


\item
Rewriting $y=e^x+c(1+x^2)$ as
$\dst\frac{y}{1+x^2}=\dst\frac{e^x}{1+x^2}+c$ and differentiating yields
$\dst\frac{y'}{1+x^2}-\dst\frac{2xy}{(1+x^2)^2}=
\dst\frac{e^x}{1+x^2}-\dst\frac{2xe^x}{(1+x^2)^2}$, so
$(1+x^2)y'-2xy=(1-x)^2e^x$.


\item
If (A) $y=f+cg$, then (B) $y'=f'+cg'$. Multiplying (A) by $g'$
and (B) by $g$ yields
(C) $yg'=fg'+cgg'$ and  (D) $y'g=f'g+c g'g$, and subtracting (C)
from (D) yields
$y'g-yg'=f'g-fg'$.


\item Let $(x_0,y_0)$ be the center
and $r$ be the radius of a circle in the family. Since $(-1,0)$ and
$(1,0)$ are on the circle, $(x_0+1)^2+y_0^2=(x_0-1)^2+y_0^2$, which
implies that $x_0=0$. Therefore, the equation of the circle is (A)
$x^2+(y-y_0)^2=r^2$. Since $(1,0)$ is on the circle, $r^2=1+y_0^2$.
Substituting this into (A) shows that the equation of the circle is
$x^2+y^2-2yy_0=1$, so $2y_0=\dst\frac{x^2+y^2-1}{ y}$. Differentiating
$y(2x+2yy')-y'(x^2+y^2-1)=0$, so $y'(y^2-x^2+1)+2xy=0$.


\item
From Example~4.5.6 the equation of the line tangent to the
parabola at $(x_0,x_0^2)$ is (A) $y=-x_0^2+2x_0x$.

\begin{alist}
\item From (A), $(x,y)=(5,9)$ is on the tangent line through
$(x_0,x_0^2)$ if and only if $9=-x_0^2+10x_0$, or
$x_0^2-10x_0+9=(x_0-1)(x_0-9)=0$. Letting $x_0=1$ in (A) yields the
line $y=-1+2x$, tangent to the parabola at $(x_0,x_0^2)=(1,1)$.
Letting $x_0=9$ in (A) yields the line $y=-81+18x$, tangent to the
parabola at $(x_0,x_0^2)=(9,81)$.

\item From (A), $(x,y)=(6,11)$ is on the tangent line through
$(x_0,x_0^2)$ if and only if $11=-x_0^2+12x_0$, or
$x_0^2-12x_0+11=(x_0-1)(x_0-11)=0$. Letting $x_0=1$ in (A) yields the
line $y=-1+2x$, tangent to the parabola at $(x_0,x_0^2)=(1,1)$.
Letting $x_0=11$ in (A) yields the line $y=-121+22x$, tangent to the
parabola at $(x_0,x_0^2)=(11,21)$.

\item From (A), $(x,y)=(-6,20)$ is on the tangent line through
$(x_0,x_0^2)$ if and only if $20=-x_0^2-12x_0$, or
$x_0^2+12x_0+20=(x_0+2)(x_0+10)=0$. Letting $x_0=-2$ in (A) yields the
line $y=-4-4x$, tangent to the parabola at $(x_0,x_0^2)=(-2,4)$.
Letting $x_0=-10$ in (A) yields the line $y=-100-20x$, tangent to the
parabola at $(x_0,x_0^2)=(-10,100)$.

\item From (A), $(x,y)=(-3,5)$ is on the tangent line through
$(x_0,x_0^2)$ if and only if $5=-x_0^2-6x_0$, or
$x_0^2+6x_0+5=(x_0+1)(x_0+5)=0$. Letting $x_0=-1$ in (A) yields the
line $y=-1-2x$, tangent to the parabola at $(x_0,x_0^2)=(-1,1)$.
Letting $x_0=-5$ in (A) yields the line $y=-25-10x$, tangent to the
parabola at $(x_0,x_0^2)=(-5,25)$.
\end{alist}

\addtocounter{enumi}{-1}
\item%4.5.15
\begin{alist}
\item If $(x_0,y_0)$ is any point on the circle such that
$x_0\ne\pm1$ (and therefore $y_0\ne0$), then differentiating (A)
yields $2x_0+2y_0y_0'=0$, so $y_0'=-\dst\frac{x_0}{ y_0}$. Therefore,the
equation of the tangent line is $y=y_0-\dst\frac{x_0}{ y_0}(x-x_0)$.
Since $x_0^2+y_0^2=1$, this is equivalent to (B).

\item
Since $y'=-\dst\frac{x_0}{ y_0}$ on the tangent line, we can rewrite (B)
as $y-xy'=\dst\frac{1}{ y_0}$. Hence (F) $\dst\frac{1}{(y-xy')^2}=y_0^2$
and (G) $x_0^2=1-y_0^2=\dfrac{(y-xy')^2-1}{(y-xy')^2}$. Since
$(y')^2=\dst\frac{x_0^2}{ y_0^2}$, (F) and (G) imply that
$(y')^2=(y-xy')^2-1$, which implies (C).

\item
Using the quadratic formula to solve (C) for $y'$
yields
\[
y'=\dfrac{xy\pm\sqrt{x^2+y^2-1}}{ x^2-1}
\tag{H}
\]
if $(x,y)$ is on a tangent  line with slope $y'$.
If $y=\dst\frac{1-x_0x}{ y_0}$, then
$x^2+y^2-1=x^2+\dst{\left(\frac{1-x_0x}{
y_0}\right)^2}-1=\dst{\left(\frac{x-x_0}{ y_0}\right)^2}$
(since $x_0^2+y_0^2=1$). Since $y'=-\dst\frac{x_0}{ y_0}$, this implies
that (H) is equivalent to
$-\dst\frac{x_0}{ y_0}=\dst{\frac{1}{ x^2-1}\left[\frac{x(1-x_0x)}{ y_0}
\pm\left|\frac{x-x_0}{ y_0}\right|\right]}$,
which holds if and only if we choose the ``$\pm$" so that
 $\pm\dst{\left|\frac{x-x_0}{ y_0}\right|}=-\dst{\left(\frac{x-x_0}{
y_0}\right)}$.  Therefore, we must choose $\pm=-$ if $\dst\frac{x-x_0}{
y_0}>0$, so (H) reduces to (D),
or  $\pm=+$ if $\dst\frac{x-x_0}{
y_0}<0$, so (H) reduces to (E).

\item Differentiating (A)
yields $2x+2yy'=0$, so $y'=-\dst\frac{x}{ y}$ on either semicircle.
Since (D) and (E) both reduce to $y'=\dst\frac{xy}{1-x^2}=-\dst\frac{x}{ y
}$ (since $x^2+y^2=1$) on both semicircles, the conclusion follows.

\item From (D) and (E) the slopes of tangent lines from (5,5)
tangent to the circle are
$y'=\dst\frac{25\pm\sqrt49}{24}=\dst\frac{3}{4},\dst\frac{4}{3}$. Therefore,
tangent lines are $y=5+\dst\frac{3}{4}(x-5)=\dfrac{1+3x/5}{4/5}$ and
$y=5+\dst\frac{4}{3}(x-5)=\dfrac{1-4x/5}{-3/5}$, which intersect the
circle at $(-3/5,4/5)$ $(4/5,-3/5)$, respectively. (See (B)).
\end{alist}

\addtocounter{enumi}{-1}
\item%4.5.16
\begin{alist}
\item If $(x_0,y_0)$ is any point on the parabola such that
$x_0>0$ (and therefore $y_0\ne0$), then differentiating (A) yields
$1=2y_0y_0'$, so $y_0'=\dst\frac{1}{ 2y_0}$. Therefore,the equation of
the tangent line is $y=y_0+\dst\frac{1}{2y_0}(x-x_0)$. Since
$x_0=y_0^2$, this is equivalent to (B).

\item
Since $y'=\dst\frac{1}{2y_0}$ on the tangent line, we can rewrite (B) as
$\dst\frac{y_0}{2}=y-xy'$. Substituting this into (B) yields
$y=(y-xy')+\dst\frac{x}{4(y-xy')}$, which implies~(C).

\item
Using the quadratic formula to solve (C) for $y'$ yields
\[
y'=\dfrac{y\pm\sqrt{y^2-x}}{2x}
\tag{F}
\]
if $(x,y)$ is on a tangent line with slope $y'$. If
$y=\dst\frac{y_0}{2}+\dst\frac{x}{2y_0}$, then
$y^2-x=\dst{\frac{1}{4}\left(y_0-\dst\frac{x}{ y_0}\right)^2}$ so (F) is
equivalent to $\dst\frac{1}{2y_0}= \dfrac{y_0+\dst\frac{x}{
y_0}\pm\left|y_0-\dst\frac{x}{ y_0}\right|}{4x}$ which holds if and
only if we choose the ``$\pm$" so that $\pm\dst{\left|y_0-\frac{x}{ y_0}
\right|}=-\dst{\left(y_0-\frac{x}{ y_0}\right)}$. Therefore,we must
choose $\pm=+$ if $x>y_0^2=x_0$, so (F) reduces to (D), or $\pm=-$ if
$x<y_0^2=x_0$, so (F) reduces to (E).

\item Differentiating (A) yields $1=2yy'$, so $y'=\dst\frac{1}{2y}$
on either half of the parabola. Since (D) and (E) both reduce to this
if $x=y^2$, the conclusion follows.
\end{alist}

\item
The equation of the line tangent to the curve at $(x_0,y(x_0))$
is $y=y(x_0)+y'(x_0)(x-x_0)$. Since $y(x_0/2)=0$,
$y(x_0)-\dfrac{y'(x_0)x_0}{2}=0$. Since $x_0$ is arbitrary,
it follows that $y'=\dst\frac{2y}{ x}$, so
$\dst\frac{y'}{ y}=\dst\frac{2}{ x}$, $\ln|y|=2\ln|x|+k$, and
$y=cx^2$. Since $(1,2)$ is on the curve, $c=2$. Therefore,$y=2x^2$.


\item
The equation of the line tangent to the curve at $(x_0,y(x_0))$ is
$y=y(x_0)+y'(x_0)(x-x_0)$. Since $(x_1,y_1)$ is on the line,
$y(x_0)+y'(x_0)(x_1-x_0)=y_1$. Since $x_0$ is arbitrary, it follows
that $y+y'(x_1-x)=y_1$, so $\dst\frac{y'}{ y-y_1}=\dst\frac{1}{ x-x_1}$,
$\ln|y-y_1|=\ln|x-x_1|+k$, and $y-y_1=c(x-x_1)$.


\item
The equation of the line tangent to the curve at $(x_0,y(x_0))$ is
$y=y(x_0)+y'(x_0)(x-x_0)$. Since $y(0)=x_0$, $x_0=y(x_0)-y'(x_0)x_0$.
Since $x_0$ is arbitrary, it follows that $x=y-xy'$, so (A)
$y'-\dst\frac{y}{ x}=-1$. The solutions of (A) are of the form $y=ux$,
where $u'x=-1$, so $u'=-\dst\frac{1}{ x}$. Therefore,$u=-\ln|x|+c$ and
$y=-x\ln|x|+cx$.


\item
The equation of the line normal to the curve at $(x_0,y_0)$ is
$y=y(x_0)-\dst\frac{x-x_0}{ y'(x_0)}$. Since $y(0)=2y(x_0)$,
$y(x_0)+\dst\frac{x_0}{ y'(x_0)}=2y(x_0)$. Since $x_0$ is arbitrary, it
follows that $y'y=x$, so (A)
$\dst\frac{y^2}{2}=\dst\frac{x^2}{2}+\dst\frac{c}{2}$ and $y^2=x^2+c$. Now
$y(2)=1\leftrightarrow c=-3$. Therefore,$y=\sqrt{x^2-3}$.


\item
Differentiating the given equation yields $2x+4y+4xy'+2yy'=0$, so
$y'=-\dst\frac{x+2y}{2x+y}$ is a differential equation for the given
family, and (A) $y'=\dst\frac{2x+y}{ x+2y}$ is a differential equation
for the orthogonal trajectories. Substituting $y=ux$ in (A) yields
$u'x+u=\dst\frac{2+u}{1+2u}$, so $u'x=-\dfrac{2(u^2-1)}{1+2u}$ and
$\dst\frac{1+2u}{(u-1)(u+1)}u'=-\dst\frac{2}{ x}$, or $\dst{\left[\frac{3}{
u-1}+\frac{1}{ u+1}\right]}u'=-\dst\frac{4}{ x}$. Therefore,
$3\ln|u-1|+\ln|u+1|=-4\ln|x|+K$, so $(u-1)^3(u+1)=\dst\frac{k}{ x^4}$.
Substituting $u=\dst\frac{y}{ x}$ yields the orthogonal trajectories
$(y-x)^3(y+x)=k$.


\item
Differentiating yields $ye^{x^2}(1+2x^2)+xe^{x^2}y'=0$, so
$y'=\dfrac{y(1+2x^2)}{ x}$ is a differential equation for the given
family. Therefore,(A) $y'=-\dst\frac{x}{ y(1+2x^2)}$ is a differential
equation for the orthogonal trajectories. From (A),
$yy'=-\dst\frac{x}{1+2x^2}$, so
$\dst\frac{y^2}{2}=-\dst\frac{1}{4}\ln(1+2x^2)+\dst\frac{k}{2}$, and the
orthogonal trajectories are given by
$y^2=-\dst\frac{1}{2}\ln(1+2x^2)+k$.


\item
Differentiating (A) $y=1+cx^2$ yields (B) $y'=2cx$. From (C),
$c=\dst\frac{y-1}{ x^2}$. Substituting this into (B) yields the
differential equation $y'=\dfrac{2(y-1)}{ x}$ for the given family of
parabolas. Therefore,$y'=-\dst\frac{x}{2(y-1)}$ is a differential
equation for the orthogonal trajectories. Separating variables yields
$2(y-1)y'=-x$, so $(y-1)^2=-\dst\frac{x^2}{2}+k$. Now
$y(-1)=3\leftrightarrow k=\dst\frac{9}{2}$, so
$(y-1)^2=-\dst\frac{x^2}{2}+\dst\frac{9}{2}$. Therefore,(D)
$y=\dst{1+\sqrt\frac{9-x^2}{2}}$. This curve intersects the parabola
(A) if and only if the equation (C) $cx^2=\dst\sqrt\frac{9-x^2}{2}$ has
a solution $x^2$ in $(0,9)$. Therefore,$c>0$ is a necessary condition
for intersection. We will show that it is also sufficient. Squaring
both sides of (C) and simplifying yields $2c^2x^4+x^2-9=0$. Using the
quadratic formula to solve this for $x^2$ yields
$x^2=\dfrac{-1+\sqrt{1+72c^2}}{4c^2}$. The condition $x^2<9$ holds if
and only if $-1+\sqrt{1+72c^2}<36c^4$, which is equivalent to
$1+72c^2<(1+36c^2)^2=1+72c^2+1296c^4$, which holds for all $c>0$.


\item
The angles $\theta$ and $\theta_1$  from the $x$-axis to the tangents
to $C$ and $C_1$ satisfy $\tan\theta=f(x_0,y_0)$ and
$\tan\theta_1=\dfrac{f(x_0,y_0)+\tan\alpha}{1-f(x_0,y_0)\tan\alpha}=
\dst\frac{\tan\theta+\tan\alpha}{1-\tan\theta\tan\alpha}=\tan(\theta+\alpha)$.
Therefore, assuming $\theta$ and $\theta_1$ are both in $[0,2\pi)$,
$\theta_1=\theta+\alpha$.

\item
Circles centered at the origin are given by $x^2+y^2=r^2$.
Differentiating yields $2x+2yy'=0$, so $y'=-\dst\frac{x}{ y}$ is a
differential equation for the given family, and
$y'=\dfrac{-(x/y)+\tan\alpha}{1+(x/y)\tan\alpha}$ is a differential
equation for the desired family. Substituting $y=ux$ yields
$u'x+u=\dfrac{-1/u+\tan\alpha}{1+(1/u)\tan\alpha}=\dst\frac{-1+u\tan\alpha}{
u+\tan\alpha}$. Therefore,$u'x=-\dst\frac{1+u^2}{ u+\tan\alpha}$,
$\dst\frac{u+\tan\alpha}{ 1+u^2}u'=-\dst\frac{1}{ x}$ and
$\dst\frac{1}{2}\ln(1+u^2)+\tan\alpha\tan^{-1}u=-\ln|x|+k$. Substituting
$u=\dst\frac{y}{ x}$ yields $\dst{\frac{1}{2}\ln (x^2+y^2)+(\tan\alpha)
\tan^{-1}\frac{y}{ x}=k}$.
\end{sectionSolutions}

\chapter{Linear Second Order Equations}

\section{Homogeneous Linear Equations}

\begin{sectionSolutions}
\item
\begin{alist}
\item If $y_1=e^x\cos x$, then $y_1'=e^x(\cos x -\sin x)$  and
$y_1''=e^x(\cos x-\sin x-\sin x-\cos x)=-2e^x\sin x$, so
$y_1''-2y_1'+2y_1=e^x(-2\sin x-2\cos x+2\sin x+2\cos x)=0$. If
$y_2=e^x\sin x$, then
$y_2'=e^x(\sin x+\cos x)$ and
$y_2''=e^x(\sin x+\cos x+\cos x-\sin x)=2e^x\cos x$, so
$y_2''-2y_2'+2y_2=e^x(2\cos x-2\sin x-2\cos x+2\sin x)=0$.

\item If (B) $y=e^x(c_1\cos x+c_2\sin x)$, then
\[
y'=e^x\left(c_1(\cos x-\sin x)+c_2(\sin x+\cos x)\right)
\tag{C}
\]
 and
\begin{align*}
y''&=c_1e^x(\cos x-\sin x-\sin x-\cos x)\\
&+c_2e^x(\sin x+\cos x+\cos x-\sin x)\\
&=2e^x(-c_1\sin x+c_2\cos x),
\end{align*}
 so
\begin{align*}
y''-2y'+2y&=c_1e^x(-2\sin x-2\cos x+2\sin x+2\cos x)\\
&+c_2e^x(2\cos x-2\sin x-2\cos x+2\sin x)=0.
\end{align*}

\item We must choose $c_1$ and $c_2$ in (B) so that $y(0)=3$
and $y'(0)=-2$. Setting $x=0$ in (B) and (C) shows that
$c_1=3$ and $c_1+c_2=-2$, so $c_2=-5$. Therefore,
 $y=e^x(3\cos x-5\sin x)$.

\item We must choose $c_1$ and $c_2$ in (B) so that $y(0)=k_0$
and $y'(0)=k_1$. Setting $x=0$ in (B) and (C) shows that
$c_1=k_0$ and $c_1+c_2=k_1$, so $c_2=k_1-k_0$.  Therefore,
$y=e^x\left(k_0\cos x+(k_1-k_0)\sin x\right)$.
\end{alist}

\item
\begin{alist}
\item If $y_1=\dst\frac{1}{ x-1}$, then $y_1'=-\dst\frac{1}{(x-1)^2}$
and
$y_1''=\dst\frac{2}{(x-1)^3}$, so
\begin{align*}
(x^2-1)y_1''+4xy_1'+2y_1&=\dfrac{2(x^2-1)}{(x-1)^3}-\dst\frac{4x}{(x-1)^2}
+\dst\frac{2}{ x-1}\\[2\jot]
&=\dfrac{2(x+1)-4x+2(x-1)}{(x-1)^2}=0.
\end{align*}
Similar manipulations show that
$(x^2-1)y_2''+4xy_2'+2y_2=0$. The general solution on each of
the intervals $(-\infty,-1)$, $(-1,1)$, and $(1,\infty)$ is
(B) $y=\dst\frac{c_1}{ x-1}+\dst\frac{c_2}{ x+1}$.

\item Differentiating (B) yields (C)
$y'=-\dst\frac{c_1}{(x-1)^2}-\dst\frac{c_2}{(x+1)^2}$.
 We must choose $c_1$ and $c_2$ in (B) so that $y(0)=-5$
and $y'(0)=1$. Setting $x=0$ in (B) and (C) shows that
$-c_1+c_2=-5,\ -c_1-c_2=1$. Therefore,$c_1=2$ and $c_2=-3$,
so $y=\dst{\frac{2}{ x-1}-\frac{3}{ x+1}}$ on $(-1,1)$.

\stepcounter{enumXi}
\item
The Wronskian of  $\{y_1,y_2\}$ is
\[
W(x)=\begin{vmatrix}
\dst \phantom{-}\frac{1}{x-1} & \dst \phantom{-}\frac{1}{x+1}\\
\dst -\frac{1}{(x-1)^2} & \dst - \frac{1}{(x+1)^2}
\end{vmatrix}
=\frac{2}{(x^2-1)^2},
\tag{D}
\]
so $W(0)=2$. Since
$p(x)=\dst\frac{4x}{ x^2-1}$, so $\dst{\int_0^x p(t)\,dt}=
\dst{\int_0^x \frac{4t}{ t^2-1}\,dt}=\ln(x^2-1)^2$, Abel's
formula implies that
$W(x)=W(0)e^{-\ln(x^2-1)^2}=\dst\frac{2}{(x^2-1)^2}$, consistent
with (D).
\end{alist}

\item From Abel's formula,
$W(x)=W(\pi)e^{-3\int_\pi^x (t^2+1)\,dt}=0\cdot
e^{-3\int_\pi^x (t^2+1)\,dt}=0$.


\item
$p(x)=\dst\frac{1}{ x}$; therefore $\dst{\int_1^x
p(t)\,dt}=\dst{\int_1^x\frac{dt}{ t}}=\ln x$, so Abel's
formula yields $W(x)=W(1)e^{-\ln x}=\dst\frac{1}{ x}$.

\item
$p(x)=-2$;\;
$P(x)=-2x$;\;
$y_2=uy_1=ue^{3x}$;\;
$u'=\dfrac{Ke^{-P(x)}}{ y_1^2(x)}=\dfrac{Ke^{2x}}{
e^{6x}}=Ke^{-4x}$;\;
$u=-\dst\frac{K}{4}e^{-4x}$.
Choose $K=-4$;  then
$y_2=e^{-4x}e^{3x}=e^{-x}$.


\item
$p(x)=-2a$;\;
$P(x)=-2ax$;\;
$y_2=uy_1=ue^{ax}$;\;
$u'=\dfrac{Ke^{-P(x)}}{ y_1^2(x)}=\dfrac{Ke^{2ax}}{ e^{2ax}}=K$;\;
$u=Kx$.
Choose $K=1$; then
$y_2=xe^{ax}$.

\item
$p(x)=-\dst\frac{1}{ x}$;\;
$P(x)=-\ln x$;\;
$y_2=uy_1=ux$;\;
$u'=\dfrac{Ke^{-P(x)}}{ y_1^2(x)}=\dst\frac{Kx}{
x^2}=\dst\frac{K}{ x}$;\;
$u=K\ln x$.
Choose $K=1$; then
$y_2=x\ln x$.


\item
$p(x)=-\dst\frac{1}{ x}$;\;
$P(x)=-\ln|x|$;\;
$y_2=uy_1=ux^{1/2}e^{2x}$;\;
$u'=\dfrac{Ke^{-P(x)}}{ y_1^2(x)}=\dst\frac{Kx}{ xe^{4x}}=e^{-4x}$;\;
$u=-\dfrac{Ke^{-4x}}{4}$.
Choose $K=-4$; then
$y_2=e^{-4x}(x^{1/2}e^{2x})=x^{1/2}e^{-2x}$.


\item
$p(x)=-\dst\frac{2}{ x}$;\;
$P(x)=-2\ln|x|$;\;
$y_2=uy_1=ux\cos x$;\;
$u'=\dfrac{Ke^{-P(x)}}{ y_1^2(x)}=\dst\frac{Kx^2}{ x^2\cos^2
x}=K\sec^2x$;\;
$u=K\tan x$.
Choose $K=1$; then
$y_2=\tan x(x\cos x)=x\sin x$.

\item
$p(x)=-\dst\frac{3x+2}{3x-1}=-1-\dst\frac{3}{3x-1}$;\;
$P(x)=-x-\ln|3x-1|$;\;
$y_2=uy_1=ue^{2x}$;\;
$u'=\dfrac{Ke^{-P(x)}}{ y_1^2(x)}=\dfrac{K(3x-1)e^x}{
e^{4x}}=K(3x-1)e^{-3x}$;\;
$u=-Kxe^{-3x}$.
Choose $K=-1$; then
$y_2=xe^{-3x}e^{2x}=xe^{-x}$.


\item
$p(x)=-\dfrac{2(2x^2-1)}{
x(2x+1)}=-2-\dst\frac{2}{2x+1}+\dst\frac{2}{ x}$;\;
$P(x)=-2x-\ln|2x+1|+2\ln|x|$;\;
$y_2=uy_1=\dst\frac{u}{ x}$;\;
$u'=\dfrac{Ke^{-P(x)}}{ y_1^2(x)}=\dfrac{K(2x+1)e^{2x}}{
x^2}x^2=K(2x+1)e^{2x}$;\;
$u=Kxe^{2x}$.
Choose $K=1$; then
$y_2=\dfrac{xe^{2x}}{ x}=e^{2x}$.


\item
Suppose that $y\equiv0$ on $(a,b)$. Then $y'\equiv0$ and $y''\equiv0$
on $(a,b)$, so $y$ is a solution of (A) $y''+p(x)y'+q(x)y=0,\
y(x_0)=0,\ y'(x_0)=0$ on $(a,b)$. Since Theorem~5.1.1
implies that (A) has only one solution on $(a,b)$, the conclusion
follows.


\item
If $\{z_1,z_2\}$ is a fundamental set of solutions of
(A) on $(a,b)$, then every solution $y$ of
(A) on $(a,b)$ is a linear combination of
$\{z_1,z_2\}$; that is, $y=c_1z_1+c_2z_2=c_1(\alpha y_1+\beta
y_2)+c_2(\gamma y_1+\delta y_2)
=(c_1\alpha+c_2\gamma)y_1+(c_1\beta+c_2\delta)y_2$, which shows that
every solution of (A) on $(a,b)$ can be written as
a linear combination of $\{y_1,y_2\}$. Therefore,$\{y_1,y_2\}$ is a
fundamental set of solutions of (A) on $(a,b)$.


\item
The Wronskian of  $\{y_1,y_2\}$ is
\[
W=\begin{vmatrix}
y_1 & y_2 \\
y'_1 & y'_2
\end{vmatrix}=
\begin{vmatrix}
y_1 & ky_1 \\
y'_1 & ky'_1
\end{vmatrix}=k(y_1y_1'-y_1'y_1)=0.
\]

nor $y_2$ can be a solution of $y''+p(x)y'+q(x)y=0$ on  $(a,b)$.

\item
$W(x_0)=\left(y_1(x_0)y_2'(x_0)-y_1'(x_0)y_2(x_0)\right)=0$
if either $y_1(x_0)=y_2(x_0)=0$ or
$y_1'(x_0)=y_2'(x_0)=0$, and  Theorem~5.1.6
implies that  $\{y_1,y_2\}$ is linearly dependent on  $(a,b)$.


\item
Let $x_0$ be an arbitrary point in  $(a,b)$. By the motivating
argument preceding  Theorem~5.1.4,
(B) $W(x_0)=y_1(x_0)y_2'(x_0)-y_1'(x_0)y_2(x_0)\ne0$. Now let $y$ be
the solution of $y''+p(x)y'+q(x)y=0,\ y(x_0)=y_1(x_0),\
y'(x_0)=y_1'(x_0)$. By assumption,  $y$ is a linear combination
of
 $\{y_1,y_2\}$ on  $(a,b)$; that is,  $y=c_1y_1+c_2y_2$, where
\begin{align*}
c_1y_1(x_0)+c_2y_2(x_0)&=y_1(x_0)\\
c_1y_1'(x_0)+c_2y_2'(x_0)&=y_1'(x_0).
\end{align*}
Solving this system by Cramer's rule yields
\[
c_1=\frac{1}{ W(x_0)}
\begin{vmatrix}
y_1(x_0) & y_2(x_0) \\
y_1'(x_0) & y'_2(x_0)
\end{vmatrix}=1
\text{ and }
c_2=\frac{1}{ W(x_0)}
\begin{vmatrix}
y_1(x_0) & y_1(x_0)\\
y'_1(x_0) &y_1(x_0)
\end{vmatrix}=0.
\]
Therefore,$y=y_1$, which
shows that $y_1$ is a solution of (A). A similar argument shows that
$y_2$ is a solution of (A).


\item
Expanding  the determinant by cofactors of its first column shows that
the first equation in the exercise can be written as
\[
\frac{y}{ W}\begin{vmatrix}
y_1' & y_2' \\
y_1'' & y_2''
\end{vmatrix}-
\frac{y'}{ W}\begin{vmatrix}
y_1 & y_2 \\
y_1'' & y_2''
\end{vmatrix}+
\frac{y''}{ W}\begin{vmatrix}
y_1 & y_2 \\
y'_1 & y'_2
\end{vmatrix}=0,
\]
which is of the form (A) with
\[
p=-\frac{1}{ W}\begin{vmatrix}
y_1 & y_2 \\
y_1'' & y_2''
\end{vmatrix}
\quad\text{and}\quad
q=\frac{1}{ W}\begin{vmatrix}
y_1' & y_2' \\
y_1'' & y_2''
\end{vmatrix}.
\]


\item Theorem~5.1.6 implies that
that there are constants $c_1$ and $c_2$ such that
(B) $y=c_1y_1+c_2y_2$
on $(a,b)$. To see that $c_1$ and $c_2$ are unique, assume that (B)
holds, and let $x_0$ be a point in  $(a,b)$. Then (C)
$y'=c_1y_1'+c_2y_2'$. Setting $x=x_0$ in (B) and (C) yields
\begin{align*}
c_1y_1(x_0)+c_2y_2(x_0)&=y(x_0)\\
c_1y_1'(x_0)+c_2y_2'(x_0)&=y'(x_0).
\end{align*}
Since  Theorem~5.1.6  implies that
$y_1(x_0)y_2'(x_0)-y_1'(x_0)y_2(x_0)\ne0$, the argument preceding
 Theorem~5.1.4 implies that $c_1$ and $c_2$ are given
uniquely by
\[
c_1=\frac{y_2'(x_0)y(x_0)-y_2(x_0)y'(x_0)}{
y_1(x_0)y_2'(x_0)-y_1'(x_0)y_2(x_0)}\qquad
c_2=\frac{y_1(x_0)y'(x_0)-y_1'(x_0)y(x_0)}{
y_1(x_0)y_2'(x_0)-y_1'(x_0)y_2(x_0)}\,.
\]


\item
The general solution of $y''=0$ is $y=c_1+c_2x$, so $y'=c_2$. Imposing
the stated initial conditions on $y_1=c_1+c_2x$  yields $c_1+c_2x_0=1$
and $c_2=0$; therefore $c_1=1$, so $y_1=1$.
Imposing the stated initial conditions on $y_2=c_1+c_2x$  yields
$c_1+c_2x_0=0$ and $c_2=1$; therefore $c_1=-x_0$, so $y_2=x-x_0$.
The solution of the general initial value problem is
$y=k_0+k_1(x-x_0)$.


\item
Let $y_1=a_1\cos\omega x+a_2\sin\omega x$ and
 $y_2=b_1\cos\omega x+b_2\sin\omega x$. Then
\begin{align*}
\phantom{(-}a_1\cos\omega x_0+a_2\sin\omega x_0\phantom{)}&=1\\
\omega(-a_1\sin\omega x_0+a_2\cos\omega x_0)&=0
\end{align*}
and
\begin{align*}
\phantom{(-}b_1\cos\omega x_0+b_2\sin\omega x_0\phantom{)}&=0\\
\omega(-b_1\sin\omega x_0+b_2\cos\omega x_0)&=1.
\end{align*}
Solving these systems yields $a_1=\cos\omega x_0$,
$a_2=\sin\omega x_0$, $b_1=-\dst\frac{\sin\omega x_0}{\omega}$, and
$b_2=\dst\frac{\cos\omega x_0}{\omega}$. Therefore,
$y_1=\cos\omega x_0\cos\omega x+\sin\omega x_0\sin\omega
x=\cos\omega(x-x_0)$ and
 $y_2=\dst\frac{1}{\omega}(-\sin\omega x_0\cos\omega x+\cos\omega x_0
\sin\omega x)=\dst\frac{1}{\omega}\sin\omega(x-x_0)$.
 The solution of the general
initial value problem is
$y=k_0\cos\omega(x-x_0)+\dst\frac{k_1}{\omega}\sin\omega(x-x_0)$.


\item
\begin{alist}
\item
 If $y_1=x^2$, then $y_1'=2x$ and $y_1''=2$, so
$x^2y_1''-4xy_1'+6y_1=x^2(2)-4x(2x)+6x^2=0$
for $x$ in $(-\infty,\infty)$.
 If $y_2=x^3$, then $y_2'=3x^2$ and $y_2''=6x$, so
$x^2y_2''-4xy_2'+6y_2=x^2(6x)-4x(3x^2)+6x^3=0$
for $x$ in $(-\infty,\infty)$. If $x\ne0$, then $y_2(x)/y_1(x)=x$,
which is nonconstant on $(-\infty,0)$ and $(0,\infty)$, so
Theorem~5.1.6 implies that  $\{y_1,y_2\}$ is a fundamental
set of solutions of (A) on each of these intervals.

\item Theorem~5.1.6 and \part{a} imply that $y$
satisfies (A) on $(-\infty,0)$ and on $(0,\infty)$ if and only if
$y=\dst{\begin{cases}
a_1x^2+a_2x^3,&x> 0,\\
b_1x^2+b_2x^3,&x<0.
\end{cases}}$
 Since $y(0)=0$ we can  complete the proof that $y$ is
a solution of
(A) on $(-\infty,\infty)$ by showing
that $y'(0)$ and $y''(0)$ both exist if and only if $a_1=b_1$. Since
\[
\dst{\frac{y(x)-y(0)}{
x-0}=\begin{cases}
a_1x+a_2x^2,&\text{
if }x>0,\\ b_1x+b_2x^2,&\text{ if }x<0,
\end{cases}}
\]
 it follows
that $y'(0)=\dst{\lim_{x\to0}\frac{y(x)-y(0)}{ x-0}}=0$. Therefore,
$y'=\dst{\begin{cases}
2a_1x+3a_2x^2,&x\ge 0,\\
2b_1x+3b_2x^2,&x<0.
\end{cases}}$
Since
$\dfrac{y'(x)-y'(0)}{
x-0}=\dst{\begin{cases}
2a_1+3a_2x,&\text{ if }x>0,\\
2b_1+3b_2x,&\text{ if }x<0,
\end{cases}}$ it
follows
that $y''(0)=\dst{\lim_{x\to0}\frac{y'(x)-y'(0)}{ x-0}}$ exists
if and only if $a_1=b_1$.  By renaming $a_1=b_1=c_1$, $a_2=c_2$,
and $b_2=c_3$  we see that $y$ is a solution of
(A) on $(-\infty,\infty)$ if and only if
$y=\dst{\begin{cases}
c_1x^2+c_2x^3,&x\ge 0,\\
c_1x^2+c_3x^3,&x<0.
\end{cases}}$

\item We have shown that $y(0)=y'(0)=0$ for any choice
of $c_1$ and $c_2$ in (C). Therefore,the given initial value problem
has a solution if and only if $k_0=k_1=0$, in which case every
function of the form (C) is a solution.

\item  If $x_0>0$, then $c_1$ and $c_2$ in (C) are uniquely
determined by $k_0$ and $k_1$, but $c_3$ can be chosen arbitrarily.
Therefore,(B) has a unique solution on
$(0,\infty)$, but infinitely many solutions on $(-\infty,\infty)$.
  If $x_0<0$, then $c_1$ and $c_3$ in (C) are uniquely
determined by $k_0$ and $k_1$, but $c_2$ can be chosen arbitrarily.
Therefore,(B) has a unique solution on
$(-\infty,0)$, but infinitely many solutions on $(-\infty,\infty)$.
\end{alist}


\item
\begin{alist}
\item
 If $y_1=x^3$, then $y_1'=3x^2$ and $y_1''=6x$, so
$x^2y_1''-6xy_1'+12y_1=x^2(6x)-6x(3x^2)+12x^3=0$
for $x$ in $(-\infty,\infty)$.
 If $y_2=x^4$, then $y_2'=4x^3$ and $y_2''=12x^2$, so
$x^2y_2''-6xy_2'+12y_2=x^2(12x^2)-6x(4x^3)+12x^4=0$
for $x$ in $(-\infty,\infty)$.
If $x\ne0$, then $y_2(x)/y_1(x)=x$, which
is nonconstant on $(-\infty,0)$ and $(0,\infty)$, so
Theorem~5.1.6 implies that  $\{y_1,y_2\}$ is a fundamental
set of solutions of (A) on each of these intervals.

\item
 Theorem~5.1.2 and \part{a} imply that $y$ satisfies (A) on
$(-\infty,0)$ and on $(0,\infty)$ if and only if
(C) $y=\dst{\begin{cases}
a_1x^3+a_2x^4,&x> 0,\\
b_1x^3+b_2x^4,&x<0.
\end{cases}}$
Since $y(0)=0$ we can  complete the proof that $y$ is
a solution of
(A) on $(-\infty,\infty)$ by showing
that $y'(0)$ and $y''(0)$ both exist for any choice of $a_1$, $a_2$,
$b_1$, and $b_2$. Since
$\dst{\frac{y(x)-y(0)}{
x-0}=\begin{cases}
a_1x^2+a_2x^3,&\text{ if }x>0,\\
b_1x^2+b_2x^3,&\text{ if }x<0,
\end{cases}}$ it
follows
that $y'(0)=\dst{\lim_{x\to0}\frac{y(x)-y(0)}{ x-0}}=0$. Therefore,
$y'=\dst{\begin{cases}
3a_1x^2+4a_2x^3,&x\ge 0,\\
3b_1x^2+4b_2x^3,&x<0.
\end{cases}}$
Since
$\dfrac{y'(x)-y'(0)}{
x-0}=\dst{\begin{cases}
3a_1x+4a_2x^2,&\text{ if }x>0,\\
3b_1x+4b_2x^2,&\text{ if }x<0,
\end{cases}}$ it
follows that $y''(0)=\dst{\lim_{x\to0}\frac{y'(x)-y'(0)}{ x-0}}=0$.
Therefore,(B) is a solution of (A) on $(-\infty,\infty)$.

\item We have shown that $y(0)=y'(0)=0$ for any choice
of $a_1$, $a_2$, $b_1$, and $b_2$ in (B). Therefore,the given initial
value
problem has a solution if and only if $k_0=k_1=0$, in which case every
function of the form (B) is a solution.

\item  If $x_0>0$, then $a_1$ and $a_2$ in (B) are uniquely
determined by $k_0$ and $k_1$, but $b_1$ and $b_2$ can be chosen
arbitrarily. Therefore,(C) has a unique solution on
$(0,\infty)$, but infinitely many solutions on $(-\infty,\infty)$.
  If $x_0<0$, then $b_1$ and $b_2$ in (B) are uniquely
determined by $k_0$ and $k_1$, but $a_1$ and $a_2$ can be chosen
arbitrarily. Therefore,(C) has a unique solution on
$(-\infty,0)$, but infinitely many solutions on $(-\infty,\infty)$.
\end{alist}

\end{sectionSolutions}

\section{Constant Coefficient Homogeneous Equations}

\begin{sectionSolutions}
\item
$p(r)=r^2-4r+5=(r-2)^2+1$;\;
$y=e^{2x}(c_1 \cos x+c_2 \sin x)$.

\item
$p(r)=r^2-4r+4=(r-2)^2$;\;
 $y=e^{2x}(c_1+c_2x)$.

\item
$p(r)=r^2+6r+10=(r+3)^2+1$;\;
$y=e^{-3x}(c_1 \cos x+c_2 \sin x)$.


\item
$p(r)=r^2+r=r(r+1)$;\;
 $y=c_1+c_2e^{-x}$.


\item
$p(r)=r^2+6r+13y=(r+3)^2+4$;\;
 $y=e^{-3x}(c_1\cos 2x+c_2\sin 2x)$.


\item
$p(r)=10r^2-3r-1=(2r-1)(5r+1)=10(r-1/2)
(r+1/5)$;\;
$y=c_1e^{-x/5}+c_2e^{x/2}$.


\item
$p(r)=6r^2-r-1=(2r-1)(3r+1)=6(r-1/2) (r+1/3)$;\;
$y=c_1e^{-x/3}+c_2e^{x/2}$;\;
$y'=-\dst\frac{c_1}{3}e^{-x/3}+\dst\frac{c_2}{2}e^{x/2}$;\;
 $y(0)=10\Rightarrow c_1+c_2=10,\ y'(0)=0\Rightarrow -\dst\frac{c_1}{3}
+\dst\frac{c_2}{2}=0$;\ $c_1=6, c_2=4$;\;
$y=4e^{x/2}+6e^{-x/3}$.



\item
$p(r)=4r^2-4r-3=(2r-3)(2r+1)= 4(r-3/2)(r+1/2)$;
$y=c_1e^{-x/2} +c_2e^{3x/2}$;\;
$y'=-\dst\frac{c_1}{2}e^{-x/2} +\dst\frac{3c_2}{2}e^{3x/2}$;\;
 $y(0)=\dst\frac{13}{ 12}\Rightarrow c_1+c_2=\dst\frac{13}{12},\
y'(0)=\dst\frac{23 }{ 24}\Rightarrow
-\dst\frac{c_1}{2}+\dst\frac{3c_2}{2}=\dst\frac{23}{24}$;\;
$c_1=\dst\frac{1}{3}, c_2=\dst\frac{3}{4}$;\;
$y=\dfrac{e^{-x/2}}{ 3}+\dfrac{3e^{3x/2}}{4}$.


\item
$p(r)=r^2+7r+12=(r+3)(r+4)$;\;
$y=c_1e^{-4x}+c_2e^{-3x}$;\;
$y'=-4c_1e^{-4x}-3c_2e^{-3x}$;\;
 $y(0)=-1\Rightarrow c_1+c_2=-1,\ y'(0)=0\Rightarrow -4c_1-3c_2=0$;\;
$c_1=3$, $c_2=-4$;\;
$y=3e^{-4x}-4e^{-3x}$.






\item
$p(r)=36r^2-12r+1=(6r-1)^2=36(r-1/6)^2$;\;
$y=e^{x/6}(c_1+c_2x)$;\;
$y'=\dfrac{e^{x/6}}{6}(c_1+c_2x)+c_2e^{x/6}$;\;
$y(0)=3\Rightarrow c_1=3,\ y'(0)=\dst\frac{5}{2}\Rightarrow
\dst\frac{c_1}{6}+c_2=\dst\frac{5}{2}\Rightarrow c_2=2$;\;
$y=e^{x/6}(3+2x)$.



\item
\begin{alist}
\item
From (A), $ay''(x)+by'(x)+cy(x)=0$ for all $x$. Replacing $x$ by
$x-x_0$ yields (C) $ay''(x-x_0)+by'(x-x_0)+cy(x-x_0)=0$. If
$z(x)=y(x-x_0)$, then the chain rule implies that $z'(x)=y'(x-x_0)$
and $z''(x)=y''(x-x_0)$, so (C) is equivalent to $az''+bz'+cz=0$.


\item  If  $\{y_1,y_2\}$  is a fundamental set of solutions of
(A) then Theorem~5.1.6  implies that
$y_2/y_1$ is nonconstant. Therefore,
$\dfrac{z_2(x)}{ z_1(x)}=\dfrac{y_2(x-x_0)}{ y_1(x-x_0)}$ is also
nonconstant, so Theorem~5.1.6 implies that  $\{z_1,z_2\}$ is a
fundamental set of solutions of (A).

\item
\it Let $p(r)=ar^2+br+c$ be the characteristic polynomial of (A).
Then:
\begin{itemize}
\item
If $p(r)=0$  has distinct real roots $r_1$ and $r_2$, then the general
solution of (A) is
\[
y=c_1e^{r_1(x-x_0)}+c_2e^{r_2(x-x_0)}.
\]
\item
If $p(r)=0$  has a repeated root  $r_1$, then
the general solution of (A) is
\[
y=e^{r_1(x-x_0)}(c_1+c_2(x-x_0)).
\]
\item % (c)
If $p(r)=0$  has complex conjugate roots $r_1=\lambda+i\omega$
and
$r_2=\lambda-i\omega$ $($where $\omega>0)$, then the general solution
of (A) is
\[
y=e^{\lambda (x-x_0)}(c_1\cos\omega(x-x_0)+c_2\sin\omega(x-x_0)).
\]
\end{itemize}
\end{alist}

\item
$p(r)=r^2-6r-7=(r-7)(r+1)$;
\begin{align*}
y&=c_1e^{-(x-2)}+c_2e^{7(x-2)};\\
y'&=-c_1e^{-(x-2)}+7c_2e^{7(x-2)};
\end{align*}
$y(2)=-\dst\frac{1}{3}\Rightarrow c_1+c_2=-\dst\frac{1}{3},\
y'(2)=-5\Rightarrow -c_1+7c_2=-5$;\;
$c_1=\dst\frac{1}{3}, c_2=-\dst\frac{2}{3}$;\;
$y=\dst\frac{1}{3}e^{-(x-2)}-\frac{2}{3}e^{7(x-2)}$.

\item
$p(r)=9r^2+6r+1=(3r+1)^2=9(r+1/3)^2$;
\begin{align*}
y&=e^{-(x-2)/3}\left(c_1+c_2(x-2)\right);\\
y'&=-\frac{1}{3}e^{-(x-2)/3}\left(c_1+c_2(x-2)\right)
+c_2e^{-(x-2)/3};
\end{align*}
$y(2)=2\Rightarrow c_1=2,\
y'(2)=-\dst\frac{14}{3}\Rightarrow -\dst\frac{c_1}{3}+c_2=-\dst\frac{14}{3}
\Rightarrow c_2=-4$;\;
$y=e^{-(x-2)/3}\left(2-4(x-2)\right)$.

\item
$p(r)=r^2+3$;
\begin{align*}
y&=c_1\dst{\cos \sqrt3\left(x-\frac{\pi}{3}\right)+c_2\sin
\sqrt3\left(x-\frac{\pi}{3}\right)};\\
y'&=-\sqrt3c_1\dst{\sin
\sqrt3\left(x-\frac{\pi}{3}\right)+\sqrt3c_2\cos
\sqrt3\left(x-\frac{\pi}{3}\right)};
\end{align*}
$y(\pi/3)=2\Rightarrow c_1=2,\ y'(\pi/3)=-1\Rightarrow
c_2=-\dst\frac{1}{\sqrt3}$;\;
\[
y=2\cos \sqrt3\left(x-\frac{\pi}{3}\right)-\frac{1}{\sqrt3}\sin
\sqrt3\left(x-\frac{\pi}{3}\right).
\]


\item
$y$ is a solution of $ay''+by'+cy=0$  if and only if
\begin{align*}
y&=\phantom{r_1}c_1e^{r_1(x-x_0)}+\phantom{r_2}e^{r_2(x-x_0)}\\
y'&=r_1c_1e^{r_1(x-x_0)}+r_2e^{r_2}(x-x_0).
\end{align*}
Now $y_1(x_0)=k_0$ and $y_1'(x_0)=k_1\Rightarrow c_1+c_2=k_0,
r_1c_1+r_2c_2=k_1$. Therefore,$c_1=\dst\frac{r_2k_0-k_1}{ r_2-r_1}$
and $c_2=\dst\frac{k_1-r_1k_0}{ r_2-r_1}$. Substituting $c_1$ and $c_2$
into  the above equations for $y$ and y' yields
\begin{align*}
y&=\frac{r_2k_0-k_1}{ r_2-r_1}e^{r_1(x-x_0)}+
\frac{k_1-r_1k_0}{ r_2-r_1}e^{r_2(x-x_0)}\\[2\jot]
&=\frac{k_0}{
r_2-r_1}\left(r_2e^{r_1(x-x_0)}-r_1e^{r_2(x-x_0)}\right)+\frac{k_1}{
r_2-r_1} \left(e^{r_2(x-x_0)}-e^{r_1(x-x_0)}\right).
\end{align*}



\item
\setcounter{equation}{0}
$y$ is a solution of $ay''+by'+cy=0$  if and only if
\[
y=\phantom{\lambda}e^{\lambda(x-x_0)}\left(c_1\cos\omega(x-x_0)+c_2\sin\omega(x-x_0)\right)
\tag{A}
\]
and
\begin{align*}
y'&=\lambda
e^{\lambda(x-x_0)}\left(c_1\cos\omega(x-x_0)+c_2\sin\omega(x-x_0)\right)\\
&
+\omega
e^{\lambda(x-x_0)}\left(-c_1\sin\omega(x-x_0)+c_2\cos\omega(x-x_0)\right).
\end{align*}
Now $y_1(x_0)=k_0\Rightarrow c_1=k_0$ and $y_1'(x_0)=k_1\Rightarrow
\lambda c_1+\omega c_2=k_1$,  so $c_2=\dfrac{k_1-\lambda
k_0}{\omega}$.
 Substituting $c_1$ and $c_2$
into
(A) yields
\[
y=e^{\lambda(x-x_0)}\left[k_0\cos\omega(x-x_0)+\left(\frac{k_1-\lambda k_0
}{\omega}\right)\sin\omega(x-x_0)\right].
\]

\item
\begin{alist}[start=2]
\item
\begin{align*}
e^{i\theta_1}e^{i\theta_2}&=(\cos\theta_1+i\sin\theta_1)
(\cos\theta_2+i\sin\theta_2)\\
&=(\cos\theta_1\cos\theta_2-\sin\theta_1\sin\theta_2)
+i(\sin\theta_1\cos\theta_2+\cos\theta_1\sin\theta_2)\\
&=\cos(\theta_1+\theta_2)+i\sin(\theta_1+\theta_2)=e^{i(\theta_1+\theta_2)}.
\end{align*}

\item
\begin{align*}
e^{z_1+z_2}&=e^{(\alpha_1+i\beta_1)+(\alpha_2+i\beta_2)}
=e^{(\alpha_1+\alpha_2)+i(\beta_1+\beta_2)}\\
&=e^{(\alpha_1+\alpha_2)}e^{i(\beta_1+\beta_2)} \text{ (from
(F) with $\alpha=\alpha_1+\alpha_2$ and $\beta=\beta_1+\beta_2$)}\\
&=e^{\alpha_1}e^{\alpha_2}e^{i(\beta_1+\beta_2)}
\text{ (property of the real-valued exponential function)}\\
&=e^{\alpha_1}e^{\alpha_2}e^{i\beta_1}e^{i\beta_2}
\text{ (from \part{b})}\\
&=e^{\alpha_1}e^{i\beta_1}e^{\alpha_2}e^{i\beta_2}
=e^{\alpha_1+i\beta_1}e^{\alpha_2+i\beta_2}=e^{z_1}e^{z_2}.
\end{align*}

\item
The real and imaginary parts of $z_1=e^{(\lambda+i\omega)x}$ are
$u_1=e^{\lambda x}\cos\omega x$ and $v_1=e^{\lambda x}\sin\omega x$,
which are both solutions of $ay''+by'+cy=0$, by
Theorem~5.2.1(c). Similarly, the real and imaginary parts of
$z_2=e^{(\lambda-i\omega)x}$ are $u_2=e^{\lambda x}\cos(-\omega
x)=e^{\lambda x}\cos\omega x$ and $v_1=e^{\lambda x}\sin(-\omega
x)=-e^{\lambda x}\sin\omega x$, which are both solutions of
$ay''+by'+cy=0$, by Theorem~5.2.1,(c).
\end{alist}
\end{sectionSolutions}

\section{Nonhomogeneous Linear Equations}

\begin{sectionSolutions}
\item
The characteristic polynomial of the complementary equation is
$p(r)=r^2-4r+5=(r-2)^2+1$, so
$\{e^{2x}\cos x,e^{2x}\sin x\}$
is a fundamental set of solutions for the complementary equation.
Let $y_p=A+Bx$; then
$y_p''-4y_p'+5y_p=-4B+5(A+Bx)=1+5x$.
Therefore,$5B=5,\  -4B+5A=1$, so $B=1$, $A=1$.
Therefore,$y_p=1+x$ is a particular solution and
$y=1+x+e^{2x}(c_1\cos x+c_2\sin x)$ is the general solution.

\item
The characteristic polynomial of the complementary equation is
$p(r)=r^2-4r+4=(r-2)^2$, so  $\{e^{2x},xe^{2x}\}$
is a fundamental set of solutions for the complementary equation.
Let $y_p=A+Bx+Cx^2$; then
$y_p''-4y_p'+4y_p=2C-4(B+2Cx)+4(A+Bx+Cx^2)=(2C-4B+4A)+(-8C+4B)x+4Cx^2
=2+8x-4x^2$. Therefore,$4C=-4,\ -8C+4B=8,\
2C-4B+4A=2$, so $C=-1$, $B=0$, and $A=1$.
Therefore,$y_p=1-x^2$ is a particular solution and
 $y=1-x^2+e^{2x}(c_1+c_2x)$ is the general solution.


\item
The characteristic polynomial of the complementary equation is
$p(r)=r^2+6r+10=(r+3)^2+1$, so
$\{e^{-3x}\cos x,e^{-3x}\sin x\}$
is a fundamental set of solutions for the complementary equation.
Let $y_p=A+Bx$; then
$y_p''+6y_p'+10y_p=6B+10(A+Bx)=22+20x$.
Therefore,$10B=20,\  6B+10A=22$, so $B=2$, $A=1$.
Therefore,$y_p=1+2x$ is a particular solution and
(A) $y=1+2x+e^{-3x}(c_1\cos x+c_2\sin x)$ is the general solution.
Now $y(0)=2\Rightarrow 2=1+c_1\Rightarrow c_1=1$. Differentiating (A)
yields
$y'=2-3e^{-3x}(c_1\cos x+c_2\sin x) +e^{-3x}(-c_1\sin x+c_2\cos x)$,
so $y'(0)=-2\Rightarrow
-2=2-3c_1+c_2\Rightarrow c_2=-1$.
 $y=1+2x+e^{-3x}(\cos x-\sin x)$ is the solution of the initial
value problem.



\item
If $y_p=\dst\frac{A}{ x}$, then $x^2y_p''+7xy_p'+8y_p=
A\dst{\left(x^2\left(\frac{2}{ x^3}\right)+7x\left(\frac{-1}{
x^2}\right)+\left(\frac{8}{ x}\right)\right)}=\dst\frac{3A}{ x}=\dst\frac{6}{
x}$ if $A=2$. Therefore,$y_p=\dst\frac{2}{ x}$ is a particular solution.


\item
If $y_p=Ax^3$, then  $x^2y_p''-xy_p'+y_p=
A\left(x^2(6x)-x(3x^2)+x^3\right)=4Ax^3=2x^3$ if $A=\dst\frac{1}{2}$.
 Therefore,$y_p=\dst\frac{x^3}{2}$ is a particular solution.


\item
If $y_p=Ax^{1/3}$, then $x^2y_p''+xy_p'+y_p=
A\dst{\left(x^2\left(\frac{-2x^{-5/3}}{9}\right)+x\left(\frac{x^{-2/3}}{3}\right)
+x^{1/3}\right)}=\dst\frac{10A}{9}x^{1/3}=10x^{1/3}$
if $A=9$. Therefore,$y_p=9x^{1/3}$ is a particular solution.


\item
If $y_p=\dst\frac{A}{ x^3}$, then  $x^2y_p''+3xy_p'-3y_p=
A\dst{\left(x^2\left(\frac{12}{ x^5}\right)+3x\left(\frac{-3}{ x^4}\right)
+\frac{3}{ x^3}\right)}=0$. Therefore,$y_p$ is not a solution of the
given equation for any choice of $A$.


\item
The characteristic polynomial of the complementary equation is
$p(r)=r^2+5r-6=(r+6)(r-1)$, so
 $\{e^{-6x},e^{x}\}$
is a fundamental set of solutions for the complementary equation.
Let $y_p=Ae^{3x}$; then
$y_p''+5y_p'-6y_p=p(3)Ae^{3x}=18Ae^{3x}=6e^{3x}$ if
$A=\dst\frac{1}{3}$.
Therefore,$y_p=\dfrac{e^{3x}}{3}$ is a particular solution and
 $y=\dfrac{e^{3x}}{3}+c_1e^{-6x}+c_2e^{x}$ is the general solution.


\item
The characteristic polynomial of the complementary equation is
$p(r)=r^2+8r+7=(r+1)(r+7)$, so
 $\{e^{-7x},e^{-x}\}$
is a fundamental set of solutions for the complementary equation.
Let $y_p=Ae^{-2x}$; then
$y_p''+8y_p'+7y_p=p(-2)Ae^{-2x}=-5Ae^{-2x}=10e^{-2x}$ if $A=-2$.
Therefore,$y_p=-2e^{-2x}$ is a particular solution  and
(A) $y=-2e^{-2x}+c_1e^{-7x}+c_2e^{-x}$ is the general solution.
Differentiating (A) yields
 $y'=4e^{-2x}-7c_1e^{-7x}-c_2e^{-x}$.
Now $y(0)=-2 \Rightarrow -2=-2+c_1+c_2$ and $y'(0)=10\Rightarrow
10=4-7c_1-c_2$. Therefore,$c_1=-1$ and $c_2=1$, so
$y=-2e^{-2x}-e^{-7x}+e^{-x}$
 is the solution of the initial value problem.


\item
The characteristic polynomial of the complementary equation is
$p(r)=r^2+2r+10=(r+1)^2+9$, so
$\{e^{-x}\cos3x,e^{-x}\sin3x\}$
is a fundamental set of solutions for the complementary equation.
If $y_p=Ae^{x/2}$, then
$y_p''+2y_p'+10y_p=p(1/2)Ae^{x/2}=\dst\frac{45}{4}Ae^{x/2}=e^{x/2}$
if $A=\dst\frac{4}{45}$.
Therefore, $y_p=\dst\frac{4}{45}e^{x/2}$ is a particular solution and
$y=\dst\frac{4}{45}e^{x/2}+e^{-x}(c_1\cos3x+c_2\sin3x)$
 is the general solution.



\item
The characteristic polynomial of the complementary equation is
$p(r)=r^2-7r+12=(r-4)(r-3)$. If $y_p=Ae^{4x}$, then
$y_p''-7y_p'+12y_p=p(4)Ae^{4x}=0\cdot e^{4x}=0$, so
$y_p''-7y_p'+12y_p\ne5e^{4x}$ for any choice of $A$.


\item
The characteristic polynomial of the complementary equation is
$p(r)=r^2-8r+16=(r-4)^2$, so
 $\{e^{4x},xe^{4x}\}$
is a fundamental set of solutions for the complementary equation.
If $y_p=A\cos x+B\sin x$, then
$y_p''-8y_p'+16y_p=-(A\cos x+B\sin x)-8(-A\sin x+B\cos x)+
16(A\cos x+B\sin x)=(15A-8B)\cos x+(8A+15B)\sin x$, so
$15A-8B=23,\ 8A+15B=-7$, which implies that $A=1$ and $B=-1$.
Hence $y_p=\cos x-\sin x$ and
 $y=\cos x-\sin x+e^{4x}(c_1+c_2x)$ is the general solution.


\item
The characteristic polynomial of the complementary equation is
$p(r)=r^2-2r+3=(r-1)^2+2$, so
$\{e^x\cos \sqrt{2}x,e^{x}\sin \sqrt{2} x\}$
is a fundamental set of solutions for the complementary equation.
If $y_p=A\cos3x+B\sin3x$, then
$y_p''-2y_p'+3y_p=-9(A\cos3x+B\sin3x)-6(-A\sin3x+B\cos3x)+
3(A\cos3x+B\sin3x)=-(6A+6B)\cos3x+(6A-6B)\sin3x$, so
$-6A-6B=-6,\ 6A-6B=6$, which implies that $A=1$ and $B=0$.
Hence $y_p=\cos3x$ is a particular solution and
$y=\cos3x+e^x(c_1\cos \sqrt{2}x+c_2\sin \sqrt{2} x)$
 is the general solution.


\item
The characteristic polynomial of the complementary equation is
$p(r)=r^2+7r+12=(r+3)(r+4)$, so $\{e^{-4x},e^{-3x}\}$ is a fundamental
set of solutions for the complementary equation. If
$y_p=A\cos2x+B\sin2x$, then
$y_p''+7y_p'+12y_p=-4(A\cos2x+B\sin2x)+14(-A\sin2x+B\cos2x)+ 12(A\cos
x+B\sin x)=(8A+14B)\cos2x+(8B-14a)\sin2x$, so $8A+14B=-2,\
-14A+8B=36$, which implies that $A=-2$ and $B=1$. Hence
$y_p=-2\cos2x+\sin2x$ is a particular solution and (A)
$y=-2\cos2x+\sin2x+c_1e^{-4x}+c_2e^{-3x}$ is the general solution.
Differentiating (A) yields
$y'=2\sin2x+2\cos2x-4c_1e^{-4x}-3c_2e^{-3x}$. Now $y(0)=-3 \Rightarrow
-3=-2+c_1+c_2$ and $y'(0)=3\Rightarrow 3=2-4c_1-3c_2$. Therefore,
$c_1=2$ and $c_2=-3$, so $y=-2\cos2x+\sin2x+2e^{-4x}-3e^{-3x}$ is the
solution of the initial value problem.


\item
$\{\cos\omega_0x,\sin\omega_0x\}$ is a fundamental set of
solutions of the complementary equation.
If $y_p=A\cos\omega x+B\sin\omega x$, then $y_p''+\omega_0^2y_p=
-\omega^2(A\cos\omega x+B\sin\omega x)+\omega_0^2(A\cos\omega
x+B\sin\omega x)=(\omega_0^2-\omega^2)(A\cos\omega x+B\sin\omega x)
=M\cos\omega x+N\sin\omega x$ if
$A=\dst\frac{M}{\omega_0^2-\omega^2}$ and
$B=\dst\frac{N}{\omega_0^2-\omega^2}$. Therefore,
\[
y_p= \dst\frac{1}{\omega_0^2-\omega^2}(M\cos\omega
x+N\sin\omega x)
\]
is a particular solution of the given equation and
\[
y=\dst\frac{1}{\omega_0^2-\omega^2}(M\cos\omega x+N\sin\omega
x)+c_1\cos\omega_0x+c_2\sin\omega_0x
\]
 is the general solution.


\item
If $y_p=A\cos\omega x+B\sin\omega x$, then
$ay_p''+by_p'+cy_p=-a\omega^2(A\cos\omega
x+B\sin\omega x)+b\omega(-A\sin\omega x+B\cos\omega x)+c(A\cos\omega
x+B\sin\omega x)=\left[(c-a\omega^2)A+b\omega B\right]\cos\omega
x +\left[-b\omega A+(c-a\omega^2)B\right]\sin\omega x$. Therefore,
$y_p$ is a solution of (A) if and only if the set of equations (B)
$(c-a\omega^2)A+b\omega B=M,\ -b\omega A+(c-a\omega^2)B=N$
has a solution. If $(c-a\omega^2)^2+(b\omega)^2\ne0$, then (B)
has the solution
$A=\dfrac{(c-a\omega^2)M-b\omega N}{(c-a\omega^2)^2+(b\omega)^2}$,
$B=\dfrac{(c-a\omega^2)N+b\omega M}{(c-a\omega^2)^2+(b\omega)^2}$,
 and $y_p=A\cos\omega x+B\sin\omega x$ is a solution of (A).
If $(c-a\omega^2)^2+(b\omega)^2=0$  (which is true if and only if
the left side of (A) is of the form $a(y''+\omega^2y)$, then
the coefficients of $A$ and $B$ in (B) are all zero, so (B) does not
have a solution, so (A) does not have a solution of the
form $y_p=A\cos\omega x+B\sin\omega x$.


\item
From Exercises~5.3.2 and 5.3.17,
$y_{p_1}=1+x$ and $y_{p_2}=e^{2x}$  are particular solutions of
$y''-4y'+5y=1+5x$
and
$y''-4y'+5y=e^{2x}$
respectively, and $\{e^{2x}\cos x,e^{2x}\sin x\}$
is a fundamental set of solutions of the  complementary equation.
Therefore,$y_p=y_{p_1}+y_{p_2}=1+x+e^{2x}$
is a particular solution of the given equation, and
$y=1+x+e^{2x}(1+c_1\cos x+c_2\sin x)$
is the general solution.


\item
From Exercises~5.3.4 and 5.3.19,
$y_{p_1}=1-x^2$ and $y_{p_2}=e^{x}$  are particular solutions of
$y''-4y'+4y=2+8x-4x^2$
and
$y''-4y'+4y=e^{x}$
respectively, and $\{e^{2x},xe^{2x}\}$
is a fundamental set of solutions of the complementary equation.
Therefore,$y_p=y_{p_1}+y_{p_2}=1-x^2+e^{x}$
is a particular solution of the given equation, and
$y=1-x^2+e^{x}+e^{2x}(c_1+c_2x)$
is the general solution.


\item
From Exercises~5.3.6 and 5.3.21,
$y_{p_1}=1+2x$ and $y_{p_2}=e^{-3x}$  are particular solutions of
$y''+6y'+10y=22+20x$
and
$y''+6y'+10y=e^{-3x}$
respectively, and $\{e^{-3x}\cos x,e^{-3x}\sin x\}$
is a fundamental set of solutions of the complementary equation.
Therefore,$y_p=y_{p_1}+y_{p_2}=1+2x+e^{-3x}$
is a particular solution of the given equation, and
$y=1+2x+e^{-3x}(1+c_1\cos x+c_2\sin x)$
is the general solution.

\item
Letting $c_1=c_2=0$ shows that (A) $y_p''+p(x)y_p'+q(x)y_p=f$.
Letting $c_1=1$ and $c_2=0$ shows that  (B) $(y_1+y_p)''+
p(x)(y_1+y_p)'+q(x)(y_1+y_p)=f$. Now subtract (A) from (B)
to see that $y_1''+p(x)y_1'+q(x)y_1=0$.
Letting $c_1=0$ and $c_2=1$ shows that  (C) $(y_2+y_p)''+
p(x)(y_2+y_p)'+q(x)(y_2+y_p)=f$. Now subtract (A) from (C)
to see that $y_2''+p(x)y_2'+q(x)y_2=0$.
\end{sectionSolutions}

\section{The Method of Undetermined Coefficients I}

\begin{sectionSolutions}
\item
If $y=ue^{-3x}$, then
$y''-6y'+5y=e^{-3x}\left[(u''-6u'+9u)-6(u'-3u)+5u\right]=
e^{-3x}(35-8x)$, so $u''-12u'+32u =35-8x$ and $u_p=A+Bx$, where
$-12B+32(A+Bx)=35-8x$. Therefore,$32B=-8$, $32A-12B=35$, so
$B=-\dst\frac{1}{4}$, $A=1$, and $u_p=\dst{1-\frac{x}{4}}$. Therefore,
$y_p=\dst{e^{-3x}\left(1-\frac{x}{4}\right)}$.

\item
If $y=ue^{2x}$, then
$y''+2y'+y=e^{2x}\left[(u''+4u'+4u)+2(u'+2u)+u\right]=
e^{2x}(-7-15x+9x^2)$ so $u''+6u'+9u =-7-15x+9x^2$ and $u_p=A+Bx+Cx^2$,
where $2C+6(B+2Cx)+9(A+Bx+Cx^2)=-7-15x+9x^2$. Therefore,$9C=9$,
$9B+12C=-15$, $9A+6B+2C=-7$, so $C=1$, $B=-3$, $A=1$, and
$u_p=1-3x+x^2$. Therefore,$y_p=e^{2x}(1-3x+x^2)$.


\item
If $y=ue^x$, then $y''-y'-2y=e^x\left[(u''+2u'+u)-(u'+u)-2u\right]=
e^x(9+2x-4x^2)$ so $u''+u'-2u =9+2x-4x^2$, and $u_p=A+Bx+Cx^2$, where
$2C+(B+2Cx)-2(A+Bx+Cx^2)=9+2x-4x^2$. Therefore,$-2C=-4$, $-2B+2C=2$,
$-2A+B+2C=9$, so $C=2$, $B=1$, $A=-2$, and $u_p=-2+x+2x^2$. Therefore,
$y_p=e^x(-2+x+2x^2)$.


\item
If $y=ue^x$, then $y''-3y'+2y=e^x\left[(u''+2u'+u)-3(u'+u)+2u\right]=
e^x(3-4x)$, so $u''-u'=3-4x$ and $u_p=Ax+Bx^2$, where
$2B-(A+2Bx)=3-4x$. Therefore,$-2B=-4$, $-A+2B=3$, so $B=2$, $A=1$, and
$u_p=x(1+2x)$. Therefore,$y_p=xe^x(1+2x)$.


\item
If $y=ue^{2x}$, then
$2y''-3y'-2y=e^{2x}\left[2(u''+4u'+4u)-3(u'+2u)-2u\right]=
e^{2x}(-6+10x)$, so $2u''+5u'=-6+10x$ and $u_p=Ax+Bx^2$, where
$2(2B)+5(A+2Bx)=-6+10x$. Therefore,$10B=10$, $5A+4B=-6$, so $B=1$,
$A=-2$, and $u_p=x(-2+x)$. Therefore,$y_p=xe^{2x}(-2+x)$.


\item
If $y=ue^x$, then $y''-2y'+y=e^x\left[(u''+2u'+u)-2(u'+u)+u\right]=
e^x(1-6x)$, so $u''=1-6x$ Integrating twice and taking the constants
of integration to be zero yields
$u_p=\dst{x^2\left(\frac{1}{2}-x\right)}$. Therefore,
$y_p=\dst{x^2e^x\left(\frac{1}{2}-x\right)}$.


\item
If $y=ue^{-x/3}$, then
$9y''+6y'+y=e^{-x/3}\dst{\left[9\left(u''-\frac{2u'}{3}+\frac{u}{9}\right)
+6\left(u'-\frac{u}{3}\right)+u\right]}= e^{-x/3}(2-4x+4x^2)$, so
$9u''=2-4x+4x^2$, or $u''=\dst\frac{1}{9}(2-4x+4x^2)$. Integrating twice
and taking the constants of integration to be zero yields
$u_p=\dst{\frac{x^2}{27}(3-2x+x^2)}$. Therefore,
$y_p=\dst{\frac{x^2e^{-x/3}}{27}(3-2x+x^2)}$.

\item
If $y=ue^x$, then $y''-6y'+8y=e^x\left[(u''+2u'+u)-6(u'+u)+8u\right]=
e^x(11-6x)$, so $u''-4u'+3u =11-6x$ and $u_p=A+Bx$, where
$-4B+3(A+Bx)=11-6x$. Therefore,$3B=-6$, $3A-4B=11$, so $B=-2$, $A=1$
and $u_p=1-2x$. Therefore,$y_p=e^x(1-2x)$. The characteristic
polynomial of the complementary equation is
$p(r)=r^2-6r+8=(r-2)(r-4)$, so $\{e^{2x},e^{4x}\}$ is a fundamental
set of solutions of the complementary equation. Therefore,
$y=e^x(1-2x)+c_1e^{2x}+c_2e^{4x}$ is the general solution of the
nonhomogeneous equation.


\item
If $y=ue^x$, then $y''+2y'-3y=e^x\left[(u''+2u'+u)+2(u'+u)-3u\right]=
-16xe^x$, so $u''+4u'=-16x$ and $u_p=Ax+Bx^2$, where
$2B+4(A+2Bx)=-16x$. Therefore,$8B=-16$, $4A+2B=0$, so $B=-2$, $A=1$,
and $u_p=x(1-2x)$. Therefore,$y_p=xe^x(1-2x)$. The characteristic
polynomial of the complementary equation is
$p(r)=r^2+2r-3=(r+3)(r-1)$, so $\{e^x,e^{-3x}\}$ is a fundamental set
of solutions of the complementary equation. Therefore,
$y=xe^x(1-2x)+c_1e^x+c_2e^{-3x}$ is the general solution of the
nonhomogeneous equation.


\item
If $y=ue^{2x}$, then
$y''-4y'-5y=e^{2x}\left[(u''+4u'+4u)-4(u'+2u)-5u\right]=
9e^{2x}(1+x)$, so $u''-9u =9+9x$ and $u_p=A+Bx$, where
$-9(A+Bx)=9+9x$. Therefore,$-9B=-9$, $-9A=9$, so $B=-1$, $A=-1$, and
$u_p=-1-x$. Therefore,$y_p=-e^{2x}(1+x)$. The characteristic
polynomial of the complementary equation is
$p(r)=r^2-4r-5=(r-5)(r+1)$, so $\{e^{-x},e^{5x}\}$ is a fundamental
set of solutions of the complementary equation. Therefore,(A)
$y=-e^{2x}(1+x)+c_1e^{-x}+c_2e^{5x}$ is the general solution of the
nonhomogeneous equation. Differentiating (A) yields
$y'=-2e^{2x}(1+x)-e^{2x}-c_1e^{-x}+5c_2e^{5x}$. Now $y(0)=0,\
y'(0)=-10\Rightarrow 0=-1+c_1+c_2,\ -10=-3-c_1+5c_2$, so $c_1=2$,
$c_2=-1$. Therefore,$y=-e^{2x}(1+x)+2e^{-x}-e^{5x}$ is the solution of
the initial value problem.


\item
If $y=ue^{-x}$, then
$y''+4y'+3y=e^{-x}\left[(u''-2u'+u)+4(u'-u)+3u\right]= -e^{-x}(2+8x)$,
so $u''+2u'=-2-8x$ and $u_p=Ax+Bx^2$, where $2B+2(A+2Bx)=-2-8x$.
Therefore,$4B=-8$, $2A+2B=-2$, so $B=-2$, $A=1$, and $u_p=x(1-2x)$.
Therefore,$y_p=xe^{-x}(1-2x)$. The characteristic polynomial of the
complementary equation is $p(r)=r^2+4r+3=(r+3)(r+1)$, so
$\{e^{-x},e^{-3x}\}$ is a fundamental set of solutions of the
complementary equation. Therefore,(A)
$y=xe^{-x}(1-2x)+c_1e^{-x}+c_2e^{-3x}$ is the general solution of the
nonhomogeneous equation. Differentiating (A) yields
$y'=-xe^{-x}(1-2x)+e^{-x}(1-4x)-c_1e^{-x}-3c_2e^{-3x}$. Now $y(0)=1,\
y'(0)=2\Rightarrow 1=c_1+c_2,\ 2=1-c_1-3c_2$, so $c_1=2$, $c_2=-1$.
Therefore,$y=e^{-x}(2+x-2x^2)-e^{-3x}$ is the solution of the initial
value problem.


\item
We must find particular solutions $y_{p_1}$ and $y_{p_2}$ of (A)
$y''+y'+y=xe^x$ and (B) $y''+y'+y=e^{-x}(1+2x)$, respectively. To find
a particular solution of (A) we write $y=ue^x$. Then
$y''+y'+y=e^x\left[(u''+2u'+u)+(u'+u)+u\right]= xe^x$ so $u''+3u'+3u
=x$ and $u_p=A+Bx$, where $3B+3(A+Bx)=x$. Therefore,$3B=1$, $3A+3B=0$,
so $B=\dst\frac{1}{3}$, $A=-\dst\frac{1}{3}$, and $u_p=-\dst\frac{1}{3}(1-x)
$, so $y_{p_1}=-\dst\frac{e^x}{3}(1-x)$. To find a particular solution
of (B) we write $y=ue^{-x}$. Then
$y''+y'+y=e^{-x}\left[(u''-2u'+u)+(u'-u)+u\right]= e^{-x}(1+2x)$, so
$u''-u'+u =1+2x$ and $u_p=A+Bx$, where $-B+(A+Bx)=1+2x$. Therefore,
$B=2$, $A-B=1$, so $A=3$, and $u_p=2+3x$, so $y_{p_2}=e^{-x}(3+2x)$.
Now $y_p=y_{p_1}+y_{p_2}=-\dst{\frac{e^x}{3}(1-x)+e^{-x}(3+2x)}$.


\item
We must find particular solutions $y_{p_1}$ and $y_{p_2}$ of (A)
$y''-8y'+16y=6xe^{4x}$ and (B) $y''-8y'+16y=2+16x+16x^2$,
respectively. To find a particular solution of (A) we write
$y=ue^{4x}$. Then
$y''-8y'+16y=e^x\left[(u''+8u'+16u)-8(u'+4u)+16u\right]= 6xe^{4x}$, so
$u'' =6x$, $u_p=x^3$. and $y_{p_1}=x^3e^{4x}$. To find a particular
solution of (B) we write $y_p=A+Bx+Cx^2$. Then
$y_p''-8y_p'+16y_p=2C-8(B+2Cx)+16(A+Bx+Cx^2)=(16A-8B+2C)
+(16B-16C)x+16Cx^2=2+16x+16x^2$ if $16C=16$, $16B-16C=16$,
$16A-8B+2C=2$. Therefore,$C=1$, $B=2$, $A=1$, and $y_{p_2}=1+2x+x^2$.
Now $y_p=y_{p_1}+y_{p_2}=x^3e^{4x}+1+2x+x^2$.


\item
We must find particular solutions $y_{p_1}$ and $y_{p_2}$ of (A)
$y''-2y'+2y=e^x(1+x)$ and (B) $y''-2y'+2y=e^{-x}(2-8x+5x^2)$,
respectively. To find a particular solution of (A) we write $y=ue^x$.
Then $y''-2y'+2y=e^x\left[(u''+2u'+u)-2(u'+u)+2u\right]= e^x(1+x)$, so
$u''+u=1+x$ and $u_p=1+x$, so $y_{p_1}=e^x(1+x)$. To find a particular
solution of (B) we write $y=ue^{-x}$. Then
$y''-2y'+2y=e^{-x}\left[(u''-2u'+u)-2(u'-u)+2u\right]=
e^{-x}(2-8x+5x^2)$, so $u''-4u'+5u =2-8x+5x^2$ and $u_p=A+Bx+Cx^2$,
where $2C-4(B+2Cx)+5(A+Bx+Cx^2)=2-8x+5x^2$. Therefore,$5C=5$,
$5B-8C=-8$, $5A-4B+2C=2$, so $C=1$, $B=0$, $A=0$, and $u_p=x^2$.
Therefore,$y_{p_2}=x^2e^{-x}$. Now
$y_p=y_{p_1}+y_{p_2}=e^x(1+x)+x^2e^{-x}$.


\item
\begin{alist}
\item
If $y=ue^{\alpha x}$, then $ay''+by'+cy=e^{\alpha x}\left[a(u''+2\alpha
u'+\alpha^2u)+b(u'+\alpha u)+cu\right]= e^{\alpha x}
\left[au''+(2a\alpha+b))u'+(a\alpha^2+b\alpha+c)u\right] =e^{\alpha
x}(au''+p'(\alpha)u'+p(\alpha)u)$. Therefore,$ay''+by'+cy=e^{\alpha
x}G(x)$ if and only if $au''+p'(\alpha)u'+p(\alpha)u=G(x)$.

\item
Substituting $u_p=A+Bx+Cx^2+Dx^3$ into (B) yields
\begin{multline*}
a(2C+6Dx)+p'(\alpha)(B+2Cx+3Dx^2)+p(\alpha)(A+Bx+Cx^2+Dx^3)\\
=[p(\alpha)A+p'(\alpha)B+2aC]+[p(\alpha)B+2p'(\alpha)C+6aD]x\\
+[p(\alpha)C +3p'(\alpha)D]x^2+p(\alpha)Dx^3=g_0+g_1x+g_2x^2+g_3x^3
\end{multline*}
if
\[
\begin{aligned}
p(\alpha)D&=g_3\\
p(\alpha)C+3p'(\alpha)D&=g_2\\
p(\alpha)B+2p'(\alpha)C+6aD&=g_1\\
p(\alpha)A+p'(\alpha)B+2aC\phantom{+6aD}&=g_0.
\end{aligned}
\tag{C}
\]
Since $e^{\alpha x}$ is not a solution of the complementary equation,
$p(\alpha)\ne0$. Therefore,the triangular system
(C) can be solved successively for $D$, $C$, $B$
and $A$.

\item
Since $e^{\alpha x}$ is a solution of the complementary equation while
$xe^{\alpha x}$ is not, $p(\alpha)=0$ and $p'(\alpha)\ne0$. Therefore,
(B) reduces to  (D) $au''+p'(\alpha)u=G(x)$.
Substituting $u_p=Ax+Bx^2+Cx^3+Dx^4$ into (D) yields
\begin{align*}
&a(2B+6Cx+12Dx^2)+p'(\alpha)(A+2Bx+3Cx^2+4Dx^3)\\
&=(p'(\alpha)A+2aB)+(2p'(\alpha)B+6aC)x
+(3p'(\alpha)C +12aD)x^2\\
&\quad+4p'(\alpha)Dx^3=g_0+g_1x+g_2x^2+g_3x^3
\end{align*}
 if
\begin{align*}
4p'(\alpha)D&=g_3\\
3p'(\alpha)C+12aD&=g_2\\
2p'(\alpha)B+6aC\phantom{+12aD}&=g_1\\
p'(\alpha)A+2aB\phantom{+6aC+12aD}&=g_0.
\end{align*}
Since $p'(\alpha)\ne0$ this triangular system can be solved
successively for $D$, $C$, $B$ and $A$.

\item
Since $e^{\alpha x}$ and $xe^{\alpha x}$ are solutions of the
complementary equation,
 $p(\alpha)=0$ and $p'(\alpha)=0$. Therefore,
(B) reduces to  (D) $au''=G(x)$, so $u''=\dfrac{G(x)}{ a}$.
Integrating this twice and taking the constants of integration yields
the particular solution $u_p=x^2\dst{\left(\frac{g_0}{2}+
\frac{g_1}{6}x+\frac{g_2}{12}x^2+\frac{g_3}{20}x^3\right)}$.
\end{alist}


\item
If $y_p=Axe^{4x}$, then
$y_p''-7y_p'+12y_p=[(8+16x)-7(1+4x)+12x]Ae^{4x}=Ae^{4x}=5e^{4x}$
if $A=1$, so $y_p=5xe^{4x}$.


\item
If $y_p=e^{3x}(A+Bx+Cx^2)$, then
\begin{align*}
y_p''-3y_p'+2y_p&=e^{3x}[(9A+6B+2C)+(9B+12C)x+9Cx^2]\\
&\quad-3e^{3x}[(3A+B)+(3B+2C)x+3Cx^2]\\
&\quad+2e^{3x}(A+Bx+Cx^2)\\
&=e^{3x}[(2A+3B+2C)+(2B+6C)x+2Cx^2]\\
&=e^{3x}(-1+2x+x^2)
\end{align*}
if $2C=1,\ 2B+6C=2,\ 2A+3B+2C= -1$. Therefore,$C=\dst\frac{1}{2}$,
$B=-\dst\frac{1}{2}$, $A=-\dst\frac{1}{4}$, and
$y_p=-\dfrac{e^{3x}}{4}(1+2x-2x^2)$.


\item
If $y_p=e^{-x/2}(Ax^2+Bx^3+Cx^4)$, then
\begin{align*}
4y_p''+4y_p'+y_p&=e^{-x/2}[8A-(8A-24B)x+(A-12B+48C)x^2]\\
&\quad+e^{-x/2}[(B-16C)x^3+Cx^4]\\
&+e^{-x/2}[8Ax-(2A-12B)x^2-(2B-16C)x^3-2Cx^4]\\
&\quad+e^{-x/2}(Ax^2+Bx^3+Cx^4)\\
&=e^{-x/2}(8A+24Bx+48Cx^2)=e^{-x/2}(-8+48x+144x^2)
\end{align*}
if $48C=144$, $24B=48$, and $8A=-8$. Therefore,$C=3$, $B=2$, $A=-1$,
and $y_p=x^2e^{-x/2}(-1+2x+3x^2)$.


\item If $y=\int e^{\alpha x}P(x)\,dx$, then
$y'=e^{\alpha x}P(x)$. Let $y=ue^{\alpha x}$; then $(u'+\alpha
u)e^{\alpha x}=e^{\alpha x}P(x)$, which implies (A). We must show that
it is possible to choose $A_0,\dots,A_k$ so that (B)
$(A_0+A_1x\cdots+A_kx^k)'+\alpha(A_0+A_1x\cdots+A_kx^k)=
p_0+p_1x+\cdots+p_kx^k$. By equating the coefficients of $x^k,
x^{k-1},\dots,1$ (in that order) on the two sides of (B), we see that
(B) holds if and only if $\alpha A_k=p_k$ and $(k-j+1)A_{k-j+1}+\alpha
A_k=p_{k-j},\ 1\le j\le k$.


\item
If $y=\int x^ke^{\alpha x}\,dx$, then $y'=x^ke^{\alpha x}$. Let
$y=ue^{\alpha x}$; then $(u'+\alpha u)e^{\alpha x}=x^ke^{\alpha x}$,
so $u'+\alpha u=x^k$. This equation has a
particular solution $u_p=A_0+A_1x\cdots+A_kx^k$, where (A)
$(A_0+A_1x\cdots+A_kx^k)'+\alpha(A_0+A_1x\cdots+A_kx^k)=
x^k$. By equating the coefficients of $x^k,
x^{k-1},\dots,1$  on the two sides of (A), we see that
(A) holds if and only if $\alpha A_k=1$ and
$(k-j+1)A_{k-j+1}+\alpha
A_k-j=0,\ 1\le j\le k$.  Therefore, $A_k=\dst\frac{1}{\alpha}$,
$A_{k-1}=-\dst\frac{k}{\alpha^2}$, $A_{k-2}=\dfrac{k(k-1)}{\alpha^3}$,
and, in general,
$A_{k-j}=(-1)^j\dfrac{k(k-1)\cdots(k-j+1)}{\alpha^{j+1}}=
\dfrac{(-1)^jk!}{\alpha^{j+1}(k-j)!},\
1\le
j\le k$. By introducing the index $r=k-j$ we can rewrite this as
$A_r=\dfrac{(-1)^{k-r}k!}{\alpha^{k-r+1}r!},\ 0\le r\le k$.
Therefore,
$u_p=\dst{\frac{(-1)^kk!}{ \alpha^{k+1}}\sum_{r=0}^k\frac{(-\alpha
x)^r}{ r!}}$ and
$y=\dst{\frac{(-1)^kk!e^{\alpha x}}{ \alpha^{k+1}}\sum_{r=0}^k\frac{(-\alpha
x)^r}{ r!}+c}$.
\end{sectionSolutions}

\section{The Method of Undetermined Coefficients II}

\begin{sectionSolutions}
\item
Let
\begin{align*}
y_p&=(A_0+A_1x)\cos x+(B_0+B_1x)\sin x; \text{ then}\\
y_p'&=(A_1+B_0+B_1x)\cos x+(B_1-A_0-A_1x)\sin x\\
y_p''&=(2B_1-A_0-A_1x)\cos x-(2A_1+B_0+B_1x)\sin x, \text{ so}\\
y_p''+3y_p'+y_p&=(3A_1+3B_0+2B_1+3B_1x)\cos x \\
&+(3B_1-3A_0-2A_1-3A_1x)\sin x \\
&=(2-6x)\cos x-9\sin x
\end{align*}
if $3B_1=-6$, $-3A_1=0$, $3B_0+3A_1+2B_1=2$, $-3A_0+3B_1+2A_1=-9$.
Therefore,$A_1=0$,   $B_1=-2$, $A_0=1$, $B_0=2$, and
 $y_p=\cos x+(2-2x)\sin x$.


\item
Let $y=ue^{2x}$. Then
\begin{align*}
y''+3y'-2y&=e^{2x}\left[(u''+4u'+4u)+3(u'+2u)-2u\right]\\
&=e^{2x}(u''+7u'+8u)=-e^{2x}(5\cos2x+9\sin2x)
\end{align*}
if $u''+7u'+8u=-5\cos2x-9\sin2x$. Now let $u_p=A\cos2x+B\sin2x$. Then
\begin{align*}
u_p''+7u_p'+8u_p&=-4(A\cos2x+B\sin2x)+14(-A\sin2x+B\cos2x)\\
&+8(A\cos2x+B\sin2x)\\
&=(4A+14B)\cos2x-(14A-4B)\sin2x\\
&=-5\cos2x-9\sin2x
\end{align*}
if $4A+14B=-5$, $-14A+4B=-9$. Therefore,$A=\dst\frac{1}{2}$,
$B=-\dst\frac{1}{2}$,
\[
 u_p=\dst{\frac{1}{2}(\cos2x-\sin2x)}, \text{ and }
y_p=\dst{\frac{e^{2x}}{2}(\cos2x-\sin2x)}.
\]


\item
Let $y=ue^{-2x}$. Then
\begin{align*}
y''+3y'-2y&=e^{-2x}\left[(u''-4u'+4u)+3(u'-2u)-2u\right]\\
&=e^{-2x}(u''-u'-4u)\\
&=e^{-2x}\left[(4+20x)\cos 3x+(26-32x)\sin 3x\right]
\end{align*}
if $u''-u'-4u=(4+20x)\cos 3x+(26-32x)\sin 3x$.
Let
\begin{align*}
u_p&=(A_0+A_1x)\cos3x+(B_0+B_1x)\sin3x; \text{ then}\\
u_p'&=(A_1+3B_0+3B_1x)\cos3x+(B_1-3A_0-3A_1x)\sin3x\\
u_p''&=(6B_1-9A_0-9A_1x)\cos3x-(2A_1+9B_0+9B_1x)\sin3x,\text{ so}\\
u_p''-u_p'-4u_p&=-\left[13A_0+A_1+3B_0-6B_1+(13A_1+3B_1)x\right]\cos3x\\
&-\left[13B_0+B_1-3A_0+6A_1+(13B_1-3A_1)x\right]\sin3x\\
&=(4+20x)\cos 3x+(26-32x)\sin 3x\qquad\text{if}\\
&\begin{cases}
-13A_1-\phantom{1}3B_1&=20\\ \phantom{-1}3A_1-13B_1&=-32
\end{cases}
\quad\text{and}\quad
\begin{cases}{rcr}
-13A_0-\phantom{1}3B_0-\phantom{6}A_1+6B_1&=4\\
\phantom{-1}3A_0-13B_0-6A_1-\phantom{6}B_1&=26.
\end{cases}
\end{align*}
From the first two equations, $A_1=-2$, $B_1=2$. Substituting these in
the last two equations yields $-13A_0-3B_0=-10$, $3A_0-13B_0=16$.
Solving this pair yields $A_0=1$, $B_0=-1$. Therefore,
\[
 u_p=(1-2x)(\cos 3x-\sin 3x) \text{ and }
y_p=e^{-2x}(1-2x)(\cos 3x-\sin 3x).
\]


\item
Let
\begin{align*}
y_p&=(A_0x+A_1x^2)\cos x+(B_0x+B_1x^2)\sin x;\text{ then}\\
y_p'&=\left[A_0+(2A_1+B_0)x+B_1x^2\right]\cos x
+\left[B_0+(2B_1-A_0)x-B_1x^2\right]\sin x\\
y_p''&=
\left[2A_1+2B_0-(A_0-4B_1)x-A_1x^2\right]\cos x\\
&+\left[2B_1-2A_0-(B_0+4A_1)x-B_1x^2\right]\sin x, \text{ so}\\
y_p''+y_p&=(2A_1+2B_0+4B_1x)\cos x+(2B_1-2A_0-4A_1x)\sin x\\
&=(-4+8x)\cos x+(8-4x)\sin x
\end{align*}
if $4B_1=8$, $-4A_1=-4$, $2B_0+2A_1=-4$, $-2A_0+2B_1=8$. Therefore,
$A_1=1$, $B_1=2$, $A_0=-2$, $B_0=-3$, and $y_p=-x\left[(2-x)\cos
x+(3-2x)\sin x\right]$.



\item
Let $y=ue^{-x}$. Then
\begin{align*}
y''+2y'+2y&=e^{-x}\left[(u''-2u'+u)+2(u'-u)+2u\right]\\
&=e^{-x}(u''+u)= e^{-x}(8\cos x-6\sin x)
\end{align*}
 if $u''+u=8\cos x-6\sin x$.
Now let
\begin{align*}
u_p&=Ax\cos x+Bx\sin x; \text{ then}\\
u_p'&=(A+Bx)\cos x+(B-Ax)\sin x\\
u_p''&=(2B-Ax)\cos x-(2A+Bx)\sin x, \text{ so}\\
u_p''+u_p&=2B\cos x-2A\sin x=8\cos x-6\sin x
\end{align*}
if $2B=8$, $-2A=-6$. Therefore,$A=3$, $B=4$, $u_p=x(3\cos x+4\sin x)$,
and $y_p=xe^{-x}(3\cos x+4\sin x)$.


\item
Let
\begin{align*}
y_p&=(A_0+A_1x+A_2x^2)\cos x +(B_0+B_1x+B_2x^2)\sin x;\text{ then}\\
y_p'&=\left[A_1+B_0+(2A_2+B_1)x+B_2x^2\right]\cos x\\
&+\left[B_1-A_0+(2B_2-A_1)x-A_2x^2\right]\sin x, \\
y_p''&=\left[-A_0+2A_2+2B_1-(A_1-4B_2)x-A_2x^2\right]\cos x\\
&+\left[-B_0+2B_2-2A_1-(B_1+4A_2)x-B_2x^2\right]\sin x,\text{ so}\\
y_p''+2y_p'+y_p&=2\left[A_1+A_2+B_0+B_1+(2A_2+B_1+2B_2)x+
B_2x^2\right]\cos x\\
& +2\left[B_1+B_2-A_0-A_1+(2B_2-A_1-2A_2)x-A_2x^2\right]\sin x\\
&=8x^2\cos x-4x\sin x \qquad\text{if}\\
\text{(i) }& \begin{cases} 2B_2&=8\\-2A_2&=0\end{cases},\qquad
\text{(ii)} \begin{cases} 2B_1+4A_2+4B_2&=0\\-2A_1-4A_2+4B_2&=-4\end{cases},\qquad\text{and}
\\
\text{(iii)}& \begin{cases}2B_0+2A_1+2B_1+2A_2&=0
\\-2A_0-2A_1+2B_1+2B_2&=0\end{cases}.
\end{align*}
From (i), $A_2=0$, $B_2=4$. Substituting these into (ii) and solving
for $A_1$ and $B_1$ yields $A_1=10$, $B_1=-8$. Substituting the known
coefficients into (iii) and solving for $A_0$ and $B_0$ yields
$A_0=-14$, $B_0=-2$. Therefore,$y_p=-(14-10x)\cos x-(2+8x-4x^2)\sin
x$.


\item
Let
\begin{align*}
y_p&=(A_0+A_1x+A_2x^2)\cos2x +(B_0+B_1x+B_2x^2)\sin2x;\text{ then}\\
y_p'&=\left[A_1+2B_0+(2A_2+2B_1)x+2B_2x^2\right]\cos2x\\
&+\left[B_1-2A_0+(2B_2-2A_1)x-2A_2x^2\right]\sin2x\\
y_p''&=\left[-4A_0+2A_2+4B_1 -(4A_1-8B_2)x-4A_2x^2\right]\cos2x\\
&+ \left[-4B_0-2B_2-4A_1-(4B_1+8A_2)x-4B_2x^2\right]\sin2x,\text{ so}\\
y_p''+3y_p'+2y_p&=\left[-2A_0+3A_1+4A_2+6B_0+4B_1\right.\\
&\left.-(2A_1-6A_2-6B_1-8B_2)x-(2A_2-6B_2)x^2\right]\cos2x\\
&+\left[-2B_0+3B_1+4B_2-6A_0-4A_1\right.\\
&\left.-(2B_1-6B_2+6A_1+8A_2)x-(2B_2+6A_2)x^2\right]\sin2x\\
&=(1-x-4x^2)\cos2x-(1+7x+2x^2)\sin2x \qquad\text{if}\\
\text{(i)}&\begin{cases}-2A_2+6B_2&=-4
\\-6A_2-2B_2&=-2\end{cases},\qquad
\text{ (ii) } \begin{cases} -2A_1+6B_1+6A_2+8B_2&=-1
\\-6A_1-2B_1-8A_2+6B_2&=-7\end{cases},\qquad\text{and}
\\
\text{(iii)}&
\begin{cases}-2A_0+6B_0+3A_1+4B_1+2A_2&=1
\\-6A_0-2B_0-4A_0+3B_1+2B_2&=-1\end{cases}.
\end{align*}
From (i), $A_2=\dst\frac{1}{2}$, $B_2=-\dst\frac{1}{2}$. Substituting
these into (ii) and solving for $A_1$ and $B_1$ yields $A_1=0$,
$B_1=0$. Substituting the known coefficients into (iii) and solving
for $A_0$ and $B_0$ yields $A_0=0$, $B_0=0$. Therefore,
$y_p=\dst\frac{x^2}{2}(\cos2x-\sin2x)$.

\item
Let $y=ue^x$. Then
\begin{align*}
y''-2y'+y&=e^x\left[(u''+2u'+u)-2(u'+u)+u\right]=e^xu''\\
&=-e^x\left[(3+4x-x^2)\cos x+(3-4x-x^2)\sin x\right]
\end{align*}
 if $u''=-(3+4x-x^2)\cos x-(3-4x-x^2)\sin x$.
Now let
\begin{align*}
u_p&=(A_0+A_1x+A_2x^2)\cos x +(B_0+B_1x+B_2x^2)\sin x;\text{ then}\\
u_p'&=\left[A_1+B_0+(2A_2+B_1)x+B_2x^2\right]\cos x\\
&+\left[B_1-A_0+(2B_2-A_1)x-A_2x^2\right]\sin x, \\
u_p ''&=\left[-A_0+2A_2+2B_1-(A_1-4B_2)x-A_2x^2\right]\cos x\\
&+\left[-B_0+2B_2-2A_1-(B_1+4A_2)x-B_2x^2\right]\sin x\\
&=-(3+4x-x^2)\cos x-(3-4x-x^2)\sin x\qquad\text{if}\\
\text{(i)}& \begin{cases}-A_2&=1\\ -B_2&=1 \end{cases},\qquad
\text{(ii)} \begin{cases} -A_1+4B_2&=-4
\\ -B_1-4A_2&=4\end{cases},\qquad\text{and}
\\
\text{(iii)}&
\begin{cases}-A_0+2B_1+2A_2&=-3
\\ -B_0-2A_1+2B_2&=-3\end{cases}.
\end{align*}
From (i), $A_2=-1$, $B_2=-1$. Substituting these into (ii) and solving
for $A_1$ and $B_1$ yields $A_1=0$, $B_1=0$. Substituting the known
coefficients into (iii) and solving for $A_0$ and $B_0$ yields
$A_0=1$, $B_0=1$. Therefore,$u_p=(1-x^2)(\cos x+\sin x)$ and
$y_p=e^x(1-x^2)(\cos x+\sin x)$.


\item
Let $y=ue^{-x}$. Then
\begin{align*}
y''+2y'+y&=e^{-x}\left[(u''-2u'+u)+2(u'-u)+u\right]\\
&=e^{-x}u''=e^{-x}\left[(5-2x)\cos x-(3+3x)\sin x\right]
\end{align*}
 if $u''=(5-2x)\cos x-(3+3x)\sin x$. Let
\begin{align*}
u_p&=(A_0+A_1x)\cos x+(B_0+B_1x)\sin x; \text{ then}\\
u_p'&=(A_1+B_0+B_1x)\cos x+(B_1-A_0-A_1x)\sin x\\
u_p''&=(2B_1-A_0-A_1x)\cos x-(2A_1+B_0+B_1x)\sin x\\
&=(5-2x)\cos x-(3+3x)\sin x
\end{align*}
if $-A_1=-2$, $-B_1=-3$, $-A_0+2B_1=5$, $-B_0-2A_1$=-3. Therefore,
$A_1=2$, $B_1=3$, $A_0=1$, $B_0=-1$, $u_p=e^{-x}\left[(1+2x)\cos
x-(1-3x)\sin x\right]$, and $y_p=e^{-x}\left[(1+2x)\cos x-(1-3x)\sin
x\right]$.


\item
Let
\begin{align*}
y_p&=(A_0x+A_1x^2+A_2x^3)\cos x +(B_0x+B_1x^2+B_2x^3)\sin x;\text{ then}
\\ y_p'&=\left[A_0+(2A_1+B_0)x+(3A_2+B_1)x^2+B_2x^3\right]\cos x\\
&+\left[B_0+(2B_1-A_0)x+(3B_2-A_1)x^2-A_2x^3\right]\sin x\\
y_p''&=\left[2A_1+2B_0-(A_0-6A_2-4B_1)x-(A_1-6B_2)x^2-A_2x^3\right]\cos x\\
&+\left[2B_1-2A_0-(B_0+6B_2+4A_1)x-(B_1+6A_2)x^2-B_2x^3\right]\sin x, \text{ so}\\
y_p''+y_p&=\left[2A_1+2B_0+(6A_2+4B_1)x+6B_2x^2\right]\cos x\\
&+\left[2B_1-2A_0+(6B_2-4A_1)x-6A_2x^2\right]\sin x\\
&=(2+2x)\cos x+(4+6x^2)\sin x \qquad\text{if}\\
\text{(i)}&\begin{cases} 6B_2&=0\\-6A_2&=6\end{cases},\quad
\text{(ii)}\quad\begin{cases}\phantom{-}4B_1+6A_2&=2
\\-4A_1+6B_2&=0\end{cases},\quad\text{and}\quad
\text{(iii)}\begin{cases}\phantom{-}2B_0+2A_1&=2
\\-2A_0+2B_1&=4\end{cases}.
\end{align*}
From (i), $A_2=-1$, $B_2=0$. Substituting these into (ii) and solving
for $A_1$ and $B_1$ yields $A_1=0$, $B_1=2$. Substituting the known
coefficients into (iii) and solving for $A_0$ and $B_0$ yields
$A_0=0$, $B_0=2$. Therefore,$y_p=-x^3\cos x+(x+2x^2)\sin x$.




\item
Let $y=ue^x$. Then
\begin{align*}
y''-7y'+6y&=e^x\left[(u''+2u'+u)-7(u'+u)+6u\right]\\
&=e^x(u''-5u')=-e^x(17\cos x-7\sin x)
\end{align*}
 if $u''-5u'=-17\cos x+7\sin x$.
Now let  $u_p=A\cos x+B\sin x$. Then
\begin{align*}
u_p''-5u_p'&=-(A\cos x+B\sin x)-5(-A\sin x+B\cos x)\\
&= (-A-5B)\cos x-(B-5A)\sin x=-17\cos x+7\sin x
\end{align*}
if $-A-5B=-17$, $5A-B=7$. Therefore,$A=2$, $B=3$, $u_p=2\cos x+3\sin
x$, and $y_p=e^x(2\cos x+3\sin x)$. The characteristic polynomial of
the complementary equation is $p(r)=r^2-7r+6=(r-1)(r-6)$, so
$\{e^x,e^{6x}\}$ is a fundamental set of solutions of the
complementary equation. Therefore,
(A) $y=e^x(2\cos x+3\sin x)+c_1e^x+c_2e^{6x}$
is the general solution of the nonhomogeneous
equation. Differentiating (A) yields
$y'=e^x(2\cos x+3\sin x)+e^x(-2\sin x+3\cos x)+c_1e^x+6c_2e^{6x}$,
so $y(0)=4,\ y'(0)=2\Rightarrow 4=2+c_1+c_2,\
2=2+3+c_1+6c_2\Rightarrow c_1+c_2=2,\ c_1+6c_2=-3$, so $c_1=3$,
$c_2=-1$, and
 $y=e^x(2\cos x+3\sin x)+3e^x-e^{6x}$.


\item
Let $y=ue^x$. Then
\begin{align*}
y''+6y'+10y&=e^x\left[(u''+2u'+u)+6(u'+u)+10u\right]\\
&=e^x(u''+8u'+17u)=
-40e^x\sin x
\end{align*}
 if  $u''+8u'+17u=-40\sin x$.
Let $u_p=A\cos x+B\sin x$. Then
\begin{align*}
u_p''+6u_p'+17u_p&=-(A\cos x+B\sin x)+8(-A\sin x+B\cos x)\\
&+17(A\cos x+B\sin x)\\
&=(16A+8B)\cos x-(8A-16B)\sin x=-40\sin x
\end{align*}
if $16A+8B=0$, $-8A+16B=-40$. Therefore,$A=1$, $B=-2$, and
$y_p=e^x(\cos x-2\sin x)$. The characteristic polynomial of the
complementary equation is $p(r)=r^2+6r+10=(r+3)^2+1$, so
$\{e^{-3x}\cos x,e^{-3x}\sin x\}$ is a fundamental set of solutions
of the complementary equation, and (A) $y=e^x(\cos x-2\sin
x)+e^{-3x}(c_1\cos x+c_2\sin x)$ is the general solution of the
nonhomogeneous equation. Therefore,$y(0)=2\Rightarrow 2=1+c_1$, so
$c_1=1$. Differentiating (A) yields $y'=e^x(\cos x-2\sin x)-e^x(\sin
x+2\cos x) -3e^{-3x}(c_1\cos x+c_2\sin x) +e^{-3x}(-c_1\sin x+c_2\cos
x)$. Therefore,$y'(0)=-3\Rightarrow -3=1-2-3c_1+c_2$, so $c_2=1$, and
$y=e^x(\cos x-2\sin x)+e^{-3x}(\cos x+\sin x)$.


\item
Let $y=ue^{3x}$. Then
\begin{align*}
y''-3y'+2y&=e^{3x}\left[(u''+6u'+9u)-3(u'+3u)+2u\right]\\
&=e^{3x}(u''+3u'+2u)=e^{3x}\left[21\cos x-(11+10x)\sin x\right]
\end{align*}
if $u''+3u'+2u=21\cos x-(11+10x)\sin x$.
Now let
\begin{align*}
u_p&=(A_0+A_1x)\cos x+(B_0+B_1x)\sin x; \text{ then}\\
u_p'&=(A_1+B_0+B_1x)\cos x+(B_1-A_0-A_1x)\sin x\\
u_p''&=(2B_1-A_0-A_1x)\cos x-(2A_1+B_0+B_1x)\sin x, \text{ so}\\
u''+3u'+2u&=\left[A_0+3A_1+3B_0+2B_1+(A_1+3B_1)x\right]\cos x\\
&+\left[B_0+3 B_1-3A_0-2A_1+(B_1-3A_1)x\right]\sin x \\
&=21\cos x-(11+10x)\sin x \qquad\text{if}\\
\begin{cases}
\phantom{-3}A_1+3B_1&=0\\ -3A_1+\phantom{3}B_1&=-10
\end{cases}
&\qquad\text{and}\qquad
\begin{cases}
\phantom{-3} A_0+3B_0+3A_1+2B_1&=21\\
-3A_0+\phantom{3}B_0-2A_1+3B_1&=-11
\end{cases}.
\end{align*}
From the first two equations $A_1=3$, $B_1=-1$. Substituting these in
last two equations %yields
and solving for $A_0$ and $B_0$ yields
$A_0=2$, $B_0=4$. Therefore,
$u_p=(2+3x)\cos x+(4-x)\sin x$ and
$y_p=e^{3x}\left[(2+3x)\cos x+(4-x)\sin x\right]$.
The characteristic polynomial of the complementary equation is
$p(r)=r^2-3r+2=(r-1)(r-2)$, so $\{e^x,e^{2x}\}$ is a fundamental
set of solutions of the complementary equation, and
(A) $y=e^{3x}\left[(2+3x)\cos x+(4-x)\sin x\right]+c_1e^x+c_2e^{2x}$
is the general solution of the nonhomogeneous
equation. Differentiating (A) yields
\begin{align*}
y'&=3e^{3x}\left[(2+3x)\cos x+(4-x)\sin x\right]\\
&+e^{3x}\left[(7-x)\cos x-(3+3x)\sin x\right]+c_1e^x+2c_2e^{2x}.
\end{align*}
 Therefore,$y(0)=0,\ y'(0)=6\Rightarrow 0=2+c_1+c_2,\
6=6+7+c_1+2c_2$, so $c_1+c_2=-2,\ c_1+2c_2=-7$.  Therefore,
$c_1=3$, $c_2=-5$, and
 $y=e^{3x}\left[(2+3x)\cos x+(4-x)\sin x\right]+3e^x-5e^{2x}$.


\item
We must find particular solutions $y_{p_1}$, $y_{p_2}$, and $y_{p_3}$
of
(A) $y''+y=4\cos x-2\sin x$ and
(B) $y''+y=xe^x$, and
(C) $y''+y=e^{-x}$, respectively.
To find a particular solution of (A) we write
\begin{align*}
y_{p_1}&=Ax\cos x+Bx\sin x; \text{ then}\\
y_{p_1}'&=(A+Bx)\cos x+(B-Ax)\sin x\\
y_{p_1}''&=(2B-Ax)\cos x-(2A+Bx)\sin x, \text{ so}
\end{align*}
$y_{p_1}''+y_{p_1}=2B\cos x-2A\sin x=4\cos x-2\sin x$ if $2B=4$,
$-2A=-2$. Therefore, $A=1$, $B=2$, and $y_{p_1}=x(\cos x+2\sin x)$.
To find a particular solution of (B) we write $y=ue^x$. Then
\begin{align*}
y''+y&=e^x\left[(u''+2u'+u)+u\right]\\
&=e^x(u''+2u'+2u)=xe^x
\end{align*}
if $u''+2u'+2u=x$. Now $u_p=A+Bx$, where $2B+2(A+Bx)=x$. Therefore,
$2B=1$, $2A+2B=0$, so $B=\dst\frac{1}{2}$, $A=-\dst\frac{1}{2}$,
$u_p=-\dst\frac{1}{2}(1-x)$, and $y_{p_2}=-\dst\frac{e^x}{2}(1-x)$. To
find a particular solution of (C) we write $y_{p_3}=Ae^{-x}$. Then
$y_{p_3}''+y_{p_3}=2Ae^{-x}=e^{-x}$ if $2A=1$, so $A=\dst\frac{1}{2}$
and $y_{p_3}=\dfrac{e^{-x}}{2}$ Now $y_p=y_{p_1}+y_{p_2}+y_{p_3}=
x(\cos x+2\sin x)-\dst\frac{e^x}{2}(1-x)+\dfrac{e^{-x}}{2}$.


\item
We must find particular solutions $y_{p_1}$, $y_{p_2}$ and $y_{p_3}$
of
(A) $y''-2y'+2y=4xe^x\cos x$,
(B) $y''-2y'+2y=xe^{-x}$, and
(C) $y''-2y'+2y=1+x^2$, respectively.
To find a particular solution of (A) we write $y=ue^x$. Then
$y''-2y'+2y=e^x\left[(u''+2u'+u)-2(u'+u)+2u\right]=e^x(u''+u)=
4xe^x\cos x$ if $u''+u=4x\cos x$. Now let
\begin{align*}
u_p&=(A_0x+A_1x^2)\cos x+(B_0x+B_1x^2)\sin x;\text{ then}\\
u_p'&=\left[A_0+(2A_1+B_0)x+B_1x^2\right]\cos x
+\left[B_0+(2B_1-A_0)x-B_1x^2\right]\sin x\\
u_p''&=\left[2A_1+2B_0-(A_0-4B_1)x-A_1x^2\right]\cos x\\
&+\left[2B_1-2A_0-(B_0+4A_1)x-B_1x^2\right]\sin x, \text{ so}\\
u_p''+u_p&=(2A_1+2B_0+4B_1x)\cos x+(2B_1-2A_0-4A_1x)\sin x\\
&=4x\cos x
\end{align*}
if $4B_1=4$, $-4A_1=0$, $2B_0+2A_1=0$, $-2A_0+2B_1=0$. Therefore,
$A_1=0$, $B_1=1$, $A_0=1$, $B_0=0$, $u_p=x(\cos x+x\sin x)$, and
$y_{p_1}=xe^x(\cos x+x\sin x)$. To find a particular solution of (B)
we write $y=ue^{-x}$. Then
\begin{align*}
y''-2y'+2y&=e^{-x}\left[(u''-2u'+u)-2(u'-u)+2u\right]\\
&=e^{-x}(u''-4u'+5u)=xe^{-x}
\end{align*}
if $u''-4u'+5u=x$. Now $u_p=A+Bx$ where $-4B+5(A+Bx)=x$. Therefore,
$5B=1$, $5A-4B=0$, $B=\dst\frac{1}{5}$, $A=\dst\frac{4}{25}$,
$u_p=\dst\frac{1}{25}(4+5x)$, and $y_{p_2}=\dfrac{e^{-x}}{25}(4+5x)$.
To find a particular solution of (C) we write $y_{p_3}=A+Bx+Cx^2$.
Then
\begin{align*}
y_{p_3}''-2y_{p_3}'+2y_{p_3}&=2C-2(B+2Cx)+2(A+Bx+Cx^2)\\
&=(2A-2B+2C)+(2B-4C)x+2Cx^2=1+x^2
\end{align*}
if $2A-2B+2C=1$, $2B-4C=0$, $2C=1$. Therefore,$C=\dst\frac{1}{2}$,
$B=1$, $A=1$, and $y_{p_3}=1+x+\dst\frac{x^2}{2}$. Now
$y_p=y_{p_1}+y_{p_2}+y_{p_3}= xe^x(\cos x+x\sin
x)+\dfrac{e^{-x}}{25}(4+5x)+1+x+\dst\frac{x^2}{2}$.


\item
We must find particular solutions $y_{p_1}$ and $y_{p_2}$ of
(A) $y''-4y'+4y=6e^{2x}$ and
(B) $y''-4y'+4y=25\sin x$, respectively.
To find a particular solution of (A), let
 $y=ue^{2x}$. Then
\begin{align*}
y''-4y'+4y&=e^{2x}\left[(u''+4u'+4u)-4(u'+2u)+4u\right]\\
&=e^{2x}u''=6e^{2x}
\end{align*}
if $u''=6$. Integrating twice and taking the constants of integration
to be zero yields $u_p=3x^2$, so $y_{p_1}=3x^2e^{2x}$.
To find a particular solution of (B), let
 $y_{p_2}=A\cos x+B\sin x$. Then
\begin{align*}
y_{p_2}''-4y_{p_2}'+4y_{p_2}&=-(A\cos x+B\sin x)-4(-A\sin x+B\cos x)\\
&+4(A\cos x+B\sin x)\\
&=(3A-4B)\cos x+(4A+3B)\sin x=25\sin x
\end{align*}
if $3A-4B=0$, $4A+3B=25$. Therefore,$A=4$, $B=3$, and $y_{p_2}=4\cos
x+3\sin x$. Now $y_p=y_{p_1}+y_{p_2}=3x^2e^{2x}+4\cos x+3\sin x$. The
characteristic polynomial of the complementary equation is
$p(r)=r^2-4r+4=(r-2)^2$, so $\{e^{2x},xe^{2x}\}$ is a fundamental set
of solutions of the complementary equation. Therefore,(C)
$y=3x^2e^{2x}+4\cos x+3\sin x+e^{2x}(c_1+c_2x)$ is the general
solution of the nonhomogeneous equation. Now $y(0)=5\Rightarrow
5=4+c_1$, so $c_1=1$. Differentiating (C) yields
$y'=6e^{2x}(x+x^2)-4\sin x+3\cos x+2e^{2x}(c_1+c_2x)+c_2e^{2x}$, so
$y'(0)=3\Rightarrow 3=3+2+c_2$. Therefore,$c_2=-2$, and
$y=(1-2x+3x^2)e^{2x}+4\cos x+3\sin x$.



\item
We must find particular solutions $y_{p_1}$ and $y_{p_2}$ of
(A) $y''+4y'+4y=2\cos2x+3\sin2x$ and
(B) $y''+4y'+4y=e^{-x}$,  respectively.
To find a particular solution of (A) we write
 $y_{p_1}=A\cos2x+B\sin2x$. Then
\begin{align*}
y_{p_1}''+4y_{p_1}'+4y_{p_1}
&=-4(A\cos2x+B\sin2x)+8(-A\sin2x+B\cos2x)
%\\
%&\qquad
+4(A\cos2x+B\sin2x)\\
&=-8A\sin2x+8B\cos2x=2\cos2x+3\sin2x
\end{align*}
if $8B=2$, $-8A=3$. Therefore, $A=-\dst\frac{3}{8}$, $B=\dst\frac{1}{4}$,
and $y_{p_1}=-\dst\frac{3}{8}\cos2x+\dst\frac{1}{4}\sin2x$. To find a
particular solution of (B) we write $y_{p_2}=Ae^{-x}$. Then
$y_{p_2}''+4y_{p_2}'+4y_{p_2}=A(1-4+4)e^{-x}=Ae^{-x}=e^{-x}$ if $A=1$.
Therefore,$y_{p_2}=e^{-x}$. Now $y_p=y_{p_1}+y_{p_2}=
-\dst\frac{3}{8}\cos2x+\dst\frac{1}{4}\sin2x+e^{-x}$. The characteristic
polynomial of the complementary equation is $p(r)=r^2+4r+4=(r-2)^2$,
so $\{e^{2x},xe^{2x}\}$ is a fundamental set of solutions of the
complementary equation. Therefore,(C)
$y=-\dst\frac{3}{8}\cos2x+\dst\frac{1}{4}\sin2x+e^{-x}+e^{-2x}(c_1+c_2x)$
is the general solution of the nonhomogeneous equation. Now
$y(0)=-1\Rightarrow -1=-\dst\frac{3}{8}+1+c_1$, so
$c_1=-\dst\frac{13}{8}$. Differentiating (C) yields
$y'=\dst\frac{3}{4}\sin2x+\dst\frac{1}{2}\cos2x-e^{-x}-2e^{-2x}(c_1+c_2x)
+c_2e^{-2x}$, so $y'(0)=2\Rightarrow 2=\dst\frac{1}{2}-1-2c_1+c_2$.
Therefore,$c_2=-\dst\frac{3}{4}$, and
$y=\dst{-\frac{3}{8}\cos2x+\frac{1}{4}\sin2x+e^{-x}-\frac{13}{8}e^{-2x}-
\frac{3}{4}xe^{-2x}}$.


\item
\part{a}, \part{b}, and \part{c} require only routine
manipulations.
\part{d}
The coefficients of $\sin\omega x$ in
$y_p'$, $y_p''$, $ay_p''+by_p'+cy_p$,
and $y_p''+\omega^2 y_p$  can be obtained by replacing $A$ by $B$
and $B$ by $-A$ in the corresponding coefficients of $\cos\omega x$.



\item
Let $y=ue^{\lambda x}$. Then
\begin{align*}
ay''+by'+cy&=e^{\lambda x}\left[a(u''+2\lambda u'+\lambda^2u)
+b(u'+\lambda u)+cu\right]\\
&=e^{\lambda x}\left[au''+(2a\lambda+b)u'+(a\lambda^2+b\lambda+c)u\right]\\
&=e^{\lambda x}\left[au''+p'(\lambda)u'+p(\lambda)u\right]\\
&=e^{\lambda x}\left(P(x)\cos \omega x+Q(x)\sin \omega x\right)\text{ if}
\end{align*}
(A) $au''+p'(\lambda)u'+p(\lambda)u
P(x)\cos \omega x+Q(x)\sin \omega x$,
where $p(r)=ar^+br+c$ is that characteristic polynomial
of the  complementary equation (B) $ay''+by'+cy=0$.
If $e^{\lambda x}\cos\omega x$ and $e^{\lambda
x}\sin\omega x$ are not solutions of (B), then
$\cos\omega x$ and $\sin\omega x$ are not solutions of the
complementary equation for (A).  Then  Theorem~5.5.1
implies that  (A) has a particular solution
\[
u_p=(A_0+A_1x+\cdots+A_kx^k)\cos\omega x+
(B_0+B_1x+\cdots+B_kx^k)\sin\omega x,
\]
 and $y_p=u_pe^{\lambda x}$
is a particular solution of the stated form for the given equation.
If $e^{\lambda x}\cos\omega x$ and $e^{\lambda
x}\sin\omega x$ are  solutions of (B), then
$\cos\omega x$ and $\sin\omega x$ are solutions of the
complementary equation for (A).  Then  Theorem~5.5.1
implies that  (A) has a particular solution
\[
u_p=(A_0x+A_1x^2+\cdots+A_kx^{k+1})\cos\omega x+
(B_0x+B_1x^2+\cdots+B_kx^{k+1})\sin\omega x,
\]
and $y_p=u_pe^{\lambda x}$ is a particular solution of the stated form
for the given equation.

\item
\begin{alist}
\item Let $y=\int x^2\cos x\,dx$; then
$y'=x^2\cos x$ Now let
\begin{align*}
y_p&=(A_0+A_1x+A_2x^2)\cos x +(B_0+B_1x+B_2x^2)\sin x;\text{ then}\\
y_p'&=\left[A_1+B_0+(2A_2+B_1)x+B_2x^2\right]\cos x\\
&\qquad+\left[B_1-A_0+(2B_2-A_1)x-A_2x^2\right]\sin x
=x^2\cos x \qquad\text{if}\\
\text{(i)}
\begin{cases} \phantom{-}B_2&=1\\ -A_2&=0 \end{cases},&
\quad
\text{(ii)}
\begin{cases}
\phantom{-}B_1+2A_2&=0\\ -A_1+2B_2&=0\end{cases},\quad\text{and}\quad
\text{(iii)}
\begin{cases} \phantom{-}B_0+A_1&=0\\ -A_0+B_1&=0
\end{cases}.
\end{align*}
Solving these equations yields $A_2=0$, $B_2=1$, $A_1=2$, $B_1=0$,
$A_0=0$, $B_0=-2$. Therefore,$y_p=2x\cos x-(2-x^2)\sin x$ and
$y=2x\cos x-(2-x^2)\sin x+c$.

\item Let $y=\int x^2e^x\cos x\,dx=ue^x$; then
$y'=(u'+u)e^x=x^2e^x\cos x$ if $u'+u=x^2\cos x$. Now let
\begin{align*}
u_p&=(A_0+A_1x+A_2x^2)\cos x +(B_0+B_1x+B_2x^2)\sin x;\text{ then}\\
u_p'&=\left[A_1+B_0+(2A_2+B_1)x+B_2x^2\right]\cos x\\
&+\left[B_1-A_0+(2B_2-A_1)x-A_2x^2\right]\sin x, \text{ so}\\
u_p''+u_p&=\left[A_0+A_1+B_0+(A_1+2A_2+B_1)x+(A_2+B_2)x^2\right]\cos x\\
& +\left[B_0+B_1-A_0+(B_1+2B_2-A_1)x+(B_2-A_2)x^2\right]\sin x\\
&=x^2\cos x \qquad\text{if}\\
\text{(i)}&
\begin{cases} \phantom{-}A_2+B_2&=1\\ -A_2+B_2&=0 \end{cases},\qquad
\text{(ii)}
\begin{cases}
\phantom{-}A_1+B_1+2A_2&=0\\ -A_1+B_1+2B_2&=0\end{cases},\qquad\text{and}
\\
\text{(iii)}&
\begin{cases} \phantom{-}A_0+B_0+A_1&=0\\ -A_0+B_0+B_1&=0
\end{cases}.
\end{align*}
From (i), $A_2=\dst\frac{1}{2}$, $B_2=\dst\frac{1}{2}$. Substituting these
into (ii) and solving for $A_1$ and $B_1$ yields $A_1=0$, $B_1=-1$.
Substituting these into (iii) and solving for $A_0$ and $B_0$ yields
$A_0=-\dst\frac{1}{2}$, $B_0=\dst\frac{1}{2}$. Therefore,
$u_p=-\dst\frac{1}{2}\left[(1-x^2)\cos x-(1-x)^2\sin x\right]$ and
$y=-\dst\frac{e^x}{2}\left[(1-x^2)\cos x-(1-x)^2\sin x\right]$.

\item Let $y=\int xe^{-x}\sin2x\,dx=ue^{-x}$; then
$y'=(u'-u)e^{-x}=xe^{-x}\sin2x$ if $u'-u=x\sin2x$. Now let
\begin{align*}
u_p&=(A_0+A_1x)\cos2x +(B_0+B_1x)\sin2x;\text{ then}\\
u_p'&=\left[(A_1+2B_0)+2B_1x\right]\cos2x
+\left[(B_1-2A_0)-2A_1x\right]\sin2x,\text{ so}\\
u_p''-u_p&=\left[-A_0+A_1+2B_0-(A_1-2B_1)x\right]\cos2x\\
&+\left[-B_0+B_1-2A_0-(B_1+2A_1)x\right]\sin2x=x\sin2x\qquad\text{if}\\
\text{(i)}& \begin{cases} -A_1+2B_1&=0\\
-2A_1-\phantom{2}B_1&=1 \end{cases}
\quad\text{and}\quad \text{(ii)} \begin{cases}
-\phantom{2}A_0+2B_0+A_1&=0\\ -2A_0-\phantom{2}B_0+B_1&=0
\end{cases}.
\end{align*}
 From (i), $A_1=-\dst\frac{2}{5}$, $B_1=-\dst\frac{1}{5}$.
Substituting these into (ii) and solving for $A_0$ and $B_0$ yields
$A_0=-\dst\frac{4}{25}$, $B_0=\dst\frac{3}{25}$. Therefore,
\begin{gather*}
u_p=-\dst\frac{1}{25}\left[(4+10x)\cos 2x-(3-5x)\sin2x\right]+c \text{
and} \\
  y_p=-\dfrac{e^{-x}}{25}\left[(4+10x)\cos
2x-(3-5x)\sin2x\right]+c. 
\end{gather*}

\item Let $y=\int x^2e^{-x}\sin x\,dx=ue^{-x}$; then
$y'=(u'-u)e^{-x}=x^2e^{-x}\sin x$ if $u'-u=x^2\sin x$.
Now let
\begin{align*}
u_p&=(A_0+A_1x+A_2x^2)\cos x +(B_0+B_1x+B_2x^2)\sin x;\text{ then}\\
u_p'&=\left[A_1+B_0+(2A_2+B_1)x+B_2x^2\right]\cos x\\
&+\left[B_1-A_0+(2B_2-A_1)x-A_2x^2\right]\sin x, \text{ so}\\
u_p''-u_p&=\left[-A_0+A_1+B_0-(A_1-2A_2-B_1)x-(A_2-B_2)x^2\right]\cos x\\
& +\left[-B_0+B_1-A_0-(B_1-2B_2+A_1)x-(B_2+A_2)x^2\right]\sin x\\
&=x^2\sin x \qquad\text{if}\\
\text{(i)}&
\begin{cases} -A_2+B_2&=0\\ -A_2-B_2&=1\end{cases},\qquad
\text{(ii)}
\begin{cases}
-A_1+B_1+2A_2&=0\\ -A_1-B_1+2B_2&=0\end{cases},\qquad\text{and}
\\
\text{(iii)}&
\begin{cases} -A_0+B_0+A_1&=0\\ -A_0-B_0+B_1&=0
\end{cases}.
\end{align*}
From (i), $A_2=-\dst\frac{1}{2}$, $B_2=-\dst\frac{1}{2}$. Substituting
these into (ii) and solving for $A_1$ and $B_1$ yields $A_1=-1$,
$B_1=0$. Substituting these into (iii) and solving for $A_0$ and $B_0$
yields $A_0=-\dst\frac{1}{2}$, $B_0=\dst\frac{1}{2}$. Therefore,
\begin{gather*}
 u_p=-\frac{e^{-x}}{2}\left[(1+x)^2\cos x-(1-x^2)\sin
x \right] \text{ and}
\\
y=-\frac{e^{-x}}{2}\left[(1+x)^2\cos x-(1-x^2)\sin
x \right]+c.
\end{gather*}

\item Let $y=\int x^3e^x\sin x\,dx=ue^x$; then
$y'=(u'+u)e^x=x^3e^x\sin x$ if $u'+u=x^3\sin x$. Now let
\begin{align*}
u_p&=(A_0+A_1x+A_2x^2+A_3x^3)\cos x +(B_0+B_1x+B_2x^2+B_3x^3)\sin x;\text{ then}\\
u_p'&=\left[A_1+B_0+(2A_2+B_1)x+(3A_3+B_2)x^2+B_3x^3\right]\cos x\\
&+\left[B_1-A_0+(2B_2-A_1)x+(3B_3-A_2)x^2-A_3x^3\right]\sin x,\text{ so}\\
u_p''+u_p&=\left[A_0+A_1+B_0+(A_1+2A_2+B_1)x\right.\\
&\left.+(A_2+3A_3+B_2)x^2+(A_3+B_3)x^3\right]\cos x\\
& +\left[B_0+B_1-A_0+(B_1+2B_2-A_1)x\right.\\
&\left.+(B_2+3B_3-A_2)x^2+(B_3-A_3)x^3\right]\sin x=x^3\sin x \qquad\text{if}\\
\text{(i)}&
\begin{cases} \phantom{-}A_3+B_3&=0\\ -A_3+B_3&=1\end{cases},\qquad
\text{(ii)}
\begin{cases} \phantom{-}A_2+B_2+3A_3&=0\\ -A_2+B_2+3B_3&=0
\end{cases},
\\
\text{(iii)}&
\begin{cases}
\phantom{-}A_1+B_1+2A_2&=0\\ -A_1+B_1+2B_2&=0\end{cases},\quad\text{and}\quad
\text{(iv)}
\begin{cases} \phantom{-}A_0+B_0+A_1&=0\\ -A_0+B_0+B_1&=0
\end{cases}.
\end{align*}
From (i), $A_3=-\dst\frac{1}{2}$, $B_3=\dst\frac{1}{2}$. Substituting
these into (ii) and solving for $A_2$ and $B_2$ yields
$A_2=\dst\frac{3}{2}$, $B_2=0$. Substituting these into (iii) and
solving for $A_1$ and $B_1$ yields $A_1=-\dst\frac{3}{2}$,
$B_1=-\dst\frac{3}{2}$. Substituting these into (iv) and solving for
$A_0$ and $B_0$ yields $A_0=0$, $B_0=\dst\frac{3}{2}$. Therefore,
\begin{gather*}
u_p=-\dst\frac{1}{2}\left[x(3-3x+x^2)\cos
x-(3-3x+x^3)\sin x\right] \text{ and}
\\
y=-\dst\frac{e^x}{2}\left[x(3-3x+x^2)\cos x-(3-3x+x^3)\sin x\right]+c.
\end{gather*}

\item Let $y=\int e^x\left[x\cos x-(1+3x)\sin x\right]\,dx=ue^x$;
then $y'=(u'+u)e^x=e^x\left[x\cos x-(1+3x)\sin x\right]$ if
$u'+u=x\cos x-(1+3x)\sin x$. Now let
\begin{align*}
u_p&=(A_0+A_1x)\cos x +(B_0+B_1x)\sin x;\text{ then}\\
u_p'&=\left[A_1+B_0+B_1x\right]\cos x+\left[B_1-A_0-A_1x\right]\sin x,\text{ so}\\
u_p''+u_p&=\left[A_0+A_1+B_0+(A_1+B_1)x\right]\cos x\\
&+\left[B_0+B_1-A_0+(B_1-A_1)x\right]\sin x\\
&=x\cos x-(1+3x)\sin x\qquad\text{if}\\
\text{(i)}&
\begin{cases}\phantom{-}A_1+B_1&=1\\ -A_1+B_1&=-3
\end{cases},\quad\text{and}\quad
\text{(ii)}
\begin{cases} \phantom{-}A_0+B_0+A_1&=0\\ -A_0+B_0+B_1&=-1
\end{cases}.
\end{align*}
From (i), $A_1=2$, $B_1=-1$. Substituting
these into (ii) and solving for $A_0$ and $B_0$ yields $A_0=-1$,
$B_0=-1$. Therefore,
 $u_p=-\left[(1-2x)\cos x+(1+x)\sin x\right]$ and
$y=-e^x\left[(1-2x)\cos x+(1+x)\sin x\right]+c$.

\item Let $y=\int e^{-x}\left[(1+x^2)\cos x+(1-x^2)\sin
x\right]\,dx =ue^{-x}$; then 
\[
y'=(u'-u)e^{-x}=e^{-x}\left[(1+x^2)\cos x+(1-x^2)\sin x\right] 
\]
 if
$u'-u=(1+x^2)\cos x+(1-x^2)\sin x$. Now let
\begin{align*}
u_p&=(A_0+A_1x+A_2x^2)\cos x +(B_0+B_1x+B_2x^2)\sin x;\text{ then}\\
u_p'&=\left[A_1+B_0+(2A_2+B_1)x+B_2x^2\right]\cos x\\
&+\left[B_1-A_0+(2B_2-A_1)x-A_2x^2\right]\sin x, \text{ so}\\
u_p''-u_p&=\left[-A_0+A_1+B_0-(A_1-2A_2-B_1)x-(A_2-B_2)x^2\right]\cos x\\
& +\left[-B_0+B_1-A_0-(B_1-2B_2+A_1)x-(B_2+A_2)x^2\right]\sin x\\
&= (1+x^2)\cos x+(1-x^2)\sin x \qquad\text{if}\\
\text{(i)}& \begin{cases} -A_2+B_2&=1\\ -A_2-B_2&=-1
\end{cases},\qquad
\text{(ii)} \begin{cases} -A_1+B_1+2A_2&=0\\
-A_1-B_1+2B_2&=0\end{cases},\qquad\text{and} \\
  \text{(iii)}& \begin{cases}
-A_0+B_0+A_1&=1\\ -A_0-B_0+B_1&=1 \end{cases}.
\end{align*}
 From (i), $A_2=0$,
$B_2=1$. Substituting these into (ii) and solving for $A_1$ and $B_1$
yields $A_1=1$, $B_1=1$. Substituting these into (iii) and solving for
$A_0$ and $B_0$ yields $A_0=0$, $B_0=0$. Therefore,$u_p=x\cos
x+x(1+x)\sin x$ and $y=e^{-x}\left[x\cos x+x(1+x)\sin x\right]+c$.
\end{alist}
\end{sectionSolutions}


\section{Reduction of Order}

\note{The term $uy_1''$ is indicated by ``$\cdots$'' in some of the
following solutions, where $y_1''$ is complicated. Since this term
always drops out of the differential equation for $u$, it is not
necessary to include it.}

\begin{sectionSolutions}

\item
If $y=ux$, then $y'=u'x+u$ and $y''=u''x+2u'$, so
$x^2y''+xy'-y=x^3u''+3x^2u'=\dst\frac{4}{ x^2}$ if $u'=z$, where (A)
$z'+\dst\frac{3}{ x}z=\dst\frac{4}{ x^5}$. Since $\dst{\int\frac{3}{
x}\,dx}=3\ln|x|$, $z_1=\dst\frac{1}{ x^3}$ is a solution of the
complementary equation for (A). Therefore,the solutions of (A) are of
the form (B) $z=\dst\frac{v}{ x^3}$, where $\dst\frac{v'}{
x^3}=\dst\frac{4}{ x^5}$, so $v'=\dst\frac{4}{ x^2}$. Hence,
$v=-\dst\frac{4}{ x}+C_1$;\ $u'=z=-\dst\frac{4}{ x^4}+\dst\frac{C_1}{ x^3}$
(see (B));\ $u=\dst\frac{4}{ 3x^3}-\dst\frac{C_1}{2x^2}+C_2$;\;
$y=ux=\dst\frac{4}{ 3x^2}-\dst\frac{C_1}{2x}+C_2x$, or $y=\dst\frac{4}{
3x^2}+c_1x+\dst\frac{c_2}{ x}$. As a byproduct, $\{x,1/x\}$ is a
fundamental set of solutions of the complementary equation.


\item
If $y=ue^{2x}$, then $y'=(u'+2u)e^{2x}$ and $y''=(u''+4u'+4u)e^{2x}$,
so $y''-3y'+2y=(u''+u')e^{2x}=\dst\frac{1}{1+e^{-x}}$ if $u'=z$, where
(A) $z'+z=\dfrac{e^{-2x}}{1+e^{-x}}$. Since $z_1=e^{-x}$ is a
solution of the complementary equation for (A), the solutions of (A)
are of the form (B) $z=ve^{-x}$, where
$v'e^{-x}=\dfrac{e^{-2x}}{1+e^{-x}}$, so
$v'=\dfrac{e^{-x}}{1+e^{-x}}$. Hence, $v=-\ln(1+e^{-x})+C_1$;\;
$u'=z=-e^{-x}\ln(1+e^{-x})+C_1e^{-x}$ (see (B));\;
$u=(1+e^{-x})\ln(1+e^{-x})-1-e^{-x}-C_1e^{-x}+C_2$;\;
$y=ue^{2x}=(e^{2x}+e^x)\ln(1+e^{-x})-(C_1+1)e^x+(C_2-1)e^{2x}$, or
$y=(e^{2x}+e^x) \ln (1+e^{-x})+c_1e^{2x}+c_2e^x$. As a byproduct,
$\{e^{2x},e^x\}$ is a fundamental set of solutions of the
complementary equation.


\item
If $y=ux^{1/2}e^x$, then
$y'=u'x^{1/2}e^x+u\dst{\left(x^{1/2}+\frac{x^{-1/2}}{2}\right)}e^x$ and
$y''=u''x^{1/2}e^x+2u'\dst{\left(x^{1/2}+\frac{x^{-1/2}}{2}\right)}e^x+\cdots$
so $4x^2y''+(4x-8x^2)y'+(4x^2-4x-1)y=e^x(4x^{5/2}u''+8x^{3/2}u')
=4x^{1/2}e^x(1+4x)$ if $u'=z$, where (A) $z'+\dst\frac{2}{
x}z=\dst\frac{1+4x}{ x^2}$. Since $\dst{\int\dst\frac{2}{ x}\,dx}=
2\ln|x|$, $z_1=\dst\frac{2}{ x^2}$ is a solution of the complementary
equation for (A). Therefore,the solutions of (A) are of the form (B)
$z=\dst\frac{v}{ x^2}$, where $\dst\frac{v'}{ x^2}=\dst\frac{1+4x}{ x^2}$,
so $v'=1+4x$. Hence, $v=x+2x^2+C_1$;\ $u'=z=\dst\frac{1}{
x}+2+\dst\frac{C_1}{ x^2}$ (see (B));\ $u=\ln x+2x-\dst\frac{C_1}{
x}+C_2$;\ $y=ux^{1/2}e^x=e^x(2x^{3/2}+x^{1/2}\ln
x-C_1x^{-1/2}+C_2x^{1/2})$, or $y=e^x(2x^{3/2}+x^{1/2}\ln
x+c_1x^{1/2}+c_2x^{-1/2})$. As a byproduct,
$\{x^{1/2}e^x,x^{-1/2}e^{-x}\}$ is a fundamental set of solutions of
the complementary equation.


\item
If $y=ue^{-x^2}$, then $y'=u'e^{-x^2}-2xue^{-x^2}$ and
$y''=u''e^{-x^2}-4x u'e^{-x^2}+\cdots$, so
$y''+4xy'+(4x^2+2)y=u''e^{-x^2}=8e^{-x(x+2)}=8e^{-x^2}e^{-2x}$ if
$u''=8e^{-2x}$. Therefore,$u'=-4e^{-2x}+C_1$;\ $u=2e^{-2x}+C_1x+C_2$,
and $y=ue^{-x^2}=e^{-x^2}(2e^{-2x}+C_1x+C_2)$, or
$y=e^{-x^2}(2e^{-2x}+c_1+c_2x)$. As a byproduct,
$\{e^{-x^2},xe^{-x^2}\}$ is a fundamental set of solutions of the
complementary equation.


\item
If $y=uxe^{-x}$, then $y'=u'x e^{-x}-ue^{-x}(x-1)$ and
$y''=u''xe^{-x}-2u'e^{-x}(x-1)+\cdots$, so
$x^2y''+2x(x-1)y'+(x^2-2x+2)y=x^3u''=x^3e^{2x}$ if $u''=e^{3x}$.
Therefore,$u'=\dfrac{e^{3x}}{3}+C_1$;\;
$u=\dfrac{e^{3x}}{9}+C_1x+C_2$, and
$y=uxe^{-x}=\dfrac{xe^{2x}}{9}+xe^{-x}(C_1x+C_2)$, or
$y=\dfrac{xe^{2x}}{9}+xe^{-x}(c_1+c_2x)$. As a byproduct,
$\{xe^{-x},x^2e^{-x}\}$ is a fundamental set of solutions of the
complementary equation.


\item
If $y=ue^x$, then $y'=(u'+u)e^x$ and $y''=(u''+2u'+u)e^x$, so
$(1-2x)y''+2y'+(2x-3)y=e^x\left[(1-2x)u''+(4-4x)u'\right]=(1-4x+4x^2)e^x$
if $u'=z$, where (A) $z'+\dst\frac{4-4x}{1-2x}z=1-2x$. Since
$\dst{\int\frac{4-4x}{1-2x}\,dx}=
\dst{\int{\left(2+\frac{2}{1-2x}\right)}\,dx}=2x-\ln|1-2x|$,
$z_1=(1-2x)e^{-2x}$ is a solution of the complementary equation for
(A). Therefore,the solutions of (A) are of the form (B)
$z=v(1-2x)e^{-2x}$, where $v'(1-2x)e^{-2x}=(1-2x)$, so $v'=e^{2x}$.
Hence, $v=\dfrac{e^{2x}}{2}+C_1$;\;
$u'=z=\dst{\left(\frac{1}{2}+C_1e^{-2x}\right)(1-2x)}$ (see (B));\;
$u=-\dfrac{(2x-1)^2}{8}+C_1xe^{-2x}+C_2$;\;
$y=ue^x=-\dfrac{(2x-1)^2e^x}{8}+C_1xe^{-x}+C_2e^x$, or
$y=-\dfrac{(2x-1)^2e^x}{8}+c_1e^x+c_2xe^{-x}$. As a byproduct,
$\{e^x,xe^{-x}\}$ is a fundamental set of solutions of the
complementary equation.


\item
If $y=ue^{-x}$, then $y'=(u'-u)e^{-x}$ and $y''=(u''-2u'+u)e^{-x}$, so
$2xy''+(4x+1)y'+(2x+1)y=e^{-x}(2xu''+u')=3x^{1/2}e^{-x}$ if $u'=z$,
where (A) $z'+\dst\frac{1}{2x}z=\dst\frac{3}{2}x^{-1/2}$. Since $\dst{\int
\frac{1}{2x}\,dx}=\dst\frac{1}{2}\ln|x|$, $z_1=x^{-1/2}$ is a solution of
the complementary equation for (A). Therefore,the solutions of (A) are
of the form (B) $z=vx^{-1/2}$, where
$v'x^{-1/2}=\dst\frac{3}{2}x^{-1/2}$, so $v'=\dst\frac{3}{2}$. Hence,
$v=\dst\frac{3x}{2}+C_1$;\ $u'=z=\dst\frac{3}{2}x^{1/2}+C_1x^{-1/2}$ (see
(B));\ $u=x^{3/2}+2C_1x^{1/2}+C_2$;\;
$y=ue^{-x}=e^{-x}(x^{3/2}+2C_1x^{1/2}+C_2)$, or
$y=e^{-x}(x^{3/2}+c_1+c_2x^{1/2})$ As a byproduct, is a
$\{e^{-x},x^{1/2}e^{-x}\}$ fundamental set of solutions of the
complementary equation.


\item
If $y=ux^{1/2}$, then $y'=u'x^{1/2}+\dst\frac{u}{2x^{1/2}}$ and
$y''=u''x^{1/2}+\dst\frac{u'}{ x^{1/2}}+\cdots$ so
$4x^2y''-4x(x+1)y'+(2x+3)y=4x^{5/2}(u''-u')=4x^{5/2}e^{2x}$ if $u'=z$,
where (A) $z'-z=e^{2x}$. Since $z_1=e^x$ is a solution of the
complementary equation for (A), the solutions of (A) are of the form
(B) $z=ve^x$, where $v'e^x=e^{2x}$, so $v'=e^x$. Hence, $v=e^x+C_1$;\;
$u'=z=e^{2x}+C_1e^x$ (see (B));\ $u=\dfrac{e^{2x}}{2}+C_1e^x+C_2$;\;
$y=ux^{1/2}=x^{1/2}\dst{\left(\frac{e^{2x}}{2}+C_1e^x+C_2\right)}$, or
$y=x^{1/2}\dst{\left(\frac{e^{2x}}{2}+c_1+c_2e^x\right)}$. As a
byproduct, $\{x^{1/2},x^{1/2}e^x\}$ is a fundamental set of solutions
of the complementary equation.


\item
If $y=ue^x$, then $y'=(u'+u)e^x$ and $y''=(u''+2u'+u)e^x$, so
$xy''+(2-2x)y'+(x-2)y=e^x(xu''+2u')=0$ if $\dst\frac{u''}{
u'}=-\dst\frac{2}{ x}$;\ $\ln|u'|=-2\ln|x|+k$; $u'=\dst\frac{C_1}{ x^2}$;\;
$u=-\dst\frac{C_1}{ x}+C_2$. Therefore,$y=ue^x=e^x\dst{\left(-\frac{C_1}{
x}+C_2\right)}$ is the general solution, and $\{e^x,e^x/x\}$ is a
fundamental set of solutions.


\item
If $y=u\ln|x|$, then $y'=u'\ln|x|+\dst\frac{u}{ x}$ and
$y''=u''\ln|x|+\dst\frac{2u'}{ x}\cdots$, so $x^2(\ln |x|)^2y''-(2x \ln
|x|)y'+(2+\ln |x|)y=x^2(\ln|x|)^3u''=0$ if $u''=0$;\ $u'=C_1$;\;
$u=C_1x+C_2$. Therefore,$y=u\ln|x|=(C_1x+C_2)\ln|x|$ is the general
solution, and $\{\ln|x|,x\ln|x|\}$ is a fundamental set of solutions.


\item
If $y=ue^x$, then $y'=u'e^x+ue^x$ and $y''=u''e^x+2u'e^x+ue^x$, so
$xy''-(2x+2)y'+(x+2)y=e^x(xu''-2u')=0$ if $\dst\frac{u''}{
u'}=\dst\frac{2}{ x}$;\ $\ln|u'|=2\ln|x|+k$;\ $u'=C_1x^2$;\;
$u=\dst\frac{C_1x^3}{3}+C_2$. Therefore,
$y=ue^x=\dst\left(\frac{C_1x^3}{3}+C_2\right)e^x$ is the general
solution, and $\{e^x,x^3e^x\}$ is a fundamental set of solutions.


\item
If $y=ux\sin x$, then $y'=u'x\sin x+u(x\cos x+\sin x)$ and
$y''=u''x\sin x+2u'(x\cos x+\sin x)+\cdots$, so
$x^2y''-2xy'+(x^2+2)y=(x^3\sin x)u''+2(x^3\cos x)u'=0$ if
$\dst\frac{u''}{ u'}=-\dst\frac{2\cos x}{\sin x}$;\ $\ln|u'|=-2\ln|\sin
x|+k$;\ $u'=\dst\frac{C_1}{\sin^2x}$;\ $u=-C_1\cot x+C_2$. Therefore,
$y=ux\sin x=x(-C_1\cos x+C_2\sin x)$ is the general solution, and
$\{x\sin x,x\cos x\}$ is a fundamental set of solutions.


\item
If $y=ux^{1/2}$, then $y'=u'x^{1/2}+\dst\frac{u}{2x^{1/2}}$ and
$y''=u''x^{1/2}+\dst\frac{u'}{ x^{1/2}}+\cdots$ so $4x^2(\sin
x)y''-4x(x\cos x+\sin x)y'+(2x\cos x+3\sin x)y=4x^{5/2}(u''\sin
x-u'\cos x)=0$ if $\dst\frac{u''}{ u'}=\dst\frac{\cos x}{\sin x}$;\;
$\ln|u'|=\ln|\sin x|+k$;\ $u'=C_1\sin x$;\ $u=-C_1\cos x+C_2$.
Therefore,$y=ux^{1/2}=(-C_1\cos x+C_2)x^{1/2}$ is the general
solution, and $\{x^{1/2},x^{1/2}\cos x\}$ is a fundamental set of
solutions.


\item
If $y=\dst\frac{u}{ x}$, then $y'=\dst\frac{u'}{ x}-\dst\frac{u}{ x^2}$ and
$y''=\dst\frac{u''}{ x}-\dst\frac{2u'}{ x^2}+\cdots$, so
$(2x+1)xy''-2(2x^2-1)y'-4(x+1)y=(2x+1)u''-(4x+4)u'=0$ if
$\dst\frac{u''}{ u'}=\dst\frac{4x+4}{2x+1}=2+\dst\frac{2}{2x+1}$;\;
$\ln|u'|=2x+\ln|2x+1|+k$; $u'=C_1(2x+1)e^{2x}$;\ $u=C_1xe^{2x}+C_2$.
Therefore,$y=\dst\frac{u}{ x}=C_1e^{2x}+\dst\frac{C_2}{ x}$ is the general
solution, and $\{1/x,e^{2x}\}$ is a fundamental set of solutions.


\item
If $y=ue^{2x}$, then $y'=(u'+2u)e^{2x}$ and $y''=(u''+4u'+4u)e^{2x}$,
so $xy''-(4x+1)y'+(4x+2)y =e^{2x}(xu''-u')=0$ if $\dst\frac{u''}{
u'}=\dst\frac{1}{ x}$;\ $\ln|u'|=\ln|x|+k$; $u'=C_1x$;\;
$u=\dst\frac{C_1x^2}{ 2}+C_2$. Therefore,
$y=ue^{2x}=e^{2x}\dst{\left(\frac{C_1x^2}{2}+C_2\right)}$ is the general
solution, and $\{e^{2x},x^2e^{2x}\}$ is a fundamental set of
solutions.




\item
If $y=ue^{2x}$, then $y'=(u'+2u)e^{2x}$ and $y''=(u''+4u'+4u)e^{2x}$,
so
$(3x-1)y''-(3x+2)y'-(6x-8)y=e^{2x}\left[(3x-1)u''+(9x-6)u'\right]=0$
if $\dst\frac{u''}{ u'}=-\dst\frac{9x-6}{3x-1}=-3+\dst\frac{3}{3x-1}$.
Therefore,$\ln|u'|=-3x+\ln|3x-1|+k$, so $u'=C_1(3x-1)e^{-3x}$,
$u=-C_1xe^{-3x}+C_2$. Therefore,the general solution is
$y=ue^{2x}=-C_1xe^{-x}+C_2e^{2x}$, or (A) $y=c_1e^{2x}+c_2xe^{-x}$.
Now $y(0)=2\Rightarrow c_1=2$. Differentiating (A) yields
$y'=2c_1e^{2x}+c_2(e^{-x}-xe^{-x})$. Now $y'(0)=3\Rightarrow
3=2c_1+c_2$, so $c_2=-1$ and $y=2e^{2x}-xe^{-x}$.


\item
If $y=ux$, then $y'=u'x+u$  and $y''=u''x+2u'$, so
$x^2y''+2xy'-2y=x^3u''+4x^2u'=x^2$
if $u'=z$, where  (A) $z'+\dst\frac{4}{ x}z=\dst\frac{1}{ x}$.
Since $\dst{\int\frac{4}{ x}\,dx}=4\ln|x|$,
$z_1=\dst\frac{1}{ x^4}$
is a solution of the complementary equation for  (A). Therefore,the
solutions of (A) are of the form (B) $z=\dst\frac{v}{ x^4}$, where
$\dst\frac{v'}{ x^4}=\dst\frac{1}{ x}$, so $v'=x^3$.
Hence, $v=\dst\frac{x^4}{ 4}+C_1$;\;
$u'=z=\dst\frac{1}{4}+\dst\frac{C_1}{ x^4}$ (see (B));\;
$u=\dst\frac{x}{4}-\dst\frac{C_1}{3x^3}+C_2$.
 Therefore,the general solution is
$y=ux=\dst\frac{x^2}{4}-\dst\frac{C_1}{3x^2}+C_2x$, or
(C) $y=\dst\frac{x^2}{4}+c_1x+\dst\frac{c_2}{ x^2}$.
Differentiating (C) yields
 $y'=\dst\frac{x}{2}+c_1-2\dst\frac{c_2}{ x^3}.$
Now $y(1)=\dst\frac{5}{4},\ y'(1)=\dst\frac{3}{2}\Rightarrow
c_1+c_2=1,\ c_1-2c_2=1$, so $c_1=1$, $c_2=0$ and
$y=\dst\frac{x^2}{4}+x$.




\item
If $y=uy_1$, then $y'=u'y_1+uy_1'$  and $y''=u''y_1+2u'y_1'+uy_1''$,
so  $y''+p_1(x)y'+p_2(x)y=y_1u''+(2y_1'+p_1y_1)u'=0$ if $u$
is any function such that
(B) $\dst\frac{u''}{ u'}=-2\dst\frac{y_1'}{ y_1}-p_1$. If
 $\ln|u'(x)|=-2\ln|y_1(x)|-\dst{\int_{x_0}^x} p_1(t)\,dt$,
then $u$ satisfies (B); therefore, if
(C) $u'(x)=\dst\frac{1}{ y_1^2(x)}\exp\left(-\int_{x_0}^x
p_1(s)\,ds\right)$, then $u$ satisfies (B).
Since $u(x)=\dst{\int^x_{x_0}\frac{1}{ y^2_1(t)} \exp
\left(-\int^t_{x_0}p_1(s) ds\right)}$ satisfies (C), $y_2=uy_1$ is a
solution of (A) on $(a,b)$. Since $\dst\frac{y_2}{ y_1}=u$ is
nonconstant, Theorem~5.1.6 implies that $\{y_1,y_2\}$ is a
fundamental set of solutions of (A) on $(a,b)$.


\item
\begin{alist}
\item The associated linear equation is (A) $z''+k^2z=0$, with
characteristic polynomial $p(r)=r^2+k^2$. The general solution of (A)
is $z=c_1\cos kx+c_2\sin kx$. Since $z'=-kc_1\sin kx+kc_2\cos kx$,
$y=\dst\frac{z'}{ z}=\dst\frac{-kc_1\sin kx+kc_2\cos kx}{ c_1 \cos
kx+c_2\sin kx}$.

\item The associated linear equation is (A) $z''-3z'+2z=0$, with
characteristic polynomial $p(r)=r^2-3r+2=(r-1)(r-2)$. The general
solution of (A) is $z=c_1e^x+c_2e^{2x}$. Since $z'=c_1e^x+c_2e^{2x}$,
$y=\dst\frac{z'}{ z}=\dst\frac{c_1+2c_2e^x}{ c_1+c_2e^x}$.

\item The associated linear equation is (A) $z''+5z'-6z=0$, with
characteristic polynomial $p(r)=r^2+5r-6=(r+6)(r-1)$. The general
solution of (A) is $z=c_1e^{-6x}+c_2e^x$. Since
$z'=-6c_1e^{-6x}+c_2e^x$, $y=\dst\frac{z'}{ z}=\dfrac{-6c_1+c_2e^{7x}}{
c_1+c_2e^{7x}}$.

\item The associated linear equation is (A) $z''+8z'+7z=0$, with
characteristic polynomial $p(r)=r^2+8r+7=(r+7)(r+1)$. The general
solution of (A) is $z=c_1e^{-7x}+c_2e^{-x}$. Since
$z'=-7c_1e^{-7x}-2c_2e^{-x}$, $y=\dst\frac{z'}{
z}=-\dfrac{7c_1+c_2e^{6x}}{ c_1+c_2e^{6x}}$.

\item The associated linear equation is (A) $z''+14z'+50z=0$, with
characteristic polynomial $p(r)=r^2+14r+50=(r+7)^2+1$. The general
solution of (A) is $z=e^{-7x}(c_1\cos x+c_2\sin x)$. Since
$z'=-7e^{-7x}(c_1\cos x+c_2\sin x)+ e^{-7x}(-c_1\sin x+c_2\cos
x)=-(7c_1-c_2)\cos x-(c_1+7c_2)\sin x$, $y=\dst\frac{z'}{
z}=-\dfrac{(7c_1-c_2)\cos x+(c_1+7c_2)\sin x}{ c_1\cos x+c_2\sin x}$.

\item The given equation is equivalent to (A)
$y'+y^2-\dst\frac{1}{6}y-\dst\frac{1}{6}=0$. The associated linear
equation is (B) $z''-\dst\frac{1}{6}z'-\dst\frac{1}{6}z=0$, with
characteristic polynomial
$p(r)=r^2-\dst\frac{1}{6}r-\dst\frac{1}{6}=\dst{\left(r+\dst\frac{1}{3}\right)
\left(r-\dst\frac{1}{2}\right)}$. The general solution of (B) is
$z=c_1e^{-x/3}+c_2e^{x/2}$. Since
$z'=-\dst\frac{c_1}{3}e^{-x/3}+\dst\frac{c_2}{2}e^{x/2}$, $y=\dst\frac{z'}{
z} =\dfrac{-2c_1+3c_2e^{5x/6}}{ 6(c_1+c_2e^{5x/6})}$.

\item The given equation is equivalent to (A)
$y'+y^2-\dst\frac{1}{3}y+\dst\frac{1}{36}=0$. The associated linear
equation is (B) $z''-\dst\frac{1}{3}z'+\dst\frac{1}{36}z=0$, with
characteristic polynomial $p(r)=r^2-\dst\frac{1}{3}r+\dst\frac{1}{36}=
\dst{\left(r-\dst\frac{1}{6}\right)^2}$. The general solution of (B) is
$z=e^{x/6}(c_1+c_2x)$. Since
$z'=\dfrac{e^{x/6}}{6}(c_1+c_2x)+c_2e^{x/6}=\dfrac{e^{x/6}}{6}
(c_1+c_2(x+6))$, $y=\dst\frac{z'}{
z}=\dfrac{c_1+c_2(x+6)}{6(c_1+c_2x)}$.
\end{alist}

\item
\begin{alist}
\item Suppose that $z$ is a solution of (B)
and let $y=\dst\frac{z'}{ rz}$. Then
(D) $\dst\frac{z''}{ rz}+\left[p(x)-\frac{r'(x)}{ r(x)}\right]
y+q(x)=0$
and  $y'=\dst\frac{z''}{ rz}-\dst{\frac{1}{ r}\left(\frac{z'}{
z}\right)^2}-\dst\frac{r'z'}{ r^2z}=\dst\frac{z''}{
rz}-ry^2-\dst\frac{r'}{ r}y$, so
$\dst\frac{z''}{ rz}=y'+ry^2+\dst\frac{r'}{ r}y$. Therefore,
(D) implies that
$y$
satisfies (A). Now suppose that $y$ is a solution of (A) and let $z$
be any function such that $z'=ryz$. Then $z''=r' yz+r y'z+ryz'=
\dst\frac{r'}{ r}z'+
(y'+ry^2)rz
=\dst\frac{r'}{ r}z'-(p(x)y+q(x))rz$, so
$z''-\dst\frac{r'}{ r}z'+p(x)ryz+q(x)rz=0$,
which implies that $z$ satisfies (B), since $ryz=z'$.

\item  If $\{z_1,z_2\}$ is a fundamental set of solutions of
(B) on  $(a,b)$, then $z=c_1z_1+c_2z_2$ is the general solution of
(B) on $(a,b)$. This and \part{a} imply that (C) is the general
solution of (A) on $(a,b)$.
\end{alist}
\end{sectionSolutions}

\section{Variation of Parameters}

\begin{sectionSolutions}
\item
\begin{align*}
y_p&=u_1\cos2x+u_2\sin2x \tag{A}\\
\phantom{-2}u_1'\cos2x+\phantom{2}u_2'\sin2x&=0\tag{B}\\
-2u_1'\sin2x+2u_2'\cos 2x&=\sin2x\sec^2x.\tag{C}
\end{align*}
Multiplying (B) by $2\sin2x$ and (C)
by $\cos2x$ and adding the resulting equations yields $2u_2'=\tan2x$,
so $u_2'=\dst\frac{\tan2x}{2}$. Then (B) implies that
$u_1'=-u_2'\tan(2x)=-\dst\frac{\tan^22x}{2}=\dst\frac{1-\sec^22x}{2}$.
Therefore,$u_1=\dst\frac{x}{2}-\dst\frac{\tan2x}{4}$ and
$u_2=-\dfrac{\ln|\cos2x|}{4}$. Now (A) yields $y_p=-\dfrac{\sin
2x\ln|\cos 2x|}{4} +\dst\frac{x\cos 2x}{2}-\dst\frac{\sin2x}{4}$. Since
$\sin2x$ satisfies the complementary equation we redefine
$y_p=-\dfrac{\sin 2x\ln|\cos 2x|}{4} +\dst\frac{x\cos 2x}{2}$.


\item
\begin{align*}
y_p&=u_1e^x\cos x+u_2e^x\sin x\tag{A}\\
u_1'e^x\cos x+u_2'e^x\sin x&=0\tag{B} \\
u_1'(e^x\cos x-e^x\sin x)+u_2'(e^x\sin x+e^x\cos x)&=3e^x\sec x. \tag{C}
\end{align*}
Subtracting (B)  from (C) and
cancelling  $e^x$ from the resulting equations yields
\begin{align*}
\phantom{-}u_1'\cos x+u_2'\sin x&=0\tag{D}\\
-u_1'\sin x+u_2'\cos x&=3\sec x.\tag{E}
\end{align*}
 Multiplying (D) by $\sin x$ and
 (E) by $\cos x$
and adding the results yields $u_2'=3$. From (D),
$u_1'=-u_2'\tan x=-3\tan x$. Therefore $u_1=3\ln|\cos x|$, $u_2=3x$.
Now (A) yields $y_p=3e^x(\cos x \ln |\cos x|+x\sin x)$.


\item
\begin{align*}
y_p&=u_1e^x+u_2e^{-x}\tag{A}\\
u_1'e^x+u_2'e^{-x}&=0\tag{B}\\
u_1'e^x-u_2'e^{-x}&=\dfrac{4e^{-x}}{ 1+e^{-2x}}.\tag{C}
\end{align*}
Adding (B) to (C) yields
$2u_1'e^x=\dfrac{4e^{-x}}{ 1+e^{-2x}}$, so $u_1'=\dfrac{2e^{-2x}}{
1-e^{-2x}}$. From (B),
$u_2'=-e^{2x}u_1'=-\dst\frac{2}{1-e^{-2x}}=\dfrac{2e^{2x}}{1-e^{2x}}$.
Using the substitution $v=e^{-2x}$ we integrate $u_1'$ to obtain
$u_1=\ln(1-e^{-2x})$. Using the substitution $v=e^{2x}$ we integrate
$u_2'$ to obtain $u_1=\ln(1-e^{2x})$. Now (A) yields
$y_p=e^x\ln(1-e^{-2x})-e^{-x}\ln(e^{2x}-1)$.


\item
\begin{align*}
y_p&=u_1e^x+u_2\dst\frac{e^x}{ x}\tag{A}\\
u_1'e^x+u_2'\frac{e^x}{ x}&=0\tag{B}\\
u_1'e^x+u_2'\left(\frac{e^x}{ x}-\frac{e^x}{ x^2}\right)&=\frac{e^{2x}}{
x}.\tag{C}
\end{align*}
Subtracting (B) from (C) yields
$-\dst\frac{u_2'e^x}{ x^2}=\dfrac{e^{2x}}{ x}$, so $u_2'=-xe^x$. From
(B), $u_1'=\dst\frac{u_2'}{ x}=e^x$. Therefore
$u_1=e^x$, $u_2=-xe^x+e^x$. Now (A) yields $y_p=\dfrac{e^{2x}}{ x}$.


\item
\begin{align*}
y_p&=u_1e^{-x^2}+u_2xe^{-x^2}\tag{A}\\
u_1'e^{-x^2} +u_2'xe^{-x^2}&=0\tag{B}\\
-2xu_1'e^{-x^2}+u_2'(e^{-x^2}-2x^2e^{-x^2})&=
4e^{-x(x+2)}.\tag{C}
\end{align*}
Multiplying (B) by $2x$ and adding the result to
(C) yields $u_2'e^{-x^2}=4e^{-x(x+2)}$, so
$u_2'=4e^{-2x}$. From (B), $u_1'=-u_2'x=-4xe^{-2x}$.
Therefore $u_1=(2x+1)e^{-2x}$, $u_2=-2e^{-2x}$. Now (A) yields
$y_p=e^{-x(x+2)}$.


\item
\begin{align*}
y_p&=u_1x+u_2x^3\tag{A}\\
u_1'x+\phantom{3}u_2'x^3&=0\tag{B}\\
u_1'\phantom{x}+3u_2'x^2&=\dst\frac{2x^4\sin x}{
x^2}=2x^2\sin x \tag{C}.
\end{align*}
Multiplying (B) by $\dst\frac{1}{ x}$ and subtracting
the result from (C) yields $2x^2u_2'=2x^2\sin x$, so
$u_2'=\sin x$. From (B), $u_1'=-u_2'x^2=-x^2\sin x$.
Therefore $u_1=(x^2-2)\cos x-2x\sin x$, $u_2=-\cos x$. Now (A) yields
$y_p=-2x^2\sin x-2x\cos x$.


\item
\begin{align*}
y_p&=u_1\cos\sqrt x+u_2\sin\sqrt x\tag{A}\\
u_1'\cos\sqrt x+u_2'\sin\sqrt x&=0\tag{B}\\
-u_1'\frac{\sin\sqrt x}{2\sqrt
x}+u_2'\frac{\cos\sqrt x}{2\sqrt x}&=\frac{\sin\sqrt x}{4x}\tag{C}
\end{align*}
Multiplying (B) by $\dst\frac{\sin\sqrt x}{2\sqrt x}$
and (C) by $\cos\sqrt x$ and adding the resulting
equations yields $\dst\frac{u_2'}{2\sqrt x}=\dfrac{\sin\sqrt x\cos\sqrt
x}{4x}$, so $u_2'=\dst\frac{\sin\sqrt x\cos\sqrt x}{2\sqrt x}$. From
(B), $u_1'=-u_2'\tan\sqrt x=-\dfrac{\sin^2\sqrt
x}{2\sqrt x}$. Therefore, $u_1=\dfrac{\sin\sqrt x\cos\sqrt
x}{2}-\dst\frac{\sqrt x}{2}$, $u_2=\dst\frac{\sin^2\sqrt x}{2}$. Now
(A) yields $y_p=\dst\frac{\sin\sqrt x}{2}-\dfrac{\sqrt x\cos\sqrt
x}{2}$. Since $\sin\sqrt x$ satisfies the complementary equation we
redefine $y_p=-\dst\frac{\sqrt x\cos\sqrt x}{2}$.


\item
\begin{align*}
y_p&=u_1x^a+u_2x^a\ln x\tag{A}\\
u_1'x^a+u_2'x^a\ln x&=0\tag{B}\\
au_1'x^{a-1}+u_2'(ax^{a-1}\ln x+x^{a-1})&=\frac{x^{a+1}}{
x^2}=x^{a-1}\tag{C}
\end{align*}
Multiplying (B) by $\dst\frac{a}{ x}$ and subtracting
the result from (C) yields $u_2'x^{a-1}=x^{a-1}$, so
$u_2'=1$. From (B), $u_1'=-u_2'\ln x=-\ln x$.
Therefore, $u_1=x-\ln x$, $u_2=x$. Now (A) yields $y_p=x^{a+1}$.


\item
\begin{align*}
y_p&=u_1e^{x^2}+u_2e^{-x^2}\tag{A}\\
u_1'e^{x^2} +u_2'e^{-x^2}&=0\tag{B}\\
2u_1'xe^{x^2}-2u_2'xe^{-x^2}&=\frac{8x^5}{x}=8x^4. \tag{C}
\end{align*}
Multiplying (B) by $2x$ and adding the result to
(C) yields $4u_1'xe^{x^2}=8x^4$, so
$u_1'=2x^3e^{-x^2}$. From (B),
$u_2'=-u_1'e^{2x^2}=-2x^3e^{x^2}$. Therefore $u_1=-e^{-x^2}(x^2+1)$,
$u_2=-e^{x^2}(x^2-1)$. Now (A) yields $y_p=-2x^2$.


\item
\begin{align*}
y_p&=u_1\sqrt xe^{2x}+u_2\sqrt xe^{-2x}\tag{A}\\
u_1'\sqrt xe^{2x}+u_2'\sqrt xe^{-2x}&=0\tag{B}\\
u_1'e^{2x}\left(2\sqrt x+\frac{1}{2\sqrt x}\right)-u_2'e^{-2x}
\left(2\sqrt x-\frac{1}{2\sqrt x}\right)&=\frac{8x^{5/2}}{4x^2}
=2\sqrt x\tag{C}
\end{align*}
Multiplying (B) by $\dst\frac{1}{2x}$, subtracting the
result from (C), and cancelling common factors from
the resulting equations yields
\begin{align*}
u_1'e^{2x}+u_2'e^{-2x}&=0\tag{D}\\
u_1'e^{2x}-u_2'e^{-2x}&=1.\tag{E}
\end{align*}
Adding (D) to (E) yields
$2u_1'e^{2x}=1$, so $u_1'=\dfrac{e^{-2x}}{2}$. From
(D), $u_2'=-u_1'e^{4x}=-\dfrac{e^{2x}}{2}$.
Therefore, $u_1=-\dfrac{e^{-2x}}{4}$, $u_2=-\dfrac{e^{2x}}{4}$. Now
(A) yields $y_p=-\dst\frac{\sqrt x}{2}$.


\item
\begin{align*}
y_p&=u_1xe^x+u_2xe^{-x}\tag{A}\\
u_1'xe^x+u_2'xe^{-x}&=0\tag{B}\\
u_1'(x+1)e^x-u_2'(x-1)e^{-x}&=\frac{3x^4}{ x^2}=3x^2\tag{C}
\end{align*}
Multiplying (B) by $\dst\frac{1}{ x}$, subtracting the
resulting equation from (C), and cancelling common
factors yields
\begin{align*}
u_1'e^x+u_2'e^{-x}&=0\tag{D}\\
u_1'e^x-u_2'e^{-x}&=3x.\tag{E}
\end{align*}
Adding (D) to (E) yields
$2u_1'e^x=3x$, so $u_1'=\dfrac{3xe^{-x}}{2}$. From
(D), $u_2'=-u_1'e^{2x}=-\dst\frac{3xe^x}{2}$.
Therefore $u_1=-\dfrac{3e^x(x+1)}{2}$, $u_2=-\dfrac{3e^x(x-1)}{2}$.
Now (A) yields $y_p=-3x^2$.


\item
\begin{align*}
y_p&=\dst\frac{u_1}{ x}+u_2x^3\tag{A}\\
\phantom{-}\dst\frac{u_1'}{x}+\phantom{3}u_2'x^3&=0\tag{B}\\
-\dst\frac{u_1'}{ x^2}+3u_2'x^2&=\frac{x^{3/2}}{x^2}=x^{-1/2}. \tag{C}
\end{align*}
Multiplying (B) by $\dst\frac{1}{ x}$ and
adding the result to (C) yields
$4u_2'x^2=x^{-1/2}$, so $u_2'=\dfrac{x^{-5/2}}{4}$. From
(B), $u_1'=-u_2'x^4=-\dfrac{x^{3/2}}{4}$.
Therefore $u_1=-\dfrac{x^{5/2}}{10}$, $u_2=-\dfrac{x^{-3/2}}{6}$.
Now (A) yields $y_p=-\dfrac{4x^{3/2}}{15}$.


\item
\begin{align*}
y_p&=u_1x^2e^x+u_2x^3e^x\tag{A}\\
u_1'x^2e^x+u_2'x^3e^x&=0\tag{B}\\
u_1'(x^2e^x+2xe^x)+u_2'(x^3e^x+3x^2e^x)&=\frac{2xe^x}{
x^2}=\dst\frac{2e^x}{ x}.\tag{C}
\end{align*}
Subtracting (B) from (C)
and cancelling common factors in the resulting equations
yields
\begin{align*}
\phantom{2x}u_1'+\phantom{3}u_2'x&=0\tag{D}\\
2u_1'x+3u_2'x^2&=\frac{2}{ x}.\tag{E}
\end{align*}
Multiplying (D) by $2x$ and subtracting the result
from (E) yields $x^2u_2'=\dst\frac{2}{ x}$, so
$u_2'=\dst\frac{2}{ x^3}$. From (D),
$u_1'=-u_2'x=-\dst\frac{2}{ x^2}$. Therefore $u_1=\dst\frac{2}{ x}$,
$u_2=-\dst\frac{1}{ x^2}$. Now (A) yields $y_p=xe^x$.


\item
\begin{align*}
y_p&=u_1x+u_2e^x\tag{A}\\
u_1'x+u_2'e^x&=0\tag{B}\\
u_1'+u_2'e^x&=\frac{2(x-1)^2e^x}{x-1}=2(x-1)e^x.\tag{C}
\end{align*}
Subtracting (B) from (C)
yields $u_1'(1-x)=2(x-1)e^x$, so $u_1'=-2e^x$. From
(B), $u_2'=-u_1'xe^{-x}=2x$. Therefore,
$u_1=-2e^x$, $u_2=x^2$. Now (A) yields $y_p=xe^x(x-2)$.



\item
\begin{align*}
y_p&=u_1e^{2x}+u_2xe^{-x}\tag{A}\\
u_1'e^{2x}+u_2'xe^{-x}&=0\tag{B}\\
2u_1'e^{2x}+u_2'(e^{-x}-xe^{-x})&=\frac{(3x-1)^2e^{2x}}{
3x-1}=(3x-1)e^{2x}.\tag{C}
\end{align*}
Multiplying (B) by $2$ and subtracting the result
from (C) yields $u_2'(1-3x)e^{-x}=(3x-1)e^{2x}$, so
$u_2'=-e^{3x}$. From (B), $u_1'=-u_2'xe^{-3x}=x$.
Therefore $u_1=\dst\frac{x^2}{2}$, $u_2=-\dfrac{e^{3x}}{3}$. Now (A)
yields $y_p=\dfrac{xe^{2x}(3x-2)}{3}$. The general solution of the
given equation is $y=\dfrac{xe^{2x}(3x-2)}{3}+c_1e^{2x}+c_2xe^{-x}$.
Differentiating this yields
$y'=\dfrac{e^{2x}(3x^2+x-1)}{3}+2c_1e^{2x}+c_2(1-x)e^{-x}$. Now
$y(0)=1,\ y'(0)=2\Rightarrow c_1=1,\ 2=-\dst\frac{1}{3} +2c_1+c_2$, so
$c_2=\dst\frac{1}{3}$, and
$y=\dfrac{e^{2x}(3x^2-2x+6)}{6}+\dfrac{xe^{-x}}{3}$.

\item
\begin{align*}
y_p&=u_1(x-1)e^x+u_2(x-1)\tag{A}\\
u_1'(x-1)e^x+u_2'(x-1)&=0 \tag{B}\\
u_1'xe^x+u_2'&=\frac{(x-1)^3e^x}{(x-1)^2}=(x-1)e^x.\tag{C}
\end{align*}
From (B), $u_1'=-u_2'e^{-x}$. Substituting this into
(C) yields $-u_2'(x-1)=(x-1)e^x$, so $u_2'=-e^x$,
$u_1'=1$. Therefore $u_1=x$, $u_2=e^x$. Now (A) yields
$y_p=e^x(x-1)^2$. The general solution of the given equation is
$y=(x-1)^2e^x+c_1(x-1)e^x+c_2(x-1)$. Differentiating this yields
$y'=(x^2-1)e^x+c_1xe^x+c_2$. Now $y(0)=4,\ y'(0)=-6\Rightarrow
4=1-c_1-c_2,\ -6=-1+c_2$, so $c_1=2, c_2=-5$ and
$y=(x^2-1)e^x-5(x-1)$.


\item
\begin{align*}
y_p&=u_1x+\dst\frac{u_2}{ x^2}\tag{A}\\
u_1'x+\phantom{2}\frac{u_2'}{ x^2}&=0\tag{B}\\
u_1'\phantom{x}-\frac{2u_2'}{ x^3}&=-\frac{2x^2}{x^2}=-2.\tag{C}
\end{align*}
Multiplying (B) by $\dst\frac{2}{ x}$ and adding the
result to (C) yields $3u_1'=-2$, so
$u_1'=-\dst\frac{2}{3}$. From (B),
$u_2'=-u_1'x^3=\dst\frac{2x^3}{3}$. Therefore $u_1=-\dst\frac{2x}{3}$,
$u_2=\dst\frac{x^4}{6}$. Now (A) yields $y_p=-\dst\frac{x^2}{2}$. The
general solution of the given equation is
$y=-\dst\frac{x^2}{2}+c_1x+\dst\frac{c_2}{ x^2}$. Differentiating this
yields $y'=-x+c_1-\dst\frac{2c_2}{ x^3}$. Now $y(1)=1,\
y'(1)=-1\Rightarrow 1=-\dst\frac{1}{2}+c_1+c_2,\ -1=-1+c_1-2c_2$, so
$c_1=1, c_2=\dst\frac{1}{2}$, and
$y=-\dst\frac{x^2}{2}+x+\dst\frac{1}{2x^2}$.




\item
Since $\overline y=y_p-a_1y_1-a_2y_2$,
\begin{align*}
P_0(x)\overline y''+P_1(x)\overline y'+P_2(x)\overline y&=
P_0(x)(y_p-a_1y_1-a_2y_2)''\\
&+P_1(x)(y_p-a_1y_1-a_2y_2)'\\
&+P_2(x)(y_p-a_1y_1-a_2y_2)\\
&=(P_0(x)y_p''+P_1(x)y_p'+P_2(x)y_p)\\
&-a_1\left[P_0(x)y_1''+P_1(x)y_1'+P_2(x)y_1\right]\\
&-a_2\left[P_0(x)y_2''+P_1(x)y_2'+P_2(x)y_2\right]\\
&=F(x)-a_1\cdot0-a_2\cdot0=F(x);
\end{align*}
 hence $\overline y$
is a particular solution of (A).

\item
\setcounter{equation}{3}
\begin{alist}
\item $y_p=u_1e^x+u_2e^{-x}$ is a solution of (A) on
$(a,\infty)$ if $u_1'e^x+u_2'e^{-x}=0$ and $u_1'e^x-u_2'e^{-x}
=f(x)$. Solving these two equations yields
$u_1'=\dfrac{e^{-x}f}{2}$,
$u_2'=-\dst\frac{e^xf}{2}$. The functions
$u_1(x)=\dst\frac{1}{2}\dst{\int_0^x}e^{-t}f(t)\,dt$
and $u_2(x)=-\dst\frac{1}{2}\dst{\int_0^x}e^tf(t)\,dt$ satisfy
these conditions. Therefore,
\begin{align*}
y_p(x)&=\frac{e^x}{2}\int_0^xe^{-t}f(t)\,dt
-\frac{e^{-x}}{2}\int_0^xe^tf(t)\,dt \\
&=\frac{1}{2}\int_0^xf(t)\left(e^{(x-t)-e^{-(x-t)}}\right)\,dt
=\int_0^x f(t)\sinh(x-t)\,dt.
\end{align*}
is a particular solution of $y''-y=f(x)$. Differentiating
$y_p$ yields
\begin{align*}
y_p'(x)&=\frac{e^x}{2}\int_0^xe^{-t}f(t)\,dt+\frac{e^{x}}{2}e^{-x}
+\frac{e^{-x}}{2}\int_0^xe^tf(t)\,dt-\frac{e^{-x}}{2}e^x\\
&=\frac{e^x}{2}\int_0^xe^{-t}f(t)\,dt
+\frac{e^{-x}}{2}\int_0^xe^tf(t)\,dt\\
&=\frac{1}{2}\int_0^xf(t)\left(e^{(x-t)}+e^{-(x-t)}\right)\,dt
=\int_0^x f(t)\cosh(x-t)\,dt.
\end{align*}
Since $y_p(x_0)=y_p'(x_0)=0$, the solution of the initial value
problem is
\begin{align*}
y&=y_p+k_0\cosh x+k_1\sinh x\\
&=k_0\cosh x+k_1\sinh x+\int_0^x\sinh (x-t)f(t)\,dt.
\end{align*}
The derivative of the solution is
\begin{align*}
y'&=y_p'+k_0\sinh x+k_1\cosh x\\
&=k_0\sinh x+k_1\cosh x+\int_0^x\cosh (x-t) f(t)\,dt.
\end{align*}
\end{alist}
\end{sectionSolutions}

\chapter{Applications of Linear Second Order Equations}

\section{Spring Problems I}
\begin{sectionSolutions}
\item
Since $\dst\frac{k}{ m}=\dst\frac{g}{\Delta l}=\dst\frac{32}{.1}=320$ the
equation of motion is (A) $y''+320y=0$. The general solution of (A) is
$y=c_1\cos8\sqrt5t+c_2\sin8\sqrt5t$, so $y'=8\sqrt5(-c_1\sin8\sqrt5
t+c_2\cos8\sqrt5t)$. Now $y(0)=-\dst\frac{1}{4}\Rightarrow
c_1=-\dst\frac{1}{4}$ and $y'(0)=-2\Rightarrow
c_2=-\dst\frac{1}{4\sqrt5}$, so $y=\dst{
-\frac{1}{4}\cos8\sqrt{5}t-\frac{1}{ 4\sqrt{5}}\sin8\sqrt{5}t}$ ft.


\item
Since $\dst\frac{k}{ m}=\dst\frac{g}{\Delta l}=\dst\frac{32}{.5}=64$ the
equation of motion is (A) $y''+64y=0$. The general solution of (A) is
$y=c_1\cos8t+c_2\sin8t$, so $y'=8(-c_1\sin8t+c_2\cos8t)$. Now
$y(0)=\dst\frac{1}{4}\Rightarrow c_1=\dst\frac{1}{4}$ and
$y'(0)=-\dst\frac{1}{2}\Rightarrow c_2=-\dst\frac{1}{16}$, so
$y=\dst{\frac{1}{4}\cos8t-\frac{1}{16}\sin8t}$ ft;\;
$R=\dfrac{\sqrt{17}}{16}$ ft;\ $\omega_0=8$ rad/s;\ $T=\pi/4$ s;
$\phi\approx\qty{-.245}{\radian}\approx\qty{-14.04}{\degree}$.


\item
Since $k=\dst\frac{mg}{\Delta l}=\dfrac{(9.8)10}{.7}=140$, the equation
of motion of the 2 kg mass is (A) $y''+70y=0$. The general solution of
(A) is $y=c_1\cos\sqrt{70}t+c_2\sin\sqrt{70}t$, so
$y'=\sqrt{70}(-c_1\sin\sqrt{70}t+c_2\cos\sqrt{70}t)$. Now
$y(0)=-\dst\frac{1}{4}\Rightarrow c_1-\dst\frac{1}{4}=$ and
$y'(0)=2\Rightarrow c_2=\dst\frac{2}{\sqrt{70}}$, so
$y=\dst{-\frac{1}{4}\cos\sqrt{70}\;t+\frac{2
}{\sqrt{70}}\sin\sqrt{70}\;t}$ m;\ $R= \dst{\frac{1}{4}\sqrt\frac{67}{
35}}$ m;\ $\omega_0=\sqrt{70}$ rad/s; $T=2\pi/\sqrt{70}$ s;\;
$\phi\approx\qty{2.38}{\radian}\approx\qty{136.28}{\degree}$.


\item
Since $\dst\frac{k}{ m}=\dst\frac{g}{\Delta l}=\dst\frac{32}{1/2}=64$ the
equation of motion is (A) $y''+64y=0$. The general solution of (A) is
$y=c_1\cos8t+c_2\sin8t$, so $y'=8(-c_1\sin8t+c_2\cos8t)$. Now
$y(0)=\dst\frac{1}{2}\Rightarrow c_1=\dst\frac{1}{2}$ and
$y'(0)=-3\Rightarrow c_2=-\dst\frac{3}{8}$, so
$y=\dst{\frac{1}{2}\cos8t-\frac{3}{8}\sin 8t}\unit{\feet}$.


\item
$m=\dst\frac{64}{32}=2$, so the equation of motion is $2y''+8y=2\sin t$,
or (A) $y''+4y=\sin t$. Let $y_p=A\cos t+B\sin t$; then $y_p''=-A\cos
t-B\sin t$, so $y_p''+4y_p=3A\cos t+3B\sin t= \sin t$ if $3A=0$,
$3B=1$. Therefore,$A=0$, $B=\dst\frac{1}{3}$, and $y_p=\dst\frac{1}{3}\sin
t$. The general solution of
(A) is (B) $y=\dst\frac{1}{3}\sin
t+c_1\cos2t+c_2\sin2t$, so $y(0)=\dst\frac{1}{2}\Rightarrow
c_1=\dst\frac{1}{2}$. Differentiating (B) yields $y'=\dst\frac{1}{3}\cos
t-2c_1\sin2t+2c_2\cos2t$, so $y'(0)=2\Rightarrow
2=\dst\frac{1}{3}+2c_2$, so $c_2=\dst\frac{5}{6}$. Therefore,
$y=\dst{\frac{1}{3}\sin t+\frac{1}{2}\cos2t+\frac{5}{6}\sin2t}$ ft.


\item
$m=\dst\frac{4}{32}=\dst\frac{1}{8}$ and $k=\dst\frac{mg}{\Delta l}=4$, so
the equation of motion is $\dst\frac{1}{8}y''+4y=\dst\frac{1}{4}\sin8t$,
or (A) $y''+32y=2\sin8t$. Let $y_p=A\cos8t+B\sin8t$; then
$y_p''=-64A\cos t-64B\sin8t$, so $y_p''+32y_p=-32A\cos8t-32B\sin8t=
2\sin8t$ if $-32A=0$, $-32B=2$. Therefore,$A=0$, $B=-\dst\frac{1}{16}$,
and $y_p=-\dst\frac{1}{16}\sin8t$. The general solution of (A) is (B)
$y=-\dst\frac{1}{16}\sin8t+c_1\cos4\sqrt2t+c_2\sin4\sqrt2t$, so
$y(0)=\dst\frac{1}{3}\Rightarrow c_1=\dst\frac{1}{3}$. Differentiating (B)
yields $y'=-\dst\frac{1}{2}\cos
8t+4\sqrt2(-c_1\sin4\sqrt2t+c_2\cos4\sqrt2t)$, so $y'(0)=-1\Rightarrow
-1=-\dst\frac{1}{2}+4\sqrt2c_2$, so $c_2=-\dst\frac{1}{8\sqrt2}$.
Therefore,$y=\dst{-\frac{1}{16}\sin8t+\frac{1}{3}\cos4\sqrt2t
-\frac{1}{8\sqrt2 }\sin 4\sqrt2t}$~ft.




\item
Since $T=\dst{\frac{2\pi}{\omega_0}=2\pi\sqrt\frac{m}{ k}}$
the period is proportional to the square root of the mass.
Therefore, doubling the mass multiplies the period by $\sqrt2$;
hence the period of the system with the \qty{20}{\gram} mass is
 $T=4\sqrt{2}$~s.


\item
$m=\dst\frac{6}{32}=\dst\frac{3}{16}$ and $k=\dst\frac{mg}{\Delta
l}=\dst\frac{6}{1/3}=18$ so the equation of motion is
$\dst\frac{3}{16}y''+18y=4\sin\omega t-6\cos\omega t$, or (A)
$y''+96y=\dst\frac{64}{3}\sin\omega t-32\cos\omega t$. The displacement
will be unbounded if $\omega=\sqrt{96}=4\sqrt6$, in which case (A)
becomes (B) $y''+96y=\dst\frac{64}{3}\sin4\sqrt6t-32\cos4\sqrt6t$. Let
\begin{align*}
y_p&=At\cos4\sqrt6t+Bt\sin4\sqrt6t; \text{ then}\\
y_p'&=(A+4\sqrt6Bt)\cos4\sqrt6t+(B-4\sqrt6At)\sin4\sqrt6t\\
y_p''&=(8\sqrt6B-96At)\cos4\sqrt6t-(8\sqrt6A+96Bt)\sin4\sqrt6t,
\qquad\text{ so}\\
 y_p''+96y_p&=8\sqrt6B\cos4\sqrt6t-8\sqrt6A\sin4\sqrt6t=
\frac{64}{3}\sin4\sqrt6t-32\cos4\sqrt6t
\end{align*}
if $8\sqrt6B=-32$, $-8\sqrt6A=\dst\frac{64}{3}$. Therefore,
$A=-\dst\frac{8}{3\sqrt6}$, $B=-\dst\frac{4}{\sqrt6}$, and
$y_p=-\dst\frac{t}{\sqrt6}\dst\left(\frac{8}{3}\cos4\sqrt6t+
4\sin4\sqrt6t\right)$.
The general solution of (B) is
\[
y=-\dst\frac{t}{\sqrt6}\left(\frac{8}{3}\cos4\sqrt6t+4\sin4\sqrt6t\right)
+c_1\cos4\sqrt6t+c_2\sin4\sqrt6t,
\tag{C}
\]
so $y(0)=0\Rightarrow c_1=0$. Differentiating
(C) yields
\begin{align*}
y'&=-\dst\left(\frac{8}{3\sqrt6}\cos4\sqrt6t+
\frac{4}{\sqrt6}\sin4\sqrt6t\right)
-4t\dst\left(-\frac{8}{3}\sin4\sqrt6t+
4\cos4\sqrt6t\right)\\
&+4\sqrt6(-c_1\sin4\sqrt6t+c_2\cos\sqrt6t),
\end{align*}
so $y'(0)=0\Rightarrow 0=-\dst\frac{8}{3\sqrt6}+4\sqrt6c_2$, and
$c_2=\dst\frac{1}{9}$. Therefore,
\[
y=-\dst{\frac{t}{
\sqrt6}\left(\frac{8}{3}\cos4\sqrt6t+4\sin
4\sqrt6t\right)+\frac{1}{9}\sin 4\sqrt6t}\unit{\feet}.
\]


\item
The equation of motion is (A) $y''+\omega_0^2y=0$. The general
solution of (A) is $y=c_1\cos\omega_0t+c_2\sin\omega_0t$.
Now $y(0)=y_0\Rightarrow c_1=y_0$. Since
 $y'=\omega_0(-c_1\sin\omega_0t+c_2\cos\omega_0t)$,
$y'(0)=v_0\Rightarrow c_2=\dst\frac{v_0}{\omega_0}$. Therefore,
$y=y_0\cos\omega_0 t+\dst\frac{v_0}{\omega_0}\sin\omega_0t$;
\[
R=\frac{1}{\omega_0} \sqrt{(\omega_0y_0)^2+(v_0)^2};\;
\cos\phi=\dst\frac{y_0\omega_0}{
\sqrt{(\omega_0y_0)^2+(v_0)^2}}\,;\ \sin\phi=\dst\frac{v_0}{
\sqrt{(\omega_0y_0)^2+(v_0)^2}}\,.
\]

\textbf{Discussion~6.1.1}
In Exercises~19, 20, and 21 we use
the fact that in a spring-mass system with mass $m$ and spring
constant $k$ the period of the motion is $T=2\pi\dst{\sqrt\frac{m}{
k}}$. Therefore, if we have two systems with masses $m_1$ and $m_2$
and spring constants $k_1$ and $k_2$, then the periods are related by
$\dst\frac{T_2}{ T_1}=\dst{\sqrt\frac{m_2k_1}{ m_1k_2}}$. We will use this
formula in the solutions of these exercises.


\item
Let $m_2=2m_1$. Since $k_1=k_2$, $\dst\frac{T_2}{
T_1}=\sqrt{\dst\frac{2m_1}{ m_1}}=\sqrt2$, so $T_2=\sqrt2T_1$.

\addtocounter{enumi}{-1}
\item%6.1.21
Suppose that $T_2=3T_1$. Since $m_1=m_2$, $\dst{\sqrt\frac{k_1}{
k_2}}=3$, $k_1=9k_2$.
\end{sectionSolutions}


\section{Spring Problems II}

\begin{sectionSolutions}
\item
Since $k=\dst\frac{mg}{\Delta l}=\dst\frac{16}{3.2}=5$ the equation of
motion is $\dst\frac{1}{2}y''+y'+5y=0$, or (A) $y''+2y'+10y=0$. The
characteristic polynomial of (A) is $p(r)=r^2+2r+10=(r+1)^2+9$.
Therefore,the general solution of (A) is
$y=e^{-t}(c_1\cos3t+c_2\sin3t)$, so
$y'=-y+3e^{-t}(-c_1\sin3t+c_2\cos3t)$. Now $y(0)=-3$ and
$y'(0)=2\Rightarrow c_1=-3$ and $2=3+3c_2$, or $c_2=-\dst\frac{1}{3}$.
Therefore,$y=\dst{-e^{-t}\left(3\cos3t+\frac{1}{3}\sin3t\right)}$ ft.
The time--varying amplitude is $\dfrac{\sqrt{82}}{3}e^{-t}$ ft.



\item
Since $k=\dst\frac{mg}{\Delta l}=\dst\frac{96}{3.2}=30$ the equation of
motion is $3y''+18y'+30y=0$, or (A) $y''+6y'+10y=0$. The
characteristic polynomial of (A) is $p(r)=r^2+6r+10=(r+3)^2+1$.
Therefore,the general solution of (A) is $y=e^{-3t}(c_1\cos t+c_2\sin
t)$, so $y'=-3y+e^{-3t}(-c_1\sin t+c_2\cos t)$. Now
$y(0)=-\dst\frac{5}{4}$ and $y'(0)=-12\Rightarrow c_1=-\dst\frac{5}{4}$
and $-12=\dst\frac{15}{4}+c_2$, or $c_2=-\dst\frac{63}{4}$. Therefore,
$y=-\dfrac{e^{-3t}}{4}(5\cos t+63\sin t)$ ft.

\item
Since $k=\dst\frac{mg}{\Delta l}=\dst\frac{8}{.32}=25$ the equation of
motion is $\dst\frac{1}{4}y''+\dst\frac{3}{2}y'+25y=0$, or (A)
$y''+6y'+100y=0$. The characteristic polynomial of (A) is
$p(r)=r^2+6r+100=(r+3)^2+91$. Therefore,the general solution of (A) is
$y=e^{-3t}(c_1\cos\sqrt{91}t+c_2\sin\sqrt{91}t)$, so
$y'=-3y+\sqrt{91}e^{-3t}(-c_1\sin\sqrt{91}t+c_2\cos\sqrt{91}t)$. Now
$y(0)=\dst\frac{1}{2}$ and $y'(0)=4\Rightarrow c_1=\dst\frac{1}{2}$ and
$4=-\dst\frac{3}{2}+\sqrt{91}c_2$, or $c_2=\dst\frac{11}{2\sqrt91}$.
Therefore,$y=\dst{\frac{1}{2}e^{-3t}\left(\cos\sqrt{91}t+\frac{11}{
\sqrt{91}}\sin\sqrt{91}t\right)}$ ft.


\item
Since $k=\dst\frac{mg}{\Delta l}=\dst\frac{20\cdot980}{5}=3920$ the
equation of motion is $20y''+400y'+3920y=0$, or (A) $y''+20y'+196y=0$.
The characteristic polynomial of (A) is
$p(r)=r^2+20r+196=(r+10)^2+96$. Therefore,the general solution of (A)
is $y=e^{-10t}(c_1\cos4\sqrt6t+c_2\sin4\sqrt6t)$, so
$y'=-10y+4\sqrt6e^{-10t}(-c_1\sin4\sqrt6t+c_2\cos4\sqrt6t)$. Now
$y(0)=9$ and $y'(0)=0\Rightarrow c_1=9$ and $0=-90+4\sqrt6c_2$, or
$c_2=\dst\frac{45}{2\sqrt6}$. Therefore,
$y=\dst{e^{-10t}\left(9\cos4\sqrt{6}t+\frac{45}{2\sqrt{6}}\sin4\sqrt{6}
t\right)}$ cm.


\item
Since $k=\dst\frac{mg}{\Delta l}=\dst\frac{32}{1}=32$ the equation of
motion is (A) $y''+3y'+32y=0$. The characteristic polynomial of (A) is
$p(r)=r^2+3r+32=\left(r+\dst\frac{3}{2}\right)^2+\dst\frac{119}{4}$.
Therefore,the general solution of (A) is
$y=e^{-3t/2}\dst{\left(c_1\cos\frac{\sqrt{119}}{2}t
+c_2\sin\frac{\sqrt{119}}{2}t\right)}$, so $y'=-\dst\frac{3}{2}y+
\dfrac{\sqrt{119}}{2}e^{-3t/2}
\dst{\left(-c_1\sin\frac{\sqrt{119}}{2}t+c_2\cos\frac{\sqrt{119}}{2}t\right)}$.
Now $y(0)=\dst\frac{1}{2}$ and $y'(0)=-3\Rightarrow c_1=\dst\frac{1}{2}$
and $-3=-\dst\frac{3}{4}+\dfrac{\sqrt{119}}{2}c_2$, or
$c_2=-\dst\frac{9}{2\sqrt{119}}$. Therefore,
$y=\dst{e^{-\frac{3}{2}t}\left(\frac{1}{2}\cos\frac{\sqrt{119}}{2} t
-\frac{9}{2\sqrt{119}}\sin\frac{\sqrt{119}}{2}t\right)}$ ft.


\item
Since $k=\dst\frac{mg}{\Delta l}=\dst\frac{2}{.32}=\dst\frac{25}{4}$ the
equation of motion is
$\dst\frac{1}{16}y''+\dst\frac{1}{8}y'+\dst\frac{25}{4}y=0$, or (A)
$y''+2y'+100y=0$. The characteristic polynomial of (A) is
$p(r)=r^2+2r+100=(r+1)^2+99$. Therefore,the general solution of (A) is
$y=e^{-t}(c_1\cos3\sqrt{11}t+c_2\sin3\sqrt{11}t)$, so
$y'=-y+3\sqrt{11}e^{-t}(-c_1\sin3\sqrt{11}t+c_2\cos3\sqrt{11}t)$. Now
$y(0)=-\dst\frac{1}{3}$ and $y'(0)=5\Rightarrow c_1=-\dst\frac{1}{3}$ and
$5=\dst\frac{1}{3}+3\sqrt{11}c_2$, or $c_2=\dst\frac{14}{9\sqrt{11}}$.
Therefore,$\dst{y=e^{-t}\left(-\frac{1}{3}\cos3\sqrt{11} t+
\frac{14}{9\sqrt{11}}\sin3\sqrt{11}t\right)}$ ft.



\item
Since $k=\dst\frac{mg}{\Delta l}=32$ the
equation of motion is
(A) $y''+12y'+32y=0$.
The characteristic polynomial of (A) is
$p(r)=r^2+12r+32=(r+8)(r+4)$. Therefore,
the general solution of (A) is $y=c_1e^{-8t}+c_2e^{-4t}$, so
$y'=-8c_1e^{-8t}-4c_2e^{-4t}$. Now $y(0)=-\dst\frac{2}{3}$ and
$y'(0)=0\Rightarrow c_1+c_2=\dst\frac{2}{3},\-8c_1-4c_2=0$, so
$c_1=-\dst\frac{2}{3}$,  $c_2=\dst\frac{4}{3}$, and
$y=-\dst\frac{2}{3}(e^{-8t}-2e^{-4t})$.


\item
Since $k=\dst\frac{mg}{\Delta l}=\dst\frac{100\cdot980}{98}=100$ the
equation of motion is
$100y''+600y'+1000y=0$, or
(A) $y''+6y'+10y=0$.
The characteristic polynomial of (A) is
$p(r)=r^2+6r+10=(r+3)^2+1$. Therefore,
the general solution of (A) is
$y=e^{-3t}(c_1\cos t+c_2\sin t)$, so
$y'=-3y+e^{-t}(-c_1\sin t+c_2\cos t)$.
Now $y(0)=10$ and $y'(0)=-100\Rightarrow c_1=10$ and
$-100=-30+c_2$, or $c_2=-70$. Therefore,
$y=e^{-3t}(10\cos t-70\sin t)$ cm.


\item
The equation of motion is
(A) $2y''+4y'+20y=3\cos4t-5\sin4t$.
The steady  state component of the solution of (A)
is of the form
$y_p=A\cos4t+B\sin4t$; therefore $y_p'=-4A\sin4t+4B\cos4t$
and $y_p''=-16A\cos4t-16B\sin4t$, so
$2y_p''+4y_p'+20y_p=(-12A+16B)\cos4t-(16A+12B)\sin4t=3\cos4t-5\sin4t$
if $-12A+16B=3$, $-16A-12B=-5$; therefore $A=\dst\frac{11}{100}$,
$B=\dst\frac{27}{100}$, and
 $y_p=\dst{\frac{11}{100}\cos4t+\frac{27}{100}\sin4t}$ cm.



\item
Since $k=\dst\frac{mg}{\Delta l}=\dst\frac{9.8}{.49}=20$ the
equation of motion is
(A) $y''+4y'+20y=8\sin2t-6\cos2t$.
The steady  state component of the solution of (A)
is of the form
$y_p=A\cos2t+B\sin2t$; therefore $y_p'=-2A\sin2t+2B\cos2t$
and $y_p''=-4A\cos2t-4B\sin2t$, so
$y_p''+4y_p'+20y_p=(16A+8B)\cos2t-(8A-16B)\sin2t=8\sin2t-6\cos2t$
if $16A+8B=-6$, $-8A+16B=8$; therefore $A=-\dst\frac{1}{2}$,
$B=\dst\frac{1}{4}$, and
$y=\dst{-\frac{1}{2}\cos2t+\frac{1}{4}\sin2t}$ m.


\item
If $e^{r_1t}(c_1+c_2t)=0$, then (A) $c_1+c_2t=0$. If $c_2=0$, then
$c_1\ne0$ (by assumption), so (A) is impossible. If $c_1\ne0$, then the
left side of (A) is strictly monotonic and therefore cannot have the
same value for two distinct values of $t$.


\item
If $y=e^{-ct/2m}(c_1\cos\omega_1t+c_2\sin\omega_1t)$, then
$y'=-\dst\frac{c}{2m}y+\omega_1
e^{-ct/2m}(-c_1\sin\omega_1t+c_2\cos\omega_1t)$, so $y(0)=y_0$ and
$y'(0)=v_0\Rightarrow c_1=y_0$ and
$v_0=-\dst\frac{cy_0}{2m}+c_2\omega_1$, so $c_2=\dst{\frac{1}{\omega_1}
\left(v_0+\frac{cy_0}{2m}t\right)}$ and $y=\dst{e^{-ct/2m}
\left(y_0\cos\omega_1t+\frac{1}{\omega_1}\left(v_0+\frac{cy_0}{
2m}\right)\sin\omega_1t\right)}$.

\item
If $y=e^{r_1t}(c_1+c_2t)$, then $y'=r_1y+c_2e^{r_1t}$, so $y(0)=y_0$
and $y'(0)=v_0\Rightarrow c_1=y_0$ and $v_0=r_1y_0+c_2$, so
$c_2=v_0-r_1y_0$. Therefore,
$y=e^{r_1t}\dst{\left(y_0+(v_0-r_1y_0)t\right)}$.
\end{sectionSolutions}

\section{The RLC Circuit}

\begin{sectionSolutions}
\item
$\dst\frac{1}{20}Q''+2Q'+100Q=0$;\, $Q''+40Q'+2000Q=0$;\,
$r^2+40r+2000=(r+20)^2+1600=0$;\,
$r=-20\pm40i$;\,
$Q=e^{-20t}(2\cos40t+c_2\sin40t)$ (since $Q_0=2$);\,
$I=Q'=e^{-20t}\left((40c_2-40)\cos40t-(20c_2+80)\sin40t\right)$;\,
$I_0=2\Rightarrow 40c_2-40=2\Rightarrow c_2=\dst\frac{21}{20}$,
so $20c_2+80=101$;\,
$I=e^{-20t}(2\cos40t-101\sin40t)$.



\item
$\dst\frac{1}{10}Q''+6Q'+250Q=0$;\, $Q''+60Q'+2500Q=0$;\,
$r^2+60r+2500=(r+30)^2+1600=0$;\,
$r=-30\pm40i$;\,
$Q=e^{-30t}(3\cos40t+c_2\sin40t)$ (since $Q_0=3$);\,
$I=Q'=e^{-30t}\left((40c_2-90)\cos40t-(30c_2+120)\sin40t\right)$;\,
$I_0=-10\Rightarrow 40c_2-90=-10\Rightarrow c_2=2$,
so $-30c_2-120=-180$;\,
$I=-10e^{-30t}(\cos40t+18\sin40t)$.


\item
$Q_p=A\cos10t+B\sin10t$;
$Q_p'=10B\cos10t-10A\sin10t$;
$Q_p''=-100A\cos10t-100B\sin10t$;
$\dst\frac{1}{10}Q_p''+3Q_p'+100Q_9=
(90A+30B)\cos10t-(30A-90B)\sin10t=5\cos10t-5\sin10t$,
so $90A+30B=5$, $-30A+90B=-5$. Therefore, $A=1/15$, $B=-1/30$,
$Q_p=\dst\frac{\cos10t}{15}-\frac{\sin10t}{30}$, and
$I_p=\dst-\frac{1}{3}(\cos10t+2\sin10t)$.


\item
$Q_p=A\cos 50t+B\sin 50t$;
$Q_p'=50B\cos 50t-50A\sin 50t$;
$Q_p''=-2500A\cos 50t-2500B\sin 50t$;
$\dst\frac{1}{10}Q_p''+2Q_p'+100Q_p
(-150A+100B)\cos 50t-(100A+150B)\sin 50t=3\cos50t-6\sin50t$,
so $-150A+100B=3$, $-100A+1500B=-6$. Therefore,$A=3/650$,
$B=12/325$, $Q_p=\dst\frac{3}{650}(\cos50t+8\sin50t)$, and
$I_p= \dst\frac{3}{13}(8\cos50t-\sin50t)$.


\item
$Q_p=A\cos 30t+B\sin 30t$;
$Q_p'=30B\cos 30t-30A\sin 30t$;
$Q_p''=-900A\cos 30t-900B\sin 30t$;
$\dst\frac{1}{20}Q_p''+4Q_p'+125Q_p=
(80A+120B)\cos 30t-(120A-80B)\sin 30t=15\cos30t-30\sin30t$, so
$80A+120B=15$, $-120A+80B=-30$,
$A=3/13$, $B=-3/104$,
$Q_p=
\dst\frac{3}{104}(8\cos 30t-\sin 30t)$, and
$I_p=-\dst\frac{45}{52}(\cos30t+8\sin30t)$.


\item
Let $\sigma=\sigma(\omega)$ be the amplitude of $I_p$.
From the solution of  Exercise~6.3.11,
$Q_p=A\cos\omega t+B\cos\omega t$, where
$A=\dfrac{(1/C-L\omega^2)U-R\omega V}{\Delta}$,
$B=\dfrac{R\omega U+(1/C-L\omega^2)V}{\Delta}$,  and
$\Delta=(1/C-L\omega^2)^2+R^2\omega^2$. Since
$I_p=Q'_p=\omega(-A\sin\omega t+B\cos\omega t)$, it follows that
$\sigma^2(\omega)=\dst{\omega^2(A^2+B^2)}=\dst\frac{U^2+V^2}{\rho(\omega)}$,
with
$\rho(\omega)=\dst\frac{\Delta}{\omega^2}=(1/C\omega-L\omega)^2+R^2$,
which attains it minimum value $R^2$ when
$\omega=\omega_0=\dst\frac{1}{\sqrt{LC}}$. The maximum amplitude of
$I_p$ is
$\sigma(\omega)=\dfrac{\sqrt{U^2+V^2}}{ R}$.
\end{sectionSolutions}

\section{Motion Under a Central Force}

\begin{sectionSolutions}
\item
Let $h=r_0^2\theta_0'$; then $\rho=\dst\frac{h^2}{ k}$. Since
$r=\dst\frac{\rho}{ 1+e\cos(\theta-\phi)}$, it follows that (A)
$e\cos(\theta-\phi)=\dst\frac{\rho}{ r}-1$. Differentiating this with
respect to $t$ yields $-e\sin(\theta-\phi)\theta'=-\dst\frac{\rho r'}{
r^2}$, so (B) $e\sin(\theta-\phi)=\dst\frac{\rho r'}{ h}$, since
$r^2\theta'\equiv h$.
Squaring and adding (A) and (B) and setting $t=0$ in the result yields
$e=\dst{\left[\left(\frac{\rho}{ r_0}-1\right)^2+\left(\frac{\rho r_0'}{
h}\right)^2\right]^{1/2}}$. If $e=0$, then $\theta_0$ is undefined, but
also irrelevant; if $e\ne0$, then set $t=0$ in (A) and (B) to see that
$\phi=\theta_0-\alpha$, where $-\pi\le\alpha<\pi$,
$\cos\alpha=\dst{\frac{1}{ e}\left(\frac{\rho}{ r_0}-1\right)}$ and
$\sin\alpha=\dst\frac{\rho r_0'}{ eh}$


\item
Recall that
(A) $\dst{\frac{d^{2}u}d{\theta^{2}}}=-\dst{\frac{1}{mh^{2}u^{2}}f(1/u)}$.
Let $u=\dst\frac{1}{ r}=\dst\frac{1}{ c\theta^2}$; then $\dst\frac{d^2u}{
d\theta^2}=\dst\frac{6}{ c\theta^4}= 6cu^2$.
 $6cu^2+u=-\dst\frac{1}{ mh^2u^2}f(1/u)$, so
$f(1/u)=-mh^2(6cu^4+u^3)$ and $f(r)=-mh^2\dst\left(\frac{6c}{
r^4}+\frac{1}{ r^3}\right)$.



\item
\begin{alist}
\item With $f(r)=-\dst\frac{mk}{ r^3}$, Eqn.~6.4.11 becomes
(A) $\dst\frac{d^2u}{ d\theta^2}+\left(1-\frac{k}{ h^2}\right)u=0$. The
initial conditions imply that $u(\theta_0)=\dst\frac{1}{ r_0}$ and
$\dst\frac{du}{ d\theta}(\theta_0)=-\dst\frac{r_0'}{ h}$ (see
Eqn.~(6.4.)

\item Let $\gamma=\left|1-\dst\frac{k}{ h^2}\right|^{1/2}$. (i) If
$h^2<k$, then (A) becomes $\dst\frac{d^2u}{ d\theta^2}-\gamma^2u=0$, and
the solution of the initial value problem for $u$ is $u=\dst\frac{1}{
r_0}\cosh\gamma(\theta-\theta_0)-\dst\frac{r_0'}{\gamma
h}\sinh\gamma(\theta-\theta_0)$; therefore $r=r_0\left(
\cosh\gamma(\theta-\theta_0)-\dst\frac{r_0r_0'}{\gamma
h}\sinh\gamma(\theta-\theta_0)\right)^{-1}$. (ii) If $h^2=k$, then (A)
becomes $\dst\frac{d^2u}{ d\theta^2}=0$, and the solution of the initial
value problem for $u$ is $u=\dst\frac{1}{ r_0}-\dst\frac{r_0'}{
h}(\theta-\theta_0)$; therefore $r=r_0\left(1 -\dst\frac{r_0r_0'}{
h}(\theta-\theta_0)\right)^{-1}$. (iii) If $h^2>k$, then (A) becomes
$\dst\frac{d^2u}{ d\theta^2}+\gamma^2u=0$, and the solution of the
initial value problem for $u$ is $u=\dst\frac{1}{
r_0}\cos\gamma(\theta-\theta_0)-\dst\frac{r_0'}{\gamma
h}\sin\gamma(\theta-\theta_0)$; therefore $r=r_0\left(
\cos\gamma(\theta-\theta_0)-\dst\frac{r_0r_0'}{\gamma
h}\sin\gamma(\theta-\theta_0)\right)^{-1}$.
\end{alist}
\end{sectionSolutions}

\chapter{Series Solutions of Linear Second Equations}

\section{Review of Power Series}
\begin{sectionSolutions}
\item
From Theorem~7.1.3, $\sum_{m=0}^\infty b_mz^m$ converges
if $|z|<1/L$ and diverges if $|z|>1/L$.
 Therefore,
$\sum_{m=0}^\infty b_m(x-x_0)^2$ converges if $|x-x_0|<1/\sqrt{L}$
and diverges if $|x-x_0|>1/\sqrt{L}$.


\item
From Theorem~7.1.3, $\sum_{m=0}^\infty b_mz^m$ converges
if $|z|<1/L$ and diverges if $|z|>1/L$.
 Therefore,
$\sum_{m=0}^\infty b_m(x-x_0)^{km}$ converges if
$|x-x_0|<1/\sqrt[k]{L}$ and diverges if $|x-x_0|>1/\sqrt[k]{L}$.


\addtocounter{enumi}{6}
\item%7.1.12
$\dst(1+3x^2)y''+3x^2y'-2y=
\sum_{n=2}^\infty n(n-1)a_nx^{n-2}
+3\sum_{n=2}^\infty n(n-1)a_nx^n
+3\sum_{n=1}^\infty na_nx^{n+1}-2\sum_{n=0}^\infty a_nx^n
=\sum_{n=0}^\infty (n+2)(n+1)a_{n+2}x^n
+3\sum_{n=1}^\infty n(n-1)na_nx^n
+3\sum_{n=1}^\infty (n-1)a_{n-1}x^n-2\sum_{n=0}^\infty a_nx^n
=2a_2-2a_0+\sum_{n=1}^\infty
[(n+2)(n+1)a_{n+2}+(3n(n-1)-2)a_n+3(n-1)a_{n-1}]x^n$.

\addtocounter{enumi}{-1}
\item%7.1.13
$\dst (1+2x^2)y''+(2-3x)y'+4y=
\sum_{n=2}^\infty n(n-1)a_nx^{n-2}
+2\sum_{n=2}^\infty n(n-1)a_nx^n
+2\sum_{n=1}^\infty na_nx^{n-1}-2\sum_{n=0}^\infty a_nx^n
+4\sum_{n=0}^\infty a_nx^n
=\sum_{n=0}^\infty (n+2)(n+1)a_{n+2}x^n
+2\sum_{n=0}^\infty n(n-1)a_nx^n
+2\sum_{n=0}^\infty (n+1)a_{n+1}x^n-3\sum_{n=0}^\infty a_nx^n
+4\sum_{n=0}^\infty a_nx^n
\sum_{n=0}^\infty
\left[(n+2)(n+1)a_{n+2}+2(n+1)a_{n+1}+(2n^2-5n+4)a_n\right]x^n$.

\addtocounter{enumi}{-1}
\item%7.1.14
\begin{align*}
(1+x^2)y''+(2-x)y'+3y
&=\sum_{n=2}^\infty n(n-1)a_nx^{n-2}+\sum_{n=2}^\infty n(n-1)a_nx^n\\
&\qquad+2\sum_{n=1}^\infty na_nx^{n-1}-\sum_{n=1}^\infty na_nx^n
+3\sum_{n=0}^\infty a_nx^n\\
&=\sum_{n=0}^\infty (n+2)(n+1)a_{n+2}x^n+\sum_{n=0}^\infty n(n-1)a_nx^n\\
&\qquad+2\sum_{n=0}^\infty (n+1)a_{n+1}x^n-\sum_{n=0}^\infty na_nx^n
+3\sum_{n=0}^\infty a_nx^n\\
&=\sum_{n=0}^\infty
\left[(n+2)(n+1)a_{n+2}+2(n+1)a_{n+1}+(n^2-2n+3)a_n\right]x^n.
\end{align*}

\item
Let $t=x+1$; then $\dst xy''+(4+2x)y'+(2+x)y=(-1+t)y''+(2+2t)y'+(1+t)y
=-\sum_{n=2}^\infty n(n-1)a_nt^{n-2}
+\sum_{n=2}^\infty n(n-1)a_nt^{n-1}
+2\sum_{n=1}^\infty n a_nt^{n-1}
+2\sum_{n=1}^\infty n a_nt^n
+\sum_{n=0}^\infty  a_nt^n
+\sum_{n=0}^\infty  a_nt^{n+1}
=-\sum_{n=0}^\infty (n+2)(n+1)a_{n+2}t^n
+\sum_{n=0}^\infty (n+1)na_{n+1}t^n
+2\sum_{n=0}^\infty (n+1) a_{n+1}t^n
+2\sum_{n=0}^\infty n a_nt^n
+\sum_{n=0}^\infty  a_nt^n
+\sum_{n=1}^\infty  a_{n-1}t^n
=(-2a_2+2a_1+a_0)
+\sum_{n=1}^\infty
 \left[-(n+2)(n+1)a_{n+2}+(n+1)(n+2)a_{n+1}+(2n+1)a_n+a_{n-1}\right]
(x+2)^n$.

\addtocounter{enumi}{2}
\item%7.1.20
$y'(x)=\dst x^r\sum_{n=0}^\infty na_nx^{n-1}+rx^{r-1}\sum_{n=0}^\infty
a_nx^n=\sum_{n=0}^\infty (n+r)x^{n+r-1}$

$y''=\dst\frac{d}{ dx}y'(x)=\frac{d}{ dx}\left[x^{r-1}\sum_{n=0}^\infty
(n+r)a_nx^n\right]=x^{r-1}\sum_{n=0}^\infty (n+r)na_nx^{n-1}+
(r-1)x^{r-2}\sum_{n=0}^\infty (n+r)a_nx^n=\sum_{n=0}^\infty
(n+r)(n+r-1)a_nx^{n+r-2}$.



\item
$x^2(1+x)y''+x(1+2x)y'-(4+6x)y=
(x^2y''+xy'-4y)+x(x^2y''+2xy'-6y)=\dst
\sum_{n=0}^\infty [(n+r)(n+r-1)+(n+r)-4]a_nx^{n+r}
+\sum_{n=0}^\infty [(n+r)(n+r-1)+2(n+r)-6]a_nx^{n+r+1}
=\sum_{n=0}^\infty (n+r-2)(n+r+2)a_nx^{n+r}
+\sum_{n=0}^\infty (n+r+3)(n+r-2)a_nx^{n+r+1}
=\sum_{n=0}^\infty (n+r-2)(n+r+2)a_nx^{n+r}
+\sum_{n=1}^\infty (n+r+2)(n+r-3)a_{n-1}x^{n+r}
=x^r\sum_{n=0}^\infty  b_nx^n$ with
$b_0=(r-2)(r+2)a_0$ and
$b_n=(n+r-2)(n+r+2)a_n+(n+r+2)(n+r-3)a_{n-1}$, $n\ge1$.


\item
$x^2(1+3x)y''+x(2+12x+x^2)y'+2x(3+x)y=
(x^2y''+2xy')+x(3x^2y''+12xy'+6y)+x^2(xy'+2y)=\dst
\sum_{n=0}^\infty [(n+r)(n+r-1)+2(n+r)]a_nx^{n+r}
+\sum_{n=0}^\infty [3(n+r)(n+r-1)+12(n+r)+6]a_nx^{n+r+1}
+\sum_{n=0}^\infty[(n+r)+2]a_nx^{n+r+2 }
=\sum_{n=0}^\infty (n+r)(n+r+1)a_nx^{n+r}
+3\sum_{n=0}^\infty (n+r+1)(n+r+2)a_nx^{n+r+1}
+\sum_{n=0}^\infty (n+r+2)a_nx^{n+r+2}
=\sum_{n=0}^\infty (n+r)(n+r+1)a_nx^{n+r}
+3\sum_{n=1}^\infty (n+r)(n+r+1)a_{n-1}x^{n+r}
+\sum_{n=2}^\infty (n+r)a_{n-2}x^{n+r}
=x^r\sum_{n=0}^\infty  b_nx^n$ with
$b_0=r(r+1)a_0$,
$b_1=(r+1)(r+2)a_1+3(r+1)(r+2)a_0$,
$b_n=(n+r)(n+r+1)a_n+3(n+r)(n+r+1)a_{n-1}+(n+r)a_{n-2}$,
$n\ge2$.



\item
$x^2(2+x^2)y''+2x(5+x^2)y'+2(3-x^2)y=
(2x^2y''+10xy'+6y)+x^2(x^2y''+2xy'-2y)=\dst
\sum_{n=0}^\infty [2(n+r)(n+r-1)+10(n+r)+6]a_nx^{n+r}
+\sum_{n=0}^\infty [(n+r)(n+r-1)+2(n+r)-2]a_nx^{n+r+2}
=2\sum_{n=0}^\infty (n+r+1)(n+r+3)a_nx^{n+r}
+\sum_{n=0}^\infty (n+r-1)(n+r+2)a_nx^{n+r+2}
=2\sum_{n=0}^\infty (n+r+1)(n+r+3)a_nx^{n+r}
+\sum_{n=2}^\infty (n+r-3)(n+r)a_{n-2}x^{n+r}
=x^r\sum_{n=0}^\infty  b_nx^n$ with
$b_0=2(r+1)(r+3)a_0$,
$b_1=2(r+2)(r+4)a_1$
$b_n=2(n+r+1)(n+r+3)a_n+(n+r-3)(n+r)a_{n-2}$,
$n\ge2$.
\end{sectionSolutions}

\section{Series Solutions Near an Ordinary Point I}

\begin{sectionSolutions}
\item
$p(n)=n(n-1)+2n-2=(n+2)(n-1)$;
$a_{n+2}=-\frac{n-1}{ n+1}a_n$;
$a_{2m+2}=-\dst\frac{2m-1}{2m+1}a_{2m}$, so
$a_{2m}=\dfrac{(-1)^m}{2m-1}a_0$;
$a_{2m+3}=\dst-\frac{m}{ m+2}a_{2m+1}=0$
if  $m\ge0$;
$y=\dst{a_0\sum_{m=0}^\infty(-1)^{m+1}\frac{x^{2m}}{2m-1}+a_1x}$.

\item
$p(n)=-n(n-1)-8n-12=-(n+3)(n+4)$;
$a_{n+2}=-\dfrac{(n+3)(n+4)}{(n+2)(n+1)}a_n$;
$a_{2m+2}=-\dfrac{(m+2)(2m+3)}{(m+1)(2m+1)}a_{2m}$, so
$a_{2m}=(m+1)(2m+1)a_0$;
$a_{2m+3}=\dfrac{(m+2)(2m+5)}{(m+1)(2m+3)}a_{2m+1}$
 so $a_{2m+1}=\dfrac{(m+1)(2m+3)}{3}a_1$;
$y=\dst{a_0\sum_{m=0}^\infty(m+1)(2m+1)x^{2m}+\frac{a_1}{3}\sum_{m=0}^\infty
(m+1)(2m+3)x^{2m+1}}$.


\item
$p(n)=n(n-1)+2n+\dst\frac{1}{4}=\frac{(2n+1)^2}{4}$;
$a_{n+2}=-\dfrac{(2n+1)^2}{4(n+2)(n+1)}a_n$;
$a_{2m+2}=-\dfrac{(4m+1)^2}{8(m+1)(2m+1)}a_{2m}$, so
$a_{2m}=\dst
(-1)^m\left[\prod_{j=0}^{m-1}\frac{(4j+1)^2}{2j+1}\right]\frac{1}{8^mm!}a_0$;
$a_{2m+3}=-\dfrac{(4m+3)^2}{8(2m+3)(m+1)}a_{2m+1}$
 so $a_{2m+1}=\dst
(-1)^m\left[\prod_{j=0}^{m-1}\frac{(4j+3)^2}{2j+3}\right]
\frac{1}{8^mm!}a_1$;
$y=\dst{a_0\sum_{m=0}^\infty
(-1)^m\left[\prod_{j=0}^{m-1}\frac{(4j+1)^2}{2j+1}\right]\frac{x^{2m}}{8^mm!}
+a_1\sum_{m=0}^\infty
(-1)^m\left[\prod_{j=0}^{m-1}\frac{(4j+3)^2}{2j+3}\right]
\frac{x^{2m+1}}{8^mm!}}$




\item
$p(n)=n(n-1)-10n+28=(n-7)(n-4)$;
$a_{n+2}=-\dfrac{(n-7)(n-4)}{(n+2)
(n+1)}a_n$;
$a_{2m+2}=-\dfrac{2(2m-7)(m-2)}{2(m+1)(2m+1)}a_{2m}$, so
$a_2=-14a_0$, $a_4=-\dst\frac{5}{6}a_2=\frac{35}{3}a_0$,
$a_{2m}=0$ if $m\ge3$;
$a_{2m+3}=-\dfrac{(m-3)(2m-3)}{(2m+3)(m+1)}a_{2m+1}$,
so $a_3=-3a_1$, $a_5=-\dst\frac{1}{5}a_3=\frac{3}{5}a_1$,
$a_7=\dst\frac{1}{21}a_5=\frac{1}{35}a_1$;\\
$y=\dst{a_0\left(1-14x^2+\frac{35}{3}x^4\right)+a_1\left(x-3x^3+\frac{3}{5}x^5
+\frac{1}{35}x^7\right)}$.



\item
$p(n)=2n+3$;
$a_{n+2}=\dst-\frac{2n+3}{(n+2)(n+1)}a_n$;
$a_{2m+2}=-\dst\frac{4m+3}{2(m+1)(2m+1)}a_{2m}$, so
$a_{2m}=\dst\left[\prod_{j=0}^{m-1}\frac{4j+3}{2j+1}\right]
\frac{(-1)^m}{2^mm!}a_0$;
$a_{2m+3}=-\dst\frac{4m+5}{2(2m+3)(m+1)}a_{2m+1}$
 so $a_{2m+1}=\dst
\left[\prod_{j=0}^{m-1}\frac{4j+5}{2j+3}\right]\frac{(-1)^m}{
2^mm!}a_1$;
$y=\dst{a_0\sum_{m=0}^\infty
(-1)^m\left[\prod_{j=0}^{m-1}\frac{4j+3}{2j+1}\right]\frac{x^{2m}}{2^mm!}
+a_1\sum_{m=0}^\infty
(-1)^m\left[\prod_{j=0}^{m-1}\frac{4j+5}{2j+3}\right]\frac{x^{2m+1}}{
2^mm!}}$


\item
$p(n)=2n(n-1)-9n-6=(n-6)(2n+1)$;
$a_{n+2}=\dst-\frac{(n-6)(2n+1)}{(n+2)(n+1)}a_n$;
$a_0=y(0)=1$; $a_1=y'(0)=-1$.


\addtocounter{enumi}{-1}
\item%7.2.13
$p(n)=8n(n-1)+2=2(2n-1)^2$;
$a_{n+2}=\dst-\frac{2(2n-1)^2}{(n+2)(n+1)}a_n$;
$a_0=y(0)=2$; $a_1=y'(0)=-1$.


\addtocounter{enumi}{1}
\item%7.2.16
$p(n)=-1$;
$a_{n+2}=\dst\frac{1}{(n+2)(n+1)}a_n$;
$a_{2m+2}=\dst\frac{1}{(2m+2)(2m+1)}a_{2m}$, so
$a_{2m}=\dst\frac{1}{(2m)!}a_0$;
$a_{2m+3}=\dst\frac{1}{(2m+3)(2m+1)}a_{2m+1}$,
so $a_{2m+1}=\dst\frac{1}{(2m+1)!}a_1$;
$y=\dst{a_0\sum_{m=0}^\infty\frac{(x-3)^{2m}}{(2m)!}+a_1\sum_{m=0}^\infty
\frac{(x-3)^{2m+1}}{(2m+1)!}}$.



\item
Let $t=x-1$;  then $(1-2t^2)y''-10ty'-6y=0$;
$p(n)=-2n(n-1)-10n-6=-2(n+1)(n+3)$;
$a_{n+2}=\dst\frac{2(n+3)}{ n+2}a_n$;
$a_{2m+2}=\dst\frac{2m+3}{ m+1}a_{2m}$, so
$a_{2m}=\dst\frac{1}{ m!}\left[\prod_{j=0}^{m-1}(2j+3)\right]a_0$;
$a_{2m+3}=\dst\frac{4(m+2)}{ 2m+3}a_{2m+1}$,
so $a_{2m+1}=\dst\frac{4^m(m+1)!}{\prod_{j=0}^{m-1}(2j+3)}a_1$;
$y=\dst{a_0\sum_{m=0}^\infty \left[\prod_{j=0}^{m-1}(2j+3)\right]\frac{
(x-1)^{2m}}{ m!}}+\dst{a_1\sum_{m=0}^\infty
\frac{4^m(m+1)!}{\prod_{j=0}^{m-1}(2j+3)}(x-1)^{2m+1}}$.



\item
Let $t=x+1$;  then
$\dst\left(1+\frac{3t^2}{3}\right)y''-\frac{9t}{2}y'+\frac{3}{2}y=0$;
$p(n)=\dst\frac{3}{2}n(n-1)+\frac{9}{2}n+\frac{3}{2}=\frac{3}{2}(n+1)^2$;
$a_{n+2}=-\dst\frac{3(n+1)}{2(n+2)}a_n$;
$a_{2m+2}=-\dst\frac{3(2m+1)}{4(m+1)}a_{2m}$, so
$a_{2m}=(-1)^m\dst\left[\prod_{j=0}^{m-1}(2j+1)\right]\frac{3^m}{4^mm!}a_0$;
$a_{2m+3}=-\dst\frac{3(m+1)}{2m+3}a_{2m+1}$,
so $a_{2m+1}=\dst(-1)^m\frac{3^mm!}{\prod_{j=0}^{m-1}(2j+3)}a_1$;\\
$y=\dst{a_0\sum_{m=0}^\infty(-1)^m\left[\prod_{j=0}^{m-1}(2j+1)\right]
\frac{3^m}{4^mm!}(x+1)^{2m}+a_1\sum_{m=0}^\infty(-1)^m\frac{3^mm!}{
\prod_{j=0}^{m-1}(2j+3)}(x+1)^{2m+1}}$.


\item
$p(n)=n+3$;
$a_{n+2}=-\dst\frac{n+3}{(n+2)(n+1)}a_n$;
$a_0=y(3)=-2$; $a_1=y'(3)=3$.



\item
Let $t=x-3$; $(1+4t^2)y''+y=0$;
$p(n)=(4n(n-1)+1=2n-1)^2$;
$a_{n+2}=-\dst\frac{(2n-1)^2}{(n+2)(n+1)}a_n$;
$a_0=y(3)=4$; $a_1=y'(3)=-6$.





\item
Let $t=x+1$;
$\dst\left(1+\frac{2t^2}{3}\right)y''-\frac{20}{3}ty'+20y=0$;
$p(n)=\dst\frac{2}{3}n(n-1)-\frac{20}{3}n+20=\frac{2(n-6)(n-5)}{3}$;
$a_{n+2}=-\dst\frac{2(n-6)(n-5)}{3(n+2)(n+1)}a_n$;
$a_0=y(-1)=3$; $a_1=y'(-1)=-3$.


\item
From Theorem~7.2.2,
$a_{n+2}=\dst-\frac{p(n) }{(n+2)(n+1)}a_n$;
$a_{2m+2}=\dst-\frac{p(2m)}{(2m+2)(2m+1)}a_{2m}$, so $a_{2m}=\dst
\left[\prod^{m-1}_{j=0}p(2j)\right] \frac{(-1)^m}{(2m)!}a_0$;
$a_{2m+3}=\dst-\frac{p(2m+1)}{(2m+3)(2m+2)}a_{2m}$, so $a_{2m+1}=\dst
\left[\prod^{m-1}_{j=0}p(2j+1)\right]\frac{(-1)^m}{(2m+1)!}a_1$.




\item
\begin{alist}
\item
Here
$p(n)=-[n(n-1)+2bn-\alpha(\alpha+2b-1)]=-(n-\alpha)(n+\alpha+2b-1)$,
so Exercise~7.2.28 implies that $y_1$ and $y_2$ have the stated
forms. If $\alpha=2k$, then
\[
y_1=\sum_{m=0}^\infty \left[\prod_{j=0}^{m-1}(2j-2k)(2j+2k+2b-1)
\right]\frac{x^{2m}}{(2m)!}
\tag{C}.
\]
If $\alpha=2k+1$, then
\[
y_2=\sum_{m=0}^\infty \left[\prod_{j=0}^{m-1}(2j-2k)(2j+2k+2b)
\right]\frac{x^{2m+1}}{(2m+1)!}.
\tag{D}
\]
Since $2b$ is not a negative integer and $\prod_{j=0}^{m-1}(2j-2k)=0$
if $m>k$, $y_1$ in (C) and $y_2$ in (D) have the stated properties.
This implies the conclusions regarding $P_n$.

\item Multiplying (A) through by $(1-x^2)^{b-1}$ yields
\[
[(1-x^2)^b P_n']'=-n(n+2b-1)(1-x^2)^{b-1}P_n.
\tag{E}
\]

\item
Therefore,
\[
[(1-x^2)^b P_m']'=-m(m+2b-1)(1-x^2)^{b-1}P_m.
\tag{F}
\]
Subtract  $P_n$ times (F) from $P_m$ times (E) to obtain (B).

\item Integrating the left side of (B) by parts over $[-1,1]$
yields zero, which implies the conclusion.
\end{alist}

\item
\begin{alist}
\item
Let $Ly=(1+\alpha x^3)y''+\beta x^2y'+\gamma xy$. If
$y=\dst\sum_{n=0}^\infty a_nx^n$, then
$Ly=\dst\sum_{n=2}^\infty n(n-1)a_nx^{n-2}+\sum_{n=0}^\infty
p(n)a_nx^{n+1}=
2a_2+\sum_{n=0}^\infty [(n+3)(n+2)a_{n+3}+p(n)a_n]x^{n+1}
= 0$ if and only if $a_2=0$ and $\dst
a_{n+3}=-\frac{p(n)}{(n+3)(n+2)}a_n$ for $n\ge0$.
\end{alist}

\item $p(r)=-2r(r-1)-10r-8=-2(r+2)^2$;
(A) $\dst\prod_{j=0}^{m-1}\frac{p(3j)}{3j+2}=
\prod_{j=0}^{m-1}\frac{(-2)(3j+2)^2}{3j+2}
=(-1)^m2^m\prod_{j=0}^{m-1}(3j+2)$;
(B) $\dst\prod_{j=0}^{m-1}\frac{p(3j+1)}{3j+4}=
\prod_{j=0}^{m-1}\frac{(-2)(3j+3)^2}{3j+4}
=\frac{(-1)^m2^m(m!)^2}{\prod_{j=0}^{m-1}(3j+4)}$.
Substituting (A) and (B) into the
result of Exercise~7.2.32\part{c} yields

$y=\dst{a_0\sum_{m=0}^\infty
\left(\frac{2}{3}\right)^m\left[\prod_{j=0}^{m-1}(3j+2)\right]\frac{x^{3m}}{ m!}
+a_1\sum_{m=0}^\infty\frac{6^mm!}{\prod_{j=0}^{m-1}(3j+4)}x^{3m+1}}$.


\item $p(r)=-2r(r-1)+6r+24=-2(r-6)(r+2)$;
(A) $\dst\prod_{j=0}^{m-1}\frac{p(3j)}{3j+2}=\prod_{j=0}^{m-1}
(-6)(j-2)$.


(B) $\dst\prod_{j=0}^{m-1}\frac{p(3j+1)}{3j+4}=
\prod_{j=0}^{m-1}\frac{(-6)(j+1)(3j-5)}{3j+4}=
(-1)^m6^mm\prod_{j=0}^{m-1}\frac{3j-5}{3j+4}$.
Substituting (A) and (B) into the
result of Exercise~7.2.32\part{c} yields

$y=\dst{a_0(1-4x^3+4x^6)+a_1\sum_{m=0}^\infty
2^m\left[\prod_{j=0}^{m-1}\frac{3j-5}{3j+4}\right]x^{3m+1}}$.


\item
\begin{alist}
\item
Let $Ly=(1+\alpha x^{k+2})y''+\beta x^{k+1}y'+\gamma x^ky$. If
$y=\dst\sum_{n=0}^\infty a_nx^n$, then
$Ly=\dst\sum_{n=2}^\infty n(n-1)a_nx^{n-2}+\sum_{n=0}^\infty
p(n)a_nx^{n+k}=
\sum_{n=-k}^{-1}(n+k+2)(n+k-1)a_{n+k+2}x^{n+k}+\sum_{n=0}^\infty
[(n+k+2)(n+k+1)a_{n+k+2}+p(n)a_n]x^{n+k}
= 0$ if and only if $a_k=0$ for $2\le n\le k+1$ and (A) $\dst
a_{n+k+1}=-\frac{p(n)}{(n+k+2)(n+k+1)}a_n$ for $n\ge0$.

\item If $a_n=0$ the $a_{n+(k+2)m}=0$ for all $m\ge0$, from
(A).
\end{alist}


\item $k=2$ and
$p(r)=1$;
(A) $\dst\prod_{j=0}^{m-1}\frac{p(4j)}{4j+3}=
\frac{1}{\prod_{j=0}^{m-1}(4j+3)}$;
(B) $\dst\prod_{j=0}^{m-1}\frac{p(4j+1)}{(4j+5)}=
\frac{1}{\prod_{j=0}^{m-1}4j+5}$.
Substituting (A) and (B) into the
result of Exercise~t.2.38\part{c} yields

$y=\dst{a_0\sum_{m=0}^\infty
(-1)^m\frac{x^{4m}}{4^mm!\prod_{j=0}^{m-1}(4j+3)}+a_1\sum_{m=0}^\infty
(-1)^m\frac{x^{4m+1}}{4^mm!\prod_{j=0}^{m-1}(4j+5)}}$.



\item $k=6$ and
$p(r)=r(r-1)-16r+72=(r-9)(r-8) $;
(A) $\dst\prod_{j=0}^{m-1}\frac{p(8j)}{8j+7}=
\prod_{j=0}^{m-1}\frac{8(j-1)(8j-9)}{8j+7}$;


(B) $\dst\prod_{j=0}^{m-1}\frac{p(8j+1)}{(8j+9)}=
\prod_{j=0}^{m-1}\frac{8(j-1)(8j-7)}{8j+9}$;

Substituting (A) and (B) into the
result of Exercise~7.2.38\part{c} yields


$y=\dst{a_0\left(1-\frac{9}{7}x^8\right)+a_1\left(x-\frac{7}{9}x^9\right)}$.



\item $k=4$ and
$p(r)=r+6$;
(A) $\dst\prod_{j=0}^{m-1}\frac{p(6j)}{6j+5}=
\prod_{j=0}^{m-1}\frac{6(j+1)}{6j+5}=\frac{6^mm!}{\prod_{j=0}^{m-1}
(6j+5)}$;


(B) $\dst\prod_{j=0}^{m-1}\frac{p(6j+1)}{(6j+7)}=1$;

Substituting (A) and (B) into the
result of Exercise~7.2.38\part{c} yields


$y=\dst{a_0\sum_{m=0}^\infty(-1)^m\frac{x^{6m}}{\prod_{j=0}^{m-1}(6j+5)}+
a_1\sum_{m=0}^\infty(-1)^m\frac{x^{6m+1}}{6^mm!}}$.

\end{sectionSolutions}

\section{Series Solutions Near an Ordinary Point II}

\begin{sectionSolutions}
\item
If $y=\dst\sum_{n=0}^\infty a_nx^n$, then
$(1+x+2x^2)y''+(2+8x)y'+4y=
\dst \sum_{n=2}^\infty n(n-1)a_nx^{n-2}
+\sum_{n=2}^\infty n(n-1)a_nx^{n-1}
+2\sum_{n=2}^\infty n(n-1)a_nx^n
+2\sum_{n=1}^\infty na_nx^{n-1}
+8\sum_{n=1}^\infty na_nx^n
+4\sum_{n=0}^\infty a_nx^n
=\sum_{n=0}^\infty(n+2)(n+1)(a_{n+2}+a_{n+1}+2a_n)x^n=0$
if
$a_{n+2}=-a_{n+1}-2a_n$, $a_n\ge0$.
Starting with $a_0=-1$ and $a_1=2$ yields
$y=\dst{-1+2x-4x^3+4x^4+4x^5-12x^6+4x^7+\cdots}$.

\item
If $y=\dst\sum_{n=0}^\infty a_nx^n$, then
$(1+x+3x^2)y''+(2+15x)y'+12y=
\dst \sum_{n=2}^\infty n(n-1)a_nx^{n-2}
+\sum_{n=2}^\infty n(n-1)a_nx^{n-1}
+3\sum_{n=2}^\infty n(n-1)a_nx^n
+2\sum_{n=1}^\infty na_nx^{n-1}
+15\sum_{n=1}^\infty na_nx^n
+12\sum_{n=0}^\infty a_nx^n
=\sum_{n=0}^\infty[(n+2)(n+1)a_{n+2}+(n+1)(n+2)a_{n+1}+3(n+2)^2a_n]x^n=0$
if
$a_{n+2}=-a_{n+1}-\dst\frac{3(n+2)}{ n+1}a_n$,
$a_n\ge0$. Starting with $a_0=0$ and $a_1=1$ yields
$y=\dst{x-x^2-\frac{7}{2}x^3+\frac{15}{2}x^4+\frac{45}{8}x^5
-\frac{261}{8}x^6+\frac{207}{16}x^7+\cdots}$.


\item
If $y=\dst\sum_{n=0}^\infty a_nx^n$, then
$(3+3x+x^2)y''+(6+4x)y'+2y=
3\dst \sum_{n=2}^\infty n(n-1)a_nx^{n-2}
+3\sum_{n=2}^\infty n(n-1)a_nx^{n-1}
+\sum_{n=2}^\infty n(n-1)a_nx^n
+6\sum_{n=1}^\infty na_nx^{n-1}
+4\sum_{n=1}^\infty na_nx^n
+2\sum_{n=0}^\infty a_nx^n
=\sum_{n=0}^\infty(n+2)(n+1)[3a_{n+2}+3a_{n+1}+a_n]x^n=0$
if
$a_{n+2}=-a_{n+1}-a_n/3$,
$a_n\ge0$. Starting with $a_0=7$ and $a_1=3$ yields
$y=\dst{7+3x-\frac{16}{3}x^2+\frac{13}{3}x^3-\frac{23}{9}x^4+\frac{10}{9}x^5
-\frac{7}{27}x^6-\frac{1}{9}x^7+\cdots}$.




\item
 The equation is equivalent to
 $(1+t+2t^2)y''+(2+6t)y'+2y=0$ with $t=x-1$.
If $y=\dst\sum_{n=0}^\infty a_nt^n$, then
$(1+t+2t^2)y''+(2+6t)y'+2y=
\dst \sum_{n=2}^\infty n(n-1)a_nt^{n-2}
+\sum_{n=2}^\infty n(n-1)a_nt^{n-1}
+2\sum_{n=2}^\infty n(n-1)a_nt^n
+2\sum_{n=1}^\infty na_nt^{n-1}
+6\sum_{n=1}^\infty na_nt^n
+2\sum_{n=0}^\infty a_nt^n
=\sum_{n=0}^\infty(n+1)[(n+2)a_{n+2}+(n+2)a_{n+1}+2(n+1)a_n]t^n=0$
if
$a_{n+2}=-\dst a_{n+1}-\frac{2(n+1)}{ n+2}a_n$,
$a_n\ge0$. Starting with $a_0=1$ and $a_1=-1$ yields
$y=\dst{1-(x-1)+\frac{4}{3}(x-1)^3-\frac{4}{3}(x-1)^4-\frac{4}{5}(x-1)^5
+\frac{136}{45}(x-1)^6-\frac{104}{63}(x-1)^7+\cdots}$.



\item
 The equation is equivalent to
 $(1+t+t^2)y''+(3+4t)y'+2y=0$ with $t=x-1$.
If $y=\dst\sum_{n=0}^\infty a_nt^n$, then
$(1+t+t^2)y''+(3+4t)y'+2y=
\dst \sum_{n=2}^\infty n(n-1)a_nt^{n-2}
+\sum_{n=2}^\infty n(n-1)a_nt^{n-1}
+\sum_{n=2}^\infty n(n-1)a_nt^n
+3\sum_{n=1}^\infty na_nt^{n-1}
+4\sum_{n=1}^\infty na_nt^n
+2\sum_{n=0}^\infty a_nt^n
=\sum_{n=0}^\infty(n+1)[(n+2)a_{n+2}+(n+3)a_{n+1}+(n+2)a_n]t^n=0$
if
$a_{n+2}=-\dst\frac{n+3}{ n+2}a_{n+1}-a_n$,
$a_n\ge0$. Starting with $a_0=2$ and $a_1=-1$ yields
$y=$
\[
2-(x-1)-\frac{1}{2}(x-1)^2+\frac{5}{3}(x-1)^3-\frac{19}{12}(x-1)^4
+\frac{7}{30}(x-1)^5+\frac{59}{45}(x-1)^6-\frac{1091}{630}(x-1)^7+\cdots
\]



\item
 The equation is equivalent to
 $(1+2t+t^2)y''+(1+7t)y'+8y=0$ with $t=x-1$.
If $y=\dst\sum_{n=0}^\infty a_nt^n$, then
$(1+2t+t^2)y''+(1+7t)y'+8y=
\dst \sum_{n=2}^\infty n(n-1)a_nt^{n-2}
+2\sum_{n=2}^\infty n(n-1)a_nt^{n-1}
+\sum_{n=2}^\infty n(n-1)a_nt^n
+\sum_{n=1}^\infty na_nt^{n-1}
+7\sum_{n=1}^\infty na_nt^n
+8\sum_{n=0}^\infty a_nt^n
=\sum_{n=0}^\infty[(n+2)(n+1)a_{n+2}+(n+1)(2n+1)a_{n+1}+(n+2)(n+4)a_n]t^n=0$
if
$a_{n+2}=-\dst\frac{2n+1}{ n+2}a_{n+1}-\frac{n+4}{ n+1}a_n$,
$a_n\ge0$. Starting with $a_0=1$ and $a_1=-2$ yields
$y=\dst{1-2(x-1)-3(x-1)^2+8(x-1)^3-4(x-1)^4-\frac{42}{5}(x-1)^5
+19(x-1)^6-\frac{604}{35}(x-1)^7+\cdots}$


\addtocounter{enumi}{2}
\item%7.3.16
If $y=\dst\sum_{n=0}^\infty a_nx^n$, then
$(1-x)y''-(2-x)y'+y=
\dst \sum_{n=2}^\infty n(n-1)a_nx^{n-2}
-\sum_{n=2}^\infty n(n-1)a_nx^{n-1}
-2\sum_{n=1}^\infty na_nx^{n-1}
+\sum_{n=1}^\infty na_nx^n
+\sum_{n=0}^\infty a_nx^n
=\sum_{n=0}^\infty[(n+2)(n+1)a_{n+2}-(n+2)(n+1)a_{n+1}+(n+1)a_n]x^n=0$
if
$a_{n+2}=a_{n+1}-\dst\frac{a_n}{ n+2}$,$a_n\ge0$.


\item
If $y=\dst\sum_{n=0}^\infty a_nx^n$, then
$(1+x^2)y''+y'+2y=
\dst \sum_{n=2}^\infty n(n-1)a_nx^{n-2}
+\sum_{n=2}^\infty n(n-1)a_nx^n
+\sum_{n=1}^\infty na_nx^{n-1}
+2\sum_{n=0}^\infty a_nx^n
=\sum_{n=0}^\infty[(n+2)(n+1)a_{n+2}+(n+1)a_{n+1}+(n^2-n+2)a_n]x^n=0$
if
 $a_{n+2}=-\dst\frac{1}{ n+2}a_{n+1}-
\dst\frac{n^2-n+2}{(n+2)(n+1)}a_n$.

\item
 The equation is equivalent to
 $(3+2t)y''+(1+2t)y'-(1-2t)y=0$ with $t=x-1$.
If $y=\dst\sum_{n=0}^\infty a_nt^n$, then
$(3+2t)y''+(1+2t)y'-(1-2t)y=
3\dst \sum_{n=2}^\infty n(n-1)a_nt^{n-2}
+2\sum_{n=2}^\infty n(n-1)a_nt^{n-1}
+\sum_{n=1}^\infty na_nt^{n-1}
+2\sum_{n=1}^\infty na_nt^n
-\sum_{n=0}^\infty a_nt^n
+2\sum_{n=0}^\infty a_nt^{n+1}
=(6a_2+a_1-a_0)+
\sum_{n=1}^\infty[3(n+2)(n+1)a_{n+2}+(n+1)(2n+1)a_{n+1}+
(2n-1)a_n+2a_{n-1}]t^n=0$  if $a_2=-\dst\frac{a_1-a_0}{6}$ and
$a_{n+2}=-\dst\frac{2n+1}{3(n+2)}a_{n+1}-\frac{2n-1}{3(n+2)(n+1)}a_n
-\frac{2}{3(n+2)(n+1)}a_{n-1}$,
$n\ge1$. Starting with $a_0=1$ and $a_1=-2$ yields
$y=\dst{1-2(x-1)+\frac{1}{2}(x-1)^2-\frac{1}{6}(x-1)^3+\frac{5}{36}(x-1)^4
-\frac{73}{1080}(x-1)^5+\cdots}$.


\item
 The equation is equivalent to
 $(1+t)y''+(2-2t)y'+(3+t)y=0$ with $t=x+3$.
If $y=\dst\sum_{n=0}^\infty a_nt^n$, then
$(1+t)y''+(2-2t)y'+(3+t)y=
\dst \sum_{n=2}^\infty n(n-1)a_nt^{n-2}
+\sum_{n=2}^\infty n(n-1)a_nt^{n-1}
+2\sum_{n=1}^\infty na_nt^{n-1}
-2\sum_{n=1}^\infty na_nt^n
+3\sum_{n=0}^\infty a_nt^n
+\sum_{n=0}^\infty a_nt^{n+1}
=(2a_2+2a_1+3a_0)+
\sum_{n=1}^\infty[(n+2)(n+1)a_{n+2}+(n+2)(n+1)a_{n+1}
-(2n-3)a_n+a_{n-1}]t^n=0$  if $a_2=-\dst\frac{2a_1+3a_0}{2}$ and
$a_{n+2}=\dst-a_{n+1}+\frac{(2n-3)a_n-a_{n-1}}{(n+2)(n+1)}$,
$n\ge1$. Starting with $a_0=2$ and $a_1=-2$ yields

$y=\dst{2-2(x+3)-(x+3)^2+(x+3)^3-\frac{11}{12}(x+3)^4+
\frac{67}{60}(x+3)^5+\cdots}$.


\item
 The equation is equivalent to
 $(1+2t)y''+3y'+(1-t)y=0$ with $t=x+1$.
If $y=\dst\sum_{n=0}^\infty a_nt^n$, then
$(1+2t)y''+3y'+(1-t)y=
\dst \sum_{n=2}^\infty n(n-1)a_nt^{n-2}
+2\sum_{n=2}^\infty n(n-1)a_nt^{n-1}
+3\sum_{n=1}^\infty na_nt^{n-1}
+\sum_{n=0}^\infty a_nt^n
-\sum_{n=0}^\infty a_nt^{n+1}
=(2a_2+3a_1+a_0)+
\sum_{n=1}^\infty[(n+2)(n+1)a_{n+2}+(2n+3)(n+1)a_{n+1}
+a_n-a_{n-1}]t^n=0$  if $a_2=-\dst\frac{3a_1+a_0}{2}$ and
$a_{n+2}=\dst-\frac{2n+3}{ n+2}a_{n+1}-\frac{a_n-a_{n-1}}{(n+2)(n+1)}$,
$n\ge1$. Starting with $a_0=2$ and $a_1=-3$ yields
$y=\dst{2-3(x+1)+\frac{7}{2}(x+1)^2-5(x+1)^3+\frac{197}{24}(x+1)^4
-\frac{287}{20}(x+1)^5+\cdots}$.


\item
 The equation is equivalent to
 $(6-2t)y''+(3+t)y=0$ with $t=x-2$.
If $y=\dst\sum_{n=0}^\infty a_nt^n$, then
$(6-2t)y''+(3+t)y=
6\dst \sum_{n=2}^\infty n(n-1)a_nt^{n-2}
-2\sum_{n=2}^\infty n(n-1)a_nt^{n-1}
+3\sum_{n=0}^\infty a_nt^n
+\sum_{n=0}^\infty a_nt^{n+1}
=(12a_2+3a_0)+
\sum_{n=1}^\infty[6(n+2)(n+1)a_{n+2}-2(n+1)na_{n+1}+3a_n
+a_{n-1}]t^n=0$  if $a_2=-\dst\frac{a_0}{4}$ and
$a_{n+2}=\dst\frac{n}{3(n+2)}a_{n+1}-\frac{3a_n+a_{n-1}}{6(n+2)(n+1)}$,
$n\ge1$. Starting with $a_0=2$ and $a_1=-4$ yields
$y=\dst{2-4(x-2)-\frac{1}{2}(x-2)^2+\frac{2}{9}(x-2)^3+\frac{49}{432}(x-2)^4
+\frac{23}{1080}(x-2)^5+\cdots}$.


\item
 The equation is equivalent to
 $(2+4t)y''-(1-2t)y=0$ with $t=x+4$.
If $y=\dst\sum_{n=0}^\infty a_nt^n$, then
$(2+4t)y''-(1-2t)y=
2\dst \sum_{n=2}^\infty n(n-1)a_nt^{n-2}
+4\sum_{n=2}^\infty n(n-1)a_nt^{n-1}
-\sum_{n=0}^\infty a_nt^n
+2\sum_{n=0}^\infty a_nt^{n+1}
=(4a_2-a_0)+
\sum_{n=1}^\infty[2(n+2)(n+1)a_{n+2}+4(n+1)na_{n+1}
-a_n+2a_{n-1}]t^n=0$  if $a_2=\dst\frac{a_0}{4}$ and
$a_{n+2}=\dst-\frac{2n}{
n+2}a_{n+1}+\frac{a_n-2a_{n-1}}{2(n+2)(n+1)}$,
$n\ge1$. Starting with $a_0=-1$ and $a_1=2$ yields
$y=\dst{-1+2(x+1)-\frac{1}{4}(x+1)^2+\frac{1}{2}(x+1)^3-\frac{65}{96}(x+1)^4
+\frac{67}{80}(x+1)^5+\cdots}$.

N=5; b=zeros(N,1); b(1)=-1;b(2)=2; b(3)=b(1)/4;
for n=1:N-2
   b(n+3)=-2*n*b(n+2)/(n+2)+(b(n+1)-2*b(n))/(2*(n+2)*(n+1));
end

\addtocounter{enumi}{-1}
\item%7.3.29
Let $Ly=(1+\alpha x+\beta x^2)y''+(\gamma+\delta x)y'+\epsilon y$.
If $y=\dst\sum_{n=0}^\infty a_nx^n$, then
$Ly=\dst\sum_{n=2}^\infty n(n-1)a_nx^{n-2}+
\alpha\sum_{n=2}^\infty n(n-1)a_nx^{n-1}+
\beta\sum_{n=2}^\infty n(n-1)a_nx^n+
\gamma\sum_{n=1}^\infty na_nx^{n-1}+
\delta\sum_{n=1}^\infty na_nx^n+
\epsilon\sum_{n=0}^\infty a_nx^n=
\sum_{n=0}^\infty (n+2)(n+1)a_{n+2}x^n+
\alpha\sum_{n=0}^\infty (n+1)na_{n+1}x^n+
\beta\sum_{n=0}^\infty n(n-1)a_nx^n+
\gamma\sum_{n=0}^\infty (n+1)a_{n+1}x^n+
\delta\sum_{n=0}^\infty na_nx^n+
\epsilon\sum_{n=0}^\infty a_nx^n=\sum_{n=0}^\infty b_nx^n$,
where
$b_n=(n+1)(n+2)a_{n+2}+(n+1)(\alpha
n+\gamma)a_{n+1}+[\beta n(n-1)+\delta n+\epsilon]a_n$, which implies
the conclusion.

\addtocounter{enumi}{-1}
\item%7.3.30
\begin{alist}
\item Let $\gamma=2\alpha$, $\delta=4\beta$, and $\epsilon=2\beta$
in Exercise~7.3.29 to obtain (B).

\item If $a_n=c_1r_1^n+c_2r_2^n$, then $a_{n+2}+\alpha
a_{n+1}+\beta a_n=c_1r_1^n(r_1^2+\alpha
r+\beta)+c_2r_2^n(r_2^2+\alpha
r_2+\beta)=c_1r_1^nP_0(r_1)+c_2r_2^nP_0(r_2)=0$, so $\{a_n\}$
satisfies (B). Since $1/r_1$ and $1/r_2$ are the zeros
of $P_0$, Theorem~7.2.1 implies  that $\sum_{n=0}^\infty
(c_1r_1^n+c_2r_2^n)x^n$ is a solution of (A) on $(-\rho,\rho)$.

\item If $|x|<\rho$, then $|r_1x|<\rho$ and $|r_2x|<1$, so
$\dst\sum_{n=0}^\infty r_i^nx^n=\frac{1}{1-r_ix}=y_i$, $i=1,2$.
Therefore,
\part{b} implies that $\{y_1,y_2\}$ is a fundamental set of solutions
of (A) on $(-\rho,\rho)$.

\item (A) can written as $P_0y''+2P_0'y'+P_0''y=(P_0y)''=0$.
Therefore,$P_0y=a+bx$ where $a$ and $b$ are arbitrary constants, and
 a partial fraction expansion shows that  the general solution of
(A)
on any interval not containing $1/r_1$ or $1/r_2$ is $y=\dst\frac{a+bx}{
P_0(x)}=\frac{c_1}{1-r_1x}+\frac{c_2}{1-r_2x}=c_1y_1+c_2y_2$.

\item If $a_n=c_1r_1^n+c_2r_2^n$, then $a_{n+2}+\alpha
a_{n+1}+\beta a_n=c_1r_1^n(r_1^2+\alpha
r+\beta)+c_2r_1^n[(n+2)r_1^2+\alpha(n+1)
r_2+\beta n)=(c_1+nc_2)r_1^nP_0(r_1)+c_2r_1^nP_0'(r_1)=0$, so
$\{a_n\}$ satisfies (B). Since $1/r_1$ is the only  zero
of $P_0$, Theorem~7.2.1 implies  that $\sum_{n=0}^\infty
(c_1+c_2n)r_1^n)x^n$ is a solution of (A) on $(-\rho,\rho)$.


\item If $|x|<\rho$, then $|r_1x|<\rho$, so
$\dst\sum_{n=0}^\infty r_1^nx^n=\frac{1}{1-r_1x}=y_1$.
Differentiating this and multiplying the result by $x$
shows that
$\dst\sum_{n=0}^\infty nr_1^nx^n=\frac{r_1x}{(1-r_1x)^2}=r_1y_2$.
Therefore,
\part{e} implies that $\{y_1,y_2\}$ is a fundamental set of solutions
of (A) on $(-\rho,\rho)$.

\item The argument is the same as in \part{c}, but now the
partial fraction expansion can be written as
 $y=\dst\frac{a+bx}{
P_0(x)}=\frac{c_1}{1-r_1x}+\frac{c_2x}{(1-r_2x)^2}=c_1y_1+c_2y_2$.
\end{alist}

\item
If $y=\dst\sum_{n=0}^\infty a_nx^n$, then
$y''+2xy'+(3+2x^2)y=
\dst \sum_{n=2}^\infty n(n-1)a_nx^{n-2}
+2\sum_{n=1}^\infty na_nx^n
+3\sum_{n=0}^\infty a_nx^n
+2\sum_{n=0}^\infty a_nx^{n+2}
=(2a_2+3a_0)+(6a_3+5a_1)x+
\sum_{n=2}^\infty[(n+2)(n+1)a_{n+2}+(2n+3)a_n+2a_{n-2}]x^n=0$ if
$a_2=-3a_0/2$, $a_3=-5a_1/6$, and
$a_{n+2}=-\dst\frac{(2n+3)a_n+2a_{n-2}}{(n+2)(n+1)}$,
 $n\ge2$.
Starting with $a_0=1$ and $a_1=-2$ yields
$y=\dst{1-2x-\frac{3}{2}x^2+\frac{5}{3}x^3+\frac{17}{24}x^4-\frac{11}{20}x^5+\cdots}$.


\item
If $y=\dst\sum_{n=0}^\infty a_nx^n$, then
$y''+5xy'-(3-x^2)y=
\dst \sum_{n=2}^\infty n(n-1)a_nx^{n-2}
+5\sum_{n=1}^\infty na_nx^n
-3\sum_{n=0}^\infty a_nx^n
+\sum_{n=0}^\infty a_nx^{n+2}
=(2a_2-3a_0)+(6a_3+2a_1)x+
\sum_{n=2}^\infty[(n+2)(n+1)a_{n+2}+(5n-3)a_n+a_{n-2}]x^n=0$ if
$a_2=3a_0/2$, $a_3=-a_1/3$, and
$a_{n+2}=-\dst\frac{(5n-3)a_n+a_{n-2}}{(n+2)(n+1)}$,
 $n\ge2$.
Starting with $a_0=6$ and $a_1=-2$ yields
$y=\dst{6-2x+9x^2+\frac{2}{3}x^3-\frac{23}{4}x^4-\frac{3}{10}x^5+\cdots}$.


\item
If $y=\dst\sum_{n=0}^\infty a_nx^n$, then
$y''-3xy'+(2+4x^2)y=
\dst \sum_{n=2}^\infty n(n-1)a_nx^{n-2}
-3\sum_{n=1}^\infty na_nx^n
+2\sum_{n=0}^\infty a_nx^n
+4\sum_{n=0}^\infty a_nx^{n+2}
=(2a_2+2a_0)+(6a_3-a_1)x+
\sum_{n=2}^\infty[(n+2)(n+1)a_{n+2}-(3n-2)a_n+4a_{n-2}]x^n=0$ if
$a_2=-a_0$, $a_3=a_1/6$, and
$a_{n+2}=\dst\frac{(3n-2)a_n-4a_{n-2}}{(n+2)(n+1)}$,
 $n\ge2$.
Starting with $a_0=3$ and $a_1=6$ yields
$y=\dst{3+6x-3x^2+x^3-2x^4-\frac{17}{20}x^5+\cdots}$.


\item
If $y=\dst\sum_{n=0}^\infty a_nx^n$, then
$3y''+2xy'+(4-x^2)y=
3\dst \sum_{n=2}^\infty n(n-1)a_nx^{n-2}
+2\sum_{n=1}^\infty na_nx^n
+4\sum_{n=0}^\infty a_nx^n
-\sum_{n=0}^\infty a_nx^{n+2}
=(6a_2+4a_0)+(18a_3+6a_1)x+
\sum_{n=2}^\infty[3(n+2)(n+1)a_{n+2}+(2n+4)a_n-a_{n-2}]x^n=0$ if
$a_2=-2a_0/3$, $a_3=-a_1/3$, and
$a_{n+2}=-\dst\frac{(2n+4)a_n-a_{n-2}}{3(n+2)(n+1)}$,
 $n\ge2$.
Starting with $a_0=-2$ and $a_1=3$ yields
$y=\dst-2+3x+\frac{4}{3}x^2-x^3-\frac{19}{54}x^4+\frac{13}{60}x^5+\cdots$.




\item
If $y=\dst\sum_{n=0}^\infty a_nx^n$, then
$(1+x)y''+x^2y'+(1+2x)y=
\dst \sum_{n=2}^\infty n(n-1)a_nx^{n-2}
+ \sum_{n=2}^\infty n(n-1)a_nx^{n-1}
+\sum_{n=1}^\infty na_nx^{n+1}
+\sum_{n=0}^\infty a_nx^n
+2\sum_{n=0}^\infty a_nx^{n+1}
=(2a_2+a_0)+
\sum_{n=1}^\infty[(n+2)(n+1)a_{n+2}+(n+1)na_{n+1}+a_n+(n+1)a_{n-1}]x^n=0$
if
$a_2=-a_0/2$ and
$a_{n+2}=-\dst\frac{(n+1)na_{n+1}+a_n+(n+1)a_{n-1}}{(n+2)(n+1)}$,
 $n\ge1$.
Starting with $a_0=-2$ and $a_1=3$ yields
$y=\dst{-2+3x+x^2-\frac{1}{6}x^3-\frac{3}{4}x^4+\frac{31}{120}x^5+\cdots}$.


\item
If $y=\dst\sum_{n=0}^\infty a_nx^n$, then
$(1+x^2)y''+(2+x^2)y'+xy=
\dst \sum_{n=2}^\infty n(n-1)a_nx^{n-2}
+\dst \sum_{n=2}^\infty n(n-1)a_nx^n
+2\sum_{n=1}^\infty na_nx^{n-1}
+\sum_{n=1}^\infty na_nx^{n+1}
+\sum_{n=1}^\infty a_nx^{n+1}
=(2a_2+2a_1)+
\sum_{n=1}^\infty[(n+2)(n+1)a_{n+2}+2(n+1)a_{n+1}+n(n-1)a_n+na_{n-1}]x^n=0$
if
$a_2=-a_1$ and
$a_{n+2}=-\dst\frac{[2(n+1)a_{n+1}+n(n-1)a_n+na_{n-1}]}{(n+2)(n+1)}$,
 $n\ge1$.
Starting with $a_0=-3$ and $a_1=5$ yields
$y=\dst{-3+5x-5x^2+\frac{23}{6}x^3-\frac{23}{12}x^4+\frac{11}{30}x^5+\cdots}$.


\item
 The equation is equivalent to
 $y''+(1+3t^2)y'+(1+2t)y=0$ with $t=x-2$.
If $y=\dst\sum_{n=0}^\infty a_nt^n$, then
$y''+(1+3t^2)y'+(1+2t)y=
\dst \sum_{n=2}^\infty n(n-1)a_nt^{n-2}
+\sum_{n=1}^\infty na_nt^{n-1}
+3\sum_{n=1}^\infty na_nt^{n+1}
+\sum_{n=0}^\infty a_nt^n
+2\sum_{n=0}^\infty a_nt^{n+1}=(2a_2+a_1+a_0)+
\sum_{n=1}^\infty[(n+2)(n+1)a_{n+2}+
(n+1)a_{n+1}+a_n+(3n-1)a_{n-1}]t^n=0$ if $a_2=-(a_1+a_0)/2$ and
$a_{n+2}=-\dst\frac{[(n+1)a_{n+1}+a_n+(3n-1)a_{n-1}]}{
(n+2)(n+1)}$,
$n\ge1$. Starting with $a_0=2$ and $a_1=-3$ yields

$y=\dst{2-3(x+2)+\frac{1}{2}(x+2)^2-\frac{1}{3}(x+2)^3+\frac{31}{24}(x+2)^4-
\frac{53}{120}(x+2)^5+\cdots}$.



\item
 The equation is equivalent to
 $(1-t^2)y''-(7-8t+t^2)y'+ty=0$ with $t=x+2$.
If $y=\dst\sum_{n=0}^\infty a_nt^n$, then
$(1-t^2)y''-(7-8t+t^2)y'+ty=
\dst \sum_{n=2}^\infty n(n-1)a_nt^{n-2}
- \sum_{n=2}^\infty n(n-1)a_nt^n
-7\sum_{n=1}^\infty na_nt^{n-1}
+8\sum_{n=1}^\infty na_nt^n
-\sum_{n=1}^\infty na_nt^{n+1}
+\sum_{n=0}^\infty a_nt^{n+1}
=(2a_2-7a_1)+
\sum_{n=1}^\infty[(n+2)(n+1)a_{n+2}
-7(n+1)a_{n+1}-n(n-9)a_n-(n-2)a_{n-1}]t^n=0$ if $a_2=7a_1/2$
and
$a_{n+2}=\dst\frac{[7(n+1)a_{n+1}+n(n-9)a_n+(n-2)a_{n-1}]}{
(n+2)(n+1)}$,
$n\ge1$. Starting with $a_0=2$ and $a_1=-1$ yields

$y=\dst{2-(x+2)-\frac{7}{2}(x+2)^2-\frac{43}{6}(x+2)^3-\frac{203}{24}(x+2)^4-\frac{167
}{30}(x+2)^5+\cdots}$.



\item
 The equation is equivalent to
 $(1+3t+2t^2)y''-(3+t-t^2)y'-(3+t)y=0$ with $t=x-1$.
If $y=\dst\sum_{n=0}^\infty a_nt^n$, then
$(1+3t+2t^2)y''-(3+t-t^2)y'-(3+t)y=
\dst \sum_{n=2}^\infty n(n-1)a_nt^{n-2}
+3\sum_{n=2}^\infty n(n-1)a_nt^{n-1}
+2\sum_{n=2}^\infty n(n-1)a_nt^n
-3\sum_{n=1}^\infty na_nt^{n-1}
-\sum_{n=1}^\infty na_nt^n
+\sum_{n=1}^\infty na_nt^{n+1}
-3\sum_{n=0}^\infty a_nt^n
-\sum_{n=0}^\infty a_nt^{n+1}
=(2a_2-3a_1-3a_0)+
\sum_{n=1}^\infty[(n+2)(n+1)a_{n+2}
+3(n^2-1)a_{n+1}+(2n^2-3n-3)(n+1)a_n+(n-2)a_{n-1}]t^n=0$ if
$a_2=3(a_1+a_0)/2$ and
$a_{n+2}=-\dst\frac{[3(n^2-1)a_{n+1}+(2n^2-3n-3)(n+1)a_n+(n-2)a_{n-1}]}{
(n+2)(n+1)}$,
$n\ge1$. Starting with $a_0=1$ and $a_1=0$ yields
$y=\dst{1+\frac{3}{2}(x-1)^2+\frac{1}{6}(x-1)^3-\frac{1}{8}(x-1)^5+\cdots}$.
\end{sectionSolutions}


\section{Regular Singular Points; Euler Equations}

\begin{sectionSolutions}
\item
$p(r)=r(r-1)-7r+7=(r-7)(r-1)$; $y=c_1x+c_2x^7$.




\item
$p(r)=r(r-1)+5r+4=(r+2)^2$;$y=x^{-2}(c_1+c_2 \ln x)$



\item
$p(r)=r(r-1)-3r+13=(r-2)^2+9$;
 $y=x^2[c_1 \cos (3 \ln x)+c_2 \sin (3 \ln x)]$.



\item
$p(r)=12r(r-1)-5r+6=(3r-2)(4r-3)$;
$y=c_1x^{2/3}+c_2 x^{3/4}$.



\item
$p(r)=3r(r-1)-r+1=(r-1)(3r-1)$; $y=c_1x+c_2x^{1/3}$.




\item
$p(r)=r(r-1)+3r+5=(r+1)^2+4$;
$y=\dst{
 \frac{1}{ x}\left[c_1\cos(2 \ln x)+c_2\sin(2 \ln x\right]}$


\item
$p(r)=r(r-1)-r+10=(r-1)^2+9$;
 $y=x\left[c_1\cos(3 \ln x)+c_2\sin(3 \ln x)\right]$.



\item
$p(r)=2r(r-1)+3r-1=(r+1)(2r-1)$;
$y=\dst{\frac{c_1}{ x}+c_2 x^{1/2}}$.




\item
$p(r)=2r(r-1)+10r+9=2(r+2)^2+1$;
$y=\dst{\frac{1}{ x^2}
\left[c_1\cos\left(\frac{1}{\sqrt{2}} \ln x\right)+c_2\sin
\left(\frac{1}{\sqrt{2}} \ln x\right) \right]}$.


\item
If $p(r)=ar(r-1)+br+c=a(r-r_1)^2$, then (A) $p(r_1)=p'(r_1)=0$.
If $y=ux^{r_1}$, then $y'=u'x^{r_1}+r_1ux^{r_1-1}$
and $y''=u''x^{r_1}+2r_1u'x_1^{r_1-1}+r_1(r_1-1)x_1^{r_1-2}$, so
\begin{align*}
ax^2y''+bxy'+cy&=ax^{r_1+2}u''+(2ar_1+b)x^{r_1+1}u'
+\left(ar_1(r_1-1)+br_1+c\right)x^{r_1}u\\
&=ax^{r_1+2}u''+p'(r_1)x_1^{r_1+1}u'+p(r)x^{r_1}u=ax^{r_1+2}u'',
\end{align*}
from (A). Therefore,$u''=0$, so $u=c_1+c_2x$ and
$y=x^{r_1}(c_1+c_2x)$.


\item
\begin{alist}
\item If $t=x-1$ and $Y(t)=y(t+1)=y(x)$, then
$\dst(1-x^2)y''-2xy'+\alpha(\alpha+1)y=
-t(2+t)\frac{d^2Y}{ dt^2}-2(1+t)\frac{dY}{ dt}+\alpha(\alpha+1)Y=0$,
so $y$ satisfies Legendre's equation if and only if $Y$ satisfies
(A) $\dst t(2+t)\frac{d^2Y}{ dt^2}+2(1+t)\frac{dY}{
dt}-\alpha(\alpha+1)Y=0$. Since (A) can be rewritten as
$\dst t^2(2+t)\frac{d^2Y}{ dt^2}+2t(1+t)\frac{dY}{
dt}-\alpha(\alpha+1)tY=0$, (A) has a regular singular point at
$t=0_0$.


\item If $t=x+1$ and $Y(t)=y(t-1)=y(x)$, then
$\dst(1-x^2)y''-2xy'+\alpha(\alpha+1)y=
t(2-t)\frac{d^2Y}{ dt^2}+2(1-t)\frac{dY}{ dt}+\alpha(\alpha+1)Y$,
so $y$ satisfies Legendre's equation if and only if $Y$ satisfies
(B) $\dst t(2-t)\frac{d^2Y}{ dt^2}+2(1-t)\frac{dY}{
dt}+\alpha(\alpha+1)Y$,
 Since (B) can be rewritten as
(B) $\dst t^2(2-t)\frac{d^2Y}{ dt^2}+2t(1-t)\frac{dY}{
dt}+\alpha(\alpha+1)tY$,
(B) has a regular singular point at
$t=0_0$.
\end{alist}
\end{sectionSolutions}

\section{The Method of Frobenius I}

\begin{sectionSolutions}
\item
$p_0(r)=r(3r-1)$;
$p_1(r)=2(r+1)$;
$p_2(r)=-4(r+2)$.

$a_1(r)=\dst-\frac{2}{3r+2}$;
$a_n(r)=\dst-\frac{2a_{n-1}(r)-4a_{n-2}(r)}{3n+3r-1}$,
$n\ge1$.

$r_1=1/3$;
$a_1(1/3)=-2/3$;
$a_n(1/3)=\dst-\frac{2a_{n-1}(1/3)-4a_{n-2}(1/3)}{3n}$,
$n\ge$1;

$y_1=\dst{x^{1/3}\left(1-\frac{2}{3}x+\frac{8}{9}x^2-
\frac{40}{81}x^3+\cdots\right)}$.


$r_2=0$;
$a_1(0)=-1$;
$a_n(0)=-\dst\frac{2a_{n-1}(0)-4a_{n-2}(0)}{3n-1}$,
$n\ge1$;

$y_2=\dst{1-x+\frac{6}{5}x^2-\frac{4}{5}x^3+\cdots}$.

%N=5; b1=zeros(N,1); b1(1)=1;b1(2)=-2/3;
%b2=zeros(N,1); b2(1)=1;b2(2)=-1;
%for n=2:N-1
%   b1(n+1)=-(2*b1(n)-4*b1(n-1))/(3*n);
%   b2(n+1)=-(2*b2(n)-4*b2(n-1))/(3*n-1);
%end
%[b1 b2]



\item
$p_0(r)=(r+1)(4r-1)$;
$p_1(r)=2(r+2)$;
$p_2(r)=4r+7$.

$a_1(r)=\dst-\frac{2}{ 4r+3}$;
$a_n(r)=\dst-\frac{2}{ 4n+4r-1}a_{n-1}(r)-\frac{1}{ n+r+1}a_{n-2}(r)$,
$n\ge1$.

$r_1=1/4$;
$a_1(1/4)=-1/2$;
$a_n(1/4)=\dst-\frac{1}{ 2n}a_{n-1}(1/4)-\frac{4}{ 4n+5}a_{n-2}(1/4)$,
$n\ge$1;


$y_1=\dst{x^{1/4}\left(1-\frac{1}{2}x-\frac{19}{104}x^2+
\frac{1571}{10608}x^3+\cdots\right)}$.

$r_2=-1$;
$a_1(-1)=2$;
$a_n(-1)=\dst-\frac{2}{ 4n-5}a_{n-1}(-1)-\frac{1}{ n}a_{n-2}(-1)$,
$n\ge1$;

$y_2=\dst{x^{-1}\left(1+2x-\frac{11}{6}x^2-\frac{1}{7}x^3+\cdots\right)}$.


\item
$p_0(r)=r(5r-1)$;
$p_1(r)=(r+1)^2$;
$p_2(r)=2(r+2)(5r+9)$.


$a_1(r)=\dst-\frac{r+1}{5r+4}\dst$;
$a_n(r)=\dst-\frac{n+r}{5n+5r-1}a_{n-1}(r)-2a_{n-2}(r)$,
$n\ge1$.

$r_1=1/5$;
$a_1(1/5)=-6/25$;
$a_n(1/5)=\dst-\frac{5n+1}{25n}a_{n-1}(1/5)-2a_{n-2}(1/5)$,
$n\ge$1;

$y_1=\dst{x^{1/5}\left(1-\frac{6}{25}x-\frac{1217}{625}x^2+
\frac{41972}{46875}x^3
+\cdots\right)}$.


$r_2=0$;
$a_1(0)=-1/4$;
$a_n(0)=\dst-\frac{n}{5n-1}a_{n-1}(0)-2a_{n-2}(0)$,
$n\ge1$;

$y_2=\dst{x-\frac{1}{4}x^2-\frac{35}{18}x^3+\frac{11}{12}x^4+\cdots}$.


\item
$p_0(r)=(3r-1)(6r+1)$;
$p_1(r)=(3r+2)(6r+1)$;
$p_2(r)=3r+5$.

$a_1(r)=\dst-\frac{6r+1}{6r+7}$;
$a_n(r)=\dst-\frac{6n+6r-5}{6n+6r+1}a_{n-1}(r)-\frac{1}{6n+6r+1}a_{n-2}(r)$,
$n\ge1$.

$r_1=1/3$;
$a_1(1/3)=-1/3$;
$a_n(1/3)=\dst-\frac{2n-1}{2n+1}a_{n-1}(1/3)-\frac{1}{6n+3}a_{n-2}(1/3)$,
$n\ge$1;

$y_1=\dst{x^{1/3}\left(1-\frac{1}{3}x+\frac{2}{15}x^2-
\frac{5}{63}x^3+\cdots\right)}$.


$r_2=-1/6$;
$a_1(-1/6)=0$;
$a_n(-1/6)=\dst-\frac{n-1}{ n}a_{n-1}(-1/6)-\frac{1}{6n}a_{n-2}(-1/6)$,
$n\ge1$;

$y_2=\dst{x^{-1/6}\left(1-\frac{1}{12}x^2+\frac{1}{18}x^3+\cdots\right)}$.


\item
$p_0(r)=(2r+1)(5r-1)$;
$p_1(r)=(2r-1)(5r+4)$;
$p_2(r)=2(2r+5)(5r-1)$.

$a_1(r)=\dst-\frac{2r-1}{2r+3}$;
$a_n(r)=\dst-\frac{2n+2r-3}{2n+2r+1}a_{n-1}(r)-
\frac{10n+10r-22}{5n+5r-1}a_{n-2}(r)$,
$n\ge1$.

$r_1=1/5$;
$a_1(1/5)=3/17$;
$a_n(1/5)=\dst-\frac{10n-13}{10n+7}a_{n-1}(1/5)-\frac{2n-4}{ n}a_{n-2}(1/5)$,
$n\ge$1;


$y_1=\dst{x^{1/5}\left(1+\frac{3}{17}x-\frac{7}{153}x^2-
\frac{547}{5661}x^3+\cdots\right)}$.

$r_2=-1/2$;
$a_1(-1/2)=1$;
$a_n(-1/2)=\dst-\frac{n-2}{ n}a_{n-1}(-1/2)-\frac{20n-54}{10n-7}a_{n-2}(-1/2)$,
$n\ge1$;

$y_2=\dst{x^{-1/2}\left(1+x+\frac{14}{13}x^2-\frac{556}{897}x^3+\cdots\right)}$.

\addtocounter{enumi}{2}
\item%7.5.14
$p_0(r)=(r+1)(2r-1)$;
$p_1(r)=2r+1$;
$a_n(r)=-\dst\frac{1}{ n+r+1}a_{n-1}(r)$.

$r_1=1/2$; $a_n(1/2)=\dst-\frac{2}{ 2n+3}a_{n-1}(1/2)$;
$y_1=\dst{x^{1/2}\sum_{n=0}^\infty\frac{(-2)^n}{\prod_{j=1}^n(2j+3)}x^n}$.

$r_2=-1$; $a_n(-1)=\dst-\frac{1}{ n}a_{n-1}(-1)$;
$y_2=\dst{x^{-1}\sum_{n=0}^\infty\frac{(-1)^n}{ n!}x^n}$.


\item
$p_0(r)=(r+2)(2r-1)$;
$p_1(r)=r+3$;
$a_n(r)=-\dst\frac{1}{ 2n+2r-1}a_{n-1}(r)$.

$r_1=1/2$;
$a_n(1/2)= -\dst\frac{1}{ 2n}a_{n-1}(1/2)$;
$y_1=\dst{x^{1/2}\sum_{n=0}^\infty \frac{(-1)^n}{2^nn!}x^n}$.


$r_2=-2$;
$a_n(-2)=-\dst\frac{1}{ 2n-5}a_{n-1}(-2)$;
$y_2=\dst{\frac{1}{ x^2}\sum_{n=0}^\infty\frac{(-1)^n}{\prod_{j=1}^n(2j-5)}
x^n}$.


\item
$p_0(r)=(r-1)(2r-1)$;
$p_1(r)= -2$;
$a_n(r)=\dst\frac{2}{ (n+r-1)(2n+2r-1)}a_{n-1}(r)$.

$r_1=1$;
$a_n(1)=\dst\frac{2}{ n(2n+1})a_{n-1}(1)$;
$y_1=\dst{x\sum_{n=0}^\infty\frac{2^n}{
n!\prod_{j=1}^n(2j+1)}x^n}$.

$r_2=1/2$;
$a_n(1/2)=\dst\frac{2}{ n(2n-1})a_{n-1}(1/2)$;
$y_2=\dst{x^{1/2}\sum_{n=0}^\infty\frac{2^n}{
n!\prod_{j=1}^n(2j-1)}x^n}$.


\item
$p_0(r)=(r-1)(3r+1)$;
$p_1(r)=r-3$;
$a_n(r)=-\dst\frac{n+r-4}{ (n+r-1)(3n+3r+1)}a_{n-1}(r)$.

$r_1=$;
$a_n(1)=-\dst\frac{n-3}{ n(3n+4)}a_{n-1}(1)$;
$y_1=\dst{x\left(1+\frac{2}{7}x+\frac{1}{70}x^2\right)}$.

$r_2=-1/3$;
$a_n(-1/3)=-\dst\frac{3n-13}{ 3n(3n-4)}a_{n-1}(-1/3)$;
$y_2=\dst{x^{-1/3}\sum_{n=0}^\infty\frac{(-1)^n}{3^nn!}\left(\prod_{j=1}^n
\frac{3j-13}{3j-4}\right)x^n}$.


\item
$p_0(r)=(r-1)(4r-1)$;
$p_1(r)=r(r+2)$;
$a_n(r)=-\dst\frac{n+r+1}{ 4n+4r-1}a_{n-1}(r)$.

$r_1=1$;
$a_n(1)= -\dst\frac{n+2}{ 4n+3}a_{n-1}(1)$;
$y_1=\dst{x\sum_{n=0}^\infty\frac{(-1)^n(n+2)!}{2\prod_{j=1}^n(4j+3)}
x^n}$.

$r_2=1/4$;
$a_n(1/4)=-\dst\frac{4n+5}{ 16n}a_{n-1}(1/4)$;
$y_2=\dst{x^{1/4}\sum_{n=0}^\infty\frac{(-1)^n}{16^nn!}\prod_{j=1}^n(4j+5)
x^n}$


\item
$p_0(r)=(r+1)(3r-1)$;
$p_1(r)=2(r+2)(2r+3)$;
$a_n(r)=-2\dst\frac{2n+2r+1}{ 3n+3r-1}a_{n-1}(r)$.

$r_1=1/3$;
$a_n(1/3)=-2\dst\frac{6n+5}{ 9n}a_{n-1}(1/3)$;
$y_1=\dst{x^{1/3}\sum_{n=0}^\infty\frac{(-1)^n}{
n!}\left(\frac{2}{9}\right)^n\left(\prod_{j=1}^n(6j+5)\right)
x^n}$;

$r_2=-1$;
$a_n(-1)=-2\dst\frac{2n-1}{ 3n-4}a_{n-1}(-1)$;
$y_2=\dst{x^{-1}\sum_{n=0}^\infty(-1)^n2^n\left(\prod_{j=1}^n
\frac{2j-1}{3j-4}\right) x^n}$

\addtocounter{enumi}{2}
\item%7.5.28
$p_0(r)=(2r-1)(4r-1)$;
$p_1(r)=(r+1)^2$;
$a_n(r)=\dst-\frac{(n+r)^2}{(2n+2r-1)(4n+4r-1)} a_{n-1}(r)$.

$r_1=1/2$;
$a_n(1/2)=\dst-\frac{4n^2+4n+1}{8n(4n+1)} a_{n-1}(1/2)$;
$y_1=\dst{x^{1/2}\left(1-\frac{9}{40}x+\frac{5}{128}x^2-\frac{245}{39936}x^3
+\cdots\right)}$.

$r_2=1/4$;
$a_n(1/4)=\dst-\frac{16n^2+8n+1}{32n(4n-1)} a_{n-1}(1/4)$;
$y_2=\dst{x^{1/4}\left(1-\frac{25}{96}x+\frac{675}{14336}x^2-
\frac{38025}{5046272}x^3
+\cdots\right)}$.


\item
$p_0(r)=(2r-1)(2r+1)$;
$p_1(r)=(2r+1)(3r+1)$;
$a_n(r)=\dst-\frac{(3n+3r-2)}{(2n+2r+1)} a_n(r)$.

$r_1=1/2$;
$a_n(1/2)=\dst-\frac{6n-1}{4(n+1)} a_{n-1}(1/2)$;
$y_1=\dst{x^{1/2}\left(1-\frac{5}{8}x+\frac{55}{96}x^2
-\frac{935}{1536}x^3+\cdots\right)}$.

$r_2=-1/2$;
$a_n(-1/2)=\dst-\frac{6n-7}{4n} a_{n-1}(-1/2)$;
$y_2=\dst{x^{-1/2}\left(1+\frac{1}{4}x-\frac{5}{32}x^2
-\frac{55}{384}x^3+\cdots\right)}$.


\item
$p_0(r)=(2r+1)(3r+1)$;
$p_1(r)=(r+1)(r+2)$;
$a_n(r)=\dst\frac{(n+r)(n+r+1)}{(2n+2r+1)(3n+3r+1)} a_n(r)$.

$r_1=-1/3$;
$a_n(-1/3)=\dst-\frac{(3n-1)(3n+2)}{9n(6n+1)}a_{n-1}(-1/3)$;
$y_1=\dst{x^{-1/3}\left(1-\frac{10}{63}x+\frac{200}{7371}x^2-
\frac{17600}{3781323}x^3
+\cdots\right)}$.

$r_2=-1/2$;
$a_n(-1/2)=\dst-\frac{(2n-1)(2n+1)}{4n(6n-1)} a_{n-1}(-1/2)$;
$y_2=\dst{x^{-1/2}\left(1-\frac{3}{20}x+\frac{9}{352}x^2
-\frac{105}{23936}x^3
+\cdots\right)}$.



\item
$p_0(r)=(2r-1)(4r-1)$;
$p_2(r)=-(2r+3)(4r+3)$;
$a_{2m}(r)=\dst\frac{8m+4r-5}{8m+4r-1}a_{2m-2}(r)$.

$r_1=1/2$;
$a_{2m}(1/2)=\dst\frac{8m-3}{8m+1}a_{2m-2}(1/2)$;
$y_1=\dst{x^{1/2}\sum_{m=0}^\infty\left(\prod_{j=1}^m\frac{8j-3}{8j+
1}\right)x^{2m}}$.

$r_2=1/4$;
$a_{2m}(1/4)=\dst\frac{2m-1}{2m}a_{2m-2}(1/4)$;
$y_2=\dst{x^{1/4}\sum_{m=0}^\infty\frac{1}{2^mm!}\left(\prod_{j=1}^m(2j-1)
\right)x^{2m}}$

\item
$p_0(r)=r(3r-1)$;
$p_2(r)=(r-4)(r+2)$;
$a_{2m}(r)=-\dst\frac{2m+r-6}{6m+3r-1}a_{2m-2}(r)$.

$r_1=1/3$;
$a_{2m}(1/3)=-\dst\frac{6m-17}{18m}a_{2m-2}(1/3)$;
$y_1=\dst{x^{1/3}\sum_{m=0}^\infty\frac{(-1)^m}{18^mm!}\left(\prod_{j=1}^m
(6j-17)\right) x^{2m}}$.

$r_2=0$;
$a_{2m}(0)=-\dst\frac{2m-6}{6m-1}a_{2m-2}(0)$;
$y_2=\dst{1+\frac{4}{5}x^2+\frac{8}{55}x^4}$


\item
$p_0(r)=(2r-1)(3r-1)$;
$p_2(r)=-(r+1)(3r+5)$;
$a_{2m}(r)=\dst\frac{2m+r-1}{4m+2r-1}a_{2m-2}(r)$.

$r_1=1/2$;
$a_{2m}(1/2)=\dst\frac{4m-1}{8m}a_{2m-2}(1/2)$;
$y_1=\dst{x^{1/2}\sum_{m=0}^\infty\frac{1}{8^mm!}\left(\prod_{j=1}^m(4j-1)
 \right)x^{2m}}$.

$r_2=1/3$;
$a_{2m}(1/3)=\dst\frac{6m-2}{12m-1}a_{2m-2}(1/3)$;
$y_2=\dst{x^{1/3}\sum_{m=0}^\infty2^m\left(\prod_{j=1}^m\frac{3j-1}{12j-1}
\right) x^{2m}}$.


\item
$p_0(r)=(2r-1)(2r+1)$;
$p_1(r)=(r+1)(2r+3)$;
$a_{2m}(r)=-\dst\frac{2m+r-1}{4m+2r+1}a_{2m-2}(r)$.

$r_1=1/2$;
$a_{2m}(1/2)=-\dst\frac{4m-1}{4(2m+1)}a_{2m-2}(1/2)$;
$y_1=\dst{x^{1/2}\sum_{m=0}^\infty\frac{(-1)^m}{4^m}\left(\prod_{j=1}^m
\frac{4j-1}{2j+1}\right)x^{2m}}$.

$r_2=-1/2$;
$a_{2m}(-1/2)=-\dst\frac{4m-3}{8m}a_{2m-2}(-1/2)$;
$y_2=\dst{x^{-1/2}\sum_{m=0}^\infty\frac{(-1)^m}{8^mm!}\left(\prod_{j=1}^m
(4j-3)\right)x^{2m}}$


\item
$p_0(r)=(r+1)(3r-1)$;
$p_1(r)=(r-1)(3r+5)$;
$a_{2m}(r)=-\dst\frac{2m+r-3}{2m+r+1}a_{2m-2}(r)$.

$r_1=1/3$;
$a_{2m}(1/3)=-\dst\frac{3m-4}{3m+2}a_{2m-2}(1/3)$;
$y_1=\dst{x^{1/3}\sum_{m=0}^\infty(-1)^m\left(\prod_{j=1}^m
\frac{3j-4}{3j+2}\right)x^{2m}}$.

$r_2=-1$;
$a_{2m}(-1)=-\dst\frac{m-2}{ m}a_{2m-2}(-1)$;
 $y_2=\dst{x^{-1}(1+x^2)}$


\item
$p_0(r)=(r+1)(2r-1)$;
$p_1(r)=r^2$;
$a_{2m}(r)=-\dst\frac{(2m+r-2)^2}{(2m+r+1)(4m+2r-1)}a_{2m-2}(r)$.

$r_1=1/2$;
$a_{2m}(1/2)=-\dst\frac{(4m-3)^2}{8m(4m+3)}a_{2m-2}(1/2)$;
$y_1=\dst{x^{1/2}\sum_{m=0}^\infty\frac{(-1)^m}{8^mm!}\left(\prod_{j=1}^m
\frac{(4j-3)^2}{4j+3}\right)x^{2m}}$.

$r_2=-1$;
$a_{2m}(-1)=-\dst\frac{(2m-3)^2}{2m(4m-3)}a_{2m-2}(-1)$;
$y_2=\dst{x^{-1}\sum_{m=0}^\infty\frac{(-1)^m}{2^mm!}\left(\prod_{j=1}^m
\frac{(2j-3)^2}{4j-3}\right)x^{2m}}$.


\item
$p_0(r)=(3r-1)(3r+1)$;
$p_1(r)=3r+5$;
$a_{2m}(r)=-\dst\frac{1}{6m+3r+1}a_{2m-2}(r)$.

$r_1=1/3$;
$a_{2m}(1/3)=-\dst\frac{1}{2(3m+1)}a_{2m-2}(1/3)$;
$y_1=\dst{x^{1/3}\sum_{m=0}^\infty\frac{(-1)^m}{2^m\prod_{j=1}^m(3j+1)}
x^{2m}}$.

$r_2=-1/3$;
$a_{2m}(-1/3)=-\dst\frac{1}{6m}a_{2m-2}(-1/3)$;
$y_2=\dst{x^{-1/3}\sum_{m=0}^\infty\frac{(-1)^m}{6^mm!} x^{2m}}$


\item
$p_0(r)=2(r+1)(4r-1)$;
$p_2(r)=(r+3)^2$;
$a_{2m}(r)=\dst-\frac{2m+r+1}{2(8m+4r-1)} a_{2m-2}(r)$.

$r_1=1/4$;
$a_{2m}(1/4)=\dst-\frac{8m+5}{64m} a_{2m-2}(1/4)$;
$y_1=\dst{x^{1/4}\left(1-\frac{13}{64}x^2+\frac{273}{8192}x^4-
\frac{2639}{524288}x^6 +\cdots\right)}$.


$r_2=-1$;
$a_{2m}(-1)=\dst-\frac{m}{8m-5} a_{2m-2}(-1)$;
$y_2=\dst{x^{-1}\left(1-\frac{1}{3}x^2+\frac{2}{33}x^4-\frac{2}{209}x^6
+\cdots\right)}$.


\item
$p_0(r)=(2r-1)(2r+1)$;
$p_2(r)=(2r+5)^2$;
$a_{2m}(r)=\dst-\frac{4m+2r+1}{4m+2r-1} a_{2m-2}(r)$.

$r_1=1/2$;
$a_{2m}(1/2)=\dst-\frac{2m+1}{2m} a_{2m-2}(1/2)$;
$y_1=\dst{x^{1/2}\left(1-\frac{3}{2}x^2+\frac{15}{8}x^4-\frac{35}{16}x^6
+\cdots\right)}$.


$r_2=-1/2$;
$a_{2m}(-1/2)=\dst-\frac{2m}{2m-1} a_{2m-2}(-1/2)$;
$y_2=\dst{x^{-1/2}\left(1-2x^2+\frac{8}{3}x^4-\frac{16}{5}x^6
+\cdots\right)}$.


\item
\begin{alist}
\item
Multiplying (A) $c_1y_1+c_2y_2\equiv0$ by $x^{-r_2}$ yields
$c_1x^{r_1-r_2}\sum_{n=0}^\infty a_nx^n+
c_2\sum_{n=0}^\infty b_nx^n=0$, $0<x<\rho$.
Letting $x\to0+$ shows that $c_2=0$, since $b_0=1$. Now (A)
reduces to $c_1y_1\equiv0$, so $c_1=0$. Therefore,$y_1$
and $y_2$ are linearly independent on $(0,\rho)$.

\item Since
$y_1=\sum_{n=0}^\infty a_n(r_1)x^n$ and
$y_2=\sum_{n=0}^\infty a_n(r_2)x^n$ are linearly independent solutions
of $Ly=0$  $(0,\rho)$, $\{y_1,y_2\}$ is a fundamental
set of solutions of $Ly=0$ on $(0,\rho)$, by Theorem~5.1.6.
\end{alist}


\item
\begin{alist}
\item  If $x>0$, then $|x|^rx^n=x^{n+r}$, so the assertions
are obvious. If $x<0$, then $|x|^r=(-x)^r$, so
$\dst\frac{d}{ dx}|x|^r=-r(-x)^{r-1}=\frac{r(-x)^r}{ x}=\frac{r|x|^r}{ x}$.
Therefore,(A) $\dst\frac{d}{ dx}(|x|^rx^n)=\frac{r|x|^r}{
x}x^n+|x|^r(nx^{n-1})=(n+r)|x|^rx^{n-1}$ and
$\dst\frac{d^2}{ dx^2}(|x|^rx^n)=(n+r)\frac{d}{ dx}(|x|^rx^{n-1})=
(n+r)(n+r-1)|x|^rx^{n-2}$, from (A) with $n$ replaced by $n-1$.
\end{alist}

 \item
\begin{alist}
\item Here $p_1\equiv0$, so Eqn.~(7.5.12) reduces to
$a_0(r)=1$, $a_1(r)=0$, $a_n(r)=-\dst\frac{p_2(n+r-2)}{
p_0(n+r)}a_{n-2}(r)$, $r\ge0$, which implies that $a_{2m+1}(r)=0$
for $m=1,2,3,\dots$. Therefore,Eqn.~(7.5.12) actually reduces to
$a_0(r)=1$, $a_{2m}(r)=\dst-\frac{p_2(2m+r-2)}{ p_0(2m+r)}$, which
holds because of condition (A).

\item Similar to the proof of Exercise~7.5.55\part{a}.

\item
$p_0(2m+r_1)=2m\alpha_0(2m+r_1-r_2)$,
which is nonzero if $m>0$, since $r_1-r_2\ge0$. Therefore, the
assumptions of Theorem~7.5.2 hold with $r=r_1$, and
$Ly_1=p_0(r_1)x^{r_1}=0$.
If $r_1-r_2$ is not an even integer, then
$p_0(2m+r_2)=2m\alpha_0(2m-r_1+r_2)\ne0$, $m=1,2,\cdots$.
Hence, the assumptions of Theorem~7.5.2 hold with $r=r_2$ and
 $Ly_2=p_0(r_2)x^{r_2}=0$.
From Exercise~7.5.52,
$\{y_1,y_2\}$ is a fundamental set of solutions.

\item
 Similar to the proof of Exercise~7.5.55\part{c}.
\end{alist}



\item
\begin{alist}
\item From Exercise~7.5.57, $b_n=0$ for $n\ge1$.
\end{alist}

\item
\begin{alist}
\item
$\dst(\alpha_0+\alpha_1x+\alpha_2x^2)\sum_{n=0}^\infty a_nx^n=
\alpha_0a_0+ (\alpha_0a_1+\alpha_1a_0)x+
\sum_{n=2}^\infty(\alpha_0a_n+\alpha_1a_{n-1}+\alpha_2a_{n-2})x^n=1$,
so $\dst\sum_{n=0}^\infty
a_nx^n=\frac{\alpha_0a_0}{\alpha_0+\alpha_1x+\alpha_2x^2}$.

\item  If
$\dst\frac{p_1(r-1)}{ p_0(r)}=
\frac{\alpha_1}{\alpha_0}$ and
$\dst\frac{p_2(r-2)}{ p_0(r)}=
\frac{\alpha_2}{\alpha_0}$, then Eqn.~(7.5.12) is equivalent to
$a_0(r)=1$, $\alpha_0a_1(r)+\alpha_1a_0(r)=0$, $\alpha_0a_n(r)
+\alpha_1a_{n-1}(r)+\alpha_2a_{n-2}(r)=0$,\quad  $n\ge2$.
Therefore,Theorem~7.5.2 implies the conclusion.
\end{alist}

\item
$p_0(r)=(2r-1)(3r-1)$;
$p_1(r)=0$;
$p_2(r)=2(2r+3)(3r+5)$;
$\dst\frac{p_1(r-1)}{ p_0(r)}=0=\frac{\alpha_1}{\alpha_0}$;
$\dst\frac{p_2(r-2)}{ p_0(r)}=2=\frac{\alpha_2}{\alpha_0}$;
$y_1=\dst\frac{x^{1/3}}{1+2x^2}$; $y_2=\dst\frac{x^{1/2}}{1+2x^2}$.


\item
$p_0(r)=5(3r-1)(3r+1)$;
$p_1(r)=(3r+2)(3r+4)$;
$p_2(r)=0$;
$\dst\frac{p_1(r-1)}{ p_0(r)}=\frac{1}{5}=\frac{\alpha_1}{\alpha_0}$;
$\dst\frac{p_2(r-2)}{ p_0(r)}=0=\frac{\alpha_2}{\alpha_0}$;
$y_1=\dst\frac{x^{1/3}}{5+x}$; $y_2=\dst\frac{x^{-1/3}}{5+x}$.


\item
$p_0(r)=(2r-3)(2r-1)$;
$p_1(r)=3(2r-1)(2r+1)$;
$p_2(r)=(2r+1)(2r+3)$;
$\dst\frac{p_1(r-1)}{ p_0(r)}=3=\frac{\alpha_1}{\alpha_0}$;
$\dst\frac{p_2(r-2)}{ p_0(r)}=1=\frac{\alpha_2}{\alpha_0}$;
$y_1=\dst\frac{x^{1/2}}{1+3x+x^2}$;  $y_2=\dst\frac{x^{3/2}}{1+3x+x^2}$.

\item
$p_0(r)=3(r-1)(4r-1)$;
$p_1(r)=2r(4r+3)$;
$p_2(r)=(r+1)(4r+7)$;
$\dst\frac{p_1(r-1)}{ p_0(r)}=\frac{2}{3}=\frac{\alpha_1}{\alpha_0}$;
$\dst\frac{p_2(r-2)}{ p_0(r)}=\frac{1}{3}=\frac{\alpha_2}{\alpha_0}$;
$y_1=\dst\frac{x}{3+2x+x^2}$; $y_2=\dst\frac{x^{1/4}}{3+2x+x^2}$.
\end{sectionSolutions}

\section{The Method of Frobenius II}

\begin{sectionSolutions}
\item
$p_0(r)=(r+1)^2$;
$p_1(r)=(r+2)(r+3)$;
$p_2(r)=(r+3)(2r-1)$;

$a_1(r)=\dst-\frac{r+3}{ r+2}$;
$a_n(r)=\dst-\frac{n+r+2}{ n+r+1}a_{n-1}(r)-\frac{2n+2r-5}{
n+r+1}a_{n-2}(r)$,
$n\ge2$.

$a_1'(r)=\dst\frac{1}{(r+2)^2}$;
$a_n'(r)=\dst-\frac{n+r+2}{ n+r+1}a_{n-1}'(r)-\frac{2n+2r-5}{
n+r+1}a_{n-2}'(r)+\frac{1}{(n+r+1)^2}a_{n-1}(r)-\frac{7}{(n+r+1)^2}a_{n-2}(r)$,
$n\ge2$.

$r_1=-1$;
$a_1(-1)=-2$;
$a_n(-1)=\dst-\frac{n+1}{ n}a_{n-1}(-1)-\frac{2n-7}{ n}a_{n-2}(-1)$,
$n\ge2$;

$\dst{y_1=x^{-1}\left(1-2x+\frac{9}{2}x^2-\frac{20}{3}x^3+\cdots\right)}$;

$a_1'(-1)=1$;
$a_n'(-1)=\dst-\frac{n+1}{ n}a_{n-1}(-1)-\frac{2n-7}{ n}a_{n-2}(-1)
+\frac{1}{ n^2}a_{n-1}(-1)-\frac{7}{ n^2}a_{n-2}(-1)$,
$n\ge2$;

$\dst{y_2=
y_1\ln x+1-\frac{15}{4}x+\frac{133}{18}x^2+\cdots}$.


\item
$p_0(r)=(2r-1)^2$;
$p_1(r)=(2r+1)(2r+3)$;
$p_2(r)=(2r+1)(2r+3)$;

$a_1(r)=\dst-\frac{2r+3}{2r+1}$;
$a_n(r)=\dst
-\frac{(2n+2r+1)a_{n-1}(r)-(2n+2r-3)a_{n-2}(r)}{2n+2r-1}$,
$n\ge2$.

$a_1'(r)=\dst\dst\frac{4}{(2r+1)^2}$;
$a_n'(r)=\dst
-\frac{(2n+2r+1)a_{n-1}'(r)-(2n+2r-3)a_{n-2}'(r)}{2n+2r-1}
+\frac{4(a_{n-1}(r)-a_{n-2}(r))}{(2n+2r-1)^2}$;
$n\ge2$.

$r_1=1/2$;
$a_1(1/2)=-2$;
$a_n(1/2)=\dst
-\frac{(n+1)a_{n-1}(1/2)+(n-1)a_{n-2}(1/2)}{ n}$;
$n\ge2$;

$\dst{y_1=x^{1/2}\left(1-2x+\frac{5}{2}x^2-2x^3+\cdots\right)}$;

$a_1'(1/2)=1$;
$a_n'(1/2)=\dst
-\frac{(n+1)a_{n-1}'(1/2)+(n-1)a_{n-2}'(1/2)}{ n}
+\frac{a_{n-1}(1/2)-a_{n-2}(1/2)}{ n^2}$,
$n\ge2$;

$\dst{y_2=
y_1\ln x+x^{3/2}\left(1-\frac{9}{4}x+\frac{17}{6}x^2+
\cdots\right)}$.


\item
$p_0(r)=(3r+1)^2$;
$p_1(r)=3(3r+4)$;
$p_2(r)=-2(3r+7)$;

$a_1(r)=\dst-\frac{3}{3r+4}$;
$a_n(r)=\dst
\frac{-3a_{n-1}(r)+2a_{n-2}(r)}{3n+3r+1}$;
$n\ge2$;

$a_1'(r)=\dst\frac{9}{(3r+4)^2}$;
$a_n'(r)=\dst
\frac{-3a_{n-1}'(r)+2a_{n-2}'(r)}{3n+3r+1}
+\frac{9a_{n-1}(r)-6a_{n-2}(r)}{(3n+3r+1)^2}$;
$n\ge2$.

$r_1=-1/3$;
$a_1(-1/3)=-1$;
$a_n(-1/3)=\dst
\frac{-3a_{n-1}(-1/3)+2a_{n-2}(-1/3)}{3n}$,
$n\ge2$;

$\dst{y_1=x^{-1/3}\left(1-x+\frac{5}{6}x^2-\frac{1}{2}x^3+\cdots\right)}$;

$a_1'(-1/3)=1$;
$a_n'(-1/3)=\dst
\frac{-3a_{n-1}'(r)+2a_{n-2}'(r)}{3n}
+\frac{3a_{n-1}(r)-2a_{n-2}(r)}{3n^2}$;
$n\ge2$;

$\dst{y_2=
y_1\ln x+x^{2/3}\left(1-\frac{11}{12}x+\frac{25}{36}x^2+
\cdots\right)}$.


\item
$p_0(r)=(r+2)^2$;
$p_1(r)=2(r+3)^2$;
$p_2(r)=3(r+4)$;

$a_1(r)=-2$;
$a_n(r)=\dst
-2a_{n-1}(r)-\frac{3a_{n-2}(r)}{ n+r+2}$;
$n\ge2$;


$a_1'(r)=0$;
$a_n'(r)=\dst
-2a_{n-1}'(r)-\frac{3a_{n-2}'(r)}{ n+r+2}
+\frac{3a_{n-2}(r)}{(n+r+2)^2}$,
$n\ge2$.

$r_1=-2$;
$a_1(-2)=-2$;
$a_n(-2)=\dst
-2a_{n-1}(-2)-\frac{3a_{n-2}(-2)}{ n}$,
$n\ge2$;

$\dst{y_1=x^{-2}\left(1-2x+\frac{5}{2}x^2-3x^3+\cdots\right)}$;

$a_1'(-2)=0$;
$a_n'(-2)=\dst
-2a_{n-1}(-2)-\frac{3a_{n-2}(-2)}{ n}+\frac{3a_{n-2}(-2)}{ n^2}$;
$n\ge2$;

$\dst{y_2=y_1\ln x+\frac{3}{4}-\frac{13}{6}x+\cdots}$.


\item
$p_0(r)=(4r+1)^2$;
$p_1(r)=4r+5$;
$p_2(r)=2(4r+9)$;

$a_1(r)=\dst-\frac{1}{4r+5}$;
$a_n(r)=\dst
-\frac{a_{n-1}(r)+2a_{n-2}(r)}{4n+4r+1}$;
$n\ge2$;


$a_1'(r)=\dst\frac{4}{(4r+5)^2}$;
$a_n'(r)=\dst
-\frac{a_{n-1}'(r)+2a_{n-2}'(r)}{4n+4r+1}
+\frac{4a_{n-1}(r)+8a_{n-2}(r)}{(4n+4r+1)^2}$;
$n\ge2$.


$r_1=-1/4$;
$a_1(-1/4)=-1/4$;
$a_n(-1/4)=\dst
-\frac{a_{n-1}(-1/4)+2a_{n-2}(-1/4)}{4n}$;
$n\ge2$;

$\dst{y_1=x^{-1/4}\left(1-\frac{1}{4}x-\frac{7}{32}x^2+\frac{23}{384}x^3
+\cdots\right)}$;

$a_1'(-1/4)=1/4$;
$a_n'(-1/4)=\dst
-\frac{a_{n-1}'(-1/4)+2a_{n-2}'(-1/4)}{4n}
+\frac{a_{n-1}(-1/4)+2a_{n-2}(-1/4)}{4n^2}$;
$n\ge2$;

$\dst{y_2=
y_1\ln x+x^{3/4}\left(\frac{1}{4}+\frac{5}{64}x-\frac{157}{2304}x^2
+\cdots\right)}$.


\item
$p_0(r)=(2r-1)^2$;
$p_1(r)=4$;

$a_n(r)=\dst-\frac{4}{(2n+2r-1)^2}a_{n-1}(r)$;

$a_n(r)=\dst\frac{(-4)^n}{\prod_{j=1}^n(2j+2r-1)}$.

By logarithmic differentiation,
$a_n'(r)=\dst a_n(r)\sum_{j=1}^n\frac{2}{2j+2r-1}$;

$r_1=1/2$;
$a_n(1/2)=\dst\frac{(-1)^n}{(n!)^2}$;

$a_n'(1/2)=a_n(1/2)\dst\left(-2\sum_{j=1}^n\frac{1}{ j}\right)$;

$\dst{y_1=x^{1/2}\sum_{n=0}^\infty\frac{(-1)^n}{(n!)^2}x^n}$;

$\dst{y_2=y_1\ln x-2x^{1/2}\sum_{n=1}^\infty
\frac{(-1)^n}{(n!)^2}\left(\sum_{j=1}^n\frac{1}{ j}\right)x^n}$;



\item
$p_0(r)=(r-2)^2$;
$p_1(r)=r^2$;
$a_n(r)=\dst-\frac{(n+r-1)^2}{(n+r-2)^2}a_{n-1}(r)$;
 $a_n(r)=\dst(-1)^n\frac{(n+r-1)^2}{(r-1)^2}$;
$a_n'(r)=(-1)^{n+1}\dst\frac{2n(r+n-1)}{(r-1)^3}$;
$r_1=2$;
$a_n(2)=(-1)^n(n+1)^2$;
$a_n'(2)=(-1)^{n+1}2n(n+1)$;
$\dst{y_1=x^2\sum_{n=0}^\infty} (-1)^n(n+1)^2x^n$;
$\dst{y_2=y_1\ln
x-2x^2\sum_{n=1}^\infty(-1)^nn(n+1)x^n}$.



\item
$p_0(r)=(5r-1)^2$;
$p_1(r)=r+1$;

$a_n(r)=\dst-\frac{(n+r)}{(5n+5r-1)^2}
a_{n-1}(r)$;

 $a_n(r)=(-1)^n\dst\prod_{j=1}^n\frac{(j+r)}{(5j+5r-1)^2}$;

By logarithmic differentiation,

$a_n'(r)=-\dst a_n(r)\sum_{j=1}^n\frac{(5j+5r+1)}{(j+r)(5j+5r-1)}$;

$r_1=1/5$;
$a_n(1/5)=\dst(-1)^n\prod_{j=1}^n\frac{(5j+1)}{125^n(n!)^2}$;

$a_n'(1/5)=a_n(1/5)\dst
\sum_{j=1}^n\frac{5j+2}{ j(5j+1)}$;

$\dst{y_1=x^{1/5}\sum_{n=0}^\infty\frac{(-1)^n\prod_{j=1}^n(5j+1)}{
125^n(n!)^2} x^n}$;

$\dst{y_2=
y_1\ln
x-x^{1/5}\sum_{n=1}^\infty\frac{(-1)^n\prod_{j=1}^n(5j+1)}{125^n(n!)^2}
\left(\sum_{j=1}^n\frac{5j+2}{ j(5j+1)}\right)x^n}$.


\item
$p_0(r)=(3r-1)^2$;
$p_1(r)=(2r-1)^2$;

$a_n(r)=\dst-\frac{(2n+2r-3)^2}{(3n+3r-1)^2}a_{n-1}(r)$;

 $a_n(r)=(-1)^n\dst\prod_{j=1}^n\frac{(2j+2r-3)^2}{(3j+3r-1)^2}$;

By logarithmic differentiation,

$a_n'(r)=14\dst a_n(r)\sum_{j=1}^n\frac{1}{(2j+2r-3)(3j+3r-1)}$;

$r_1=1/3$;
$a_n(1/3)=\dst\frac{(-1)^n\prod_{j=1}^n(6j-7)^2}{81^n(n!)^2}$;

$a_n'(1/3)=\dst14a_n(1/3)
\sum_{j=1}^n\frac{1}{ j(6j-7)})$;

$\dst{y_1=x^{1/3}\sum_{n=0}^\infty\frac{(-1)^n\prod_{j=1}^n(6j-7)^2}{
81^n(n!)^2}
x^n}$;

$\dst{y_2=y_1\ln
x+14x^{1/3}\sum_{n=1}^\infty\frac{(-1)^n\prod_{j=1}^n(6j-7)^2}{
81^n(n!)^2}\left(\sum_{j=1}^n\frac{1}{ j(6j-7)})\right)x^n}$.


\item
$p_0(r)=(r+1)^2$;
$p_1(r)=-2(r+2)(2r+3)$;

$a_n(r)=\dst\frac{2(2n+2r+1)}{ n+r+1}
a_{n-1}(r)$, $n\ge1$;
 $a_n(r)=2^n\dst\prod_{j=1}^n\frac{2j+2r+1}{ j+r+1}$;

By logarithmic differentiation,

$a_n'(r)=\dst a_n(r)\sum_{j=1}^n\frac{1}{(j+r+1)(2j+2r+1)}$;

$r_1=-1$;
$a_n(-1)=\dst\frac{2^n\prod_{j=1}^n(2j-1)}{ n!}$;

$a_n'(-1)=a_n(-1)\dst\sum_{j=1}^n\frac{1}{ j(2j-1)}$;

$\dst{y_1=\frac{1}{ x}\sum_{n=0}^\infty \frac{2^n\prod_{j=1}^n(2j-1)}{ n!}
x^n}$;

$\dst{y_2=y_1\ln
x+\frac{1}{ x}\sum_{n=1}^\infty\frac{2^n\prod_{j=1}^n(2j-1)}{ n!}
\left(\sum_{j=1}^n\frac{1}{ j(2j-1)}\right)x^n}$.



\item
$p_0(r)=2(r-2)^2$;
$p_1(r)=(r-1)(2r+1)$;

$a_n(r)=\dst-\frac{2n+2r-1}{2(n+r-2)}a_{n-1}(r)$;

 $a_n(r)=\dst\frac{(-1)^n}{ 2^n}\prod_{j=1}^n\frac{2j+2r-1}{ j+r-2}$;

By logarithmic differentiation,

$a_n'(r)=-3\dst a_n(r)\sum_{j=1}^n\frac{1}{(j+r-2)(2j+2r-1)}$;

$r_1=2$;
$a_n(2)=\dst\frac{(-1)^n\prod_{j=1}^n(2j+3)}{2^nn!}$;

$a_n'(2)=-3a_n(2)\dst\sum_{j=1}^n\dst\frac{1}{ j(2j+ 3)}$;

$\dst{y_1=x^2\sum_{n=0}^\infty
\frac{(-1)^n\prod_{j=1}^n(2j+3)}{2^nn!}x^n}$;

$\dst{y_2=y_1\ln x-3x^2\sum_{n=0}^\infty
\frac{(-1)^n\prod_{j=1}^n(2j+3)}{2^nn!}
\left(\sum_{j=1}^n\frac{1}{ j(2j+ 3)}\right)x^n}$.


\item
$p_0(r)=(r-3)^2$;
$p_1(r)=-2(r-1)(r+2)$;

$a_n(r)=\dst
\frac{2(n+r-2)(n+r+1)}{(n+r-3)^2}a_{n-1}(r)$;


$a_n'(r)=\dst
\frac{2(n+r-2)(n+r+1)}{(n+r-3)^2}a_{n-1}'(r)
-\frac{2(5n+5r-7)}{(n+r-3)^3}a_{n-1}(r)$;

$r_1=3$;
$a_n(3)=\dst
\frac{2(n+1)(n+4)}{ n^2}a_{n-1}(3)$;

$\dst{y_1=x^3(1+20x+180x^2+1120x^3+\cdots}$;

$a_n'(3)=\dst
\frac{2(n+1)(n+4)}{ n^2}a_{n-1}'(3)
-\frac{2(5n+8)}{ n^3}a_{n-1}(3)$;

$\dst{y_2=y_1\ln
x-x^4\left(26+324x+\frac{6968}{3}x^2+\cdots\right)}$


\item
$p_0(r)=r^2$;
$p_1(r)=r^2+r+1$;

$a_n(r)=\dst
-\frac{(n^2+n(2r-1)+r^2-r+1)}{(n+r)^2}a_{n-1}(r)$;


$a_n'(r)=\dst
-\frac{(n^2+n(2r-1)+r^2-r+1)}{(n+r)^2}a_{n-1}'(r)
-\frac{(n+r-2)}{(n+r)^3}a_{n-1}(r)$;

$r_1=0$;
$a_n(0)=\dst-\frac{(n^2-n+1)}{ n^2}a_{n-1}(0)$;

$\dst{y_1=1-x+\frac{3}{4}x^2-\frac{7}{12}x^3+\cdots}$;

$a_n'(0)=\dst
-\frac{(n^2-n+1)}{ n^2}a_{n-1}'(0)-\frac{(n-2)}{ n^3}a_{n-1}(0)$;

$\dst{y_2=y_1\ln x+x\left(1-\frac{3}{4}x+\frac{5}{9}x^2+\cdots\right)}$.



\item
$p_0(r)=(r-1)^2$;
$p_2(r)=r+1$;

$a_{2m}(r)=\dst-\frac{1}{ 2m+r-1}
a_{2m-2}(r)$, $n\ge1$;
 $a_{2m}(r)=\dst\frac{(-1)^m}{\prod_{j=1}^m(2j+r-1)}$;

By logarithmic differentiation,

$a_{2m}'(r)=\dst a_{2m}(r)\sum_{j=1}^m$;

$r_1=1$;
$a_{2m}(1)=\dst\frac{(-1)^m}{2^mm!}$;

$a_{2m}'(1)=-\dst\frac{1}{2}a_{2m}(1)\sum_{j=1}^m\frac{1}{ j}$;

$\dst{y_1=x\sum_{m=0}^\infty\frac{(-1)^m}{2^mm!}x^{2m}}$;

$\dst{y_2=y_1\ln x-\frac{x}{2}\sum_{m=1}^\infty\frac{(-1)^m}{
2^mm!}\left(\sum_{j=1}^m\frac{1}{ j}\right)x^{2m}}$.


\item
$p_0(r)=(2r-1)^2$;
$p_2(r)=2r+3$;

$a_{2m}(r)=\dst-\frac{1}{4m+2r-1}
a_{2m-2}(r)$;

 $a_{2m}(r)=\dst\frac{(-1)^m}{\prod_{j=1}^m(4j+2r-1)}$;

By logarithmic differentiation,

$a_{2m}'(r)=-2\dst a_{2m}(r)\sum_{j=1}^m\frac{1}{4j+2r-1}$;

$r_1=1/2$;
$a_{2m}(1/2)=\dst\frac{(-1)^m}{4^mm!}$;

$a_{2m}'(1/2)=-\dst\frac{1}{2}a_{2m}(1/2)\sum_{j=1}^m\frac{1}{ j}$;

$\dst{y_1=x^{1/2}\sum_{m=0}^\infty\frac{(-1)^m}{4^mm!}x^{2m}}$;

$\dst{y_2=y_1\ln x-\frac{x^{1/2}}{2}\sum_{m=1}^\infty\frac{(-1)^m}{
4^mm!}\left(\sum_{j=1}^m\frac{1}{ j}\right) x^{2m}}$.

\item
$p_0(r)=(2r-1)^2$;
$p_2(r)=(r+1)(2r+3)$;
$a_{2m}(r)=\dst-\frac{2m+r-1}{4m+2r-1}
a_{2m-2}(r)$;

 $a_{2m}(r)=(-1)^m\dst\prod_{j=1}^m\frac{2j+r-1}{4j+2r-1}$;

By logarithmic differentiation,

$a_{2m}'(r)=\dst a_{2m}(r)\sum_{j=1}^m\frac{1}{(2j+r-1)(4j+2r-1)}$;

$r_1=1/2$;
$a_{2m}(1/2)=\dst\frac{(-1)^m\prod_{j=1}^m(4j-1)}{8^mm!}$;

$a_{2m}'(1/2)=a_{2m}(1/2)\dst\sum_{j=1}^m\frac{1}{2j(4j-1)}$;


$\dst{y_1=x^{1/2}\sum_{m=0}^\infty\frac{(-1)^m\prod_{j=1}^m(4j-1)}{8^mm!}
x^{2m}}$;

$\dst{y_2=y_1\ln
x+\frac{x^{1/2}}{2}\sum_{m=1}^\infty\frac{(-1)^m\prod_{j=1}^m(4j-1)}{8^mm!}
\left(\sum_{j=1}^m\frac{1}{ j(4j-1)}\right)x^{2m}}$.


\item
$p_0(r)=(4r+1)^2$;
$p_2(r)=(r-1)(4r+9)$;

$a_{2m}(r)=\dst-\frac{2m+r-3}{8m+4r+1}
a_{2m-2}(r)$;

 $a_{2m}(r)=(-1)^m\dst\prod_{j=1}^m\frac{2j+r-3}{8j+4r+1}$;

By logarithmic differentiation,

$a_{2m}'(r)=\dst a_{2m}(r)\sum_{j=1}^m\frac{13}{(2j+r-3)(8j+4r+1)}$;

$r_1=-1/4$;
$a_{2m}(-1/4)=\dst\frac{(-1)^m\prod_{j=1}^m(8j-13)}{(32)^mm!}$;

$a_{2m}'(-1/4)=a_{2m}(-1/4)\dst\sum_{j=1}^m\frac{13}{2j(8j-13)}$;


$\dst{y_1=x^{-1/4}\sum_{m=0}^\infty\frac{(-1)^m\prod_{j=1}^m(8j-13)}{(32)^mm!}
x^{2m}}$;

$\dst{y_2=y_1\ln
x+\frac{13}{2}x^{-1/4}\sum_{m=1}^\infty\frac{(-1)^m\prod_{j=1}^m(8j-13)}{(32)^mm!}
\left(\sum_{j=1}^m\frac{1}{ j(8j-13)}\right)x^{2m}}$.


\item
$p_0(r)=(2r-1)^2$;
$p_2(r)=16r(r+1)$;

$a_{2m}(r)=\dst-\frac{16(2m+r-2)(2m+r-1)}{(4m+2r-1)^2}
a_{2m-2}(r)$;


$a_{2m}(r)=(-16)^m\dst\prod_{j=1}^m\frac{(2j+r-2)(2j+r-1)}{(4j+2r-1)^2}$;

By logarithmic differentiation,

$a_{2m}'(r)=\dst
a_{2m}(r)\sum_{j=1}^m\frac{8j+4r-5}{(2j+r-2)(2j+r-1)(4j+2r-1)}$;

$r_1=1/2$;
$a_{2m}(1/2)=\dst\frac{(-1)^m\prod_{j=1}^m(4j-3)(4j-1)}{4^m(m!)^2}$;

$a_{2m}'(1/2)=a_{2m}(1/2)\dst\sum_{j=1}^m\frac{8j-3}{ j(4j-3) (4j-1)}$;


$\dst{y_1=x^{1/2}\sum_{m=0}^\infty\frac{(-1)^m\prod_{j=1}^m(4j-3)(4j-1)}{4^m(m!)^2}
x^{2m}}$;

$\dst{y_2=y_1\ln
x+x^{1/2}\sum_{m=1}^\infty\frac{(-1)^m\prod_{j=1}^m(4j-3)(4j-1)}{4^m(m!)^2}
\left(\sum_{j=1}^m\frac{8j-3}{ j(4j-3) (4j-1)}\right)x^{2m}}$.


\item
$p_0(r)=(r+1)^2$;
$p_2(r)=(r+3)(2r-1)$;

$a_{2m}(r)=\dst-\frac{4m+2r-5}{2m+r+1}
a_{2m-2}(r)$;

 $a_{2m}(r)=(-1)^m\dst\prod_{j=1}^m\frac{4j+2r-5}{2j+r+1}$;

By logarithmic differentiation,

$a_{2m}'(r)=\dst a_{2m}(r)\sum_{j=1}^m\frac{7}{(2j+r+1)(4j+2r-5)}$;

$r_1=-1$;
$a_{2m}(-1)=\dst\frac{(-1)^m\prod_{j=1}^m(4j-7)}{2^mm!}$;

$a_{2m}'(-1)=a_{2m}(-1)\dst\sum_{j=1}^m\frac{7}{2j(4j-7)}$;

$\dst{y_1=\frac{1}{ x
}\sum_{m=0}^\infty\frac{(-1)^m\prod_{j=1}^m(4j-7)}{2^mm!} x^{2m}}$;

$\dst{y_2=y_1\ln
x+\frac{7}{2x}\sum_{m=1}^\infty\frac{(-1)^m\prod_{j=1}^m(4j-7)}{2^mm!}
\left(\sum_{j=1}^m\frac{1}{ j(4j-7)}\right)x^{2m}}$.


\item
$p_0(r)=(r-1)^2$;
$p_2(r)=r+1$;

$a_{2m}(r)=\dst-\frac{1}{2m+r-1}a_{2m-2}(r)$;


$a_{2m}'(r)=\dst
-\frac{1}{2m+r-1}a_{2m-2}'(r)+\frac{1}{(2m+r-1)^2}a_{2m-2}(r)$;

$r_1=1$;
$a_{2m}(1)=\dst-\frac{1}{2m}a_{2m-2}(1)$;

$\dst{y_1=x\left(1-\frac{1}{2}x^2+\frac{1}{8}x^4-\frac{1}{48}x^6+\cdots\right)}$;

$a_{2m}'(1)=\dst
-\frac{1}{2m}a_{2m-2}'(1)+\frac{1}{4m^2}a_{2m-2}(1)$,
$m\ge1$;


$\dst{y_2=y_1 \ln x+x^3\left(\frac{1}{4}-\frac{3}{32}x^2+
\frac{11}{576}x^4+\cdots\right)}$.


\item
$p_0(r)=2(r+3)^2$;
$p_2(r)=r^2-2r+2$;

$a_{2m}(r)=\dst
-\frac{4m^2+4m(r-3)+r^2-6r+10}{2(2m+r+3)^2}a_{2m-2}(r)$;


$a_{2m}'(r)=\dst
-\frac{4m^2+4m(r-3)+r^2-6r+10}{2(2m+r+3)^2}a_{2m-2}'(r)
-\frac{12m+6r-19}{(2m+r+3)^3}a_{2m-2}(r)$;

$r_1=-3$;
$a_{2m}(-3)=\dst-\frac{4m^2-24m+37}{8m^2}a_{2m-2}(-3)$;

$\dst{y_1=x^{-3}\left(1-\frac{17}{8}x^2+\frac{85}{256}x^4-
\frac{85}{18432}x^6+\cdots\right)}$;


$a_{2m}'(-3)=\dst
-\frac{4m^2-24m+37}{8m^2}a_{2m-2}'(-3)+\frac{37-12m}{8m^3}a_{2m-2}(-3)$,
$m\ge1$;

$\dst{y_2=y_1\ln x
+x^{-1}\left(\frac{25}{8}-\frac{471}{512}x^2+\frac{1583}{110592}x^4
+\cdots\right)}$.


\item
$p_0(r)=(r+1)^2$;
$p_1(r)=2(2-r)(r+1)$;
$r_1=-1$.

$a_n(r)=\dst\frac{2(n+r)(n+r-3)}{(n+r+1)^2}
a_{n-1}(r)$;
 $a_n(r)=2^n\dst\prod_{j=1}^n\frac{(j+r)(j+r-3)}{(j+r+1)^2}$,
$n\ge0$. Therefore,$a_n(-1)=0$ if $n\ge1$ and $y_1=1/x$.
If $n\ge4$, then $a_n(r)=(r+1)^2b_n(r)$, where $b_n'(-1)$
exists; therefore $a_n'(-1)=0$ if $n\ge4$. For $r=1,2,3$,
$a_n(r)=(r+1)c_n(r)$, where
$c_1(r)=\dst\frac{2(r-2)}{(r+2)^2}$,
$c_2(r)=\dst\frac{4(r-2)(r-1)}{(r+2)(r+3)^2}$,
$c_3(r)=\dst
\frac{8r(r-2)(r-1)}{(r+2)(r+3)(r+4)^2}$. Hence,
$a_1'(-1)=c_1(-1)=-6$, $a_2'(-1)=c_2(-1)=6$, $a_3'(-1)=c_3(-1)=-8/3$,
and
 $\dst{y_2=y_1\ln x-6+6x-\frac{8}{3}x^2}$.


\item
$p_0(r)=(r+1)^2$;
$p_1(r)=-(r-1)(r+2)$;
$r_1=-1$.


$a_n(r)=\dst\frac{n+r-2}{ n+r+1}
a_{n-1}(r)$;
 $a_n(r)=\dst\prod_{j=1}^m\frac{j+r-2}{ j+r+1}$,
$n\ge0$. Therefore,$a_1(-1)=-2$, $a_2(-1)=1$, and $a_n(-1)=0$
if $n\ge3$, so
$\dst{y_1=\frac{(x-1)^2}{ x}}$.

$a_1(r)=\dst\frac{r-1}{ r+2}$, $a_1'(r)=\dst\frac{3}{(r+2)^2}$,
$a_1'(-1)=3$;
$a_2(r)=\dst
\frac{r(r-1)}{(r+2)(r+3)}$,
$a_2'(r)=\dst
\frac{6(r^2+2r-1)}{(r+2)^2(r+3)^2}$,
$a_2'(-1)=-3$;
if $n\ge3$
$a_n(r)=(r+1)c_n(r)$ where
$c_n(r)=\dst\frac{r(r-1)}{(n+r)(n+r-1)(n+r+1)}$, so $a_n'(-1)=c_n(-1)=
\dst\frac{2}{ n(n-2)(n-1)}$ and
 $\dst{y_2=y_1\ln x+3-3x+2\sum_{n=2}^\infty
\frac{1}{ n(n^2-1)}x^n}$.


\item
$p_0(r)=(r-2)^2$;
$p_1(r)=-(r-5)(r-1)$;
$r_1=2$.


$a_n(r)=\dst \frac{n+r-6}{ n+r-2}
a_{n-1}(r)$;

 $a_n(r)=\dst\prod_{j=1}^m\frac{j+r-6}{ j+r-2}$,
$n\ge0$. Therefore,$a_1(2)=-3$, $a_2(2)=3$, $a_3(2)=-1$, and
$a_n(2)=0$ if $n\ge4$, so $y_1=x^2(1-x)^3$.


$a_1(r)=\dst\frac{r-5}{ r-1}$, $a_1'(r)=\dst\frac{4}{(r-1)^2}$,
$a_1'(2)=4$;

$a_2(r)=\dst\frac{(r-5)(r-4)}{ r(r-1)}$,
$a_2'(r)=\dst\frac{4(2r^2-10r+5)}{ r^2(r-1)^2}$,
$a_2'(2)=-7$;

$a_3(r)=\dst\frac{(r-5)(r-4)(r-3)}{ r(r-1)(r+1)}$,
$a_3'(r)=\dst\frac{12(r^4-8r^3+16r^2-5)}{ r^2(r-1)^2(r+1)^2}$,
$a_3'(2)=11/3$; if $n\ge4$, then
$a_n(r)=(r-2)c_n(r)$ where
$c_n(r)=\dst
\frac{(r-5)(r-4)(r-3)}{(n+r-5)(n+r-4)(n+r-3)(n+r-2)}$, so
$a_n'(2)=c_n(2)=-\dst\frac{6}{ n(n-2)(n^2-1)}$ and

$\dst{y_2=y_1\ln
x+x^3\left(4-7x+\frac{11}{3}x^2-6\sum_{n=3}^\infty\frac{1}{
n(n-2)(n^2-1)}x^n\right)}$.


\item
$p_0(r)=(3r-1)^2$;
$p_2(r)=7-3r$;
$r_1=1/3$.

$a_{2m}(r)=\dst\frac{6m+3r-13}{(6m+3r-1)^2}
a_{2m-2}(r)$;

 $a_{2m}(r)=\dst\prod_{j=1}^m\frac{6j+3r-13}{(6j+3r-1)^2}$,
$m\ge0$. Therefore,$a_2(1/3)=1/6$ and $a_{2m}(1/3)=0$ if $m\ge2$,
so $\dst{y_1=x^{1/3}\left(1-\frac{1}{6}x^2\right)}$.


$a_2(r)=\dst\frac{3r-7}{(3r+5)^2}$;
$a_2'(r)=\dst\frac{3(19-3r)}{(3r+5)^3}$; $a_2'(1/3)= 1/4$.
If $m\ge2$, then $a_{2m}(r)=(r-1/3)c_{2m}(r)$ where
$c_{2m}(r)=\dst\frac{3(3r-7)}{(6m+3r-7)(6m+3r-1)\prod_{j=1}^m(6j+3r-1)}$,
so
$a_{2m}'(1/3)=c_{2m}(1/3)=-\dst\frac{1}{12}\frac{1}{6^{m-1}(m-1)m\,m!}$,
and

 $\dst{y_2=y_1\ln x
+x^{7/3}\left(\frac{1}{4}-\frac{1}{12}\sum_{m=1}^\infty
\frac{1}{6^mm(m+1)(m+1)!}x^{2m}\right)}$.


\item
$p_0(r)=(2r+1)^2$;
$p_2(r)=7-2r$;
$r_1=-1/2$.

$a_{2m}(r)=\dst\frac{4m+2r-11}{(4m+2r+1)^2}
a_{2m-2}(r)$;

 $a_{2m}(r)=\dst\prod_{j=1}^m\frac{4j+2r-11}{(4j+2r+1)^2}$,
$m\ge0$. Therefore,$a_2(-1/2)=-1/2$, $a_4(-1/2)=1/32$, and
$a_{2m}(-1/2)=0$ if $m\ge3$, so
$\dst{y_1=x^{-1/2}\left(1-\frac{1}{2}x^2+\frac{1}{32}x^4\right)}$.



$a_2(r)=\dst\frac{2r-7}{(2r+5)^2}$,
$a_2'(r)=\dst\frac{2(19-2r)}{(2r+5)^3}$,
$a_2'(-1/2)=5/8$,


$a_4(r)=\dst\frac{(2r-7)(2r-3)}{(2r+5)^2(2r+9)^2}$,
$a_4'(r)=-\dst\frac{4(8r^3-60r^2-146r+519)}{(2r+5)^3(2r+9)^3}$,
$a_4'(-1/2)=-9/128$; if $m\ge3$, then $a_{2m}(r)=(r+1/2)c_{2m}(r)$
where

$c_{2m}(r)=\dst\frac{2(2r-7)(2r-3)}{(4m+2r-7)(4m+2r-3)(4m+2r+1)\prod_{j=1}^m
(4j+2r+1)}$, so
$a_{2m}'(-1/2)=c_{2m}(-1/2)=\dst\frac{1}{4^m(m-2)(m-1)m\,m!}$, and

$\dst{y_2=y_1\ln x+x^{3/2}\left(\frac{5}{8}-\frac{9}{128}x^2
+\sum_{m=2}^\infty\frac{1}{4^{m+1}(m-1)m(m+1)(m+1)!}x^{2m}\right)}$.

\item
\begin{alist}
\item If $p_0(r)=\alpha_0(r-r_1)^2$, then
(A) $a_n(r)=\dst\frac{(-1)^n}{\alpha_0^n}
\prod_{j=1}^n\frac{p_1(j+r-1)}{(j+r-r_1)^2}$.
Therefore,
$a_n(r_1)=\dst\frac{(-1)^n}{\alpha_0^n(n!)^2}\prod_{j=1}^np_1(j+r_1-1)$.
Theorem~7.6.2  implies  $Ly_1=0$.

\item
From (A),
$\dst\ln|a_n(r)|=-n\ln|\alpha_0|+\sum_{j=1}^n\left(\ln|p_1(j+r-1)|
-2\ln|j+r-r_1|\right)$, so
$a_n'(r)=\dst a_n(r)\sum_{j=1}^n\left(\frac{p_1'(j+r-1)}{
p_1(j+r-1)}-\frac{2}{ j+r-r_1}\right)$ and
$a_n'(r_1)=\dst a_n(r_1)\sum_{j=1}^n\left(\frac{p_1'(j+r_1-1)}{
p_1(j+r_1-1)}-\frac{2}{ j}\right)$.
Theorem~7.6.2  implies that $Ly_2=0$.

\item  Since $p_1(r)=\gamma_1$, $y_1$ and $y_2$ reduce to the
stated forms.  If $\gamma_1=0$, then $y_1=x^{r_1}$ and $y_2=x^{r_1}\ln
x$, which are solutions of the Euler equation $\alpha_0
x^2y''+\beta_0 xy'+\gamma_0 y$.
\end{alist}

\item
\begin{alist}
\item $Ly_1=p_0(r_1)x^{r_1}=0$. Now use the fact that
$p_0(j+r_1)=\alpha_0j^2$, so
$\prod_{j=1}^np_0(j+r_1)=\alpha_0^n(n!)^2$.

\item From Theorem~7.6.2, $y_2=\dst y_1\ln
x+x^{r_1}\sum_{n=1}^\infty a_n'(r_1)x^n$ is a second solution
of $Ly=0$. Since
$a_n(r)=\dst\frac{(-1)^n}{\alpha_0^n}\prod_{j=1}^n\frac{p_1(j+r-1)}{
(j+r-r_1)^2}$, (A)
$\ln|a_n(r)|=-n\ln|\alpha_0|+\dst\sum_{j=1}^n\ln|p_1(j+r-1)|-2\sum_{j=1}^n
\ln|j+r-r_1|$, provided that $p_1(j+r-1)$ and $j+r-r_1$
are nonzero for all positive integers $j$. Differentiating (A)
and then setting $r=r_1$ yields
$\dst\frac{a_n'(r_1)}{ a_n(r_1)}=\sum_{j=1}^n\frac{p_1'(j+r_1-1)}{
p_2(j+r_1-1)}-2\sum_{j=1}^n\frac{1}{ j}$, which implies the conclusion.

\item In this case $p_1(r)=\gamma_1$ and $p_1'(r)=0$, so
$a_n(r_1)=\dst\frac{(-1)^n}{(n!)^2}\left(\frac{\gamma_1}{\alpha_0}\right)^n$
and $\dst J_n=-2\sum_{j=1}^n\frac{1}{ j}$.
If $\gamma_1=0$, then $y_1=x^{r_1}$ and $y_2=x^{r_1}\ln x$,
while the differential equation
is an Euler equation with  indicial polynomial $\alpha_0(r-r_1^2)$.
See Theorem~7.4.3.
\end{alist}

\item
$p_0(r)=r^2$; $p_1(r)=1$; $r_1=0$.
$a_{2m}(r)=\dst-\frac{a_{2m-1}(r)}{(2m+r)^2}$, $m\ge1$;
$a_{2m}(r)=\dst\frac{(-1)^m}{\prod_{j=1}^m(2j+r)^2}$,  $m\ge0$.
Therefore,$a_{2m}(0)=\dst\frac{(-1)^m}{4^m(m!)^2}$, so
$\dst{y_1=\sum_{m=0}^\infty\frac{(-1)^m}{4^m(m!)^2}x^{2m}}$.

By logarithmic differentiation, $a_{2m}'(r)=-2a_{2m}(r)\dst
\sum_{j=1}^m\frac{1}{2m+r}$, so
 $a_{2m}'(0)=-a_{2m}(0)\dst\sum_{j=1}^m\frac{1}{ j}$ and
$\dst{y_2=y_1\ln x-
\sum_{m=1}^\infty\frac{(-1)^m}{4^m(m!)^2}
\left(\sum_{j=1}^m\frac{1}{ j}\right)x^{2m}}$.


\item
$p_0(r)=(2r-1)^2$;
$p_1(r)=(2r+1)^2$;
$p_2(r)=0$;
$\dst\frac{p_1(r-1)}{ p_0(r)}=1=\frac{\alpha_1}{\alpha_0}$;
$\dst\frac{p_2(r-2)}{ p_0(r)}=0=\frac{\alpha_2}{\alpha_0}$;
$y_1=\dst\frac{x^{1/2}}{1+x}$; $y_2=\dst\frac{x^{1/2}\ln x}{1+x}$.

\item
$p_0(r)=2(r-1)^2$;
$p_1(r)=0$;
$p_2(r)=-(r+1)^2$;
$\dst\frac{p_1(r-1)}{ p_0(r)}=0=\frac{\alpha_1}{\alpha_0}$;
$\dst\frac{p_2(r-2)}{ p_0(r)}=-1/2=\frac{\alpha_2}{\alpha_0}$;
$y_1=\dst\frac{x}{2-x^2}$; $y_2=\dst\frac{x\ln x}{2-x^2}$.


\item
$p_0(r)=4(r-1)^2$;
$p_1(r)=3r^2$;
$p_2(r)=0$;
$\dst\frac{p_1(r-1)}{ p_0(r)}=3/4=\frac{\alpha_1}{\alpha_0}$;
$\dst\frac{p_2(r-2)}{ p_0(r)}=0=\frac{\alpha_2}{\alpha_0}$;
$y_1=\dst\frac{x}{4+3x}$; $y_2=\dst\frac{x\ln x}{4+3x}$.

\item
$p_0(r)=(r-1)^2$;
$p_1(r)=-2r^2$;
$p_2(r)=(r+1)^2$;
$\dst\frac{p_1(r-1)}{ p_0(r)}=-2=\frac{\alpha_1}{\alpha_0}$;
$\dst\frac{p_2(r-2)}{ p_0(r)}=1=\frac{\alpha_2}{\alpha_0}$;
$y_1=\dst\frac{x}{(1-x)^2}$; $y_2=\dst\frac{x\ln x}{(1-x)^2}$.


\item
See the proofs of Theorems~7.6.1 and 7.6.2.
\end{sectionSolutions}

\section{The Method of Frobenius III}

\begin{sectionSolutions}
\item
$p_0(r)=r(r-1)$;
$p_1(r)=1$;
$r_1=1$; $r_2=0$; $k=r_1-r_2=1$;

$a_n(r)=\dst-\frac{1}{(n+r)(n+r-1)}
a_{n-1}(r)$;


 $a_n(r)=\dst\frac{(-1)^n}{\prod_{j=1}^n(j+r)(j+r-1)}$;

$a_n(1)=\dst\frac{(-1)^n}{ n!(n+1)!}$;

$\dst{y_1=x\sum_{n=0}^\infty\frac{(-1)^n}{ n!(n+1)!}x^n}$;

$z=1$; $C=-p_1(0)a_0(0)=-1$.

By logarithmic differentiation,

$a_n'(r)=-\dst a_n(r)\sum_{j=1}^n\frac{2n+2r-1}{(n+r)(n+r-1)}$;

$a_n'(1)=-\dst a_n(1)\sum_{j=1}^n\frac{2j+1}{ j(j+1)}$;

$\dst{y_2=1-y_1\ln x
+x\sum_{n=1}^\infty\frac{(-1)^n}{
n!(n+1)!}\left(\sum_{j=1}^n\frac{2j+1}{ j(j+1)}\right)x^n}$.


\item
$p_0(r)=r(r-1)$;
$p_1(r)=r+1$;
$r_1=1$; $r_2=0$; $k=r_1-r_2=1$;
$a_n(r)=\dst-\frac{a_{n-1}(r)}{ n+r-1}$;
 $a_n(r)=\dst\frac{(-1)^n}{\prod_{j=1}^n(j+r-1)}$;
$a_n(1)=\dst\frac{(-1)^n}{ n!}$;
$\dst{y_1=x\sum_{n=0}^\infty\frac{(-1)^n}{ n!}x^n=xe^{-x}}$;
$z=1$;
$C=-p_1(0)a_0(0)=-1$.
By logarithmic differentiation,
$a_n'(r)=-\dst a_n(r)\sum_{j=1}^n\frac{1}{ j+r-1}$;
$a_n'(1)=-\dst a_n(1)\sum_{j=1}^n\frac{1}{ j}$;
$\dst{y_2=1-y_1\ln x+x\sum_{n=1}^\infty\frac{(-1)^n}{
n!}\left(\sum_{j=1}^n\frac{1}{ j}\right)x^n}$.


\item
$p_0(r)=(r-1)(r+2)$;
$p_1(r)=r+3$;
$r_1=1$; $r_2=-2$; $k=r_1-r_2=3$.
$a_n(r)=\dst-\frac{1}{ n+r-1}
a_{n-1}(r)$;
 $a_n(r)=\dst\frac{(-1)^n}{\prod_{j=1}^n(j+r-1)}$;
$a_n(1)=\dst\frac{(-1)^n}{ n!}$;
$\dst{y_1=x\sum_{n=0}^\infty\frac{(-1)^n}{ n!}x^n=xe^{-x}}$;
$z=x^{-2}\left(\dst1+\frac{1}{2}x+\frac{1}{2}x^2\right)$;
$C=-\dst\frac{p_1(0)}{ 3}a_2(-2)=-1/2$.
By logarithmic differentiation,
$a_n'(r)=-\dst a_n(r)\sum_{j=1}^n\frac{1}{ j+r-1}$;
$a_n'(1)=\dst a_n(1)\sum_{j=1}^n\frac{1}{ j}$;
 $\dst{y_2=x^{-2}\left(1+\frac{1}{2}x+\frac{1}{2}x^2\right)-\frac{1}{2}
\left(y_1\ln x-x\sum_{n=1}^\infty\frac{(-1)^n}{
n!}\left(\sum_{j=1}^n\frac{1}{ j}\right)x^n\right)}$;


\item
$p_0(r)=(r+2)(r+7)$;
$p_1(r)=1$;
$r_1=-2$; $r_2=-7$; $k=r_1-r_2=5$;
$a_n(r)=\dst-\frac{a_{n-1}(r)}{(n+r+2)(n+r+7)}$;
$a_n(r)=\dst\frac{(-1)^n}{\prod_{j=1}^n(j+r+2)(j+r+7)}$;
$a_n(-2)=120\dst\prod_{j=1}^n\frac{(-1)^n}{ n!(n+5)!}$;
$\dst{y_1=\frac{120}{ x^2}\sum_{n=0}^\infty\frac{(-1)^n}{
n!(n+5)!}x^n}$;

$z=\dst
x^{-7}\left(1+\frac{1}{4}x+\frac{1}{24}x^2+\frac{1}{144}x^3
+\frac{1}{576}x^4\right)$;
$C=-\dst\frac{p_1(-3)}{5}a_4(-7)=-1/2880$.
By logarithmic differentiation,
$a_n'(r)=-\dst a_n(r)\sum_{j=1}^n-\frac{2j+2r+9}{(j+r+2)(j+r+7)}$;
$a_n'(-2)=-\dst a_n(-2)\sum_{j=1}^n\frac{2j+5}{ j(j+5)}$;
\begin{multline*}
y_2=x^{-7}\left(1+\frac{1}{4}x+\frac{1}{24}x^2+\frac{1}{144}x^3
+\frac{1}{576}x^4\right)\\
-\frac{1}{2880}\left(y_1\ln x-\frac{120}{ x^2}
\sum_{n=1}^\infty\frac{(-1)^n}{
n!(n+5)!}\left(\sum_{j=1}^n\frac{2j+5}{ j(j+5)}\right)x^n\right).
\end{multline*}


\item
$p_0(r)=r(r-4)$;
$p_1(r)=(r-6)(r-5)$;
$r_1=4$; $r_2=0$; $k=r_1-r_2=4$;
$a_n(r)=\dst
-\frac{(n+r-7)(n+r-6)}{(n+r)(n+r-4)}
a_{n-1}(r)$;
 $a_n(r)=(-1)^n\dst\prod_{j=1}^n\frac{(j+r-7)(j+r-6)}{(j+r)(j+r-4)}$.
Setting $r=4$ yields
$\dst{y_1=x^4\left(1-\frac{2}{5}x\right)}$.
$z=\dst1+10x+50x^2+200x^3$;
$C=-\dst\frac{p_1(3)}{ 4}a_3(0)=-300$.
$a_1(r)=-\dst\frac{(r-6)(r-5)}{(r-3)(r+1)}$;
$a_1'(r)=-\dst\frac{3(3r^2-22r+31)}{(r-3)^2(r+1)^2}$;
$a_1'(4)=27/25$.
$a_2(r)=(r-4)c_2(r)$, with
$c_2(r)=\dst\frac{(r-6)(r-5)^2}{(r-3)(r-2)(r+1)(r+2)}$, so
$a_2'(4)=c_2(4)=-1/30$.
If $n\ge3$, then $a_n(r)=(r-4)^2b_n(r)$ where $b_n'(4)$ exists, so
$a_n'(4)=0$  and
$\dst{y_2=1+10x+50x^2+200x^3-300\left(y_1\ln
x+\frac{27}{25}x^5-\frac{1}{30}x^6\right)}$.


\item
$p_0(r)=(r-2)(r+2)$;
$p_1(r)=-2r-1$;
$r_1=2$; $r_2=-2$; $k=r_1-r_2=4$;
$a_n(r)=\dst\frac{2j+2r-1}{(j+r-2)(j+r+2)}
a_{n-1}(r)$;
 $a_n(r)=\dst\prod_{j=1}^n\frac{2n+2r-1}{(n+r-2)(n+r+2)}$;
$a_n(2)=\frac{1}{ n!}\left(\prod_{j=1}^n\frac{2j+3}{ j+4}\right)$;
$\dst{y_1=x^2\sum_{n=0}^\infty\frac{1}{ n!}\left(\prod_{j=1}^n\frac{2j+3}{
j+4}\right)x^n}$;
$z=\dst x^{-2}\left(1+x+\frac{1}{4}x^2-\frac{1}{12}x^3\right)$;
$C=-\dst\frac{p_1(1)}{ 4}a_3(-2)=-1/16$.
By logarithmic differentiation,
$a_n'(r)=-2\dst a_n(r)\sum_{j=1}^n
\frac{j^2+j(2r-1)+r^2-r+4}{(j+r-2)(j+r+2)(2j+2r-1)}$;
$a_n'(2)=-2\dst a_n(2)\sum_{j=1}^n\frac{(j^2+3j+6)}{
j(j+4)(2j+3)}$; $y_2=$
\[
x^{-2}\left(1+x+\frac{1}{4}x^2-\frac{1}{12}x^3\right)-\frac{1}{16}
y_1\ln x+ \frac{x^2}{8}\sum_{n=1}^\infty\frac{1}{
n!}\left(\prod_{j=1}^n\frac{2j+3}{ j+4}
\right)\left(\sum_{j=1}^n\frac{(j^2+3j+6)}{ j(j+4)(2j+3)}\right)x^n.
\]


\item
$p_0(r)=(r+1)(r+7)$;
$p_1(r)=(r+5)(2r+1)$;
$r_1=-1$; $r_2=-7$; $k=r_1-r_2=6$;
$a_n(r)=\dst -\frac{(n+r+4)(2n+2r-1)}{(n+r+1)(n+r+7)} a_{n-1}(r)$;
 $a_n(r)=(-1)^n\dst\prod_{j=1}^n
\frac{(j+r+4)(2j+2r-1)}{(j+r+1)(j+r+7)}$;
$a_n(-1)=
\dst\frac{(-1)^n}{
n!}\left(\prod_{j=1}^n\frac{(j+3)(2j-3)}{ j+6}\right)$;
$y_1=\dst{\frac{1}{ x}\sum_{n=0}^\infty\frac{(-1)^n}{
n!}\left(\prod_{j=1}^n\frac{(j+3)(2j-3)}{ j+6}\right)x^n}$;
$z=\dst{x^{-7}\left(1+\frac{26}{5}x+\frac{143}{20}x^2\right)}$;
$C=-\dst\frac{p_1(-2)}{ 6}a_5(-7)=0$;
$y_2=\dst{x^{-7}\left(1+\frac{26}{5}x+\frac{143}{20}x^2\right)}$.


\item
$p_0(r)=(3r-10)(3r+2)$;
$p_1(r)=r(3r-4)$;
$r_1=10/3$; $r_2=-2/3$; $k=r_1-r_2=4$;
$a_n(r)=\dst-\frac{(n+r-1)(3n+3r-7)}{(3n+3r-10)(3n+3r+2)}
a_{n-1}(r)$;
$a_n(r)=(-1)^n\dst\prod_{j=1}^n\frac{(j+r-1)(3j+3r-7)}{(3j+3r-10)(3j+3r+2)}$;
$a_n(10/3)=\dst\frac{(-1)^n(n+1)}{9^n}\left(\prod_{j=1}^n\frac{3j+7}{
j+4}\right)$;
$y_1=\dst{x^{10/3}\sum_{n=0}^\infty\frac{(-1)^n(n+1)}{9^n}
\left(\prod_{j=1}^n\frac{3j+7}{ j+4}\right)x^n}$;
$z=\dst{x^{-2/3}\left(1+\frac{4}{27}x-\frac{1}{243}x^2\right)}$;
$C=-\dst\frac{p_1(7/3)}{36}a_3(-2/3)=0$;
$y_2=\dst{x^{-2/3}\left(1+\frac{4}{27}x-\frac{1}{243}x^2\right)}$.



\item
$p_0(r)=(r-3)(r+2)$;
$p_1(r)=(r+1)^2$;
$r_1=3$; $r_2=-2$; $k=r_1-r_2=5$;
$a_n(r)=\dst-\frac{(n+r)^2}{(n+r-3)(n+r+2)} a_{n-1}(r)$;
 $a_n(r)=(-1)^n\dst\prod_{j=1}^n\frac{(j+r)^2}{(j+r-3)(j+r+2)}$;
$a_n(3)= \frac{(-1)^n}{ n!}\left(\prod_{j=1}^n\frac{(j+3)^2}{
j+5}\right)$;
$y_1=\dst{x^3\sum_{n=0}^\infty\frac{(-1)^n}{
n!}\left(\prod_{j=1}^n\frac{(j+3)^2}{
j+5}\right)x^n}$;
$z=\dst{x^{-2}\left(1+\frac{1}{4}x\right)}$;
$C=-\dst\frac{p_1(2)}{5}a_4(-2)=0$;
$y_2=\dst{x^{-2}\left(1+\frac{1}{4}x\right)}$.



\item
$p_0(r)=(r-6)(r-1)$;
$p_1(r)=(r-8)(r-4)$;
$r_1=6$; $r_2=1$; $k=r_1-r_2=5$;
$a_n(r)=\dst-\frac{(n+r-9)(n+r-5)}{(n+r-6)(n+r-1)}
a_{n-1}(r)$;
$y_1=\dst{x^6\left(1+\frac{2}{3}x+\frac{1}{7}x^2\right)}$;
$z=\dst{x\left(1+\frac{21}{4}x+\frac{21}{2}x^2+\frac{35}{4}x^3\right)}$;
$C=-\dst\frac{p_1(5)}{6}a_5(1)=0$;
$y_2=\dst{x\left(1+\frac{21}{4}x+\frac{21}{2}x^2+\frac{35}{4}x^3\right)}$.


\item
$p_0(r)=r(r-10)$;
$p_1(r)=2(r-6)(r+1)$;
$r_1=10$; $r_2=0$; $k=r_1-r_2=10$;
$a_n(r)=\dst -\frac{2(n+r-7)}{ n+r-10}a_{n-1}(r)$;
 $a_n(r)=(-2)^n\dst\frac{(n+r-9)(n+r-8)(n+r-7)}{(r-9)(r-8)(r-7)}$;
$a_n(10)=\dst\frac{(-1)^n2^n(n+1)(n+2)(n+3)}{6}$;
$y_1=\dst{\frac{x^{10}}{6}\sum_{n=0}^\infty(-1)^n2^n(n+1)(n+2)(n+3)x^n}$;
$z=\dst\left(1-\frac{4}{3}x+\frac{5}{3}x^2-\frac{40}{21}x^3
+\frac{40}{21}x^4-\frac{32}{21}x^5+\frac{16}{21}x^6\right)$;
$C=-\dst\frac{p_1(9)}{10}a_9(0)=0$;
$y_2=\dst\left(1-\frac{4}{3}x+\frac{5}{3}x^2-\frac{40}{21}x^3
+\frac{40}{21}x^4-\frac{32}{21}x^5+\frac{16}{21}x^6\right)$.


\textbf{Note: in the solutions to Exercises~7.7.23--7.7.40,
$z=x^{r_2}\sum_{m=0}^{k-1}a_{2m}(r_2)x^{2m}$.}




\item
$p_0(r)=(r-6)(r-2)$;
$p_2(r)=r$;
$r_1=6$; $r_2=2$; $k=(r_1-r_2)/2=2$;
$a_{2m}(r)=\dst-\frac{a_{2m-2}(r)}{2m+r-6}
$;
 $a_{2m}(r)=\dst\frac{(-1)^m}{\prod_{j=1}^m(2j+r-6)}$;
$a_{2m}(6)=\dst\frac{(-1)^m}{2^mm!}$;
$\dst{y_1=x^6\sum_{m=0}^\infty\frac{(-1)^m}{2^mm!}x^{2m}=x^6e^{-x^2/2}}$;
$z=\dst x^2\left(1+\frac{1}{2}x^2\right)$;
$C=-\dst\frac{p_2(4)}{4}a_2(2)=-1/2$.
By logarithmic differentiation,
$a_{2m}'(r)=-\dst a_{2m}(r)\sum_{j=1}^m\frac{1}{2j+r-6}$;
$a_{2m}'(6)=-\dst a_{2m}(6)\sum_{j=1}^m\frac{1}{2j}$;
$\dst{y_2=x^2\left(1+\frac{1}{2}x^2\right)-\frac{1}{2}y_1\ln
x+\frac{x^6}{4}\sum_{m=1}^\infty\frac{(-1)^m}{2^mm!}\left(\sum_{j=1}^m
\frac{1}{ j}\right)x^{2m}}$.


\item
$p_0(r)=(r-1)(r+1)$;
$p_2(r)=2r+10$;
$r_1=1$; $r_2=-1$; $k=(r_1-r_2)/2=1$;
$a_{2m}(r)=\dst-\frac{2(2m+r+3)}{(2m+r-1)(2m+r+1)}a_{2m-2}(r)$;
 $a_{2m}(r)=(-2)^m\dst\prod_{j=1}^m\frac{2j+r+3}{(2j+r-1)(2j+r+1)}$;
$a_{2m}(1)=\dst\frac{(-1)^m(m+2)}{2\, m!}$;
$\dst{y_1=\frac{x}{2}\sum_{m=0}^\infty\frac{(-1)^m(m+2)}{ m!}x^{2m}}$;
$z=x^{-1}$;
$C=-\dst\frac{p_2(-1)}{ 2}a_0(-1)=-4$.
By logarithmic differentiation,

$a_{2m}'(r)=-\dst a_{2m}(r)\sum_{j=1}^m
\frac{(4j^2+4j(r+3)+r^2+6r+1)}{(2j+r-1)(2j+r+1)(2j+r+3)}$;

$a_{2m}'(1)=-\dst a_{2m}(1)\sum_{j=1}^m\frac{j^2+4j+2}{2j(j+1)(j+2)}$;

$\dst{y_2=x^{-1}-4y_1\ln x+x\sum_{m=1}^\infty\frac{(-1)^m(m+2)}{
m!}\left(\sum_{j=1}^m\frac{j^2+4j+2}{ j(j+1)(j+2)} \right)x^{2m}}$.




\item
$p_0(r)=(2r+1)(2r+5)$;
$p_2(r)=2r+3$;
$r_1=-1/2$; $r_2=-5/2$; $k=(r_1-r_2)/2=1$;
$a_{2m}(r)=\dst-\frac{(4m+2r-1)}{(4m+2r+1)(4m+2r+5)}a_{2m-2}(r)$;
$a_{2m}(r)=(-1)^m\dst\prod_{j=1}^m\frac{(4j+2r-1)}{(4j+2r+1)(4j+2r+5)}$;
$a_{2m}(-1/2)=\dst\frac{(-1)^m\prod_{j=1}^m(2j-1)}{8^mm!(m+1)!}$;
$\dst{y_1=x^{-1/2}\sum_{m=0}^\infty
\frac{(-1)^m\prod_{j=1}^m(2j-1)}{8^mm!(m+1)!}
x^{2m}}$;
$z=x^{-5/2}$;
$C=-\dst\frac{p_2(-5/2)}{8}a_0(-5/2)=1/4$.
By logarithmic differentiation,

$a_{2m}'(r)=-2\dst a_{2m}(r)\sum_{j=1}^m
\frac{(16j^2+8j(2r-1)+4r^2-4r-11)}{(4j+2r-1)(4j+2r+1)(4j+2r+5)}$;
$a_{2m}'(-1/2)=-\dst a_{2m}(-1/2)\sum_{j=1}^m\frac{2j^2-2j-1}{
2j(j+1)(2j-1)}$;

$\dst{y_2=x^{-5/2}+\frac{1}{4}y_1\ln x
-x^{-1/2}\sum_{m=1}^\infty\frac{(-1)^m\prod_{j=1}^m(2j-1)}{8^{m+1}m!(m+1)!}
\left(\sum_{j=1}^m\frac{2j^2-2j-1}{ j(j+1)(2j-1)}\right) x^{2m}}$.


\item
$p_0(r)=(r-2)(r+2)$;
$p_2(r)=-2(r+4)$;
$r_1=2$; $r_2=-2$; $k=(r_1-r_2)/2=2$;
$a_{2m}(r)=\dst\frac{2}{2m+r-2}a_{2m-2}(r)$;
 $a_{2m}(r)=\dst\frac{2^m}{\prod_{j=1}^m(2j+r-2)}$;
$a_{2m}(2)=\dst\frac{1}{ m!}$;
$\dst{y_1=x^2\sum_{m=0}^\infty\frac{1}{ m!}x^{2m}=x^2e^{x^2}}$;
$z=x^{-2}(1-x^2)$;
$C=-\dst\frac{p_2(0)}{4}a_2(-2)=-2$.
By logarithmic differentiation,
$a_{2m}'(r)=-\dst a_{2m}(r)\sum_{j=1}^m\frac{1}{2j+r-2} $;
$a_{2m}'(2)=-\dst a_{2m}(2)\sum_{j=1}^m\frac{1}{2j}$;
$\dst{y_2=x^{-2}(1-x^2)-2y_1\ln x+x^2\sum_{m=1}^\infty\frac{1}{
m!}\left(\sum_{j=1}^m\frac{1}{ j}\right)x^{2m}}$.


\item
$p_0(r)=(3r-13)(3r-1)$;
$p_2(r)=2(5-3r)$;
$r_1=13/3$; $r_2=1/3$; $k=(r_1-r_2)/2=2$;
$a_{2m}(r)=\dst\frac{2(6m+3r-11)}{(6m+3r-13)(6m+3r-1)}a_{2m-2}(r)$;
$a_{2m}(r)=2^m\dst\prod_{j=1}^m\frac{(6j+3r-11)}{(6j+3r-13)(6j+3r-1)}$;
$a_{2m}(13/3)=\dst\frac{\prod_{j=1}^m(3j+1) }{9^mm!(m+2)!}$;
$\dst{y_1=2x^{13/3}\sum_{m=0}^\infty\frac{\prod_{j=1}^m(3j+1)
}{9^mm!(m+2)!}x^{2m}}$;
$z=\dst x^{1/3}\left(1+\frac{2}{9}x^2\right)$;
$C=-\dst\frac{p_2(7/3)}{36}a_2(1/3)=2/81$.
By logarithmic differentiation,
$a_{2m}'(r)=-9\dst a_{2m}(r)\sum_{j=1}^m
\frac{(12j^2+4j(3r-11)+3r^2-22r+47)}{(6j+3r-13)(6j+3r-11)(6j+3r-1)}$;
$a_{2m}'(13/3)=-\dst
a_{2m}(13/3)\sum_{j=1}^m\frac{3j^2+2j+2}{2j(j+2)(3j+1)}$;

$\dst{y_2=x^{1/3}\left(1+\frac{2}{9}x^2\right)+\frac{2}{81}\left(
y_1\ln
x-x^{13/3}\sum_{m=0}^\infty\frac{\prod_{j=1}^m(3j+1)}{9^mm!(m+2)!}
\left(\sum_{j=1}^m\frac{3j^2+2j+2}{ j(j+2)(3j+1)}\right)x^{2m}\right)}$.


\item
$p_0(r)=(r-2)(r+2)$;
$p_2(r)=-3(r-4)$;
$r_1=2$; $r_2=-2$; $k=(r_1-r_2)/2=2$;
$a_{2m}(r)=\dst\frac{3(2m+r-6)}{(2m+r-2)(2m+r+2)}
a_{2m-2}(r)$;
$\dst{y_1=x^2\left(1-\frac{1}{2}x^2\right)}$;
$z=\dst x^{-2}\left(1+\frac{9}{2}x^2\right)$;
$C=-\dst\frac{p_2(0)}{4}a_2(-2)=-27/2$;
$a_2(r)=\dst\frac{3(r-4)}{ r(r+4)}$,
$a_2'(r)=-\dst\frac{3(r^2-8r-16)}{ r^2(r+4)^2}$,
$a_2'(2)=7/12$. If $m\ge2$, then $a_{2m}(r)=(r-2)c_{2m}(r)$ where
$c_{2m}(r)=\dst\frac{3^m(r-4)}{(2m+r-4)(2m+r-2)\prod_{j=1}^m(2j+r+2)}$,
so $a_{2m}'(2)=c_{2m}(2)=\dst -\frac{\left(\frac{3}{2}\right)^m }{
m(m-1)(m+2)!}$;
$\dst{y_2=x^{-2}\left(1+\frac{9}{2}x^2\right)-\frac{27}{2}\left(
y_1\ln x+\frac{7}{12}x^4-x^2\sum_{m=2}^\infty\frac{\left(\frac{3}{2}\right)^m
}{ m(m-1)(m+2)!}x^{2m}\right)}$.



\item
$p_0(r)=(2r-5)(2r+7)$;
$p_2(r)=(2r-1)^2$;
$r_1=5/2$; $r_2=-7/2$; $k=(r_1-r_2)/2=3$;
$a_{2m}(r)=\dst-\frac{4m+2r-5}{4m+2r+7}a_{2m-2}(r)$;
 $a_{2m}(r)=\dst
\frac{(2r-1)(2r+3)(2r+7)}{(4m+2r-1)(4m+2r+3)(4m+2r+7)}$;
$a_{2m}(5/2)=\dst\frac{(-1)^m}{(m+1)(m+2)(m+3)}$;
$y_1=\dst{x^{5/2}\sum_{m=0}^\infty\frac{(-1)^m}{(m+1)(m+2)(m+3)}x^{2m}}$;
$z=\dst{x^{-7/2}(1+x^2)^2}$
$C=-\dst\frac{p_2(1/2)}{24}a_4(-7/2)=0$;
$y_2=\dst{x^{-7/2}(1+x^2)^2}$.


\item
$p_0(r)=(r-3)(r+7)$;
$p_2(r)=r(r+1)$;
$r_1=3$; $r_2=-7$; $k=(r_1-r_2)/2=5$;
$a_{2m}(r)=\dst-\frac{(2m+r-2)(2m+r-1)}{(2m+r-3)(2m+r+7)}
a_{2m-2}(r)$;
$a_{2m}(r)=(-1)^m\dst\prod_{j=1}^m\frac{(2j+r-2)(2j+r-1)}{(2j+r-3)(2j+r+7)}$;
$a_{2m}(3)=\dst (-1)^m\frac{m+1}{2^m}\left(\prod_{j=1}^m\frac{2j+1}{
j+5}\right)$;
$y_1=\dst{x^3\sum_{m=0}^\infty(-1)^m\frac{m+1}{2^m}\left(\prod_{j=1}^m\frac{2j+1}{
j+5}\right)x^{2m}}$;
$z=\dst{x^{-7}\left(1+\frac{21}{8}x^2+\frac{35}{16}x^4+\frac{35}{64}x^6\right)}$
$C=-\dst\frac{p_2(1)}{10}a_8(-7)=0$;
$y_2=\dst{x^{-7}\left(1+\frac{21}{8}x^2+\frac{35}{16}x^4+\frac{35}{64}x^6\right)}$.


\item
$p_0(r)=(2r-3)(2r+5)$;
$p_2(r)=(2r-1)(2r+1)$;
$r_1=3/2$; $r_2=-5/2$; $k=(r_1-r_2)/2=2$;
$a_{2m}(r)=\dst-\frac{4m+2r-5}{4m+2r+5}
a_{2m-2}(r)$;
 $a_{2m}(r)=(-1)^m\dst\prod_{j=1}^m\frac{4j+2r-5}{4j+2r+5}$;
$a_{2m}(3/2)=\dst\frac{(-1)^m \prod_{j=1}^m(2j-1)}{2^{m-1}(m+2)!}$;
$y_1=\dst{x^{3/2}\sum_{m=0}^\infty\frac{(-1)^m
\prod_{j=1}^m(2j-1)}{2^{m-1}(m+2)!}x^{2m}}$;
$z=\dst{x^{-5/2}\left(1+\frac{3}{2}x^2\right)}$
$C=-\dst\frac{p_2(-1/2)}{16}a_2(-5/2)=0$;
$y_2=\dst{x^{-5/2}\left(1+\frac{3}{2}x^2\right)}$.


\item $p_0(r)=r^2-\nu^2$; $p_2(r)=1$; $r_1=\nu$;
$r_2=-\nu$; $k=(r_1-r_2)/2=\nu$;
$a_{2m}(r)=-\dst\frac{a_{2m-2}(r)}{(2m+r+\nu)(2m+r-\nu)}$;
$a_{2m}(r)=\dst\frac{(-1)^m}{\prod_{j=1}^m(2j+r+\nu)(2j+r-\nu)}$;
$a_{2m}(\nu)=\dst\frac{(-1)^m}{4^mm!\prod_{j=1}^m(j+\nu)}$;
$a_{2m}(-\nu)=\dst\frac{(-1)^m}{4^mm!\prod_{j=1}^m(j-\nu)}$,
$j=0,\dots,\nu-1$; $\dst{y_1=x^\nu\sum_{m=0}^\infty\frac{(-1)^m}{4^mm!
\prod_{j=1}^m(j+\nu)}x^{2m}}$; $\dst
z=x^{-\nu}\sum_{m=0}^{\nu-1}\frac{(-1)^m}{4^mm!
\prod_{j=1}^m(j-\nu)}x^{2m}$; $C=-\dst\frac{p_2(\nu-2)}{
2\nu}a_{2\nu-2}(-\nu)= -\frac{a_{2\nu-2}(-\nu)}{ 2\nu}=
-\frac{2}{4^\nu\nu!(\nu-1)!}$. By logarithmic differentiation,
$a_{2m}'(r)=-2\dst
a_{2m}(r)\sum_{j=1}^m\frac{2j+\nu}{(2j+r+\nu)(2j+r-\nu)}$;
$a_{2m}'(\nu)=-\dst a_{2m}(\nu)\sum_{j=1}^m\frac{2j+\nu}{2j(j+\nu)}$;
\begin{multline*}
y_2=x^{-\nu}\sum_{m=0}^{\nu-1}\frac{(-1)^m}{4^mm!\prod_{j=1}^m(j-\nu)}x^{2m}\\
-\frac{2}{4^\nu\nu!(\nu-1)!}\left(y_1 \ln x-\frac{x^\nu}{2}
\sum_{m=1}^\infty\frac{(-1)^m}{4^mm!\prod_{j=1}^m(j+\nu)}
\left(\sum_{j=1}^m\frac{2j+\nu}{ j(j+\nu)}\right) x^{2m}\right).
\end{multline*}



\item
Since
$a_n(r_2)=-\dst\frac{p_1(n+r_2-1)}{ p_0(n+r_2)}a_{n-1}(r_2)$,
$1\le n\le k-1$, $a_{k-1}(r_2)=(-1)^{k-1}\dst\prod_{j=1}^{k-1}
\frac{p_1(r_2+j-1)}{ p_0(r_2+j)}$. But $C=-\dst\frac{p_1(r_1-1)}{
k\alpha_0}a_{k-1}(r_2)=-\frac{p_1(r_2+k-1)}{ k\alpha_0}a_{k-1}(r_2)
=(-)^k\frac{\prod_{j=1}^kp_1(r_2+j-1)}{ k\alpha_0\prod_{j=1}^{k-1}
p_0(r_2+j)}=0$ if and only if $\dst\prod_{j=1}^kp_1(r_2+j-1)=0$.



\item Since $p_1(r)=\gamma_1$,
$a_n(r)=\dst-\frac{\gamma_1}{\alpha_0(n+r-r_1)(n+r-r_2)}a_{n-1}(r)$ and
(A) $a_n(r)=\dst(-1)^n\left(\frac{\gamma_1}{\alpha_0}\right)^n
\frac{1}{\prod_{j=1}^n(j+r-r_1)(j+r-r_2)}$. Therefore,
$a_n(r_1)=\dst\frac{(-1)^n}{ n!}\left(\frac{\gamma_1}{\alpha_0}\right)^n
\frac{1}{\prod_{j=1}^n(j+k)}$ for $n\ge0$ (so $Ly_1=0$) and
$a_n(r_2)=\dst\frac{(-1)^n}{ n!}\left(\frac{\gamma_1}{\alpha_0}\right)^n
\frac{1}{\prod_{j=1}^n(j-k)}$ for $n=0,\dots,k-1$. $Ly_2=0$ if $y_2=\dst
x^{r_2}\sum_{n=0}^{k-1} a_n(r_2)x^n+C\left(y_1\ln x+
x^{r_1}\sum_{n=1}^\infty a_n'(r_1)x^n\right)$ if
$C=\dst-\frac{\gamma_1}{ k\alpha_0}a_{k-1}(r_2)=-\frac{\gamma_1}{
k\alpha_0}\frac{(-1)^{k-1}}{(k-1)!}
\left(\frac{\gamma_1}{\alpha_0}\right)^{k-1}\frac{(-1)^{k-1}}{(k-1)!}
=-\frac{1}{ k!(k-1)!} \left(\frac{\gamma_1}{\alpha_0}\right)^k$. From (A),
$\ln|a_n(r)|=\dst-n\ln\left|\frac{\gamma_1}{\alpha_0}\right|
-\sum_{j=1}^n(\left(\ln|j+r-r_1|+\ln|j+r-r_2|\right)$, so
$a_n'(r)=\dst-a_n(r)\sum_{j=1}^n\left(\frac{1}{ j+r-r_1}+ \frac{1}{
j+r-r_2}\right)$ and $a_n'(r_1)=\dst-a_n(r_1)\sum_{j=1}^n \frac{2j+k}{
j(j+k)}$.



\item
\begin{alist}
\item From Exercise~7.6.66\part{a} of
Section~7.6, $\dst L\left(\frac{\partial y}{ \partial
r}(x,r)\right)=p'_0(r)x^r+x^rp_0(r)\ln x$. Setting $r=r_1$ yields
$\dst L\left(y_1\ln x+x^{r_1}\sum_{n=1}^\infty
a_n'(r_1)\right)=p_0'(r_1)x^{r_1}$. Since
$p_0'(r)=\alpha_0(2r-r_1-r_2)$, $p_0'(r_1)~=~k\alpha_0$.

\item From Exercise~7.5.57 of Section~7.5, $\dst
L\left(x^{r_2}\sum_{n=0}^\infty a_n(r_2)x^n\right)=x^{r_2}
\sum_{n=0}^\infty b_nx^n$, where $b_0=p_0(r_2)=0$ and
$b_n=\dst\sum_{j=0}^np_j(n+r_2-j)a_{n-j}(r_2)$ if $n\ge1$. From the
definition of $\{a_n(r_2)\}$, $b_n=0$ if $n\ne k$, while
$b_k=\dst\sum_{j=0}^kp_k(k+r_2-j)a_{k-j}(r_2)=
\sum_{j=1}^kp_j(r_1-j)a_{k-j}(r_2)$.

\stepcounter{enumXi}
\item Let $\{\tilde a_n(r_2)\}$ be the coefficients that would
obtained if $\tilde a_k(r_2)=0$. Then $a_n(r_2)=\tilde a_n(r_2)$ if
$n=0,\dots,k-1$, and (A) $a_n(r_2)-\tilde a_n(r_2)=-\dst\frac{1}{
p_0(n+r_2)}\sum_{j=0}^{n-k} p_j(n+r_2-j)(a_{n-j}(r_2)-\tilde
a_{n-j}{r_2})$ if $n>k$. Now let $c_m=a_{k+m}(r_2)-\tilde
a_{k+m}(r_2)$. Setting $n=m+k$ in (A) and recalling the $k+r_2=r_1$
yields (B) $c_m=\dst-\frac{1}{ p_0(m+r_1)}\sum_{j=0}^m
p_j(m+r_1-j)c_{m-j}$. Since $c_k=a_k(r_2)$, (B) implies that
$c_m=a_k(r_2)a_m(r_1)$ for all $m\ge0$, which implies the conclusion.
\end{alist}
\end{sectionSolutions}

\chapter{Laplace Transforms}

\section{Introduction to the Laplace Transform}
\begin{sectionSolutions}
\item
\begin{alist}
\item $\cosh t\sin t=\dst{\frac{1}{
2}\left(e^t\sin t+e^{-t}\sin t\right)}
\leftrightarrow\dst{\frac{1}{2}\left[\frac{1}{(s-1)^2+1}+\frac{1}{(s+1)^2+1}\right]}
=\dst\frac{s^2+2}{[(s-1)^2+1][(s+1)^2+1]}$.

\item $\dst
\sin^2t=\frac{1-\cos2t}{2}\leftrightarrow\frac{1}{2}\left[\frac{1}{ s}
-\frac{s}{(s^2+4)}\right]=\frac{2}{ s(s^2+4)}$.

\item  $\dst \cos^22t=\frac{1}{2}\left[\frac{1}{ s}+\frac{s}{
s^2+16}\right]=
\frac{s^2+8}{ s(s^2+16)}$.

\item $\cosh^2
t=\dst\frac{(e^t+e^{-t})^2)}{4}=\dst\frac{(e^{2t}+2+e^{-2t})}{4}
\leftrightarrow\dst{\frac{1}{4}\left(\frac{1}{ s-2}+\frac{2}{ s}+\frac{1}{
s+2}\right)} =\dst\frac{s^2-2}{ s(s^2-4)}$.

\item $t\sinh2t=\dst\frac{te^{2t}-te^{-2t}}{2}
\leftrightarrow\dst{\frac{1}{2}\left(\frac{1}{(s-2)^2}-\frac{1}{(s+2)^2}\right)}
=\dst\frac{4s}{(s^2-4)^2}$.

\item  $\dst\sin t\cos t=\frac{\sin2t}{2}\leftrightarrow\frac{1}{
s^2+4}$.

\item  $\dst\sin(t+\pi/4)=\sin
t\cos(\pi/4)+\cos
t\cos(\pi/4)\leftrightarrow\frac{1}{\sqrt{2}}\,\frac{s+1}{ s^2+1}$.

\item $\cos 2t -\cos 3t\leftrightarrow\dst\frac{s}{
s^2+4}-\dst\frac{s}{ s^2+9}=\dst\frac{5s}{(s^2+4)(s^2+9)}$.

\item $\sin 2t +\cos4t\leftrightarrow\dst\frac{2}{
s^2+4}+\dst\frac{s}{ s^2+16} =\dst\frac{s^3+2s^2+4s+32}{(s^2+4)(s^2+16)}$.
\end{alist}

\addtocounter{enumi}{2}
\item%8.1.6
If $F(s)=\dst\int_0^\infty e^{-st}f(t)\,dt$, then
$F'(s)=\dst\int_0^\infty(-te^{-st})f(t)\,dt=-\int_0^\infty
e^{-st}(tf(t))\,dt$. Applying this argument repeatedly yields
the assertion.



\item
Let $f(t)=1$ and $F(s)=1/s$. From
Exercise~8.1.6,
$t^n\leftrightarrow(-1)^nF^{(n)}(s)=n!/s^{n+1}$.



\item If $|f(t)|\le Me^{s_0t}$ for  $t\ge t_0$, then
$|f(t)e^{-st}|\le Me^{-(s-s_0)t}$ for  $t\ge t_0$. Let
 $g(t)=e^{-st}f(t)$, $w(t)=Me^{-(s-s_0)t}$, and $\tau=t_0$.
Since $\int_{t_0}^\infty w(t)\,dt$ converges if $s>s_0$, $F(s)$
is defined for $s>s_0$.


\item
$\dst\int_0^Te^{-st}\left(\int_0^tf(\tau)\,d\tau\right)\,dt
=-\frac{e^{-st}}{ s}\int_0^t f(\tau)\,d\tau\bigg|_0^T+\frac{1}{
s}\int_0^Te^{-st}f(t)\,dt =-\frac{e^{-sT}}{ s}\int_0^T
f(\tau)\,d\tau+\frac{1}{ s}\int_0^Te^{-st}f(t)\,dt$. Since $f$ is of
exponential order $s_0$, the second integral on the right converges to
$\dst\frac{1}{ s}\laplace (f)$ as $T\to\infty$ (Exercise~8.1.10).
Now it suffices to show that (A)
$\dst\lim_{T\to\infty}e^{-sT}\int_0^Tf(\tau)\,d\tau=0$ if $s>s_0$.
Suppose that $|f(t)|\le Me^{s_0t}$ if $t\ge t_0$ and $|f(t)|\le K$ if
$0\le t\le t_0$, and let $T> t_0$. Then
$\dst\left|\int_0^Tf(\tau)\,d\tau\right|
\le\left|\int_0^{t_0}f(\tau)\,d\tau\right|+
\left|\int_{t_0}^Tf(\tau)\,d\tau \right|<Kt_0+M\int_{t_0}^T
e^{s_0\tau}\,d\tau<Kt_0+\frac{Me^{s_0T}}{ s_0}$, which proves (A).


\item
\begin{alist}
\item If $T>0$, then
$\dst\int_0^Te^{-st}f(t)\,dt=\int_0^Te^{-(s-s_0)t}(e^{-s_0t}f(t))\,dt$.
Use integration by parts with $u=e^{-(s-s_0)t}$,
$dv=e^{-s_0t}f(t)\,dt$, $du=-(s-s_0)e^{-(s-s_0)t}$, and $v=g$ t obtain
$\dst\int_0^T e^{-st}f(t)\,dt = e^{-(s-s_0)t}g(t)\lims0T
+(s-s_0)\int_0^Te^{-(s-s_0)t}g(t)\,dt$. Since $g(0)=0$ this reduces to
$\dst\int_0^T e^{-st}f(t)\,dt = e^{-(s-s_0)T}g(T)
+(s-s_0)\int_0^Te^{-(s-s_0)t}g(t)\,dt$. Since $|g(t)|\le M$ for all
$t\ge0$, we can let $t\to\infty$ to conclude that $\dst\int_0^\infty
e^{-st}f(t)\,dt = (s-s_0)\int_0^\infty e^{-(s-s_0)t}g(t)\,dt$ if
$s>s_0$.

\item  If $F(s_0)$ exists, then $g(t)$ is bounded on $[0,\infty)$.
Now apply \part{a}.

\item  Since  $f(t)=\dst\frac{1}{2}\frac{d}{ dt}\sin(e^{t^2})$,
$\dst\left|\int_0^tf(\tau)\,d\tau\right|=\frac{|\sin(e^{t^2})-\sin(1)|}{2}\le1$
for all $t\ge0$. Now apply \part{a} with $s_0=0$.
\end{alist}

\item
\begin{alist}
\item
$\dst\Gamma (\alpha)=\int_0^\infty x^{\alpha-1}e^{-x}\,dx=
\frac{x^\alpha e^{-x}}{\alpha}\lims0\infty+\frac{1}{\alpha}
\int_0^\infty x^\alpha e^{-x}\,dx=\frac{\Gamma(\alpha+1)}{\alpha}$.

\item  Use induction. $\Gamma(1)=\dst\int_0^\infty e^{-x}\,dx=1$.
If (A) $\Gamma(n+1)=n!$, then $\Gamma(n+2)=(n+1)\Gamma(n+1)$ (from
\part{a}) $=(n+1)n!$ (from (A)) $=(n+1)!$.


\item $\Gamma(\alpha+1)=\dst\int_0^\infty x^\alpha e^{-x}\,dt$. Let
$x=st$. Then $\Gamma(\alpha+1)=\dst\int_0^\infty (st)^\alpha
e^{-st}s\,dt$, so $\dst\int_0^\infty
e^{-st}t^\alpha\,dt=\frac{\Gamma(\alpha+1)}{\alpha}$.
\end{alist}

\item
\begin{alist}
\item
$\dst\int_0^2e^{-st}f(t)\,dt=\int_0^1e^{-st}t\,dt+\int_1^2e^{-st}(2-t)\,dt=
\left(\frac{1}{ s^2}-\frac{e^{-s}(s+1)}{ s^2}\right)
+\left(\frac{e^{-s}(s-1)}{ s^2}+\frac{e^{-2s}}{ s^2}\right)=
-\frac{2e^{-s}}{ s^2}+\frac{e^{-2s}}{ s^2}+\frac{1}{ s^2}
=\frac{(1-e^{-s})^2}{ s^2}$. Therefore,$F(s)=\dst\frac{(1-e^{-s})^2}{
s^2(1-e^{-2s})}=\frac{1-e^{-s}}{ s^2(1+e^{-s})}=
\frac{1}{ s^2}\tanh\frac{s}{ 2}$.

\item
$\dst\int_0^1e^{-st}f(t)\,dt=\int_0^{1/2}e^{-st}\,dt-\int_{1/2}^1e^{-st}\,dt=
\frac{1}{ s}-\frac{e^{-s/2}}{ s}+
\frac{e^{-s}}{ s}-\frac{e^{-s/2}}{ s}=
-\frac{2e^{-s/2}}{ s}+\frac{e^{-s}}{ s}+\frac{1}{ s}=
\frac{(1-e^{-s/2})^2}{ s}$. Therefore,$F(s)=\dst\frac{(1-e^{-s/2})^2}{
s(1-e^{-s})}=\frac{1-e^{-s/2}}{ s(1+e^{-s/2})}=
\frac{1}{ s}\tanh\frac{s}{ 4}$.

\item
$\dst\int_0^\pi e^{-st}f(t)\,dt=\int_0^\pi e^{-st}\sin t\,dt=
\frac{1+e^{-\pi s}}{(s^2+1)}$. Therefore,$F(s)=\dst\frac{1+e^{-\pi
s}}{(s^2+1)(1-e^{-\pi s})}
\frac{1}{ s^2+1}\coth \frac{\pi s}{2}$.

\item
$\dst\int_0^{2\pi}e^{-st}f(t)\,dt=\int_0^\pi
e^{-st}\sin t\,dt=
\frac{1+e^{-\pi s}}{(s^2+1)}$. Therefore,$F(s)=\dst\frac{1+e^{-\pi s}}{(s^2+1)
1+e^{-2\pi s}}=
\frac{1}{(s^2+1)(1-e^{-\pi s})}$.
\end{alist}
\end{sectionSolutions}

\section{The Inverse Laplace Transform}

\begin{sectionSolutions}
\item
\begin{alist}
\item
$\dst\frac{2s+3}{(s-7)^4}=\dst\frac{2(s-7)+17}{(s-7)^4}=
\dst\frac{2}{(s-7)^3}+\frac{17}{(s-7)^4}
=\dst\frac{2!}{(s-7)^3}+\dst\frac{17}{6}\frac{3!}{(s-7)^4}
\leftrightarrow e^{7t}\dst{\left(t^2+\frac{17}{6}t^3\right)}$.

\item
$\dst\frac{s^2-1}{(s-2)^6}=\dst\frac{[(s-2)+2]^2-1}{(s-2)^6}
=\dst\frac{(s-2)^2+4(s-2)+3}{(s-2)^6}
=\dst\frac{1}{(s-2)^4}+\dst\frac{4}{(s-2)^5}+\dst\frac{3}{(s-2)^6}
=\dst{\frac{1}{6}\frac{3!}{(s-2)^4}}+\dst{\frac{1}{6}\frac{4!}{(s-2)^5}}
+\dst{\frac{1}{40} \frac{5!}{(s-2)^6}}\leftrightarrow\dst{\left(\frac{1}{
6}t^3+\frac{1}{6}t^4+\frac{1}{40}t^5\right)e^{2t}}$.

\item $\dst\frac{s+5}{ s^2+6s+18}=\frac{(s+3)}{(s+3)^2+9}
+\frac{2}{3}\frac{3}{(s+3)^2+9}
\leftrightarrow e^{-3t}\left(\cos 3t+\frac{2}{3}\sin 3t\right)$.

\item  $\dst\frac{2s+1}{ s^2+9}=2\frac{s}{
s^2+9}+\frac{1}{3}\frac{3}{ s^2+9}\leftrightarrow
\dst{2\cos 3t+\frac{1}{3}\sin 3t}$.

\item  $\dst\frac{s}{ s^2+2s+1}=\frac{(s+1)-1}{(s+1)^2}=
\frac{1}{ s+1}-\frac{1}{ (s+1)^2}\leftrightarrow
(1-t)e{^{-t}}$.

\item $\dst\frac{s+1}{ s^2-9}=\frac{s}{ s^2-9}+\frac{1}{3}\frac{3}{
s^2-9}\leftrightarrow {\cosh 3t+\frac{1}{3}\sinh 3t}$.

\item Expand the numerator in powers of $s+1$:
$s^3+2s^2-s-3=[(s+1)-1]^3+2[(s+1)-1]^2-[(s+1)-1]-3=
(s+1)^3-(s+1)^2-2(s+1)-1$; therefore
$\dst\frac{s^3+2s^2-s-3}{(s+1)^4}=\frac{1}{ s+1}-\frac{1}{(s+1)^2}
-\frac{2}{(s+1)^3}-\frac{1}{6}\frac{6}{(s+1)^4}\leftrightarrow
{\left(1-t-t^2-\frac{1}{6}t^3\right)e^{-t}}$.

\item  $\dst\frac{2s+3}{(s-1)^2+4}=2\frac{(s-1)}{(s-1)^2+4}
+\frac{5}{2}\frac{2}{(s-1)^2+4}\leftrightarrow
{e^t\left(2\cos 2t+\frac{5}{2}\sin 2t\right)}$.

\item $\dst\frac{1}{ s}-\frac{s}{ s^2+1}\leftrightarrow 1-\cos t$.

\item
$\dst\frac{3s+4}{ s^2-1}=\dst\frac{3s}{ s^2-1}+\frac{4}{ s^2-1}
\leftrightarrow3\cosh t+4\sinh t$.
Alternatively,
$\dst\frac{3s+4}{
s^2-1}=\dst\frac{3s+4}{(s-1)(s+1)}=\dst{\frac{1}{2}\left[\frac{7}{
s-1}-\frac{1}{ s+1}\right]}
\leftrightarrow\dst\frac{7e^t-e^{-t}}{2}$.

\item  $\dst\frac{3}{ s-1}+\frac{4s+1}{ s^2+9}=3\frac{1}{ s-1}
+4\frac{s}{ s^2+9}+\frac{1}{3}\frac{3}{ s^2+9}\leftrightarrow
\dst{3e^t+4\cos 3t+\frac{1}{3}\sin 3t}$.

\item  $\dst\frac{3}{(s+2)^2}-\frac{2s+6}{ s^2+4}=
3\frac{1}{(s+2)^2}-2\frac{s}{ s^2+4}-3\frac{2}{ s^2+4}\leftrightarrow
3te^{-2t}-2\cos 2t-3\sin 2t$.
\end{alist}

\item
\begin{alist}
\item
\[
\frac{2+3s}{(s^2+1)(s+2)(s+1)}=
\frac{A}{ s+2}+\frac{B}{ s+1}+\frac{Cs+D}{ s^2+1},
\]
where
\begin{gather*}
A(s^2+1)(s+1)+
B(s^2+1)(s+2)+(Cs+D)(s+2)(s+1)=2+3s.
\\
\begin{array}{rc rl}
-5A&=&-4&(\text{when }s=-2);\\
2B&=&-1& (\text{when }s=-1);\\
A+2B+2D&=&2&(\text{when }s=0);\\
A+B+C&=&0&(\text{equate coefficients of }s^3).
\end{array}
\end{gather*}
Solving this system yields $A=\dst\frac{4}{5}$, $B=-\dst\frac{1}{2}$,
$C=-\dst\frac{3}{10}$, $D=\dst\frac{11}{10}$. Therefore,
\begin{align*}
\frac{2+3s}{(s^2+1)(s+2)(s+1)}&=
\frac{4}{5}\frac{1}{ s+2}-\frac{1}{2}\frac{1}{ s+1}-\frac{1}{10}\frac{3s-11}{
s^2+1}\\
&\leftrightarrow
\frac{4}{5}e^{-2t}-\frac{1}{2}e^{-t}-\frac{3}{10}\cos t
+\frac{11}{10}\sin t.
\end{align*}

\item
\[
\frac{3s^2+2s+1}{(s^2+1)(s^2+2s+2)}=
\frac{As+B}{ s^2+1}+\frac{C(s+1)+D}{(s+1)^2+1},
\]
where
\begin{gather*}
(As+B)((s+1)^2+1)+(C(s+1)+D)(s^2+1)=3s^2+2s+1.
\\
\begin{array}{rc rl}
2B+C+D&=&1&(\text{when }s=0);\\
-A+B+2D&=&2& (\text{when }s=-1);\\
2B+C+D&=&1&(\text{when }s=0);\\
A+C&=&0&(\text{equate coefficients of }s^3).
\end{array}
\end{gather*}
Solving this system yields
$A=6/5$, $B=2/5$, $C=-6/5$, $D=7/5$.
Therefore,
\begin{align*}
\frac{3s^2+2s+1}{(s^2+1)(s^2+2s+2)}&=
\frac{1}{5}\left[\frac{6s+2}{ s^2+1}-\frac{6(s+1)-7}{(s+1)^2+1}\right]\\
&\leftrightarrow
 \frac{6}{5}\cos t+\frac{2}{5}\sin t
-\frac{6}{5}e^{-t}\cos t+\frac{7}{5}e^{-t}\sin t.
\end{align*}

\item
$s^2+2s+5=(s+1)^2+4$;
\[
\frac{3s+2}{(s-2)((s+1)^2+4)}=\frac{A}{ s-2}+
\frac{B(s+1)+C}{(s+1)^2+4},
\]
where
\begin{gather*}
A\left((s+1)^2)+4\right)+
\left(B(s+1)+C\right)(s-2)=3s+2.
\\
\begin{array}{rc rl}
13A&=&8&(\text{when }s=2);\\
4A-3C&=&-1&(\text{when }s=-1);\\
A+B&=&0&(\text{equate coefficients of }s^2).
\end{array}
\end{gather*}
Solving this system yields $A=\dst\frac{8}{13}$, $B=-\dst\frac{8}{13}$,
$C=\dst\frac{15}{13}$. Therefore,
\begin{align*}
\frac{3s+2}{(s-2)((s+1)^2+4)}&=\frac{1}{13}\left[\frac{8}{
s-2}-
\frac{8(s-1)-15}{(s+1)^2+4}\right]\\
&\leftrightarrow
\frac{8}{13}e^{2t}-\frac{8}{13}e^{-t}\cos 2t+\frac{15}{26}
e^{-t}\sin 2t.
\end{align*}

\item
\[
\frac{3s^2+2s+1}{(s-1)^2(s+2)(s+3)}=
\frac{A}{ s-1}+\frac{B}{(s-1)^2}+\frac{C}{ s+2}+\frac{D}{ s+3},
\]
where
\begin{gather*}
(A(s-1)+B)(s+2)(s+3)+(C(s+3)+D(s+2))(s-1)^2=3s^2+2s+1.
\\
\begin{array}{rc rl}
12B&=&6&(\text{when }s=1);\\
9C&=&9& (\text{when }s=-2);\\
-16D&=&22&(\text{when }s=-3);\\
A+C+D&=&0&(\text{equate coefficients of }s^3).
\end{array}
\end{gather*}
Solving this system yields
$A=3/8$, $B=1/2$, $C=1$, $D=-11/8$.
Therefore,
\begin{align*}
\frac{3s^2+2s+1}{(s-1)^2(s+2)(s+3)}&=
\frac{3}{8}\frac{1}{ s-1}+\frac{1}{2}\frac{1}{(s-1)^2}+\frac{1}{
s+2}-\frac{11}{8}\frac{1}{ s+3}\\
&\leftrightarrow
\frac{3}{8}e^t+\frac{1}{2}te^t+e^{-2t}-\frac{11}{8}e^{-3t}.
\end{align*}

\item
\[
\frac{2s^2+s+3}{(s-1)^2(s+2)^2}=
\frac{A}{ s-1}+\frac{B}{(s-1)^2}+\frac{C}{ s+2}+\frac{D}{(s+2)^2},
\]
where
\begin{gather*}
(A(s-1)+B)(s+2)^2+(C(s+2)+D)(s-1)^2=2s^2+s+3.
\\
\begin{array}{rc rl}
9B&=&6&(\text{when }s=1);\\
9D&=&9& (\text{when }s=-2);\\
-4A+4B+2C+D&=&3&(\text{when }s=0);\\
A+C&=&0&(\text{equate coefficients of }s^3).
\end{array}
\end{gather*}
Solving this system yields
$A=1/9$, $B=2/3$, $C=-1/9$, $D=1$.
Therefore,
\begin{align*}
\frac{2s^2+s+3}{(s-1)^2(s+2)^2}&=
\frac{1}{9}\frac{1}{ s-1}+\frac{2}{3}\frac{1}{(s-1)^2}-\frac{1}{9}\frac{1}{
s+2}+\frac{1}{(s+2)^2}\\
&\leftrightarrow
\frac{1}{9}e^t+\frac{2}{3}te^t-\frac{1}{9}e^{-2t}+te^{-2t}.
\end{align*}

\item
\[
\frac{3s+2}{(s^2+1)(s-1)^2}=\frac{A}{ s-1}+\frac{B}{(s-1)^2}+
\frac{Cs+D}{ s^2+1},
\]
where
\[
A(s-1)(s^2+1)+B(s^2+1)+(Cs+D)(s-1)^2=3s+2.
\tag{A}
\]
Setting $s=1$ yields $2B=5$, so $B=\dst\frac{5}{2}$.
Substituting this into (A) shows that
\begin{align*}
A(s-1)(s^2+1)+(Cs+D)(s-1)^2&=3s+2-\frac{5}{2}(s^2+1)\\
&=-\frac{5s^2-6s+1}{2}=-\frac{(s-1)(5s-1)}{2}.
\end{align*}
Therefore,
\begin{gather*}
A(s^2+1)+(Cs+D)(s-1)=\frac{1-5s}{2}.
\\
\begin{array}{rc rl}
2A&=&-2&(\text{when }s=1);\\
A-D&=&1/2&(\text{when }s=0);\\
A+C&=&0&(\text{equate coefficients of }s^2).
\end{array}
\end{gather*}
Solving this system yields $A=-1$, $C=1$,
$D=-\dst\frac{3}{2}$. Therefore,
\begin{align*}
\frac{3s+2}{(s^2+1)(s-1)^2}&=-\frac{1}{ s-1}+\frac{5}{2}\frac{1}{(s-1)^2}
+\frac{s-3/2}{ s^2+1}\\
&\leftrightarrow
-e^t+\frac{5}{2}te^t+\cos t-\frac{3}{2}\sin t.
\end{align*}
\end{alist}

\item
\begin{alist}
\item
\[
\frac{17 s-15}{(s^2-2s+5)(s^2+2s+10)}
=\frac{A(s-1)+B}{(s-1)^2+4}+\frac{C(s+1)+D}{(s+1)^2+9}
\]
where
\begin{gather*}
(A(s-1)+B)((s+1)^2+9)+(C(s+1)+D)((s-1)^2+4)=17 s-15.
\\
\begin{array}{rc rl}
13B+8C+4D&=&2&(\text{when }s=1);\\
-18A+9B+8D&=&-32& (\text{when }s=-1);\\
-10A+10B+5C+5D&=&-15&(\text{when }s=0);\\
A+C&=&0&(\text{equate coefficients of }s^3).
\end{array}
\end{gather*}
Solving this system yields $A=1$, $B=2$,
$C=-1$, $D=-4$. Therefore,
\begin{align*}
\frac{17 s-15}{(s^2-2s+5)(s^2+2s+10)}
&=
\frac{(s-1)+2}{(s-1)^2+4}-\frac{(s+1)+4}{(s+1)^2+9}
\\&\leftrightarrow
 e^t(\cos 2t+\sin 2t)-e^{-t}\left(\cos
3t+\frac{4}{3}\sin 3t\right).
\end{align*}

\item
\[\frac{8s+56}{(s^2-6s+13)(s^2+2s+5)}
=\frac{A(s-3)+B}{(s-3)^2+4}+\frac{C(s+1)+D}{(s+1)^2+4}
\]
where
\begin{gather*}
(A(s-3)+B)((s+1)^2+4)+(C(s+1)+D)((s-3)^2+4)=8s+56.
\\
\begin{array}{rc rl}
20B+16C+4D&=&80&(\text{when }s=3);\\
-16A+4B+20D&=&48& (\text{when }s=-1);\\
-15A+5B+13C+13D&=&56&(\text{when }s=0);\\
A+C&=&0&(\text{equate coefficients of }s^3).
\end{array}
\end{gather*}
Solving this system yields $A=-1$, $B=3$,
$C=1$, $D=1$. Therefore,
\begin{align*}
\frac{8s+56}{(s^2-6s+13)(s^2+2s+5)}
&=\frac{-(s-3)+3}{(s-3)^2+4}+\frac{(s+1)+1}{(s+1)^2+4}
\\&\leftrightarrow
{e^{3t}\left(-\cos 2t+\frac{3}{2}\sin
2t\right)+e^{-t}\left(\cos 2t+\frac{1}{2}\sin 2t\right)}.
\end{align*}

\item
\[\frac{s+9}{(s^2+4s+5)(s^2-4s+13)}
=\frac{A(s+2)+B}{(s+2)^2+1}+\frac{C(s-2)+D}{(s-2)^2+9}
\]
where
\begin{gather*}
(A(s+2)+B)((s-2)^2+9)+(C(s-2)+D)((s+2)^2+1)=s+9.
\\
\begin{array}{rc rl}
25B-4C+D&=&7&(\text{when }s=-2);\\
36A+9B+17D&=&11& (\text{when }s=2);\\
26A+13B-10C+5D&=&9&(\text{when }s=0);\\
A+C&=&0&(\text{equate coefficients of }s^3).
\end{array}
\end{gather*}
Solving this system yields $A=1/8$, $B=1/4$,
$C=-1/8$, $D=1/4$. Therefore,
\begin{align*}
\frac{s+9}{(s^2+4s+5)(s^2-4s+13)}
&=
=\left[\frac{1}{8}\frac{(s+2)+2}{(s+2)^2+1}-\frac{(s-2)-2}{(s-2)^2+3}\right]
\\&\leftrightarrow
{e^{-2t}\left(\frac{1}{8}\cos t+\frac{1}{
4}\sin t\right)-e^{2t}\left(\frac{1}{8}\cos 3t-\frac{1}{ 12}\sin
3t\right)}.
\end{align*}

\item
\[\frac{3s-2}{(s^2-4s+5)(s^2-6s+13)}
=\frac{A(s-2)+B}{(s-2)^2+1}+\frac{C(s-3)+D}{(s-3)^2+4}
\]
where
\begin{gather*}
(A(s-2)+B)((s-3)^2+4)+(C(s-3)+D)((s-2)^2+1)=3s-2.
\\
\begin{array}{rc rl}
5B-C+D&=&4&(\text{when }s=2);\\
4A+4B+2D&=&7& (\text{when }s=3);\\
-26A+13B-15C+5D&=&-2&(\text{when }s=0);\\
A+C&=&0&(\text{equate coefficients of }s^3).
\end{array}
\end{gather*}
Solving this system yields $A=1$, $B=1/2$,
$C=-1$, $D=1/2$. Therefore,
\begin{align*}
\frac{3s-2}{(s^2-4s+5)(s^2-6s+13)}
&=
=\frac{1}{2}\left[\frac{2(s-2)+1}{(s-2)^2+1}-\frac{2(s-3)-1}{(s-3)^2+4}\right]
\\&\leftrightarrow
e^{2t}\left(\cos t+\frac{1}{2}\sin
t\right)-e^{3t}\left(\cos 2t-\frac{1}{4}\sin 2t\right).
\end{align*}

\item
\[
\frac{3s-1}{(s^2-2s+2)(s^2+2s+5)}
=\frac{A(s-1)+B}{(s-1)^2+1}+\frac{C(s+1)+D}{(s+1)^2+4}
\]
where
\begin{gather*}
(A(s-1)+B)((s+1)^2+4)+(C(s+1)+D)((s-1)^2+1)=3s-1.
\\
\begin{array}{rc rl}
8B+2C+D&=&2&(\text{when }s=1);\\
-8A+4B+5D&=&-4& (\text{when }s=-1);\\
-5A+5B+2C+2D&=&-1&(\text{when }s=0);\\
A+5B+C&=&0&(\text{equate coefficients of }s^3).
\end{array}
\end{gather*}
Solving this system yields $A=1/5$, $B=2/5$,
$C=-1/5$, $D=-4/5$. Therefore,
\begin{align*}
\frac{3s-1}{(s^2-2s+2)(s^2+2s+5)}
&=
\frac{1}{5}\left[\frac{(s-1)+2}{(s-1)^2+1}-\frac{(s+1)+4}{(s+1)^2+4}\right].
\\&\leftrightarrow
e^t\left(\frac{1}{5}\cos t+\frac{2}{5} \sin t\right)
-e^{-t}\left(\frac{1}{5}\cos 2t+\frac{2}{5}\sin 2t\right).
\end{align*}

\item
\[
\frac{20s+40}{(4s^2-4s+5)(4s^2+4s+5)}
=\frac{A(s-1/2)+B}{(s-1/2)^2+1}+
\frac{C(s+1/2)+D}{(s+1/2)^2+1}
\]
where
\begin{gather*}
(A(s-1/2)+B)((s+1/2)^2+1)+(C(s+1/2)+D)((s-1/2)^2+1)=\frac{5s+10}{4}.
\\
\begin{array}{rc rl}
2B+C+D&=&25/8&(\text{when }s=1/2);\\
-A+B+2D&=&15/8& (\text{when }s=-1/2);\\
-5A+10B+5C+10D&=&20&(\text{when }s=0);\\
A+C&=&0&(\text{equate coefficients of }s^3).
\end{array}
\end{gather*}
Solving this system yields $A=-1$, $B=9/8$,
$C=1$, $D=-1/8$. Therefore,
\begin{align*}
\frac{20s+40}{(4s^2-4s+5)(4s^2+4s+5)}
&=
\frac{1}{8}\left[\frac{-8(s-1/2)+9}{(s-1/2)^2+1}+\frac{8(s+1/2)
-1}{(s+1/2)^2+1}\right]
\\&\leftrightarrow
e^{t/2}\left(-\cos t+\frac{9}{8}\sin
t\right)+e^{-t/2}\left(\cos t-\frac{1}{8}\sin t\right).
\end{align*}
\end{alist}

\item
\begin{alist}
\item
\[
\frac{2s+1}{(s^2+1)(s-1)(s-3)}
=\frac{A}{ s-1}+\frac{B}{ s-3}+\frac{Cs+D}{ s^2+1}
\]
where
\begin{gather*}
(A(s-3)+B(s-1))(s^2+1)+(Cs+D)(s-1)(s-3)=2s+1.
\\
\begin{array}{rc rl}
-4A&=&3&(\text{when }s=1);\\
20B&=&7& (\text{when }s=3);\\
-3A-B+3D&=&1& (\text{when }s=0);\\
A+B+C&=&0&(\text{equate coefficients of }s^3).
\end{array}
\end{gather*}
Solving this system yields $A=-3/4$, $B=7/20$,
$C=2/5$, $D=-3/10$. Therefore,
\begin{align*}
\frac{2s+1}{(s^2+1)(s-1)(s-3)}
&=
-\frac{3}{4}\frac{1}{ s-1}+\frac{7}{20}\frac{1}{ s-3}+\frac{2}{5}\frac{s}{
s^2+1}-\frac{3}{10}\frac{1}{ s^2+1}
\\&\leftrightarrow
-\frac{3}{4}e^t+\frac{7}{20}e^{3t}+\frac{2}{5}\cos t-\frac{3}{ 10}\sin t.
\end{align*}

\item
\[
\frac{s+2}{(s^2+2s+2)(s^2-1)}
=\frac{A}{ s-1}+\frac{B}{ s+1}+\frac{C(s+1)+D}{(s+1)^2+1}
\]
where
\begin{gather*}
(A(s+1)+B(s-1))((s+1)^2+1)+(C(s+1)+D)(s^2-1)=s+2.
\\
\begin{array}{rc rl}
10A&=&3&(\text{when }s=1);\\
-2B&=&1& (\text{when }s=-1);\\
2A-2B-C-D&=&2& (\text{when }s=0);\\
A+B+C&=&0&(\text{equate coefficients of }s^3).
\end{array}
\end{gather*}
Solving this system yields $A=3/10$, $B=-1/2$,
$C=1/5$, $D=-3/5$. Therefore,
\begin{align*}
\frac{s+2}{(s^2+2s+2)(s^2-1)}
&=
\frac{3}{ 10}\frac{1}{ s-1}-\frac{1}{2}\frac{1}{ s+1}+\frac{1}{5}\frac{s+1}{(s+1)^2+1}
-\frac{3}{5}\frac{1}{(s+1)^2+1}
\\&\leftrightarrow
\frac{3}{ 10}e^{t}-\frac{1}{2}e^{-t}
+\frac{1}{5}e^{-t}\cos te^{-t}\sin t.
\end{align*}

\item
\[
\frac{2s-1}{(s^2-2s+2)(s+1)(s-2)}
=\frac{A}{ s-2}+\frac{B}{ s+1}+\frac{C(s-1)+D}{(s-1)^2+1}
\]
where
\begin{gather*}
(A(s+1)+B(s-2))((s-1)^2+1)+(C(s-1)+D)(s-2)(s+1)=2s-1.
\\
\begin{array}{rc rl}
6A&=&3&(\text{when }s=2);\\
-15B&=&-3& (\text{when }s=-1);\\
2A-4B+2C-2D&=&-1& (\text{when }s=0);\\
A+B+C&=&0&(\text{equate coefficients of }s^3).
\end{array}
\end{gather*}
Solving this system yields $A=1/2$, $B=1/5$,
$C=-7/10$, $D=-1/10$. Therefore,
\begin{align*}
\frac{2s-1}{(s^2-2s+2)(s+1)(s-2)}
&=
=\frac{1}{2}\frac{1}{ s-2}+\frac{1}{5}\frac{1}{ s+1}-\frac{7}{ 10}\frac{s-1}{(s-1)^2+1}
-\frac{1}{ 10}\frac{1}{(s-1)^2+1}
\\&\leftrightarrow
\frac{1}{2}e^{2t}+\frac{1}{5} e^{-t}
-\frac{7}{ 10}e^t\cos t-\frac{1}{10}e^t\sin t.
\end{align*}

\item
\[
\frac{s-6}{(s^2-1)(s^2+4)}
=\frac{A}{ s-1}+\frac{B}{ s+1}+\frac{Cs+D}{ s^2+4}
\]
where
\begin{gather*}
(A(s+1)+B(s-1))(s^2+4)+(Cs+D)(s^2-1)=s-6.
\\
\begin{array}{rc rl}
10A&=&-5&(\text{when }s=1);\\
-10B&=&-7& (\text{when }s=-1);\\
4A-4B-D&=&-6& (\text{when }s=0);\\
A+B+C&=&0&(\text{equate coefficients of }s^3).
\end{array}
\end{gather*}
Solving this system yields $A=-1/2$, $B=7/10$,
$C=-1/5$, $D=6/5$. Therefore,
\begin{align*}
\frac{s-6}{(s^2-1)(s^2+4)}
&=
=-\frac{1}{2}\frac{1}{ s-1}+\frac{7}{ 10}\frac{1}{ s+1}-\frac{1}{5}\frac{s}{ s^2+4}
+\frac{3}{5}+\frac{1}{ s^2+4}
\\&\leftrightarrow
-\frac{1}{2}e^t+\frac{7}{ 10}e^{-t}-\frac{1}{5}\cos 2t+\frac{3}{5}\sin 2t.
\end{align*}

\item
\[
\frac{2s-3}{ s(s-2)(s^2-2s+5)}
=\frac{A}{ s}+\frac{B}{ s-2}+\frac{C(s-1)+D}{(s-1)^2+4}
\]
where
\begin{gather*}
(A(s-2)+Bs)((s-1)^2+4)+(C(s-1)+D)s(s-2)=2s-3.
\\
\begin{array}{rc rl}
-10A&=&-3&(\text{when }s=0);\\
10B&=&1& (\text{when }s=2);\\
-4A+4B-D&=&-1& (\text{when }s=1);\\
A+B+C&=&0&(\text{equate coefficients of }s^3).
\end{array}
\end{gather*}
Solving this system yields $A=3/10$, $B=1/10$,
$C=-2/5$, $D=1/5$. Therefore,
\begin{align*}
\frac{2s-3}{ s(s-2)(s^2-2s+5)}
&=
=\frac{3}{ 10s}+\frac{1}{ 10}\frac{1}{ s-2}-\frac{2}{5}\frac{s-1}{(s-1)^2+4}
+\frac{1}{ 5}\frac{1}{(s-1)^2+4}
\\&\leftrightarrow
\frac{3}{ 10}+\frac{1}{ 10}e^{2t}-\frac{2}{5}e^t\cos 2t
+\frac{1}{ 10}e^t\sin2t.
\end{align*}

\item
\[
\frac{5s-15}{(s^2-4s+13)(s-2)(s-1)}
=\frac{A}{ s-1}+\frac{B}{ s-2}+\frac{C(s-2)+D}{(s-2)^2+9}
\]
where
\begin{gather*}
(A(s-2)+B(s-1))((s-2)^2+9)+(C(s-2)+D)(s-1)(s-2)=5s-15.
\\
\begin{array}{rc rl}
-10A&=&-10&(\text{when }s=1);\\
9B&=&-5& (\text{when }s=2);\\
-26A-13B-4C+2D&=&-15& (\text{when }s=0);\\
A+B+C&=&0&(\text{equate coefficients of }s^3).
\end{array}
\end{gather*}
Solving this system yields $A=1$, $B=-5/9$,
$C=-4/9$, $D=1$. Therefore,
\begin{align*}
\frac{5s-15}{(s^2-4s+13)(s-2)(s-1)}
&=
=\frac{1}{ s-1}-\frac{5}{9}\frac{1}{ s-2}-\frac{4}{9}\frac{s-2}{(s-2)^2+9}
+\frac{1}{(s-2)^2+9}
\\&\leftrightarrow
e^t-\frac{5}{9}e^{2t}-\frac{4}{9}e^{2t}\cos 3t+\frac{1}{3}e^{2t}\sin 3t.
\end{align*}
\end{alist}

\item
\begin{alist}
\item Let $i=1$. (The proof for $i=2,\dots,n$) is similar.
Multiplying the given equation through by $s-s_1$ yields
\[
\frac{P(s)}{(s-s_2)\cdots(s-s_n)}=
A_1+(s-s_1)\left[\frac{A_2}{ s-s_2}+\cdots+\frac{A_n}{ s-s_n}\right],
\]
and setting $s=s_1$ yields
$A_1=\dst\frac{P(s_1)}{(s_1-s_2)\cdots(s_2-s_n)}$.

\item  From calculus we know that $F$  has a partial
fraction expansion of the form
$\dst\frac{P(s)}{(s-s_1)Q_1(s)}=\frac{A}{ s-s_1}+G(s)$ where $G$ is
continuous at $s_1$. Multiplying through by $s-s_1$ shows that
$\dst\frac{P(s)}{ Q_1(s)}=A+(s-s_1)G(s)$. Now set $s=s_1$ to obtain
$A=\dst\frac{P(s_1)}{ Q(s_1)}$.

\item The result in \part{b} is generalization of the result
in \part{a}, since it shows that if $s_1$ is a simple zero
of the denominator of the rational function, then Heaviside's
method can be used to determine the coefficient of $1/(s-s_1)$
in the partial fraction expansion even if some of the other zeros
of the denominator are repeated or complex.
\end{alist}
\end{sectionSolutions}

\section{Solution of Initial Value Problems}
\begin{sectionSolutions}
\item
\[
(s^2-s-6)Y(s)=\frac{2}{ s}+s-1=\frac{2+s(s-1)}{ s}.
\]
Since $(s^2-s-6)=(s-3)(s+2)$,
\[
Y(s)=\frac{2+s(s-1)}{ s(s-3)(s+2)}
=-\frac{1}{3s}+\frac{8}{15}\frac{1}{ s-3}+\frac{4}{5}\frac{1}{
s+2}
\]
and $y=-\dst\frac{1}{3}+\dst\frac{8}{15}e^{3t}+\dst\frac{4}{5}e^{-2t}$.


\item
\[
(s^2-4)Y(s)=\frac{2}{ s-3}+(-1+s)=\frac{2+(s-1)(s-3)}{ s-3}.
\]
Since $s^2-4=(s-2)(s+2)$,
\[
Y(s)=\frac{2+(s-1)(s-3)}{(s-2)(s+2)(s-3)}=
-\frac{1}{4}\frac{1}{ s-2}+\frac{17}{20}\frac{1}{ s+2}+\frac{2}{5}\frac{1}{ s-3}
\]
and   $y=\dst-\frac{1}{4}e^{2t}+\frac{17}{
20}e^{-2t}+\frac{2}{5}e^{3t}$.


\item
\[
(s^2+3s+2)Y(s)=\frac{6}{ s-1}+(-1+s)+3=\frac{6+(s-1)(s+2)}{ s-1}.
\]
Since $s^2+3s+2=(s+2)(s+1)$,
\[
Y(s)=\frac{6+(s-1)(s+2)}{(s-1)(s+2)(s+1)}=
\frac{1}{ s-1}+\frac{2}{ s+2}-\frac{2}{ s+1}
\]
and
 $y=e^t+2e^{-2t}-2e^{-t}$.


\item
\[
 (s^2-3s+2)Y(s)=\frac{2}{ s-3}+(-1+s)-3=\frac{2+(s-3)(s-4)}{ s-3}.
\]
Since $s^2-3s+2=(s-1)(s-2)$,
\[
Y(s)=\frac{2+(s-3)(s-4)}{(s-1)(s-2)(s-3)}=
\frac{4}{ s-1}-\frac{4}{ s-2}+\frac{1}{ s-3}
\]
and
  $y=4e^t-4e^{2t}+e^{3t}$.



\item
\[
(s^2-3s+2)Y(s)=\frac{1}{ s-3}+(-4-s)+3=\frac{1-(s-3)(s+1)}{ s-3}.
\]
Since $s^2-3s+2=(s-1)(s-2)$,
\[
Y(s)=\frac{1-(s-3)(s+1)}{(s-1)(s-2)(s-3)}=
\frac{5}{2}\frac{1}{ s-1}-\frac{4}{ s-2}+\frac{1}{2}\frac{1}{ s-3}
\]
and
  $y=\dst{\frac{5}{2}e^t-4e^{2t}+\frac{1}{
2}e^{3t}}$.



\item
\[
(s^2+s-2)Y(s)=-\frac{4}{ s}+(3+2s)+2=\frac{-4+s(5+2s)}{ s}.
\]
Since $(s^2+s-2)=(s+2)(s-1)$,
\[
Y(s)=\frac{-4+s(5+2s)}{ s(s+2)(s-1)}
=\frac{2}{ s}-\frac{1}{ s+2}+\frac{1}{ s-1},
\]
and  $y=2-e^{-2t}+e^t$.



\item
\[
 (s^2-s-6)Y(s)=\frac{2}{ s}+s-1=\frac{2+s(s-1)}{ s}.
\]
Since $s^2-s-6=(s-3)(s+2)$,
\[
Y(s)=\frac{2+s(s-1)}{ s(s-3)(s+2)}=
-\frac{1}{3s}+\frac{8}{15}\frac{1}{ s-3}+\frac{4}{5}\frac{1}{ s+2}
\]
and
 $y=\dst{-\frac{1}{3}+\frac{8}{15}e^{3t}
+\frac{4}{5}e^{-2t}}$.


\item
\[
(s^2-1)Y(s)=\frac{1}{ s}+s=\frac{1+s^2}{ s}.
\]
Since $s^2-1=(s-1)(s+1)$,
\[
Y(s)=\frac{1+s^2}{ s(s-1)(s+1)}=
-\frac{1}{ s}+\frac{1}{ s-1}+\frac{1}{ s+1}
\]
and
 $y=-1+e^t+e^{-t}$.


\item
\[
(s^2+s)Y(s)=\frac{2}{ s-3}+(4-s)-1=\frac{2-(s-3)^2}{ s-3}.
\]
Since $s^2+s=s(s+1)$,
\[
Y(s)=\frac{2-(s-3)^2}{ s(s+1)(s-3)}=\frac{7}{3s}-\frac{7}{2}\frac{1}{ s+1}
+\frac{1}{6}\frac{1}{ s-3}
\]
and $y=\dst{\frac{7}{3}-\frac{7}{2}e^{-t} +\frac{1}{6}e^{3t}}$.

\item
\[
(s^2+1)Y(s)=\frac{1}{ s^2}+2,\quad\text{so}\quad
Y(s)=\frac{1}{(s^2+1)s^2}+\frac{2}{ s^2+1}.
\]
Substituting $x=s^2$ into
\[
\frac{1}{(x+1)x}=\frac{1}{ x+1}-\frac{1}{ x}
\quad\text{yields}\quad
\frac{1}{(s^2+1)s^2}=\frac{1}{ s^2}-\frac{1}{ s^2+1},
\]
so $Y(s)=\dst\frac{1}{ s^2}+\dst\frac{1}{ s^2+1}$ and $y=t+\sin t$.

\item
\[
 (s^2+5s+6)Y(s)=\frac{2}{ s+1}+(3+s)+5=\frac{2+(s+1)(s+8)}{ s+1}.
\]
Since $s^2+5s+6=(s+2)(s+3)$,
\[
Y(s)=\frac{2+(s+1)(s+8)}{(s+1)(s+2)(s+3)}=
\frac{1}{ s+1}+\frac{4}{ s+2}-\frac{4}{ s+3}
\]
and
 $y=e^{-t}+4e^{-2t}-4e^{-3t}$.

\item
\[
(s^2-2s-3)Y(s)=\frac{10s}{ s^2+1}+(7+2s)-4=
\frac{10s}{ s^2+1}+(2s+3).
\]
Since $s^2-2s-3=(s-3)(s+1)$,
\begin{gather*}
Y(s)=
\frac{10s}{(s-3)(s+1)(s^2+1)}+\frac{2s+3}{(s-3)(s+1)}.
\tag{A}\\
\frac{2s+3}{(s-3)(s+1)}=\frac{9}{4}\frac{1}{ s-3}-\frac{1}{4}\frac{1}{ s+1}
\leftrightarrow\frac{9}{4}e^{3t}-\frac{1}{4}e^{-t}.
\tag{B}\\
\frac{10s}{(s-3)(s+1)(s^2+1)}=\frac{A}{ s-3}+\frac{B}{ s+1}+\frac{Cs+D}{ s^2+1}
\end{gather*}
 where
\begin{gather*}
(A(s+1)+B(s-3))(s^2+1)+(Cs+D)(s-3)(s+1)=10s.
\\
\begin{array}{rc rl}
40A&=&30&(\text{when }s=3);\\
-8B&=&-10& (\text{when }s=-1);\\
A-3B-3D&=&0& (\text{when }s=0);\\
A+B+C&=&0&(\text{equate coefficients of }s^3).
\end{array}
\end{gather*}
Solving this system yields $A=3/4$, $B=5/4$,
$C=-2$, $D=-1$. Therefore,
\begin{align*}
\frac{10s}{(s-3)(s+1)(s^2+1)}
&=\frac{3}{4}\frac{1}{ s-3}+\frac{5}{4}\frac{1}{ s+1}-\frac{2s+1}{ s^2+1}
\\&\leftrightarrow
\frac{3}{4}e^{3t}+\frac{5}{4}e^{-t}-2\cos t-\sin t.
\end{align*}
From this, (A), and (B),
 $y=-\sin t-2\cos t+3e^{3t}+e^{-t}$.



\item
\begin{gather*}
(s^2+4)Y(s)=\frac{16}{ s^2+4}+\frac{9s}{ s^2+1}+s,\quad\text{so}
\\
Y(s)=\frac{16}{ (s^2+4)^2}+\frac{9s}{(s^2+4)(s^2+1)}+\frac{s}{
s^2+4}.
\end{gather*}
From the table of Laplace transforms,
\begin{align*}
t\cos2t&\leftrightarrow\frac{s^2-4}{(s^2+4)^2}=\frac{s^2+4}{(s^2+4)^2}
-\frac{8}{(s^2+4)^2}\\
&=\frac{1}{ s^2+4}-\frac{8}{(s^2+4)^2}.
\end{align*}
Therefore,
\[
\frac{8}{(s^2+4)^2}=\frac{1}{ s^2+4}-\laplace (t\cos2t),
\text{ so }
\frac{16}{(s^2+4)^2}\leftrightarrow\sin2t-2t\cos2t.
\tag{A}
\]
Substituting $x=s^2$ into
\[
\frac{9}{(x+4)(x+1)}=\frac{3}{ x+1}-\frac{3}{ x+4}
\]
and multiplying  by $s$ yields
\[
\frac{9s}{(s^2+4)(s^2+1)}=\frac{3s}{ s^2+1}-\frac{3s}{ s^2+4}
\leftrightarrow 3\cos t-3\cos2t.
\tag{B}
\]
Finally,
\[
\frac{s}{ s^2+4}\leftrightarrow\cos2t.
\tag{C}
\]
Adding (A), (B), and (C) yields
$y=-(2t+2)\cos2t+\sin2t+3\cos t$.


\item
\[
(s^2+2s+2)Y(s)=\frac{2}{ s^2}+(-7+2s)+4.
\]
Since $(s^2+2s+2)=(s+1)^2+1$,
\begin{gather*}
Y(s)=\frac{2}{ s^2((s+1)^2+1)}+\frac{2s-3}{ (s+1)^2+1}.
\tag{A}\\
\frac{2s-3}{ (s+1)^2+1}=\frac{2(s+1)-5}{ (s+1)^2+1}
\leftrightarrow e^{-t}(2\cos t-5\sin t).
\tag{B}\\
\frac{2}{ s^2((s+1)^2+1)}=\frac{A}{ s}+\frac{B}{
s^2}+\frac{C(s+1)+D}{(s+1)^2+1},
\end{gather*}
where
$(As+B)\left((s+1)^2+1\right)+s^2\left(C(s+1)+D\right)=2$.
\[
\begin{array}{rc ll}
2B&=&2&(\text{when }s=0);\\
-A+B+D&=&2&(\text{when }s=-1);\\
A+C&=&0&(\text{equate coefficients of }s^3);\\
2A+B+C+D&=&0&(\text{equate coefficients of }s^2).
\end{array}
\]
Solving this system yields $A=-1$, $B=1$, $C=1$, $D=0$. Therefore,
\[
\frac{2}{ s^2((s+1)^2+1)}=-\frac{1}{ s}+\frac{1}{
s^2}+\frac{(s+1)}{(s+1)^2+1}\leftrightarrow-1+t+e^{-t}\cos t.
\]
From this, (A), and (B),
$y=-1+t+e^{-t}(\cos t-5\sin t)$.


\item
$(s^2+4s+5)Y(s)=\dst\frac{(s+1)+3}{ (s+1)^2+1}+4$.
Since $(s^2+4s+5)=(s+2)^2+1$,
\begin{gather*}
Y(s)=\frac{s+4}{ ((s+1)^2+1)((s+2)^2+1)}+\frac{4}{ (s+2)^2+1}.
\tag{A}\\
\frac{4}{ (s+2)^2+1}
\leftrightarrow 4e^{-2t}\sin t.
\tag{B}\\
\frac{s+4}{ ((s+1)^2+1)((s+2)^2+1)}=\frac{A(s+1)+B}{
(s+1)^2+1}+\frac{C(s+2)+D}{ (s+2)^2+1},
\end{gather*}
where
$\left(A(s+1)+B\right)\left((s+2)^2+1\right)+
\left(C(s+2)+D\right)\left((s+1)^2+1\right)=4+s$.
\[
\begin{array}{rc ll}
5A+5B+4C+2D&=&4&(\text{when }s=0);\\
2B+C+D&=&3&(\text{when }s=-1);\\
-A+B+2D&=&2&(\text{when }s=-2);\\
A+C&=&0&(\text{equate coefficients of }s^3).\\
\end{array}
\]
Solving this system yields $A=-1$, $B=1$, $C=1$, $D=0$. Therefore,
\[
\frac{s+4}{ ((s+1)^2+1)((s+2)^2+1)}=\frac{-(s+1)+1}{
(s+1)^2+1}+\frac{s+2}{ (s+2)^2+1},
\leftrightarrow e^{-t}(-\cos t+\sin t)+e^{-2t}\cos t.
\]
From this, (A), and (B),
$y=e^t(-\cos t+\sin t)+e^{-2t}(\cos t+4\sin t)$.



\item
\[
(2s^2-3s-2)Y(s)=\frac{4}{ s-1}+2(-2+s)-3=\frac{4+(2s-7)(s-1)}{ s-1}
\]

Since $2s^2-3s-2=(s-2)(2s+1)$,
\[
Y(s)=\frac{4+(2s-7)(s-1)}{2(s-2)(s-1)(s+1/2)}=
\frac{1}{5}\frac{1}{ s-2}-\frac{4}{3}\frac{1}{ s-1}+\frac{32}{15}\frac{1}{ s+1/2}
\]
and
 $y=\dst{\frac{1}{5}e^{2t}-\frac{4}{3}e^t
+\frac{32}{15}e^{-t/2}}$.


\item
\[
(2s^2+2s+1)Y(s)=\frac{2}{ s^2}+2(-1+s)+2=\frac{2}{ s^2}+2s.
\]
Since $2s^2+2s+1=2((s+1/2)^2+1/4)$,
\begin{gather*}
Y(s)=\frac{1}{ s^2((s+1/2)^2+1/4)}+\frac{s}{((s+1/2)^2+1/4)}.
\tag{A}\\
\frac{s}{((s+1/2)^2+1/4)}\leftrightarrow
e^{-t/2}(\cos(t/2)-\sin(t/2)).
\tag{B}\\
 \frac{1}{ s^2((s+1/2)^2+1/4)}=\frac{A}{ s}+\frac{B}{ s^2}
+\frac{C(s+1/2)+D}{((s+1/2)^2+1/4)}
\end{gather*}
where
\begin{gather*}
(As+B)((s+1/2)^2+1/4)+(C(s+1/2)+D)s^2=1.
\\
\begin{array}{rc rl}
B&=&2&(\text{when }s=0);\\
-A+2B+2D&=&8& (\text{when }s=-1/2);\\
5A+10B+2C+2D&=&8& (\text{when }s=1/2);\\
A+C&=&0&(\text{equate coefficients of }s^3).
\end{array}
\end{gather*}
Solving this system yields $A=-4$, $B=2$,
$C=4$, $D=0$. Therefore,
\begin{align*}
\frac{1}{ s^2((s+1/2)^2+1/4)}
&=
-\frac{4}{ s}+\frac{2}{ s^2}
+\frac{4(s+1/2)}{(s+1/2)^2+1/4}
\\&\leftrightarrow
-4+2t+4e^{-t/2}\cos(t/2).
\end{align*}
This, (A), and (B) imply that
$y=\dst e^{-t/2}(5\cos(t/2)-\sin(t/2))+2t-4$.


\item
\[
(4s^2+4s+1)Y(s)=\frac{3+s}{ s^2+1}+4(-1+2s)+8=
\frac{3+s}{ s^2+1}+4(-1+2s)+8s+4.
\]
Since $4s^2+4s+1=4(s+1/2)^2$,
\begin{gather*}
Y(s)=
\frac{3+s}{4(s+1/2)^2(s^2+1)}+\frac{2}{ s+1/2}.
\tag{A}\\
\frac{3+s}{4(s+1/2)^2(s^2+1)}=\frac{A}{ s+1/2}+\frac{B}{(s+1/2)^2}
+\frac{Cs+D}{ s^2+1}
\end{gather*}
where
\begin{gather*}
(A(s+1/2)+B)(s^2+1)+(Cs+D)(s+1/2)^2=\frac{3+s}{4}.
\\
\begin{array}{rc rl}
10B&=&5&(\text{when }s=-1/2);\\
2A+4B+D&=&3& (\text{when }s=0);\\
12A+8B+9C+9D&=&4& (\text{when }s=1);\\
A+C&=&0&(\text{equate coefficients of }s^3).
\end{array}
\end{gather*}
Solving this system yields $A=3/5$, $B=1/2$,
$C=-3/5$, $D=-1/5$. Therefore,
\begin{align*}
\frac{3+s}{4(s+1/2)^2(s^2+1)}
&=
\frac{3}{5}\frac{1}{ s+1/2}+\frac{1}{2}\frac{1}{(s+1/2)^2}
-\frac{1}{5}\frac{3s+1}{ s^2+1}.
\\&\leftrightarrow
\frac{3}{5}e^{-t/2}+\frac{1}{2}te^{-t/2}-\frac{1}{5}(3\cos t+\sin t).
\end{align*}
Since $\dst\frac{2}{ s+1/2}\leftrightarrow2e^{-t/2}$, this and (A)
imply that
 $y=\dst\frac{e^{-t/2}}{10}(5t+26)-\frac{1}{5}(3\cos
t+\sin t)$.


\item
Transforming the initial value problem
\[
ay''+by'+cy=0,\  y(0)=1,\ y'(0)=0
\]
yields
$(as^2+bs+c)Y(s)=as+b$,
so $Y(s)=\dst\frac{as+b}{ as^2+bs+c}$. Therefore,
$\dst y_1=\laplace ^{-1}\left(\frac{as+b}{ as^2+bs+c}\right)$
satisfies the initial conditions $y_1(0)=1$, $y_1'(0)=0$.

Transforming the initial value problem
\[
ay''+by'+cy=0,\  y(0)=0,\ y'(0)=1
\]
yields
$(as^2+bs+c)Y(s)=a$,
so $Y(s)=\dst\frac{a}{ as^2+bs+c}$. Therefore,
$\dst y_2=\laplace ^{-1}\left(\frac{a}{ as^2+bs+c}\right)$
satisfies the initial conditions $y_1(0)=0$, $y_1'(0)=1$.
\end{sectionSolutions}

\section{The Unit Step Function}

\begin{sectionSolutions}
\item
\[
\laplace (f)  =\int_0^\infty e^{-st}f(t)\, dt
 =\int_0^1 e^{-st}t\,dt+\int_1^\infty e^{-st}\,dt.
\tag{A}
\]
To relate the first term to a Laplace transform  we   add and subtract
$\int_1^\infty e^{-st}t\, dt$
  in   (A) to obtain
\[
\laplace (f)=\int_0^\infty e^{-st}t\,dt+
\int_1^\infty e^{-st}(1-t)\, dt
=\laplace (t)-\int_1^\infty e^{-st}(t-1)\,dt.
\tag{B}
\]
Letting $t=x+1$  in the last integral yields
\[
\int_1^\infty e^{-st}(t-1)\,dt=-\int_0^\infty e^{-s(x+1)}x\,dx
 =e^{-s}\laplace (t).
\]
This and (B) imply that
$\dst\laplace (f)=(1-e^{-s})\laplace (t)
= \frac{1-e^{-s}}{ s^2}$.

Alternatively,
$f(t)=\dst t-u(t-1)(t-1)\leftrightarrow(1-e^{-s}){\cal
L}(t)= \frac{1-e^{-s}}{ s^2}$.



\item
\[
\laplace (f)  =\int_0^\infty e^{-st}f(t)\, dt
 =\int_0^1 e^{-st}\,dt+\int_1^\infty e^{-st}(t+2)\,dt.
\tag{A}
\]
To relate the first term to a Laplace transform  we   add and subtract
$\int_1^\infty e^{-st}\, dt$
  in   (A) to obtain
\[
\laplace (f)=\int_0^\infty e^{-st}\,dt+
\int_1^\infty e^{-st}(t+1)\, dt
=\laplace (t)+\int_1^\infty e^{-st}(t+1)\,dt.
\tag{B}
\]
Letting $t=x+1$  in the last integral yields
\[
\int_1^\infty e^{-st}(t+1)\,dt=\int_0^\infty e^{-s(x+1)}(x+2)\,dx
 =e^{-s}\laplace (t+2).
\]
This and (B) imply that
$\dst\laplace (f)=\laplace (1)+e^{-s}\laplace (t+2)
= \frac{1}{ s}+e^{-s}\left(\frac{1}{ s^2}+\frac{2}{ s}\right)$.

Alternatively,

 $f(t)=\dst1+u(t-1)(t+1)\leftrightarrow \laplace (1)+e^{-s}{\cal
L}(t+2)=
\frac{1}{ s}+e^{-s}\left(\frac{1}{ s^2}+\frac{2}{ s}\right)$.


\item
\[
\laplace (f)  =\int_0^\infty e^{-st}f(t)\, dt
 =\int_0^1 e^{-st}t^2=\laplace (t^2)-\int_1^\infty t^2\,dt.
\tag{A}
\]
Letting $t=x+1$  in the last integral yields
\[
\int_1^\infty e^{-st}t^2\,dt=\int_0^\infty e^{-s(x+1)}(t^2+2t+1)\,dx
 =e^{-s}\laplace (t^2+2t+1).
\]
This and (A) imply that
\[
\laplace (f)=\laplace (t^2)+e^{-s}\laplace  (t^2+2t+1)
=\frac{2}{ s^3}-e^{-s}\left(\frac{2}{ s^3}+\frac{2}{ s^2}+\frac{1}{
s}\right).
\]

Alternatively,
\[
f(t)=t^2\left(1-u(t-1)\right)\leftrightarrow
\laplace (t^2)+e^{-s}\laplace  (t^2+2t+1)=
\frac{2}{ s^3}-e^{-s}\left(\frac{2}{ s^3}+\frac{2}{ s^2}+\frac{1}{ s}\right)
\]


\item
$f(t)=t^2+2+u(t-1)(t-t^2-2)$.
Since $t^2+2\leftrightarrow\dst\frac{2}{ s^3}+\frac{2}{ s}$  and
\begin{align*}
\laplace \left(u(t-1)(t-t^2-2)\right)&=e^{-s}{\cal
L}\left((t+1)-(t+1)^2-2\right)\\&=-e^{-s}{\cal
L}(t^2+t+2)=-e^{-s}\dst\left(\frac{2}{ s^3}+\frac{1}{
s^2}+\frac{2}{ s}\right),
\end{align*}
it follows that $F(s)=\dst{\frac{2}{ s^3} +\frac{2}{
s}-e^{-s}\left(\frac{2}{ s^3}+\frac{1}{ s^2}+\frac{2 }{ s}\right)}$.


\item
$f(t)\dst=e^{-t}+u(t-1)(e^{-2t}-e^{-t})
\leftrightarrow \laplace (e^{-t})+e^{-s}{\cal
L}(e^{-2(t+1)})-e^{-s}\laplace (e^{-t-1})=
 \laplace (e^{-t})+e^{-(s+2)}{\cal
L}(e^{-2t})-e^{-(s+1)}\laplace (e^{-t})=
\frac{1-e^{-(s+1)}}{ s+1}+\frac{e^{-(s+2)}}{ s+2}$.


\item
$f(t)=\dst\left[u(t-1)-u(t-2)\right]t\leftrightarrow
e^{-s}\laplace (t+1)-e^{-2s}\laplace (t+2)=
e^{-s}\left(\frac{1}{ s^2}+\frac{1}{ s}\right)-e^{-2s}\left(\frac{1}{ s^2}
+\frac{2}{ s}\right)$.


\item
\begin{align*}
f(t)&=t-2u(t-1)(t-1)+u(t-2)(t+4)\leftrightarrow
   \frac{1}{ s^2}-2e^{-s}\laplace (t)+e^{-2s}\laplace (t+6)\\
  &=\frac{1}{ s^2}-\frac{2e^{-s}}{ s^2}
+e^{-2s}\left(\frac{1}{ s^2}+\frac{6}{ s}\right).
\end{align*}




\item
$f(t)=\dst2-2u(t-1)t+u(t-3)(5t-2)
\leftrightarrow \laplace (2)-2e^{-s}\laplace (t+1)+e^{-3s}{\cal
L}(5t+13)=
\frac{2}{ s}-e^{-s}\left(\frac{2}{ s^2}+\frac{2}{
s}\right)+e^{-3s}\left(\frac{5}{ s^2}+\frac{13}{ s}\right)$.



\item
$f(t)=\dst(t+1)^2+u(t-1)\left((t+2)^2-(t+1)^2\right)
=t^2+2t+1+u(t-1)(2t+3)
\leftrightarrow
\laplace (t^2+2t+1)+e^{-s}\laplace (2t+5)=
\frac{2}{ s^3 }+\frac{2}{ s^2}+\frac{1}{ s}+e^{-s}\left(\frac{2}{
s^2}+\frac{5}{ s}\right)$.


\item
$\dst\frac{1}{ s(s+1)}=\frac{1}{ s}-\frac{1}{ s+1}\leftrightarrow
 1-e^{-t}
\Rightarrow
e^{-s}\frac{1}{ s(s+1)}\leftrightarrow
u(t-1)\left(1-e^{-(t-1)}\right)=
\dst\begin{cases} 0,&0\le t<1,\\
1-e^{-(t-1)},&t\ge1.\end{cases}$


\item
\begin{gather*}
\frac{3}{ s}-\frac{1}{ s^2}\leftrightarrow 3-t\Rightarrow
e^{-s}\left(\frac{3}{ s}-\frac{1}{ s^2}\right)\leftrightarrow
u(t-1)\left(3-(t-1)\right)=u(t-1)(4-t);
\\
\frac{1}{ s}+\frac{1}{ s^2}\leftrightarrow 1+t\Rightarrow
e^{-3s}\left(\frac{1}{ s}+\frac{1}{ s^2}\right)\leftrightarrow
u(t-3)\left(1+(t-3)\right)=u(t-3)(t-2);
\end{gather*}
therefore
\[
h(t)=2+t+u(t-1)(4-t)+u(t-3)(t-2)
=\begin{cases}
2+t,&0\le t<1,\\
6,&1\le t<3,\\[2\jot] t+4,&t\ge3.
\end{cases}
\]



\item
\[
\frac{1-2s}{ s^2+4s+5}=\frac{5-2(s+2)}{(s+2)^2+1}\leftrightarrow
e^{-2t}(5\sin t-2\cos t);
\]
therefore,
\begin{align*}
h(t)&=u(t-\pi)e^{-2(t-\pi)} \left(5\sin(t-\pi)-2\cos(t-\pi)\right)\\
&=u(t-\pi)e^{-2(t-\pi)} (2\cos t-5\sin t)\\
&=\begin{cases} 0, &0\le t<\pi,\\
e^{-2(t-\pi)}(2\cos t-5\sin t),&t\ge\pi.\end{cases}
\end{align*}


\item
Denote $F(s)=\dst\frac{3(s-3)}{(s+1)(s-2)}-\frac{s+1}{(s-1)(s-2)}$.
Since $\dst\frac{3(s-3)}{(s+1)(s-2)}=\frac{4}{ s+1}-\frac{1}{
s-2}$ and
$\dst\frac{s+1}{(s-1)(s-2)}=\frac{3}{ s-2}-\frac{2}{ s-1}$,
  $F(s)=\dst\frac{4}{ s+1}-\frac{4}{ s-2}+\frac{2}{ s-1}
\leftrightarrow4e^{-t}-4e^{2t}+2e^t$.
Therefore,
\[
e^{-2s}F(s)\leftrightarrow
u(t-2)\left
(4e^{-(t-2)}-4e^{2(t-2)}+2e^{(t-2)}\right)=
\begin{cases} 0,&0\le t<2,\\
4e^{-(t-2)}-4e^{2(t-2)}+2e^{(t-2)},&t\ge2.
\end{cases}
\]


\item
\begin{gather*}
\frac{3}{ s}-\frac{1}{ s^3}\leftrightarrow 3-\frac{t^2}{2}\Rightarrow
e^{-2s}\left(\frac{3}{ s}-\frac{1}{ s^3}\right)\leftrightarrow
u(t-2)\dst{\left(3-\frac{(t-2)^2}{2}\right)}=u(t-2)
\dst{\left(-\frac{t^2}{2}+2t+1\right)};
\\
\frac{1}{ s^2}\leftrightarrow t\Rightarrow
\frac{e^{-4s}}{ s^2}\leftrightarrow u(t-4)(t-4);
\end{gather*}
therefore
\begin{align*}
h(t)&=\dst{1-t^2+u(t-2)\left(-\frac{t^2}{2}+2t+1\right) +u(t-4)(t-4)}\\
&=
\begin{cases} 1-t^2,&0\le t<2\\
-\dst\frac{3t^2}{2} +2t+2,&2\le t<4,\\
-\dst\frac{3t^2}{2}+3t-2,&t\ge 4.
\end{cases}
\end{align*}


\item
Let $T$ be an arbitrary positive number.
Since $\lim_{m\to\infty}t_m=\infty$, only finitely many
members of  $\{t_m\}$ are in $[0,T]$.
Since $f_m$ is continuous on $[t_m,\infty)$ for each $m$,
$f$ is piecewise continuous on $[0,T]$.
If $t_M\le t<t_{M+1}$, then   $u(t-t_m)=1$ if $m\le M$, while
$u(t-t_m)=0$ if $m>M$. Therefore,
$f(t)=\dst f_0(t)+\sum_{m=1}^M(f_m(t)-f_{m-1}(t))=f_M(t)$


\item Since $\dst\sum_{m=0}^\infty e^{-\rho Km}$
converges if $\rho>0$, $\dst\sum_{m=0}^\infty e^{-\rho t_m}$ converges
if $\rho>0$, by the comparison test. Therefore,(C) of
Exercise~8.3.31 holds if $s>s_0+\rho$ if $\rho$ is any positive
number. This implies that it holds if $s>s_0$.


\item
Let $t_m=m$ and $f_m(t)=(-1)^m$, $m=0,1,2,\dots$.
Then $f_m(t)-f_{m-1}(t)=(-1)^m2$, so
$f(t)=\dst1+2\sum_{m=1}^\infty (-1)^mu(t-m)$ and
$F(s)=\dst\frac{1}{ s}\left(1+2\dst\sum_{m=1}^\infty(-1)^me^{-ms}
\right)$
Substituting $x=e^{-s}$ in the identity $\dst\sum_{m=1}^\infty
(-1)^mx^m=-\frac{x}{1+x}$ ($|x|<1$) yields
$F(s)=\dst\frac{1}{ s}\left(1-\frac{2e^{-s}}{1+e^{-s}}\right)=
\frac{1}{ s}\ \frac{1-e^{-s}}{ 1+e^{-s}}$.



\item
Let $t_m=m$ and $f_m(t)=(-1)^mm$, $m=0,1,2,\dots$.
Then $f_m(t)-f_{m-1}(t)=(-1)^m(2m-1)$, so
$f(t)=\dst\sum_{m=1}^\infty(-1)^m(2m-1)u(t-m)$and
$F(s)=\dst\frac{1}{ s}\sum_{m=1}^\infty(-1)^m(2m-1)e^{-ms}$.
Substituting $x=e^{-s}$ in the identities $\dst\sum_{m=1}^\infty
(-1)^mx^m=-\frac{x}{1+x}$ and $\dst\sum_{m=1}^\infty
(-1)^mmx^m=-\frac{x}{(1+x)^2}$
($|x|<1$) yields
$F(s)=\dst\frac{1}{
s}\left[\frac{e^{-s}}{1+e^{-s}}-\frac{2e^{-s}}{(1+e^{-s})^2}\right]=
\frac{1}{ s}\frac{(1-e^s)}{(1+e^s)^2}$.
\end{sectionSolutions}

\section[running={Piecewise Continuous Forcing Functions}]
{Constant Coefficient Equations with Piecewise Continuous Forcing Functions}

\begin{sectionSolutions}
\item
 $y''+y=3+u(t-4)(2t-8),\ y(0)=1,\ y'(0)=0$.
Since
\begin{gather*}
\laplace \left(u(t-4)(2t-8)\right)=e^{-4s}\laplace \left(2(t+4)-8\right)
=e^{-4s}\laplace (2t)=\frac{2e^{-4s}}{ s^2},
\\
(s^2+1)Y(s)=\frac{3}{ s}+\frac{2e^{-4s}}{ s^2}+s.
\end{gather*}
\begin{align*}
Y(s)&=\frac{3}{ s(s^2+1)}+\frac{2e^{-4s}}{ s^2(s^2+1)}+\frac{s}{ s^2+1}\\
&=3\left(\frac{1}{ s}-\frac{s}{ s^2+1}\right)+2e^{-4s}
\left(\frac{1}{ s^2}-\frac{1}{ s^2+1}\right)+\frac{s}{ s^2+1}\\
&=\frac{3}{ s}-\frac{2s}{ s^2+1}+2e^{-4s}\left(\frac{1}{ s^2}-\frac{1}{
s^2+1}\right).
\end{align*}
Since
\[
\frac{1}{ s^2}-\frac{1}{ s^2+1}\leftrightarrow t-\sin t\Rightarrow
e^{-4s}\left(\frac{1}{ s^2}-\frac{1}{ s^2+1}\right)\leftrightarrow
u(t-4)\left(t-4-\sin(t-4)\right),
\]
$y=3-2\cos t+2u(t-4)\left(t-4-\sin(t-4)\right)$.


\item
 $y''-y=e^{2t}+u(t-2)(1-e^{2t}),\ y(0)=3,\ y'(0)=-1$.
Since
\begin{gather*}
\laplace (u(t-2)(1-e^{2t}))=e^{-2s}\laplace (1-e^{2(t+2)})=e^{-2s}
\left(\frac{1}{ s}-\frac{e^4}{ s-2}\right),
\\
(s^2-1)Y(s)=\frac{1}{ s-2}+e^{-2s}\left(\frac{1}{ s}-\frac{e^4}{ s-2}\right)
+(-1+3s).
\end{gather*}
Therefore,
\begin{gather*}
\begin{aligned}
Y(s)&=\frac{1}{(s-1)(s+1)(s-2)}+\frac{3s-1}{(s-1)(s+1)}\\
&\quad+e^{-2s}\left(\frac{1}{s(s-1)(s+1)}-\frac{e^4}{(s-1)(s+1)(s-2)}\right).
\end{aligned}\\
\begin{aligned}
\frac{1}{(s-1)(s+1)(s-2)}&=\dst-\frac{1}{2}\frac{1}{s-1}+\frac{1}{6}\frac{1}{ s+1}
+\frac{1}{3}\frac{1}{ s-2}\\
&\leftrightarrow&-\frac{1}{2}e^t+\frac{1}{6}e^{-t}+\frac{1}{3}e^{2t};
\end{aligned}
\\
\frac{e^{-2s}e^4}{(s-1)(s+1)(s-2)}
\leftrightarrow u(t-2)\left(
-\frac{1}{2}e^{t+2}+\frac{1}{6}e^{-(t-6)}+\frac{1}{3}e^{2t}\right);
\\
\frac{1}{ s(s-1)(s+1)}=-\frac{1}{ s}+\frac{1}{2}\frac{1}{ s-1}
+\frac{1}{2}\frac{1}{ s+1}
\leftrightarrow -1+\frac{1}{2}e^t+\frac{1}{2}e^{-t};
\\
\frac{e^{-2s}}{ s(s-1)(s+1)}
\leftrightarrow
u(t-2)\left(-1+\frac{1}{2}e^{t-2}+\frac{1}{2}e^{-(t-2)}\right);
\\
\frac{3s-1}{(s-1)(s+1)}=\frac{1}{ s-1}+\frac{2}{ s+1}\leftrightarrow
e^t+2e^{-t}.
\end{gather*}
Therefore,
\[
y=\frac{1}{2}e^t+\frac{13}{6}e^{-t}+ \frac{1}{ 3}e^{2t}
+u(t-2)\left(-1+\frac{1}{ 2}e^{t-2}+\frac{1}{2}e^{-(t-2)}+
\frac{1}{2}e^{t+2}-\frac{1}{6} e^{-(t-6)}-\frac{1}{3}e^{2t}\right).
\]



\item
Note that $|\sin t|=\sin t$ if $0\le t<\pi$, while
 $|\sin t|=-\sin t$ if $\pi\le t<2\pi$. Rewrite the initial value
problem as
\[
y''+4y=\sin t-2u(t-\pi)\sin t+u(t-2\pi)\sin t,\ y(0)=-3,\ y'(0)=1.
\]
Since
\[
\laplace \left(u(t-\pi)\sin t\right)=e^{-\pi s}{\cal
L}(\sin(t+\pi))=-e^{-\pi s}\laplace (\sin t)
\]
 and
\begin{gather*}
\laplace \left(u(t-2\pi)\sin t\right)=e^{-2\pi s}{\cal
L}(\sin(t+2\pi))=e^{-2\pi s}\laplace (\sin t),
\\
(s^2+4)Y(s)=\frac{1+2e^{-\pi s}+e^{-2\pi s}}{(s^2+1)}+1-3s,\text{ so }
Y(s)=\frac{1+2e^{-\pi s}+e^{-2\pi
s}}{(s^2+1)(s^2+4)}+\frac{1-3s}{ s^2+4}.
\\
\frac{1}{(s^2+1)(s^2+4)}=\left(\frac{1}{ s^2+1}-\frac{1}{ s^2+4}\right)
\leftrightarrow \frac{1}{3}\sin t-\frac{1}{6}\sin2t;
\end{gather*}
therefore
\begin{align*}
\frac{e^{-\pi s}}{(s^2+1)(s^2+4)}
&\leftrightarrow
u(t-\pi)\left(\frac{1}{3}\sin(t-\pi)-\frac{1}{6}\sin2(t-\pi)\right)\\
&=-u(t-\pi)\left(\frac{1}{3}\sin t+\frac{1}{6}\sin2t\right)
\end{align*}
and
\begin{align*}
\frac{e^{-2\pi s}}{(s^2+1)(s^2+4)}
&\leftrightarrow
u(t-2\pi)\left(\frac{1}{3}\sin(t-2\pi)-\frac{1}{6}\sin2(t-2\pi)\right)\\
&=u(t-2\pi)\left(\frac{1}{3}\sin t-\frac{1}{6}\sin2t\right);
\end{align*}
therefore
\[
y=\frac{1}{3}\sin2t-3\cos2t+\frac{1}{3}\sin t
- 2 u(t-\pi)\left(\frac{1}{3}\sin t+\frac{1}{6}\sin2t\right)
+u(t-2\pi)\left(\frac{1}{3}\sin
t-\frac{1}{6}\sin2t\right).
\]


\item
 $y''+9y=\cos t+u(t-3\pi/2)(\sin t-\cos t),\ y(0)=0,\
y'(0)=0$. Since
\begin{align*}
\laplace \left(u(t-3\pi/2)(\sin t-\cos t)\right)&=
e^{-3\pi s/2}{\cal
L}\left(\sin(t+3\pi/2)-\cos(t+3\pi/2)\right)\\
&\qquad-e^{-3\pi s/2}\laplace (\cos t+\sin t),\\
(s^2+9)Y(s)&=\frac{1}{ s^2+1}-e^{-3\pi s/2}\frac{s+1}{
s^2+1},\qquad\text{so}\\
Y(s)&=\frac{1}{(s^2+1)(s^2+9)}-e^{-3\pi s/2}\frac{s+1}{
(s^2+1)(s^2+9)}.
\\
\frac{1}{(s^2+1)(s^2+9)}&=\frac{1}{8}\left(\frac{1}{ s^2+1}-\frac{1}{
s^2+9}\right)\leftrightarrow\frac{1}{8}\left(\sin
t-\frac{1}{3}\sin3t\right)\quad\text{and}
\\
\frac{s}{(s^2+1)(s^2+9)}&=\frac{1}{8}\left(\frac{s}{ s^2+1}-\frac{s}{
s^2+9}\right)\leftrightarrow\frac{1}{8}\left(\cos
t-\cos3t\right).\\
\frac{s+1}{(s^2+1)(s^2+9)}&=\frac{s+1}{(s^2+1)(s^2+9)}=
\frac{s+1}{8}\left(\frac{s+1}{ s^2+1}-\frac{s+1}{s^2+9}\right)\\
&\leftrightarrow\frac{1}{8}\left(\cos t+\sin
t-\cos3t-\frac{1}{3}\sin3t\right),\qquad\text{so}\\
e^{-3\pi s/2}\frac{s+1}{(s^2+1)(s^2+9)}&\leftrightarrow
\frac{u(t-3\pi/2)}{8}\left(\cos(t-3\pi/2)+\sin(t-3\pi/2)\right.\\
&\left. -\cos3(t-3\pi/2)-\frac{1}{3}\sin3(t-\pi/2)\right)\\
&=\frac{u(t-3\pi/2)}{8}\left(\sin t-\cos t+
\sin3t-\frac{1}{3}\cos3t\right).
\end{align*}
Therefore,
$y=\dst{\frac{1}{ 8}\left(\cos t-\cos3t\right)
-\frac{1}{8} u(t-3\pi/2)\left(\sin t-\cos t+
\sin 3t-\frac{1}{3}\cos 3t\right)}$.


\item
 $y''+y=t-2u(t-\pi)t,\ y(0)=0,\ y'(0)=0$.
Since
\begin{gather*}
\laplace \left(u(t-\pi)t\right)=e^{-\pi s}\laplace \left(t+\pi\right)
=e^{-\pi s}\left(\frac{1}{ s^2}+\frac{\pi}{ s}\right),
\\
(s^2+1)Y(s)=\frac{1}{ s^2}-2e^{-\pi s}\left(\frac{1}{ s^2}+\frac{\pi}{
s}\right);
\end{gather*}
\begin{align*}
Y(s)&= \frac{1}{ s^2(s^2+1)}-2e^{-\pi s}\left(\frac{1}{
s^2(s^2+1)}+\frac{\pi}{ s(s^2+1)}\right)\\
&=\left(\frac{1}{ s^2}-\frac{1}{
s^2+1}\right)-2e^{-\pi s} \left(\frac{1}{ s^2}-\frac{1}{
s^2+1}\right)-2\pi e^{-\pi s}\left(\frac{1}{ s}-\frac{s}{
s^2+1}\right).
\end{align*}
Since
\begin{align*}
\frac{1}{ s^2}-\frac{1}{ s^2+1}&\leftrightarrow t-\sin t\Rightarrow
e^{-\pi s}\left(\frac{1}{ s^2}-\frac{1}{ s^2+1}\right)\\
&\leftrightarrow
u(t-\pi)\left(t-\pi-\sin(t-\pi)\right)=u(t-\pi)(t-\pi+\sin t)
\end{align*}
and
\begin{align*}
\frac{1}{ s}-\frac{s}{ s^2+1}&\leftrightarrow 1-\cos t\Rightarrow
e^{-\pi s}\left(\frac{1}{ s}-\frac{s}{ s^2+1}\right)\\
&\leftrightarrow u(t-\pi)\left(1-\cos(t-\pi)\right)=u(t-\pi)(1+\cos t),
\end{align*}
 $y=t-\sin t-2u(t-\pi)(t+\sin t+\pi\cos t)$.


\item
$y''+y=t-3u(t-2\pi)t,\ y(0)=1,\ y'(0)=2;$
\begin{gather*}
\laplace (u(t-2\pi)t)=e^{-2\pi s}\laplace (t+2\pi)=e^{-2\pi s}
\left(\frac{1}{ s^2}+\frac{2\pi}{ s}\right);
\\
(s^2+1)Y(s)=\frac{1-3e^{-2\pi s}}{ s^2}-\frac{6\pi e^{-2\pi s}}{ s}+2+s;
\\
Y(s)=\frac{1-3e^{-2\pi s}}{ s^2(s^2+1)}-\frac{6\pi e^{-2\pi s}}{ s(s^2+1)}
+\frac{2+s}{ s^2+1};
\\
\frac{1}{ s^2(s^2+1)}=\frac{1}{ s^2}-\frac{1}{ s^2+1}\leftrightarrow
t-\sin t;
\\
\frac{e^{-2\pi s}}{ s^2(s^2+1)}\leftrightarrow
u(t-2\pi)((t-2\pi-\sin(t-2\pi))=u(t-2\pi)(t-2\pi-\sin t);
\\
\frac{1}{ s(s^2+1)}=\frac{1}{ s}-\frac{s}{ s^2+1}\leftrightarrow
1-\cos t;
\\
\frac{e^{-2\pi s}}{ s(s^2+1)}\leftrightarrow u(t-2\pi)
(1-\cos(t-2\pi))=u(t-2\pi)(1-\cos t);
\\
\frac{2+s}{ s^2+1}\leftrightarrow 2\sin t+\cos t;
\end{gather*}
$y=\dst t+\sin t+\cos t-u(t-2\pi)(3t-3\sin t-6\pi\cos t)$.


\item
 $y''-4y'+3y=-1+2u(t-1),\ y(0)=0,\ y'(0)=0$;

$\dst(s^2-4s+3)Y(s)=\frac{-1+2e^{-s}}{ s}$;
$\dst Y(s)=\frac{-1+2e^{-s}}{ s(s-1)(s-3)}$;

$\dst\frac{1}{ s(s-1)(s-3)}=\frac{1}{3s}+\frac{1}{6}\frac{1}{
s-3}-\frac{1}{2}\frac{1}{ s-1}
\leftrightarrow \frac{1}{3}+\frac{1}{6}e^{3t}-\frac{1}{2}e^t$;

$\dst\frac{e^{-s}}{ s(s-1)(s-3)}
\leftrightarrow u(t-1)\left(
\frac{1}{3}+\frac{1}{6}e^{3(t-1)}-\frac{1}{2}e^{t-1}\right)$;

$y=\dst-\frac{1}{3}-\frac{1}{6}e^{3t}
+\frac{1}{2}e^t+u(t-1)\left(\frac{2}{3}+\frac{1}{3}e^{3(t-1)}
-e^{t-1}\right)$.



\item
 $y''+2y'+y=4e^t-4u(t-1)e^t,\ y(0)=0,\ y'(0)=0$. Since
\begin{gather*}
\laplace \left(4u(t-1)e^t\right)=e^{-s}{\cal
L}\left(4e^{(t+1)}\right)=\frac{4e^{-s+1}}{ s-1},
\\
(s^2+2s+1)Y(s)=\frac{4}{ s-1}-\frac{4e^{-s+1}}{ s-1},\quad\text{so}
\\
Y(s)=\frac{4}{(s-1)(s+1)^2}-\frac{4e^{-s+1}}{(s-1)(s+1)^2}.
\\
\frac{1}{(s-1)(s+1)^2}=\frac{A}{ s-1}+\frac{B}{ s+1}+
\frac{C}{(s+1)^2},
\end{gather*}
where
\begin{gather*}
A(s+1)^2+B(s-1)(s+1)+C(s-1)=4.
\\
\begin{array}{rc ll}
A&=&1&(\text{when }s=1);\\
C&=&-2&(\text{when }s=-1);\\
A+B&=&0&(\text{equate coefficients of }s^2).
\end{array}
\end{gather*}
Solving this system yields $A=1$, $B=-1$, $C=-2$. Therefore,
\[
\frac{1}{(s-1)(s+1)^2}=\frac{1}{ s-1}-\frac{1}{ s+1}-
\frac{2}{(s+1)^2}\quad\text{and}
\]
\begin{align*}
y&=e^t-e^{-t}-2te^{-t}-eu(t-1)\left(e^{t-1}-e^{-(t-1)}-
2(t-1)e^{-(t-1)}\right)\\
&=e^t-e^{-t}-2te^{-t}-u(t-1)\left(e^t-e^{-(t-2)}-2(t-1)e^{-(t-2)}\right).
\end{align*}


\item
 $y''-4y'+4y=e^{2t}-2u(t-2)e^{2t},\ y(0)=0,\ y'(0)=-1$.
Since
\begin{gather*}
\laplace \left(u(t-2)e^{2t}\right)=e^{-2s}{\cal
L}\left(e^{2t+4}\right)=\frac{e^{-2s+4}}{ s-2},
\\
(s^2-4s+4)Y(s)=\frac{1}{ s-2}-\frac{2e^{-2s+4}}{ s-2}-1,\quad\text{so}
\\
Y(s)=\frac{1}{(s-2)^3}-\frac{2e^{-2s+4}}{(s-2)^3}-\frac{1}{(s-2)^2}.
\\
\frac{1}{(s-2)^3}\leftrightarrow \frac{t^2e^{2t}}{2}\Rightarrow
\frac{e^{-2s+4}}{(s-2)^3}\leftrightarrow\frac{e^4}{2}u(t-2)e^{2(t-2)}(t-2)^2
=u(t-2)\frac{(t-2)^2e^{2t}}{2};
\end{gather*}
therefore
$y=\dst\frac{t^2e^{2t}}{2}-te^{2t}-u(t-2)(t-2)^2e^{2t}$.


\item
$y''+2y'+2y=1+u(t-2\pi)(t-1)-u(t-3\pi)(t+1),\ y(0)=2,\ y'(0)=-1$;
\begin{gather*}
\laplace (u(t-2\pi)(t-1))=e^{-2\pi s}\laplace ((t+2\pi-1))
=e^{-2\pi s}\left(\frac{1}{ s^2}+\frac{2\pi-1}{ s}\right);
\\
\laplace (u(t-3\pi)(t+1))=e^{-3\pi s}\laplace ((t+3\pi+1))
=e^{-3\pi s}\left(\frac{1}{ s^2}+\frac{3\pi+1}{ s}\right);
\\
(s^2+2s+2)Y(s)=\frac{1}{ s}+
e^{-2\pi s}\left(\frac{1}{ s^2}+\frac{2\pi-1}{ s}\right)
-e^{-3\pi s}\left(\frac{1}{ s^2}+\frac{3\pi+1}{ s}\right)+(-1+2s)+4.
\end{gather*}
Let $G(s)=\dst\frac{1}{ s(s^2+2s+2)}$,
 $H(s)=\dst\frac{1}{ s(s^2+2s+2)}$; then
\[
Y(s)=Y_1(s)+e^{-2\pi s}Y_2(s)-e^{-3\pi s}Y_3(s),\tag{A}
\]
where
\begin{gather*}
Y_1(s)=\dst G(s)+\frac{2s+3}{ s^2+2s+2},\tag{B}\\
Y_2(s)=H(s)+(2\pi-1)G(s),\tag{C}\\
Y_3(s)=H(s)+(3\pi+1)G(s). \tag{D}
\end{gather*}
Let $y_i(t)=\laplace ^{-1}(Y_i(s))$, $(i=1,2,3)$. From (A),
\[
y(t)=y_1(t)+u(t-2\pi)y_2(t-2\pi)-u(t-3\pi)y_3(t-3\pi).\tag{E}
\]

Find $\laplace ^{-1}(G(s))$:
\[
G(s)=\frac{A}{ s}+\frac{B(s+1)+C}{(s+1)^2+1}
\]
where
$A((s+1)^2+1)+(B(s+1)+C)s=1$. Setting $s=0$ yields $A=1/2$;
setting $s=-1$ yields $A-C=1$, so $C=-1/2$; since $A+B=0$
(coefficient of $x^2$), $B=-1/2$. Therefore,
\[
G(s)=\frac{1}{2}\left(\frac{1}{ s}-\frac{(s+1)+1}{(s+1)^2+1)}\right)
\leftrightarrow \frac{1}{2}-\frac{1}{2}e^{-t}(\cos t+\sin t).
\tag{F}
\]


Find $\laplace ^{-1}(H(s))$:
\[
H(s)=\frac{A}{ s}+\frac{B}{ s+2}+\frac{C(s+1)+D}{(s+1)^2+1}
\]
where
$(As+B)((s+1)^2+1)+(C(s+1)+D)s^2=1$.
\[
\begin{array}{rc rl}
2B&=&1&(\text{when }s=0);\\
-A+B+D&=&1& (\text{when }s=-1);\\
5A+5B+2C+D&=&1& (\text{when }s=1);\\
A+B=0&=&0&(\text{equate coefficients of }s^3).
\end{array}
\]
Solving this system yields  $A=-1/2$, $B=1/2$, $C=1/2$, $D=0$;
therefore
\[
H(s)=-\frac{1}{2}\left(\frac{1}{
s}-\frac{1}{ s^2}-\frac{s+1}{(s+1)^2+1)}\right)
\leftrightarrow -\frac{1}{2}(1-t-e^{-t}\cos t). \tag{G}
\]
Since
\[
\frac{2s+3}{ s^2+2s+2}=\frac{2(s+1)+1}{(s+1)^2+1}
\leftrightarrow e^{-t}(2\cos t+\sin t),
\]
(B) and (F)) imply that
\[
y_1(t)=\frac{1}{2}e^{-t}(3\cos t+\sin t)+\frac{1}{2}.
\tag{H}
\]
From (C), (F), and (G),
\[
y_2(t)=\pi-1+\frac{t}{2}+(\pi-1) e^{-t}\cos t-\frac{2\pi-1}{2}e^{-t}\sin
t,
\]
so
\[
y_2(t-2\pi)=-\left(
e^{-(t-2\pi)}\left((\pi-1)\cos t+\frac{2\pi-1}{2}\sin t\right)+
1-\frac{t}{2}\right). \tag{I}
\]
From (D), (F)), and (G),
\[
y_3(t)=\frac{1}{2}\left(-e^{-t}(3\pi\cos t+(3\pi+1)\sin
t+t+3\pi)\right),
\]
so
\[
y_3(t-3\pi)=\frac{1}{2}\left(e^{-(t-3\pi)}(3\pi\cos t+(3\pi+1)\sin
t+t)\right).  \tag{J}
\]
Now (E), (20), (I), and
(J)

$y=\dst\frac{1}{2}e^{-t}(3\cos t+\sin t)+\frac{1}{2}$ imply that

$\dst{-u(t-2\pi)\left(
e^{-(t-2\pi)}\left((\pi-1)\cos t+\frac{2\pi-1}{2}\sin t\right)+
1-\frac{t}{2}\right)}$

$\dst{-\frac{1}{2}u(t-3\pi)\left(e^{-(t-3\pi)}(3\pi\cos
t+(3\pi+1)\sin t
)+t\right)}$.



\item
\begin{alist}
\item
$f(t)=\dst\sum_{n=0}^\infty u(t-n\pi)$; $F(s)=\dst\frac{1}{
s}\sum_{n=0}^\infty e^{-n\pi s}$; $Y(s)=\dst\frac{1}{
s(s^2+1)}\sum_{n=0}^\infty e^{-n\pi s}$;
 $\dst\frac{1}{ s(s^2+1)}=\frac{1}{ s}-\frac{s}{ s^2+1}
\leftrightarrow 1-\cos t$;
 $\dst\frac{e^{-n\pi s}}{ s(s^2+1)}\leftrightarrow
u(t-n\pi)(1-\cos(t-n\pi))=u(t-n\pi)(1-(-1)^n\cos t)$;
$y(t)=\dst
\sum_{n=0}^\infty u(t-n\pi)(1-(-1)^n\cos t)$. If $m\pi\le t<(m+1)\pi$,
$y(t)=\dst
\sum_{n=0}^m (1-(-1)^n\cos t)$. Therefore,


 $y(t)=\dst\begin{cases}2m+1-\cos t,&2m\pi\le
t<(2m+1)\pi\quad(m=0,1,\dots)\\
2m,&(2m-1)\pi\le t<2m\pi\quad(m=1,2,\dots).
\end{cases}$

\item
$f(t)=\dst\sum_{n=0}^\infty u(t-2n\pi)t$;
 $F(s)=\dst\sum_{n=0}^\infty e^{-2n\pi s}\laplace (t+2n\pi s)=
\sum_{n=0}^\infty e^{-2n\pi s}\left(\frac{1}{ s^2}+\frac{2n\pi}{
s}\right)$;
$Y(s)=\dst\sum_{n=0}^\infty e^{-2n\pi s}Y_n(s)$, where
$Y_n(s)=\dst\frac{1}{ s^2(s^2+1)}+\frac{2n\pi}{ s(s^2+1)}
=\frac{1}{ s^2}-\frac{1}{ s^2+1}+\frac{2n\pi}{ s}-\frac{2n\pi}{ s^2+1}
\leftrightarrow y_n(t)=t-\sin t+2n\pi-2n\pi\cos t$.
Since $\cos(t-2n\pi)=\cos t$ and $\sin(t-2n\pi)=\sin t$,
$\dst e^{-2n\pi s}Y_n(s)\leftrightarrow
u(t-2n\pi)y_n(t)=u(t-2n\pi)(t-\sin t-2n\pi\cos t)$; therefore
$y(t)=\dst\sum_{n=0}^\infty u(t-2n\pi)(t-\sin t-2n\pi\cos t)$.
If $2m\pi\le t<2(m+1)\pi$, then
\[
y(t)=\sum_{n=0}^m (t-\sin t-2n\pi\cos t)=
(m+1)(t-\sin t-m\pi\cos t).
\]

\item
$f(t)=\dst1+2\sum_{n=1}^\infty(-1)^n u(t-n\pi)$; $F(s)=\dst\frac{1}{
s}\left(1+2\sum_{n=1}^\infty(-1)^n e^{-n\pi s}\right)$;
$Y(s)=\dst\frac{1}{ s(s^2+1)}\left(1+2\sum_{n=1}^\infty(-1)^n e^{-n\pi
s}\right)$;
 $\dst\frac{1}{ s(s^2+1)}=\frac{1}{ s}-\frac{s}{ s^2+1}
\leftrightarrow 1-\cos t$;
 $\dst\frac{e^{-n\pi s}}{ s(s^2+1)}\leftrightarrow
u(t-n\pi)(1-\cos(t-n\pi))=u(t-n\pi)(1-(-1)^n\cos t)$;
$y(t)=\dst1-\cos t+
2\sum_{n=1}^\infty(-1)^n u(t-n\pi)(1-(-1)^n\cos t)$. If
$m\pi\le t<(m+1)\pi$,
\[
y(t)=1-\cos t+
2\sum_{n=1}^m (-1)^n(1-(-1)^n\cos t)
=(-1)^m-(2m+1)\cos t.
\]

\item
$f(t)=\dst\sum_{n=0}^\infty u(t-n)$; $F(s)=\dst\frac{1}{
s}\sum_{n=0}^\infty e^{-ns}$; $Y(s)=\dst\frac{1}{
s(s^2-1)}\sum_{n=0}^\infty e^{-ns}$;
 $\dst\frac{1}{ s(s^2-1)}=\frac{1}{2}\frac{1}{ s-1}+\frac{1}{2}\frac{1}{
s+1}-\frac{1}{ s}
\leftrightarrow \frac{1}{2}(e^t+e^{-t}-2)$;
 $\dst\frac{e^{-ns}}{ s(s^2-1)}\leftrightarrow
\frac{u(t-n)}{2}\left(e^t+e^{-t}-2\right)$;
$y(t)=\dst\frac{1}{2}
\sum_{n=0}^\infty u(t-n)\left(e^{t-n}+e^{-(t-n)}-2\right)$. If $m\le
t<(m+1)$,
\begin{align*}
y(t)&=\frac{1}{2}
\sum_{n=0}^m\left(e^{t-n}+e^{-(t-n)}-2\right)=
\frac{1}{2}(e^{t-m}+e^{-t})\sum_{n=0}^me^n-m-1\\
&=\frac{1-e^{m+1}}{2(1-e)}(e^{t-m}+e^{-t})-m-1.
\end{align*}

\item
$f(t)=\dst(\sin t+2\cos t)\sum_{n=0}^\infty u(t-2n\pi)$;
$F(s)=\dst\frac{1+2s}{ s^2+1}\sum_{n=0}^\infty e^{-2n\pi s}$;
$Y(s)=\dst\frac{1+2s}{(s^2+1)(s^2+2s+2)}\sum_{n=0}^\infty e^{-2n\pi s}$;
\[
\frac{1+2s}{(s^2+1)(s^2+2s+2)}=\frac{As+B}{
s^2+1}+\frac{C(s+1)+D}{(s+1)^2+1}
\]
where
\begin{gather*}
(As+B)((s+1)^2+1)+(C(s+1)+D)(s^2+1)=1+2s.
\\
\begin{array}{rc rl}
2B+C+D&=&1&(\text{when }s=0);\\
-A+B+2D&=&-1& (\text{when }s=-1);\\
5A+5B+4C+2D&=&3&(\text{when }s=1);\\
A+C&=&0&(\text{equate coefficients of }s^3).
\end{array}
\end{gather*}
Solving this system yields $A=0$, $B=1$,
$C=0$, $D=-1$. Therefore,
\begin{align*}
\frac{1+2s}{(s^2+1)(s^2+2s+2)}
&=\frac{1}{ s^2+1}-\frac{1}{(s+1)^2+1}\\
&\leftrightarrow\left(1-e^{-t}\right)\sin t.
\end{align*}
Since $\sin(t-2n\pi)=\sin t$,
\[
e^{-2n\pi s}\frac{1+2s}{(s^2+1)(s^2+2s+2)}\leftrightarrow
u(t-2n\pi)\left(1-e^{-(t-2n\pi)}\right)\sin t,
\]
so
\[
y(t)=\sin t\sum_{n=0}^\infty u(t-2n\pi)\left(1-e^{-(t-2n\pi)}\right).
\]
If $2m\pi\le t<2(m+1)\pi$,
\[
y(t)=\sin t\sum_{n=0}^m\left(1-e^{-(t-2n\pi)}\right)=
\left(m+1-\left(\frac{1-e^{2(m+1)\pi}}{
1-e^{2\pi}}\right)e^{-t}\right)\sin t.
\]

\item
$f(t)=\dst\sum_{n=0}^\infty u(t-n)$; $F(s)=\dst\frac{1}{
s}\sum_{n=0}^\infty e^{-ns}$; $Y(s)=\dst\frac{1}{ s(s-1)(s-2)}$;
\begin{gather*}
\frac{1}{ s(s-1)(s-2)}=\frac{1}{2s}-\frac{1}{ s-1}+\frac{1}{2}\frac{1}{ s-2}
\leftrightarrow \frac{1}{2}\left(1-2e^t+e^{2t}\right);
\\
\frac{e^{-ns}}{ s(s-1)(s-2)}
\leftrightarrow \frac{1}{2}u(t-n)\left(1-2e^{t-n}+e^{2(t-n)}\right);
\\
y(t)=\frac{1}{2}\sum_{n=0}^\infty
u(t-n)\left(1-2e^{t-n}+e^{2(t-n)}\right).
\end{gather*}
If $m\le t< m+1$,
\begin{align*}
y(t)&=\sum_{n=0}^m\left(1-2e^{t-n}+e^{2(t-n)}\right)
=\frac{m+1}{2}-e^{t-m}\sum_{n=0}^m e^n
 +\frac{1}{2}e^{2(t-m)}\sum_{n=0}^m e^{2n}\\
&=\frac{m+1}{2}-e^{t-m}\frac{1-e^{m+1}}{
1-e}+\frac{1}{2}e^{2(t-m)}\frac{1-e^{2m+2}}{1-e^2}.
\end{align*}
\end{alist}

\item
\begin{alist}
\item
The assumptions imply that
$y''(t)=\dst\frac{f(t)-by'(t)-cy(t)}{ a}$  on $(\alpha,t_0)$
and $(t_0,\beta)$,
$y''(t_0+)=\dst\frac{f(t_0+)-by'(t_0)-cy(t_0)}{ a}$, and
$y''(t_0-)=\dst\frac{f(t_0-)-by'(t_0)-cy(t_0)}{ a}$. This implies the
conclusion.

\item
Since $y''$ has a jump discontinuity at $t_0$, applying
Exercise~8.4.23\part{c} to $y'$  shows that $y'$
is not differentiable at $t_0$. Therefore,$y$ cannot satisfy
(A) on $(\alpha,\beta)$ if $f$ has a jump discontinuity at some
$t_0$ in $(\alpha,\beta)$.
\end{alist}

\item
If $0\le t<t_0$, then $y(t)=z_0(t)$. Therefore,
$y(0)=z_0(0)=k_0$ and $y'(0)=z'(0)=k_1$, and
\[
ay''+by'+cy=az_0''+bz_0'+cz_0=f_0(t)=f(t),\quad 0<t<t_0.
\]

Now suppose that $1\le m\le n$.
For convenience, define $t_{n+1}=\infty$.
If $t_m\le t<t_{m+1}$, then $y(t)=\sum_{k=0}^m z_m(t)$, so
\[
ay''+by+cy=\sum_{k=0}^m(az_k''+bz_k'+cz_k)
=f_0+\sum_{k=1}^m(f_k-f_{k-1})=f_m=f,\quad t_m<t<t_{m+1}.
\]
Thus, $y$ satisfies $ay''+by'+cy=f$ on any open interval that does not
contain any of the points $t_1$, $t_2$,\dots, $t_n$.

Since $z(t_m)=z'(t_m)$ for $m=1,2,\dots$, $y$ and $y'$
are continuous on $[0,\infty)$. Since
$y''(t)=-(by'(t)+cy(t))/a$ if $t\ne t_m$ ($m=1,2,\dots$),
$y''$ has limits from the left at  $t_1,\dots,t_n$.
\end{sectionSolutions}

\section{Convolution}

\begin{sectionSolutions}
\item
\begin{alist}
\item
$\sin at\leftrightarrow\dst\frac{a}{ s^2+a^2}$ and
$\cos bt\leftrightarrow\dst\frac{s}{ s^2+b^2}$, so
   $H(s)=\dst\frac{a s}{ (s^2+a^2)(s^2+b^2)}$.

\item $\dst e^t\leftrightarrow\frac{1}{ s-1}$ and
$\dst\sin at\leftrightarrow\frac{a}{ s^2+a^2}$, so
$H(s)=\dst\frac{a}{(s-1)(s^2+a^2)}$.

\item $\dst\sinh at\leftrightarrow\frac{a}{ s^2-a^2}$ and
$\dst\cosh at\leftrightarrow\frac{1}{ s^2-a^2}$, so
$H(s)=\dst\frac{as}{(s^2-a^2)^2}$.

\item
 $t\sin\omega t\leftrightarrow\dst\frac{2\omega
s}{(s^2+\omega^2)^2}$ and $t\cos\omega
t\leftrightarrow\dst\frac{s^2-\omega^2}{(s^2+\omega^2)^2}$, so
$H(s)=\dst\frac{2\omega s (s^2-\omega^2)}{ (s^2+\omega^2)^4}$.

\item
 $\dst{e^t\int_0^t\sin\omega\tau\cos\omega
(t-\tau)\,d\tau}=\dst{\int_0^t
\left(e^\tau\sin\omega\tau\right)\left(e^{(t-\tau)}\cos\omega
(t-\tau)\right)\,d\tau}$; $e^t\sin\omega
t\leftrightarrow\dst\frac{\omega}{(s-1)^2+\omega^2}$ and $e^t\cos\omega
t\leftrightarrow\dst\frac{s-1}{(s-1)^2+\omega^2}$, so
$H(s)=\dst\frac{(s-1)\omega}{\left((s-1)^2+\omega^2\right)^2}$.

\item
$\dst e^t\int_0^t\tau^2
(t-\tau)e^\tau\,d\tau=\int_0^t\tau^2e^{2\tau}(t-\tau)e^{(t-\tau)}\,d\tau$;
$\dst t^2e^{2t}\leftrightarrow\frac{2}{(s-2)^3}$ and
$\dst te^t\leftrightarrow\frac{1}{(s-1)^2}$, so
 $H(s)=\dst\frac{2}{(s-2)^3 (s-1)^2}$.



\item $\dst{e^{-t}\int_0^t e^{-\tau}\tau\cos\omega
(t-\tau)\,d\tau}=
\int_0^t\tau e^{-2\tau}e^{-(t-\tau)}\cos\omega(t-\tau)\,d\tau$;
$\dst te^{-2t}\leftrightarrow\frac{1}{(s+2)^2}$ and
$\dst e^{-t}\cos\omega
 t\leftrightarrow\frac{s+1}{(s+1)^2+\omega^2}$, so
$H(s)=\dst\frac{s+1}{(s+2)^2\left[(s+1)^2+
\omega^2\right]}$.

\item $\dst{e^t\int_0^t e^{2\tau}\sinh (t-\tau)\,d\tau}=
\dst{\int_0^t e^{3\tau}\left(e^{(t-\tau)}\sinh
(t-\tau)\right)\,d\tau}$; $e^{3t}\leftrightarrow\dst\frac{1}{ s-3}$ and
$e^t\sinh t\leftrightarrow\dst\frac{1}{(s-1)^2-1}$, so
$H(s)=\dst\frac{1}{ (s-3)\left((s-1)^2-1\right)}$.

\item $\dst te^{2t}\leftrightarrow\frac{1}{(s-2)^2}$ and
$\dst\sin2t\leftrightarrow\frac{2}{ s^2+4}$, so
 $H(s)=\dst\frac{2}{(s-2)^2(s^2+4)}$.

\item $\dst t^3\leftrightarrow\frac{6}{ s^4}$  and
$\dst e^t\leftrightarrow\frac{1}{ s-1}$, so
$H(s)=\dst\frac{6}{ s^4(s-1)}$.

\item $\dst t^6\leftrightarrow \frac{6!}{ s^7}$ and
$\dst e^{-t}\sin3t\leftrightarrow\frac{3}{(s+1)^2+9}$, so
$H(s)=\dst{\frac{3\cdot 6!}{ s^7\left[(s+1)^2+
9\right]}}$.

\item $\dst t^2\leftrightarrow\frac{2}{ s^3}$ and
$\dst t^3\leftrightarrow\frac{6}{ s^4}$, so
 $H(s)=\dst\frac{12}{ s^7}$.

\item $\dst t^7\leftrightarrow\frac{7!}{ s^8}$ and
$\dst e^{-t}\sin2t\leftrightarrow\frac{2}{(s+1)^2+4}$, so
  $H(s)=\dst{\frac{2\cdot 7!}{ s^8\left[(s+1)^2+
4\right]}}$.

\item $\dst t^4\leftrightarrow\frac{24}{ s^5}$ and
$\dst \sin2t\leftrightarrow\frac{2}{ s^2+4}$, so
 $H(s)=\dst\frac{48}{ s^5(s^2+4)}$.
\end{alist}

\item
\begin{alist}
\item
$Y(s)=\dst\frac{1}{ s^2}-\dst\frac{Y(s)}{ s^2}$;
$Y(s)\dst{\left(1+\frac{1}{
s^2}\right)}=\dst\frac{1}{ s^2}$; $Y(s)\dst\frac{s^2+1}{
s^2}=\dst\frac{1}{ s^2}$; $Y(s)=\dst\frac{1}{ s^2+1}$, so $y=\sin
t$.

\item
$Y(s)=\dst\frac{1}{ s^2+1}-\dst\frac{2sY(s)}{ s^2+1}$;
$Y(s)\dst{\left(1+\frac{2s}{
s^2+1}\right)}=\dst\frac{1}{ s^2+1}$; $Y(s)\dst\frac{(s+1)^2}{
s^2+1}=\dst\frac{1}{ s^2+1}$; $Y(s)=\dst\frac{1}{(s+1)^2}$, so
$y=te^{-t}$.

\item
$Y(s)=\dst\frac{1}{ s}+\dst\frac{2sY(s)}{ s^2+1}$;
$Y(s)\dst{\left(1-\frac{2s}{
s^2+1}\right)}=\dst\frac{1}{ s}$; $Y(s)\dst\frac{(s-1)^2}{
s^2+1}=\dst\frac{1}{ s}$;
$Y(s)=\dst\frac{(s^2+1)}{ s(s-1)^2}=\dst\frac{A}{ s}+\dst\frac{B}{ s^2}
+\dst\frac{C}{(s-1)^2}$, where $A(s-1)^2+Bs(s-1)+Cs=s^2+1$.
Setting $s=0$ and $s=1$ shows that $A=1$ and $C=2$; equating
coefficients of $s^2$ yields $A+B=1$, so $B=0$.  Therefore,
$Y(s)=\dst\frac{1}{ s}
+\dst\frac{1}{(s-1)^2}$, so $y=1+2te^t$.

\item
$Y(s)=\dst\frac{1}{ s^2}+\dst\frac{Y(s)}{ s+1}$;
$Y(s)\dst{\left(1-\frac{1}{ s+1}\right)}=\dst\frac{1}{ s^2}$;
$Y(s)\dst\left(\frac{s}{ s+1}\right)=\dst\frac{1}{ s^2}$;
$Y(s)=\dst\frac{s+1}{ s^3}
=\dst\frac{1}{ s^2}+\dst\frac{1}{ s^3}$, so
$y=t+\dst\frac{t^2}{ 2}$.

\item
$sY(s)-4=\dst\frac{1}{ s^2}+\dst\frac{sY(s)}{ s^2+1}$;
$Y(s)\dst{\left(s-\frac{s}{
s^2+1}\right)}=4+\dst\frac{1}{ s^2}$; $Y(s)\dst\frac{s^3}{
s^2+1}=\dst\frac{4s^2+1}{ s^2}$;
$Y(s)=\dst\frac{(4s^2+1)(s^2+1)}{ s^5}=\dst\frac{4s^4+5s^2+1}{ s^5}
=\dst\frac{4}{ s}+\dst\frac{5}{ s^3}+\dst\frac{1}{ s^5}$, so
$y=\dst{4+\frac{5}{ 2}t^2+\frac{1}{
24}t^4}$.

\item
$Y(s)=\dst\frac{s-1}{ s^2+1}+\dst\frac{Y(s)}{ s^2+1}$;
$Y(s)\dst{\left(1-\frac{1}{
s^2+1}\right)}=\dst\frac{s-1}{ s^2+1}$; $Y(s)\dst\frac{s^2}{
s^2+1}=\dst\frac{s-1}{ s^2+1}$; $Y(s)=\dst\frac{s-1}{ s^2}
=\dst\frac{1}{ s}-\dst\frac{1}{ s^2}$, so
 $y=1-t$.
\end{alist}

\item Substituting $x=t-\tau$ yields
$\dst\int_0^tf(t-\tau)g(\tau)\,d\tau=-\int_t^0f(x)g(t-x)(-\,dx)=
\int_0^tf(x)g(t-x)\,dx=\int_0^tf(\tau)g(t-\tau)\,d\tau$.


\item
$p(s)Y(s)=F(s)+a(k_1+k_0s)+bk_0$, so
(A) $Y(s)=\dst\frac{F(s)}{ p(s)}+\frac{k_0(as+b)+k_1a}{ p(s)}$.
Since $p(s)=a(s-r_1)(s-r_2)$ and therefore $b=-a(r_1+r_2)$,
(A) can be rewritten as
\begin{gather*}
Y(s)=\frac{F(s)}{ a(s-r_1)(s-r_2)}+\frac{k_0(s-r_1-r_2)}{(s-r_1)(s-r_2)}
+\frac{k_1}{(s-r_1)(s-r_2)}.
\\
\frac{1}{(s-r_1)(s-r_2)}=\frac{1}{ r_2-r_1}\left(\frac{1}{ s-r_2}
-\frac{1}{ s-r_1}\right)\leftrightarrow\frac{e^{r_2t}-e^{r_1t}}{
r_2-r_1},
\end{gather*}
so the convolution theorem implies that
\begin{gather*}
\frac{F(s)}{ a(s-r_1)(s-r_2)}\leftrightarrow
\frac{1}{
a}\int_0^t\frac{e^{r_2\tau}-e^{r_1\tau}}{ r_2-r_1}f(t-\tau)\,d\tau.
\\
\frac{s-r_1-r_2}{(s-r_1)(s-r_2)}=\frac{r_2}{ r_2-r_1}\frac{1}{ s-r_1}
-\frac{r_1}{ r_2-r_1}\frac{1}{ s-r_2}\leftrightarrow
\frac{r_2e^{r_1t}-r_1e^{r_2}t}{ r_2-r_1}.
\end{gather*}
Therefore,
\[
y(t)=\ k_0\frac{r_2e^{r_1t}-r_1e^{r_2t}}{ r_2-r_1}+k_1\frac{e^{r_2t}-e^{r_1t}
}{ r_2-r_1}
+\frac{1}{ a}
a\int_0^t\frac{e^{r_2\tau}-e^{r_1\tau}}{ r_2-r_1}f(t-\tau)\,d\tau.
\]


\item
$p(s)Y(s)=F(s)+a(k_1+k_0s)+bk_0$, so
(A) $Y(s)=\dst\frac{F(s)}{ p(s)}+\frac{k_0(as+b)+k_1a}{ p(s)}$.
Since $p(s)=a(s-\lambda)^2+\omega^2$ and therefore $b=-2a\lambda$,
(A) can be rewritten as
\[
Y(s)=\frac{F(s)}{
a[(s-\lambda)^2+\omega^2}]+\frac{k_0(s-2\lambda)}{(s-\lambda)^2+\omega^2}
+\frac{k_1}{(s-\lambda)^2+\omega^2}.
\]
$\dst\frac{1}{(s-\lambda)^2+\omega^2}\leftrightarrow
\frac{1}{\omega}e^{\lambda t}\sin\omega t$, so the convolution theorem
implies that
\begin{gather*}
\frac{F(s)}{ a[(s-\lambda)^2+\omega^2]}\leftrightarrow
\frac{1}{ a\omega}\int_0^te^{\lambda t}f(t-\tau)\sin\omega\tau\,
d\tau.
\\
\frac{s-2\lambda}{(s-\lambda)^2+\omega^2}=
\frac{(s-\lambda)-\lambda}{(s-\lambda)^2+\omega^2}\leftrightarrow
e^{\lambda t}\left(\cos\omega
t-\frac{\lambda}{\omega}\sin\omega t\right).
\end{gather*}
Therefore,
\[
y(t)=e^{\lambda t}\left[k_0\left(\cos\omega
t-\frac{\lambda}{\omega}\sin\omega
t\right)+\frac{k_1}{\omega}\sin\omega t\right]
+\frac{1}{ a\omega}\int_0^te^{\lambda t}f(t-\tau)\sin\omega\tau\,
d\tau.
\]


\item
\begin{alist}
\item
\begin{gather*}
ay''+by'+cy=f_0(t)+u(t-t_1)(f_1(t)-f_0(t)),\,y(0)=0,\,y'(0)=0;\\
p(s)Y(s)=F_0(s)+\laplace (u(t-t_1)(f_1(t)-f_0(t)))
=F_0(s)+e^{-st_1}\laplace (g);\\
Y(s)=\frac{F_0(s)+e^{-st_1}G(s)}{ p(s)}.
\tag{B}
\end{gather*}

\item Since $F_0(s)\leftrightarrow f_0(t)$, $G(s)\leftrightarrow
g(t)$, and $\dst\frac{1}{ p(s)}\leftrightarrow w(t)$, the convolution
theorem implies that
\[
\frac{F_0(s)}{ p(s)}\leftrightarrow
\int_0^tw(t-\tau)f_0(\tau)\,d\tau\quad\text{and}\quad
\frac {G(s)}{ p(s)}\leftrightarrow
\int_0^tw(t-\tau)g(\tau)\,d\tau.
\]
 Now Theorem~8.4.2 implies that
 $\dst\frac{e^{-st_1}G(s)}{ p(s)}\leftrightarrow
u(t-t_1)\int_0^tw(t-t_1-\tau)g(\tau)\,d\tau$, and (B) implies
that
\[
y(t)=\int_0^t
w(t-\tau)f_0(\tau)\,d\tau+u(t-t_1)\int_0^{t-t_1}
w(t-t_1-\tau)g(\tau)\,d\tau.
\]

\item Let $z_0(t)=\int_0^tw(t-\tau)f_0(\tau)\,d\tau$
and $z_1(t)=\int_0^tw(t-\tau)g(\tau)\,d\tau$.
Then $y(t)=z_0(t)+u(t-t_1)z_1(t-t_1)$. Using Leibniz's rule as in the
solution of Exercise~8.6.11\part{b} shows that
\begin{gather*}
z_0'(t)=\int_0^tw'(t-\tau)f_0(\tau)\,d\tau,\quad
z_1'(t)=\dst\int_0^tw'(t-\tau)g(\tau)\,d\tau,\quad t>0,
\\
z_0''(t)=\frac{f_0(t)}{ a}+\int_0^tw''(t-\tau)f_0(\tau)\,d\tau,
z_1''(t)=\dst\frac{g(t)}{ a}+\int_0^tw''(t-\tau)g(\tau)\,d\tau,\quad
t>0,
\end{gather*}
 if $t>0$, and that
\[
az_0''+bz_0'+cz_0=f_0(t)\quad\text{and}\quad
az_1''+bz_1'+cz_1=f_1(t+t_1)-f_0(t+t_1),\quad  t>0.
\]
This implies the stated conclusion for $y'$ and $y''$  on $(0,t)$ and
$(t,\infty)$, and that $ay''+by'+cy=f(t)$ on these intervals.

\item Since  the functions $z_0(t)$ and  $h(t)=u(t-t_1)z_1(t-t_1)$
are both continuous on $[0,\infty)$ and $h(t)=0$ if $0\le t\le t_1$,
$y$ is continuous on $[0,\infty)$. From \part{c}, $y'$ is continuous
on $[0,t_1)$ and $(t_1,\infty)$, so we need only show that $y'$
is continuous at $t_1$.
For this it suffices to show  that $h'(t_1)=0$.
 Since $h(t_1)=0$ if $t\le t_1$,
(B) $\dst\lim_{t\to t_1-}\frac{h(t)-h(t_1)}{ t-t_1}=0$.
If $t>t_1$, then $h(t)=\dst\int_0^{t-t_1}w(t-t_1\tau)g(\tau)\,d\tau$.
Since $h(t_1)=0$,
\[
\left|\frac{h(t)-h(t_1)}{ t-t_1}\right|\le\int_0^{t-t_1}
|w(t-t_1-\tau)g(\tau)|\,d\tau.
\tag{B}
\]
 Since $g$ is continuous from the right at
$0$, we can choose constants $T>0$  and $M>0$ so that
$|g(\tau)|<M$ if $0\le \tau\le T$. Then (B) implies that
\[
\left|\frac{h(t)-h(t_1)}{ t-t_1}\right|\le M\int_0^{t-t_1}
|w(t-t_1-\tau)|\,d\tau,\quad t_1<t<t_1+T.
\tag{C}
\]
Now suppose $\epsilon>0$. Since $w(0)=0$, we can choose $T_1$
such that $0<T_1<T$ and $|w(x)|<\epsilon/M$ if $0\le x<T_1$.
If $t_1<t<t_1+T_1$ and $0\le\tau\le t-t_1$, then $0\le t-t_1-\tau<T_1$,
so (C) implies that
\[
\left|\frac{h(t)-h(t_1)}{ t-t_1}\right|< \epsilon,\quad
t_1<t<t_1+T.
\]
Therefore,
 $\dst\lim_{t\to t_1+}\frac{h(t)-h(t_1)}{ t-t_1}=0$. This and
(B) imply that $h'(t_1)=0$.
\end{alist}
\end{sectionSolutions}

\section{Constant Coefficient Equations with Impulses}

\begin{sectionSolutions}
\item
 $(s^2+s-2)\hat Y(s)=\dst-\frac{10}{ s+1}+(-9+7s)+7$;
$\hat Y(s)=\dst\frac{-10+(s+1)(7s-2)}{(s-1)(s+2)(s+1)}=\frac{2}{
s+2}+\frac{5}{ s+1}$;
$\hat y=\dst2e^{-2t}+5e^{-t}$;
$\dst\frac{1}{ p(s)}=\frac{1}{(s+2)(s-1)}=\frac{1}{3}\left(\frac{1}{
s-1}-\frac{1}{ s+2}\right)$;
$w=\dst\laplace ^{-1}\left(\frac{1}{ p(s)}\right)=\frac{e^t-e^{-2t}}{3}$;
$y=\dst2e^{-2t}+5e^{-t}+\frac{5}{3}u(t-1)\left(e^{(t-1)}-e^{-2(t-1)}\right)$.


\item
 $(s^2+1)\hat Y(s)=\dst\frac{3}{ s^2+9}-1+s$;

$\hat Y(s)=\dst\frac{3}{(s^2+1)(s^2+9)}+\frac{s-1}{ s^2+1}=
\frac{3}{8}\left(\frac{1}{ s^2+1}-\frac{1}{ s^2+9}\right)+\frac{s-1}{ s^2+1}
=\frac{1}{8}\left(\frac{8s-5}{ s^2+1}-\frac{3}{ s^2+9}\right)$;


$\hat y=\dst\frac{1}{8}(8\cos t-5\sin t-\sin3t)$;
$\dst\frac{1}{ p(s)}=\frac{1}{ s^2+1}$;
$w=\dst\laplace ^{-1}\left(\frac{1}{ p(s)}\right)=\sin t$;
$y=\dst\frac{1}{8}(8\cos t-5\sin t-\sin3t)-2u(t-\pi/2)\cos t$.


\item
 $(s^2-1)\hat Y(s)=\dst\frac{8}{ s}+1-s$;
$\hat Y(s)=\dst\frac{8+s(1-s)}{ s(s-1)(s+1)}=
\frac{4}{ s-1}+\frac{3}{ s+1}-\frac{8}{ s}$;
$\hat y=4e^t+3e^{-t}-8$;
$\dst\frac{1}{ p(s)}=\frac{1}{(s-1)(s+1)}=\frac{1}{2}\frac{1}{ s-1}
-\frac{1}{ s+1}$;
$w=\dst\laplace ^{-1}\left(\frac{1}{ p(s)}\right)=\frac{e^t+e^{-t}}{2}=
\sinh t$;
$y=4e^t+3e^{-t}-8+2u(t-2)\sinh(t-2)$;


\item
 $(s^2+4)\hat Y(s)=\dst\frac{8}{ s-2}+8s$;
(A) $\hat Y(s)=\dst\frac{8}{(s-2)(s^2+4)}+\frac{8s}{ s^2+4}$;
$\dst\frac{8}{(s-2)(s^2+4)}=\frac{A}{ s-2}+\frac{Bs+C}{ s^2+4}$
where $A(s^2+4)+(Bs+C)(s-2)=8$. Setting $s=2$ yields $A=1$;
setting $s=0$ yields $4A-2C=8$, so $C=-2$; $A+B=0$
(coefficient of $x^2$), so $B=-A=-1$; therefore
$\dst\frac{8}{(s-2)(s^2+4)}=\frac{1}{ s-2}-\frac{s+2}{ s^2+4}$, so
(A) implies that
$\hat y=\dst e^{2t}+7\cos2t-\sin2t$;
$\dst\frac{1}{ p(s)}=\frac{1}{ s^2+4}$;
$w=\dst\laplace ^{-1}\left(\frac{1}{ p(s)}\right)=\frac{1}{2}\sin2t$.
Since $\sin(2t-\pi)=-\sin2t$,
$y=\dst e^{2t}+7\cos2t-\sin2t-\frac{1}{2}u(t-\pi/2)\sin2t$.


\item
 $(s^2+2s+1)\hat Y(s)=\dst\frac{1}{ s-1}+(2-s)-2$;
$\hat Y(s)=\dst\frac{1-s(s-1)}{(s-1)(s+1)^2}=
\frac{A}{ s-1}+\frac{B}{ s+1}+\frac{C}{(s+1)^2}$ where
$A(s+1)^2+(B(s+1)+C)(s-1)=1-s(s-1)$. Setting $s=1$
yields $A=1/4$; setting $s=-1$ yields $C=1/2$; since $A+B=-1$
(coefficient of $s^2$), $B=-1-A=-5/4$. Therefore,
$\hat Y(s)=\dst\frac{1}{4}\frac{1}{ s-1}-\frac{5}{4}\frac{1}{ s+1}+\frac{1}{2}
\frac{1}{(s+1)^2}$;
 $\hat y=\dst\frac{1}{4}e^t+\frac{1}{4}e^{-t}(2t-5)$;
$\dst\frac{1}{ p(s)}=\frac{1}{(s+1)^2}$;
$w=\dst\laplace ^{-1}\left(\frac{1}{ p(s)}\right)=te^{-t}$;
 $y=\dst\frac{1}{4}e^t+\frac{1}{4}e^{-t}(2t-5)+2u(t-2)(t-2)e^{-(t-2)}$.


\item
 $(s^2+2s+2)\hat Y(s)=\dst(2-s)-2$;
$Y(s)=\dst\frac{-(s+1)+1}{(s+1)^2+1}$;
$\hat y=\dst e^{-t}(\sin t-\cos t)$;
$\dst\frac{1}{ p(s)}=\frac{1}{(s+1)^2+1}$;
$w=\dst\laplace ^{-1}\left(\frac{1}{ p(s)}\right)=e^{-t}\sin t$.
Since $\sin(t-\pi)=-\sin t$ and $\sin(t-2\pi)=\sin t$,
$y=\dst e^{-t}(\sin t-\cos t)-e^{-(t-\pi)}u(t-\pi)\sin
t-3u(t-2\pi)e^{-(t-2\pi)}\sin t$.


\item
$(2s^2-3s-2)\hat Y(s)=\dst\frac{1}{ s}+2(2-s)+3$;
$\hat Y(s)=\dst\frac{1+s(7-2s)}{2s(s+1/2)(s-2)}=\frac{7}{10}\frac{1}{
s-2}-\frac{6}{5}\frac{1}{ s+1/2}-\frac{1}{2s}$;
$\hat y=\dst\frac{7}{10}e^{2t}-\frac{6}{5}e^{-t/2}-\frac{1}{2}$;
$\dst\frac{1}{
p(s)}=\frac{1}{2(s+1/2)(s-2)}=\frac{1}{5}\left(\frac{1}{
s-2}-\frac{1}{ s+1/2}\right)$;
$w=\dst\laplace ^{-1}\left(\frac{1}{
p(s)}\right)=\frac{1}{5}(e^{2t}-e^{-t/2})$;
$y=\dst\frac{7}{10}e^{2t}-\frac{6}{5}e^{-t/2}-\frac{1}{2}+\frac{1}{5}u(t-2)
\left(e^{2(t-2)}-e^{-(t-2)/2}\right)$;


\item
$(s^2+1)\hat Y(s)=\dst\frac{s}{ s^2+4}-1$;
$\hat Y(s)=\dst\frac{s}{(s^2+1)(s^2+4)}-\frac{1}{ s^2+1}
=\frac{1}{3}\left(\frac{s}{ s^2+1}-\frac{s}{ s^2+4}\right)-\frac{1}{ s^2+1}$;
$\hat y=\dst\frac{1}{3}(\cos t-\cos2t-3\sin t)$;
$\dst\frac{1}{ p(s)}=\frac{1}{ s^2+1}$;
$w=\dst\laplace ^{-1}\left(\frac{1}{ p(s)}\right)=\sin t$.
Since $\sin(t-\pi/2)=-\cos t$ and $\sin(t-\pi)=-\sin t$,
\[
y=\frac{1}{3}(\cos t-\cos2t-3\sin t)-2u(t-\pi/2)\cos t+3u(t-\pi)\sin
t.
\]


\item
$(s^2+2s+1)\hat Y(s)=\dst\frac{1}{ s-1}-1$;
(A) $\hat Y(s)=\dst\frac{1}{(s-1)(s+1)^2}-\frac{1}{(s+1)^2}$;
\[
\frac{1}{(s-1)(s+1)^2}=\frac{A}{ s-1}+\frac{B}{ s+1}+\frac{C}{(s+1)^2}
\]
where
$A(s+1)^2+(B(s+1)+C)(s-1)=1$. Setting $s=1$ yields $A=1/4$;
setting $s=-1$ yields $C=-1/2$; since $A+B=0$ (coefficient of $s^2$),
$B=-A=-1/4$. Therefore,
\[
\frac{1}{(s-1)(s+1)^2}=\frac{1}{4}\frac{1}{ s-1}-\frac{1}{4}\frac{1}{
s+1}-\frac{1}{2}\frac{1}{(s+1)^2}.
\]
This and (A) imply that
\[
\hat Y(s)=\frac{1}{4}\frac{1}{ s-1}-\frac{1}{4}\frac{1}{
s+1}-\frac{3}{2}\frac{1}{(s+1)^2};
\]
$\hat y=\dst\frac{1}{4}\left(e^t-e^{-t}(1+6t)\right)$;
$\dst\frac{1}{ p(s)}=\frac{1}{(s+1)^2}$;
$w=\dst\laplace ^{-1}\left(\frac{1}{ p(s)}\right)=te^{-t}$;

$y=\dst{\frac{1}{4}\left(e^t-e^{-t}(1+6t)\right)-
u(t-1)(t-1)e^{-(t-1)}+2u(t-2)(t-2)e^{-(t-2)}}$.


\item
$y''+4y=1-2u(t-\pi/2)+\delta(t-\pi)-3\delta(t-3\pi/2),\ y(0)=1,\
y'(0)=-1$.
$\dst (s^2+4)\hat Y(s)=\frac{1-2e^{-\pi s/2}}{ s}+s-1$;
$\hat Y(s)=\dst\frac{1-2e^{-\pi s/2}}{ s(s^2+4)}+\frac{s-1}{ s^2+4}$.
Since
$\dst\frac{1}{ s(s^2+4)}=
\frac{1}{4}\left(\frac{1}{ s}-\frac{s}{ s^2+4}\right)$,
$\hat Y(s)=
\dst\frac{1}{4s}+\frac{3}{4}\frac{s}{ s^2+4}-\frac{1}{ s^2+4}
-2e^{-\pi s/2}\left(\frac{1}{ s}-\frac{s}{ s^2+4}\right)$.
$\hat
y=\dst\frac{3}{4}\cos2t-\frac{1}{2}\sin2t+\frac{1}{4}+\frac{1}{4}u(t-\pi/2)
(1+\cos2t)$.
$w=\dst\laplace ^{-1}\left(\frac{1}{ p(s)}\right)=\frac{1}{2}\sin2t$.
Since $\sin2(t-\pi)=\sin2t$ and $\sin2(t-3\pi/2)=-\sin2t$,


$y=\dst\frac{3}{4}\cos2t-\frac{1}{2}\sin2t+\frac{1}{4}+\frac{1}{4}u(t-\pi/2)
(1+\cos2t)+\frac{1}{2}u(t-\pi)\sin2t+\frac{3}{2}u(t-3\pi/2)\sin2t$.

\addtocounter{enumi}{4}
\item%8.7.26
$w(t)=e^{-t}\sin t$;
$f_h(t)=\dst\frac{u(t-t_0)-u(t-t_0-h)}{ h}$;
$(s^2+2s+2)Y_h(s)=\dst\frac{1}{ h}\dst\frac{e^{-st_0}-e^{-s(t_0+h)}}{ s}$;
$Y_h(s)=\dst\frac{1}{ h}\dst\frac{e^{-st_0}-e^{-s(t_0+h)}}{ s(s^2+2s+2)}$;
$\dst\frac{1}{ s(s^2+2s+2)}=-\dst\frac{(s+1)+1}{2\left((s+1)^2+1\right)}
+\dst\frac{1}{2s}\leftrightarrow
\dst\frac{1}{2}\left(1-e^{-t}(\cos t+\sin t)\right)$;
\[
\hspace{-25pt}
y_h(t)=\begin{cases} 0,&0\le t<t_0,\\
\dst\frac{1}{2h}\left[1-e^{-(t-t_0)}(\cos(t-t_0)+\sin(t-t_0))\right],
&\hspace{-10pt}t_0\le t< t_0+h,\\
\dst\frac{e^{-(t-t_0)}}{2h}\left[e^h\left(\cos(t-t_0-h)+\sin(t-t_0-h)\right)
-\cos(t-t_0)-\sin(t-t_0)\right] ,&t\ge t_0+h.\end{cases}
\]

\item
$w(t)=e^{-t}-e^{-2t}$;
$f_h(t)=\dst\frac{u(t-t_0)-u(t-t_0-h)}{ h}$;
$(s^2+3s+2)Y_h(s)=\dst\frac{1}{ h}\dst\frac{e^{-st_0}-e^{-s(t_0+h)}}{ s}$;
$Y_h(s)=\dst\frac{1}{ h}\dst\frac{e^{-st_0}-e^{-s(t_0+h)}}{ s(s+1)(s+2)}$;
$\dst\frac{1}{ s(s+1)(s+2)}=\dst\frac{1}{2(s+2)}-\dst\frac{1}{
s+1}+\dst\frac{1}{2s}\leftrightarrow
\dst\frac{e^{-2t}}{2}-e^{-t}+\dst\frac{1}{2}=\dst\frac{(e^{-t}-1)^2}{2}$;
\[
y_h(t)=
\begin{cases} 0,&0\le t<t_0,\\
\dst\frac{(e^{-(t-t_0)}-1)^2}{2h}, &t_0\le t< t_0+h,\\
\dst\frac{(e^{-(t-t_0)}-1)^2-(e^{-(t-t_0-h)}-1)^2}{2h}, &t\ge t_0+h.
\end{cases}
\]

\item
\begin{alist}
\item
$(s^2-1)\hat Y(s)=1$, so $\hat y=w=\dst\laplace ^{-1}\left(\frac{1}{
s^2-1}\right)=\dst\frac{1}{2}\left(e^t-e^{-t}\right)$;
$y=\dst\hat y+\sum_{k=0}^\infty u(t-k)w(t-k)=
\dst\frac{1}{2}\sum_{k=0}^\infty u(t-k)\left(e^{t-k}-e^{-t-k}\right)$.
If $m\le t< m+1$, then
$y=\dst\frac{1}{2}\sum_{k=0}^m
\left(e^{t-k}-e^{-t-k}\right)=\frac{1}{2}(e^{t-m}-e^{-t})
\sum_{k=0}^m e^k
\frac{e^{m+1}-1}{2(e-1)}(e^{t-m}-e^{-t})$.

\item
$(s^2+1)\hat Y(s)=1$, so $\hat y=w=\dst\laplace ^{-1}\left(\frac{1}{
s^2+1}\right)=\sin t$;
$y=\dst\hat y+\sum_{k=0}^\infty u(t-2k\pi)w(t-2k\pi)=
\sin t\sum_{k=0}^\infty u(t-2k\pi)$.
If $2m\pi\le t< 2(m+1)\pi$, then
$y=(m+1)\sin t$.

\item
$(s^2-3s+2)\hat Y(s)=1$, so $\hat y=w=\dst\laplace ^{-1}\left(\frac{1}{
(s-1)(s-2)}\right)=(e^{2t}-e^t)$;
$y=\dst\hat y+\sum_{k=0}^\infty u(t-k)w(t-k)=
\sum_{k=0}^\infty u(t-k)\left(e^{2(t-k)}-e^{t-k}\right)$.
If $m\le t< m+1$, then
$y=\dst\sum_{k=0}^m
\left(e^{2(t-k)}-e^{t-k}\right)=
e^{2(t-m)}\sum_{k=0}^me^{2k}-e^{t-m}\sum_{k=0}^me^k=
e^{2(t-m)}\frac{e^{2m+2}-1}{
e^2-1}-e^{(t-m)}\frac{e^{m+1}-1}{ e-1}$.

\item
 $w=\dst\laplace ^{-1}\left(\frac{1}{ s^2+1}\right)=\sin t$;
$y=\dst\sum_{k=1}^\infty u(t-k\pi)w(t-k\pi)=
\sin t\sum_{k=1}^\infty(-1)^k u(t-k\pi)$, so
  $y=\dst{\begin{cases}0,&2m\pi\le t<(2m+1)\pi,\\
-\sin t,&(2m+1)\pi\le t<(2m+2)\pi,\end{cases}
\quad
(m=0,1,\dots)}$.
\end{alist}
\end{sectionSolutions}

\chapter{Linear Higher Order Equations}

\section{Introduction to Linear Higher Order Equations}
\begin{sectionSolutions}
\item From Example~9.1.1,
$y=\dst c_1x^2+c_2x^3+\dst\frac{c_3}{ x}$
$y'=\dst2c_1x+3c_2x^2-\dst\frac{c_3}{ x^2}$, and
$y''=\dst2c_1+6c_2x+\dst\frac{2c_3}{ x^3}$, where
\begin{align*}
c_1-c_2-c_3&=4\phantom{,}\\
-2c_1+3c_2-c_3&=-14\phantom{,}\\
2c_1-6c_2-2c_3&=20,
\end{align*}
so $c_1=2$, $c_2=-3$, $c_3=1$, and
$y=\dst2x^2-3x^3+\frac{1}{ x}$.

\item The general solution of $y^{(n)}=0$
can be written as $y(x)=\dst\sum_{m=0}^{n-1}c_m(x-x_0)^m$. Since
$y^{(j)}(x)=\dst\sum_{m=j}^{n-1}m(m-1)\cdots(m-j+1)c_m(x-x_0)^{m-j}$,
$y^{(j)}(x_0)=j!c_j$. Therefore,
$\dst y_i=\frac{(x-x_0)^{i-1}}{(i-1)!},\,1\le i
\le n$.


\item
We omit the verification that the given functions are solutions of the
given equations.

\begin{alist}
\item
The equation is normal on $(-\infty,\infty)$.
$W(x)=\begin{vmatrix}e^x&e^{-x}&xe^{-x}\\
e^x&-e^{-x}&e^{-x}(1-x)\\e^x&e^{-x}&e^{-x}
(x-2)\end{vmatrix}$;
$W(0)=\begin{vmatrix}1&1&0\\1&-1&1\\1&1&-2\end{vmatrix}=4$.
Apply  Theorem~9.1.4.

\item
The equation is normal on $(-\infty,\infty)$.

$W(x)=
\begin{vmatrix}
e^x&e^x\cos 2x&e^x\sin 2x\\
e^x&e^x(\cos 2x-2\sin 2x)&e^x(2\cos2x
+\sin 2x)\\e^x&-e^x(3\cos 2x+4\sin 2x)&e^x(4\cos 2x-3\sin2x)
\end{vmatrix}$;

$W(0)=
\begin{vmatrix}1&1&0\\1&1&2\\1&-3&4\end{vmatrix}=
8$. Apply  Theorem~9.1.4.

\item
The equation is normal on $(-\infty,0)$ and $(0,\infty)$.

$W(x)=
\begin{vmatrix}
e^x&e^{-x}&x\\e^x&-e^{-x}&1\\
e^x&e^{-x}&0
\end{vmatrix}=2x$.
Apply  Theorem~9.1.4.

\item The equation is normal on $(-\infty,0)$ and $(0,\infty)$.

$W(x)=\begin{vmatrix}
e^x/x&e^{-x}/x&1\\
e^x(1/x-1/x^2)&-e^{-x}(x+1)/x^2&0\\
e^x(1/x-2/x^2+2/x^3)&e^{-x}(x^2+2x+2)/x^3&0
\end{vmatrix}=2/x^2$.
Apply  Theorem~9.1.4.

\item
The equation is normal on $(-\infty,\infty)$.
$W(x)=\begin{vmatrix}x&x^2&e^x\\1&2x&e^x\\0&2&e^x\end{vmatrix}$;
$W(0)=\begin{vmatrix}0&0&1\\1&0&1\\0&2&1\end{vmatrix}=2$;
Apply  Theorem~9.1.4.

\item
The equation is normal on $(-\infty,1/2)$ and $(1/2,\infty)$.

$W(x)=\begin{vmatrix}
x&e^x&e^{-x}&e^{2x}\\
1&e^x&-e^{-x}&2e^{2x}\\
0&e^x&e^{-x}&4e^(2x)\\
0&e^x&-e^{-x}&8e^{2x}
\end{vmatrix}
=e^{2x}(12x-6)$.
Apply  Theorem~9.1.4.


\item
The equation is normal on $(-\infty,0)$ and $(0,\infty)$.

$W(x)=\begin{vmatrix}
1&x^2&e^{2x}&e^{-2x}\\
0&2x&2e^{2x}&-2e^{-2x}\\0&2&4e^{2x}&4
e^{-2x}\\0&0&8e^{2x}&-8e^{-2x}
\end{vmatrix}=-128x$.
Apply  Theorem~9.1.4.
\end{alist}


\item
From Abel's formula, (A) $W(x)=\dst
W(\pi/2)\exp\left(-\int_{\pi/4}^x\tan t\,dt\right)$;
$\dst\int_{\pi/4}^x \tan t\,dt=-\ln \cos
t\lims{\pi/4}{x}=-\ln(\sqrt2\cos x)$; therefore (A) implies that
$W(x)=\sqrt2K\cos x$.


\item
\begin{alist}
\item $W(x)=
\begin{vmatrix}1&e^x&e^{-x}\\0&e^x&-e^{-x}\\0&e^x&e^{-x}\end{vmatrix}
=(e^x)(e^{-x})\begin{vmatrix}1&1&1\\0&1&-1\\0&1&1\end{vmatrix}=2$.

\item
\begin{align*}
W(x)&=\begin{vmatrix}
e^x&e^x\sin x&e^x\cos x\\
e^x&e^x(\cos x+\sin x)&e^x(\cos x-\sin x)\\
e^x&2e^x\cos x&-2e^x\sin x\end{vmatrix}=\\
&=e^{3x}\begin{vmatrix}
1&\sin x&\cos x\\
1&\cos x+\sin x&\cos x-\sin x\\
1&2\cos x&-2\sin x\end{vmatrix}\\
&=e^{3x}\begin{vmatrix}
1&\sin x&\phantom{-}\cos x\\
0&\cos x&-\sin x\\
1&2\cos x-\sin x&-2\sin x-\cos x\end{vmatrix}\\
&=e^{3x}\begin{vmatrix}
1&\sin x&\cos x\\
0&\cos x&-\sin x\\
0&-\sin x&-\cos x
\end{vmatrix}=-e^{3x}
\end{align*}

\item
$W(x)=\begin{vmatrix}2&x+1&x^2+2\\0&1&2x\\0&0&2\end{vmatrix}=4$.

\item
\begin{align*}
W(x)&=\begin{vmatrix}
x&x\ln|x|&1/x\\
1&\ln|x|+1&-1/x^2\\
0&1/x&2/x^3\end{vmatrix}
=\begin{vmatrix}
1&\ln|x|&1/x^2\\
1&\ln|x|+1&-1/x^2\\
0&1&2/x^2\end{vmatrix}\\
&=\frac{1}{ x^2}\begin{vmatrix}
1&\ln|x|&1\\
1&\ln|x|+1&-1\\
0&1&2\end{vmatrix}
=\frac{1}{ x^2}\begin{vmatrix}
1&\ln|x|&1\\
0&1&-2\\0&1&2
\end{vmatrix}\\
&=\frac{1}{ x^2}\begin{vmatrix}1&\ln|x|&1\\0&1&-2\\0&0&4\end{vmatrix}
=4/x^2.
\end{align*}

\item
$W(x)=
\begin{vmatrix}
1&x&x^2/2&x^3/3&\cdots&x^n/n!\\
0&1&x&x^2/2&\cdots&x^{n-1}/(n-1)!\\
0&0&1&x&\cdots&x^{n-2}/(n-2)!\\
0&0&0&1&\cdots&x^{n-3}/(n-3)!\\
\vdots&\vdots&\vdots&\vdots&\ddots&\vdots\\
0&0&0&0&\cdots&1
\end{vmatrix}=1$.

\item
$W(x)=\begin{vmatrix}
e^x&e^{-x}&x\\e^x&-e^{-x}&1\\
e^x&e^{-x}&0\end{vmatrix}
=\begin{vmatrix}1&1&x\\1&-1&1
\\1&1&0\end{vmatrix}
=\begin{vmatrix}1&1&x\\0&-2&1-x\\0&0&-x\end{vmatrix}=2x$.

\item
\begin{align*}
W(x)&=\begin{vmatrix}
e^x/x&e^{-x}/x&1\\
e^x/x-e^x/x^2&-e^{-x}/x-e^{-x}/x^2&0\\
e^x/x-2e^x/x^2+2e^x/x^3&e^{-x}/x+2e^{-x}/x^2+2e^{-x}/x^3&0
\end{vmatrix}\\
&=\begin{vmatrix}
1/x&1/x&1\\
1/x-1/x^2&-1/x-1/x^2&0\\
1/x-2/x^2+2/x^3&1/x+2/x^2+2/x^3&0
\end{vmatrix}\\
&=\begin{vmatrix}
1/x-1/x^2&-1/x-1/x^2\\
1/x-2/x^2+2/x^3&1/x+2/x^2+2/x^3
\end{vmatrix}\\
&=\begin{vmatrix}
1/x-1/x^2&-2/x\\
1/x-2/x^2+2/x^3&4/x^2
\end{vmatrix}=2/x^2.
\end{align*}

\item
$W(x)=\begin{vmatrix}
x&x^2&e^x\\
1&2x&e^x\\
0&2&e^x
\end{vmatrix}
=e^x\begin{vmatrix}
x&x^2&1\\
1&2x&1\\
0&2&1
\end{vmatrix}=\dst
e^x\left(
x\begin{vmatrix}2x&1\\2&1\end{vmatrix}-
\begin{vmatrix}x^2&1\\2&1\end{vmatrix}
\right)
=e^x(x^2-2x+2)$.

\item
\begin{align*}
W(x)&=\begin{vmatrix}
x&x^3&1/x&1/x^2\\
1&3x^2&-1/x^2&-2/x^3\\
0&6x&2/x^3&6/x^4\\
0&6&-6/x^4&-24/x^5
\end{vmatrix}
=\begin{vmatrix}
0&-2x^3&2/x&3/x^2\\
1&3x^2&-1/x^2&-2/x^3\\
0&6x&2/x^3&6/x^4\\
0&6&-6/x^4&-24/x^5
\end{vmatrix}\\
&=-\begin{vmatrix}
-2x^3&2/x&3/x^2\\
6x&2/x^3&6/x^4\\
6&-6/x^4&-24/x^5
\end{vmatrix}
=-x^4\begin{vmatrix}
-2&2/x^4&3/x^5\\
6&2/x^4&6/x^5\\
6&-6/x^4&-24/x^5
\end{vmatrix}\\
&=-x^4\begin{vmatrix}
-2&2/x^4&3/x^5\\
0&8/x^4&15/x^5\\
0&0&-15/x^5
\end{vmatrix}=-240/x^5.
\end{align*}

\item
\begin{align*}
W(x)&=\begin{vmatrix}
e^x&e^{-x}&x&e^{2x}\\
e^x&-e^{-x}&1&2e^{2x}\\
e^x&e^{-x}&0&4e^{2x}\\
e^x&-e^{-x}&0&8e^{2x}
\end{vmatrix}
=e^{2x}\begin{vmatrix}
1&1&x&1\\
1&-1&1&2\\
1&1&0&4\\
1&-1&0&8
\end{vmatrix}\\
&=e^{2x}\begin{vmatrix}
1-x&1+x&0&1-2x\\
1&-1\phantom{-}&1&2\\
1&1&0&4\\
1&-1\phantom{-}&0&8
\end{vmatrix}
=-e^{2x}\begin{vmatrix}
1-x&1+x&1-2x\\
1&1&4\\
1&-1\phantom{-}&8
\end{vmatrix}\\
&=-e^{2x}\begin{vmatrix}
2&1+x&1-2x\\
2&1&4\\
0&-1\phantom{-}&8
\end{vmatrix}
=-e^{2x}\begin{vmatrix}
2&1+x&1-2x\\
0&-x&3+2x\\
0&-1\phantom{-}&8
\end{vmatrix}=6e^{2x}(2x-1).
\end{align*}

\item
\begin{align*}
W(x)&=\begin{vmatrix}
e^{2x}&e^{-2x}&1&x^2\\
2e^{2x}&-2e^{-2x}&0&2x\\
4e^{2x}&4e^{-2x}&0&2\\
8e^{2x}&-8e^{-2x}&0&0
\end{vmatrix}=
\begin{vmatrix}
1&1&1&x^2\\
2&-2&0&2x\\
4&4&0&2\\
8&-8&0&0
\end{vmatrix}\\
&=\begin{vmatrix}
2&-2&2x\\
4&4&2\\
8&-8&0
\end{vmatrix}
=\begin{vmatrix}
0&-2&2x\\
8&4&2\\
0&-8&0
\end{vmatrix}=-128x.
\end{align*}
\end{alist}

\item
Let $y$ be an arbitrary solution of $Ly=0$ on $(a,b)$. Since
$\{z_1,\dots,z_n\}$ is a fundamental set of solutions of $Ly=0$
on $(a,b)$, there are constants $c_1$, $c_2$,\dots, $c_n$
such that $y=\dst\sum_{i=1}^nc_iy_i$. Therefore,
$y=\dst\sum_{i=1}^nc_i\sum_{j=1}^n a_{ij}y_j=\sum_{j=1}^n C_jy_j$,
with $C_j=\dst\sum_{i=1}^na_{ij}c_i$.
Hence $\{y_1,\dots,y_n\}$ is a fundamental set of solutions of
$Ly=0$ on $(a,b)$.


\item
Let $y$ be a given solution of $Ly=0$ and
$z=\dst\sum_{j=1}^ny^{(j-1)}(x_0)y_j$ Then
$z^{(r)}(x_0)=y^{(r)}(x_0)$,
$r=0,\dots,n-1$. Since the solution of every initial value problem
is unique (Theorem~9.1.1), $z=y$.


\item
If $\{y_1,y_2,\dots,y_n\}$ is linearly dependent on $(a,b)$
there are constants $c_1,\dots,c_n$, not all zeros, such that
$c_1y_1+c_2y_2+\cdots+c_ny_n$. Let $k$ be the smallest integer
such that $c_k\ne0$.  If $k=1$, then $y_1=\dst\frac{1}{
c_1}(c_2y_2+\cdots + c_ny_n)$; if $1<k<n$, then $y_k=0\cdot y_1+
\cdots+ 0\cdot y_{k-1}+\dst\frac{1}{
c_k}(c_{k+1}y_{k+1}+\cdots+c_ny_n)$;
 if $k=n$, then $y_n=0$, so $y_n=0\cdot y_1+0\cdot
y_2+\cdots+ 0\cdot y_n$.


\item
Since
$F=\sum\pm f_{1i_1}f_{2i_2},\dots,f_{ni_n}$,
\begin{align*}
F'&=\sum\pm f'_{1i_1}f_{2i_2},\dots,f_{ni_n}
+\sum\pm f_{1i_1}f'_{2i_2},\dots,f_{ni_n}+\dots+
\sum\pm f_{1i_1}f_{2i_2},\dots,f'_{ni_n}\\
&=F_1+F_2+\dots+F_n.
\end{align*}


\item
Since $y_j^{(n)}=-\dst\sum_{k=1}^n(P_k/P_0)y_j^{(n-k)}$,
Exercise~9.1.19 implies that
\[
W'=-
\begin{vmatrix}
y_1&y_2&\cdots&y_n\\
y'_1&y'_2&\cdots&y'_n\\
\vdots&\vdots&\ddots&\vdots\\
y_1^{(n-2)}&y_2^{(n-2)}&\cdots&y_n^{(n-2)}\\
\sum_{k=1}^n(P_k/P_0)y_1^{(n-k)}&
\sum_{k=1}^n(P_k/P_0)y_2^{(n-k)}&\cdots&
\sum_{k=1}^n(P_k/P_0)y_n^{(n-k)}
\end{vmatrix},
\]
so Exercise~9.1.17 implies that
\[
W'=-\sum_{k=1}^n\frac{P_k}{ P_0}
\begin{vmatrix}
y_1&y_2&\cdots&y_n\\[2\jot]
y'_1&y'_2&\cdots&y'_n\\[2\jot]
\vdots&\vdots&\ddots&\vdots\\[2\jot]
y_1^{(n-2)}&y_2^{(n-2)}&\cdots&y_n^{(n-2)}\\[2\jot]
y_1^{(n-k)}&
y_2^{(n-k)}&\cdots&
y_n^{(n-k)}
\end{vmatrix}.
\]
However, the determinants on the right each have two identical rows
if $k=2,\dots,n$. Therefore,$W'=-\dst\frac{P_1W}{ P_0}$.  Separating
variables yields $\dst\frac{W'}{ W}=-\frac{P_1}{ P_0}$; hence
$\dst\ln \frac{W(x)}{ W(x_0)}=-\int_{x_0}^x \frac{P_1(t)}{
P_0(t)}\,dt$, which implies Abel's formula.



\item
See the proof of Theorem~5.3.3.


\item
\begin{alist}
\item
\[
P_0(x)=
-\begin{vmatrix}
x&x^2-1&x^2+1\\
1&2x&2x\\0&2&2
\end{vmatrix}=
-\begin{vmatrix}
x&x^2-1&x^2+1\\
1&0&0\\0&2&2
\end{vmatrix}=
\begin{vmatrix}
x^2-1&x^2+1\\2&2
\end{vmatrix}=-4;
\]
$P_1(x)=\begin{vmatrix}x&x^2-1&x^2+1\\1&2x&2x\\0&0&0\end{vmatrix}=0$;
$P_2(x)=-\begin{vmatrix}x&x^2-1&x^2+1\\0&2&2\\0&0&0\end{vmatrix}=0$;
$P_3(x)=\begin{vmatrix}1&2x&2x\\0&2&2\\0&0&0\end{vmatrix}=0$.
Therefore, $-4y'''=0$, which is equivalent
to $y'''=0$.

\item
\begin{gather*}
P_0=-
\begin{vmatrix}
e^x&e^{-x}&x\\
e^x&-e^{-x}&1\\
e^x&e^{-x}&0
\end{vmatrix}=
-\begin{vmatrix}
1&1&x\\1&-1&1\\1&1&0
\end{vmatrix}=
-\begin{vmatrix}
0&0&x\\0&-2&1\\1&1&0
\end{vmatrix}=-2x;
\\
P_1=
\begin{vmatrix}
e^x&e^{-x}&x\\e^x&-e^{-x}&1\\e^x&-e^{-x}&0
\end{vmatrix}=
\begin{vmatrix}
1&1&x\\1&-1&1\\1&-1&0
\end{vmatrix}=
\begin{vmatrix}
0&2&x\\0&0&1\\1&-1&0
\end{vmatrix}=2;
\\
P_2=
-\begin{vmatrix}
e^x&e^{-x}&x\\e^x&e^{-x}&0\\e^x&-e^{-x}&0
\end{vmatrix}=
-\begin{vmatrix}
1&1&x\\1&1&0\\1&-1&0
\end{vmatrix}=
-\begin{vmatrix}
0&0&x\\1&1&0\\0&-2&0
\end{vmatrix}=2x;
\\
P_3=
\begin{vmatrix}
e^x&-e^{-x}&1\\e^x&e^{-x}&0\\e^x&-e^{-x}&0
\end{vmatrix}=
\begin{vmatrix}
1&-1&1\\1&1&0\\1&-1&0
\end{vmatrix}=
\begin{vmatrix}
0&0&1\\1&2&0\\1&-1&0
\end{vmatrix}=-2.
\end{gather*}
Therefore,
$-2xy'''+2y''+2xy'-2y=0$,
which is equivalent to
$xy'''-y''-xy'+y=0$.

\item
\begin{gather*}
P_0(x)=
-\begin{vmatrix}
e^x&xe^{-x}&1\\
e^x&e^{-x}(1-x)&0\\
e^x&e^{-x}(x-2)&0
\end{vmatrix}=
-\begin{vmatrix}
1&x&1\\1&1-x&0\\1&x-2&0
\end{vmatrix}=
-\begin{vmatrix}
1&1-x\\1&x-2
\end{vmatrix}=3-2x;
\\
P_1(x)=
\begin{vmatrix}
e^x&xe^{-x}&1\\
e^x&e^{-x}(1-x)&0\\
e^x&e^{-x}(3-x)&0
\end{vmatrix}=
\begin{vmatrix}
1&x&1\\
1&1-x&0\\
1&3-x&0
\end{vmatrix}=
\begin{vmatrix}
1&1-x\\
1&3-x
\end{vmatrix}=2;
\\
P_2(x)=
-\begin{vmatrix}
e^x&xe^{-x}&1\\
e^x&e^{-x}(x-2)&0\\
e^x&e^{-x}(3-x)&0
\end{vmatrix}=
-\begin{vmatrix}
1&x&1\\
1&x-2&0\\
1&3-x&0
\end{vmatrix}=
-\begin{vmatrix}
1&x-2\\
1&3-x
\end{vmatrix}=2x-5;
\\
P_3(x)=
\begin{vmatrix}
e^x&e^{-x}(1-x)&0\\
e^x&e^{-x}(x-2)&0\\
e^x&e^{-x}(3-x)&0
\end{vmatrix}=0.
\end{gather*}
Therefore,
$(3-2x)y'''+2y''+(2x-5)y'=0$.

\item
\begin{align*}
P_0(x)&=
-\begin{vmatrix}
x&x^2&e^x\\
1&2x&e^x\\
0&2&e^x
\end{vmatrix}=
-e^x\begin{vmatrix}
x&x^2&1\\
1&2x&1\\
0&2&1
\end{vmatrix}=
-e^x\begin{vmatrix}
x&x^2-2&0\\
1&2x-2&0\\
0&2&1
\end{vmatrix}\\
&=-e^x\begin{vmatrix}
x&x^2-2\\
1&2x-2
\end{vmatrix}=-e^x(x^2-2x+2);\\
P_1(x)&=
\begin{vmatrix}
x&x^2&e^x\\
1&2x&e^x\\
0&0&e^x
\end{vmatrix}=
e^x\begin{vmatrix}
x&x^2\\
1&2x
\end{vmatrix}=x^2e^x;
\\
P_2(x)&=
-\begin{vmatrix}
x&x^2&e^x\\
0&2&e^x\\
0&0&e^x
\end{vmatrix}=
-e^x\begin{vmatrix}
x&x^2\\
0&2&
\end{vmatrix}=-2xe^x;
\\
P_3(x)&=
\begin{vmatrix}
1&2x&e^x\\
0&2&e^x\\
0&0&e^x
\end{vmatrix}=
e^x\begin{vmatrix}
1&2x\\
0&2\\
\end{vmatrix}=2e^x.
\end{align*}
Therefore,
$-e^x(x^2-2x+2)y'''+x^2e^xy''-2xe^xy'+2e^xy=0$;
which is equivalent to
 $(x^2-2x+2)y'''-x^2y''+2xy'-2y=0$.

\item
\begin{align*}
P_0(x)&=
-\begin{vmatrix}
x&x^2&1/x\\
1&2x&-1/x^2\\
0&2&2/x^3
\end{vmatrix}=
-x\begin{vmatrix}
1&x&1/x^2\\
1&2x&-1/x^2\\
0&2&2/x^3
\end{vmatrix}=
-x\begin{vmatrix}
1&x&1/x^2\\
0&x&-2/x^2\\
0&2&2/x^3
\end{vmatrix}\\
&=-x\begin{vmatrix}
x&-2/x^2\\
2&2/x^3
\end{vmatrix}=-\frac{6}{ x};\\
P_1(x)&=
\begin{vmatrix}
x&x^2&1/x\\
1&2x&-1/x^2\\
0&0&-6/x^4
\end{vmatrix}=
-\frac{6}{ x^4}\begin{vmatrix}
x&x^2\\
1&2x
\end{vmatrix}=-\frac{6}{ x^2};
\\
P_2(x)&=
-\begin{vmatrix}
x&x^2&1/x\\
0&2&2/x^3\\
0&0&-6/x^4
\end{vmatrix}=
\frac{6}{ x^4}\begin{vmatrix}
x&x^2\\
0&2\\
\end{vmatrix}=\frac{12}{ x^3};
\\
P_3(x)&=
\begin{vmatrix}
1&2x&-1/x^2\\
0&2&2/x^3\\
0&0&-6/x^4
\end{vmatrix}=
-\frac{6}{ x^4}\begin{vmatrix}
1&2x\\
0&2\\
\end{vmatrix}=-\frac{12}{ x^4}.
\end{align*}
Therefore,
$\dst-\frac{6}{ x}y'''-\frac{6}{ x^2}y''+\frac{12}{ x^3}y'-\frac{12}{
x^4}y=0$,
which is equivalent to $x^3y'''+x^2y''-2xy'+2y=0$.

\item
\begin{align*}
P_0(x)&=
-\begin{vmatrix}
x+1&e^x&e^{3x}\\
1&e^x&3e^{3x}\\
0&e^x&9e^{3x}
\end{vmatrix}=
-e^{4x}\begin{vmatrix}
x+1&1&1\\
1&1&3\\
0&1&9
\end{vmatrix}=
-e^{4x}\begin{vmatrix}
x+1&0&-8\\
1&0&-6\\
0&1&9
\end{vmatrix}\\
&=e^{4x}\begin{vmatrix}
x+1&-8\\
1&-6\\
\end{vmatrix}
=2e^{4x}(1-3x);\\
P_1(x)&=
\begin{vmatrix}
x+1&e^x&e^{3x}\\
1&e^x&3e^{3x}\\
0&e^x&27e^{3x}
\end{vmatrix}=
e^{4x}\begin{vmatrix}
x+1&1&1\\
1&1&3\\
0&1&27
\end{vmatrix}=
e^{4x}\begin{vmatrix}
x+1&0&-26\\
1&0&-24\\
0&1&27
\end{vmatrix}\\
&=-e^{4x}\begin{vmatrix}
x+1&-26\\
1&-24\\
\end{vmatrix}=2e^{4x}(12x-1);\\
P_2(x)&=
-\begin{vmatrix}
x+1&e^x&e^{3x}\\
0&e^x&9e^{3x}\\
0&e^x&27e^{3x}
\end{vmatrix}=
-e^{4x}\begin{vmatrix}
x+1&1&1\\
0&1&9\\
0&1&27
\end{vmatrix}\\
&=-e^{4x}\begin{vmatrix}
x+1&1&1\\
0&1&9\\
0&0&18
\end{vmatrix}=-18e^{4x}(x+1);\\
P_3(x)=&
\begin{vmatrix}
1&e^x&3e^{3x}\\
0&e^x&9e^{3x}\\
0&e^x&27e^{3x}
\end{vmatrix}=
e^{4x}\begin{vmatrix}
1&1&3\\
0&1&9\\
0&1&27
\end{vmatrix}=
e^{4x}\begin{vmatrix}
1&1&3\\
0&1&9\\
0&0&18
\end{vmatrix}=18e^{4x}.
\end{align*}
Therefore,
\[
2e^{4x}(1-3x)y'''+2e^{4x}(12x-1)y''-18e^{4x}(x+1)y'+18e^{4x}y=0,
\]
which is equivalent to
\[
(3x-1)y'''-(12x-1)y''+9(x+1)y'-9y=0.
\]

\item
\begin{align*}
P_0(x)&=
\begin{vmatrix}
x&x^3&1/x&1/x^2\\
1&3x^2&-1/x^2&-2/x^3\\
0&6x&2/x^3&6/x^4\\
0&6&-6/x^4&-24/x^5
\end{vmatrix}=
x\begin{vmatrix}
1&x^2&1/x^2&1/x^3\\
1&3x^2&-1/x^2&-2/x^3\\
0&6x&2/x^3&6/x^4\\
0&6&-6/x^4&-24/x^5
\end{vmatrix}\\
&=x\begin{vmatrix}
1&x^2&1/x^2&1/x^3\\
0&2x^2&-2/x^2&-3/x^3\\
0&6x&2/x^3&6/x^4\\
0&6&-6/x^4&-24/x^5
\end{vmatrix}=
x\begin{vmatrix}
2x^2&-2/x^2&-3/x^3\\
6x&2/x^3&6/x^4\\
6&-6/x^4&-24/x^5
\end{vmatrix}\\
&=x^2\begin{vmatrix}
2x&-2/x^3&-3/x^4\\
6x&2/x^3&6/x^4\\
6&-6/x^4&-24/x^5
\end{vmatrix}
=x^2\begin{vmatrix}
2x&-2/x^3&-3/x^4\\
0&8/x^3&15/x^4\\
0&0&-15/x^5
\end{vmatrix}=-\frac{240}{ x^5}\\
P_1(x)&=
-\begin{vmatrix}
x&x^3&1/x&1/x^2\\
1&3x^2&-1/x^2&-2/x^3\\
0&6x&2/x^3&6/x^4\\0&0&24/x^5&120/x^6
\end{vmatrix}=
-x\begin{vmatrix}
1&x^2&1/x^2&1/x^3\\
1&3x^2&-1/x^2&-2/x^3\\
0&6x&2/x^3&6/x^4\\0&0&24/x^5&120/x^6
\end{vmatrix}\\
&=-x\begin{vmatrix}
1&x^2&1/x^2&1/x^3\\
0&2x^2&-2/x^2&-3/x^3\\
0&6x&2/x^3&6/x^4\\0&0&24/x^5&120/x^6
\end{vmatrix}
=-x^2\begin{vmatrix}
2x&-2/x^3&-3/x^4\\
6x&2/x^3&6/x^4\\
0&24/x^5&120/x^6
\end{vmatrix}\\
&=-x^2\begin{vmatrix}
2x&-2/x^3&-3/x^4\\
0&8/x^3&15/x^4\\
0&24/x^5&120/x^6
\end{vmatrix}=
-2x^3\begin{vmatrix}
8/x^3&15/x^4\\
24/x^5&120/x^6
\end{vmatrix}=-\frac{1200}{ x^6};\\
P_2(x)&=
\begin{vmatrix}
x&x^3&1/x&1/x^2\\
1&3x^2&-1/x^2&-2/x^3\\
0&6&-6/x^4&-24/x^5\\
0&0&24/x^5&120/x^6
\end{vmatrix}=
x\begin{vmatrix}
1&x^2&1/x^2&1/x^3\\
1&3x^2&-1/x^2&-2/x^3\\
0&6&-6/x^4&-24/x^5\\
0&0&24/x^5&120/x^6
\end{vmatrix}\\
&=x\begin{vmatrix}
1&x^2&1/x^2&1/x^3\\
0&2x^2&-2/x^2&-3/x^3\\
0&6&-6/x^4&-24/x^5\\
0&0&24/x^5&120/x^6
\end{vmatrix}
=x^3\begin{vmatrix}
2&-2/x^4&-3/x^5\\
6&-6/x^4&-24/x^5\\
0&24/x^5&120/x^6
\end{vmatrix}\\
&=x^3\begin{vmatrix}
2&-2/x^4&-3/x^5\\
0&0&-15/x^5\\
0&24/x^5&120/x^6
\end{vmatrix}
=2x^3\begin{vmatrix}
0&-15/x^5\\
24/x^5&120/x^6
\end{vmatrix}=\frac{720}{ x^7};\\
P_3(x)&=
-\begin{vmatrix}
x&x^3&1/x&1/x^2\\
0&6x&2/x^3&6/x^4\\
0&6&-6/x^4&-24/x^5\\0&0&24/x^5&120/x^6
\end{vmatrix}=
-x^2\begin{vmatrix}
6&2/x^4&6/x^5\\
6&-6/x^4&-24/x^5\\
0&24/x^5&120/x^6
\end{vmatrix}\\
&=-x^2\begin{vmatrix}
6&2/x^4&6/x^5\\
0&-8/x^4&-30/x^5\\
0&24/x^5&120/x^6
\end{vmatrix}
=-6x^2\begin{vmatrix}
-8/x^4&-30/x^5\\
24/x^5&120/x^6
\end{vmatrix}=\frac{1440}{ x^6};\\
P_4(x)&=
\begin{vmatrix}
1&3x^2&-1/x^2&-2/x^3\\
0&6x&2/x^3&6/x^4\\
0&6&-6/x^4&-24/x^5\\0&0&24/x^5&120/x^6
\end{vmatrix}=
x\begin{vmatrix}
6&2/x^4&6/x^5\\
6&-6/x^4&-24/x^5\\
0&0&24/x^5&120/x^6
\end{vmatrix}\\
&=x\begin{vmatrix}
6&2/x^4&6/x^5\\
0&-8/x^4&-30/x^5\\
0&0&24/x^5&120/x^6
\end{vmatrix}
=6x\begin{vmatrix}
-8/x^4&-30/x^5\\
24/x^5&120/x^6
\end{vmatrix}=-\frac{1440}{ x^9}.
\end{align*}
Therefore,
\[
-\frac{240}{ x^5}y^{(4)}-\frac{1200}{ x^6}y'''+\frac{720}{
x^7}y''+\frac{1440}{ x^8}y'-\frac{1440}{ x^9}y=0,
\]
which is equivalent to
$x^4y^{(4)}+5x^3y'''-3x^2y''-6xy'+6y=0$.

\item
\begin{align*}
P_0(x)&=
\begin{vmatrix}
x&x\ln|x|&1/x&x^2\\
1&\ln|x|+1&-1/x^2&2x\\
0&1/x&2/x^3&2\\
0&-1/x^2&-6/x^4&0
\end{vmatrix}=
x\begin{vmatrix}
1&\ln|x|&1/x^2&x\\
1&\ln|x|+1&-1/x^2&2x\\
0&1/x&2/x^3&2\\
0&-1/x^2&-6/x^4&0
\end{vmatrix}\\
&=x\begin{vmatrix}
1&\ln|x|&1/x^2&x\\
0&1&-2/x^2&x\\
0&1/x&2/x^3&2\\
0&-1/x^2&-6/x^4&0
\end{vmatrix}
=x\begin{vmatrix}
1&-2/x^2&x\\
1/x&2/x^3&2\\
-1/x^2&-6/x^4&0
\end{vmatrix}\\
&=x\begin{vmatrix}
1&-2/x^2&x\\
1/x&2/x^3&2\\
-1/x^2&-6/x^4&0
\end{vmatrix}
=x\begin{vmatrix}
1&-2/x^2&x\\
0&4/x^3&1\\
0&-8/x^4&1/x
\end{vmatrix}
=x\begin{vmatrix}
4/x^3&1\\
-8/x^4&1/x
\end{vmatrix}=\frac{12}{ x^3};\\
P_1(x)&=
-\begin{vmatrix}
x&x\ln|x|&1/x&x^2\\
1&\ln|x|+1&-1/x^2&2x\\
0&1/x&2/x^3&2\\
0&2/x^3&24/x^5&0
\end{vmatrix}=
-x\begin{vmatrix}
1&\ln|x|&1/x^2&x\\
1&\ln|x|+1&-1/x^2&2x\\
0&1/x&2/x^3&2\\
0&2/x^3&24/x^5&0
\end{vmatrix}\\
&=-x\begin{vmatrix}
1&\ln|x|&1/x^2&x\\
0&1&-2/x^2&x\\
0&1/x&2/x^3&2\\
0&2/x^3&24/x^5&0
\end{vmatrix}
=-x\begin{vmatrix}
1&-2/x^2&x\\
1/x&2/x^3&2\\
2/x^3&24/x^5&0
\end{vmatrix}\\
&=-x\begin{vmatrix}
1&-2/x^2&x\\
0&4/x^3&1\\
0&28/x^5&-2/x^2
\end{vmatrix}
=-x\begin{vmatrix}
4/x^3&1\\
28/x^5&-2/x^2
\end{vmatrix}=\frac{36}{ x^4};\\
P_2(x)&=
\begin{vmatrix}
x&x\ln|x|&1/x&x^2\\
1&\ln|x|+1&-1/x^2&2x\\
0&-1/x^2&-6/x^4&0\\
0&2/x^3&24/x^5&0
\end{vmatrix}=
x\begin{vmatrix}
1&\ln|x|&1/x^2&x\\
1&\ln|x|+1&-1/x^2&2x\\
0&-1/x^2&-6/x^4&0\\
0&2/x^3&24/x^5&0
\end{vmatrix}\\
&=x\begin{vmatrix}
1&\ln|x|&1/x^2&x\\
0&1&-2/x^2&x\\
0&-1/x^2&-6/x^4&0\\
0&2/x^3&24/x^5&0
\end{vmatrix}=
x\begin{vmatrix}
1&-2/x^2&x\\
-1/x^2&-6/x^4&0\\
2/x^3&24/x^5&0
\end{vmatrix}\\
&=x\begin{vmatrix}
1&-2/x^2&x\\
0&-8/x^4&1/x\\
0&28/x^5&-2/x^2
\end{vmatrix}=
x\begin{vmatrix}
-8/x^4&1/x\\
28/x^5&-2/x^2
\end{vmatrix}=-\frac{12}{ x^5};\\
P_3(x)&=
-\begin{vmatrix}
x&x\ln|x|&1/x&x^2\\
0&1/x&2/x^3&2\\
0&-1/x^2&-6/x^4&0\\
0&2/x^3&24/x^5&0
\end{vmatrix}=
-x\begin{vmatrix}
1/x&2/x^3&2\\
-1/x^2&-6/x^4&0\\
2/x^3&24/x^5&0
\end{vmatrix}\\
&=-x\begin{vmatrix}
1/x&2/x^3&2\\
0&-4/x^4&2/x\\
0&20/x^5&-4/x^2
\end{vmatrix}=
-\begin{vmatrix}
-4/x^4&2/x\\
20/x^5&-4/x^2
\end{vmatrix}=
\frac{24}{ x^6};\\
P_4(x)&=
\begin{vmatrix}
1&\ln|x|+1&-1/x^2&2x\\
0&1/x&2/x^3&2\\
0&-1/x^2&-6/x^4&0\\
0&2/x^3&24/x^5&0
\end{vmatrix}=
\begin{vmatrix}
1/x&2/x^3&2\\
-1/x^2&-6/x^4&0\\
2/x^3&24/x^5&0
\end{vmatrix}\\
&=2\begin{vmatrix}
-1/x^2&-6/x^4\\
2/x^3&24/x^5
\end{vmatrix}=-\frac{24}{ x^7}.
\end{align*}
Therefore,
\[
\frac{12}{ x^3}y^{(4)}+\frac{36}{ x^4}y'''-\frac{12}{ x^5}y''+\frac{24}{
x^6}y'-\frac{24}{ x^7}y=0,
\]
which is equivalent to
 $x^4y^{(4)}+3x^2y'''-x^2y''+2xy'-2y=0$.

\item
\begin{align*}
P_0(x)&=
\begin{vmatrix}
e^x&e^{-x}&x&e^{2x}\\
e^x&-e^{-x}&1&2e^{2x}\\
e^x&e^{-x}&0&4e^{2x}\\
e^x&-e^{-x}&0&8e^{2x}
\end{vmatrix}=
e^{2x}\begin{vmatrix}
1&1&x&1\\
1&-1&1&2\\
1&1&0&4\\
1&-1&0&8
\end{vmatrix}\\
&=e^{2x}\begin{vmatrix}
2&1&x&1\\
0&-1&1&2\\
2&1&0&4\\
0&-1&0&8
\end{vmatrix}=
e^{2x}\begin{vmatrix}
2&1&x&1\\
0&-1&1&2\\
0&0&-x&3\\
0&-1&0&8
\end{vmatrix}\\
&=2e^{2x}\begin{vmatrix}
-1&1&2\\
0&-x&3\\
-1&0&8
\end{vmatrix}
=2e^{2x}\begin{vmatrix}
-1&1&2\\
0&-x&3\\
0&-1&6
\end{vmatrix}\\
&=-2e^{2x}\begin{vmatrix}
-x&3\\
-1&6
\end{vmatrix}=6e^{2x}(2x-1);\\
P_1(x)&=
-\begin{vmatrix}
e^x&e^{-x}&x&e^{2x}\\
e^x&-e^{-x}&1&2e^{2x}\\
e^x&e^{-x}&0&4e^{2x}\\
e^x&e^{-x}&0&16e^{2x}
\end{vmatrix}=
-e^{2x}\begin{vmatrix}
1&1&x&1\\
1&-1&1&2\\
1&1&0&4\\
1&1&0&16
\end{vmatrix}\\
&=-e^{2x}\begin{vmatrix}
0&1&x&1\\
2&-1&1&2\\
0&1&0&4\\
0&1&0&16
\end{vmatrix}
=2e^{2x}\begin{vmatrix}
1&x&1\\
1&0&4\\
1&0&16
\end{vmatrix}\\
&=2e^{2x}\begin{vmatrix}
1&x&1\\
0&-x&3\\
0&-x&15
\end{vmatrix}
=2e^{2x}\begin{vmatrix}
-x&3\\
-x&15
\end{vmatrix}=-24xe^{2x};\\
P_2(x)&=
\begin{vmatrix}
e^x&e^{-x}&x&e^{2x}\\
e^x&-e^{-x}&1&2e^{2x}\\
e^x&-e^{-x}&0&8e^{2x}\\
e^x&e^{-x}&0&16e^{2x}
\end{vmatrix}=
e^{2x}\begin{vmatrix}
1&1&x&1\\
1&-1&1&2\\
1&-1&0&8\\
1&1&0&16
\end{vmatrix}\\
&=e^{2x}\begin{vmatrix}
0&1&x&1\\
2&-1&1&2\\
2&-1&0&8\\
0&1&0&16
\end{vmatrix}
=e^{2x}\begin{vmatrix}
0&1&x&1\\
2&-1&1&2\\
0&0&-1&6\\
0&1&0&16
\end{vmatrix}
=-2e^{2x}\begin{vmatrix}
1&x&1\\
0&-1&6\\
1&0&16
\end{vmatrix}\\
&=-2e^{2x}\begin{vmatrix}
1&x&1\\
0&-1&6\\
0&-x&15
\end{vmatrix}
=-2e^{2x}\begin{vmatrix}
-1&6\\
-x&15
\end{vmatrix}=6e^{2x}(5-2x);\\
P_3(x)&=
-\begin{vmatrix}
e^x&e^{-x}&x&e^{2x}\\
e^x&e^{-x}&0&4e^{2x}\\
e^x&-e^{-x}&0&8e^{2x}\\
e^x&e^{-x}&0&16e^{2x}
\end{vmatrix}=
-e^{2x}\begin{vmatrix}
1&1&x&1\\
1&1&0&4\\
1&-1&0&8\\
1&1&0&16
\end{vmatrix}\\
&=-e^{2x}\begin{vmatrix}
0&1&x&1\\
0&1&0&4\\
2&-1&0&8\\
0&1&0&16
\end{vmatrix}
=-2e^{2x}\begin{vmatrix}
1&x&1\\
0&-x&3\\
0&-x&15
\end{vmatrix}
=-2e^{2x}\begin{vmatrix}
-x&3\\
-x&15
\end{vmatrix}=24xe^{2x};\\
P_4(x)&=
\begin{vmatrix}
e^x&-e^{-x}&1&2e^{2x}\\
e^x&e^{-x}&0&4e^{2x}\\
e^x&-e^{-x}&0&8e^{2x}\\
e^x&e^{-x}&0&16e^{2x}
\end{vmatrix}=
e^{2x}\begin{vmatrix}
1&-1&1&2\\
1&1&0&4\\
1&-1&0&8\\
1&1&0&16
\end{vmatrix}\\
&=e^{2x}\begin{vmatrix}
0&-1&1&2\\
2&1&0&4\\
0&-1&0&8\\
2&1&0&16
\end{vmatrix}
=e^{2x}\begin{vmatrix}
0&-1&1&2\\
2&1&0&4\\
0&-1&0&8\\
0&0&0&12
\end{vmatrix}\\
&=-2e^{2x}\begin{vmatrix}
-1&1&2\\
-1&0&8\\
0&0&12
\end{vmatrix}
=-2e^{2x}\begin{vmatrix}
-1&1&2\\
0&-1&6\\
0&0&12
\end{vmatrix}=-24e^{2x}.
\end{align*}
Therefore,
\[
6e^{2x}(2x-1)y^{(4)}-24xe^{2x}y'''
+6e^{2x}(5-2x)y''+24xe^{2x}y'-24e^{2x}y=0,
\]
which is equivalent to
 $(2x-1)y^{(4)}-4xy'''+(5-2x)y''+4xy'-4y=0$.

% 9.1.24
\item
\begin{align*}
P_0(x)&=
\begin{vmatrix}
e^{2x}&e^{-2x}&1&x^2\\
2e^{2x}&-2e^{-2x}&0&2x\\
4e^{2x}&4e^{-2x}&0&2\\
8e^{2x}&-8e^{-2x}&0&0
\end{vmatrix}=
\begin{vmatrix}
1&1&1&x^2\\
2&-2&0&2x\\
4&4&0&2\\
8&-8&0&0
\end{vmatrix}\\
&=\begin{vmatrix}
2&-2&2x\\
4&4&2\\
8&-8&0
\end{vmatrix}
=\begin{vmatrix}
2&-2&2x\\
0&8&2-4x\\
0&0&-8x
\end{vmatrix}=-128x\\
P_1(x)&=
-\begin{vmatrix}
e^{2x}&e^{-2x}&1&x^2\\
2e^{2x}&-2e^{-2x}&0&2x\\
4e^{2x}&4e^{-2x}&0&2\\
16e^{2x}&16e^{-2x}&0&0
\end{vmatrix}=
-\begin{vmatrix}
1&1&1&x^2\\
2&-2&0&2x\\
4&4&0&2\\
16&16&0&0
\end{vmatrix}\\
&=-\begin{vmatrix}
2&-2&2x\\
4&4&2\\
16&16&0
\end{vmatrix}
=-\begin{vmatrix}
2&-2&2x\\
0&8&2-4x\\
0&0&-8
\end{vmatrix}=128;\\
P_2(x)&=
\begin{vmatrix}
e^{2x}&e^{-2x}&1&x^2\\
2e^{2x}&-2e^{-2x}&0&2x\\
8e^{2x}&-8e^{-2x}&0&0\\
16e^{2x}&16e^{-2x}&0&0
\end{vmatrix}=
\begin{vmatrix}
1&1&1&x^2\\
2&-2&0&2x\\
8&-8&0&0\\
16&16&0&0
\end{vmatrix}\\
&=\begin{vmatrix}
2&-2&2x\\
8&-8&0\\
16&16&0
\end{vmatrix}
=2x\begin{vmatrix}
8&-8\\
16&16
\end{vmatrix}=512x;\\
P_3(x)&=
-\begin{vmatrix}
e^{2x}&e^{-2x}&1&x^2\\
4e^{2x}&4e^{-2x}&0&2\\
8e^{2x}&-8e^{-2x}&0&0\\
16e^{2x}&16e^{-2x}&0&0
\end{vmatrix}=
-\begin{vmatrix}
1&1&1&x^2\\
4&4&0&2\\
8&-8&0&0\\
16&16&0&0
\end{vmatrix}\\
&=-\begin{vmatrix}
4&4&2\\
8&-8&0\\
16&16&0
\end{vmatrix}
=-2\begin{vmatrix}
8&-8\\
16&16
\end{vmatrix}=-512;\\
P_4(x)&=
\begin{vmatrix}
2e^{2x}&-2e^{-2x}&0&2x\\
4e^{2x}&4e^{-2x}&0&2\\
8e^{2x}&-8e^{-2x}&0&0\\
16e^{2x}&16e^{-2x}&0&0
\end{vmatrix}=0.
\end{align*}
Therefore,
$-128xy^{(4)}+128y'''+512xy''-512y=0$,
which is equivalent to
 $xy^{(4)}-y'''-4xy''+4y'=0$.
\end{alist}
\end{sectionSolutions}



\section{Higher Order Constant Coefficient Homogeneous Equations}
\begin{sectionSolutions}
\item
$p(r)=r^4+8r^2-9=(r-1)(r+1)(r^2+9)$;
 $y=c_1e^x+c_2e^{-x}+c_3\cos3x+c_4\sin3x$.




\item
$p(r)=2r^3+3r^2-2r-3=(r-1)(r+1)(2r+3)$;
$y=c_1e^x+c_2e^{-x}+c_3e^{-3x/2}$.


\item
$p(r)=4r^3-8r^2+5r-1=(r-1)(2r-1)^2$;
$y=c_1e^x+e^{x/2}(c_2+c_3x)$.

\item
$p(r)=r^4+r^2=r^2(r^2+1)$;
 $y=c_1+c_2x+c_3\cos x+c_4\sin x$.


\item
$p(r)=r^4+12r^2+36=(r^2+6)^2$;
 $y=(c_1+c_2x)\cos\sqrt{6} x+(c_3+c_4x)\sin\sqrt{6} x$.


\item
$p(r)=6r^4+5r^3+7r^2+5r+1=(2r+1)(3r+1)(r^2+1)$;
$y=c_1e^{-x/2}+c_2e^{-x/3}+c_3\cos x+c_4\sin x$.


\item
$p(r)=r^4-4r^3+7r^2-6r+2=(r-1)^2(r^2-2r+2)$;
$y=e^x(c_1+c_2x+c_3\cos x+c_4\sin x)$.


\item
$p(r)=r^3+3r^2-r-3=(r-1)(r+1)(r+3)$;
\[
\begin{aligned}
y&=c_1e^x+c_2e^{-x}+c_3e^{-3x}\\
y'&=c_1e^x-c_2e^{-x}-3c_3e^{-3x}\\
y''&=c_1e^x+c_2e^{-x}+9c_3e^{-3x}
\end{aligned};\quad
\begin{aligned}
c_1+c_2+c_3&=0\\
c_1-c_2-3c_3&=14\\
c_1+c_2+9c_3&=-40
\end{aligned};
\]
$c_1=2$, $c_2=3$, $c_3=-5$;
 $y=2e^x+3e^{-x}-5e^{-3x}$.


\item
$p(r)=r^3-2r-4=(r-2)(r^2+2r+2)$;
\[
\begin{aligned}
y&=e^{-x}(c_1\cos x+c_2\sin x)+c_3e^{2x}\\
y'&=-e^{-x}((c_1-c_2)\cos x+(c_1+c_2)\sin x)+2c_3e^{2x}\\
y''&=e^{-x}(2c_1\sin x-2c_2\cos x)+4c_3e^{2x}
\end{aligned};
\begin{aligned}
c_1+c_3&=6\\
-c_1+c_2+2c_3&=3\\
-2c_2+4c_3&=22
\end{aligned};
\]
$c_1=2$, $c_2=-3$, $c_3=4$;
 $y=2e^{-x}\cos x-3e^{-x}\sin x+4e^{2x}$.


\item
$p(r)=r^3-6r^2+12r-8=(r-2)^3$;
\[
\begin{aligned}
y&=e^{2x}(c_1+c_2x+c_3x^2)\\
y'&=e^{2x}(2c_1+c_2+(2c_2+2c_3)x+2c_3x^2)\\
y''&=2e^{2x}(2c_1+2c_2+c_3+2(c_2+2c_3)x+2c_3x^2)
\end{aligned};
\begin{aligned}
c_1&=1\\
2c_1+c_2&=-1\\
4c_1+4c_2+2c_3&=-4
\end{aligned}
\]
$c_1=1$, $c_2=-3$, $c_3=2$;
 $y=e^{2x}(1-3x+2x^2)$.


\item
$p(r)=8r^3-4r^2-2r+1=(2r+1)(2r-1)^2$;
\[
\begin{aligned}
y&=e^{x/2}(c_1+c_2x)+c_3e^{-x/2}\\
y'&=\frac{1}{2}e^{x/2}(c_1+2c_2+c_2x)-\frac{1}{2}c_3e^{-x/2}\\
y''&=\frac{1}{4}e^{x/2}(c_1+4c_2+c_2x)+\frac{1}{4}c_3e^{-x/2}
\end{aligned};
\begin{aligned}
c_1+c_3&=4\\
\frac{1}{2}c_1+c_2-\frac{1}{2}c_3&=-3\\
\frac{1}{4}c_1+c_2+\frac{1}{4}c_3&=-1
\end{aligned};
\]
$c_1=1$, $c_2=-2$, $c_3=3$;
 $y=e^{x/2}(1-2x)+3e^{-x/2}$.


\item
$p(r)=r^4-6r^3+7r^2+6r-8=(r-1)(r-2)(r-4)(r+1)$;
\[
\begin{aligned}
y&=c_1e^x+c_2e^{2x}+c_3e^{4x}+c_4e^{-x}\\
y'&=c_1e^x+2c_2e^{2x}+4c_3e^{4x}-c_4e^{-x}\\
y''&=c_1e^x+4c_2e^{2x}+16c_3e^{4x}+c_4e^{-x}\\
y'''&=c_1e^x+8c_2e^{2x}+64c_3e^{4x}-c_4e^{-x}
\end{aligned};
\begin{aligned}
c_1+c_2+c_3+c_4&=-2\\
c_1+2c_2+4c_3-c_4&=-8\\
c_1+4c_2+16c_3+c_4&=-14\\
c_1+8c_2+64c_3-c_4&=-62
\end{aligned};
\]
$c_1=-4$, $c_2=1$, $c_3=-1$, $c_4=2$;
 $y=-4e^x+e^{2x}-e^{4x}+2e^{-x}$.


\item
$p(r)=r^4+2r^3-2r^2-8r-8=(r-2)(r+2)(r^2+2r+2)$;
\begin{align*}
y&=c_1e^{2x}+c_2e^{-2x}+e^{-x}(c_3\cos x+c_4\sin x)\\
y'&=2c_1e^{2x}-2c_2e^{-2x}-e^{-x}((c_3-c_4)\cos x+(c_3+c_4)\sin x)\\
y''&=4c_1e^{2x}+4c_2e^{-2x}+e^{-x}(2c_3\sin x-2c_4\cos x)\\
y'''&=8c_1e^{2x}-8c_2e^{-2x}+e^{-x}((2c_3+2c_4)\cos x+2(c_4-c_3)\sin
x);
\end{align*}
\begin{align*}
c_1+c_2+c_3&=5\\
2c_1-2c_2-c_3+c_4&=-2\\
4c_1+4c_2-2c_4&=6\\
8c_1-8c_2+2c_3+2c_4&=8;
\end{align*}
$c_1=1$, $c_2=1$, $c_3=3$, $c_4=1$;
$y=e^{2x}+e^{-2x}+e^{-x}(3\cos x+\sin x)$.


\item
\begin{alist}
\item
$
W(x)=\begin{vmatrix}
e^x&xe^x&e^{2x}\\
e^x&e^x(x+1)&2e^{2x}\\
e^x&e^x(x+2)&4e^{2x}
\end{vmatrix}$;
$
W(0)=\begin{vmatrix}
1&0&1\\1&1&2\\1&2&4
\end{vmatrix}=
\begin{vmatrix}
1&0&0\\1&1&1\\1&2&3
\end{vmatrix}
=\begin{vmatrix}
1&1\\2&3\end{vmatrix}=1.
$

\item
$
W(x)=\begin{vmatrix}
\cos 2x&\sin 2x&e^{3x}\\
-2\sin 2x&2\cos 2x&3e^{3x}\\
-4\cos 2x&-4\sin 2x&9e^{3x}
\end{vmatrix};
$
$
W(0)=\begin{vmatrix}
1&0&1\\0&2&3\\-4&0&9
\end{vmatrix}=
\begin{vmatrix}
1&0&0\\0&2&3\\-4&0&13
\end{vmatrix}=
\begin{vmatrix}
2&3\\0&13\end{vmatrix}=26.
$

\item
$
W(x)=\begin{vmatrix}
e^{-x}\cos x&e^{-x}\sin x&e^x\\
-e^{-x}(\cos x+\sin x)&e^{-x}(\cos x-\sin x)&e^x\\
2e^{-x}\sin x&-2e^{-x}\cos x&
e^x\end{vmatrix};
$
$
W(0)=\begin{vmatrix}
1&0&1\\-1&1&1\\0&-2&1
\end{vmatrix}=
\begin{vmatrix}
1&0&1\\0&1&2\\0&-2&1
\end{vmatrix}=
\begin{vmatrix}
1&2\\-2&1\end{vmatrix}=5.
$

\item
$
W(x)=\begin{vmatrix}
1&x&x^2&e^x\\
0&1&2x&e^x\\0&0&2&e^x\\0&0&0&e^x\end{vmatrix};
$
$
W(0)=\begin{vmatrix}
1&0&0&1\\0&1&0&1\\0&0&2&1\\0&0&0&1
\end{vmatrix}=1.
$

\item
$
W(x)=\begin{vmatrix}e^x&e^{-x}&\cos x&\sin x\\
e^x&-e^{-x}&-\sin x&\cos x\\e^x&e^{-x}&-\cos x&-\sin
x\\e^x&-e^{-x}&\sin x&-\cos x\end{vmatrix};
$
\begin{align*}
W(0)&=\begin{vmatrix}
1&1&1&0\\1&-1&0&1\\1&1&-1&0\\1&-1&0&-1
\end{vmatrix}
=\begin{vmatrix}
1&1&1&0\\1&-1&0&1\\1&1&-1&0\\0&0&0&-2
\end{vmatrix}
=-2\begin{vmatrix}
1&1&1\\1&-1&0\\1&1&-1\\
\end{vmatrix}\\
&=-2\begin{vmatrix}
1&1&1\\1&-1&0\\0&0&-2\\
\end{vmatrix}
=4\begin{vmatrix}
1&1\\1&-1\\\end{vmatrix}=-8.
\end{align*}

\item
$
W(x)=\begin{vmatrix}\cos x&\sin x&e^x\cos x&e^x\sin x\\
-\sin x&\cos x&e^x(\cos x-\sin x)&e^x(\cos x+\sin x)\\
-\cos x&-\sin x&-2e^x\sin x&2e^x\cos x\\
\sin x&-\cos x&-e^x(2\cos x+2\sin x)&e^x(2\cos x-2\sin x)
\end{vmatrix};
$

\begin{align*}
W(0)&=\begin{vmatrix}
1&0&1&0\\0&1&1&1\\-1&0&0&2\\0&-1&-2&2
\end{vmatrix}=
\begin{vmatrix}
1&0&1&0\\0&1&1&1\\0&0&1&2\\0&-1&-2&2
\end{vmatrix}=
\begin{vmatrix}
1&1&1\\0&1&2\\-1&-2&2
\end{vmatrix}\\
&=\begin{vmatrix}
1&1&1\\0&1&2\\0&-1&3
\end{vmatrix}=
\begin{vmatrix}
1&1&1\\0&1&2\\0&0&5
\end{vmatrix}=5.
\end{align*}
\end{alist}

\item
\begin{alist}
\item
Since $y=Q_1(D)P_1(D)y+Q_2(D)P_2(D)y$ and $P_1(D)y=P_2(D)y=0$, it
follows that $y=0$.

\item Suppose that
(A) $a_1u_1+\cdots+a_ru_r+b_1v_1+\cdots+b_sv_s=0$,
where $a_1,\dots,a_r$ and $b_1,\dots,b_s$ are constants. Denote
$u=a_1u_1+\cdots+a_ru_r$ and $v=b_1v_1+\cdots+b_sv_s$. Then
(B) $P_1(D)u=0$ and (C) $P_2(D)v=0$. Since $u+v=0$, $P_2(D)(u+v)=0$.
Therefore,$0=P_2(D)(u+v)=P_2(D)u+P_2(D)v$. Now (C) implies that
$P_2(D)u=0$. This, (B), and \part{a} imply that
$u=a_1u_1+\cdots+a_ru_r=0$, so $a_1=\dots=a_r=0$, since
$u_1,\dots,u_r$ are linearly independent. Now (A) reduces to
$b_1v_1+\cdots+b_sv_s=0$,
 so $b_1=c\dots=b_s=0$, since
$v_1,\dots,v_s$ are linearly independent.
Therefore,$u_1,\dots,u_r,
v_1,\dots,v_r$ are linearly independent.

\item It suffices to show that
 $\{y_1,y_2,\dots,y_n\}$ is linearly independent. Suppose that
$c_1y_1+\cdots+c_ny_n=0$. We may assume that $y_1,\dots,y_r$
are linearly independent solutions of $p_1(D)y=0$ and
$y_{r+1},\dots,y_n$ are solutions of
$P_2(D)=p_2(D)\cdots p_k(D)y=0$. Since $p_1(r)$ and $P_2(r)$ have no
common factors,
\part{b} implies
that (A) $c_1y_1+\cdots+c_ry_r=0$ and
(B) $c_{r+1}y_{r+1}+\cdots+c_ny_n=0$.
Now (A) implies that $c_1=\cdots=c_r=0$, since $y_1,\dots,y_r$
are linearly independent. If $k=2$, then $y_{r+1},\dots,y_n$
are linearly independent, so $c_{r+1}=\cdots=c_n=0$, and the proof is
complete. If $k>2$ repeat this argument, starting from (B), with $p_1$
replaced by $p_2$, and $P_2$ replaced by $P_3=p_3\cdots p_n$.
\end{alist}

\item
\begin{alist}
\item
\begin{align*}
(\cos A+i\sin A)(\cos B+i\sin B)&=(\cos A\cos B-\sin A\sin B)\\
&+(\cos A\sin B+\sin A\cos B)\\
&=\cos(A+B)+i\sin(A+B).
\end{align*}

\item Obvious for $n=0$. If $n=-1$ write
\begin{align*}
\frac{1}{\cos\theta+i\sin\theta}
&=\frac{1}{\cos\theta+i\sin\theta}
\frac{\cos\theta-i\sin\theta}{\cos\theta-i\sin\theta}\\
&=\frac{\cos\theta-i\sin\theta}{\cos^2\theta+\sin^2\theta}=
\cos\theta-i\sin\theta=\cos(-\theta)+i\sin(-\theta).
\end{align*}

\stepcounter{enumXi}
\item If $n$ is a negative integer, then
(B)
$(\cos\theta+i\sin\theta)^n=\dst\frac{1}{(\cos\theta+i\sin\theta)^{|n|}}$.
From the hint,
(C)
$\dst\frac{1}{(\cos\theta+i\sin\theta)^{|n|}}=(\cos\theta-i\sin\theta)^{|n|}
=(\cos(-\theta)+i\sin(-\theta))^{|n|}$.
Replacing $\theta$ by $-\theta$ and $n$ by $|n|$ in (A) shows that
(D)
$(\cos(-\theta)+i\sin(-\theta))^{|n|}=\cos(-|n|\theta)+i\sin(-|n|\theta)$.
Since $|n|=-n$, (E)
$\cos(-|n|\theta)+i\sin(-|n|\theta)=\cos n\theta+i\sin n\theta$. Now
(B), (C), (D), and (E) imply (A).

\item From (A), $z_k^n=\cos2k\pi+i\sin2k\pi=1$ and
$\zeta_k^n=\cos(2k+1)\pi+i\sin(2k+1)\pi=\cos(2k+1)\pi=\cos\pi=-1$.

\item From \part{e}, $\rho^{1/n}z_0,\dots,\rho^{1/n}z_{n-1}$
are all zeros of $z^n-\rho$. Since they are distinct numbers,
$z^n-\rho$ has the stated factorization.

 From \part{e}, $\rho^{1/n}\zeta_0,\dots,\rho^{1/n}\zeta_{n-1}$
are all zeros of $z^n+\rho$. Since they are distinct numbers,
$z^n+\rho$ has the stated factorization.
\end{alist}

\addtocounter{enumi}{-1}
\item%9.2.43
\begin{alist}
\item
$p(r)=r^3-1=(r-z_0)(r-z_1)(r-z_2)$ where
$z_k=\dst\cos\frac{2k\pi}{3}+i\sin\frac{2k\pi}{3}$, $k=0,1,2$. Hence,
$z_0=1$, $z_1=\dst-\frac{1}{2}+i\frac{\sqrt3}{2}$, and
$z_2=\dst-\frac{1}{2}-i\frac{\sqrt3}{2}$. Therefore,
$p(r)=(r-1)\dst\left(\left(r+\frac{1}{2}\right)^2+\frac{3}{4}\right)$,
so
$\dst\left\{e^x,\,e^{-x/2}\cos\left(\frac{\sqrt3}{2}x\right),
\,e^{-x/2}\sin\left(\frac{\sqrt3}{2}x\right)\right\}$
is a fundamental set of solutions.

\item
$p(r)=r^3+1=(r-\zeta_0)(r-\zeta_1)(r-\zeta_2)$ where
$\zeta_k=\dst\cos\frac{(2k+1)\pi}{3}+i\sin\frac{(2k+1)\pi}{3}$, $k=0,1,2$.
Hence,
 $\zeta_0=\dst\frac{1}{2}+i\frac{\sqrt3}{2}$, $\zeta_1=-1$
$\zeta_2=\dst\frac{1}{2}-i\frac{\sqrt3}{2}$. Therefore,
$p(r)=(r+1)\dst\left(\left(r-\frac{1}{2}\right)^2+\frac{3}{4}\right)$,
so
 $\dst\left\{e^{-x},\,e^{x/2}\cos\left(\frac{\sqrt3}{2}x\right),
\,e^{x/2}\sin\left(\frac{\sqrt3}{2}x\right)\right\}$
is a fundamental set of solutions.

\item
$p(r)=r^4+64=(r-2\sqrt2\zeta_0)(r-2\sqrt2\zeta_1)
(r-2\sqrt2\zeta_2)(r-2\sqrt2\zeta_3)$,
where
$\zeta_k=\dst\cos\frac{(2k+1)\pi}{4}+i\sin\frac{(2k+1)\pi}{4}$,
$k=0,1,2,3$. Therefore,
$\zeta_0=\dst\frac{1+i}{\sqrt2}$,
$\zeta_1=\dst\frac{-1+i}{\sqrt2}$,
$\zeta_2=\dst\frac{-1-i}{\sqrt2}$,
and $\zeta_3=\dst\frac{1-i}{\sqrt2}$, so
 $p(r)=((r-2)^2+4)((r+2)^2+4)$ and
$\{e^{2x}\cos2x,\,e^{2x}\sin2x,\,e^{-2x}\cos2x,\,e^{-2x}\sin2x\}$
is a fundamental set of solutions.

\item
$p(r)=r^6-1=(r-z_0)(r-z_1)(r-z_2)(r-z_3)(r-z_4)(r-z_5)$ where
$z_k=\dst\cos\frac{2k\pi}{6}+i\sin\frac{2k\pi}{6}$,
$k=0,1,2,3,4,5$. Therefore,
$z_0=1$,
$z_1=\dst\frac{1}{2}+i\frac{\sqrt3}{2}$,
$z_2=\dst-\frac{1}{2}+i\frac{\sqrt3}{2}$,
$z_3=-1$,
$z_4=\dst-\frac{1}{2}-i\frac{\sqrt3}{2}$,
and $z_5=\dst\frac{1}{2}-i\frac{\sqrt3}{2}$, so
$p(r)=(r-1)(r+1)\dst\left(\left(r-\frac{1}{2}\right)^2+\frac{3}{4}\right)
\left(\left(r+\frac{1}{2}\right)^2+\frac{3}{4}\right)$ and

 $\dst\left\{e^x,\,e^{-x}
,\,e^{x/2}\cos\left(\frac{\sqrt3}{2}x\right),
\,e^{x/2}\sin\left(\frac{\sqrt3}{2}x\right),
\,e^{-x/2}\cos\left(\frac{\sqrt3}{2}x\right),
\,e^{-x/2}\sin\left(\frac{\sqrt3}{2}x\right)\right\}$
is a fundamental set of solutions.

\item
$p(r)=r^6+64=(r-2\zeta_0)(r-2\zeta_1)(r-2\zeta_2)(r-2\zeta_3)
(r-2\zeta_4)(r-2\zeta_5)$ where
$\zeta_k=\dst\cos\frac{(2k+1)\pi}{6}+i\sin\frac{(2k+1)\pi}{6}$,
$k=0,1,2,3,4,5$. Therefore,
$\zeta_0=\dst\frac{\sqrt3}{2}+\frac{i}{2}$,
$\zeta_1=i$,
$\zeta_2=\dst-\frac{\sqrt3}{2}+\frac{i}{2}$,
$\zeta_3=\dst-\frac{\sqrt3}{2}-\frac{i}{2}$,
$\zeta_4=-i$, and
$\zeta_5=\dst\frac{\sqrt3}{2}-\frac{i}{2}$,  so
$p(r)=(r^2+4)((r-\sqrt3)^2+1)((r+\sqrt3)^2+1)$ and
 $\{\cos2x,\,\sin2x,\,e^{-\sqrt3x}\cos x,\,e^{-\sqrt3x}\sin x
,\,e^{\sqrt3x}\cos x,\,e^{\sqrt3x}\sin x\}$
is a fundamental set of solutions.

\item
$p(r)=(r-1)^6-1=(r-1-z_0)(r-1-z_1)(r-1-z_2)(r-1-z_3)(r-1-z_4)(r-1-z_5)$
where
$z_k=\dst\cos\frac{2k\pi}{6}+i\sin\frac{2k\pi}{6}$,
$k=0,1,2,3,4,5$. Therefore,
$z_0=1$,
$z_1=\dst\frac{1}{2}+i\frac{\sqrt3}{2}$,
$z_2=\dst-\frac{1}{2}+i\frac{\sqrt3}{2}$,
$z_3=-1$,
$z_4=\dst-\frac{1}{2}-i\frac{\sqrt3}{2}$,
and $z_5=\dst\frac{1}{2}-i\frac{\sqrt3}{2}$, so
$p(r)=r(r-2)\dst\left(\left(r-\frac{3}{2}\right)^2+\frac{3}{4}\right)
\left(\left(r-\frac{1}{2}\right)^2+\frac{3}{4}\right)$ and
$\dst\left\{1,\,e^{2x}
,\,e^{3x/2}\cos\left(\frac{\sqrt3}{2}x\right),
\,e^{3x/2}\sin\left(\frac{\sqrt3}{2}x\right),
\,e^{x/2}\cos\left(\frac{\sqrt3}{2}x\right),
\,e^{x/2}\sin\left(\frac{\sqrt3}{2}x\right)\right\}$
is a fundamental set of solutions.

\item
$p(r)=r^5+r^4+r^3+r^2+r+1=\dst\frac{r^6-1}{ r-1}$. Therefore,
from the solution of \part{d}
$p(r)=(r+1)\dst\left(\left(r-\frac{1}{2}\right)^2+\frac{3}{4}\right)
\left(\left(r+\frac{1}{2}\right)^2+\frac{3}{4}\right)$ and

$\dst\left\{e^{-x}
,\,e^{x/2}\cos\left(\frac{\sqrt3}{2}x\right),
\,e^{x/2}\sin\left(\frac{\sqrt3}{2}x\right),
\,e^{-x/2}\cos\left(\frac{\sqrt3}{2}x\right),
\,e^{-x/2}\sin\left(\frac{\sqrt3}{2}x\right)\right\}$
is a fundamental set of solutions.
\end{alist}
\end{sectionSolutions}

\section{Undetermined Coefficients for Higher Order  Equations}

\begin{sectionSolutions}
\item If $y=u^{-3x}$, then $y'''-2y''-5y'+6y=e^{-3x}[
(u'''-11u''+34u'-24u)-2(u''-6u'+9u)-5(u'-3u)+6u]
=e^{-3x}(u'''-11u''+34u'-24u)$. Let $u_p=A+Bx+Cx^2$, where
$(-24A+34B-22C)+(-24B+68C)x-24Cx^2=32-23x+6x^2$. Then $C=-1/4$,
$B=1/4$, $A=-3/4$ and $y_p=-\dst\frac{e^{-3x}}{4}(3-x+x^2)$.



\item If $y=ue^{-2x}$, then $y'''+3y''-y'-3y=e^{-2x}[
(u'''-6u''+12u'-8u)+ 3(u''-4u'+4u) -(u'-2u) -3u]
=e^{-2x}(u'''-3u''-u'+3u)$. Let $u_p=A+Bx+Cx^2$, where $(3A-B-6C)+
(3B-2C)x +3Cx^2=2-17x+3x^2$. Then $C=1$, $B=-5$, $A=1$, and
$y_p=e^{-2x}(1-5x+x^2)$.




\item If $y=ue^x$, then $y'''+y''-2y=e^x[
(u'''+3u''+3u'+u) +(u''+2u'+u) -2u]= e^x(u'''+4u''+5u')$. Let
$u_p=x(A+Bx+Cx^2)$, where $(5A+8B+6C) +(10B+24C)x +15Cx^2
=14+34x+15x^2$. Then $C=1$, $B=1$, $A=0$, and $y_p=x^2e^x(1+x)$.


\item If $y=ue^x$, then $y'''-y''-y'+y=e^x[
(u'''+3u''+3u'+u) -(u''+2u'+u) -(u'+u) +u]= e^x(u'''+2u'')$. Let
$u_p=x^2(A+Bx)$ where $(4A+6B)+12Bx= 7+6x$. Then $B=1/2$, $A=1$, and
$y_p=\dst\frac{x^2e^x}{2}(2+x)$.


\item If $y=ue^{3x}$, then $y'''-5y''+3y'+9y=e^{3x}[
(u'''+9u''+27u'+27u) -5(u''+6u'+9u) +3(u'+3u) +9u]=
e^{3x}(u'''+4u'')$.
Let $u_p=x^2(A+Bx+Cx^2)$, where $(8A+6B) +(24B+24C)x+ 48Cx^2=
22-48x^2$. Then $C=-1$, $B=1$, $A=2$, and $y_p=x^2e^{3x}(2+x-x^2)$.


\item If $y=ue^{x/2}$, then
$8y'''-12y''+6y'-y=e^{x/2}[ 8(u'''+3u''/2+3u'/4+u/8) -12(u''+u'+u/4)
+6(u'+u/2) -u] =8e^{x/2}u'''$, so $u'''=\dst\frac{1+4x}{8}$. Integrating
three times and taking the constants of integration to be zero yields
$u_p=\dst\frac{x^3}{48}(1+x)$. Therefore,
$y_p=\dst\frac{x^3e^{x/2}}{48}(1+x)$.


\item If $y=ue^{2x}$, then
$y^{(4)}+3y'''+y''-3y'-2y=e^{2x}[ (u^{(4)}+8u'''+24u''+32u'+16u)
+3(u'''+6u''+12u'+8u) +(u''+4u'+4u) -3(u'+2u) -2u]
=e^{2x}(u^{(4)}+11u'''+43u''+69u'+36u)$. Let $u_p=A+Bx$ where
$(36A+69B)+36Bx=-33-36x$. Then $B=-1$, $A=1$, and $y_p=e^{2x}(1-x)$.



\item If $y=ue^x$, then $4y^{(4)}-11y''-9y'-2y=e^x[
4(u^{(4)}+4u'''+6u''+4u'+u) -11(u''+2u'+u) -9(u'+u) -2u]
=e^x(4u^{(4)}+16u'''+13u''-15u'-18u)$. Let $u_p=A+Bx$ where
$-(18A+15B)-18Bx=-1+6x$. Then $B=-1/3$, $A=1/3$, and
$y_p=\dst\frac{e^x}{3}(1-x)$.


\item If $y=ue^x$, then
$y^{(4)}-4y'''+6y''-4y'+2y=e^x[ (u^{(4)}+4u'''+6u''+4u'+u)
-4(u'''+3u''+3u'+u) +6(u''+2u'+u) -4(u'+u) +2u]= e^x(u^{(4)}+u)$. Let
$u_p=A+Bx+Cx^2+Dx^3+Ex^4$ where $ (A+24E)+Bx+Cx^2+Dx^3+Ex^4=24+x+x^4$.
Then $E=1$ $D=0$, $C=0$ $B=1$, $A=0$, and $y_p=xe^x(1+x^3)$.


\item If $y=ue^{2x}$, then
$y^{(4)}+y'''-2y''-6y'-4y=e^{2x}[ (u^{(4)}+8u'''+24u''+32u'+16u)
+(u'''+6u''+12u'+8u) -2(u''+4u'+4u) -6(u'+2u) -4u]=
e^{2x}(u^{(4)}+9u'''+28u''+30u')$. Let $u_p=x(A+Bx+Cx^2)$ where
$(30A+56B+54C)+ (60B+168C)x +90Cx^2=-(4+28x+15x^2)$. Then $C=-1/6$,
$B=0$, $A=1/6$, and $y_p=\dst\frac{xe^{2x}}{6}(1-x^2)$.


\item If $y=ue^x$, then $y^{(4)}-5y''+4y=e^x[
(u^{(4)}+4u'''+6u''+4u'+u) -5(u''+2u'+u) +4u]
=e^x(u^{(4)}+4u'''+u''-6u')$. Let $u_p=x(A+Bx+Cx^2)$ where
$(-6A+2B+24C)+(-12B+6C)x-18Cx^2=3+x-3x^2$, so $C=1/6$, $B=0$, $A=1/6$.
Then $y_p=\dst\frac{xe^x}{6}(1+x^2)$.


\item If $y=ue^{2x}$, then
$y^{(4)}-3y'''+4y'=e^{2x}[ (u^{(4)}+8u'''+24u''+32u'+16u)
-3(u'''+6u''+12u'+8u) +4(u'+2u)] =e^{2x}(u^{(4)}+5u'''+6u'')$. Let
$u_p=x^2(A+Bx+Cx^2)$ where $(12A+30B+24C)+
(36B+120C)x+72Cx^3=15+26x+12x^2$. Then $C=1/6$, $B=1/6$, $A=1/2$, and
$y_p=\dst\frac{x^2e^{2x}}{6}(3+x+x^2)$.



\item
If $y=ue^x$, then
$2y^{(4)}-5y'''+3y''+y'-y=e^x[
2(u^{(4)}+4u'''+6u''+4u'+u) -5(u'''+3u''+3u'+u) +3(u''+2u'+u) +(u'+u)
-u]=e^x(2u^{(4)}+3u''')$.
Let $u_p=x^3(A+Bx)$, where
$(18A+48B)+72Bx=11+12x$.
Then  $B=1/6$, $A=1/6$, and
  $y_p=\dst\frac{x^3e^x}{6}(1+x)$.



\item
If $y=ue^{2x}$, then
$y^{(4)}-7y'''+18y''-20y'+8y=e^{2x}[
(u^{(4)}+8u'''+24u''+32u'+16u) -7(u'''+6u''+12u'+8u) +18(u''+4u'+4u)
-20(u'+2u) +8u]
=e^{2x}(u^{(4)}+u''')$.
Let $u_p=x^3(A+Bx+Cx^2)$ where
$(6A+24B) +(24B+120C)x+ 60Cx^2=3-8x-5x^2$. Then
  so $C=-1/12$  $B=1/12$, $A=1/6$, and
  $y_p=\dst\frac{x^3e^{2x}}{12}(2+x-x^2)$.


\item
If $y=ue^{-x}$, then
$y'''+y''-4y'-4y=e^{-x}[
(u'''-3u''+3u'-u) +(u''-2u'+u) -4(u'-u) -4u]
=e^{-x}(u'''-2u''-3u')$.
Let $u_p=(A_0+A_1x)\cos2x+(B_0+B_1x)\sin2x$, where
\begin{align*}
8A_1-14B_1&=-22\\
14A_1+8B_1&=-6\\
8A_0-14B_0-15A_1-8B_1&=1\\
14A_0+8B_0+8A_1-15B_1&=-1.
\end{align*}
Then
$A_1=-1$, $B_1=1$, $A_0=1$, $B_0=1$, and
$y_p=e^{-x}\left[(1-x)\cos2x+(1+x)\sin2x\right]$.



\item
If $y=ue^x$, then
$y'''-2y''+y'-2y=e^x[
(u'''+3u''+3u'+u) -2(u''+2u'+u) +(u'+u) -2u]
=e^x(u'''+u''-2u)$.
Let $u_p=(A_0+A_1x+A_2x^2)\cos2x+(B_0+B_1x+B_2x^2)\sin2x$ where
\begin{align*}
-6A_2-8B_2&=-4\\
8A_2-6B_2&=-3\\
-6A_1-8B_1-24A_2+8B_2&=5\\
8A_1-6B_1-8A_2-24B_2&=-5\\
-6A_0-8B_0-12A_1+4B_1+2A_2+12B_2&=-9\\
8A_0-6B_0-4A_1-12B_1-12A_2+2B_2&=6.
\end{align*}
Then $A_2=0$, $B_2=1/2$; $A_1=1/2$, $B_1=-1/2$; $A_0=1/1$, $B_0=1/2$;
and $y_p=\dst\frac{e^x}{2}[(1+x)\cos2x+(1-x+x^2)\sin2x]$.



\item
If $y=ue^x$, then
$y'''-y''+2y=e^x[
(u'''+3u''+3u'+u) -(u''+2u'+u) +2u]
=e^x(u'''+2u''+u'+2u)$.
Since $\cos x$ and $\sin x$ satisfy  $u'''+2u''+u'+2u=0$,
let  $u_p=x[(A_0+A_1x)\cos x+(B_0+B_1x)\sin x]$ where
\begin{align*}
-4A_1+8B_1&=4\\
-8A_1-4B_1&=-12\\
-2A_0+4B_0+4A_1+6B_1&=20\\
-4A_0-2B_0-6A_1+4B_1&=-12.
\end{align*}
Then $A_1=1$, $B_1=1$; $A_0=1$, $B_0=3$; and
 $y_p=xe^x[(1+x)\cos x+(3+x)\sin x]$.


\item
If $y=ue^{3x}$, then
$=e^{3x}[
(u'''+9u''+27u'+27u) -6(u''+6u'+9u) +18(u'+3u)]
=e^{3x}(u'''+3u''+9u'+27u)$.
Since $\cos3x$ and $\sin3x$ satisfy  $u'''+3u''+9u'+27u=0$,
let  $u_p=x[(A_0+A_1x)\cos3x+(B_0+B_1x)\sin3x]$ where
\begin{align*}
-36A_1+36B_1&=3\\
-36A_1-36B_1&=3\\
-18A_0+18B_0+6A_1+18B_1&=-2\\
-18A_0-18B_0-18A_1+6B_1&=3.
\end{align*}
Then  $A_1=-1/12$, $B_1=0$; $A_0=$, $B_0=-1/12$; and
 $y_p=-\dst\frac{xe^{3x}}{12}(x\cos3x+\sin3x)$.


\item
If $y=ue^x$, then
$y^{(4)}-3y'''+2y''+2y'-4y=e^x[
(u^{(4)}+4u'''+6u''+4u'+u) -3(u'''+3u''+3u'+u) +2(u''+2u'+u) +4(u'+u)
+u]=e^x(u^{(4)}+u'''-u''+u'-2u)$.
Let $u_p=A\cos2x+B\sin2x$ where $18A-6B=2$ and $6A+18B=-1$.
Then $A=1/12$, $B=-1/12$, and
$y_p=\dst\frac{e^x}{12}(\cos2x-\sin2x)$.


\item
If $y=ue^{-x}$, then
$y^{(4)}+6y'''+13y''+12y'+4y=e^{-x}[
(u^{(4)}-4u'''+6u''-4u'+u) +6(u'''-3u''+3u'-u) +13(u''-2u'+u) +12u'-u)
+4u]=e^{-x}(u^{(4)}+2u'''+u'')$.
Let $u_p=(A_0+A_1x)\cos x+(B_0+B_1x)\sin x$ where
\begin{align*}
-2B_1&=-1\\
2A_1&=-1\\
-2B_0-6A_1-2B_1&=4\\
2A_0+2A_1-6B_1&=-5.
\end{align*}
Then $A_1=-1/2$, $B_1=1/2$, $A_0=-1/2$, $B_0=-1$, and
 $y_p=-\dst\frac{e^{-x}}{2}[(1+x)\cos x+(2-x)\sin x]$.


\item
If $y=ue^x$, then
$y^{(4)}-5y'''+13y''-19y'+10y=e^x[
(u^{(4)}+4u'''+6u''+4u'+u) -5(u'''+3u''+3u'+u) +13(u''+2u'+u)
-19(u'+u) +10u]=e^x(u^{(4)}-u'''+4u''-4u')$. Since $\cos2x$ and $\sin2x$
satisfy $u^{(4)}-u'''+4u''-4u'=0$,
let $u_p=x(A\cos2x+B\sin2x)$ where $8A-16B=1$ and $16A+8B=1$. Then
$A=3/40$, $B_0=-1/40$, and
 $y_p=\dst\frac{xe^x}{40}(3\cos2x-\sin2x)$.


\item
If $y=ue^x$, then
$y^{(4)}-5y'''+13y''-19y'+10y=e^x[
(u^{(4)}+4u'''+6u''+4u'+u) -5(u'''+3u''+3u'+u) +13(u''+2u'+u)
-19(u'+u) +10u]=e^x(u^{(4)}-u'''+4u''-4u')$.
Since $\cos2x$ and $\sin2x$ satisfy $u^{(4)}-u'''+4u''-4u'=0$,
let $u_p=x[(A_0+A_1x)\cos2x+(B_0+B_1x)\sin2x)]$ where
\begin{align*}
16A_1-32B_1&=8\\
32A_1+16B_1&=-4\\
8A_0-16B_0-40A_1-12B_1&=7\\
16A_0+8B_0+12A_1-40B_1&=8.
\end{align*}
Then $A_1=0$, $B_1=-1/4$; $A_0=0$, $B_0=-1/4$, and
 $y_p=-\dst\frac{xe^x}{4}(1+x)\sin2x$.



\item
If $y=ue^{2x}$, then
$y^{(4)}-8y'''+32y''-64y'+64y+4y=e^{2x}[
(u^{(4)}+8u'''+24u''+32u'+16u) -8(u'''+6u''+12u'+8u) +32(u''+4u'+4u)
-64(u'+2u) +64u]=e^{2x}(u^{(4)}+8u''+16u)$. Since $\cos2x$,  $\sin2x$, $x\cos2x$,
and $x\sin2x$ satisfy $u^{(4)}+8u''+16u=0$,
let $u_p=x^2(A\cos2x+B\sin2x)$ where $-32A=1$ and $-32B=-1$. Then
$A=-1/32$, $B=1/32$, and
 $y_p=-\dst\frac{x^2e^{2x}}{32}(\cos2x-\sin2x)$.



\item
Find particular solutions of
\part{a} $y'''-4y''+5y'-2y=-4e^x$,
\part{b} $y'''-4y''+5y'-2y=e^{2x}$,
and  \part{c} $y'''-4y''+5y'-2y=-2\cos x+4\sin x$.

\begin{alist}
\item If $y=ue^x$, then
$y'''-4y''+5y'-2y=e^x[(u'''+3u''+3u'+u) -4(u''+2u'+u) +5(u'+u)
-2u]=e^x(u'''-u'')$. Let $u_{1p}=Ax^2$ where $-2A=-4$. Then $A=2$,
and $y_{1p}=2x^2e^x$.


\item If $y=ue^{2x}$, then
$y'''-4y''+5y'-2y=e^{2x}[(u'''+6u''+12u'+8u) -4(u''+4u'+4u) +5(u'+2u)
-2u ]=e^{2x}(u'''+2u''+u')$. Let $u_{2p}=x$. Then $y_{2p}=xe^{2x}$.

\item If $y_{3p}=A\cos x+B\sin x$, then
$y_{3p}'''-4y_{3p}''+5y_{3p}'-2y_{3p}=(2A+4B)\cos x+(-4A+2B)\sin x=
-2\cos x+4\sin x$
if  $A=-1$ and $B=0$, so
$y_{3p}=-\cos x$.
\end{alist}

From the principle of superposition,
 $y_p=2x^2e^x+xe^{2x}-\cos x$.



\item Find particular solutions of \part{a}
$y'''-y'=-2(1+x)$, \part{b} $y'''-y'=4e^x$, \part{c}
$y'''-y'=-6e^{-x}$, and \part{d} $y'''-y'=96e^{3x}$

\begin{alist}
\item Let $y_{1p}=x(A+Bx)$. Then $y_{1p}'''-y_{1p}'=-A-2Bx=-2(1+x)$
if $A=2$ and $B=1$; therefore $y_{1p}=2x+2x^2$.

\item If $y=ue^x$, then $y'''-y'=e^x [(u'''+3u''+3u'+u)
-(u'+u)]=e^x(u'''+3u''+2u')$. Let $u_{2p}=4x$. Then $y_{2p}=4xe^x$.

\item If $y=ue^{-x}$, then $y'''-y'=e^{-x} [(u'''-3u''+3u'-u)
-(u'-u)]=e^{-x}(u'''-3u''+2u')$. Let $u_{2p}=-3x$. Then
$y_{2p}=-6xe^{-x}$.

\item Since $e^{3x}$ does not satisfy the complementary equation,
let $y_{4p}=Ae^{3x}$. Then $y_{4p}'''-y_{4p}' =24Ae^{3x}$. Let $A=4$;
then $y_{4p}=4e^{4x}$.
\end{alist}

From the principle of superposition,
$y_p=2x+x^2+2xe^x-3xe^{-x}+4e^{3x}$

\item
 Find particular solutions of
\part{a} $y'''+3y''+3y'+y=12e^{-x}$ and \part{b}
 $y'''+3y''+3y'+y=9\cos2x-13\sin2x$.

\begin{alist}
\item If $y=ue^{-2x}$, then $y'''+3y''+3y'+y=e^{-2x}[
(u'''-3u''+3u'-u) +3(u''-2u'+u) +3(u'-u) +u]
=e^{-x}u'''$. Let $u_{1p}'''=12$. Integrating three times and
taking the constants of integration to be zero yields
$u_{1p}=2x^3$. Therefore,$y_{1p}=2x^3$.

\item Let $y_{2p}=A\cos2x+B\sin2x$ where
$-11A-2B=9$ and $2A-11B=-13$.
Then $A=-1$, $B=1$, and
 $y_{2p}=-\cos2x+\sin2x$.

From the principle of superposition,
 $y_p=2x^3e^{-2x}-\cos2x+\sin2x$.
\end{alist}

\item
 Find particular solutions of
\part{a} $y^{(4)}-5y''+4y=-12e^x$,
\part{b} $y^{(4)}-5y''+4y=6e^{-x}$, and
\part{c} $y^{(4)}-5y''+4y=10\cos x$.

\begin{alist}
\item If $y=ue^x$, then $y^{(4)}-5y''+4y=e^x[
(u^{(4)}+4u'''+6u''+4u'+u) -5(u''+2u'+u) +4u]
=e^x(u^{(4)}+4u'''+u''-6u')$. Let $u_{1p}=2x$. Then $y_{1p}=2xe^x$.


\item If $y=ue^{-x}$, then $y^{(4)}-5y''+4y=e^{-x}[
(u^{(4)}-4u'''+6u''-4u'+u) -5(u''-2u'+u) +4u]
=e^{-x}(u^{(4)}-4u'''+u''+6u')$. Let $u_{2p}=x$. Then
$y_{2p}=xe^{-x}$.

\item Let $y_{3p}=A\cos x+B\sin x$ where $10A=10$ and $10B=0$. Then
$A=1$, $B=0$, and $y_{3p}=\cos x$.
\end{alist}

From the principle of superposition, $y_p=2xe^x+xe^{-x}+\cos x$.


\item
 Find particular solutions of
\part{a} $y^{(4)}+2y'''-3y''-4y'+4y=2e^x(1+x)$  and
\part{b} $y^{(4)}+2y'''-3y''-4y'+4y=e^{-2x}$.

\begin{alist}
\item
If $y=ue^x$, then $y^{(4)}+2y'''-3y''-4y'+4y=e^x[
(u^{(4)}+4u'''+6u''+4u'+u) +2(u'''+3u''+3u'+u) -3(u''+2u'+u)
-4(u'+u)+4u]=e^x(u^{(4)}+6u'''+9u'')$. Let $u_{1p}=x^2(A+Bx)$
where $(18A+36B)+54Bx=2+2x$. Then $B=1/27$, $A=1/27$, and
$y_{1p}=\dst\frac{x^2}{27}(1+x)e^x$.

\item
If $y=ue^{-2x}$, then $y^{(4)}+2y'''-3y''-4y'+4y=e^{-2x}[
(u^{(4)}-8u'''+24u''-32u'+16u) +2(u'''-6u''+12u'-8u) -3(u''-4u'+4u)
-4(u'-2u) +4u]=e^{-2x}(u^{(4)}-6u'''+9u'')$. Let $u_{2p}=Ax^2$
where $18A=1$. Then $A=1/18$ and
$y_p=\dst\frac{x^2}{18}e^{-2x}$.
\end{alist}

From the principle of superposition,
$y_p=\dst\frac{x^2}{54}[(2+2x)e^x+3e^{-2x}]$.


\item
 Find particular solutions of
\part{a} $y^{(4)}+5y'''+9y''+7y'+2y=e^{-x}(30+24x)$ and
\part{b} $y^{(4)}+5y'''+9y''+7y'+2y=-e^{-2x}$.

\begin{alist}
\item
If $y=ue^{-x}$, then $y^{(4)}+5y'''+9y''+7y'+2y=e^{-x}[
(u^{(4)}-4u'''+6u''-4u'+u) +5(u'''-3u''+3u'-u) +9(u''-2u'+u) +7(u'-u)
+2u]=e^{-x}(u^{(4)}+u''')$. Let $u_{1p}=x^3(A+Bx)$ where
$(6A+24B)+24Bx=30+24x$. The $B=1$, $A=1$, and
 $y_{1p}=x^3(1+x)e^{-x}$.

\item
If $y=ue^{-2x}$, then $y^{(4)}+5y'''+9y''+7y'+2y=e^{-2x}[
(u^{(4)}-8u'''+24u''-32u'+16u) +5(u'''-6u''+12u'-8u) +9(u''-4u'+4u)
+7(u'-2u) +2u]=e^{-2x}(u^{(4)}-3u'''+3u''-u')$. Let $u_{2p}=x$. Then
 $y_{2p}=xe^{-2x}$.
\end{alist}

From the principle of superposition,
 $y_p=x^3(1+x)e^{-x}+xe^{-2x}$.



\item
 If $y=ue^{2x}$, then $y'''-y''-y'+y=e^{2x}[
(u'''+6u''+12u'+8u) -(u''+4u'+4u) -(u'+2u) +u]
=e^{2x}(u'''+5u''+7u'+3u)$. Let $u_p=A+Bx$, where
$(3A+7B)+3x=10+3x$. Then $B=1$,
 $A=1$ and $y_p=e^{2x}(1+x)$. Since $p(r)=(r+1)(r-1)^2$,
 $y=e^{2x}(1+x)+c_1e^{-x}+e^x(c_2+c_3x)$


\item
 If $y=ue^{2x}$, then $y'''-6y''+11y'-6y=e^{2x}[
(u'''+6u''+12u'+8u) -6(u''+4u'+4u) +11(u'+2u) -6u]
=e^{2x}(u'''-u')$. Let $u_p=x(A+Bx+Cx^2)$ where
$(-A+6C)-2Bx-3Cx^2=5-4x-3x^2$. Then $C=1$, $B=2$,
 $A=1$, and $y_p=xe^{2x}(1+x)^2$.
 Since $p(r)=(r-1)(r-2)(r-3)$,
 $y=xe^{2x}(1+x)^2+c_1e^x+c_2e^{2x}+c_3e^{3x}$.


\item
 If $y=ue^x$, then $y'''-3y''+3y'-y=e^x[
(u'''+3u''+3u'+u) -3(u''+2u'+u) +3(u'+u) -u]
=e^xu'''$. Let $u'''=1+x$. Integrating three times and taking
the constants of integration to be zero yields
$u=\dst\frac{x^3}{24}(4+x)$. Therefore,
 $y_p=\dst\frac{x^3e^x}{24}(4+x)$.
 Since $p(r)=(r-1)^3$,
 $y=\dst\frac{x^3e^x}{24}(4+x)
+e^x(c_1+c_2x+c_3x^2)$.



\item
If $y=ue^{-2x}$, then
$y'''+2y''-y'-2y=e^{-2x}[
(u'''-6u''+12u'-8u) +2(u''-4u'+4u) -(u'-2u) -2u]=
e^{-2x}(u'''-4u''+3u')$.
Let $u_p=(A_0+A_1x)\cos x+(B_0+B_1x)\sin x$ where
\begin{align*}
4A_1+2B_1&=-2\phantom{.}\\
-2A_1+4B_1&=-9\phantom{.}\\
4A_0+2B_0-8B_1 &=23\phantom{.}\\
-2A_0+4B_0+8A_1&=8.
\end{align*}
Then $A_1=1/2$, $B_1=-2$; $A_0=1$, $B_0=3/2$, and
 $y_p=e^{-2x}\dst\left[\left(1+\frac{x}{2}\right)\cos
x+\left(\frac{3}{2}-2x\right)\sin x\right]$.
Since $p(r)=(r-1)(r+1)(r+2)$,
 $y=e^{-2x}\dst{\left[\left(1+\frac{x}{2}\right)
\cos x+\left(\frac{3}{2}-2x\right)\sin x\right]}
+c_1e^x+c_2e^{-x}+c_3e^{-2x}$.


\item
If $y=ue^x$, then
$y^{(4)}-4y'''+14y''-20y'+25y=e^x[
(u^{(4)}+4u'''+6u''+4u'+u) -4(u'''+3u''+3u'+u) +14(u''+2u'+u)
-20(u'+u) +25u
]=e^x(u^{(4)}+8u''+16u)$.
Since $\cos2x$, $\sin2x$, $x\cos2x$, and $x\sin2x$
satisfy $u^{(4)}+8u''+16u=0$,
let $u_p=x^2[(A_0+A_1x)\cos2x+(B_0+B_1x)\sin2x]$ where
\begin{align*}
-96A_1&=6\phantom{.}\\
-96B_1&=0\phantom{.}\\
-32A_0+48B_1&=2\phantom{.}\\
-48A_1-32B_0&=3.
\end{align*}
Then $A_1=-1/16$, $B_1=0$; $A_0=-1/16$, $B_0=0$, and
 $y_[=-\dst\frac{x^2e^x}{16}(1+x)\cos2x$. Since
$p(r)=[(r-1)^2+1]^2$,
 $y=-\dst\frac{x^2e^x}{16}(1+x)\cos2x
+e^x\left[(c_1+c_2x)\cos2x+(c_3+c_4x)\sin2x\right]$.



\item
 If $y=ue^{-x}$, then $y'''-y''-y'+y=e^{-x}[
(u'''-3u''+3u'-u) -(u''-2u'+u) -(u'-u) +u]
=e^{-x}(u'''-4u''+4u')$. Let $u_p=x(A+Bx)$, where
$(4A-8B)+8Bx=-4+8x$. Then $B=1$,
 $A=1$, and $y_p=x(1+x)e^{-x}$. Since $p(r)=(r+1)(r-1)^2$
the general solution is
 $y=x(1+x)e^{-x}+c_1e^{-x}+c_2e^x+c_3xe^x$.
 Therefore,
\[
\begin{bmatrix}y\\y'\\y''\end{bmatrix}=
\begin{bmatrix}
x(1+x)e^{-x}\\
-e^{-x}(x^2-x-1)\\
e^{-x}(x^2-3x)\\
\end{bmatrix}+
\begin{bmatrix}
e^{-x}&e^x&xe^x\\
-e^{-x}&e^x&e^x(x+1)\\
e^{-x}&e^x&e^x(x+2)
\end{bmatrix}
\threecol{c_1}{c_2}{c_3}.
\]
Setting $x=0$ and imposing the initial conditions yields
\[
\threecol200=\threecol010+
\begin{bmatrix}1&1&0\\-1&1&1\\1&1&2\end{bmatrix}
\threecol{c_1}{c_2}{c_3},
\]
so $c_1=1$, $c_2=1$, $c_3=-1$, and
 $y=e^{-x}(1+x+x^2) +(1-x)e^x$.


\item
 If $y=ue^{-x}$, then $y'''-2y''-5y'+6y=e^{-x}[
(u^{(4)}-4u'''+6u''-4u'+u) +2(u'''-3u''+3u'-u) +2(u''-2u'+u) +2(u'-u)
+u=e^{-x}(u^{(4)}-2u'''+2u'')$. Let $u_p=x^2(A+Bx)$, where
$(4A-12B)+12Bx=20-12x$. Then $B=-1$,
 $A=2$, and $y_p=x^2(2-x)e^{-x}$. Since $p(r)=(r+1)^2(r^2+1)$,
the general solution is
$y=x^2(2-x)e^{-x}+e^{-x}(c_1+c_2x)+c_3\cos x+c_4\sin x$.
 Therefore,
\[
\begin{bmatrix}y\\y'\\y''\\y'''\end{bmatrix}=
\begin{bmatrix}
x^2(2-x)e^{-x}\\
x(x^2-5x+4)e^{-x}\\
-(x^3-8x^2+14x-4)e^{-x}\\
(x^3-11x^2+30x-18)e^{-x}
\end{bmatrix}+
\begin{bmatrix}
e^{-x}&xe^{-x}&\cos x&\sin x\\
-e^{-x}&(1-x)e^{-x}&-\sin x&\cos x\\
e^{-x}&(x-2)e^{-x}&-\cos x&-\sin x\\
-e^{-x}&(3-x)e^{-x}&\sin x&-\cos x
\end{bmatrix}
\begin{bmatrix}c_1\\c_2\\c_3\\c_4\end{bmatrix}
\]
Setting $x=0$ and imposing the initial conditions yields
\[
\begin{bmatrix}3\\-4\\7\\-22\end{bmatrix}=
\begin{bmatrix}0\\0\\4\\-18\end{bmatrix}+
\begin{bmatrix}1&0&1&0\\-1&1&0&1\\1&-2&-1&0\\-1&3&0&-1\end{bmatrix}
\begin{bmatrix}c_1\\c_2\\c_3\\c_4\end{bmatrix},
\]
so $c_1=2$, $c_2=-1$, $c_3=1$, $c_4=-1$, and
$y=(2-x)(x^2+1)e^{-x}+\cos x-\sin x$.


\item
 If $y=ue^x$, then $y^{(4)}-3y'''+5y''-2y'=e^x[
(u^{(4)}+4u'''+6u''+4u'+u) -3(u'''+3u''+3u'+u) +4(u''+2u'+u) -2(u'+u)
]=e^x(u^{(4)}+u'''+u''+u')$. Since $\cos x$ and $\sin x$
satisfy $u^{(4)}+u'''+u''+u'=0$, let $u_p=x(A\cos x+B\sin x)$
where $-2A-2B=-2$ and $2A-2B=2$. Then $A=1$, $B=0$, and $y_p=xe^x\cos
x$. Since $p(r)=r(r-1)[(r-1)^2+1]$ the general solution is
$y=xe^x\cos x+c_1+e^x(c_2+c_3\cos x+c_4\sin x)$.
 Therefore,
\[
\hspace{-25pt}
\begin{bmatrix}y\\y'\\y''\\y'''\end{bmatrix}\!=\!
\begin{bmatrix}
xe^x\cos x\\
e^x((x+1)\cos x-x\sin x)\\
e^x(2\cos x-2(x+1)\sin x)\\
\!-e^x(2x\cos x+2(x+3)\sin x)\!
\end{bmatrix}+
\begin{bmatrix}
1&e^x&e^x\cos x&e^x\sin x\\
0&e^x&e^x(\cos x-\sin x)&e^x(\cos x+\sin x)\\
0&e^x&-2e^x\sin x&2e^x\cos x\\
0&e^x&\!\!-e^x(2\cos x+2\sin x)\!\!&e^x(2\cos x-2\sin x)
\end{bmatrix}\!
\begin{bmatrix}
c_1\\c_2\\c_3\\c_4
\end{bmatrix}
\]
Setting $x=0$ and imposing the initial conditions yields
\[
\begin{bmatrix}
2\\0\\-1\\-5
\end{bmatrix}=
\begin{bmatrix}
0\\1\\2\\0
\end{bmatrix}+
\begin{bmatrix}1&1&1&0\\0&1&1&1\\
0&1&0&2\\0&1&-2&2\end{bmatrix}
\begin{bmatrix}
c_1\\c_2\\c_3\\c_4
\end{bmatrix},
\]
so $c_1=2$, $c_2=-1$, $c_3=1$, $c_4=-1$, and
$y=2+e^x\left[(1+x)\cos x-\sin x-1\right]$.

\end{sectionSolutions}



\section{Variation of Parameters for Higher Order  Equations}

\begin{sectionSolutions}
\item
$W=\begin{vmatrix}
e^{-x^2}&xe^{-x^2}&x^2e^{-x^2}\\
-2xe^{-x^2}&e^{-x^2}(1-2x^2)&2x
e^{-x^2}(1-x^2)\\
e^{-x^2}(4x^2-2)&2xe^{-x^2}(2x^2-3)&2e^{-x^2}(2x^4-5x^2+1)
\end{vmatrix}=2e^{-3x^2}$;

$W_1=\begin{vmatrix}
xe^{-x^2}&x^2e^{-x^2}\\
e^{-x^2}(1-2x^2)&2xe^{-x^2}(1-x^2)
\end{vmatrix}=x^2e^{-2x^2}$;
$W_2=\begin{vmatrix}
e^{-x^2}&x^2e^{-x^2}\\
-2xe^{-x^2}&2xe^{-x^2}(1-x^2)
\end{vmatrix}=2xe^{-2x^2}$;
$W_3=\begin{vmatrix}
e^{-x^2}&xe^{-x^2}\\
-2xe^{-x^2}&e^{-x^2}(1-2x^2)
\end{vmatrix}=e^{-2x^2}$;
$u_1'=\dst\frac{FW_1}{ P_0W}=\dst \frac{1}{2}x^{5/2}$;
$u_2'=-\dst\frac{FW_2}{ P_0W}=\dst -x^{3/2}$;
$u_3'=\dst\frac{FW_2}{ P_0W}=\dst \sqrt{x}/2$;
$u_1=\dst x^{7/2}/7$;
$u_2=\dst -\frac{2}{5}x^{5/2}$;
$u_3=\dst x^{3/2}/3$;
$y_p=u_1y_1+u_2y_2+u_3y_3=\dst \frac{8}{105}e^{-x^2}x^{7/2}$.


\item
$W=\dst\begin{vmatrix}
1&\dst\frac{e^x}{ x}&\dst\frac{e^{-x}}{ x}\\
0&\dst\frac{e^x(x-1)}{ x^2}&-\dst\frac{e^{-x}(x+1)}{ x^2}\\
0&\dst\frac{e^x(x^2-2x+2)}{ x^3}&\dst\frac{e^{-x}(x^2+2x+2)}{ x^3}
\end{vmatrix}=2/x^2$;
$W_1=\begin{vmatrix}
\dst\frac{e^x}{ x}&\dst\frac{e^{-x}}{ x}
\\\dst\frac{e^x(x-1)}{ x^2}&-\dst\frac{e^{-x}(x+1)}{ x^2}
\end{vmatrix}=-\dst\frac{2}{ x^2}$;
$W_2=
\begin{vmatrix}1&\dst\frac{e^{-x}}{ x}\\
0&-\dst\frac{e^{-x}(x+1)}{ x^2}
\end{vmatrix}=-\dst\frac{e^{-x}(x+1)}{ x^2}$;
$W_3=\begin{vmatrix}
1&\dst\frac{e^x}{ x}\\
0&\dst\frac{e^x(x-1)}{ x^2}
\end{vmatrix}=\dst\frac{e^x(x-1)}{ x^2}$;
$u_1'=\dst\frac{FW_1}{ P_0W}=\dst -2$;
$u_2'=-\dst\frac{FW_2}{ P_0W}=\dst e^{-x}(x+1)$;
$u_3'=\dst\frac{FW_2}{ P_0W}=\dst e^x(x-1)$;
$u_1=\dst -2x$;
$u_2=\dst -e^{-x}(x+2)$;
$u_3=\dst e^x(x-2)$;
$y_p=u_1y_1+u_2y_2+u_3y_3=\dst -2(x^2+2)/x$.



\item
$W=\begin{vmatrix}
e^x&e^{-x}&\dst\frac{1}{ x}\\
e^x&-e^{-x}&-\dst{1/x^2}\\e^x&e^{-x}&\dst\frac{2}{ x^3}
\end{vmatrix}=\dst\frac{2(x^2-2)}{ x^3}$;
$W_1=\begin{vmatrix}
e^{-x}&\dst\frac{1}{ x}\\-e^{-x}&-\dst\frac{1}{ x^2}
\end{vmatrix}=\dst\frac{e^{-x}(x-1)}{ x^2}$;
$W_2=\begin{vmatrix}
e^x&\dst\frac{1}{ x}\\e^x&-\dst\frac{1}{ x^2}
\end{vmatrix}=-\dst\frac{e^x(x+1)}{ x^2}$;
$W_3=\begin{vmatrix}
e^x&e^{-x}\\e^x&-e^{-x}
\end{vmatrix}=-2$;
$u_1'=\dst\frac{FW_1}{ P_0W}=\dst e^{-x}(x-1)$;
$u_2'=-\dst\frac{FW_2}{ P_0W}=\dst e^x(x+1)$;
$u_3'=\dst\frac{FW_2}{ P_0W}=\dst -2x^2$;
$y_p=u_1y_1+u_2y_2+u_3y_3=\dst-2\frac{x^2}{3}$.



\item
$W=\begin{vmatrix}
\sqrt{x}&\dst\frac{1}{ \sqrt{x}}&x^2\\
\dst\frac{1}{ 2\sqrt{x}}&-\dst\frac{1}{ 2x^{3/2}}&2x\\
-\dst\frac{1}{ 4x^{3/2}}&\dst\frac{3}{ 4x^{5/2}}&2
\end{vmatrix}=-\dst\frac{15}{ 4x}$;
$W_1=\begin{vmatrix}
\dst\frac{1}{ \sqrt{x}}&x^2\\
-\dst\frac{1}{ 2x^{3/2}}&2x
\end{vmatrix}=\dst\frac{5\sqrt{x}}{ 2}$;
$W_2=\begin{vmatrix}
\sqrt{x}&x^2\\
\dst\frac{1}{ 2\sqrt{x}}&2x
\end{vmatrix}=\dst\frac{3x^{3/2}}{ 2}$;
$W_3=\begin{vmatrix}
\sqrt{x}&\dst\frac{1}{ \sqrt{x}}\\
\dst\frac{1}{ 2\sqrt{x}}&-\dst\frac{1}{ 2x^{3/2}}
\end{vmatrix}=-\dst\frac{1}{ x}$;
$u_1'=\dst\frac{FW_1}{ P_0W}=\dst-5\sqrt{x}$;
$u_2'=-\dst\frac{FW_2}{ P_0W}=\dst3x^{3/2}$;
$u_3'=\dst\frac{FW_2}{ P_0W}=\dst\dst\frac{2}{ x}$;
$u_1=\dst-\frac{10}{3}x^{3/2}$;
$u_2=\dst\frac{6}{5}x^{5/2}$;
$u_3=\dst2\ln|x|$;
$y_p=u_1y_1+u_2y_2+u_3y_3=\dst2x^2\ln|x|-\frac{32}{15}x^2$.
Since $-\dst\frac{32}{15}x^2$ satisfies the complementary equation
we take $y_p=\ln|x|$.




\item
$W=\begin{vmatrix}
x&1/x&\dst\frac{e^x}{ x}\\1&-\dst\frac{1}{ x^2}&\dst\frac{e^x(x-1)}{ x^2}\\
0&\dst\frac{2}{ x^3}&\dst\frac{e^x(x^2-2x+2)}{ x^3}
\end{vmatrix}=\dst\frac{2e^x(1-x)}{ x^3}$;
$W_1=\begin{vmatrix}
\dst\frac{1}{ x}&\dst\frac{e^x}{ x}\\-\dst\frac{1}{
x^2}&\dst\frac{e^x(x-1)}{ x^2}
\end{vmatrix}=\dst\frac{e^x}{ x^2}$;
$W_2=\begin{vmatrix}
x&\dst\frac{e^x}{ x}\\
1&\dst\frac{e^x(x-1)}{ x^2}
\end{vmatrix}=\dst\frac{e^x(x-2)}{ x}$;
$W_3=\begin{vmatrix}
x&\dst\frac{1}{ x}\\
1&-\dst\frac{1}{ x^2}
\end{vmatrix}=-\dst\frac{2}{ x}$;
$u_1'=\dst\frac{FW_1}{ P_0W}=\dst1$;
$u_2'=-\dst\frac{FW_2}{ P_0W}=\dst x(2-x)$;
$u_3'=\dst\frac{FW_2}{ P_0W}=\dst-2xe^{-x}$;
$u_1=\dst x$;
$u_2=\dst \frac{x^2(3-x)}{ 3}$;
$u_3=\dst2e^{-x}(x+1)$;
$y_p=u_1y_1+u_2y_2+u_3y_3=\dst\frac{2x^3+3x^2+6x+6}{ 3x}$.
Since $x+\dst\frac{2}{ x}$ satisfies the complementary equation
we take $y_p=\dst\frac{2x^2+6}{3}$.


\item
$W=\begin{vmatrix}x&e^x&e^{-x}\\
1&e^x&-e^{-x}\\0&e^x&e^{-x}
\end{vmatrix}=2x$;
$W_1=\begin{vmatrix}
e^x&e^{-x}\\e^x&-e^{-x}
\end{vmatrix}=-2$;
$W_2=\begin{vmatrix}
x&e^{-x}\\1&-e^{-x}
\end{vmatrix}=-e^{-x}(x+1)$;
$W_3=\begin{vmatrix}x&e^x\\1&e^x
\end{vmatrix}=e^x(x-1)$;
$u_1'=\dst\frac{FW_1}{ P_0W}=\dst-1$;
$u_2'=-\dst\frac{FW_2}{ P_0W}=\dst e^{-x}(x+1)/2$;
$u_3'=\dst\frac{FW_2}{ P_0W}=\dst e^x(x-1)/2$;
$u_1=\dst-x$;
$u_2=\dst-e^{-x}(x+2)/2$;
$u_3=\dst e^x(x-2)/2$;
$y_p=u_1y_1+u_2y_2+u_3y_3=\dst-x^2-2$.



\item
$W=\begin{vmatrix}
\sqrt{x}&1/\sqrt{x}&x^{3/2}&\dst\frac{1}{ x^{3/2}}
\\\dst\frac{1}{
2\sqrt{x}}&-\dst\frac{1}{ 2x^{3/2}}&\dst\frac{3\sqrt{x}}{
2}&-\dst\frac{3}{ 2x^{5/2}}\\
-\dst\frac{1}{
4x^{3/2}}&\dst\frac{3}{ 4x^{5/2}}&\dst\frac{3}{
4\sqrt{x}}&\dst\frac{15}{ 4x^{7/2}}\\
\dst\frac{3}{
8x^{5/2}}&-\dst\frac{15}{
8x^{7/2}}&-\dst\frac{3}{ 8x^{3/2}}&-\dst\frac{105}{ 8x^{9/2}}
\end{vmatrix}=\dst\frac{12}{ x^6}$;
$W_1=
\begin{vmatrix}
\dst\frac{1}{ \sqrt{x}}&x^{3/2}&\dst\frac{1}{ x^{3/2}}\\
-\dst\frac{1}{ 2x^{3/2}}&\dst\frac{3\sqrt{x}}{ 2}&-\dst\frac{3}{
2x^{5/2}}\\
\dst\frac{3}{ 4x^{5/2}}&\dst\frac{3}{ 4\sqrt{x}}&\dst\frac{15}{ 4x^{7/2}}
\end{vmatrix}=\dst\frac{6}{ x^{7/2}}$;
$W_2=\begin{vmatrix}
\sqrt{x}&x^{3/2}&\dst\frac{1}{ x^{3/2}}\\
\dst\frac{1}{ 2\sqrt{x}}&\dst\frac{3\sqrt{x}}{ 2}&-\dst\frac{3}{
2x^{5/2}}\\
-\dst\frac{1}{ 4x^{3/2}}&\dst\frac{3}{ 4\sqrt{x}}&\dst\frac{15}{ 4x^{7/2}}
\end{vmatrix}=\dst\frac{6}{ x^{5/2}}$;
$W_3=\begin{vmatrix}
\sqrt{x}&\dst\frac{1}{ \sqrt{x}}&\dst\frac{1}{ x^{3/2}}\\
\dst\frac{1}{ 2\sqrt{x}}&-\dst\frac{1}{ 2x^{3/2}}&-\dst\frac{3}{
2x^{5/2}}\\
-\dst\frac{1}{ 4x^{3/2}}&\dst\frac{3}{ 4x^{5/2}}&\dst\frac{15}{ 4x^{7/2}}
\end{vmatrix}=-\dst\frac{2}{ x^{9/2}}$;
$W_4=\begin{vmatrix}
\sqrt{x}&\dst\frac{1}{ \sqrt{x}}&x^{3/2}\\
\dst\frac{1}{ 2\sqrt{x}}&-\dst\frac{1}{ 2x^{3/2}}&\dst\frac{3\sqrt{x}}{ 2}\\
-\dst\frac{1}{ 4x^{3/2}}&\dst\frac{3}{ 4x^{5/2}}&\dst\frac{3}{ 4\sqrt{x}}
\end{vmatrix}=-\dst\frac{2}{ x^{3/2}}$;
$u_1'=-\dst\frac{FW_1}{ P_0W}=\dst-3x$;
$u_2'=\dst\frac{FW_2}{ P_0W}=\dst3x^2$;
$u_3'=-\dst\frac{FW_2}{ P_0W}=\dst1$;
$u_4'=\dst\frac{FW_4}{ P_0W}=\dst-x^3$;
$u_1=\dst-\frac{3x^2}{ 2}$;
$u_2=\dst x^3$;
$u_3=\dst x$;
$u_4=\dst-\frac{x^4}{ 4}$;
$y_p=u_1y_1+u_2y_2+u_3y_3=\dst \frac{x^{5/2}}{ 4}$.


\item
$W=\begin{vmatrix}x&x^2&x^3&x^4\\
1&2x&3x^2&4x^3\\0&2&6x&12x^2\\0&0&6&24x
\end{vmatrix}=12x^4$;
$W_1=\begin{vmatrix}x^2&x^3&x^4\\
2x&3x^2&4x^3\\2&6x&12x^2\end{vmatrix}=2x^6$;
$W_2=\begin{vmatrix}x&x^3&x^4\\
1&3x^2&4x^3\\0&6x&12x^2\end{vmatrix}=6x^5$;
$W_3=\begin{vmatrix}x&x^2&x^4\\
1&2x&4x^3\\0&2&12x^2\end{vmatrix}=6x^4$;
$W_4=\begin{vmatrix}x&x^2&x^3\\
1&2x&3x^2\\0&2&6x\end{vmatrix}=2x^3$;
$u_1'=-\dst\frac{FW_1}{ P_0W}=\dst-\frac{x^2}{ 6}$;
$u_2'=\dst\frac{FW_2}{ P_0W}=\dst\frac{x}{ 2}$;
$u_3'=-\dst\frac{FW_2}{ P_0W}=\dst-\frac{1}{ 2}$;
$u_4'=\dst\frac{FW_4}{ P_0W}=\dst\frac{1}{ 6x}$;
$u_1=\dst-\frac{x^3}{ 18}$;
$u_2=\dst\frac{x^2}{ 4}$;
$u_3=\dst-\frac{x}{ 2}$;
$u_4=\dst\frac{\ln|x|}{ 6}$;
$y_p=u_1y_1+u_2y_2+u_3y_3=\dst
\frac{x^4\ln|x|}{ 6}-\frac{11x^4}{ 36}$.
Since $-\dst\frac{11x^4}{ 36}$ satisfies the complementary equation we
take
$y_p=\dst\frac{x^4\ln|x|}{ 6}$.


\item
$W=\begin{vmatrix}
x&x^2&1/x&1/x^2\\
1&2x&-1/x^2&-2/x^3\\
0&2&2/x^3&6/x^4\\0&0&-6/x^4&-24/x^5
\end{vmatrix}=-72/x^6$;
$W_1=\begin{vmatrix}
x^2&1/x&1/x^2\\2x&-1/x^2&-2/x^3
\\2&2/x^3&6/x^4\end{vmatrix}=-12/x^4$;
$W_2=\begin{vmatrix}
x&1/x&1/x^2\\1&-1/x^2&-2/x^3\\
0&2/x^3&6/x^4\end{vmatrix}=-6/x^5$;
$W_3=\begin{vmatrix}
x&x^2&1/x^2\\1&2x&-2/x^3\\
0&2&6/x^4\end{vmatrix}=12/x^2$;
$W_4=\begin{vmatrix}
x&x^2&1/x\\1&2x&-1/x^2\\
0&2&2/x^3\end{vmatrix}=6/x$;
$u_1'=-\dst\frac{FW_1}{ P_0W}=\dst-2$;
$u_2'=\dst\frac{FW_2}{ P_0W}=\dst 1/x$;
$u_3'=-\dst\frac{FW_2}{ P_0W}=\dst2x^2$;
$u_4'=\dst\frac{FW_4}{ P_0W}=\dst-x^3$;
$u_1=\dst-2x$;
$u_2=\dst\ln|x|$;
$u_3=\dst2x^3/3$;
$u_4=\dst-x^4/4$;
$y_p=u_1y_1+u_2y_2+u_3y_3=\dst x^2\ln|x|-19x^2/12$.
Since $-19x^2/12$ satisfies the complementary equation we take
$y_p= x^2\ln|x|$.


\item
$W=\begin{vmatrix}
e^x&e^{2x}&e^x/x&\dst\frac{e^{2x}}{ x}\\
e^x&2e^{2x}&\dst\frac{e^x(x-1)}{ x^2}&\dst\frac{e^{2x}(2x-1)}{ x^2}\\
e^x&4e^{2x}&\dst\frac{e^x(x^2-2x+2)}{
x^3}&\dst\frac{2e^{2x}(2x^2-2x+1)}{ x^3}\\
e^x&8e^{2x}&\dst\frac{e^x(x^3-3x^2+6x-6)}{ x^4}
&\dst\frac{2e^{2x}(4x^3-6x^2+6x-3)}{ x^4}
\end{vmatrix}=-\dst\frac{e^{6x}}{ x^4}$;

$W_1=\begin{vmatrix}
e^{2x}&\dst\frac{e^x}{ x}&\dst\frac{e^{2x}}{ x}\\
2e^{2x}&\dst\frac{e^x(x-1)}{ x^2}&\dst\frac{e^{2x}(2x-1)}{ x^2}\\
4e^{2x}&\dst\frac{e^x(x^2-2x+2)}{ x^3}&\dst\frac{2e^{2x}(2x^2-2x+1)}{ x^3}
\end{vmatrix}=\dst\frac{e^{5x}}{ x^3}$;

$W_2=\begin{vmatrix}
e^x&\dst\frac{e^x}{ x}&\dst\frac{e^{2x}}{ x}\\
e^x&\dst\frac{e^x(x-1)}{ x^2}&\dst\frac{e^{2x}(2x-1)}{ x^2}\\
e^x&\dst\frac{e^x(x^2-2x+2)}{ x^3}&\dst\frac{2e^{2x}(2x^2-2x+1)}{ x^3}
\end{vmatrix}=-\dst\frac{e^{4x}}{ x^3}$;

$W_3=\begin{vmatrix}
e^x&e^{2x}&\dst\frac{e^{2x}}{ x}\\
e^x&2e^{2x}&\dst\frac{e^{2x}(2x-1)}{ x^2}\\
e^x&4e^{2x}&\dst\frac{2e^{2x}(2x^2-2x+1)}{ x^3}
\end{vmatrix}=\dst\frac{e^{5x}(2-x)}{ x^3}$;

$W_4=\begin{vmatrix}
e^x&e^{2x}&\dst\frac{e^x}{ x}\\
e^x&2e^{2x}&\dst\frac{e^x(x-1)}{ x^2}\\
e^x&4e^{2x}&\dst\frac{e^x(x^2-2x+2)}{ x^3}
\end{vmatrix}=\dst\frac{e^{4x}(x+2)}{ x^3}$;

$u_1'=-\dst\frac{FW_1}{ P_0W}=\dst3$;
$u_2'=\dst\frac{FW_2}{ P_0W}=\dst3e^{-x}$;
$u_3'=-\dst\frac{FW_2}{ P_0W}=\dst3(2-x)$;
$u_4'=\dst\frac{FW_4}{ P_0W}=\dst-3e^{-x}(x+2)$;
$u_1=\dst 3x$;
$u_2=\dst-3e^{-x}$;
$u_3=\dst3x(4-x)/2$;
$u_4=\dst3e^{-x}(x+3)$;
$y_p=u_1y_1+u_2y_2+u_3y_3=\dst\frac{3e^x(x^2+4x+6)}{ 2x}$.
Since
$\dst\frac{3e^x(2x+3)}{ x}$ is a solution of the complementary
equation we take $y_p=\dst\frac{3xe^x}{2}$.





\item
$W=\begin{vmatrix}
x&x^3&x\ln x\\1&3x^2&1+\ln x\\0&6x&1/x
\end{vmatrix}=-4x^2$;
$W_1=\begin{vmatrix}
x^3&x\ln x\\3x^2&1+\ln x
\end{vmatrix}=x^3-2x^3\ln x$;
$W_2=\begin{vmatrix}
x&x\ln x\\1&1+\ln x\end{vmatrix}=x$;
$W_3=\begin{vmatrix}x&x^3\\1&3x^2
\end{vmatrix}=2x^3$;
$u_1'=\dst\frac{FW_1}{ P_0W}=\dst2\ln x/x-\dst\frac{1}{ x}$;
$u_2'=-\dst\frac{FW_2}{ P_0W}=\dst\frac{1}{ x^3}$;
$u_3'=\dst\frac{FW_2}{ P_0W}=\dst-\frac{2}{ x}$;
$u_1=\dst(\ln x)^2-\ln x$;
$u_2=\dst-\frac{1}{ 2x^2}$;
$u_3=\dst-2\ln x$;
$y_p=u_1y_1+u_2y_2+u_3y_3=\dst-x(\ln x)^2-x\ln x-\frac{x}{ 2}$.
Since $-x\ln x-\dst\frac{x}{ 2}$ satisfies the complementary equation
we take $y_p=-x(\ln x)^2$


The general solution is $y=-x(\ln x)^2+c_1x+c_2x^3+c_3x\ln x$,
so
\[
\begin{bmatrix}y\\y'\\y''
\end{bmatrix}=
\begin{bmatrix}
-x(\ln x)^2\\ -(\ln x)^2-2\ln x\\-\dst\frac{2\ln x}{ x}-\frac{2}{ x}
\end{bmatrix}+
\begin{bmatrix}
x&x^3&x\ln x\\1&3x^2&1+\ln x\\0&6x&\dst\frac{1}{ x}
\end{bmatrix}
\begin{bmatrix}
c_{1}\\ c_{2}\\ c_{3}
\end{bmatrix}
\]
Setting $x=1$ and imposing the initial conditions yields
\[
\threecol44{-2}=
\threecol00{-2}+\threebythree110131061\threecol{c_1}{c_2}{c_3}.
\]
Solving this system yields $c_1=3$, $c_2=1$, $c_3=-2$.
Therefore,$y=-x(\ln x)^2+3x+x^3-2x\ln x$.


\item
$W=\begin{vmatrix}
e^x&e^{2x}&xe^{-x}\\e^x&2e^{2x}&e^{-x}(1-x)\\e^x&4e^{2x}&e
^{-x}(x-2)\end{vmatrix}=e^{2x}(6x-5)$;
$W_1=\begin{vmatrix}
e^{2x}&xe^{-x}\\2e^{2x}&e^{-x}(1-x)\end{vmatrix}=e^x(1-3x)$;
$W_2=\begin{vmatrix}
e^x&xe^{-x}\\e^x&e^{-x}(1-x)\end{vmatrix}=1-2x$;
$W_3=\begin{vmatrix}
e^x&e^{2x}\\e^x&2e^{2x}\end{vmatrix}=e^{3x}$;
$u_1'=\dst\frac{FW_1}{ P_0W}=\dst1-3x$;
$u_2'=-\dst\frac{FW_2}{ P_0W}=\dst e^{-x}(2x-1)$;
$u_3'=\dst\frac{FW_2}{ P_0W}=\dst e^{2x}$;
$u_1=\dst\frac {x(2-3x)}{ 2}$;
$u_2=\dst -e^{-x}(2x+1)$;
$u_3=\dst\frac {e^{2x}}{ 2}$;
$y_p=u_1y_1+u_2y_2+u_3y_3=\dst-e^x(3x^2+x+2)/2$. Since
$\dst-\frac{e^x}{2}$ is a solution of the complementary equation
we take $y_p=-\dst\frac{e^x(3x+1)x}{2}$.

The general solution is
$y=-\dst\frac{e^x(3x+1)x}{2}+c_1e^x+c_2e^{2x}+c_3xe^{-x}$, so
\[
\begin{bmatrix}
y\\y'\\y''
\end{bmatrix}=
\begin{bmatrix}
-\dst\frac{e^x(3x+1)x}{ 2}\\
-\dst\frac{e^x(3x^2+7x+1)}{ 2}\\
-\dst\frac{e^x(3x^2+13x+8)}{ 2}
\end{bmatrix}+
\begin{vmatrix}
e^x&e^{2x}&xe^{-x}\\e^x&2e^{2x}&e^{-x}(1-x)\\e^x&4e^{2x}&e
^{-x}(x-2)\end{vmatrix}
\threecol{c_1}{c_2}{c_3}.
\]
Setting $x=0$  and imposing the initial conditions yields
\[
\begin{bmatrix}
-4\\-\frac{3}{2}\\-19
\end{bmatrix}=
\begin{bmatrix}
0\\-\frac{1}{2}\\-4
\end{bmatrix}+
\threebythree11012114{-2}
\threecol{c_1}{c_2}{c_3}.
\]
Solving this system yields $c_1=-3$, $c_2=-1$, $c_3=4$.
Therefore,
$y=-\dst\frac{e^x(3x+1)x}{2}-3e^x-e^{2x}+4xe^{-x}$.





\item
$W=\begin{vmatrix}
x&x^2&e^x\\1&2x&e^x\\0&2&e^x\end{vmatrix}=e^x(x^2-2x+2)$;
$W_1=\begin{vmatrix}
x^2&e^x\\2x&e^x\end{vmatrix}=e^x(x^2-2x)$;
$W_2=\begin{vmatrix}x&e^x\\1&e^x\end{vmatrix}=e^x(x-1)$;
$W_3=\begin{vmatrix}x&x^2\\1&2x\end{vmatrix}=x^2$;
$u_1'=\dst\frac{FW_1}{ P_0W}=\dst x(x-2)$;
$u_2'=-\dst\frac{FW_2}{ P_0W}=\dst1-x$;
$u_3'=\dst\frac{FW_2}{ P_0W}=\dst x^2e^{-x}$;
$y_p=u_1y_1+u_2y_2+u_3y_3=\dst-\frac{x^4+6x^2+12x+12}{ 6}$. Since
$-\dst\frac{6x^2+12x}{6}$ is a solution of the complementary equations
we take $y_p=-\dst\frac{x^4+12}{6}$.

The general solution is
$y=-\dst\frac{x^4+12}{6}+c_1x+c_2x^2+c_3e^x$, so
\[
\begin{bmatrix}y\\y'\\y''
\end{bmatrix}=-
\begin{bmatrix}
(x^4+12)/6\\2x^3/3\\2x^2
\end{bmatrix}+
\begin{vmatrix}
x&x^2&e^x\\1&2x&e^x\\0&2&e^x\end{vmatrix}
\threecol{c_1}{c_2}{c_3}.
\]
Setting $x=0$ and imposing the initial conditions yields
\[
\threecol050=
-\threecol200+\threebythree001101021\threecol{c_1}{c_2}{c_3}.
\]
Solving this system yields $c_1=3$, $c_2=-1$, $c_3=2$.
Therefore,
$y=-\dst\frac{x^4+12}{6}+3x-x^2+2e^x$.



\item
$W=\begin{vmatrix}
x+1&e^x&e^{3x}\\1&e^x&3e^{3x}\\0&e^x&9e^{3x}
\end{vmatrix}=2e^{4x}(3x-1)$;
$W_1=\begin{vmatrix}
e^x&e^{3x}\\e^x&3e^{3x}\end{vmatrix}=2e^{4x}$;
$W_2=\begin{vmatrix}
x+1&e^{3x}\\1&3e^{3x}\end{vmatrix}=e^{3x}(3x+2)$;
$W_3=\begin{vmatrix}
x+1&e^x\\1&e^x\end{vmatrix}=xe^x$;
$u_1'=\dst\frac{FW_1}{ P_0W}=\dst 2e^x$;
$u_2'=-\dst\frac{FW_2}{ P_0W}=\dst-3x-2$;
$u_3'=\dst\frac{FW_2}{ P_0W}=\dst xe^{-2x}$;
$u_1=\dst 2e^x$;
$u_2=\dst-\frac{x(3x+4)}{2}$;
$u_3=\dst-\frac{e^{-2x}(2x+1)}{4}$;
$y_p=u_1y_1+u_2y_2+u_3y_3=\dst-\frac{e^x(6x^2+2x-7)}{4}$.
Since $\dst\frac{7e^x}{4}$ is a solution of the complementary
equation we take
$y_p=-\dst\frac{xe^x(3x+1)}{2}$

The general solution is
$y=-\dst\frac{xe^x(3x+1)}{2}+c_1(x+1)+c_2e^x+c_3e^{2x}$, so
\[
\begin{bmatrix}y\\y'\\y''
\end{bmatrix}=
\begin{bmatrix}
-xe^x(3x+1)/2\\
-e^x(3x^2+7x+1)/2\\
-e^x(3x^2+13x+8)/2
\end{bmatrix}+
\begin{vmatrix}
x+1&e^x&e^{3x}\\1&e^x&3e^{3x}\\0&e^x&9e^{3x}
\end{vmatrix}.
\]
Setting $x=0$ and imposing the initial conditions yields
\[
\begin{bmatrix}
\frac{3}{4}\\\frac{5}{4}\\\frac{1}{4}
\end{bmatrix}=
\begin{bmatrix}
0\\-\frac{1}{2}\\-4
\end{bmatrix}
+\threebythree111113019\threecol{c_1}{c_2}{c_3}.
\]
Solving this system yields $c_1=\dst\frac{1}{2}$, $c_2=-\dst\frac{1}{4}$,
$c_3=\dst\frac{1}{2}$. Therefore,
$y=-\dst\frac{xe^x(3x+1)}{2}+\frac{x+1}{2}-\frac{e^x}{4}+\frac{e^{2x}}{2}$.




\item
$W=\begin{vmatrix}
x&x^2&\dst\frac{1}{ x}&x\ln x\\1&2x&-\dst\frac{1}{ x^2}&\ln x+1\\
0&2&\dst\frac{2}{ x^3}&\dst\frac{1}{ x}\\0&0&-\dst\frac{6}{ x^2}&-\dst\frac{1}{ x^2}
\end{vmatrix}=-\dst\frac{12}{ x^3}$;
$W_1=\begin{vmatrix}
x^2&\dst\frac{1}{ x}&x\ln x\\2x&-\dst\frac{1}{ x^2}&\ln x+1\\
2&\dst\frac{2}{ x^3}&\dst\frac{1}{ x}\end{vmatrix}=6\dst\frac{\ln
x}{ x}-\dst\frac{3}{ x}$;
$W_2=\begin{vmatrix}
x&\dst\frac{1}{ x}&x\ln x\\1&-\dst\frac{1}{ x^2}&\ln x+1\\
0&\dst\frac{2}{ x^3}&\dst\frac{1}{ x}\end{vmatrix}=-\dst\frac{4}{
x^2}$;
$W_3=\begin{vmatrix}x&x^2&x\ln x\\
1&2x&\ln x+1\\0&2&\dst\frac{1}{ x}\end{vmatrix}=-x$;
$W_4=\begin{vmatrix}
x&x^2&\dst\frac{1}{ x}\\1&2x&-\dst\frac{1}{ x^2}\\0&2&\dst\frac{2}{ x^3}
\end{vmatrix}=\dst\frac{6}{ x}$;
$u_1'=-\dst\frac{FW_1}{ P_0W}=\dst9\ln x/2-\frac{9}{ 4}$;
$u_2'=\dst\frac{FW_2}{ P_0W}=\dst\frac{3}{ x}$;
$u_3'=-\dst\frac{FW_2}{ P_0W}=\dst-\frac{3x^2}{ 4}$;
$u_4'=\dst\frac{FW_4}{ P_0W}=\dst-\frac{9}{ 2}$;
$u_1=\dst\frac{9x\ln x}{ 2}-\frac{27x}{ 4}$;
$u_2=\dst3\ln x$;
$u_3=\dst-\frac{x^3}{ 4}$;
$u_4=\dst-\frac{9x}{ 2}$;
$y_p=u_1y_1+u_2y_2+u_3y_3=\dst3x^2\ln x-7x^2$.
Since $-7x^2$ satisfies the complementary equation we take
$y_p=3x^2\ln x$.

The general solution is $y=3x^2\ln x+c_1x+c_2x^2+\dst\frac{c_3}{ x}
+c_4x\ln x$, so
\[
\begin{bmatrix}y\\y'\\y''\\y'''
\end{bmatrix}=
\begin{bmatrix}
3x^2\ln x\\6x\ln x+3x\\6\ln x+9\\\dst\frac{6}{ x}
\end{bmatrix}+
\begin{vmatrix}
x&x^2&\dst\frac{1}{ x}&x\ln x\\1&2x&-\dst\frac{1}{ x^2}&\ln x+1\\
0&2&\dst\frac{2}{ x^3}&\dst\frac{1}{ x}\\0&0&-\dst\frac{6}{
x^2}&-\dst\frac{1}{ x^2}
\end{vmatrix}
\begin{bmatrix}
c_1\\c_2\\c_3\\c_4
\end{bmatrix}.
\]
Setting $x=1$ and imposing the initial conditions yields
\[
\begin{bmatrix}
-7\\-11\\-5\\6
\end{bmatrix}=
\begin{bmatrix}
0\\3\\9\\6
\end{bmatrix}+
\begin{vmatrix}1&1&1&0\\1&2&-1&1
\\0&2&2&1\\0&0&-6&-1\end{vmatrix}
\begin{bmatrix}
c_1\\c_2\\c_3\\c_4
\end{bmatrix}.
\]
Solving this system yields $c_1=0$, $c_2=-7$, $c_3=0$, $c_4=0$.
Therefore,
 $y=3x^2\ln x-7x^2$.



\item
$W=\begin{vmatrix}
x&\sqrt{x}&1/x&1/\sqrt{x}\\
1&1/2\sqrt{x}&-1/x^2&-1/2x^{3/2}\\
0&-1/4x^{3/2}&2/x^3&3/4x^{5/2}\\
0&3/8x^{5/2}&-6/x^4&-15/8x^{7/2}
\end{vmatrix}=-9/8x^6$;
$W_1=\begin{vmatrix}
\sqrt{x}&1/x&1/\sqrt{x}\\
1/2\sqrt{x}&-1/x^2&-1/2x^{3/2}\\
-1/4x^{3/2}&2/x^3&3/4x^{5/2}\end{vmatrix}=3/4x^4$;
$W_2=\begin{vmatrix}
x&1/x&1/\sqrt{x}\\1&-1/x^2&-1/2x^{3/2}\\
0&2/x^3&3/4x^{5/2}\end{vmatrix}=3/2x^{7/2}$;
$W_3=\begin{vmatrix}
x&\sqrt{x}&1/\sqrt{x}\\
1&1/2\sqrt{x}&-1/2x^{3/2}\\0&-1/4x^{3/2}&3/4x^{5/2}
\end{vmatrix}=-3/4x^2$;
$W_4=\begin{vmatrix}
x&\sqrt{x}&1/x\\1&1/2\sqrt{x}&-1/x^2\\
0&-1/4x^{3/2}&2/x^3\end{vmatrix}=-3/2x^{5/2}$;
$u_1'=-\dst\frac{FW_1}{ P_0W}=\dst 1/x$;
$u_2'=\dst\frac{FW_2}{ P_0W}=\dst-2/\sqrt{x}$;
$u_3'=-\dst\frac{FW_2}{ P_0W}=\dst-x$;
$u_4'=\dst\frac{FW_4}{ P_0W}=\dst2\sqrt{x}$;
$u_1=\dst\ln x$;
$u_2=\dst-4\sqrt{x}$;
$u_3=\dst-x^2/2$;
$u_4=\dst4x^{3/2}/3$;
$y_p=u_1y_1+u_2y_2+u_3y_3=\dst x\ln x-19x/6$.
since $-19x/6$ satisfies the complementary equation we take
$y_p=x\ln x$.

The general solution is $y=x\ln x+c_1x+c_2\sqrt x+c_3/x+c_4/\sqrt x$,
so
\[
\begin{bmatrix}y\\y'\\y''\\y'''
\end{bmatrix}=
\begin{bmatrix}
x\ln x\\\ln x+1\\1/x\\-1/x^2
\end{bmatrix}+
\begin{vmatrix}
x&\sqrt{x}&1/x&1/\sqrt{x}\\
1&1/2\sqrt{x}&-1/x^2&-1/2x^{3/2}\\
0&-1/4x^{3/2}&2/x^3&3/4x^{5/2}\\
0&3/8x^{5/2}&-6/x^4&-15/8x^{7/2}
\end{vmatrix}
\begin{bmatrix}
c_1\\c_2\\c_3\\c_4
\end{bmatrix}.
\]
Setting $x=1$ and imposing the initial conditions yields
\[
\begin{bmatrix}
2\\0\\4\\\-\frac{37}{4}
\end{bmatrix}=
\begin{bmatrix}
0\\1\\1\\-1
\end{bmatrix}+
\begin{vmatrix}1&1&1&1\\1&1/2&-1&-1/2\\
0&-1/4&2&3/4\\0&3/8&-6&-15/8\end{vmatrix}
\begin{bmatrix}
c_1\\c_2\\c_3\\c_4
\end{bmatrix}.
\]
Solving this system yields $c_1=1$, $c_2=-1$, $c_3=1$, $c_4=1$.
Therefore,  $y=x\ln x+x-\sqrt x+\dst\frac{1}{ x}+\frac{1}{\sqrt x}$.


\item
\begin{alist}
\item
Since
$u'_j=(-1)^{n-j}\dst\frac{FW_j}{ P_0W}$  ($1\le j\le n$), the argument
used in the derivation of the method of variation of parameters
implies that $y_p$ is a solution of (A).

\item Follows immediately from \part{a}, since $u_j(x_0)=0$,
$j=1,2,\dots,n$.

\item Expand the determinant in cofactors of its $n$th row.

\item Just differentiate the determinant $n-1$ times.

\item If $0\le j\le n-2$, then
$\dst\left.\frac{\partial^{j}G(x,t)}{\partial x^j}\right|_{x=t}$
has two identical rows, and is therefore zero, while
$\dst\left.\frac{\partial^{n-1}G(x,t)}{\partial x^j}\right|_{x=t}=W(t)$

\item Since
$y_p(x)=\dst\int_{x_0}^x G(x,t)F(t)\,dt$,
$y_p'(x)=G(x,x)F(x)+\dst\int_{x_0}^x\frac{\partial
G(x,t)}{\partial x}F(t)\,dt$. But $G(x,x)=0$ from \part{e}, so
$y_p'(x)=\dst\int_{x_0}^x\frac{\partial
G(x,t)}{\partial x}F(t)\,dt$. Repeating this argument for
$j=1,\dots,n$ and invoking \part{e} each time yields the conclusion.
\end{alist}


\item
\begin{align*}
\begin{vmatrix}
y_1(t)&y_1(t)&y_2(t)\\
y_1'(t)&y_1'(t)&y_2'(t)\\
y_1(x)&y_1(x)&y_2(x)\\
\end{vmatrix} &=
\begin{vmatrix}
t&t^2&1/t\\
1&2t&-1/t^2\\
x&x^2&1/x
\end{vmatrix}\\
&=
x\begin{vmatrix}t^2&1/t\\2t&-1/t^2\end{vmatrix}
-x^2\begin{vmatrix}t&1/t\\1&-1/t^2\end{vmatrix}
+\frac{1}{ x}\begin{vmatrix}t&t^2\\1&2t\end{vmatrix}\\
&=-3x+2\frac{x^2}{ t}+\frac{t^2}{ x}=\frac{(x-t)^2(2x+t)}{ xt}.
\end{align*}
Since $P_0(t)=t^3$ and
$W(t)=\dst\begin{vmatrix}
t&t^2&1/t\\
1&2t&-1/t^2\\
x&x^2&1/x
\end{vmatrix}=\frac{6}{ t}$,
$G(x,t)=\dst\frac {(x-t)^2(2x+t)}{6xt^3}$,
so
$y_p=\dst{\int_{x_0}^x \frac{(x-t)^2(2x+t)}{6xt^3}
F(t)\,dt}$.


\item
\begin{align*}
\begin{vmatrix}
y_1(t)&y_1(t)&y_2(t)\\
y_1'(t)&y_1'(t)&y_2'(t)\\
y_1(x)&y_1(x)&y_2(x)
\end{vmatrix}&=
\begin{vmatrix}
t&1/t&e^t/t\\
1&-1/t^2&e^t(1/t-1/t^2)\\
x&1/x&e^x/x\\
\end{vmatrix}\\
&=x\begin{vmatrix}1/t&e^t/t\\-1/t^2&e^t(1/t-1/t^2)\end{vmatrix}
-\frac{1}{ x}\begin{vmatrix}t&e^t/t\\1&e^t(1/t-1/t^2)\end{vmatrix}
+\frac{e^x}{ x}\begin{vmatrix}t&1/t\\1&-1/t^2\end{vmatrix}\\
&=\frac{xe^t}{ t^2}-\frac{e^t(t-2)}{ xt}-\frac{2e^x}{ xt}
=\frac{x^2e^t-e^tt(t-2)-2te^x}{ xt^2}
\end{align*}
Since $P_0(t)=t(1-t)$ and
$W(t)=\dst\begin{vmatrix}
t&1/t&e^t/t\\
1&-1/t^2&e^t(1/t-1/t^2)\\
0&2/t^3&e^t(1/t-2/t^2+2/t^3)
\end{vmatrix}=\frac{2e^t(1-t)}{ t^3}$,
$G(x,t)=\dst\frac{x^2-t(t-2)-2te^{(x-t)}}{ 2x(t-1)^2}$, so
$y_p=\dst{\int_{x_0}^x \frac{x^2-t(t-2)-2te^{(x-t)}}{ 2x(t-1)^2}
F(t)\,dt}$.


\item
\begin{align*}
\begin{vmatrix}
y_1(t)&y_1(t)&y_2(t)&y_3(t)\\
y_1'(t)&y_1'(t)&y_2'(t)&y_3'(t)\\
y_1''(t)&y_1''(t)&y_2''(t)&y_3''(t)\\
y_1(x)&y_1(x)&y_2(x)&y_3(x)
\end{vmatrix}&=
\begin{vmatrix}
1&t&t^2&1/t\\
0&1&2t&-1/t^2\\
0&0&2&2/t^3\\
1&x&x^2&1/x\\
\end{vmatrix}\\
&=-\begin{vmatrix}
t&t^2&1/t\\
1&2t&-1/t^2\\
0&2&2/t^3\end{vmatrix}
+x\begin{vmatrix}
1&t^2&1/t\\
0&2t&-1/t^2\\
0&2&2/t^3
\end{vmatrix}\\
&-x^2\begin{vmatrix}
1&t&1/t\\0&1&-1/t^2\\
0&0&2/t^3\end{vmatrix}
+\frac{1}{ x}\begin{vmatrix}
1&t&t^2\\0&1&2t\\0&0&2\end{vmatrix}\\
&=-\frac{6}{ t}+\frac{6x}{ t^2}-\frac{2x^2}{ t^3}+\frac{2}{ x}
=\frac{2(t-x)^3}{ xt^3}.
\end{align*}
Since $P_0(t)=t$ and
$W(t)=\begin{vmatrix}
1&t&t^2&1/t\\
0&1&2t&-1/t^2\\
0&0&2&2/t^3\\
0&0&0&-6/t^4
\end{vmatrix}=-\dst\frac{12}{ t^4}$,
$G(x,t)=\dst\frac{(x-t)^3}{6x}$, so
$y_p=\dst\int_{x_0}^x \frac{(x-t)^3}{6x}F(t)\,dt$.



\item
\begin{align*}
\begin{vmatrix}
y_1(t)&y_1(t)&y_2(t)&y_3(t)\\
y_1'(t)&y_1'(t)&y_2'(t)&y_3'(t)\\
y_1''(t)&y_1''(t)&y_2''(t)&y_3''(t)\\
y_1(x)&y_1(x)&y_2(x)&y_3(x)
\end{vmatrix}&=
\begin{vmatrix}
1&t^2&e^{2t}&e^{-2t}\\
0&2t&2e^{2t}&-2e^{-2t}\\
0&2&4e^{2t}&4e^{-2t}\\
1&x^2&e^{2x}&e^{-2x}
\end{vmatrix}\\
&=-\begin{vmatrix}t^2&e^{2t}&e^{-2t}\\
2t&2e^{2t}&-2e^{-2t}\\2&4e^{2t}&4e^{-2t}\end{vmatrix}
+x^2\begin{vmatrix}1&e^{2t}&e^{-2t}\\
0&2e^{2t}&-2e^{-2t}\\0&4e^{2t}&4e^{-2t}\end{vmatrix}\\
&\qquad-e^{2x}\begin{vmatrix}1&t^2&e^{-2t}\\
0&2t&-2e^{-2t}\\0&2&4e^{-2t}\end{vmatrix}
+e^{-2x}\begin{vmatrix}1&t^2&e^{2t}\\0&2t&2e^{2t}\\
0&2&4e^{2t}\end{vmatrix}\\
&=-(16t^2-8)+16x^2-e^{2(x-t)t}(8t+4)+e^{-2(x-t)}(8t-4).
\end{align*}
Since $P_0(t)=t$ and
$W(t)=\begin{vmatrix}
1&t^2&e^{2t}&e^{-2t}\\
0&2t&2e^{2t}&-2e^{-2t}\\
0&2&4e^{2t}&4e^{-2t}\\
0&0&8e^{2t}&-8e^{-2t}
\end{vmatrix}=-128t$,

$G(x,t)=\dst\frac{e^{2(x-t)}(1+2t)+e^{-2(x-t)}(1-2t)-4x^2+4t^2-2}{32t^2}$,
so

$y_p=\dst\int_{x_0}^x
\frac{e^{2(x-t)}(1+2t)+e^{-2(x-t)}(1-2t)-4x^2+4t^2-2}{32t^2}
F(t)\,dt$.
\end{sectionSolutions}

\chapter{Linear Systems of Differential Equations}

\section{Introduction to Systems of Differential Equations}

\begin{sectionSolutions}
\item
$Q_1'=(\mathrm{rate in})_1-(\mathrm{rate out})_1$ and $Q_2'=(\mathrm{rate
in})_2-(\mathrm{rate out})_2$.

The volumes of the solutions in $T_1$ and $T_2$ are $V_1(t)=100+2t$
and $V_2(t)=100+3t$, respectively. $T_1$ receives salt from the
external source at the rate of $\qty{2}{\pound\per\gallon}\times\qty{6}{\gallon\per\minute}=\qty{12}{\pound\per\minute}$, and from $T_2$ at the rate of
$T_2\unit{\pound\per\gallon}\times\qty{1}{\gallon\per\minute}=\dst\frac{1}{100+3t}Q_2\unit{\pound\per\minute}$. Therefore, (A) $\mathrm{(rate
in)}_1= 12+\dst\frac{1}{100+3t}Q_2$. Solution leaves $T_1$ at \qty{5}{\gallon\per\minute},
since \qty{3}{\gallon\per\minute} are drained and \qty{2}{\gallon\per\minute} are pumped to $T_2$; hence
(B) $(\mathrm{rate out})_1=T_1\unit{\pound\per\gallon}\times
\qty{5}{\gallon\per\minute}
=\dst\frac{1}{100+2t}Q_1\times5=\dst\frac{5}{100+2t}Q_1$. Now (A) and (B)
imply that (C) $Q_1'=\dst{12-\frac{5}{100+2t}Q_1+\frac{1}{100+3t}Q_2}$.

$T_2$ receives salt from the external source at the rate of
$\qty{1}{\pound\per\gallon}\times\qty{5}{\gallon\per\minute}=\qty{5}{\pound\per\minute}$, and from $T_1$
at the rate of $T_1\unit{\pound\per\gallon}\times\qty{2}{\gallon\per\minute}=\dst\frac{1}{100+2t}Q_1\times2=\dst\frac{1}{50+t}Q_1 \unit{\pound\per\minute}$.
Therefore, (D) $\mathrm{(rate in)}_2= 5+\dst\frac{1}{50+t}Q_1$. Solution
leaves $T_2$ at \qty{4}{\gallon\per\minute}, since
\qty{3}{\gallon\per\minute} are drained and \qty{1}{\gallon\per\minute}
is pumped to $T_1$; hence (E)
$(\mathrm{rate out})_2=T_2\unit{\pound\per\gallon}\times\qty{4}{\gallon\per\minute}
=\dst\frac{1}{100+3t}Q_2\times4=\dst\frac{4}{100+3t}Q_2$. Now (D) and (E)
imply that (F) $Q_2'=\dst{5+\frac{1}{50+t}Q_1-\frac{4}{100+3t}Q_2}$. Now
(C) and (F) form the desired system.




\item
$m{\vect X}''=-\alpha{\vect{X}'}-mgR^2\dst\frac{{\vect X}}{\|{\vect X}\|^3}$;
see Example~10.1.3.

\addtocounter{enumi}{2}
\item%10.1.8
\begin{align*}
I_{1i}&=g_1(t_i,y_{1i},y_{2i}),\\
J_{1i}&=g_2(t_i,y_{1i},y_{2i}),\\
I_{2i}&=g_1\left(t_i+h,y_{1i}+hI_{1i},y_{2i}+hJ_{1i}\right),\\
J_{2i}&=g_2\left(t_i+h,y_{1i}+hI_{1i},y_{2i}+hJ_{1i}\right),\\
y_{1,i+1}&=y_{1i}+\frac{h}{2}(I_{1i}+I_{2i}),\\
y_{2,i+1}&=y_{2i}+\frac{h}{2}(J_{1i}+J_{2i}).
\end{align*}

\end{sectionSolutions}

\section{Linear Systems of Differential  Equations}

\begin{sectionSolutions}
\addtocounter{enumi}{4}
\item%10.2.6
Let  $y_i=y^{(i-1)}$, $i=1,2,\dots,n$; then $y_i'=y_{i+1}$,
$i=1,2,\dots,n-1$ and $P_0(t)y_n'+P_1(t)y_n+\cdots+P_n(t)y_1=F(t)$, so
\[A=-\frac{1}{ P_0}
\begin{bmatrix}
0&1&0&\cdots&0\\
0&0&1&\cdots&0\\
\vdots&\vdots&\vdots&\ddots&\vdots\\
0&0&0&\cdots&1\\
P_n&P_{n-1}&P_{n-2}&\cdots&P_1
\end{bmatrix}
\quad\text{and}\quad
{\vect f}=\frac{1}{ P_0}
\begin{bmatrix}
0\\ 0\\\vdots\\0\\F
\end{bmatrix}.
\]
If $P_0,P_1,\dots,P_n$ and $F$ are continuous and $P_0$
has no zeros on $(a,b)$, then $P_1/P_0,\dots,P_n/P_0$  and
$F/P_0$ are continuous on $(a,b)$.

\addtocounter{enumi}{-1}
\item%10.2.7
\begin{alist}
\item
$(c_1P+c_2Q)'_{ij}=
(c_1p_{ij}+c_2q_{ij})'=c_1p_{ij}'+c_2q_{ij}'=(c_1P'+c_2Q')_{ij}$;
hence
$(c_1P+c_2Q)'=c_1P'+c_2Q'$.

\item Let $P$ be $k\times r$ and $Q$ be $r\times s$; then
$PQ$ is $k\times s$ and $(PQ)_{ij}=\dst\sum_{\ell=1}^rp_{i\ell}q_{\ell j}$.
Therefore,
$(PQ)_{ij}'=\dst\sum_{\ell=1}^rp_{i\ell}'q_{\ell j}
+\sum_{\ell=1}^rp_{i\ell}q_{\ell j}'=(P'Q)_{ij}+(PQ')_{ij}$. Therefore,
$(PQ)'=P'Q+PQ'$.
\end{alist}

\addtocounter{enumi}{1}
\item%10.2.10
\begin{alist}
\item
From Exercise~10.2.7\part{b} with $P=Q=X$,
$(X^2)'=(XX)'=X'X+XX'$.

\item By starting from
 Exercise~10.2.7\part{b} and using induction it can be shown
if $P_1,P_2,\dots,P_n$ are square matrices of the same order, then
 $(P_1P_2\cdots P_n)'=P_1'P_2\cdots P_n+P_1P_2'\cdots P_n+\cdots+
P_1P_2\cdots P_n'$. Taking $P_1=P_2=\cdots=P_n=X$  yields
(A) $(Y^n)'=Y'Y^{n-1}+Y Y'Y^{n-2}+Y^2Y'Y^{n-3}+\cdots+Y^{n-1}Y'=\dst{
\sum_{r=0}^{n-1}}Y^r Y'Y^{n-r-1}$.

\item
If $Y$ is a scalar function, then (A) reduces to
the familiar result $(Y^n)'=nY^{n-1}Y'$.
\end{alist}

\item
From Exercise~10.2.6, the initial value problem
(A) $P_0(x)y^{(n)}+P_1(x)y^{(n-1)}+\cdots+P_n(x)y=F(x)$,
$y(x_0)=k_0,\ y'(x_0)=k_1,\dots,\ y^{(n-1)}(x_0)=k_{n-1}$ is
equivalent to the initial value problem (B) ${\vect{ y'}}=A(t){\vect y}+{\vect
f}(t)$, with
\[A=-\frac{1}{ P_0}
\begin{bmatrix}
0&1&0&\cdots&0\\
0&0&1&\cdots&0\\
\vdots&\vdots&\vdots&\ddots&\vdots\\
0&0&0&\cdots&1\\
P_n&P_{n-1}&P_{n-2}&\cdots&P_1
\end{bmatrix},
{\vect f}=\frac{1}{ P_0}
\begin{bmatrix}
0\\ 0\\\vdots\\0\\F
\end{bmatrix},
\quad\text{and}\quad
{\vect k}=
\begin{bmatrix}
k_0\\k_1\\\vdots\\k_{n-1}
\end{bmatrix}.
\]
Since Theorem~10.2.1 implies that (B) has a unique solution on
$(a,b)$, it follows that (A) does also.
\end{sectionSolutions}

\section{Basic Theory of Homogeneous Linear System}

\begin{sectionSolutions}
\item
\begin{alist}
\item
The system equivalent of (A) is (B) ${\vect
y}'=-\dst\frac{1}{ P_0(x)}\twobytwo01{P_2(x)}{P_1(x)}{\vect y}$, where
${\vect y}=\twocol{y}{y'}$. Let ${\vect y}_1=\twocol{y_1}{y_1'}$ and
 ${\vect y}_1=\twocol{y_2}{y_2'}$. Then the Wronskian of $\{{\vect
y}_1,{\vect y}_2\}$ as defined in this section is
$\twobytwo{y_1}{y_2}{y_1'}{y_2'}=W$.

\item
The trace of the matrix in (B) is $-P_1(x)/P_0(x)$, so Eqn.~10.3.6 implies
that
$W(x)=W(x_0)\dst\exp\left\{-\int^x_{x_0}\frac{P_1(s)}{ P_0(s)}\,
ds\right\}$.
\end{alist}

\item
\begin{alist}
\item See the solution of Exercise~9.1.18.

\item
$
\begin{vmatrix}
y_{11}'&y_{12}'\\y_{21}&y_{22}
\end{vmatrix}=
\begin{vmatrix}
a_{11}y_{11}+a_{12}y_{21}&a_{11}y_{12}+a_{12}y_{22}\\
y_{21}&y_{22}
\end{vmatrix}=
a_{11}\begin{vmatrix}
y_{11}&y_{12}\\y_{21}&y_{22}
\end{vmatrix}
+a_{12}\begin{vmatrix}
y_{21}&y_{22}\\y_{21}&y_{22}
\end{vmatrix}=a_{11}W+a_{12}0=a_{11}W$.
Similarly,
$
\begin{bmatrix}
y_{11}&y_{12}\\y_{21}'&y_{22}'
\end{bmatrix}=a_{22}W$.
\end{alist}


\item
\begin{alist}
\item From  the equivalence of Theorem~10.3.3\part{b}and \part{e},
$Y(t_0)$ is invertible.

\item From  the equivalence of Theorem~10.3.3\part{a}and \part{b},
the solution of the initial value problem is ${\vect y}=Y(t){\vect
c}$, where ${\vect c}$ is a constant vector. To satisfy ${\vect
y}(t_0)={\vect k}$, we must have
$Y(t_0){\vect c}={\vect k}$, so ${\vect c}=Y^{-1}(t_0){\vect k}$
and ${\vect y}=Y^{-1}(t_0)Y(t){\vect k}$.
\end{alist}

\item
\begin{alist}[start=2]
\item
${\vect y}=\begin{bmatrix}
e^{-4t} \\ e^{-4t}\end{bmatrix}+c_2
\begin{bmatrix}-2e^{3t}
\\ 5e^{3t}\end{bmatrix}$ where
$\begin{cases}
c_1-2c_2&=10\\c_1+5c_2&=-4
\end{cases},
$
 so $c_1=6$, $c_2=-2$, and
${\vect y}=\dst\begin{bmatrix}6e^{-4t}+4e^{3t}\\6e^{-4t}-10e^{3t}
\end{bmatrix}$.

\item
$Y(t)=\begin{bmatrix}e^{-4t}&-2e^{3t}\\
e^{-4t}&5e^{3t}\end{bmatrix}$;
$Y(0)=\begin{bmatrix}1&-2\\1&5\end{bmatrix}$;
$Y^{-1}(0)=
\dst\frac{1}{7}\begin{bmatrix}5&2\\-1&1\end{bmatrix}$;
${\vect y}=Y(t)Y^{-1}(0){\vect k}=
\dst\frac{1}{7}\begin{bmatrix}5e^{-4t}+2e^{3t}&2e^{-4t}-2e^{3t}
\\5e^{-4t}-5e^{3t}&2e^{-4t}+5e^{3t}\end{bmatrix}{\vect k}$.
\end{alist}

\item
\begin{alist}[start=2]
\item
${\vect y}_1=c_1\begin{bmatrix} e^{3t} \\e^{3t}
\end{bmatrix}+c_2\begin{bmatrix}e^t \\
-e^t\end{bmatrix}$,
where
$\begin{cases}
c_1+c_2&=2\\c_1-c_2&=8
\end{cases},
$
 so $c_1=5$, $c_2=-3$, and
${\vect y}=\dst\begin{bmatrix}5e^{3t}-3e^t\\5e^{3t}+3e^t
\end{bmatrix}$.

\item $Y(t)=\begin{bmatrix}e^{3t}&e^t\\
3e^{3t}&-e^t\end{bmatrix}$;
$Y(0)=\begin{bmatrix}1&1\\1&-1\end{bmatrix}$;

$Y^{-1}(0)=
\dst\frac{1}{2}\begin{bmatrix}1&1\\1&-1\end{bmatrix}$;
${\vect y}=Y(t)Y^{-1}(0){\vect k}=
\dst\frac{1}{2}\begin{bmatrix}e^{3t}+e^t&e^{3t}-e^t
\\e^{3t}-e^t&e^{3t}+e^t\end{bmatrix}{\vect k}$.
\end{alist}

\item
\begin{alist}[start=2]
\item
${\vect y}=c_1\begin{bmatrix}-e^{-2t} \\ 0 \\ e^{-2t}
\end{bmatrix}+c_2
\begin{bmatrix}-e^{-2t} \\ e^{-2t} \\
0\end{bmatrix}+
c_3\begin{bmatrix} e^{4t} \\ e^{4t} \\ e^{4t}\end{bmatrix}$,
where
$\begin{cases}
-c_1-c_2+c_3&=0\\c_2+c_3&=-9\\c_1+c_3&=12
\end{cases},
$
 so $c_1=11$, $c_2=-10$, $c_3=1$, and
${\vect y
}=\dst\frac{1}{3}\begin{bmatrix}-e^{-2t}+e^{4t}\\-10e^{-2t}+e^{4t}\\
11e^{-2t}+e^{4t}\end{bmatrix}$.


\item $Y(t)=
\begin{bmatrix}-e^{-2t}&-e^{-2t}&e^{4t}\\
0&e^{-2t}&e^{4t}\\e^{-2t}&0&e^{4t}\end{bmatrix}$;
$Y(0)=
\begin{bmatrix}-1&-1&1\\0&1&1\\1&0&1\end{bmatrix}$;
$Y^{-1}(0)=
\dst\frac{1}{3}\begin{bmatrix}-1&-1&2\\-1&2&-1
\\1&1&1\end{bmatrix}$;
${\vect y}=Y(t)Y^{-1}(0){\vect k}=
\dst\frac{1}{3}\begin{bmatrix}2e^{-2t}+e^{4t}&-e^{-2t}+e^{4t}
&-e^{-2t}+e^{4t}
\\-e^{-2t}+e^{4t}&2e^{-2t}+e^{4t}&-e^{-2t}+e^{4t}\\
-e^{-2t}+e^{4t}&-e^{-2t}+e^{4t}&2e^{-2t}+e^{4t}
\end{bmatrix}{\vect k}$.
\end{alist}

\item
If $Y$  and $Z$ are both fundamental matrices for ${\vect y}'=A(t){\vect y
}$, then $Z=CY$, where $C$ is a constant invertible matrix. Therefore,
$ZY^{-1}=C$ and $YZ^{-1}=C^{-1}$.


\item
\begin{alist}
\item
The Wronskian  of $\{{\vect y}_1,{\vect y}_2,\dots,{\vect y}_n\}$ equals
one when $t=t_0$. Apply Theorem~10.3.3.

\item
Let $Y$ be the matrix with columns $\{{\vect y}_1,{\vect y}_2,\dots,{\vect
y}_n\}$. From \part{a}, $Y$ is a fundamental matrix for ${\vect
y}'=A(t){\vect y}$ on $(a,b)$. From Exercise~10.3.15\part{b},
so is $Z=YC$ if $C$ is any invertible constant matrix.
\end{alist}

\item
\begin{alist}
\item
$\Gamma_1'(t)=Z'(t)Z(s)=AZ(t)Z(s)=A\Gamma(t)$ and $\Gamma_1(0)=Z(s)$,
since $Z(0)=I$.
$\Gamma_2'(t)=Z'(t+s)=AZ(t+s)=A\Gamma_2(t)$ (since $A$ is constant)
and $\Gamma_2(0)=Z(s)$. Applying Theorem~10.2.1  to the columns of
$\Gamma_1$ and $\Gamma_2$  shows that $\Gamma_1=\Gamma_2$.

\item With $s=-t$, \part{a} implies that $Z(t)Z(-t)=Z(0)=I$;
therefore $(Z(t))^{-1}=Z(-t)$.

\item $e^{0\cdot A}=I$ is analogous to $e^{o\cdot a}=e^{0}=1$ when
$a$ is a scalar, while $e^{(t+s)A}=e^{tA}e^{sA}$  is analogous
to $e^{(t+s)a}=e^{ta}e^{sa}$ when $a$ is a scalar.
\end{alist}
\end{sectionSolutions}

\section{Constant Coefficient Homogeneous Systems I}

\begin{sectionSolutions}
\item
$\dst\frac{1}{4}{\begin{vmatrix}-5-4\lambda&3\\
3&-5-4\lambda\end{vmatrix}}
=(\lambda+1/2)(\lambda+2)$.
Eigenvectors  associated with $\lambda_1=-1/2$  satisfy
$\dst{\begin{bmatrix}-3&3\\-4&4
\end{bmatrix}\begin{bmatrix}
x_1\\x_2\end{bmatrix}=\begin{bmatrix} 0\\0\end{bmatrix}}$,
so $x_1=x_2$.  Taking $x_2=1$ yields
${\vect y}_1=\dst{\begin{bmatrix}1 \\
1\end{bmatrix}}e^{-t/2}$.
Eigenvectors  associated with $\lambda_2=-2$
satisfy
$\dst{\begin{bmatrix}\frac{3}{4}&\frac{3}{4}\\1&1
\end{bmatrix}\begin{bmatrix}
x_1\\x_2\end{bmatrix}=\begin{bmatrix} 0\\0\end{bmatrix}}$,
so $x_1=-x_2$.  Taking $x_2=1$ yields
${\vect y}_2=\dst{\begin{bmatrix} -1\\
1\end{bmatrix}}e^{-2t}$. Hence
 ${\vect y}=\dst{ c_1\begin{bmatrix} 1\\1\end{bmatrix}e^{-t/2}
 +c_2\begin{bmatrix}-1\\1\end{bmatrix}e^{-2t}}$.


\item
$\dst{\begin{vmatrix}-1-\lambda&-4\\-1&-1-\lambda\end{vmatrix}}
=(\lambda-1)(\lambda+3)$.
Eigenvectors  associated with $\lambda_1=-3$ satisfy
$\dst{\begin{bmatrix}2&-4\\1&-2
\end{bmatrix}\begin{bmatrix}
x_1\\x_2\end{bmatrix}=\begin{bmatrix} 0\\0\end{bmatrix}}$,
so $x_1=2x_2$.  Taking $x_2=1$ yields
${\vect y}_1=\dst{\begin{bmatrix}
2\\1\end{bmatrix}}e^{-3t}$.
Eigenvectors  associated with $\lambda_2=1$  satisfy
$\dst{\begin{bmatrix}-2&-4\\-1&-2
\end{bmatrix}\begin{bmatrix}
x_1\\x_2\end{bmatrix}=\begin{bmatrix} 0\\0\end{bmatrix}}$,
so $x_1=-2x_2$.  Taking $x_2=1$ yields
${\vect y}_2=\dst{\begin{bmatrix} -2\\1\end{bmatrix}}e^t$.
Hence  ${\vect y}=\dst{ c_1\begin{bmatrix} 2\\1\end{bmatrix}e^{-3t}
+c_2\begin{bmatrix}-2\\1\end{bmatrix}e^t}$.


\item
$\dst{\begin{vmatrix}4-\lambda&-3\\2&-1-\lambda\end{vmatrix}}
=(\lambda-2)(\lambda-1)$.
Eigenvectors  associated with $\lambda_1=2$  satisfy
$\dst{\begin{bmatrix}2&-3\\2&-3
\end{bmatrix}\begin{bmatrix}
x_1\\x_2\end{bmatrix}=\begin{bmatrix} 0\\0\end{bmatrix}}$,
so $x_1=\dst\frac{3}{2}x_2$.  Taking $x_2=2$ yields
${\vect y}_1=\dst{\begin{bmatrix}3\\2\end{bmatrix}}e^{2t}$.
Eigenvectors  associated with $\lambda_2=1$ satisfy
$\dst{\begin{bmatrix}3&-3\\2&-2
\end{bmatrix}\begin{bmatrix}
x_1\\x_2\end{bmatrix}=\begin{bmatrix} 0\\0\end{bmatrix}}$,
so $x_1=x_2$.  Taking $x_2=1$ yields
${\vect y}_2=\dst{\begin{bmatrix}1\\1\end{bmatrix}}e^t$.
Hence
 ${\vect y}=\dst{c_1\begin{bmatrix}3\\2\end{bmatrix}e^{2t}
 +c_2\begin{bmatrix}1\\1\end{bmatrix}e^t}$.


\item
$\dst{\begin{vmatrix}1-\lambda&-1&-2\\1&-2-\lambda&-3 \\
-4&1&-1-\lambda\end{vmatrix}}=-(\lambda+3)(\lambda+1)(\lambda-2)$.
The eigenvectors associated with
 with $\lambda_1=-3$ satisfy the system with  augmented matrix
$\dst{\begin{bmatrix}4&-1&-2&\vdots&0\\1&1
&-3&\vdots&0\\-4&1&2&\vdots&0
\end{bmatrix}}$,
which is row equivalent to
$\dst{\begin{bmatrix}1&0&-1&\vdots&0\\0&1&-2&
\vdots&0\\0&0&0&\vdots&0\end{bmatrix}}$.
Hence $x_1=x_3$ and $x_2=2x_3$.  Taking $x_3=1$ yields
${\vect y}_1=\dst{\begin{bmatrix} 1\\2\\1
\end{bmatrix}}e^{-3t}$.
The eigenvectors associated with
 with $\lambda_2=-1$ satisfy the system with  augmented matrix
$\dst{\begin{bmatrix}2&-1&-2&\vdots&0\\1&-1
&-3&\vdots&0\\-4&1&0&\vdots&0
\end{bmatrix}}$,
which is row equivalent to
$\dst{\begin{bmatrix}1&0&1&\vdots&0\\0&1&4&0
\vdots&0\\0&0&0&\vdots&0\end{bmatrix}}$.
Hence $x_1=-x_3$ and $x_2=-4x_3$.  Taking $x_3=1$ yields
${\vect y}_2=\dst{\begin{bmatrix}-1\\-4\\ 1
\end{bmatrix}}e^{-t}$.
The eigenvectors associated with
 with $\lambda_3=2$ satisfy the system with  augmented matrix
$\dst{\begin{bmatrix}-1&-1&-2&\vdots&0\\1&
-4&-3&\vdots&0\\-4&1&-3&\vdots&0
\end{bmatrix}}$,
which is row equivalent to
$\dst{\begin{bmatrix}1&0&1&\vdots&0\\0&1&1&
\vdots&0\\0&0&0&\vdots&0\end{bmatrix}}$.
Hence $x_1=-x_3$ and $x_2=-x_3$.  Taking $x_3=1$ yields
${\vect y}_3=\dst{\begin{bmatrix}-1\\-1\\1
\end{bmatrix}}e^{2t}$. Hence
 ${\vect y}=\dst{c_1\begin{bmatrix} 1\\2\\1
\end{bmatrix}e^{-3t}+c_2\begin{bmatrix}-1\\-4\\1
\end{bmatrix}e^{-t}+c_3\begin{bmatrix}-1\\-1\\1\end{bmatrix}e^{2t}}$.


\item
$\dst{\begin{vmatrix}3-\lambda&5&8\\1&-1-\lambda&-2\\
-1&-1&-1-\lambda\end{vmatrix}}=-(\lambda-1)(\lambda+2)(\lambda-2)$.
The eigenvectors associated with
 with $\lambda_1=1$ satisfy the system with  augmented matrix
$\dst{\begin{bmatrix}2&5&8&\vdots&0\\1&-2
&-2&\vdots&0\\-1&-1&-2&\vdots&0
\end{bmatrix}}$,
which is row equivalent to
$\dst{\begin{bmatrix}1&0&\frac{2}{3}&\vdots&0\\[1\jot]0&1&\frac{4}{3}&
\vdots&0\\0&0&0&\vdots&0\end{bmatrix}}$.
Hence $x_1=-\frac{2}{3}x_3$ and $x_2=-\frac{4}{3}x_3$.  Taking $x_3=3$
yields
${\vect y}_1=\dst{\begin{bmatrix}-2\\-4\\3
\end{bmatrix}}e^t$.
The eigenvectors associated with
 with $\lambda_2=-2$ satisfy the system with  augmented matrix
$\dst{\begin{bmatrix}5&5&8&\vdots&0\\1&1
&-2&\vdots&0\\-1&-1&1&\vdots&0
\end{bmatrix}}$,
which is row equivalent to
$\dst{\begin{bmatrix}1&1&0&\vdots&0\\0&0&1&
\vdots&0\\0&0&0&\vdots&0\end{bmatrix}}$.
Hence $x_1=-x_2$ and $x_3=0$.  Taking $x_2=1$ yields
${\vect y}_2=\dst{\begin{bmatrix} -1\\1\\0
\end{bmatrix}}e^{-2t}$.
The eigenvectors associated with
 with $\lambda_3=2$ satisfy the system with  augmented matrix
$\dst{\begin{bmatrix}1&5&8&\vdots&0\\1&-3
&-2&\vdots&0\\-1&-1&-3&\vdots&0
\end{bmatrix}}$,
which is row equivalent to
$\dst{\begin{bmatrix}1&0&\frac{7}{4}&\vdots&0\\[1\jot]0&1&\frac{5}{4}&
\vdots&0\\0&0&0&\vdots&0\end{bmatrix}}$.
Hence $x_1=-\frac{7}{4}x_3$ and $x_2=-\frac{5}{4}x_3$.  Taking $x_3=4$
yields
${\vect y}_3=\dst{\begin{bmatrix}-7\\-5\\4
\end{bmatrix}}e^{2t}$. Hence
${\vect y}=\dst{ c_1\begin{bmatrix}-2\\-4\\
3\end{bmatrix}e^t+c_2\begin{bmatrix}-1\\1\\0
\end{bmatrix}e^{-2t}+c_3\begin{bmatrix}-7\\-5\\4
\end{bmatrix}e^{2t}}$.


\item
$\dst{\begin{vmatrix}4-\lambda&-1&-4\\4&-3-\lambda&-2\\
1&-1&-1-\lambda\end{vmatrix}}=-(\lambda-3)(\lambda+2)(\lambda+1)$.
The eigenvectors associated
 with $\lambda_1=3$ satisfy the system with  augmented matrix
$\dst{\begin{bmatrix}1&-1&-4&\vdots&0\\4&-6
&-2&\vdots&0\\1&-1&-4&\vdots&0
\end{bmatrix}}$,
which is row equivalent to
$\dst{\begin{bmatrix}1&0&-11&\vdots&0\\0&1&-7&
\vdots&0\\0&0&0&\vdots&0\end{bmatrix}}$.
Hence  $x_1=11x_3$ and $x_2=7x_3$.  Taking $x_3=1$ yields
${\vect y}_1=\dst{\begin{bmatrix}11\\7\\1
\end{bmatrix}}e^{3t}$.
The eigenvectors associated
 with $\lambda_2=-2$ satisfy the system with  augmented matrix
$\dst{\begin{bmatrix}6&-1&-4&\vdots&0\\4&-1
&-2&\vdots&0\\1&-1&1&\vdots&0
\end{bmatrix}}$,
which is row equivalent to
$\dst{\begin{bmatrix}1&0&-1&\vdots&0\\0&1&-2&
\vdots&0\\0&0&0&\vdots&0\end{bmatrix}}$.
Hence $x_1=x_3$ and $x_2=2x_3$.  Taking $x_3=1$ yields
${\vect y}_2=\dst{\begin{bmatrix}1\\2\\1
\end{bmatrix}}e^{-2t}$.
The eigenvectors associated
 with $\lambda_3=-1$ satisfy the system with  augmented matrix
$\dst{\begin{bmatrix}5&-1&-4&\vdots&0\\4&-2
&-2&\vdots&0\\1&-1&0&\vdots&0
\end{bmatrix}}$,
which is row equivalent to
$\dst{\begin{bmatrix}1&0&-1&\vdots&0\\0&1&-1&
\vdots&0\\0&0&0&\vdots&0\end{bmatrix}}$.
Hence  $x_1=x_3$ and $x_2=x_3$.  Taking $x_3=1$ yields
${\vect y}_3=\dst{\begin{bmatrix}1\\1\\1
\end{bmatrix}}e^{-t}$. Hence
 ${\vect y}=\dst{c_1\begin{bmatrix}11\\\phantom{1}7
\\\phantom{1}1\end{bmatrix}e^{3t}+c_2\begin{bmatrix}1\\2
\\1\end{bmatrix}e^{-2t}+c_3\begin{bmatrix}1\\1\\1
\end{bmatrix}e^{-t}}$.


\item
$\dst{\begin{vmatrix}3-\lambda&2&-2\\-2&7-\lambda&-2\\-10
&10&-5-\lambda\end{vmatrix}}=-(\lambda+5)(\lambda-5)^2$.
The eigenvectors associated
 with $\lambda_1=-5$ satisfy the system with  augmented matrix
$\dst{\begin{bmatrix}8&2&-2&\vdots&0\\-2&12
&-2&\vdots&0\\-10&10&0&\vdots&0
\end{bmatrix}}$,
which is row equivalent to
$\dst{\begin{bmatrix}1&0&-\frac{1}{5}&\vdots&0\\[1\jot]0&1&-\frac{1}{5}&
\vdots&0\\0&0&0&\vdots&0\end{bmatrix}}$.
Hence  $x_1=\dst\frac{1}{5}x_3$ and $x_2=\dst\frac{1}{5}x_3$.  Taking
$x_3=5$ yields
${\vect y}_1=\dst{\begin{bmatrix}1\\1\\5
\end{bmatrix}}e^{-5t}$.
The eigenvectors associated
 with $\lambda_2=5$ satisfy the system with  augmented matrix
$\dst{\begin{bmatrix}-2&2&-2&\vdots&0\\-2&2
&-2&\vdots&0\\-10&10&-10&\vdots&0
\end{bmatrix}}$,
which is row equivalent to
$\dst{\begin{bmatrix}1&-1&1&\vdots&0\\0&0&0&
\vdots&0\\0&0&0&\vdots&0\end{bmatrix}}$.
Hence $x_1=x_2-x_3$.  Taking $x_2=0$ and $x_3=1$ yields
${\vect y}_2=\dst{\begin{bmatrix}-1\\0\\1
\end{bmatrix}}e^{5t}$.
 Taking $x_2=1$ and $x_3=0$ yields
${\vect y}_3=\dst{\begin{bmatrix}1\\1\\0
\end{bmatrix}}e^{5t}$. Hence
 ${\vect y}=\dst{c_1\begin{bmatrix}1\\1\\5
\end{bmatrix}e^{-5t}+c_2\begin{bmatrix}-1\\0\\
 1\end{bmatrix}e^{5t}+c_3\begin{bmatrix}1\\1\\0\end{bmatrix}e^{5t}}$.


\item
$\dst{\begin{vmatrix}-7-\lambda&4\\-6&7-\lambda\end{vmatrix}}
=(\lambda-5)(\lambda+5)$.
Eigenvectors  associated with $\lambda_1=5$  satisfy
$\dst{\begin{bmatrix}-12&4\\-6&2
\end{bmatrix}\begin{bmatrix}
x_1\\x_2\end{bmatrix}=\begin{bmatrix} 0\\0\end{bmatrix}}$,
so $x_1=\dst\frac{x_2}{3}$.  Taking $x_2=3$ yields
${\vect y}_1=\dst{\begin{bmatrix} 1\\3\end{bmatrix}}e^{5t}$.
Eigenvectors  associated with $\lambda_2=5$
satisfy
$\dst{\begin{bmatrix}-2&4\\-6&12
\end{bmatrix}\begin{bmatrix}
x_1\\x_2\end{bmatrix}=\begin{bmatrix} 0\\0\end{bmatrix}}$,
so $x_1=2x_2$.  Taking $x_2=1$ yields
${\vect y}_2=\dst{\begin{bmatrix}2\\1\end{bmatrix}}e^{-5t}$.
The general solution is
${\vect y}=c_1\dst{\begin{bmatrix}
1\\3\end{bmatrix}}e^{5t}+
c_2\dst{\begin{bmatrix}2\\1\end{bmatrix}}e^{-5t}$.
Now ${\vect y}(0)=\dst{\twocol2{-4}}\Rightarrow
c_1\dst{\begin{bmatrix}
1\\3\end{bmatrix}}+
c_2\dst{\begin{bmatrix}2\\1\end{bmatrix}}=\dst{\twocol2{-4}}$,
so $c_1=-2$ and $c_2=2$. Therefore,
${\vect y}=\dst{-\twocol26e^{5t}+\twocol42e^{-5t}}$.


\item
$\dst{\begin{vmatrix}21-\lambda&-12\\24&-15-\lambda\end{vmatrix}}
=(\lambda-9)(\lambda+3)$.
Eigenvectors  associated with $\lambda_1=9$  satisfy
$\dst{\begin{bmatrix}12&-12\\24&-24
\end{bmatrix}\begin{bmatrix}
x_1\\x_2\end{bmatrix}=\begin{bmatrix} 0\\0\end{bmatrix}}$,
so $x_1=x_2$.  Taking $x_2=1$ yields
${\vect y}_1=\dst{\begin{bmatrix}1\\1\end{bmatrix}}e^{9t}$.
Eigenvectors  associated with $\lambda_2=-3$
$\dst{\begin{bmatrix}24&-12\\24&-12
\end{bmatrix}\begin{bmatrix}
x_1\\x_2\end{bmatrix}=\begin{bmatrix} 0\\0\end{bmatrix}}$,
so $x_1=\dst\frac{1}{2}x_2$.  Taking $x_2=2$ yields
${\vect y}_2=\dst{\begin{bmatrix}1\\2\end{bmatrix}}e^{-3t}$.
The general solution is
${\vect y}=c_1\dst{\twocol11}e^{9t}+c_2\dst{\twocol12}e^{-3t}$.
Now  ${\vect y}(0)=\dst{\twocol53}\Rightarrow
c_1\dst{\twocol11}+c_2\dst{\twocol12}=\dst{\twocol53}$, so $c_1=7$
and $c_2=-2$. Therefore,
${\vect y}=\dst{\twocol77e^{9t}-\twocol24e^{-3t}}$.


\item
$\dst\begin{vmatrix}
\frac{1}{6}-\lambda&\frac{1}{3}&0\\
\frac{2}{3}&-\frac{1}{6}-\lambda&0\\
0&0&\frac{1}{2}-\lambda
\end{vmatrix}=-(\lambda+1/2)(\lambda-1/2)^2$.
The eigenvectors associated with
 with $\lambda_1=-1/2$ satisfy the system with  augmented matrix
$\dst{\begin{bmatrix}\frac{2}{3}&\frac{1}{3}&0&\vdots&0
\\[1\jot]\frac{2}{3}&\frac{1}{3}&0&\vdots&0\\0&0&1&\vdots&0
\end{bmatrix}}$,
which is row equivalent to
$\dst{\begin{bmatrix}1&\frac{1}{2}&0&\vdots&0\\[1\jot]0&0&1&
\vdots&0\\0&0&0&\vdots&0\end{bmatrix}}$.
Hence $x_1=-\dst\frac{x_2}{2}$ and $x_3=0$.  Taking $x_2=2$ yields
${\vect y}_1=\dst{\begin{bmatrix}-1\\2\\0
\end{bmatrix}}e^{-t/2}$.
The eigenvectors associated with
 with $\lambda_2=\lambda_3=1/2$ satisfy the system with  augmented
matrix
$\dst{\begin{bmatrix}-\frac{1}{3}&\frac{1}{3}&0&\vdots&0\\[1\jot]
\frac{2}{3}&-\frac{2}{3}&0&\vdots&0\\0&0&0&\vdots&0
\end{bmatrix}}$,
which is row equivalent to
$\dst{\begin{bmatrix}1&-1&0&\vdots&0\\0&0&0&
\vdots&0\\0&0&0&\vdots&0\end{bmatrix}}$.
Hence $x_1=x_2$ and $x_3$ is arbitrary.  Taking $x_2=1$ and $x_3=0$
yields
${\vect y}_2=\dst{\begin{bmatrix}1\\1\\0
\end{bmatrix}}e^{t/2}$.
Taking $x_2=0$ and $x_3=1$ yields
${\vect y}_3=\dst{\begin{bmatrix}0\\0\\1
\end{bmatrix}}e^{t/2}$.
 The general solution is
${\vect y}=\dst{c_1\threecol{-1}20e^{-t/2}+
c_2\threecol110e^{t/2}+c_3\threecol001e^{t/2}}$.
Now ${\vect y}(0)=\dst{\threecol471}\Rightarrow
\dst{c_1\threecol{-1}20+
c_2\threecol110+c_3\threecol001e^{t/2}}=\dst{\threecol471}$,
so $c_1=1$, $c_2=5$, and $c_3=1$. Hence
${\vect y}=\dst{\threecol{-1}20e^{-t/2}+
\threecol550e^{t/2}+\threecol001e^{t/2}}$.


\item
$\dst{\begin{vmatrix}6-\lambda&-3&-8\\2&1-\lambda&-2\\
3&-3&-5-\lambda\end{vmatrix}}=-(\lambda-1)(\lambda+2)(\lambda-3)$.
The eigenvectors associated
 with $\lambda_1=1$ satisfy the system with  augmented matrix
$\dst{\begin{bmatrix}5&-3&-8&\vdots&0\\2&0
&-2&\vdots&0\\3&-3&-6&\vdots&0
\end{bmatrix}}$,
which is row equivalent to
$\dst{\begin{bmatrix}1&0&-1&\vdots&0\\0&1&1&
\vdots&0\\0&0&0&\vdots&0\end{bmatrix}}$.
Hence  $x_1=x_3$ and $x_2=-x_3$.  Taking $x_3=$ yields
${\vect y}_1=\dst{\begin{bmatrix}1\\-1\\1
\end{bmatrix}}e^t$.
The eigenvectors associated
 with $\lambda_2=-2$ satisfy the system with  augmented matrix
$\dst{\begin{bmatrix}8&-3&-8&\vdots&0\\2&3
&-2&\vdots&0\\3&-3&-3&\vdots&0
\end{bmatrix}}$,
which is row equivalent to
$\dst{\begin{bmatrix}1&0&-1&\vdots&0\\0&1&0&
\vdots&0\\0&0&0&\vdots&0\end{bmatrix}}$.
Hence $x_1=x_3$ and $x_2=0$.  Taking $x_3=1$ yields
${\vect y}_2=\dst{\begin{bmatrix}1\\0\\1
\end{bmatrix}}e^{-2t}$.
The eigenvectors associated
 with $\lambda_3=3$ satisfy the system with  augmented matrix
$\dst{\begin{bmatrix}3&-3&-8&\vdots&0\\2&-2
&-2&\vdots&0\\3&-3&-8&\vdots&0
\end{bmatrix}}$,
which is row equivalent to
$\dst{\begin{bmatrix}1&-1&0&\vdots&0\\0&0&1&
\vdots&0\\0&0&0&\vdots&0\end{bmatrix}}$.
Hence  $x_1=x_2$ and $x_3=0$.  Taking $x_2=1$ yields
${\vect y}_3=\dst{\begin{bmatrix}1\\1\\0
\end{bmatrix}}e^{3t}$.
The general solution is
${\vect y}=c_1\dst{\begin{bmatrix}1\\-1\\1
\end{bmatrix}}e^t+
c_2\dst{\begin{bmatrix}1\\0\\1
\end{bmatrix}}e^{-2t}
+c_3\dst{\begin{bmatrix}1\\1\\0
\end{bmatrix}}e^{3t}$.
Now ${\vect y}(0)=\dst{\threecol0{-1}{-1}}\Rightarrow
c_1\dst{\begin{bmatrix}1\\-1\\1
\end{bmatrix}}+
c_2\dst{\begin{bmatrix}1\\0\\1
\end{bmatrix}}
+c_3\dst{\begin{bmatrix}1\\1\\0
\end{bmatrix}}=\dst{\threecol0{-1}{-1}}$, so
$c_1=2$, $c_2=-3$, and $c_3=1$. Therefore,
${\vect
y}=\dst{\threecol2{-2}2e^t-\threecol303e^{-2t}+\threecol110e^{3t}}$.


\item
$\dst{\begin{vmatrix}3-\lambda&0&1\\11&-2-\lambda&7\\
1&0&3-\lambda\end{vmatrix}}=-(\lambda-2)(\lambda+2)(\lambda-4)$.
The eigenvectors associated with
 with $\lambda_1=2$ satisfy the system with  augmented matrix
$\dst{\begin{bmatrix}1&0&1&\vdots&0\\11&-4
&7&\vdots&0\\1&0&1&\vdots&0
\end{bmatrix}}$,
which is row equivalent to
$\dst{\begin{bmatrix}1&0&1&\vdots&0\\0&1&1&
\vdots&0\\0&0&0&\vdots&0\end{bmatrix}}$.
Hence $x_1=-x_3$ and $x_2=-x_3$.  Taking $x_3=1$ yields
${\vect y}_1=\dst{\begin{bmatrix}-1\\-1\\1
\end{bmatrix}}e^{2t}$.
The eigenvectors associated with
 with $\lambda_2=-2$ satisfy the system with  augmented matrix
$\dst{\begin{bmatrix}5&0&1&\vdots&0\\11&0
&7&\vdots&0\\1&0&5&\vdots&0
\end{bmatrix}}$,
which is row equivalent to
$\dst{\begin{bmatrix}1&0&0&\vdots&0\\0&0&1&
\vdots&0\\0&0&0&\vdots&0\end{bmatrix}}$.
Hence $x_1=x_3=0$ and $x_2$ is arbitrary.  Taking $x_3=1$ yields
${\vect y}_2=\dst{\begin{bmatrix}0\\1\\0
\end{bmatrix}}e^{-2t}$.
The eigenvectors associated with
 with $\lambda_3=4$ satisfy the system with  augmented matrix
$\dst{\begin{bmatrix}-1&0&1&\vdots&0\\11&-6
&7&\vdots&0\\1&0&-1&\vdots&0
\end{bmatrix}}$,
which is row equivalent to
$\dst{\begin{bmatrix}1&0&-1&\vdots&0\\0&1&-3&
\vdots&0\\0&0&0&\vdots&0\end{bmatrix}}$.
Hence $x_1=x_3$ and $x_2=3x_3$.  Taking $x_3=1$ yields
${\vect y}_3=\dst{\begin{bmatrix}1\\3\\1
\end{bmatrix}}e^{4t}$.
 The general solution is
${\vect y}=c_1\dst{\begin{bmatrix}-1\\-1\\1
\end{bmatrix}}e^{2t}
+c_2\dst{\begin{bmatrix}0\\1\\0
\end{bmatrix}}e^{-2t}
+c_3\dst{\begin{bmatrix}1\\3\\1
\end{bmatrix}}e^{4t}$.
Now ${\vect y}(0)=\dst{\threecol276}\Rightarrow
c_1\dst{\begin{bmatrix}-1\\-1\\1
\end{bmatrix}}
+c_2\dst{\begin{bmatrix}0\\1\\0
\end{bmatrix}}
+c_3\dst{\begin{bmatrix}1\\3\\1
\end{bmatrix}}=\dst{\threecol276}$, so $c_1=2$, $c_2=-3$, and
$c_3=4$. Hence
${\vect y}=\dst{\threecol{-2}{-2}2e^{2t}-\threecol030e^{-2t}+
\threecol4{12}4e^{4t}}$


\item
$\begin{vmatrix}3-\lambda&-1&0\\4&-2-\lambda&0
\\4&-4&2-\lambda\end{vmatrix}
=-(\lambda+1)(\lambda-2)^2$.
The eigenvectors associated
 with $\lambda_1=-1$ satisfy the system with  augmented matrix
$\begin{bmatrix}4&-1&0&\vdots&0\\4&-1&0&\vdots&0\\4&-4&3&\vdots&0\end{bmatrix}$,
which is row equivalent to
$\begin{bmatrix}1&0&-\frac{1}{4}&\vdots&0\\0&1&-1&\vdots&0\\0&0&0&\vdots&0\end{bmatrix}$
Hence $x_1=x_2/4$ and $x_2=x_3$.  Taking $x_3=4$ yields
${\vect y}_1=\dst{\begin{bmatrix}1\\4\\4
\end{bmatrix}}e^{-t}$.
The eigenvectors associated with
 with $\lambda_2=\lambda_3=2$ satisfy the system with  augmented
matrix
$\begin{bmatrix}1&-1&0&\vdots&0\\4&-4&0&\vdots&0\\4&-4&0&\vdots&0\end{bmatrix}$,
which is row equivalent to
$\begin{bmatrix}1&-1&0&\vdots&0\\0&0&0&\vdots&0\\0&0&0&\vdots&0\end{bmatrix}$.
Hence $x_1=x_2$ and $x_3$ is arbitrary.  Taking $x_2=1$ and $x_3=0$
yields
${\vect y}_2=\dst{\begin{bmatrix}1\\1\\0
\end{bmatrix}}e^{2t}$.
Taking $x_2=0$ and $x_3=1$ yields
${\vect y}_3=\dst{\begin{bmatrix}0\\0\\1
\end{bmatrix}}e^{2t}$.
 The general solution is
${\vect y}=c_1\dst{\begin{bmatrix}1\\4\\4
\end{bmatrix}}e^{-t}+
c_2\dst{\begin{bmatrix}1\\1\\0
\end{bmatrix}}e^{2t}+
c_3\dst{\begin{bmatrix}0\\0\\1
\end{bmatrix}}e^{2t}$.
Now ${\vect y}(0)=\dst{\threecol7{10}{2}}\Rightarrow
c_1\dst{\begin{bmatrix}1\\4\\4
\end{bmatrix}}+
c_2\dst{\begin{bmatrix}1\\1\\0
\end{bmatrix}}+
c_3\dst{\begin{bmatrix}0\\0\\1
\end{bmatrix}}=\dst{\threecol7{10}{2}}$,
 so $c_1=1$, $c_2=6$, and
$c_3=-2$. Hence
 ${\vect y}=\dst
\threecol144e^{-t}+\threecol66{-2}e^{2t}$.


\item
\begin{alist}
\item If ${\vect y}(t_0)=0$, then ${\vect y}$ is the solution
of the initial value problem ${\vect y}'=A{\vect y},\ {\vect y}(t_0)={\vect
0}$. Since ${\vect y}\equiv0$ is a solution of this problem,
Theorem~10.2.1 implies the conclusion.

\item It is given that ${\vect y}_1'(t)=A{\vect y}_1(t)$ for all $t$.
Replacing $t$ by $t-\tau$ shows that ${\vect y}_1'(t-\tau)=A{\vect
y}_1(t-\tau)=A{\vect y}_2(t)$ for all $t$. Since ${\vect y}_2'(t)={\vect
y}_1'(t-\tau)$ by the chain rule, this implies that ${\vect
y}_2'(t)=A{\vect y}_2(t)$ for all $t$.

\item If ${\vect z}(t)={\vect y}_1(t-\tau)$, then ${\vect z}(t_2)={\vect
y}_1(t_1)={\vect y}_2(t_2)$; therefore ${\vect z}$ and ${\vect y}_2$ are
both solutions of the initial value problem ${\vect y}'=A{\vect y},\ {\vect
y}(t_2)={\vect k}$, where ${\vect k}={\vect y}_2(t_2)$.
\end{alist}

\addtocounter{enumi}{12}
\item%10.4.42
The characteristic polynomial of $A$ is
$p(\lambda)=\lambda^2-(a+b)+ab-\alpha\beta$, so the eigenvalues of $A$
are $\lambda_1=\dst\frac{a+b-\gamma}{2}$ and
$\lambda_1=\dst\frac{a+b+\gamma}{2}$, where
$\gamma=\sqrt{(a-b)^2+4\alpha\beta}$; ${\vect
x}_1=\ctwocol{b-a+\gamma}{2\beta}$ and ${\vect
x}_2=\ctwocol{b-a-\gamma}{2\beta}$ are associated eigenvectors. Since
$\gamma>|b-a|$, if $L_1$ and $L_2$ are lines through the origin
parallel to ${\vect x}_1$ and ${\vect x}_2$, then $L_1$ is in the first and
third quadrants and $L_2$ is in the second and fourth quadrants. The
slope of $L_1$ is $\rho=\dst\frac{2\beta}{ b-a+\gamma}>0$. If $Q_0=\rho
P_0$ there are three possibilities: (i) if $\alpha\beta=ab$, then
$\lambda_1=0$ and $P(t)=P_0$, $Q(t)=Q_0$ for all $t>0$; (ii) if
$\alpha\beta<ab$, then $\lambda_1>0$ and
$\lim_{t\to\infty}P(t)=\lim_{t\to\infty}Q(t)=\infty$ (monotonically);
(iii) if $\alpha\beta>ab$, then $\lambda_1<0$ and
$\lim_{t\to\infty}P(t)=\lim_{t\to\infty}Q(t)=0$ (monotonically). Now
suppose $Q_0\ne\rho P_0$, so that the trajectory cannot intersect
$L_1$, and assume for the moment that (A) makes sense for all $t>0$;
that is, even if one or the other of $P$ and $Q$ is negative. Since
$\lambda_2>0$ it follows that either $\lim_{t\to\infty}P(t)=\infty$ or
$\lim_{t\to\infty}Q(t)=\infty$ (or both), and the trajectory is
asymptotically parallel to $L_2$. Therefore,the trajectory must cross
into the third quadrant (so $P(T)=0$ and $Q(T)>0$ for some finite $T$)
if $Q_0>\rho P_0$, or into the fourth quadrant (so $Q(T)=0$ and
$P(T)>0$ for some finite $T$) if $Q_0<\rho P_0$.
\end{sectionSolutions}

\section{Constant Coefficient Homogeneous Systems II}

\begin{sectionSolutions}
\item
$\dst{\twochar{}{-1}1{-2}}=(\lambda+1)^2$.
Hence $\lambda_1=-1$.
Eigenvectors   satisfy
$\dst{\twobytwo1{-1}1{-1}\twocol{x_1}{x_2}=\twocol00}$,
so $x_1=x_2$.  Taking $x_2=1$ yields
${\vect y}_1=\dst{\twocol11}e^{-t}$.
For a second solution we need a vector ${\vect u}$
such that $\dst{\twobytwo1{-1}1{-1}\twocol{u_1}{u_2}
=\twocol11}$. Let $u_1=1$ and $u_2=0$. Then
${\vect y}_2=\dst{\twocol10}e^{-t}+\dst{\twocol11}te^{-t}$.
The general solution is
${\vect y}=\dst{
c_1\twocol11e^{-t}+c_2\left(\twocol10e^{-t}+\twocol11te^{-t}\right)}$.

\item
 $\dst{{\vect y}'=\twochar31{-1}1}=(\lambda-2)^2$.
Hence $\lambda_1=2$.
Eigenvectors   satisfy
$\dst{\twobytwo11{-1}{-1}\twocol{x_1}{x_2}=\twocol00}$,
so $x_1=-x_2$.  Taking $x_2=1$ yields
${\vect y}_1=\dst{\twocol{-1}1}e^{2t}$.
For a second solution we need a vector ${\vect u}$
such that $\dst{\twobytwo11{-1}{-1}\twocol{u_1}{u_2}
=\twocol{-1}1}$. Let $u_1=-1$ and $u_2=0$. Then
${\vect y}_2=\dst{\twocol{-1}0}e^{2t}+\dst{\twocol{-1}1}te^{2t}$.
The general solution is
 ${\vect y}=\dst{
c_1\twocol{-1}1e^{2t}+c_2\left(\twocol{-1}0
e^{2t}+\twocol{-1}1te^{2t}\right)}$.


\item
$\dst{\twochar{-10}9{-4}2}=(\lambda+4)^2$.
Hence $\lambda_1=-4$.
Eigenvectors   satisfy
$\dst{\twobytwo{-6}9{-4}6\twocol{x_1}{x_2}=\twocol00}$,
so $x_1=\frac{3}{2}x_2$.  Taking $x_2=2$ yields
${\vect y}_1=\dst{\twocol32}e^{-4t}$.
For a second solution we need a vector ${\vect u}$
such that $\dst{\twobytwo{-6}9{-4}6\twocol{u_1}{u_2}
=\twocol32}$. Let $u_1=-\frac{1}{2}$ and $u_2=0$. Then
${\vect y}_2=\dst{\twocol{-1}0}\frac{e^{-4t}}{2}+\dst{\twocol32}te^{-4t}$.
The general solution is
${\vect y}=\dst{c_1\twocol32e^{-4t}+
c_2\left(\twocol{-1}0\frac{e^{-4t}}{2}+\twocol32te^{-4t}\right)}$.


\item
$\dst{\threechar{}21{-4}61042}
=-\lambda(\lambda-4)^2$.
Hence $\lambda_1=0$ and  $\lambda_2=\lambda_3=4$.
The eigenvectors associated
 with $\lambda_1=0$ satisfy the system with  augmented matrix
$\dst{\begin{bmatrix}0&2&1&\vdots&0\\{-4}&6&1&
\vdots&0\\0&4&2&\vdots&0\end{bmatrix}}$,
which is row equivalent to
$\dst{\begin{bmatrix}1&0&\frac{1}{2}&\vdots&0\\0&1&\frac{1}{2}&
\vdots&0\\0&0&0&\vdots&0\end{bmatrix}}$.
Hence  $x_1=-\dst\frac{1}{2}x_3$ and $x_2=-\dst\frac{1}{2}x_3$.  Taking $x_3=2$
yields
${\vect y}_1=\dst{\threecol{-1}{-1}2}$.
The eigenvectors associated
 with $\lambda_2=4$ satisfy the system with  augmented matrix
$\dst{\begin{bmatrix}-4&2&1&\vdots&0\\-4&2&1&
\vdots&0\\0&4&-2&\vdots&0\end{bmatrix}}$,
which is row equivalent to
$\dst{\begin{bmatrix}1&0&-\frac{1}{2}&\vdots&0\\0&1&-\frac{1}{2}&
\vdots&0\\0&0&0&\vdots&0\end{bmatrix}}$.
Hence  $x_1=\dst\frac{1}{2}x_3$ and $x_2=\dst\frac{1}{2}x_3$.  Taking
$x_3=2$ yields
${\vect y}_2=\dst{\threecol112}e^{4t}$.
For a third solution we need a vector ${\vect u}$ such that
$\dst{\threebythree{-4}21{-4}2104{-2}\threecol{u_1}{u_2}{u_3}}=\dst{\threecol112}$.
The augmented matrix of this system is row equivalent to
$\dst{\begin{bmatrix}1&0&-\frac{1}{2}&\vdots&0\\0&1&-\frac{1}{2}&
\vdots&\frac{1}{2}\\0&0&0&\vdots&0\end{bmatrix}}$.
Let  $u_3=0$, $u_1=0$, and $u_2=\dst\frac{1}{2}$. Then
${\vect
y}_3=\dst{\threecol010}\dst\frac{e^{4t}}{2}+\dst{\threecol112}te^{4t}$.
The general solution is
\[
{\vect y}=\dst{
c_1\threecol{-1}{-1}2+c_2\threecol112e^{4t}+
c_3\left(\threecol010\frac{e^{4t}}{2}+\threecol112te^{4t}\right)}.
\]


\item
$\dst{\threechar{-1}1{-1}{-2}{}2{-1}3{-1}}
=-(\lambda-2)(\lambda+2)^2$.
Hence $\lambda_1=2$ and  $\lambda_2=\lambda_3=-2$.
The eigenvectors associated
 with $\lambda_1=2$ satisfy the system with  augmented matrix
$\dst{\begin{bmatrix}-3&1&-1&\vdots&0\\-2&-2&2&
\vdots&0\\-1&3&-3&\vdots&0\end{bmatrix}}$,
which is row equivalent to
$\dst{\begin{bmatrix}1&0&0&\vdots&0\\0&1&-1&
\vdots&0\\0&0&0&\vdots&0\end{bmatrix}}$.
Hence  $x_1=0$ and $x_2=x_3$.  Taking $x_3=1$
yields
${\vect y}_1=\dst{\threecol011}e^{2t}$.
The eigenvectors associated
 with $\lambda_2=-2$ satisfy the system with  augmented matrix
$\dst{\begin{bmatrix}1&1&-1&\vdots&0\\-2&2&2&
\vdots&0\\-1&3&1&\vdots&0\end{bmatrix}}$,
which is row equivalent to
$\dst{\begin{bmatrix}1&0&-1&\vdots&0\\0&1&0&
\vdots&0\\0&0&0&\vdots&0\end{bmatrix}}$.
Hence  $x_1=x_3$ and $x_2=0$.  Taking $x_3=1$
yields
${\vect y}_2=\dst{\threecol101}e^{-2t}$.
For a third solution we need a vector ${\vect u}$ such that
$\dst{\threebythree11{-1}{-2}22{-1}31\threecol{u_1}{u_2}{u_3}}=\dst{\threecol101}$.
The augmented matrix of this system is row equivalent to
$\dst{\begin{bmatrix}1&0&-1&\vdots&\frac{1}{2}\\0&1&0&
\vdots&\frac{1}{2}\\0&0&0&\vdots&0\end{bmatrix}}$.
Let $u_3=0$, $u_1=\dst\frac{1}{2}$, and $u_2=\dst\frac{1}{2}$. Then ${\vect
y}_3=\dst{\threecol110}\dst\frac{e^{-2t}}{2}+\dst{\threecol101}te^{-2t}$.
The general solution is
\[
{\vect y}=\dst{
c_1\threecol011e^{2t}+c_2\threecol101e^{-2t}+c_3\left
(\threecol110\frac{e^{-2t}}{2}+\threecol101te^{-2t}\right)}.
\]


\item
$\dst{\threechar6{-5}32{-1}3211}
=-(\lambda+2)(\lambda-4)^2$.
Hence $\lambda_1=-2$ and  $\lambda_2=\lambda_3=4$.
The eigenvectors associated
 with $\lambda_1=-2$ satisfy the system with  augmented matrix
$\dst{\begin{bmatrix}8&-5&3&\vdots&0\\2&1&3&
\vdots&0\\2&1&3&\vdots&0\end{bmatrix}}$,
which is row equivalent to
$\dst{\begin{bmatrix}1&0&1&\vdots&0\\0&1&1&
\vdots&0\\0&0&0&\vdots&0\end{bmatrix}}$.
Hence  $x_1=-x_3$ and $x_2=-x_3$.  Taking $x_3=1$
yields
${\vect y}_1=\dst{\threecol{-1}{-1}1}e^{-2t}$.
The eigenvectors associated
 with $\lambda_2=4$ satisfy the system with  augmented matrix
$\dst{\begin{bmatrix}2&-5&3&\vdots&0\\2&-5&3&
\vdots&0\\2&1&-3&\vdots&0\end{bmatrix}}$,
which is row equivalent to
$\dst{\begin{bmatrix}1&0&-1&\vdots&0\\0&1&-1&
\vdots&0\\0&0&0&\vdots&0\end{bmatrix}}$.
Hence  $x_1=x_3$ and $x_2=x_3$.  Taking $x_3=1$
yields
${\vect y}_2=\dst{\threecol111}e^{4t}$.
For a third solution we need a vector ${\vect u}$ such that
$\dst{\threebythree2{-5}32{-5}321{-3}\threecol{u_1}{u_2}{u_3}}=
\dst{\threecol111}$.
The augmented matrix of this system is row equivalent to
$\dst{\begin{bmatrix}1&0&-1&\vdots&\frac{1}{2}\\0&1&-1&
\vdots&0\\0&0&0&\vdots&0\end{bmatrix}}$.
Let $u_3=0$, $u_1=\dst\frac{1}{2}$, and $u_2=0$. Then ${\vect
y}_3=\dst{\threecol100}\frac{e^{4t}}{2}+\dst{\threecol111}te^{4t}$. The
general solution is
 ${\vect y}=\dst{
c_1\threecol{-1}{-1}1e^{-2t}+c_2\threecol111e^{4t}+c_3\left
(\threecol100\frac{e^{4t}}{2}+\threecol111te^{4t}\right)}$.


\item
$\dst{\twochar{15}{-9}{16}{-9}}=(\lambda-3)^2$.
Hence $\lambda_1=3$.
Eigenvectors   satisfy
$\dst{\twobytwo{12}{-9}{16}{-12}\twocol{x_1}{x_2}=\twocol00}$,
so $x_1=\dst\frac{3}{4}x_2$.  Taking $x_2=4$ yields
${\vect y}_1=\dst{\twocol34}e^{3t}$.
For a second solution we need a vector ${\vect u}$
such that $\dst{\twobytwo{12}{-9}{16}{-12}\twocol{u_1}{u_2}
=\twocol34}$. Let $u_1=\dst\frac{1}{4}$ and $u_2=0$. Then
${\vect y}_2=\dst{\twocol10}\frac{e^{3t}}{4}+\dst{\twocol34}te^{3t}$.
The general solution is
${\vect y}=c_1\dst{\twocol34}e^{3t}+
c_2\left(\dst{\twocol10}\frac{e^{3t}}{4}+\dst{\twocol34}te^{3t}\right)$.
Now ${\vect y}(0)=\dst{\twocol58}\Rightarrow
c_1\dst{\twocol34}+
c_2\dst{\twocol{\frac{1}{4}}0}=\dst{\twocol58}$,  so $c_1=2$
and $c_2=-4$. Therefore,
${\vect y}=\dst{\twocol58e^{3t}-\twocol{12}{16}te^{3t}}$.


\item
$\dst{\twochar{-7}{24}{-6}{17}}=(\lambda-5)^2$.
Hence $\lambda_1=5$.
Eigenvectors   satisfy
$\dst{\twobytwo{-12}{24}{-6}{12}\twocol{x_1}{x_2}=\twocol00}$,
so $x_1=2x_2$.  Taking $x_2=1$ yields
${\vect y}_1=\dst{\twocol21}e^{5t}$.
For a second solution we need a vector ${\vect u}$
such that $\dst{\twobytwo{-12}{24}{-6}{12}\twocol{u_1}{u_2}}
=\dst{\twocol21}$. Let $u_1=\dst\frac{1}{6}$ and $u_2=0$. Then
${\vect y}_2=\dst{\twocol10}\frac{e^{5t}}{6}+\dst{\twocol21}te^{5t}$.
The general solution is
${\vect y}=c_1\dst{\twocol21}e^{5t}
+c_2\left(\dst{\twocol10}\frac{e^{5t}}{6}+\dst{\twocol21}te^{5t}\right)$.
Now
${\vect y}(0)=\dst{\twocol31}\Rightarrow c_1\dst{\twocol21}
+c_2{\twocol{\frac{1}{6}}0}=\dst{\twocol31}$, so $c_1=1$
and $c_2=6$. Therefore,
${\vect y}=\dst{\twocol31e^{5t}-\twocol{12}6te^{5t}}$.


\item
$\dst{\threechar{-1}101{-1}{-2}{-1}{-1}{-1}}
=-(\lambda-1)(\lambda+2)^2$.
Hence $\lambda_1=1$ and  $\lambda_2=\lambda_3=-2$.
The eigenvectors associated
 with $\lambda_1=1$ satisfy the system with  augmented matrix
$\dst{\begin{bmatrix}{-2}&1&0&\vdots&0\\1&-2&-2&
\vdots&0\\-1&-1&-2&\vdots&0\end{bmatrix}}$,
which is row equivalent to
$\dst{\begin{bmatrix}1&0&\frac{2}{3}&\vdots&0\\0&1&\frac{4}{3}&
\vdots&0\\0&0&0&\vdots&0\end{bmatrix}}$.
Hence  $x_1=-\dst\frac{2}{3}x_3$ and $x_2=-\dst\frac{4}{3}x_3$.  Taking
$x_3=3$ yields
The eigenvectors associated
 with $\lambda_2=-2$ satisfy the system with  augmented matrix
$\dst{\begin{bmatrix}1&1&0&\vdots&0\\1&1&-2&
\vdots&0\\-1&-1&1&\vdots&0\end{bmatrix}}$,
which is row equivalent to
$\dst{\begin{bmatrix}1&1&0&\vdots&0\\0&0&1&
\vdots&0\\0&0&0&\vdots&0\end{bmatrix}}$.
Hence  $x_1=-x_2$ and $x_3=0$.  Taking $x_2=1$
yields
${\vect y}_2=\dst{\threecol{-1}10}e^{-2t}$.
For a third solution we need a vector ${\vect u}$ such that
$\dst{\threebythree11011{-2}{-1}{-1}1\threecol{u_1}{u_2}{u_3}}=
\dst{\threecol{-1}10}$.
The augmented matrix of this system is row equivalent to
$\dst{\begin{bmatrix}1&1&0&\vdots&-1\\0&0&1&
\vdots&-1\\0&0&0&\vdots&0\end{bmatrix}}$.
Let $u_2=0$, $u_1=-1$, and $u_3=-1$. Then ${\vect
y}_3=\dst{\threecol{-1}0{-1}}e^{-2t}+\dst{\threecol{-1}10}te^{-2t}$.
The general solution is
${\vect y}=c_1\dst{\threecol{-2}{-4}3}e^t+
c_2\dst{\threecol{-1}10}e^{-2t}
+c_3\left(\dst{\threecol{-1}0{-1}}e^{-2t}+
\dst{\threecol{-1}10}te^{-2t}\right)$.
Now ${\vect y}(0)=\dst{\threecol65{-7}}\Rightarrow
c_1\dst{\threecol{-2}{-4}3}+ c_2\dst{\threecol{-1}10}
+c_3\dst{\threecol{-1}0{-1}}=\dst{\threecol65{-7}}$,
so $c_1=-2$, $c_2=-3$, and $c_3=1$.  Therefore,
${\vect y}=\dst{\threecol48{-6}e^t+\threecol2{-3}{-1}e^{-2t}+
\threecol{-1}10te^{-2t}}$.


\item
$\dst{\threechar{-7}{-4}4{-1}01{-9}{-5}6}
=-(\lambda+3)(\lambda-1)^2$.
Hence $\lambda_1=-3$ and  $\lambda_2=\lambda_3=1$.
The eigenvectors associated
 with $\lambda_1=-3$ satisfy the system with  augmented matrix
$\dst{\begin{bmatrix}-4&-4&4&\vdots&0\\-1&3&1&
\vdots&0\\-9&-5&9&\vdots&0\end{bmatrix}}$,
which is row equivalent to
$\dst{\begin{bmatrix}1&0&-1&\vdots&0\\0&1&0&
\vdots&0\\0&0&0&\vdots&0\end{bmatrix}}$.
Hence  $x_1=x_3$ and $x_2=0$.  Taking $x_3=1$
yields
${\vect y}_1=\dst{\threecol101}e^{-3t}$.
The eigenvectors associated
 with $\lambda_2=1$ satisfy the system with  augmented matrix
$\dst{\begin{bmatrix}-8&-4&4&\vdots&0\\-1&-1&1&
\vdots&0\\-9&-5&-5&\vdots&0\end{bmatrix}}$,
which is row equivalent to
$\dst{\begin{bmatrix}1&0&0&\vdots&0\\0&1&-1&
\vdots&0\\0&0&0&\vdots&0\end{bmatrix}}$.
Hence  $x_1=0$ and $x_2=x_3$.  Taking $x_3=1$
yields
${\vect y}_2=\dst{\threecol011}e^t$.
For a third solution we need a vector ${\vect u}$ such that
$\dst{\threebythree{-8}{-4}4{-1}{-1}1{-9}{-5}5\threecol{u_1}{u_2}{u_3}}
=\dst{\threecol011}$.
The augmented matrix of this system is row equivalent to
$\dst{\begin{bmatrix}1&0&0&\vdots&1\\0&1&-1&
\vdots&-2\\0&0&0&\vdots&0\end{bmatrix}}$.
Let $u_3=0$, $u_1=1$, and $u_2=-2$. Then ${\vect
y}_3=\dst{\threecol1{-2}0}e^t+\dst{\threecol011}te^{t}$. The general
solution is
${\vect y}=c_1\dst{\threecol101}e^{-3t}+ c_2\dst{\threecol011}e^t+
c_3\left(\dst{\threecol1{-2}0}e^t+\dst{\threecol011}te^{t}\right)$.
Now ${\vect y}(0)=\dst{\threecol{-6}9{-1}}\Rightarrow
c_1\dst{\threecol101}+ c_2\dst{\threecol011}+
c_3\dst{\threecol1{-2}0}=\dst{\threecol{-6}9{-1}}$,
so $c_1=-2$, $c_2=1$, and $c_3=-4$.  Therefore,
${\vect y}=\dst{-\threecol202e^{-3t}+\threecol{-4}91e^t-
\threecol044te^t}$.


\item
$\dst{\threechar4{-8}{-4}{-3}{-1}{-3}1{-1}9}
=-(\lambda+4)(\lambda-8)^2$.
Hence $\lambda_1=-4$ and  $\lambda_2=\lambda_3=8$.
The eigenvectors associated
 with $\lambda_1=-4$ satisfy the system with  augmented matrix
$\dst{\begin{bmatrix}8&-8&-4&\vdots&0\\-3&3&-3&
\vdots&0\\1&-1&13&\vdots&0\end{bmatrix}}$,
which is row equivalent to
$\dst{\begin{bmatrix}1&-1&0&\vdots&0\\0&0&1&
\vdots&0\\0&0&0&\vdots&0\end{bmatrix}}$.
Hence  $x_1=x_2$ and $x_3=0$.  Taking $x_2=1$
yields
${\vect y}_1=\dst{\threecol110}e^{t}$.
The eigenvectors associated
 with $\lambda_2=8$ satisfy the system with  augmented matrix
$\dst{\begin{bmatrix}-4&-8&-4&\vdots&0\\-3&-9&-3&
\vdots&0\\1&-1&1&\vdots&0\end{bmatrix}}$,
which is row equivalent to
$\dst{\begin{bmatrix}1&0&1&\vdots&0\\0&1&0&
\vdots&0\\0&0&0&\vdots&0\end{bmatrix}}$.
Hence  $x_1=-x_3$ and $x_2=0$.  Taking $x_3=1$
yields
${\vect y}_2=\dst{\threecol{-1}01}e^{8t}$.
For a third solution we need a vector ${\vect u}$ such that
$\dst{\threebythree{-4}{-8}{-4}{-3}{-9}{-3}1{-1}1
\threecol{u_1}{u_2}{u_3}}=\dst{\threecol{-1}01}$.
The augmented matrix of this system is row equivalent to
$\dst{\begin{bmatrix}1&0&1&\vdots&\frac{3}{4}\\0&1&0&
\vdots&-\frac{1}{4}\\0&0&0&\vdots&0\end{bmatrix}}$.
Let $u_3=0$, $u_1=\dst\frac{3}{4}$, and $u_2=-\dst\frac{1}{4}$. Then ${\vect
y}_3=\dst{\threecol3{-1}0}\frac{e^{8t}}{4}+\dst{\threecol{-1}01}te^{8t}$.
The general solution is
$c_1\dst{\threecol110}e^{t}+
c_2\dst{\threecol{-1}01}e^{8t}+
c_3\dst\left({\threecol3{-1}0}\frac{e^{8t}}{4}+
\dst{\threecol{-1}01}te^{8t}\right)$.
Now ${\vect y}(0)=\dst{\threecol{-4}1{-3}}\Rightarrow
c_1\dst{\threecol110}+
c_2\dst{\threecol{-1}01}+
c_3\dst{\threecol{\frac{3}{4}}{-\frac{1}{4}}0}=\dst{\threecol{-4}1{-3}}$,
so $c_1=-1$, $c_2=-3$, and $c_3=-8$.  Therefore,
${\vect y}=\dst{-\threecol110e^{-4t}+\threecol{-3}2{-3}e^{8t}+
\threecol80{-8}te^{8t}}$.


\item
$\dst{\threechar5{-1}1{-1}9{-3}{-2}24}
=-(\lambda-6)^3$.
Hence $\lambda_1=6$. The eigenvectors
 satisfy the system with  augmented matrix
$\dst{\begin{bmatrix}1&-1&1&\vdots&0\\-1&3&-3&
\vdots&0\\-2&2&-2&\vdots&0\end{bmatrix}}$,
which is row equivalent to
$\dst{\begin{bmatrix}1&0&0&\vdots&0\\0&1&-1&
\vdots&0\\0&0&0&\vdots&0\end{bmatrix}}$.
Hence  $x_1=0$ and $x_2=x_3$.  Taking $x_3=1$
yields
${\vect y}_1=\dst{\threecol011}e^{6t}$.
For a second solution we need a vector ${\vect u}$ such that
$\dst{\threebythree1{-1}1{-1}3{-3}{-2}2{-2}\threecol{u_1}{u_2}{u_3}}
=\dst{\threecol011}$.
The augmented matrix of this system is row equivalent to
$\dst{\begin{bmatrix}1&0&0&\vdots&-\frac{1}{4}\\0&1&-1&
\vdots&\frac{1}{4}\\0&0&0&\vdots&0\end{bmatrix}}$.
Let $u_3=0$, $u_1=-\dst\frac{1}{4}$, and $u_2=\dst\frac{1}{4}$. Then ${\vect
y}_2=\dst{\threecol{-1}10}\frac{e^{6t}}{4}+\dst{\threecol011}te^{6t}$.
For a third solution we need a vector ${\vect v}$ such that
$\dst{\threebythree1{-1}1{-1}3{-3}{-2}2{-2}\threecol{v_1}{v_2}{v_3}}
=\dst{\threecolj{-\frac{1}{4}}{\frac{1}{4}}0}$.
The augmented matrix of this system is row equivalent to
$\dst{\begin{bmatrix}1&0&0&\vdots&\frac{1}{8}\\0&1&-1&
\vdots&\frac{1}{8}\\0&0&0&\vdots&0\end{bmatrix}}$.
Let $v_3=0$, $v_1=\dst\frac{1}{8}$, and $v_2=\dst\frac{1}{8}$. Then
${\vect y}_3= \dst{\threecol110}\frac{e^{6t}}{8}+
\dst{\threecol{-1}10}\frac{te^{6t}}{4}+\dst{\threecol011}\frac{t^2e^{6t}}{2}$.
The general solution is
${\vect y}=\dst{c_1\threecol011e^{6t}+
c_2\left(\threecol{-1}10\frac{e^{6t}}{4}+\threecol011te^{6t}\right)}$
\[+c_3\left(\threecol110\frac{e^{6t}}{8}+\threecol{-1}10\frac{t
e^{6t}}{4}+\threecol011\frac{t^2e^{6t}}{2}\right).\]


\item
$\dst{\threechar{-6}{-4}{-4}2{-1}1231}
=-(\lambda+2)^3$.
Hence $\lambda_1=-2$.
The eigenvectors
 satisfy the system with  augmented matrix
$\dst{\begin{bmatrix}-4&-4&-4&\vdots&0\\2&1&1&
\vdots&0\\2&3&3&\vdots&0\end{bmatrix}}$,
which is row equivalent to
$\dst{\begin{bmatrix}1&0&0&\vdots&0\\0&1&1&
\vdots&0\\0&0&0&\vdots&0\end{bmatrix}}$.
Hence  $x_1=0$ and $x_2=-x_3$.  Taking $x_3=1$
yields
${\vect y}_1=\dst{\threecol0{-1}1}e^{-2t}$.
For a second solution we need a vector ${\vect u}$ such that
$\dst{\threebythree{-4}{-4}{-4}211233\threecol{u_1}{u_2}{u_3}}
=\dst{\threecol0{-1}1}$.
The augmented matrix of this system is row equivalent to
$\dst{\begin{bmatrix}1&0&0&\vdots&-1\\0&1&1&
\vdots&1\\0&0&0&\vdots&0\end{bmatrix}}$.
Let $u_3=0$, $u_1=-1$, and $u_2=1$. Then
${\vect y}_2=\threecol{{-1}}10e^{-2t}+\threecol0{-1}1te^{-2t}$.
For a third solution we need a vector ${\vect v}$ such that
$\dst{\threebythree{-4}{-4}{-4}211233\threecol{v_1}{v_2}{v_3}}
=\dst{\threecol{-1}10}$.
The augmented matrix of this system is row equivalent to
$\dst{\begin{bmatrix}1&0&0&\vdots&\frac{3}{4}\\0&1&1&
\vdots&-\frac{1}{2}\\0&0&0&\vdots&0\end{bmatrix}}$.
Let $v_3=0$, $v_1=\dst\frac{3}{4}$, and $v_2=-\dst\frac{1}{2}$. Then
${\vect y}_3= \threecol3{-2}0\dst\frac{e^{-2t}}{4}+\threecol{-1}10t
e^{-2t}+\threecol0{-1}1\dst\frac{t^2e^{-2t}}{2}$.
The general solution is
${\vect y}=\dst{c_1\threecol0{-1}1e^{-2t}+
c_2\left(\threecol{{-1}}10e^{-2t}+\threecol0{-1}1te^{-2t}\right)}$
\[+c_3\left(\threecol3{-2}0\frac{e^{-2t}}{4}+\threecol{-1}10t
e^{-2t}+\threecol0{-1}1\frac{t^2e^{-2t}}{2}\right)\]


\item
$\dst{\threechar{-2}{-12}{10}2{-24}{11}2{-24}8}
=-(\lambda+6)^3$.
Hence $\lambda_1=-6$.
The eigenvectors
 satisfy the system with  augmented matrix
$\dst{\begin{bmatrix}4&-12&10&\vdots&0\\2&-18&11&
\vdots&0\\2&-24&14&\vdots&0\end{bmatrix}}$,
which is row equivalent to
$\dst{\begin{bmatrix}1&0&1&\vdots&0\\0&1&-\frac{1}{2}&
\vdots&0\\0&0&0&\vdots&0\end{bmatrix}}$.
Hence  $x_1=-x_3$ and $x_2=\dst\frac{x_3}{2}$.  Taking $x_3=2$
yields
${\vect y}_1=\dst{\threecol{-2}12}e^{-6t}$.
For a second solution we need a vector ${\vect u}$ such that
$\dst{\threebythree4{-12}{10}2{-18}{11}2{-24}{14}\threecol{u_1}{u_2}{u_3}}
=\dst{\threecol{-2}12}$.
The augmented matrix of this system is row equivalent to
$\dst{\begin{bmatrix}1&0&1&\vdots&-1\\0&1&-\frac{1}{2}&
\vdots&-\frac{1}{6}\\0&0&0&\vdots&0\end{bmatrix}}$.
Let $u_3=0$, $u_1=-1$, and $u_2=-\dst\frac{1}{6}$. Then
${\vect y}_2=-\threecol610\dst\frac{e^{-6t}}{6}+\threecol{-2}12te^{-6t}$.
For a third solution we need a vector ${\vect v}$ such that
$\dst{\threebythree4{-12}{10}2{-18}{11}2{-24}{14}\threecol{v_1}{v_2}{v_3}}
=\dst{\threecol{-1}{-\frac{1}{6}}0}$.
The augmented matrix of this system is row equivalent to
$\dst{\begin{bmatrix}1&0&1&\vdots&-\frac{1}{3}\\0&1&-\frac{1}{2}&
\vdots&-\frac{1}{36}\\0&0&0&\vdots&0\end{bmatrix}}$.
Let $v_3=0$, $v_1=-\dst\frac{1}{3}$, and $v_2=-\dst\frac{1}{36}$. Then
${\vect y}_3= -\threecol{12}10\dst\frac{e^{-6t}}{36}-\threecol610\frac{t
e^{-6t}}{6}+\threecol{-2}12\frac{t^2e^{-6t}}{2}$.
The general solution is
 ${\vect y}=\dst{c_1\threecol{-2}12e^{-6t}+
c_2\left(-\threecol610\frac{e^{-6t}}{6}+\threecol{-2}12te^{-6t}\right)}$
\[+c_3\left(-\threecol{12}10\frac{e^{-6t}}{36}-\threecol610\frac{t
e^{-6t}}{6}+\threecol{-2}12\frac{t^2e^{-6t}}{2}\right).\]


\item
$\dst{\threechar{-4}0{-1}{-1}{-3}{-1}10{-2}}
=-(\lambda+3)^3$.
Hence $\lambda_1=3$.
The eigenvectors
 satisfy the system with  augmented matrix
$\dst{\begin{bmatrix}-1&0&-1&\vdots&0\\-1&0&-1&
\vdots&0\\1&0&1&\vdots&0\end{bmatrix}}$,
which is row equivalent to
$\dst{\begin{bmatrix}1&0&1&\vdots&0\\0&0&0&
\vdots&0\\0&0&0&\vdots&0\end{bmatrix}}$.
Hence  $x_1=-x_3$ and $x_2$ is arbitrary.  Taking $x_2=0$ and $x_3=1$
yields
${\vect y}_1=\threecol{-1}01e^{-3t}$.
Taking $x_2=1$ and $x_3=0$
yields  ${\vect y}_2=\threecol010e^{-3t}$.
For a third solution we need constants $\alpha$ and $\beta$  and a
vector ${\vect u}$ such that
$\dst{\threebythree{-1}0{-1}{-1}0{-1}101\threecol{u_1}{u_2}{u_3}}
=\alpha\dst{\threecol{-1}01}+\beta\threecol010$.
The augmented matrix of this system is row equivalent to
$\dst{\begin{bmatrix}1&0&1&\vdots&\alpha\\0&0&0&
\vdots&\alpha+\beta\\0&0&0&\vdots&0\end{bmatrix}}$;
hence the system
has a solution if $\alpha=-\beta=1$, which yields the eigenvector
${\vect x}_3=\threecol{-1}{-1}1$. Taking $u_1=1$  and $u_2=u_3=0$
yields the solution
${\vect y}_3=\threecol100e^{-3t}+\threecol{-1}{-1}1te^{-3t}$.
The general solution is
${\vect
y}=\dst{c_1\threecol{-1}01e^{-3t}+c_2\threecol010e^{-3t}+c_3\left
(\threecol100e^{-3t}+\threecol{-1}{-1}1te^{-3t}\right)}$


\item
$\dst{\threechar{-3}{-1}01{-1}0{-1}{-1}{-2}}
=-(\lambda+2)^3$.
Hence $\lambda_1=-2$.
The eigenvectors
 satisfy the system with  augmented matrix
$\dst{\begin{bmatrix}-1&-1&0&\vdots&0\\1&1&0&
\vdots&0\\-1&-1&0&\vdots&0\end{bmatrix}}$,
which is row equivalent to
$\dst{\begin{bmatrix}1&1&0&\vdots&0\\0&0&0&
\vdots&0\\0&0&0&\vdots&0\end{bmatrix}}$.
Hence  $x_1=-x_2$ and $x_3$ is arbitrary.  Taking $x_2=1$ and $x_3=0$
yields
${\vect y}_1=\threecol{-1}10e^{-2t}$.
Taking  $x_2=0$ and $x_3=1$ yields
  ${\vect y}_2=\threecol001e^{-2t}$.
For a third solution we need constants $\alpha$ and $\beta$  and a
vector ${\vect u}$ such that
$\dst{\threebythree{-1}{-1}0110{-1}{-1}0\threecol{u_1}{u_2}{u_3}}
=\alpha\dst{\threecol{-1}10}+\beta\threecol001$.
The augmented matrix of this system is row equivalent to
$\dst{\begin{bmatrix}1&1&0&\vdots&\alpha\\0&0&0&
\vdots&\alpha+\beta\\0&0&0&\vdots&0\end{bmatrix}}$;
hence the system
has a solution if $\alpha=-\beta=1$, which yields the eigenvector
${\vect x}_3=\threecol{-1}1{-1}$. Taking $u_1=1$  and $u_2=u_3=0$
yields the solution
${\vect y}_3=
\threecol100e^{-2t}+\threecol{-1}1{-1}te^{-2t}$.
The general solution is
 ${\vect y}=\dst{
c_1\threecol{-1}10e^{-2t}+c_2\threecol001e^{-2t}+c_3\left
(\threecol100e^{-2t}+\threecol{-1}1{-1}te^{-2t}\right)}$.


\item
\begin{align*}
{\vect y}_3'-A{\vect y}_3&=(\lambda_1I-A){\vect v}e^{\lambda_1t}
+(\lambda_1I-A){\vect u}te^{\lambda_1t}+{\vect u}e^{\lambda_1t}\\
&+(\lambda_1I-A){\vect x}\frac{t^2e^{\lambda_1t}}{2}+{\vect
x}te^{\lambda_1t}\\
&=-{\vect u}e^{\lambda_1t}-{\vect x}te^{\lambda_1t}+{\vect
u}e^{\lambda_1t}+0+{\vect x}te^{\lambda_1t}=0.
\end{align*}
Now suppose that $c_1{\vect y}_1+c_2{\vect y}_2+c_3{\vect y}_3={\vect 0}$.
Then
\[
 c_1{\vect x}+c_2({\vect u}+t{\vect x})+c_3\dst\left({\vect
v}+t{\vect u}+\frac{t^2}{2}{\vect x}\right)={\vect 0}.
\tag{A}
\]
Differentiating this twice yields $c_3{\vect x}=0$, so $c_3=0$ since
${\vect x}\ne{\vect 0}$. Therefore,(A) reduces to
(B) $c_1{\vect x}+c_2({\vect u}+t{\vect x})={\vect 0}$.
Differentiating this yields $c_2{\vect x}=0$, so $c_2=0$ since
${\vect x}\ne{\vect 0}$. Therefore,(B) reduces to $c_3{\vect x}={\vect 0}$,
so $c_1=0$ since
${\vect x}\ne{\vect 0}$. Therefore,${\vect y}_1$, ${\vect y}_2$, and ${\vect
y}_3$ are linearly independent.
\end{sectionSolutions}

\section{Constant Coefficient Homogeneous Systems III}

\begin{sectionSolutions}
\item
$\begin{vmatrix}-11-\lambda&4\\-26&9-\lambda
\end{vmatrix}=(\lambda+1)^2+4$.
The augmented matrix of
$\left(A-\left(-1+2i\right)I\right){\vect x}={\vect 0}$ is
$\begin{bmatrix}-10-2i&4&\vdots&0\\
-26&10-2i&\vdots&0 \end{bmatrix}$,
which is row equivalent to
$\begin{bmatrix} 1&\frac{-5+i}{13}&\vdots&0\\
0&0&\vdots&0
\end{bmatrix}$.
Therefore,$x_1=(5-i)x_2/13$. Taking $x_2=13$ yields the eigenvector
${\vect x}=\begin{bmatrix}5-i\\13\end{bmatrix}$.
Taking  real and imaginary parts of
$e^{-t}(\cos2t+i\sin2t)
)\begin{bmatrix}5-i\\13\end{bmatrix}$ yields
\[
 {\vect y}=\dst{c_1e^{-t}\ctwocol{5\cos2t+\sin2t}{13\cos2t}
+c_2e^{-t}\ctwocol{5\sin2t-\cos2t}{13\sin2t}}.
\]


\item
$\begin{vmatrix} 5-\lambda&-6\\ 3&-1-\lambda
\end{vmatrix}=(\lambda-2)^2+9$.
Hence, $\lambda=2+3i$ is an eigenvalue of $A$. The associated
eigenvectors satisfy $\left(A-\left(2+3i\right)I\right){\vect x}={\vect
0}$. The augmented matrix of this system is
$\begin{bmatrix} 3-3i&-6&\vdots&0\\
3&-3-3i&\vdots&0  \end{bmatrix}$,
which is row equivalent to
$\begin{bmatrix} 1&-1-i&\vdots&0\\
0&0&\vdots&0
\end{bmatrix}$.
Therefore,$x_1=(1+i)x_2$. Taking $x_2=1$ yields $x_1=1+i$, so
${\vect x}=\begin{bmatrix}1+i\\1\end{bmatrix}$
is an eigenvector.  Taking  real and imaginary parts of
$e^{2t}(\cos3t+i\sin3t)
\begin{bmatrix}1+i\\1\end{bmatrix}$ yields
 ${\vect y}=\dst{c_1e^{2t}\ctwocol{\cos3t-\sin3t}{\cos3t}
+c_2e^{2t}\ctwocol{\sin3t+\cos3t}{\sin3t}}$.


\item
$\begin{vmatrix}-3-\lambda&3&1\\1&-5-\lambda&-3\\-3
&7&3-\lambda\end{vmatrix}=-(\lambda+1)\left((\lambda+2)^2+4\right)$.
 The augmented matrix of $(A+I){\vect x=0}$ is
$\begin{bmatrix}-2&3&1&\vdots&0\\1&-4&-3&\vdots&0\\
-3&7&4&\vdots&0\end{bmatrix}$,
which is row equivalent to
$\begin{bmatrix} 1&0&1&\vdots&0\\ 0&1&1&
\vdots&0\\ 0&0&0&\vdots&0\end{bmatrix}$.
Therefore,$x_1=x_2=-x_3$.  Taking $x_3=1$ yields
 ${\vect y}_1=\threecol{-1}{-1}1e^{-t}$.
The augmented matrix of
$\left(A-(-2+2i)I\right){\vect x=0}$ is
$\begin{bmatrix}-1-2i&3&1&\vdots&0\\1&
-3-2i&-3&\vdots&0\\-3&7&5-2i&\vdots&0
\end{bmatrix}$,
which is row equivalent to
$\begin{bmatrix} 1&0&-\frac{1+i}{2}&\vdots&0\\
0&1&\frac{1-i}{2}&
\vdots&0\\ 0&0&0&\vdots&0\end{bmatrix}$.
Therefore,$x_1=\dst\frac{(1+i)}{2}x_3$ and $x_2=-\dst\frac{(1-i)}{2}x_3$.
Taking $x_3=2$ yields the eigenvector
${\vect x}_2=\begin{bmatrix}1+i\\-1+i\\2\end{bmatrix}$.
The real and imaginary parts of
$e^{-2t}(\cos2t+i\sin2t)\begin{bmatrix}1+i\\-1+i\\2\end{bmatrix}$
 are
${\vect y}_2=e^{-2t}\cthreecol{\cos2t-\sin2t}{-\cos2t-\sin2t}{2\cos2t}$
 and
${\vect y}_3=e^{-2t}\cthreecol{\sin2t+\cos2t}{-\sin2t+\cos2t}{2\sin2t}$.
 Therefore,
\[
\dst{{\vect y}=c_1\threecol{-1}{-1}1e^{-t}+c_2e^{-2t}\cthreecol{\cos2t-\sin
2t}{-\cos2t-\sin  2t}{2\cos2t}+c_3e^{-2t}\cthreecol
{\sin2t+\cos2t}{-\sin2t+\cos2t}{2\sin2t}}.
\]


\item
$\begin{vmatrix}-3-\lambda&1&-3\\4&-1-\lambda&2\\4
&-2&3-\lambda\end{vmatrix}=-(\lambda-1)\left((\lambda+1)^2+4\right)$.
 The augmented matrix of $(A-I){\vect x=0}$ is
$\begin{bmatrix}-4&1&-3&\vdots&0\\4&-2&2&\vdots&0\\
4&-2&2&\vdots&0\end{bmatrix}$,
which is row equivalent to
$\begin{bmatrix} 1&0&1&\vdots&0\\ 0&1&1&
\vdots&0\\ 0&0&0&\vdots&0\end{bmatrix}$.
Therefore,$x_1=x_2=-x_3$.  Taking $x_3=1$ yields
${\vect y}_1=\threecol{-1}{-1}1e^t$. The augmented matrix of
$\left(A-(-1+2i)I\right){\vect x=0}$ is
$\begin{bmatrix}-2-2i&1&-3&\vdots&0\\4&-2i&2&\vdots&0\\
4&-2&4-2i&\vdots&0\end{bmatrix}$,
which is row equivalent to
$\begin{bmatrix} 1&0&\frac{1-i}{2}&\vdots&0\\
0&1&-1&\vdots&0\\ 0&0&0&\vdots&0\end{bmatrix}$.
Therefore,$x_1=-\dst\frac{1-i}{2}x_3$ and $x_2=x_3$.
Taking $x_3=2$ yields the eigenvector
${\vect x}_2=\begin{bmatrix}-1+i\\2\\2\end{bmatrix}$.
The real and imaginary parts of
$e^{-t}(\cos2t+i\sin2t)\begin{bmatrix}-1+i\\2\\2\end{bmatrix}$
are ${\vect y}_2=e^{-t}\cthreecol{-\sin2t-\cos2t}{2\cos2t}{2\cos2t}$
and ${\vect y}_3=e^{-t}\cthreecol {\cos2t-\sin2t}{2\sin2t}{2\sin2t}$.
Therefore,
\[
\dst{{\vect
y}=c_1\threecol{-1}11e^{-t}+c_2e^{-t}\cthreecol{-\sin
2t-\cos2t}{2\cos2t}{2\cos2t}+c_3e^{-t}\cthreecol
{\cos2t-\sin2t}{2\sin2t}{2\sin2t}}.
\]


\item
$\dst\frac{1}{3}\begin{vmatrix}7-3\lambda&-5\\2&5-3\lambda
\end{vmatrix}=\dst{\left(\lambda-2\right)^2+1}$.
The augmented matrix of
$\dst{\left(A-(2+i)I\right)}{\vect x}={\vect 0}$ is
$\dst\frac{1}{3}\begin{bmatrix}1-3i&-5&\vdots&0\\2&-1-3i&\vdots&0
\end{bmatrix}$, which is row equivalent to
$\begin{bmatrix} 1&-\frac{1+3i}{2}&\vdots&0\\ 0&0&\vdots&0
\end{bmatrix}$. Therefore,$x_1=\dst\frac{1+3i}{2}x_2$. Taking
$x_2=2$ yields the eigenvector
${\vect x}=\begin{bmatrix}1+3i\\2\end{bmatrix}$.
 Taking real and imaginary parts of
$e^{2t}(\cos t+i\sin t)
\begin{bmatrix}1+3i\\2\end{bmatrix}$
yields
\[
\dst{{\vect y}=c_1e^{2t}\begin{bmatrix}\cos t-3\sin t
\\2\cos t\end{bmatrix}+
c_2e^{2t}\begin{bmatrix}
\sin t+3\cos t\\ 2\sin t\end{bmatrix}}.
\]


\item
$\begin{vmatrix}34-\lambda&52\\-20&-30-\lambda
\end{vmatrix}=(\lambda-2)^2+16$. The augmented matrix of
$\left(A-\left(2+4i\right)I\right){\vect x}={\vect 0}$ is
$\begin{bmatrix}32-4i&52&\vdots&0\\-20&-32-4i&\vdots&0
\end{bmatrix}$, which is row equivalent to
$\begin{bmatrix} 1&\frac{8+i}{5}&\vdots&0\\ 0&0&\vdots&0
\end{bmatrix}$. Therefore,$x_1=-\dst\frac{(8+i)}{5}x_2$. Taking
$x_2=5$ yields the eigenvector ${\vect
x}=\begin{bmatrix}-8-i\\5\end{bmatrix}$.
Taking real and imaginary parts of $e^{2t}(\cos4t+i\sin
4t)\begin{bmatrix}-8-i\\5\end{bmatrix}$ yields $\dst{{\vect
y}=c_1e^{2t}\begin{bmatrix}\sin4t-8\cos4t\\
5\cos4t\end{bmatrix}+c_2e^{2t}\begin{bmatrix}
-\cos4t-8\sin4t\\ 5\sin4t\end{bmatrix}}$.


\item
$\begin{vmatrix}3-\lambda&-4&-2\\-5&7-\lambda&-8\\-10
&13&-8-\lambda\end{vmatrix}=-(\lambda+2)\left((\lambda-2)^2+9\right)$.
 The augmented matrix of $(A+2I){\vect x=0}$ is
$\begin{bmatrix}5&-4&-2&\vdots&0\\-5&9&-8&\vdots&0\\
-10&13&-6&\vdots&0\end{bmatrix}$,
which is row equivalent to
$\begin{bmatrix} 1&0&-2&\vdots&0\\ 0&1&-2&
\vdots&0\\ 0&0&0&\vdots&0\end{bmatrix}$.
Therefore,$x_1=x_2=2x_3$.  Taking $x_3=1$ yields
${\vect y}_1=\threecol221e^{-2t}$.
The augmented matrix of
$\left(A-(2+3i)I\right){\vect x=0}$ is
$\begin{bmatrix}1-3i&-4&-2&\vdots&0\\-5&5-3i&-8&\vdots&0\\
-10&13&-10-3i&\vdots&0\end{bmatrix}$,
which is row equivalent to
$\begin{bmatrix} 1&0&1-i&\vdots&0\\
0&1&-i&\vdots&0\\ 0&0&0&\vdots&0\end{bmatrix}$.
Therefore,$x_1=-(1-i)x_3$ and $x_2=ix_3$.
Taking $x_3=1$ yields the eigenvector
${\vect x}_2=\begin{bmatrix}-1+i\\i\\1\end{bmatrix}$.
The real and imaginary parts of $e^{2t}(\cos3t+i\sin3t)
\begin{bmatrix}-1+i\\i\\1\end{bmatrix}$ are
${\vect y}_2=e^{2t}\begin{bmatrix}-\cos3t-\sin3t\\-\sin3t\\
\cos3t\end{bmatrix}$ and ${\vect y}_3=
c_3e^{2t}\begin{bmatrix}-\sin3t+\cos3t\\\cos3t\\
\sin3t\end{bmatrix}$. Therefore,
\[
\dst{{\vect
y}=c_1\begin{bmatrix} 2\\ 2\\1 \end{bmatrix}e^{-2t}+
c_2e^{2t}\begin{bmatrix}-\cos3t-\sin3t\\-\sin3t\\
\cos3t\end{bmatrix}+c_3e^{2t}\begin{bmatrix}
-\sin3t+\cos3t\\\cos3t\\\sin3t\end{bmatrix}}.
\]


\item
$\begin{vmatrix}1-\lambda&2&-2\\0&2-\lambda&-1\\1
&0&-\lambda\end{vmatrix}=-(\lambda-2)\left((\lambda-1)^2+1\right)$.
The augmented matrix of $(A-I){\vect x=0}$ is
$\begin{bmatrix}0&2&-2&\vdots&0\\0&
1&-1&\vdots&0\\ 1&0&-1&\vdots&0
\end{bmatrix}$,
which is row equivalent to
$\begin{bmatrix} 1&0&-1&\vdots&0\\ 0&1&-1&
\vdots&0\\ 0&0&0&\vdots&0\end{bmatrix}$.
Therefore,$x_1=x_2=1$.  Taking $x_3=1$ yields
 ${\vect y}_1=\threecol111e^t$.
The augmented matrix of
$\left(A-(1+i)I\right){\vect x=0}$ is
$\begin{bmatrix}-i&2&-2&\vdots&0\\0&
1-i&-1&\vdots&0\\ 1&0&-1-i&\vdots&0
\end{bmatrix}$,
which is row equivalent to
$\begin{bmatrix} 1&0&-1-i&\vdots&0\\ 0&1&-\frac{1+i}{2}&
\vdots&0\\ 0&0&0&\vdots&0\end{bmatrix}$.
Therefore,$x_1=(1+i)x_3$ and $x_2=\dst\frac{(1+i)}{2}x_3$.
Taking $x_3=2$ yields the eigenvector
${\vect x}_2=\begin{bmatrix}2+2i\\1+i\\2\end{bmatrix}$.
The real and imaginary parts of
$e^{4t}(\cos t+i\sin t)\begin{bmatrix}2+2i
\\1+i\\2\end{bmatrix}$  are
${\vect y}_2=e^t\begin{bmatrix} 2\cos t-2\sin t\\\cos t-\sin
t\\ 2\cos t\end{bmatrix}$ and
${\vect y}_3=
c_3e^t\begin{bmatrix} 2\sin t+2\cos t\\\cos t+\sin t\\
2\sin t\end{bmatrix}$. Therefore,
\[
\dst{{\vect y}=c_1\begin{bmatrix} 1\\1\\1
\end{bmatrix}e^t+
c_2e^t\begin{bmatrix} 2\cos t-2\sin t
\\\cos t-\sin t\\ 2\cos t\end{bmatrix}+
c_3e^t\begin{bmatrix} 2\sin t+2\cos t\\\cos t+\sin t\\
2\sin t\end{bmatrix}}.
\]


\item
$\begin{vmatrix}7-\lambda&15\\-3&1-\lambda
\end{vmatrix}=(\lambda-4)^2+36$. The augmented matrix of
$\left(A-\left(4+6i\right)I\right){\vect x}={\vect 0}$ is
$\begin{bmatrix}3-6i&15&\vdots&0\\-3&-3-6i&\vdots&0
\end{bmatrix}$, which is row equivalent to
$\begin{bmatrix} 1&1+2i&\vdots&0\\ 0&0&\vdots&0
\end{bmatrix}$. Therefore,$x_1=-(1+2i)x_2$. Taking $x_2=1$ yields
the eigenvector ${\vect
x}=\begin{bmatrix}-1-2i\\1\end{bmatrix}$.
Taking real and imaginary parts of
$e^{4t}(\cos6t+i\sin6t)\begin{bmatrix}-1-2i\\1\end{bmatrix}$
yields ${\vect y}=\dst{c_1e^{4t}\ctwocol{2\sin6t-\cos6t}{\cos6t}
+c_2e^{4t}\ctwocol{-2\cos6t-\sin6t}{\sin6t}}$. Now ${\vect
y}(0)=\twocol51\Rightarrow
\twobytwo{-1}{-2}10\twocol{c_1}{c_2}=\twocol51$, so $c_1=1$, $c_2=-3$,
and $\dst{{\vect
y}=e^{4t}\begin{bmatrix}5\cos6t+5\sin6t\\\cos6t-3\sin6t
\end{bmatrix}}$.


\item
$\dst\frac{1}{6}\begin{vmatrix}4-6\lambda&-2\\5&2-6\lambda
\end{vmatrix}=\dst{\left(\lambda-\frac{1}{2}\right)^2+\frac{1}{4}}$.
The augmented matrix of
$\dst{\left(A-\frac{1+i}{2}I\right)}{\vect x}={\vect 0}$ is
$\dst\frac{1}{6}\begin{bmatrix}1-3i&-2&\vdots&0\\5&-1-3i&\vdots&0
\end{bmatrix}$, which is row equivalent to
$\begin{bmatrix} 1&-\frac{1+3i}{5}&\vdots&0\\ 0&0&\vdots&0
\end{bmatrix}$. Therefore,$x_1=\dst\frac{1+3i}{5}x_2$. Taking
$x_2=5$ yields the eigenvector
${\vect x}=\begin{bmatrix}1+3i\\5\end{bmatrix}$.
 Taking real and imaginary parts of
$e^{t/2}(\cos t/2+i\sin
t/2)\begin{bmatrix}1+3i\\5\end{bmatrix}$
yields ${\vect y}=\dst{c_1e^{t/2}\ctwocol{\cos t/2-3\sin t/2}{5\cos t/2}
+c_2e^{t/2}\ctwocol{\sin t/2+3\cos t/2}{5\sin t/2}}$. Now ${\vect
y}(0)=\twocol1{-1}\Rightarrow
\twobytwo1350\twocol{c_1}{c_2}=\twocol1{-1}$, so
$c_1=-\dst\frac{1}{5}$,
$c_2=\dst\frac{2}{5}$, and
${\vect y}=\dst{e^{t/2}\begin{bmatrix}\cos(t/2)+
\sin(t/2)\\-\cos(t/2)+2\sin(t/2)
\end{bmatrix}}$.


\item
$\begin{vmatrix}4-\lambda&4&0\\8&10-\lambda&-20\\2
&3&-2-\lambda\end{vmatrix}=-(\lambda-8)\left((\lambda-2)^2+4\right)$.
 The augmented matrix of $(A-8I){\vect x=0}$ is
$\begin{bmatrix}0&4&0&\vdots&0\\8&6&-20&\vdots&0\\
2&3&-6&\vdots&0\end{bmatrix}$,
which is row equivalent to
$\begin{bmatrix} 1&0&-2&\vdots&0\\ 0&1&-2&
\vdots&0\\ 0&0&0&\vdots&0\end{bmatrix}$.
Therefore,$x_1=x_2=2x_3$. Taking $x_3=2$ yields
${\vect y}_1=\threecol221e^{8t}$.
The augmented matrix of
$\left(A-(2+2i)I\right){\vect x=0}$ is
$\begin{bmatrix}2-2i&4&0&\vdots&0\\8&
8-2i&-20&\vdots&0\\2&3&-4-2i&\vdots&0\end{bmatrix}$,
which is row equivalent to
$\begin{bmatrix} 1&0&-2+2i&\vdots&0\\
0&1&-2i&\vdots&0\\ 0&0&0&\vdots&0\end{bmatrix}$.
Therefore,$x_1=(2-2i)x_3$ and $x_2=2ix_3$.
Taking $x_3=1$ yields the eigenvector
${\vect x}_2=\begin{bmatrix}2-2i\\2i\\1\end{bmatrix}$.
The real and imaginary parts of
$e^{2t}(\cos2t+i\sin
2t)\begin{bmatrix}2-2i\\2i\\1\end{bmatrix}$
are ${\vect y}_2=e^{2t}\begin{bmatrix}2\cos2t+2\sin2t
\\-2\sin2t\\2\cos2t\end{bmatrix}$ and
${\vect y}_3=e^{2t}\begin{bmatrix}2\sin2t+2\cos2t
\\2\cos2t\\\sin2t\end{bmatrix}$, so the general solution
is ${\vect y}=c_1\threecol221e^{8t}+c_2
e^{2t}\begin{bmatrix}2\cos2t+2\sin2t
\\-2\sin2t\\2\cos2t\end{bmatrix}+
c_3e^{2t}\begin{bmatrix}2\sin2t+2\cos2t
\\2\cos2t\\\sin2t\end{bmatrix}$.
Now ${\vect y}(0)=\threecol865\Rightarrow\threebythree22{-2}202110
\threecol{c_1}{c_2}{c_3}=\threecol865$, so $c_1=2$, $c_2=3$,
$c_3=1$, and $\dst{{\vect y}=\threecol442e^{8t}
+e^{2t}\begin{bmatrix}4\cos2t+8\sin2t\\
-6\sin2t+2\cos2t\\3\cos2t+\sin2t
\end{bmatrix}}$.


\item
$\begin{vmatrix}4-\lambda&-4&4\\-10&3-\lambda&15\\2
&-3&1-\lambda\end{vmatrix}=-(\lambda-8)(\lambda^2+16)$.
 The augmented matrix of $(A-8I){\vect x=0}$ is
$\begin{bmatrix}-4&-4&4&\vdots&0\\-10&-5&15&\vdots&0\\
2&-3&-7&\vdots&0\end{bmatrix}$,
which is row equivalent to
$\begin{bmatrix} 1&0&-2&\vdots&0\\ 0&1&1&
\vdots&0\\ 0&0&0&\vdots&0\end{bmatrix}$.
Therefore,$x_1=2x_3$ and $x_2=-x_3$. Taking $x_3=1$ yields
${\vect y}_1=\threecol2{-1}1e^{8t}$.
The augmented matrix of
$\left(A-4iI\right){\vect x=0}$ is
$\begin{bmatrix}4-4i&-4&4&\vdots&0\\-10&3-4i&15&\vdots&0\\
2&-3&1-4i&\vdots&0\end{bmatrix}$,
which is row equivalent to
$\begin{bmatrix} 1&0&-1+i&\vdots&0\\
0&1&-1+2i&\vdots&0\\ 0&0&0&\vdots&0\end{bmatrix}$.
Therefore,$x_1=(1-i)x_3$ and $x_2=(1-2i)x_3$.
Taking $x_3=1$ yields the eigenvector
${\vect x}_2=\begin{bmatrix}1-i\\1-2i\\1\end{bmatrix}$.
The real and imaginary parts of
$(\cos4t+i\sin
4t)\begin{bmatrix}1-i\\1-2i\\1\end{bmatrix}$
are ${\vect y}_2=\begin{bmatrix}\cos4t+\sin4t
\\\cos4t+2\sin4t\\\cos4t\end{bmatrix}$ and
${\vect y}_3=\begin{bmatrix}\sin4t-\cos4t
\\\sin4t-2\cos4t\\\sin4t\end{bmatrix}$, so the general solution
is ${\vect y}=c_1 \threecol2{-1}1e^{8t}+
c_2\begin{bmatrix}\cos4t+\sin4t
\\\cos4t+2\sin4t\\\cos4t\end{bmatrix}
+c_3\begin{bmatrix}\sin4t-\cos4t
\\\sin4t-2\cos4t\\\sin4t\end{bmatrix}$.
Now ${\vect
y}(0)=\threecol{16}{14}6\Rightarrow\threebythree21{-1}{-1}1{-2}110
\threecol{c_1}{c_2}{c_3}=\threecol{16}{14}6$, so $c_1=3$, $c_2=3$,
$c_3=-7$, and $\dst{{\vect y}=\threecol6{-3}3e^{8t}
+\begin{bmatrix}10\cos4t-4\sin4t\\17\cos4t-\sin4t\\3\cos4t-7\sin4t
\end{bmatrix}}$.

\addtocounter{enumi}{2}
\item%10.6.28
\begin{alist}
\item From the quadratic formula the roots are
\begin{align*}
k_1&=\frac{\|{\vect u}\|^2-\|{\vect v}^2\|+\sqrt{(\|{\vect u}\|^2-\|{\vect
v}^2\|)^2+4({\vect u},{\vect v})^2}}{2({\vect u},{\vect v})}\\
k_2&=\frac{\|{\vect u}\|^2-\|{\vect v}^2\|-\sqrt{(\|{\vect u}\|^2-\|{\vect
v}^2\|)^2+4({\vect u},{\vect v})^2}}{2({\vect u},{\vect v})}.
\end{align*}
Clearly $k_1>0$ and $k_2<0$. Moreover,
\[
k_1k_2=\frac{(\|{\vect u}\|^2-\|{\vect v}^2\|)^2-\left[(\|{\vect u}\|^2-\|{\vect
v}^2\|)^2+4({\vect u},{\vect v})^2\right]}{4({\vect u},{\vect v})^2}=-1.
\]

\item Since $k_2=-1/k_1$,
\begin{align*}
{\vect u}_1^{(2)}&={\vect u}-k_2{\vect v}={\vect u}+\frac{1}{ k_1}{\vect
v}=\frac{1}{ k_1}({\vect v}+k_1{\vect u})=\frac{1}{ k_1}{\vect v}_1^{(1)}\\
{\vect v}_1^{(2)}&={\vect v}+k_2{\vect u}={\vect v}-\frac{1}{ k_1}{\vect
u}=-\frac{1}{ k_1}({\vect u}-k_1{\vect v})=-\frac{1}{ k_1}{\vect u}_1^{(1)}.
\end{align*}
\end{alist}

\item
$\begin{vmatrix}-15-\lambda&10\\-25&15-\lambda
\end{vmatrix}=\lambda^2+25$.
The augmented matrix of $(A-5iI){\vect x}={\vect 0}$ is
$\begin{bmatrix}-15-5i&10&\vdots&0\\
-25&15-5i&\vdots&0 \end{bmatrix}$,
which is row equivalent to
$\begin{bmatrix} 1&\frac{-3+i}{5}&\vdots&0\\ 0&0&\vdots&0
\end{bmatrix}$.
Therefore,$x_1=\dst\frac{(3-i)}{5}x_2$. Taking $x_2=5$
yields the eigenvector
 ${\vect x}=\begin{bmatrix}3-i\\5\end{bmatrix}$,
so ${\vect u}=\dst{\twocol35}$ and ${\vect v}=\dst{\twocol{-1}0}$.
The quadratic equation is $-3k^2-33k+3=0$, with positive root
$k\approx.0902$.
 Routine calculations yield
 ${\vect U}\approx\dst{\twocol{.5257}{.8507}}$,
${\vect V}\approx\dst{\twocol{-.8507}{.5257}}$.


\item
$\begin{vmatrix}-3-\lambda&-15\\3&3-\lambda
\end{vmatrix}=\lambda^2+36$.
The augmented matrix of $(A-6iI){\vect x}={\vect 0}$ is
$\begin{bmatrix}-3-6i&-15&\vdots&0\\
3&3-6i&\vdots&0 \end{bmatrix}$,
which is row equivalent to
$\begin{bmatrix} 1&1-2i&\vdots&0\\ 0&0&\vdots&0
\end{bmatrix}$.
Therefore,$x_1=-(1-2i)x_2$. Taking $x_2=1$
yields the eigenvector
 ${\vect x}=\begin{bmatrix}-1+2i\\1\end{bmatrix}$,
so ${\vect u}=\dst{\twocol{-1}1}$ and ${\vect v}=\dst{\twocol20}$.
The quadratic equation is $-2k^2+2k+2=0$, with positive root
$k\approx1.6180$.
 Routine calculations yield
 ${\vect U}\approx\dst{\twocol{-.9732}{.2298}}$,
${\vect V}\approx\dst{\twocol{.2298}{.9732}}$.



\item
$\begin{vmatrix}5-\lambda&-12\\6&-7-\lambda
\end{vmatrix}=(\lambda+1)^2+36$.
The augmented matrix of $(A-(-1+6i)I){\vect x}={\vect 0}$ is
$\begin{bmatrix}6-6i&-12&\vdots&0\\
6&-6-6i&\vdots&0 \end{bmatrix}$,
which is row equivalent to
$\begin{bmatrix} 1&-(1+i)&\vdots&0\\ 0&0&\vdots&0
\end{bmatrix}$.
Therefore,$x_1=(1+i)x_2$. Taking $x_2=1$
yields the eigenvector
 ${\vect x}=\begin{bmatrix}1+i\\1\end{bmatrix}$,
so ${\vect u}=\dst{\twocol11}$ and ${\vect v}=\dst{\twocol10}$.
The quadratic equation is $k^2-k-1=0$, with positive root
$k\approx1.6180$.
 Routine calculations yield
 ${\vect U}\approx\dst{\twocol{-.5257}{.8507}}$,
${\vect V}\approx\dst{\twocol{.8507}{.5257}}$.


\item
$\begin{vmatrix}-4-\lambda&9\\-5&2-\lambda
\end{vmatrix}=(\lambda+1)^2+36$.
The augmented matrix of $(A-(-1+6i)I){\vect x}={\vect 0}$ is
$\begin{bmatrix}-3-6i&9&\vdots&0\\
-5&3-6i&\vdots&0 \end{bmatrix}$,
which is row equivalent to
$\begin{bmatrix} 1&-\frac{3-6i}{5}&\vdots&0\\ 0&0&\vdots&0
\end{bmatrix}$.
Therefore,$x_1=\dst\frac{3-6i}{5}x_2$. Taking $x_2=5$
yields the eigenvector
 ${\vect x}=\begin{bmatrix}3-6i\\5\end{bmatrix}$,
so ${\vect u}=\dst{\twocol35}$ and ${\vect v}=\dst{\twocol{-6}0}$.
The quadratic equation is $-18k^2+2k+18=0$, with positive root
$k\approx1.0571$.
 Routine calculations yield
 ${\vect U}\approx\dst{\twocol{.8817}{.4719}}$,
${\vect V}\approx\dst{\twocol{-.4719}{.8817}}$.


\item
$\begin{vmatrix}-1-\lambda&-5\\20&-1-\lambda
\end{vmatrix}=(\lambda+1)^2+100$.
The augmented matrix of $(A-(-1+10i)I){\vect x}={\vect 0}$ is
$\begin{bmatrix}-10i&-5&\vdots&0\\
20&-10i&\vdots&0 \end{bmatrix}$,
which is row equivalent to
$\begin{bmatrix} 1&-\frac{i}{2}&\vdots&0\\ 0&0&\vdots&0
\end{bmatrix}$.
Therefore,$x_1=\dst\frac{i}{2}x_2$. Taking $x_2=2$
yields the eigenvector
 ${\vect x}=\begin{bmatrix}i\\2\end{bmatrix}$,
so ${\vect u}=\dst{\twocol02}$ and ${\vect v}=\dst{\twocol10}$.
Since $({\vect u},{\vect v})=0$ we just normalize ${\vect u}$ and ${\vect v}$
to obtain
${\vect U}=\dst{\twocol01}$, ${\vect V}=\dst{\twocol10}$.


\item
$\begin{vmatrix}-7-\lambda&6\\-12&5-\lambda
\end{vmatrix}=(\lambda+1)^2+36$.
The augmented matrix of $(A-(-1+6i)I){\vect x}={\vect 0}$ is
$\begin{bmatrix}-6-6i&6&\vdots&0\\
-12&6-6i&\vdots&0 \end{bmatrix}$,
which is row equivalent to
$\begin{bmatrix} 1&-\frac{1-i}{2}&\vdots&0\\ 0&0&\vdots&0
\end{bmatrix}$.
Therefore,$x_1=\dst\frac{1-i}{2}x_2$. Taking $x_2=2$
yields the eigenvector
 ${\vect x}=\begin{bmatrix}1-i\\2\end{bmatrix}$,
so ${\vect u}=\dst{\twocol12}$ and ${\vect v}=\dst{\twocol{-1}0}$.
The quadratic equation is $-k^2-4k+1=0$, with positive root
$k\approx.2361$.
 Routine calculations yield
 ${\vect U}\approx\dst{\twocol{.5257}{.8507}}$,
${\vect V}\approx\dst{\twocol{-.8507}{.5257}}$.
\end{sectionSolutions}


\section{Variation of Parameters for Nonhomogeneous Linear Systems}

\begin{sectionSolutions}
\item
$Y=\begin{bmatrix}
-3e^{-t}&-e^{-2t}\\e^{-t}&2e^{-2t}\end{bmatrix}$;
${\vect u}'=Y^{-1}{\vect f}=
\dst\frac{1}{5}\begin{bmatrix}-2e^t&-e^t
\\e^{2t}&3e^{2t}\end{bmatrix}
\begin{bmatrix}
50e^{37}\\10e^{-3t}
\end{bmatrix}=
\begin{bmatrix}-20e^{4t}-2e^{-2t}\\
10e^{5t}+6e^{-t}\end{bmatrix}$;
${\vect u}=
\begin{bmatrix}e^{-2t}-5e^{4t}\\
2e^{5t}-6e^{-t}\end{bmatrix}$;
${\vect y}_p=Y{\vect u}=
\begin{bmatrix}13e^{3t}+3e^{-3t}\\
-e^{3t}-11e^{-3t}\end{bmatrix}$.




\item
$Y=\begin{bmatrix}
-e^{2t}&-e^{-t}\\2e^{2t}&e^{-t}\end{bmatrix}$;
${\vect u}'=Y^{-1}{\vect f}=
\begin{bmatrix}e^{-2t}&e^{-2t}\\
-2e^t&-e^t\end{bmatrix}
\begin{bmatrix}2\\-2e^t\end{bmatrix}=
\begin{bmatrix}2e^{-2t}-2e^{-t}\\2e^{2t}-4e^t\end{bmatrix}$;
${\vect u}=
\begin{bmatrix}2e^{-t}-e^{-2t}\\
e^{2t}-4e^t\end{bmatrix}$;
${\vect y}_p=Y{\vect u}=
\begin{bmatrix}5-3e^t\\5e^t-6\end{bmatrix}$.


\item
$Y=\begin{bmatrix}\sin t&-\cos t\\
\cos t&\sin t\end{bmatrix}$;
${\vect u}'=Y^{-1}{\vect f}=
\begin{bmatrix}\sin t&\cos t\\-\cos t&\sin t\end{bmatrix}
\begin{bmatrix}1\\t\end{bmatrix}=
\begin{bmatrix}t\cos t+\sin t\\t\sin t-\cos
t\end{bmatrix}$;
${\vect u}=
\begin{bmatrix}t\sin t\\-t\cos t\end{bmatrix}$;
${\vect y}_p=Y{\vect u}=
\begin{bmatrix}t\\0\end{bmatrix}$.


\item
$Y=\begin{bmatrix}
e^{3t}&e^{2t}&e^{-t}\\-e^{3t}&0&-3e^{-t}\\e^{3t}&e^{2t}&7
e^{-t}\end{bmatrix}$;
${\vect u}'=Y^{-1}{\vect f}=
\dst\frac{1}{6}\begin{bmatrix}
3e^{-3t}&-6e^{-3t}&-3e^{-3t}\\4e^{-2t}&6e^{-2t}&2e^{-2t}\\
-e^t&0&e^t\end{bmatrix}
\begin{bmatrix}1\\e^t\\e^t\end{bmatrix}=
\dst\frac{1}{6}\begin{bmatrix}
3e^{-3t}-9e^{-2t}\\8e^{-t}+4e^{-2t}\\e^{2t}-e^t\end{bmatrix}$;
${\vect u}=
\dst\frac{1}{12}\begin{bmatrix}9e^{-2t}-2e^{-3t}\\
-16e^{-t}-4e^{-2t}\\e^{2t}-2e^t\end{bmatrix}$;
${\vect y}_p=Y{\vect u}=
-\dst\frac{1}{6}\begin{bmatrix}3e^t+4\\6e^t-4\\10
\end{bmatrix}$.


\item
$Y=\begin{bmatrix}
-e^t&e^{-t}&te^{-t}\\e^t&-e^{-t}&3e^{-t}-te^{-t}\\
e^t&e^{-t}&te^{-t}\end{bmatrix}$;
${\vect u}'=Y^{-1}{\vect f}=
\dst\frac{1}{18}\begin{bmatrix}
-9e^{-t}&0&9e^{-t}\\3e^t(3-2t)&-6te^t&9e^t\\
6e^t&6e^t&0\end{bmatrix}
\begin{bmatrix}e^t\\e^t\\e^t\end{bmatrix}=
\dst\frac{1}{3}\begin{bmatrix}0\\
e^{2t}(3-2t)\\2e^{2t}\end{bmatrix}$;
${\vect u}=
\dst\frac{1}{3}\begin{bmatrix}0\\
e^{2t}(2-t)\\e^{2t}\end{bmatrix}$;
${\vect y}_p=Y{\vect u}=
\dst\frac{1}{3}\begin{bmatrix}
2e^t\\e^t\\2e^t\end{bmatrix}$.


\item
${\vect u}'=Y^{-1}{\vect f}=
\dst\frac{1}{2t}\begin{bmatrix}
e^{-t}&e^{-t}\\e^t&-e^t\end{bmatrix}
\begin{bmatrix}t\\t^2\end{bmatrix}=
\frac{1}{2}\begin{bmatrix}
e^{-t}(t+1)\\e^t(1-t)\end{bmatrix}$;
${\vect u}=
\dst\frac{1}{2}\begin{bmatrix}
-e^{-t}(t+2)\\e^t(2-t)\end{bmatrix}$;
${\vect y}_p=Y{\vect u}=
\dst\frac{t}{2}\begin{bmatrix}
e^t&e^{-t}\\e^t&-e^{-t}\end{bmatrix}
\begin{bmatrix}-e^{-t}(t+2)\\
e^t(2-t)\end{bmatrix}=
-\begin{bmatrix}t^2\\2t\end{bmatrix}$.



\item
${\vect u}'=Y^{-1}{\vect f}=
\dst\frac{1}{3}\begin{bmatrix}2&-e^{-t}\\-e^t&2\end{bmatrix}
\begin{bmatrix}e^{2t}\\e^{-2t}\end{bmatrix}=
\dst\frac{1}{3}\begin{bmatrix}
2e^{2t}-e^{-3t}\\2e^{-2t}-e^{3t}\end{bmatrix}$;
${\vect u}=
\dst\frac{1}{9}\begin{bmatrix}
3e^{2t}+e^{-3t}\\-e^{3t}-3e^{-2t}\end{bmatrix}$;
${\vect y}_p=Y{\vect u}=
\dst\frac{1}{9}\begin{bmatrix}
2&e^{-t}\\e^t&2\end{bmatrix}
\begin{bmatrix}
3e^{2t}+e^{-3t}\\-e^{3t}-3e^{-2t}\end{bmatrix}=
\frac{1}{9}\begin{bmatrix}
5e^{2t}-e^{-3t}\\e^{3t}-5e^{-2t}\end{bmatrix}$.


\item
${\vect u}'=Y^{-1}{\vect f}=
\dst\frac{1}{ t^2-1}\begin{bmatrix}t&-e^{-t}\\-e^t&t\end{bmatrix}
\begin{bmatrix}t^2-1\\t^2-1\end{bmatrix}=
\begin{bmatrix}t-e^{-t}\\t-e^t\end{bmatrix}$;
${\vect u}=
\dst\frac{1}{2}\begin{bmatrix}
2e^{-t}+t^2\\t^2-2e^t\end{bmatrix}$;
${\vect y}_p=Y{\vect u}
=\dst\frac{1}{2}\begin{bmatrix}t&e^{-t}\\e^t&t\end{bmatrix}
\begin{bmatrix}
2e^{-t}+t^2\\t^2-2e^t\end{bmatrix}=
\dst\frac{1}{2}\begin{bmatrix}
te^{-t}(t+2)+t^3-2\\te^t(t-2)+t^3+2\end{bmatrix}$.


\item
${\vect u}'=Y^{-1}{\vect f}=
\dst\frac{1}{3}\begin{bmatrix}e^{-5t}&e^{-4t}&e^{-3t}\\
2e^{-2t}&-e^{-t}&-1\\-e^{-2t}&2e^{-t}&-1\end{bmatrix}
\begin{bmatrix}e^t\\0\\0\end{bmatrix}=
\dst\frac{1}{3}\begin{bmatrix}
e^{-4t}\\2e^{-t}\\-e^{-t}\end{bmatrix}$;
${\vect u}=
\dst\frac{1}{12}\begin{bmatrix}
-e^{-4t}\\-8e^{-t}\\4e^{-t}\end{bmatrix}$;
${\vect y}_p=Y{\vect u}=
\dst\frac{1}{12}\begin{bmatrix}e^{5t}&e^{2t}&0\\
e^{4t}&0&e^t\\e^{3t}&-1&-1\end{bmatrix}
\begin{bmatrix}
-e^{-4t}\\-8e^{-t}\\4e^{-t}\end{bmatrix}=
\dst\frac{1}{4}\begin{bmatrix}-3e^t\\
1\\e^{-t}\end{bmatrix}$.


\item
${\vect u}'=Y^{-1}{\vect f}=
\dst\frac{1}{2t}\begin{bmatrix}
-e^{-2t}&te^{-t}&te^{-t}+e^{-2t}\\1&-te^t&te^t-1\\
2&0&-2\end{bmatrix}
\begin{bmatrix}e^t\\0\\e^t\end{bmatrix}=
\begin{bmatrix}1\\e^{2t}\\0\end{bmatrix}$;
${\vect u}=
\dst\frac{1}{4}\begin{bmatrix}
2t\\e^{2t}\\0\end{bmatrix}$;
${\vect y}_p=Y{\vect u}=
Y=\dst\frac{1}{4t}\begin{bmatrix}e^t&e^{-t}&t\\
e^t&-e^{-t}&e^{-t}\\e^t&e^{-t}&0\end{bmatrix}
\begin{bmatrix}2t\\e^{2t}\\0\end{bmatrix}=
\frac{e^{t}}{4t}\begin{bmatrix}
2t+1\\2t-1\\2t+1\end{bmatrix}$.


\item
\begin{alist}[start=3]
\item
If ${\vect y}_p=Y{\vect u}$, then ${\vect y}_p'=Y'{\vect u}+Y{\vect u}'=AY{\vect
u}+Y{\vect u}'$, so (E) ${\vect y}_p'=A{\vect y_p}+Y{\vect u}'$. However, from
the derivation of the method of variation of parameters in
Section~9.4, $Y{\vect u}'={\vect f}$ as defined in the solution
of \part{a}. This and (E) imply the conclusion.

\item Since $Y{\vect u}'={\vect f}$ with ${\vect f}$ as defined in the
solution of \part{a}, $u_1,u_2,\dots,u_n$ satisfy
the conditions required in the derivation of the method of
variation of parameters in Section~9.4; hence,
$y_p=c_1y_1+c_2y_2+\cdots+c_ny_n$ is a particular solution of (A).
\end{alist}
\end{sectionSolutions}

\iftoggle{boundaryvalues}{}{\end{document}}

\chapter{Boundary Value Problems and Fourier Expansions}

\section{Eigenvalue Problems for \texorpdfstring{$y''+\lambda y=0$}{y''+𝜆y=0}}

\begin{sectionSolutions}
\item
From Theorem~11.1.2 with $L=\pi$,
$\lambda_n=n^2$,
$y_n=\sin nx$, $n=1,2,3,\dots$


\item
From Theorem~11.1.4 with $L=\pi$,
$\lambda_n=\dst\frac{(2n-1)^2}{4}$, $y_n=\dst\sin\frac{(2n-1) x}{2}$,
$n=1,2,3,\dots,$

\item
From Theorem~11.1.6 with $L=\pi$, $\lambda_0=0$,
$y_0=1$, $\lambda_n=n^2$, $y_{1n}=\cos nx$, $y_{2n}=\sin nx$,
$n=1,2,3,\dots$


\item
From Theorem~11.1.5 with $L=1$,
$\lambda_n=\dst\frac{(2n-1)^2\pi^2}{4}$, $y_n=\dst\cos\frac{(2n-1)\pi
x}{2}$, $n=1,2,3,\dots$


\item
From Theorem~11.1.6 with $L=1$, $\lambda_0=0$,
$y_0=1$, $\lambda_n=n^2\pi^2$, $y_{1n}=\cos n\pi x$, $y_{2n}=\sin n\pi
x$, $n=1,2,3,\dots$


\item
From Theorem~11.1.6 with $L=2$, $\lambda_0=0$,
$y_0=1$, $\lambda_n=\dst\frac{n^2\pi^2}{4}$, $y_{1n}=\dst\cos\frac{n\pi
x}{2}$, $y_{2n}=\dst\sin\frac{n\pi x}{2}$, $n=1,2,3,\dots$


\item
From Theorem~11.1.5 with $L=3$,
$\lambda_n=\dst\frac{(2n-1)^2\pi^2}{36}$, $y_n=\dst\cos\frac{(2n-1)\pi
x}{6}$, $n=1,2,3,\dots$


\item
From Theorem~11.1.3 with $L=5$,
$\lambda_n=\dst\frac{n^2\pi^2}{ 25}$, $y_n=\dst\cos\frac{n\pi x}{5}$,
$n=1,2,3,\dots$


\item
From Theorem~11.1.1,
 any  eigenvalues of Problem~11.1.4 must be  positive.
 If $\lambda>0$, then every solution of
$y''+\lambda y=0$ is of the form
$y=c_1\cos\sqrt\lambda\, x+c_2\sin\sqrt\lambda\, x$
where $c_1$ and $c_2$ are constants. Therefore,
$y'=\sqrt\lambda\,(-c_1\sin\sqrt\lambda\, x+c_2\cos\sqrt\lambda\, x)$.
Since $y'(0)=0$, $c_2=0$. Therefore,
$y= c_1\cos\sqrt\lambda\, x$. Since
$y(L)=0$,
$c_1\cos\sqrt\lambda\, L=0$. To make $c_1\cos\sqrt\lambda\, L=0$ with
$c_1\ne0$ we must choose
 $\sqrt\lambda=\dst\frac{(2n-1)\pi}{2L}$, where $n$ is a
positive integer. Therefore,$\lambda_n=\dst\frac{(2n-1)^2\pi^2}{4L^2}$
is an eigenvalue and
$y_n=\dst\cos\frac{(2n-1)\pi x}{2L}$ is an associated eigenfunction.

\item
If $r$ is a positive integer, then
$\dst\int_{-L}^L\cos\frac{r\pi x}{ L}\,dx
=\frac{L}{ r\pi}\sin\frac{r\pi x}{ L}\bigg|_{-L}^L=0$, so
$y_0=1$ is orthogonal to all the other eigenfunctions.
If $m$ and $n$ are distinct positive integers, then
$\dst\int_0^L\cos\frac{m\pi x}{ L}\cos\frac{n\pi x}{ L}\,dx
=\frac{1}{2}\int_{-L}^L\cos\frac{m\pi x}{ L}\cos\frac{n\pi x}{ L}\,dx=0$,
from Example~11.1.4.




\item
Let $m$ and $n$ be distinct positive integers. From the identity
$\cos A\cos B=\dst\frac{1}{2}[\cos(A-B)+\cos(A+B)]$
 with $A=(2m-1)\pi x/2L$ and $B=(2n-1)\pi x/2L$,
\[
\int_0^L\cos\frac{(2m-1)\pi x}{2L}\cos\frac{(2n-1)\pi x}{2L}\,dx
=\frac{1}{2}\int_0^L\left[\cos\frac{(m-n)\pi x}{ L}+\cos\frac{(m+n-1)\pi x}{
L}\right]\,dx=0.
\]


\item
If $y=c_1+c_2x$, then  $y'(0)=0$ implies that $c_2=0$,
so $y=c_1$. Now $\dst\int_0^Ly(x)\,dx=c_1\int_0^L \,dx=c_1L=0$ only if
$c_1=0$. Therefore,zero is not an eigenvalue.

If $y=c_1\cosh kx+c_2\sinh kx$, then $y'(0)=0$ implies that $c_2=0$, so
$y=c_1\cosh kx$. Now $\dst\int_0^Ly(x)\,dx=c_1\int_0^L\cosh kx
\,dx=c_1\frac{\sinh kL}{ k}=0$ with $k>0$ only if $c_1=0$. Therefore,
there are no negative eigenvalues.

If $y=c_1\cos kx+c_2\sin kx$, then $y'(0)=0$ implies
that $c_2=0$, so $y=c_1\cos kx$. Now
$\dst\int_0^Ly(x)\,dx=c_1\int_0^L\cos kx\,dx=c_2\frac{\sin kL}{ k}=0$ if
$k=\dst\frac{n\pi}{ L}$, where $n$ is a positive integer. Therefore,
$\lambda_n=\dst\frac{n^2\pi^2}{ L^2}$
and $y_n=\dst\cos\frac{n\pi x}{ L}$, $n=1,2,3,\dots$.


\item
If $y=c_1+c_2(x-L)$, then $y'(L)=0$ implies that
$c_2=0$, so $y=c_1$. Now $\dst\int_0^Ly(x)\,dx=c_1\int_0^L
\,dx=c_1L=0$ only if $c_1=0$. Therefore,zero is not an eigenvalue.

If $y=c_1\cosh k(x-L)+c_2\sinh k(x-L)$, then $y'(L)=0$ implies that
$c_2=0$, so $y=c_1\cosh k(x-L)$. Now
$\dst\int_0^Ly(x)\,dx=c_1\int_0^L\cosh k(x-L) \,dx=c_1\frac{\sinh kL}{
k}=0$ with $k>0$ only if $c_1=0$. Therefore,there are no negative
eigenvalues.

If $y=c_1\cos k(x-L)+c_2\sin k(x-L)$, then $y'(L)=0$
implies that $c_2=0$, so $y=c_1\cos k(x-L)$. Now
$\dst\int_0^Ly(x)\,dx=c_1\int_0^L\cos k(x-L)\,dx=c_2\frac{\sin kL}{
k}=0$ if $k=\dst\frac{n\pi}{ L}$, where $n$ is a positive integer.
Therefore,
$\lambda_n=\dst\frac{n^2\pi^2}{ L^2}$ and $y_n=\dst\cos\frac{n\pi (x-L)}{
L}$, or, equivalently, $y_n=\dst\cos\frac{n\pi x}{ L}$, $n=1,2,3,\dots$.
\end{sectionSolutions}

\section{Fourier Expansions I}

\begin{sectionSolutions}
\item
\begin{align*}
a_0&=\frac{1}{2}\int_{-1}^1(2-x)\,dx=\int_0^12\,dx=2;\\
a_n&=\int_{-1}^1(2-x)\cos n\pi x\,dx=4\int_0^1\cos n\pi
x\,dx=\frac{4}{ n\pi}\sin n\pi x\lims01=0;\\
b_n&=\int_{-1}^1(2-x)\sin n\pi x\,dx=-2\int_0^1x\sin n\pi x\,dx\\
&=\frac{2}{ n\pi}\left[x\cos n\pi x\lims01-
\int_0^1\cos n\pi x\,dx\right]\\
&=\frac{2}{ n\pi}\left[\cos n\pi-\frac{1}{ n\pi}\sin n\pi
x\lims01\right]=(-1)^n\frac{2}{ n\pi};
\end{align*}
$F(x)=\dst2+\frac{2}{\pi}\sum_{n=1}^\infty\frac{(-1)^n}{
n}\sin n\pi x$. From Theorem~11.2.4,
\[
F(x)=
\begin{cases}
2,&\phantom{-}x=-1,\\2-x,&-1<x<1,\\2,&\phantom{-}x=1.
\end{cases}
\]


\item
Since $f$ is even, $b_n=0$ for $n\ge1$;
$a_0=
\dst\int_0^1(1-3x^2)\,dx=(x-x^3)\lims01=0$;
 if $n\ge1$, then
\begin{align*}
a_n&=
2\int_0^1(1-3x^2)\cos n\pi x\,dx
=\frac{2}{ n\pi}\left[(1-3x^2)\sin n\pi x\lims01
+6\int_0^1x\sin n\pi x\,dx\right]\\
&=-\frac{12}{ n^2\pi^2}\left[x\cos n\pi x\lims01
-\int_0^1\cos n\pi x\,dx\right]\\
&=-\frac{12}{ n^2\pi^2}\left[\cos n\pi-\frac{1}{ n\pi}
\sin n\pi x\lims01\right]=(-1)^{n+1}\frac{12}{ n^2\pi^2};
\end{align*}
$F(x)=-\dst\frac{12}{\pi^2}\sum_{n=1}^\infty(-1)^n\frac{\cos n\pi x}{
n^2}$.
From Theorem~11.2.4,
$F(x)=1-3x^2$,
$-1\le x\le1$.



\item
Since $f$ is odd,
$a_n=0$ if $n\ge0$;
\begin{align*}
b_1&=\dst\frac{2}{\pi}\int_0^\pi x\cos x\sin x\,dx=
\frac{1}{\pi}\int_0^\pi x\sin2x\,dx\\
&=-\frac{1}{2\pi}\left[x\cos2x\lims0\pi-\int_0^\pi\cos2x\,dx\right]
=-\frac{1}{2\pi}\left[\pi-\frac{\sin2x}{2}\lims0\pi\right]
=-\frac{1}{2}.
\end{align*}
 if $n\ge2$, then
\begin{align*}
b_n&=\frac{2}{\pi}\int_0^\pi x\cos x\sin nx\,dx=
\frac{1}{\pi}\int_0^\pi x[\sin(n+1)x+\sin(n-1)x]\,dx\\
&=-\frac{1}{\pi}\left[x\left[\frac{\cos(n+1)x}{ n+1}
+\frac{\cos(n-1)x}{ n-1}\right]\lims0\pi-
\int_0^\pi\left[\frac{\cos(n+1)x}{
n+1}+\frac{\cos(n-1)x}{ n-1}\right]\,dx\right]\\
&=(-1)^n\left[\frac{1}{ n+1}+\frac{1}{ n-1}\right]
+\frac{1}{
\pi}\left[\frac{\sin(n+1)x}{(n+1)^2}+\frac{\sin(n-1)x}{(n-1)^2}
\right]\lims0\pi
=(-1)^n\frac{2n}{
n^2-1};
\end{align*}
 $F(x)=\dst-\frac{1}{2}\sin x+2\sum_{n=2}^\infty(-1)^n
\frac{n}{ n^2-1}\sin nx$. From Theorem~11.2.4,
$F(x)=x\cos x$, $-\pi\le x\le\pi$.




\item
Since $f$ is even,
$b_n=0$ if $n\ge1$;
$a_0=\dst\frac{1}{\pi}\int_0^\pi x\sin
x\,dx=-\frac{1}{\pi}\left[x\cos x\lims0\pi-\int_0^\pi\cos
x\,dx\right]=1+\frac{\sin x}{\pi}\lims0\pi=1$;
 $a_1=\dst\frac{2}{\pi}\int_0^\pi x\sin x\cos
x\,dx=\frac{1}{\pi}\int_0^\pi x\sin2x\,dx=-\frac{1}{2\pi}
\left[x\cos2x\lims0\pi-\int_0^\pi\cos2x\,dx\right]=
-\frac{1}{2}+\frac{\sin2x}{4\pi}\lims0\pi=-\frac{1}{2}$;
 if $n\ge2$, then
\begin{align*}
a_n&=\frac{2}{\pi}\int_0^\pi x\sin x\cos nx\,dx=
\frac{1}{\pi}\int_0^\pi x[\sin(n+1)x-\sin(n-1)x]\,dx\\
&=\frac{1}{\pi}\left[x\left[\frac{\cos(n-1)x}{ n-1}
-\frac{\cos(n+1)x}{ n+1}\right]\lims0\pi
-\int_0^\pi\left[\frac{\cos(n-1)x}{
n-1}-\frac{\cos(n+1)x}{ n+1}\right]\,dx\right]\\
&=(-1)^{n+1}\left[\frac{1}{ n-1}-\frac{1}{ n+1}\right]
-\frac{1}{\pi}\left[\frac{\sin(n-1)x}{(n-1)^2}-\frac{\sin(n+1)x}{
(n+1)^2}\right]\lims0\pi=
(-1)^{n+1}\frac{2}{ n^2-1};
\end{align*}
 $F(x)=\dst1-\frac{1}{2}\cos
x-2\sum_{n=2}^\infty
\frac{(-1)^n}{ n^2-1}\cos nx$. From Theorem~11.2.4,
$F(x)=x\sin x$, $-\pi\le x\le\pi$.



\item
Since $f$ is even,
$b_n=0$ if $n\ge0$; $a_0=\dst\int_0^{1/2}\cos\pi
x\,dx=\frac{\sin\pi x}{\pi}\lims0{1/2}=\frac{1}{\pi}$;
$a_1=\dst2\int_0^{1/2}\cos^2\pi
x\,dx=\int_0^{1/2}(1-\cos2\pi
x)\,dx=\frac{1}{2}-\frac{\sin2\pi x}{2\pi}\lims0{1/2}=\frac{1}{2}$; if
$n\ge2$, then
\begin{align*}
a_n&=\dst2\int_0^{1/2}\cos\pi x\cos n\pi
x\,dx=\int_0^{1/2}\left[\cos(n+1)x+\cos(n-1)x\right]\,dx\\
&=\frac{1}{\pi}\left[\frac{\sin(n+1)\pi x}{ n+1}+
\frac{\sin(n-1)\pi x}{ n-1}\right]\lims0{1/2}=
\frac{1}{\pi}\left[\frac{1}{ n+1}-\frac{1}{ n-1}\right]\cos\frac{n\pi}{2}\\
&=-\frac{2}{(n^2-1)\pi}\cos\frac{n\pi}{2}
=\begin{cases}
(-1)^{m+1}\dst\frac{2}{(4m^2-1)\pi}&\text{ if }n=2m,\\
0&\text{ if }n=2m+1;
\end{cases}
\end{align*}
$F(x)=\dst\frac{1}{\pi}+\frac{1}{2}\cos\pi x-\frac{2}{\pi}\sum_{n=1}^\infty
\frac{(-1)^n}{4n^2-1}\cos2n\pi x$.
 From Theorem~11.2.4,
$F(x)=f(x)$, $-1\le x\le1$.





\item
Since $f$ is odd,
$a_n=0$ if $n\ge0$;
$b_1=\dst2\int_0^{1/2}\sin^22\pi
x\,dx=\int_0^{1/2}(1-\cos4\pi
x)\,dx=\frac{1}{2}-\frac{\sin4\pi x}{4\pi}\lims0{1/2}=\frac{1}{2}$; if
$n\ge2$, then
\begin{align*}
b_n&=2\int_0^{1/2}\sin\pi x\sin n\pi x\,dx=
\int_0^{1/2}[\cos(n-1)\pi x-\cos(n+1)\pi x]\,dx\\
&=\frac{1}{\pi}\left[\frac{\sin(n-1)\pi x}{ n-1}-
\frac{\sin(n+1)\pi x}{ n+1}\right]\lims0{1/2}\\
&=-\frac{2n}{(n^2-1)\pi}\cos\frac{n\pi}{2}=
\begin{cases}
(-1)^{m+1}\dst\frac{4m}{4m^2-1}&\text{ if }n=2m,\\
0&\text{ if }n=2m+1;
\end{cases}
\end{align*}
$F(x)=\dst\frac{1}{2}\sin\pi x-\frac{4}{\pi}\sum_{n=1}^\infty
(-1)^n\frac{n}{4n^2-1}\sin2n\pi x$.
 From Theorem~11.2.4,
\[
F(x)=
\begin{cases}
\phantom{-}0,&-1\le x<\frac{1}{2},\\
-\frac{1}{2},&\phantom{-}x=-\frac{1}{2},\\
\phantom{-}\sin\pi x,&-\frac{1}{2}<x<\frac{1}{2},\\
\phantom{-}\frac{1}{2},&\phantom{-}x=\frac{1}{2},\\
\phantom{-}0,&\phantom{-}\frac{1}{2}<x\le1.
\end{cases}
\]



\item
Since $f$ is even,
$b_n=0$ if $n\ge1$;
\[
a_0=\dst\int_0^{1/2}x\sin\pi x\,dx=
-\frac{1}{\pi}\left[x\cos\pi x\lims0{1/2}-\int_0^{1/2}\cos\pi
x\,dx\right]=\frac{\sin\pi x}{\pi^2}\lims0{1/2}=
\frac{1}{\pi^2};
\]
\begin{align*}
a_1&=\dst2\int_0^{1/2}x\sin\pi x\cos\pi x\,dx=\int_0^{1/2}x\sin2x\,dx\\
&=-\frac{1}{2\pi}\left[x\cos2\pi
x\lims0{1/2}-\int_0^{1/2}\cos2\pi x\,dx\right]
=\frac{1}{4\pi}
+\frac{\sin2\pi x}{4\pi^2}\lims0{1/2}=\frac{1}{4\pi};
\end{align*}
if $n\ge2$, then
\begin{align*}
a_n&=2\int_0^{1/2}x\sin\pi x\cos n\pi x\,dx=
\int_0^{1/2}x[\sin(n+1)\pi x-\sin(n-1)\pi x]\,dx\\
&=\frac{x}{\pi}\left[\frac{\cos(n-1)\pi x}{ n-1}
-\frac{\cos(n+1)\pi x}{ n+1}\right]\lims0{1/2}-
\frac1\pi\int_0^{1/2}\left[\frac{\cos(n-1)\pi x}{ n-1}-
\frac{\cos(n+1)\pi x}{ n+1}\right]\,dx\\
%&=\frac{1}{\pi}\left[x\left[\frac{\cos(n-1)\pi x}{ n-1}
%-\frac{\cos(n+1)\pi x}{ n+1}\right]\lims0{1/2}-
%\int_0^{1/2}\left[\frac{\cos(n-1)\pi x}{ n-1}-
%\frac{\cos(n+1)\pi x}{ n+1}\right]\,dx\right]\\
&=\frac{1}{2\pi}\left[\frac{\cos(n-1)\pi/2}{ n-1}-
\frac{\cos(n+1)\pi/2}{ n+1}\right]-\frac{1}{\pi^2}
\left[\frac{\sin(n-1)\pi/2}{(n-1)^2}-\frac{\sin(n+1)\pi/2
}{(n+1)^2}\right]\\
&=\frac{1}{\pi}\frac{n}{ n^2-1}\sin\frac{n\pi}{2}
+\frac{2}{\pi^2}\frac{n^2+1}{(n^2-1)^2}\cos\frac{n\pi}{2}=
\begin{cases}
(-1)^m\dst\frac{2}{\pi^2}\frac{4m^2+1}{(4m^2-1)^2}\text{ if }n=2m,\\
(-1)^m\dst\frac{1}{4\pi}\frac{2m+1}{ m(m+1)}\text{ if }n=2m+1;
\end{cases}\\
F(x)&=\dst\frac{1}{\pi^2}+\frac{1}{4\pi}\cos\pi
x+\frac{2}{\pi^2}\sum_{n=1}^\infty
(-1)^n\frac{4n^2+1}{(4n^2-1)^2}\cos2n\pi x
\\&\qquad\qquad{}
+\frac{1}{4\pi}\sum_{n=1}^\infty(-1)^n\frac{2n+1}{ n(n+1)}\cos(2n+1)\pi x.
\end{align*}
 From
Theorem~11.2.4,
\[
F(x)=
\begin{cases}
0,&-1\le x<\frac{1}{2},\\
\frac{1}{4},&\phantom{-}x=-\frac{1}{2},\\
x\sin\pi x,&-\frac{1}{2}<x<\frac{1}{2},\\
\frac{1}{4},&\phantom{-}x=\frac{1}{2},\\
0,&\phantom{-}\frac{1}{2}<x\le1.\\
\end{cases}
\]


\item
Note that $\dst\int_{-1}^0x^2g(x)\,dx=\int_0^1x^2g(-x)\,dx$;
therefore,
$a_0=\dst\frac{1}{2}\left[\int_{-1}^0x^2\,dx
+\int_0^1(1-x^2)\,dx\right]=\frac{1}{2}\int_0^1\,dx=
\frac{1}{2}$, and
if $n\ge1$, then
\[
a_n=\int_{-1}^0x^2\cos n\pi x\,dx
+\int_0^1(1-x^2)\cos n\pi x\,dx=\int_0^1\cos n\pi x\,dx=\frac{\sin n\pi x
}{ n\pi}\lims01=0
\]
and
\begin{align*}
 b_n&=\int_{-1}^0x^2\sin n\pi x\,dx
+\int_0^1(1-x^2)\sin n\pi x\,dx=\int_0^1(1-2x^2)\sin n\pi x\,dx\\
&=-\frac{1}{ n\pi}\left[(1-2x^2)\cos n\pi x\lims01+
4\int_0^1x\cos n\pi x\,dx\right]\\
&=\frac{1+\cos n\pi}{ n\pi}-\frac{4}{ n^2\pi^2}\left[x\sin n\pi x\lims01-
\int_0^1\sin n\pi x\,dx\right]\\
&=\frac{1+\cos n\pi}{ n\pi}-\frac{4\cos n\pi x}{ n^3\pi^3}\lims01
=\frac{1+\cos n\pi}{ n\pi}+\frac{4(1-\cos n\pi)}{ n^3\pi^3}\\
&=\begin{cases}
\dst\frac{1}{ m\pi}&\text{ if }n=2m,\\
\dst\frac{8}{(2m+1)^3\pi^3}&\text{ if }n=2m+1;
\end{cases}\\
F(x)&=\dst\frac{1}{2}+\frac{1}{\pi}\sum_{n=1}^\infty\frac{1}{ n}\sin2n\pi x
+\frac{8}{\pi^3}\sum_{n=0}^\infty\frac{1}{(2n+1)^3}\sin(2n+1)\pi x.
\end{align*}
From Theorem~11.2.4,
\[
 F(x)=
\begin{cases}
\frac{1}{2},&\phantom{-}x=-1,\\
x^2,&-1< x<0,\\
\frac{1}{2},&\phantom{-}x=0,\\
1-x^2,&\phantom{-}0<x<1,\\
\frac{1}{2},&\phantom{-}x=1.
\end{cases}
\]



\item
$\dst a_0=\frac{1}{6}\left[
\int_{-3}^{-2}2\,dx+\int_{-2}^23\,dx+\int_2^31\,dx\right]=\frac{5}{2}$.
If $n\ge1$, then
\begin{align*}
a_n&=\frac{1}{3}\int_{-3}^3f(x)\cos\frac{n\pi x}{3}\,dx\\
&=\frac{1}{3}\left[\int_{-3}^{-2}2\cos\frac{n\pi x}{3}\,dx+
\int_{-2}^23\cos\frac{n\pi x}{3}\,dx+\int_2^3\cos\frac{n\pi
x}{3}\,dx\right]=
\frac{3}{ n\pi}\sin\frac{2n\pi}{3},\\
b_n&=\frac{1}{3}\left[\int_{-3}^{-2}2\sin\frac{n\pi x}{3}\,dx+
\int_{-2}^23\sin\frac{n\pi x}{3}\,dx+\int_2^3\sin\frac{n\pi
x}{3}\,dx\right]=
\frac{1}{ n\pi}\left(\cos n\pi-\cos\frac{2n\pi}{3}\right);\\
F(x)&=\dst\frac{5}{2}+\frac{3}{\pi}\sum_{n=1}^\infty
\frac{1}{ n}\sin\frac{2n\pi}{3}\cos\frac{n\pi x}{3}+\frac{1}{\pi}
\sum_{n=1}^\infty \frac{1}{ n}\left(\cos
n\pi-\cos\frac{2n\pi}{3}\right)\sin\frac{n\pi x}{3}.
\end{align*}


\item
\begin{alist}
\item $a_0=\dst\frac{1}{2\pi}\int_{-\pi}^\pi
e^x\,dx=\frac{e^x}{2\pi}\lims{-\pi}\pi=\frac{\sinh\pi}{\pi}$. If  $n\ge1$
then (A) $a_n=\dst\frac{1}{\pi}\int_{-\pi}^\pi e^x\cos nx\,dx$
and\\ (B) $b_n=\dst\frac{1}{\pi}\int_{-\pi}^\pi e^x\sin nx\,dx$.
Integrating (B) by parts yields
\[
b_n=\dst\frac{1}{\pi}\left[e^x\sin nx\lims{-\pi}\pi-n
\int_{-\pi}^\pi e^x\cos nx\,dx\right]=-na_n.
\tag{C}
\]
Integrating (A) by parts yields
\[
a_n=\dst\frac{1}{\pi}\left[e^x\cos nx\lims{-\pi}\pi+n
\int_{-\pi}^\pi e^x\sin nx\,dx\right]=(-1)^n\frac{2\sinh\pi}{\pi}+nb_n
=(-1)^n\frac{2\sinh\pi}{\pi}-n^2a_n,
\]
 from (C). Therefore,
$a_n=\dst\frac{2\sinh\pi}{\pi}\frac{(-1)^n}{ n^2+1}$. Now (C) implies that
$b_n=\dst\frac{2\sinh\pi}{\pi}\frac{(-1)^{n+1}n}{ n^2+1}$.
Therefore,
\[
F(x)=\dst\frac{\sinh\pi}{\pi}\left(
1+2\sum_{n=1}^\infty\frac{(-1)^n}{ n^2+1}\cos
nx-2\sum_{n=1}^\infty\frac{(-1)^nn}{ n^2+1}\sin nx\right).
\]

\item From Theorem~11.2.4, $F(\pi)=\cosh\pi$,
so $\dst\frac{\sinh\pi}{\pi}
\left(1+2\sum_{n=1}^\infty\frac{1}{ n^2+1}\right)=\cosh\pi$,
 which implies
the stated result.
\end{alist}



\addtocounter{enumi}{2}
\item%11.2.24
Since $f$ is even,
$b_n=0$, $n\ge1$, $a_0=\dst\frac{1}{\pi}\int_0^\pi\cos kx\,dx=
\frac{\sin kx}{ k\pi}\lims0\pi=\frac{\sin k\pi}{ k\pi}$; if $n\ge1$
then
\begin{align*}
a_n&=\dst\frac{2}{\pi}\int_0^\pi\cos kx\cos nx\,dx=
\frac{1}{\pi}\int_0^\pi[\cos(n-k)x+\cos(n+k)x]\,dx\\
&=\frac{1}{\pi}\left[\frac{\sin(n-k)x}{ n-k}+
\frac{\sin(n+k)x}{ n+k}\right]\lims0\pi\\
&=\frac{\cos n\pi\sin kx}{\pi}\left[\frac{1}{ n+k}-\frac{1}{ n-k}\right]
=(-1)^{n+1}\frac{2k\sin k\pi}{(n^2-k^2)\pi};
\end{align*}
$F(x)=\dst\frac{\sin
k\pi}{\pi}\left[\frac{1}{ k}-2k\sum_{n=1}^\infty\frac{(-1)^n}{
n^2-k^2}\cos nx\right]$.


\item
Since $f$ is continuous on $[-L,L]$ and $f(-L)=f(L)$,
Theorem~11.2.4 implies that
\[
f(x)=a_0+\sum_{n=1}^\infty\left(a_n\cos\frac{n\pi x}{
L}+b_n\sin\frac{n\pi x}{ L}\right),\quad-L\le x\le L,
\]
if
$\dst a_0=\frac{1}{ 2L}\int_{-L}^L f(x)\,dx$,  and, for $n\ge1$,
\begin{align*}
a_n&= \frac{1}{ L}\int_{-L}^L f(x)\cos\frac{n\pi
x}{ L}\,dx=
\frac{1}{ n\pi}\left[f(x)\sin\frac{n\pi x}{ L}\lims{-L}L
-\int_{-L}^Lf'(x)\sin\frac{n\pi x}{ L}\,dx\right]\\
&=\frac{L}{ n^2\pi^2}\left[f'(x)\cos\frac{n\pi x}{ L}\lims{-L}L
-\int_{-L}^Lf''(x)\cos\frac{n\pi x}{ L}\,dx\right]
=-\frac{L}{ n^2\pi^2}\int_{-L}^Lf''(x)\cos\frac{n\pi x}{ L}\,dx
\end{align*}
 (since $f'(-L)=f'(L)$), and
\begin{align*}
b_n&=\frac{1}{ L}\int_{-L}^L f(x)\sin\frac{n\pi x}{ L}\,dx=
-\frac{1}{ n\pi}\left[f(x)\cos\frac{n\pi x}{ L}\lims{-L}L
-\int_{-L}^Lf'(x)\cos\frac{n\pi x}{ L}\,dx\right]\\
&=\frac{1}{ n\pi}\int_{-L}^Lf'(x)\cos\frac{n\pi x}{ L}\,dx\text{
(since $f(-L)=f(L)$)}\\
&=\frac{L}{ n^2\pi^2}\left[f'(x)\sin\frac{n\pi x}{ L}\lims{-L}L
-\int_{-L}^Lf''(x)\sin\frac{n\pi x}{ L}\,dx\right]
=-\frac{L}{ n^2\pi^2}\int_{-L}^Lf''(x)\sin\frac{n\pi x}{ L}\,dx.
\end{align*}
If $f'''$  is integrable on $[-L,L]$, then
\begin{align*}
a_n&= -\frac{L}{ n^2\pi^2}\int_{-L}^L f''(x)\cos\frac{n\pi
x}{ L}\,dx
=-\frac{L^2}{ n^3\pi^3}\left[f'''(x)\sin\frac{n\pi x}{ L}\lims{-L}L
-\int_{-L}^Lf'''(x)\sin\frac{n\pi x}{ L}\,dx\right]\\
&=\frac{L^2}{ n^3\pi^3}\int_{-L}^Lf'''(x)\sin\frac{n\pi x}{ L}\,dx.
\end{align*}



\item
The Fourier series is
$\dst a_0+\sum_{n=1}^\infty \left(a_n\cos\frac{n\pi x}{
L}+b_n\sin\frac{n\pi x}{ L}\right)$
where
\[
a_0=\dst\frac{1}{2L}\int_{-L}^L f(x)\,dx=\frac{1}{2L}
\left[\int_{-L}^0f(x)\,dx+\int_0^Lf(x)\,dx\right].
\tag{A}
\]
Since
$\dst\int_{-L}^0f(x)\,dx=-\int_{-L}^0f(x+L)\,dx=-\int_0^Lf(x)\,dx$,
(A) implies that $a_0=0$. If  $n\ge1$, then
\begin{align*}
a_n&=\dst\frac{1}{ L}\int_{-L}^L f(x)\cos\frac{n\pi
x}{ L}\,dx\\
&=\dst\frac{1}{ L}
\left[\int_{-L}^0f(x)\cos\frac{n\pi x}{
L}\,dx+\int_0^Lf(x)\cos\frac{n\pi x}{ L}\,dx\right].
\tag{B}\\
\text{Since}\qquad
\int_{-L}^0f(x)\cos\frac{n\pi x
}{ L}\,dx&=-\int_{-L}^0f(x+L)\cos\frac{n\pi x}{
L}\,dx=-\int_0^Lf(x)\cos\frac{n\pi(x+L)}{ L}\,dx\\
&=(-1)^{n+1}\int_0^L
f(x)\cos\frac{n\pi x}{ L}\,dx,
\end{align*}
(B) implies that $a_{2n-1}=A_n$ and $a_{2n}=0$, $n\ge1$.
A similar argument shows that $b_{2n-1}=B_n$ and $b_{2n}=0$, $n\ge1$.



\item
\begin{alist}[start=2]
\item
Let $\phi_0=1$ and $\phi_{2m}=\dst\cos\frac{m\pi x}{ L}$,
$\phi_{2m-1}=\dst\sin\frac{m\pi x}{ L}$, $m\ge1$. Then $c_0=a_0$
and $c_{2m}=a_m$, $c_{2m-1}=b_m$, $m\ge1$. Since
$\dst\int_{-L}^L\phi_0^2(x)\,dx=2L$ and
$\dst\int_{-L}^L\phi_{2m-1}^2(x)\,dx=\frac{1}{2}\int_{-L}^L
\left(1-\cos\frac{2m x}{ L}\right)\,dx=L$,
$\dst\int_{-L}^L\phi_{2m}^2(x)\,dx=\frac{1}{2}\int_{-L}^L
\left(1+\cos\frac{2m x}{ L}\right)\,dx=L$, $m\ge1$,
Exercise~12.2.29\part{d} implies the conclusion.
\end{alist}
\end{sectionSolutions}


\section{Fourier Expansions II}

\begin{sectionSolutions}
\item
  $a_0=\dst\int_0^1(1-x)\,dx=-\frac{(1-x)^2}{2}\lims01=\frac{1}{2}$;
if $n\ge1$,
\begin{align*}
a_n&=2\int_0^1(1-x)\cos n\pi x\,dx=
\frac{2}{ n\pi}\left[(1-x)\sin n\pi x\lims01+
\int_0^1\sin n\pi x\,dx\right]\\
&=-\frac{2}{ n^2\pi^2}\cos n\pi x\lims01
=\frac{2}{ n^2\pi^2}[1-(-1)^n]=
\begin{cases}
\dst\frac{4}{(2m-1)^2\pi^2} &\text{ if }n=2m-1,\\
0 &\text{ if }n=2m;
\end{cases}\\
C(x)&=\dst\frac{1}{2}+\frac{4}{\pi^2}\sum_{n=1}^\infty\frac{1}{(2n-1)^2}
\cos(2n-1)\pi x.
\end{align*}


\item
$a_0=\dst\frac{1}{\pi}\int_0^\pi\sin kx\,dx=-\frac{\cos kx}{
k}\lims0\pi=\frac{1-\cos k\pi}{ k\pi}$;
if $n\ge1$, then
\begin{align*}
a_n&=\frac{2}{\pi}\int_0^\pi\sin kx\cos nx\,dx
=\frac{1}{\pi}\int_0^\pi[\sin(n+k)x-\sin(n-k)x]\,dx\\
&=\frac{1}{\pi}\left[\frac{\cos(n-k)x}{ n-k}-
\frac{\cos(n+k)x}{ n+k}\right]\lims0\pi
=\frac{1}{\pi}\left[\frac{\cos(n-k)\pi-1}{ n-k}-
\frac{\cos(n+k)\pi-1}{ n+k}\right]\\
&=-\frac{2k[1-(-1)^n\cos k\pi]}{(n^2-k^2)\pi};\\
C(x)&=\dst\frac{1-\cos k\pi}{
k\pi}-\frac{2k}{\pi}\sum_{n=1}^\infty
\frac{[1-(-1)^n\cos k\pi]}{ n^2-k^2}\cos nx.
\end{align*}



\item
 $a_0=\dst\frac{1}{ L}\int_0^L(x^2-L^2)\,dx=\frac{1}{
L}\left(\frac{x^3}{3}-L^2x\right)\lims0L=-\frac{2L^2}{3}$;
if $n\ge1$,
\begin{align*}
a_n&=\frac{2}{ L}\int_0^L(x^2-L^2)\cos\frac{n\pi x}{ L}\,dx=
\frac{2}{ n\pi}\left[(x^2-L^2)\sin\frac{n\pi x}{ L}\lims0L
-2\int_0^Lx\sin\frac{n\pi x}{ L}\,dx\right]\\
&=\frac{4L}{ n^2\pi^2}\left[x\cos\frac{n\pi x}{ L}\lims0L
-\int_0^L\cos\frac{n\pi x}{ L}\,dx\right]
\\&
=(-1)^n\frac{4L^2}{ n^2\pi^2}-\frac{4L^2}{ n^3\pi^3}\sin\frac{n\pi x}{
L}\lims0L=
(-1)^n\frac{4L^2}{ n^2\pi^2};\\
C(x)&=\dst-\frac{2L^2}{3}+\frac{4L^2}{\pi^2}\sum_{n=1}^\infty\frac{(-1)^n}{
n^2}\cos\frac{n\pi x}{ L}.
\end{align*}




\item
  $a_0=\dst\int_0^\pi e^x\,dx=e^x\lims0\pi=\frac{e^\pi-1}{\pi}$;
if $n\ge1$, then
\begin{align*}
a_n&=\frac{2}{\pi}\int_0^\pi e^x\cos nx\,dx
=\frac{2}{\pi}\left[e^x\cos nx\lims0\pi-n\int_0^\pi e^x\sin
nx\,dx\right]\\
&=\frac{2}{\pi}\left[(-1)^ne^\pi-1-ne^x\sin nx\lims0\pi-
n^2\int_0^\pi e^x\cos nx\,dx\right]
=\frac{2}{\pi}[(-1)^ne^\pi-1]-n^2a_n;
\end{align*}
$(1+n^2)a_n=\dst\frac{2}{\pi}[(-1)^ne^\pi-1]$;
$a_n=\dst\frac{2}{(n^2+1)\pi}[(-1)^ne^\pi-1]$;
\[
C(x)=\dst\frac{e^\pi-1}{\pi}+\frac{2}{\pi}\sum_{n=1}^\infty
\frac{[(-1)^ne^\pi-1]}{(n^2+1)}\cos nx.
\]

\item
$a_0=\dst\frac{1}{ L}\int_0^L(x^2-2Lx)\,dx=\frac{1}{
L}\left(\frac{x^3}{3}-Lx^2
\right)\lims0L=-\frac{2L^2}{3}$;  if $n\ge1$,
\begin{align*}
a_n&=\frac{2}{ L}\int_0^L(x^2-2Lx)\cos\frac{n\pi x}{ L}\,dx
=\frac{2}{ n\pi}\left[(x^2-2Lx)\sin\frac{n\pi x}{ L}\lims0L-
2\int_0^L
(x-L)\sin\frac{n\pi x}{ L}\,dx\right]\\
&=\frac{4L}{ n^2\pi^2}\left[(x-L)\cos\frac{n\pi x}{ L}\lims0L-
\int_0^L\cos\frac{n\pi x}{ L}\,dx\right]=
\frac{4L^2}{ n^2\pi^2}-\frac{4L^3}{ n^3\pi^3}\sin\frac{n\pi x}{ L}\lims0L
=\frac{4L^2}{ n^2\pi^2};\\
C(x)&=-\frac{2L^2}{3}+\frac{4L^2}{\pi^2}\sum_{n=1}^\infty\frac{1}{
n^2}\cos\frac{n\pi x}{ L}.
\end{align*}


\item
 $b_n=\dst2\int_0^1(1-x)\sin n\pi x\,dx=
-\frac{2}{ n\pi}\left[(1-x)\cos n\pi x\lims01+\int_0^1\cos
n\pi x\,dx\right]=\frac{2}{ n\pi}+\frac{2}{ n^2\pi^2}\sin n \pi x\lims01
=\frac{2}{ n\pi}$;
$S(x)=\dst\frac{2}{\pi}\sum_{n=1}^\infty\frac{1}{ n}
\sin n\pi x$.

\item
$b_n=\dst\frac{2}{ L}\int_0^{L/2}\sin\frac{n\pi x}{ L}\,dx=
-\frac{2}{ n\pi}\cos\frac{n\pi x}{ L}\lims0{L/2}=
\frac{2}{ n\pi}\left[1-\cos\frac{n\pi}{2}\right]$;
\[
S(x)=\dst\frac{2}{\pi}\sum_{n=1}^\infty\frac{1}{ n}
\left[1-\cos\frac{n\pi}{2}\right]\sin\frac{n\pi x}{ L}.
\]


\item
\begin{align*}
b_1&=\frac{2}{\pi}\int_0^\pi x\sin^2 x\,dx=
\frac{1}{ \pi}\int_0^\pi x(1-\cos2x)\,dx
=\frac{x^2}{2\pi}\lims0\pi
-\frac{1}{\pi}\int_0^\pi x\cos2x\,dx\\
&=\frac{\pi}{2}-\frac{1}{2\pi}\left[x\sin2x\lims0\pi-\int_0^\pi
\cos2x\,dx\right]
=\frac{\pi}{2}+\frac{\sin2x}{4\pi}\lims0\pi=\frac{\pi}{2};
\end{align*}
if $n\ge2$, then
\begin{align*}
b_n&=
\frac{2}{\pi}\int_0^\pi x\sin x\sin nx\,dx
=\frac{1}{\pi}\int_0^\pi x[\cos(n-1)x-\cos(n+1)x]\,dx\\
&=\frac{1}{\pi}\left[x\left[\frac{\sin(n-1)x}{ n-1}-
\frac{\sin(n+1)x}{ n+1}\right]\lims0\pi-\int_0^\pi
\left[\frac{\sin(n-1)x}{ n-1}-\frac{\sin(n+1)x}{ n+1}\right]\,dx\right]\\
&=\frac{1}{\pi}\left[\frac{\cos(n-1)x }{(n-1)^2}
-\frac{\cos(n+1)x}{(n+1)^2}\right]\lims0\pi
=\frac{1}{\pi}\left[\frac{1}{(n-1)^2}-\frac{1}{(n+1)^2}\right]
\left[(-1)^{n+1}-1\right]\\
&=\frac{4n}{(n^2-1)^2\pi}\left[(-1)^{n+1}-1\right]=
\begin{cases}
0& \text{ if }n=2m-1,\\
-\dst\frac{16m}{(4m^2-1)\pi}&\text{ if }n=2m;
\end{cases}\\
S(x)&=\dst\frac{\pi}{2}\sin
x-\frac{16}{\pi}\sum_{n=1}^\infty\frac{n}{(4n^2-1)^2}\sin2nx.
\end{align*}


\item
$c_n=\dst\frac{2}{ L}\int_0^L\cos\frac{(2n-1)\pi x}{2L}\,dx=
\frac{4}{(2n-1)\pi}\sin\frac{(2n-1)\pi x}{ 2L}\lims0L=
(-1)^{n+1}\frac{4}{(2n-1)\pi}$;
\[
C_M(x)=\dst-\frac{4}{\pi}\sum_{n=1}^\infty\frac{(-1)^n}{2n-1}
\cos\frac{(2n-1)\pi x}{2L}.
\]



\item
\begin{align*}
d_n&=2\int_0^1x\cos\frac{(2n-1)\pi x}{2}\,dx\\
&=\frac{4}{(2n-1)\pi}\left[x\sin\frac{(2n-1)\pi x}{2}\lims01
-\int_0^1\sin\frac{(2n-1)\pi x}{2}\,dx\right]\\
&=\frac{4}{(2n-1)\pi}\left[(-1)^{n+1}
+\frac{2}{(2n-1)\pi}\cos\frac{(2n-1)\pi x}{2}\lims01\right]\\
&=-\frac{4}{(2n-1)\pi}\left[(-1)^n+\frac{2}{(2n-1)\pi}\right];\\
C_M(x)&=
-\dst\frac{4}{\pi}\sum_{n=1}^\infty\left[(-1)^n+\frac{2}{(2n-1)\pi}\right]
\cos\frac{(2n-1)\pi x}{2}.
\end{align*}



\item
\begin{align*}
 c_n&=\frac{2}{\pi}\int_0^\pi\cos x\cos\frac{(2n-1) x}{2}\,dx
=\frac{1}{\pi}\int_0^\pi\left[\frac{\cos(2n+1)x}{2}+\frac{\cos(2n-3)x}{2}\right]\,dx\\
&=\frac{2}{\pi}\left[\frac{\sin(2n+1)x/2}{2n+1}+\frac{\sin(2n-3)x/2}{2n-3}\right]\lims0\pi\\
&=(-1)^n\frac{2}{\pi}\left[\frac{1}{2n+1}+\frac{1}{2n-3}\right]
=(-1)^n\frac{4(2n-1)}{(2n-3)(2n+1)\pi};\\
C_M(x)&=
\dst\frac{4}{\pi}\sum_{n=1}^\infty(-1)^n\frac{2n-1}{(2n-3)(2n+1)}
\cos\frac{(2n-1) x}{2}.
\end{align*}


\item
\begin{align*}
c_n&=\frac{2}{ L}\int_0^L(Lx-x^2)\cos\frac{(2n-1)\pi x}{2L}\,dx\\
&=\frac{4}{(2n-1)\pi}\left[(Lx-x^2)\sin\frac{(2n-1)\pi
x}{2L}\lims0L-
\int_0^L
(L-2x)\sin\frac{(2n-1)\pi x}{2L}\,dx\right]\\
&=\frac{8L}{(2n-1)^2\pi^2}\left[(L-2x)\cos\frac{(2n-1)\pi
x}{2L}\lims0L+2
\int_0^L\cos\frac{(2n-1)\pi x}{2L}\,dx\right]\\
&=-\frac{8L^2}{(2n-1)^2\pi^2}+
\frac{32L^2}{(2n-1)^3\pi^3}\sin\frac{(2n-1)\pi x}{2L}\lims0L
=\frac{32L^2}{(2n-1)^3\pi^3}\sin\frac{(2n-1)\pi x}{2}\\
&=-\frac{8L^2}{(2n-1)^2\pi^2}
+(-1)^{n-1}\frac{32L^2}{(2n-1)^3\pi^3};\\
C_M(x)&=-\dst\frac{8L^2}{\pi^2}\sum_{n=1}^\infty\frac{1}{(2n-1)^2}\left[
1+\frac{4(-1)^n}{(2n-1)\pi}\right]\cos\frac{(2n-1)\pi
x}{2L}.
\end{align*}

\item
\begin{align*}
 d_n&=\frac{2}{ L}\int_0^Lx^2\sin\frac{(2n-1)\pi x}{2L}\,dx\\
&=-\frac{4}{(2n-1)\pi}\left[x^2\cos\frac{(2n-1)\pi x}{ 2L}\lims0L
-2\int_0^Lx\cos\frac{(2n-1)\pi x}{ 2L}\,dx\right]\\
&=\frac{16L}{(2n-1)^2\pi^2}\left[x\sin\frac{(2n-1)\pi x}{ 2L}\lims0L-
\int_0^L\sin\frac{(2n-1)\pi x}{ 2L}\,dx\right]\\
&=(-1)^{n+1}\frac{16L^2}{(2n-1)^2\pi^2}+\frac{32L^2}{(2n-1)^3\pi^3}
\cos\frac{(2n-1)\pi x}{ 2L}\lims0L\\
&=(-1)^{n+1}\frac{16L^2}{(2n-1)^2\pi^2}-\frac{32L^2}{(2n-1)^3\pi^3};\\
S_M(x)&=\dst-\frac{16L^2}{\pi^2}\sum_{n=1}^\infty\frac{1}{(2n-1)^2}\left[
(-1)^n+\frac{2}{(2n-1)\pi}\right]\sin\frac{(2n-1)\pi x}{2L}.
\end{align*}


\item
\begin{align*}
d_n&=\frac{2}{\pi}\int_0^\pi\cos x\sin\frac{(2n-1)x}{2}\,dx
=\frac{1}{\pi}\int_0^\pi\left[\frac{\sin(2n+1)x}{2}+
\frac{\sin(2n-3)x}{2}\right]\,dx\\
&=-\frac{2}{\pi}\left[\frac{\cos(2n+1)x/2}{2n-1}+\frac{\cos(2n-3)x/2}{2n-3}\right]\lims0\pi\\
&=\frac{2}{\pi}\left[\frac{1}{2n+1}+\frac{1}{2n-3}\right]
=\frac{4(2n-1)}{(2n-3)(2n+1)\pi};\\
S_M(x)&=\dst\frac{4}{\pi}\sum_{n=1}^\infty\frac{2n-1}{(2n-3)(2n+1)}\sin\frac{(2n-1)
x}{2}.
\end{align*}


\item
\begin{align*}
d_n&=\frac{2}{ L}\int_0^L(Lx-x^2)\sin\frac{(2n-1)\pi x}{2L}\,dx\\
&=-\frac{4}{(2n-1)\pi}\left[(Lx-x^2)\cos\frac{(2n-1)\pi
x}{2L}\lims0L-
\int_0^L
(L-2x)\cos\frac{(2n-1)\pi x}{2L}\,dx\right]\\
&=\frac{8L}{(2n-1)^2\pi^2}\left[(L-2x)\sin\frac{(2n-1)\pi
x}{2L}\lims0L+2
\int_0^L\sin\frac{(2n-1)\pi x}{2L}\,dx\right]\\
&=(-1)^n\frac{8L^2}{(2n-1)^2\pi^2}-
\frac{32L^2}{(2n-1)^3\pi^3}\cos\frac{(2n-1)\pi x}{2L}\lims0L\\
&=(-1)^n\frac{8L^2}{(2n-1)^2\pi^2}
+\frac{32L^2}{(2n-1)^3\pi^3};\\
S_M(x)&=\dst\frac{8L^2}{\pi^2}\sum_{n=1}^\infty\frac{1}{(2n-1)^2}\left[
(-1)^n+\frac{4}{(2n-1)\pi}\right]\sin\frac{(2n-1)\pi
x}{2L}.
\end{align*}



\item
$\dst a_0=\frac{1}{ L}\int_0^L(3x^4-4Lx^3)\,dx=\frac{1}{
L}\left(\frac{3x^5}{5}-Lx^4
\right)\lims0L=-\frac{2L^4}{5}$.
Since $f'(0)=f'(L)=0$ and  $f'''(x)=24(3x-L)$,
\begin{align*}
a_n&=
\frac{48L^2}{ n^3\pi^3}\int_0^L(3x-L)\sin\frac{n\pi x}{ L}\,dx
=-\frac{48L^3}{ n^4\pi^4}\left[(3x-L)\cos\frac{n\pi x}{ L}\lims0L
-3\int_0^L\cos\frac{n\pi x}{ L}\,dx\right]\\
&=-\frac{48L^3}{ n^4\pi^4}\left[(-1)^n2L+L\right]
+\frac{144L^4}{ n^5\pi^5}\sin\frac{n\pi x}{ L}\lims0L=
-\frac{48L^4}{ n^4\pi^4}\left[1+(-1)^n2\right],\quad n\ge1;\\
C(x)&=-\frac{2L^4}{5}-\frac{48L^4}{\pi^4}\sum_{n=1}^\infty\frac{1+(-1)^n2}{
n^4}\cos\frac{n\pi x}{ L}.
\end{align*}


\item
$\dst a_0=\frac{1}{ L}\int_0^L(x^4-2Lx^3+L^2x^2)\,dx=\frac{1}{
L}\left(\frac{x^5}{5}-\frac{Lx^4}{2}+\frac{L^2x^3}{3}
\right)\lims0L=\frac{L^4}{30}$.
Since  $f'(0)=f'(L)=0$ and  $f'''(x)=12(2x-L)$,
\begin{align*}
a_n&=
\frac{24L^2}{ n^3\pi^3}\int_0^L(2x-L)\sin\frac{n\pi x}{ L}\,dx
=-\frac{24L^3}{ n^4\pi^4}\left[(2x-L)\cos\frac{n\pi x}{ L}\lims0L
-2\int_0^L\cos\frac{n\pi x}{ L}\,dx\right]\\
&=-\frac{24L^3}{ n^4\pi^4}\left[(-1)^nL+L\right]+
\frac{48L^4}{ n^5\pi^5}\sin\frac{n\pi x}{ L}\lims0L=
-\frac{24L^4}{ n^4\pi^4}\left[1+(-1)^n\right]\\
&=
\begin{cases}
0&\text{ if }n=2m-1,\\
\dst-\frac{3L^4}{ m^4\pi^4}&\text{ if }n=2m,
\end{cases}
n\ge1.\\
C(x)&=\frac{L^4}{30}-\frac{3L^4}{\pi^4}\sum_{n=1}^\infty\frac{1}{
n^4}\cos\frac{2n\pi x}{ L}.
\end{align*}



\item
Since $f(0)=f(L)=0$ and $f''(x)=-2$,
\begin{align*}
b_n&=
\frac{4L}{ n^2\pi^2}
\int_0^L\sin\frac{n\pi x}{ L}\,dx
=-\frac{4L^2}{ n^3\pi^3}\cos\frac{n\pi x}{ L}\lims0L=
-\frac{4L^2}{ n^3\pi^2}(\cos n\pi-1)\\
&=\begin{cases}
\dst\frac{8L^2}{(2m-1)^3\pi^3},&\text{if $n=2m-1$},\\
0,&\text{if  $n=2m$};\\
\end{cases}\\
S(x)&=\frac{8L^2}{\pi^3}\sum_{n=1}^\infty\frac{1}{(2n-1)^3}
\sin\frac{(2n-1)\pi x}{ L}.
\end{align*}



\item
Since $f(0)=f(L)=0$ and $f''(x)=-6x$,
\begin{gather*}
b_n=
\frac{12L}{ n^2\pi^2}\int_0^Lx\sin\frac{n\pi x}{ L}\,dx
=-\frac{12L^2}{ n^3\pi^3}\left[x\cos\frac{n\pi x}{ L}\lims0L
-\int_0^L\cos\frac{n\pi x}{ L}\,dx\right]
=(-1)^{n+1}\frac{12L^3}{ n^3\pi^3};
\\
S(x)=\dst-\frac{12L^3}{\pi^3}\sum_{n=1}^\infty\frac{(-1)^n}{ n^3}
\sin\frac{n\pi x}{ L}.
\end{gather*}




\item
Since $f(0)=f(L)=f''(0)=f''(L)=0$ and $f^{(4)}=360x$,
\begin{align*}
b_n&=
\frac{720L^3}{ n^4\pi^4}\int_0^Lx\sin\frac{n\pi x}{ L}\,dx
=-\frac{720L^4}{ n^5\pi^5}\left[x\cos\frac{n\pi x}{
L}\lims0L-\int_0^L\cos\frac{n\pi x}{
L}\,dx\right]\\
&=(-1)^{n+1}\frac{720L^5}{ n^5\pi^5}+\frac{720L^5}{ n^6\pi^6}\sin\frac{n\pi
x}{ L}\lims0L=
(-1)^{n+1}\frac{720L^5}{ n^5\pi^5};\\
S(x)&=-\frac{720L^5}{\pi^5}\sum_{n=1}^\infty\frac{(-1)^n}{ n^5}\sin\frac{n\pi
x}{ L}.
\end{align*}




\item
\begin{alist}
\item
Since $f$ is continuous on $[0,L]$ and $f(L)=0$,
Theorem~11.3.3 implies that\\
$f(x)=\dst\sum_{n=1}^\infty c_n\cos\frac{(2n-1)\pi x}{2L}$, $-L\le
x\le L$, with
\begin{align*}
c_n&=\frac{2}{ L}\int_0^Lf(x)\cos\frac{(2n-1)\pi x}{2L}\,dx\\
&=\frac{4}{(2n-1)\pi}\left[f(x)\sin\frac{(2n-1)\pi x}{2L}\lims0L
-\int_0^Lf'(x)\sin\frac{(2n-1)\pi x}{2L}\,dx\right]\\
&=-\frac{4}{(2n-1)\pi}\int_0^Lf'(x)\sin\frac{(2n-1)\pi x}{2L}\,dx
\text{ (since $f(L)=0$)}\\
&=\frac{8L^2}{(2n-1)^2\pi^2}\left[f'(x)\cos\frac{(2n-1)\pi x}{2L}\lims0L
-\int_0^Lf''(x)\cos\frac{(2n-1)\pi x}{2L}\,dx\right]\\
&=-\frac{8L}{(2n-1)^2\pi^2}\int_0^Lf''(x)\cos\frac{(2n-1)\pi x}{2L}\,dx
\text{ (since $f'(0)=0$)}.
\end{align*}


\item Continuing the integration by parts yields
\begin{align*}
c_n&=-\frac{16L^2}{(2n-1)^3\pi^3}\left[f''(x)\sin\frac{(2n-1)\pi
x}{2L}\lims0L-\int_0^Lf'''(x)\sin\frac{(2n-1)\pi x}{2L}\,dx\right]\\
&=\frac{16L^2}{(2n-1)^3\pi^3}
\int_0^Lf'''(x)\sin\frac{(2n-1)\pi x}{2L}\,dx.
\end{align*}
\end{alist}




\item
Since $f'(0)=f(L)=0$ and $f''(x)=-2$,
\begin{align*}
c_n&=
\frac{16L}{(2n-1)^2\pi^2}
\int_0^L\cos\frac{(2n-1)\pi x}{2L}\,dx \\
&=\frac{32L^2}{(2n-1)^3\pi^3}\sin\frac{(2n-1)\pi x}{2L}\lims0L
=(-1)^{n+1}\frac{32L^2}{(2n-1)^3\pi^3};\\
C_M(x)&=-
\frac{32L^2}{\pi^3}\sum_{n=1}^\infty\frac{(-1)^n}{
(2n-1)^3}\cos\frac{(2n-1)\pi x}{2L}.
\end{align*}


\item
Since $f'(0)=f(L)=0$ and $f''(x)=6(2x+L)$,
\begin{align*}
c_n&=
-\frac{48L}{(2n-1)^2\pi^2}
\int_0^L(2x+L)\cos\frac{(2n-1)\pi x}{2L}\,dx\\
&=-\frac{96L^2}{(2n-1)^3\pi^3}\left[(2x+L)\sin\frac{(2n-1)\pi x}{2L}
\lims0L-2\int_0^L\sin\frac{(2n-1)\pi x}{2L}\right]\,dx\\
&=-\frac{96L^2}{(2n-1)^3\pi^3}
\left[(-1)^{n+1}3L-\frac{4L}{(2n-1)\pi}\cos\frac{(2n-1)\pi x}{2L}
\lims0L\right] \\
&=\frac{96L^3}{(2n-1)^3\pi^3}
\left[(-1)^n3+\frac{4}{(2n-1)\pi}\right];\\
C_M(x)&=\dst\frac{96L^3}{\pi^3}\sum_{n=1}^\infty\frac{1}{(2n-1)^3}\left[
(-1)^n3+\frac{4}{(2n-1)\pi}\right]\cos\frac{(2n-1)\pi
x}{2L}.
\end{align*}


\item
Since $f'(0)=f(L)=f''(L)=0$ and $f'''(x)=12(2x-L)$,
\begin{align*}
c_n&=
\frac{192L^2}{(2n-1)^3\pi^3}
\int_0^L(2x-L)\sin\frac{(2n-1)\pi x}{2L}\,dx\\
&=-\frac{384L^3}{(2n-1)^4\pi^4}\left[(2x-L)\cos\frac{(2n-1)\pi x}{2L}
\lims0L-2\int_0^L\cos\frac{(2n-1)\pi x}{2L}\right]\,dx\\
&=-\frac{384L^3}{(2n-1)^4\pi^4}
\left[L-\frac{4L}{(2n-1)\pi}\sin\frac{(2n-1)\pi x}{2L}
\lims0L\right] \\
&=-\frac{384L^4}{(2n-1)^4\pi^4}
\left[1+\frac{(-1)^n4}{(2n-1)\pi}\right];\\
C_M(x)&=-\dst\frac{384L^4}{\pi^4}\sum_{n=1}^\infty\frac{1}{(2n-1)^4}\left[
1+\frac{(-1)^n4}{(2n-1)\pi}\right]\cos\frac{(2n-1)\pi
x}{2L}.
\end{align*}


\item
\begin{alist}
\item
Since $f$ is continuous on $[0,L]$ and $f(0)=0$,
Theorem~11.3.4 implies that\\
$f(x)=\dst\sum_{n=1}^\infty d_n\sin\frac{(2n-1)\pi x}{2L}$, $-L\le
x\le L$,  with
\begin{align*}
d_n&=\frac{2}{ L}\int_0^Lf(x)\sin\frac{(2n-1)\pi x}{2L}\,dx\\
&=-\frac{4}{(2n-1)\pi}\left[f(x)\cos\frac{(2n-1)\pi x}{2L}\lims0L
-\int_0^Lf'(x)\cos\frac{(2n-1)\pi x}{2L}\,dx\right]\\
&=\frac{4}{(2n-1)\pi}\int_0^Lf'(x)\cos\frac{(2n-1)\pi x}{2L}\,dx
\text{ (since $f(0)=0$)}\\
&=\frac{8L}{(2n-1)^2\pi^2}\left[f'(x)\sin\frac{(2n-1)\pi x}{2L}\lims0L
-\int_0^Lf''(x)\sin\frac{(2n-1)\pi x}{2L}\,dx\right]\\
&=-\frac{8L}{(2n-1)^2\pi^2}\int_0^Lf''(x)\sin\frac{(2n-1)\pi x}{2L}\,dx
\text{ since $f'(L)=0$}.
\end{align*}

\item Continuing the integration by parts yields
\begin{align*}
d_n&=\frac{16L^2}{(2n-1)^3\pi^3}\left[f''(x)\cos\frac{(2n-1)\pi
x}{2L}\lims0L-\int_0^Lf'''(x)\cos\frac{(2n-1)\pi x}{2L}\,dx\right]\\
&=-\frac{16L^2}{(2n-1)^3\pi^3}
\int_0^Lf'''(x)\cos\frac{(2n-1)\pi x}{2L}\,dx.
\end{align*}
\end{alist}

\item
Since $f(0)=f'(L)=0$, and $f''(x)=6(L-2x)$
\begin{align*}
d_n&=-\dst\frac{48L}{(2n-1)^2\pi^2}\int_0^L(L-2x)
\sin\frac{(2n-1)\pi x}{2L}\,dx\\
&=\dst\frac{96L^2}{(2n-1)^3\pi^3}\left[(L-2x)\cos\frac{(2n-1)\pi
x}{2L}\lims0L+2\int_0^L\cos\frac{(2n-1)\pi x}{2L}\,dx\right]\\
&=\dst\frac{96L^2}{(2n-1)^3\pi^3}\left[-L+\frac{4L}{(2n-1)\pi}
\sin\frac{(2n-1)\pi x}{2L}\lims0L\right]\\
&=\dst-\frac{96L^3}{(2n-1)^3\pi^3}\left[1+(-1)^n\frac{4}{(2n-1)\pi}\right];\\
S_M(x)
&=\dst-\frac{96L^3}{\pi^3}\sum_{n=1}^\infty
\frac{1}{(2n-1)^3}\left[1+(-1)^n\frac{4}{(2n-1)\pi}\right]\sin\frac{(2n-1)\pi x
}{2L}.
\end{align*}



\item
Since $f(0)=f'(L)=f''(0)=0$ and $f'''(x)=6$,
\begin{align*}
d_n&=-\frac{96L^2}{(2n-1)^3\pi^3}\int_0^L
\cos\frac{(2n-1)\pi x}{2L}\,dx\\
&=-\frac{192L^3}{(2n-1)^4\pi^4}
\sin\frac{(2n-1)\pi x}{2L}\lims0L=
(-1)^n\frac{192L^3}{(2n-1)^4\pi^4};\\
S_M(x)&=\frac{192L^3}{\pi^4}\sum_{n=1}^\infty\frac{(-1)^n}{(2n-1)^4}
\sin\frac{(2n-1)\pi x}{2L}.
\end{align*}


\item
Since $f(0)=f'(L)=f''(0)=0$ and $f'''(x)=12(2x-L)$,
\begin{align*}
d_n&=-\frac{192L^2}{(2n-1)^3\pi^3}\int_0^L(2x-L)
\cos\frac{(2n-1)\pi x}{2L}\,dx\\
&=-\frac{384L^3}{(2n-1)^4\pi^4}
\left[(2x-L)\sin\frac{(2n-1)\pi x}{2L}\lims0L
-2\int_0^L\sin\frac{(2n-1)\pi x}{2L}\,dx\right]\\
&=-\frac{384L^3}{(2n-1)^4\pi^4}
\left[(-1)^{n+1}L+\frac{4L}{(2n-1)\pi}
\cos\frac{(2n-1)\pi x}{2L}\lims0L\right]\\
&=-\frac{384L^3}{(2n-1)^4\pi^4}
\left[(-1)^{n+1}L-\frac{4L}{(2n-1)\pi}\right]\\
&=\frac{384L^4}{(2n-1)^4\pi^4}
\left[(-1)^n+\frac{4}{(2n-1)\pi}\right];\\
S_M(x)&=\frac{384L^4}{\pi^4}\sum_{n=1}^\infty\frac{1}{(2n-1)^4}
\left[(-1)^n+\frac{4}{(2n-1)\pi}\right]\sin\frac{(2n-1)\pi x}{2L}.
\end{align*}



\item
The Fourier sine series of $f_4$ on $[0,2L]$ is $\dst
\sum_{n=1}^\infty B_n\sin\frac{n\pi x}{2L}$, where
\[
B_n=\dst\frac{1}{ L}\int_0^{2L}f_4(x)\sin\frac{n\pi x}{2L}\,dx
=\frac{1}{
L}\left[\int_0^Lf(x)\sin\frac{n\pi x}{2L}\,dx+\int_L^{2L}f(2L-x)
\sin\frac{n\pi x}{2L}\,dx\right].
\]
Replacing $x$ by $2L-x$ yields $\dst\int_L^{2L}f(2L-x)\sin\frac{n\pi
x}{2L}\,dx= \int_0^Lf(x)\sin\frac{n\pi(2L-x)}{2L}\,dx$. Since
$\dst\sin\frac{n\pi(2L-x)}{2L}=(-1)^{n+1} \sin\frac{n\pi x}{2L}$,
\[
\int_L^{2L}f(2L-x)\sin\frac{n\pi
x}{2L}\,dx=(-1)^{n+1}\int_0^Lf(x)\sin\frac{n\pi x}{2L}\,dx,
\]
so
\[
B_n=\frac{1+(-1)^{n+1}}{ L}\int_0^Lf(x)\sin\frac{n\pi x}{2L}\,dx=
\begin{cases}
\dst\frac{2}{ L}\int_0^Lf(x)\sin\frac{(2m-1)\pi x}{2L}\,dx&\text{if
$n=2m-1$},\\
0&\text{if $n=2m$}.
\end{cases}
\]
Therefore,the Fourier sine series of $f_4$ on $[0,2L]$ is
$\dst\sum_{n=1}^\infty d_n\sin\frac{(2n-1)\pi x}{2L}$ with
\[
\dst d_n=\frac{2}{ L}\int_0^Lf(x)\sin\frac{(2n-1)\pi
x}{2L}\,dx.
\]




\item
The Fourier cosine series of $f_4$ on $[0,2L]$ is $\dst
A_0+\sum_{n=1}^\infty A_n\cos\frac{n\pi x}{2L}$, where
\[
A_0=\dst\frac{1}{2L}\int_0^{2L}f_4(x)\,dx
=\frac{1}{2L}\left[\int_0^Lf(x)\,dx+\int_L^{2L}f(2L-x)\,dx\right]
=\frac{1}{ L}\int_0^L f(x)\,dx
\]
 and
\[
A_n=\dst\frac{1}{ L}\int_0^{2L}f_4(x)\cos\frac{n\pi x}{2L}\,dx
=\frac{1}{
L}\left[\int_0^Lf(x)\cos\frac{n\pi x}{2L}\,dx+\int_L^{2L}f(2L-x)
\cos\frac{n\pi x}{2L}\,dx\right].
\]
Replacing $x$ by $2L-x$   yields
\[
\int_L^{2L}f(2L-x)\cos\frac{n\pi x}{2L}\,dx=-\int_L^0f(x)
\cos\frac{n\pi(2L-x)}{2L}\,dx
=\int_0^Lf(x)\cos\frac{n\pi(2L-x)}{2L}\,dx.
\]
Since
$\dst\cos\frac{n\pi(2L-x)}{2L}=\cos n\pi\cos\frac{n\pi x}{2L}=(-1)^n
\cos\frac{n\pi x}{2L}$,
\[
A_n=\frac{1+(-1)^n}{ L}\int_0^Lf(x)\cos\frac{n\pi x}{2L}\,dx
=\begin{cases}
0&\text{if $n=2m-1$}\\
\dst\frac{2}{ L}\int_0^Lf(x)\cos\frac{m\pi x}{ L}\,dx&\text{if
$n=2m$}.
\end{cases}
\]
Therefore,the Fourier cosine series of $f_4$ on $[0,2L]$
is $A_0+\dst\sum_{n=0}^\infty A_{2n}cos\frac{n\pi x}{ L}$=
$a_0+\dst\sum_{n=0}^\infty a_n\cos\frac{n\pi x}{ L}$.
\end{sectionSolutions}

\chapter{Fourier Solutions of Partial Differential}

\section{The Heat Equation}

\begin{sectionSolutions}
\item
 $X(x)T(t)$  satisfies
$u_t=a^2u_{xx}$ if $X''+\lambda X=0$ and (A) $T'=-a^2\lambda T$
 for the same value of $\lambda$.  The product also
satisfies  the
 boundary conditions $u(0,t)= u_x(L,t)=0, \ t>0$,  if and
only if
$X(0)=
X'(L)=0$.
Since we are interested in nontrivial solutions, $X$ must be a
nontrivial solution of
(B) $X''+\lambda X=0,\ X(0)=0,\ X'(L)=0$.
 From Theorem~11.1.4,
$\lambda_n=(2n-1)^2\pi^2/4L^2$ is an eigenvalue of (B) with
associated eigenfunction $X_n=\dst\sin\frac{(2n-1)\pi x}{2L}$,
$n=1,2,3,\dots$. Substituting $\lambda=(2n-1)^2\pi^2/4L^2$  into
(A) yields
$T'=-((2n-1)^2\pi^2a^2/4L^2)T$,
which has the solution $T_n=e^{-(2n-1)^2\pi^2a^2t/4L^2}$.

We have now shown that the functions
$\dst u_n(x,t)=e^{-(2n-1)^2\pi^2a^2t/4L^2}\sin\frac{(2n-1)\pi x}{2L}$,
$n=1,2,3,\dots$
 satisfy $u_t=a^2u_{xx}$ and the boundary conditions
 $u(0,t)= u_x(L,t)=0, \ t>0$.
Any finite sum
$\dst \sum_{n=1}^m d_ne^{-(2n-1)^2\pi^2a^2t/4L^2}\sin\frac{(2n-1)\pi x}{2L}$
also has these properties.
Therefore,it is plausible to expect that that this is also true of
the infinite series
(C) $\dst u(x,t)=\sum_{n=1}^\infty
d_ne^{-(2n-1)^2\pi^2a^2t/4L^2}\sin\frac{(2n-1)\pi x}{2L}$
under suitable conditions on the coefficients $\{d_n\}$.
Since
$\dst u(x,0)=\sum_{n=1}^\infty d_n\sin\frac{(2n-1)\pi x}{2L}$,
if $\{d_n\}$ are the mixed Fourier sine coefficients of
$f$ on $[0,L]$, then
$u(x,0)=f(x)$ at all points $x$ in $[0,L]$ where the mixed Fourier
sine series
 converges to $f(x)$. In this case  (C)
is a formal solution of the initial-boundary value problem
of Definition~12.1.3.

\addtocounter{enumi}{4}
\item%12.1.8
Since $f(0)=f(1)=0$ and $f''(x)=-2$, Theorem~11.3.5\part{b}
 implies that
\begin{align*}
\alpha_n&=
\frac{4}{ n^2\pi^2}
\int_0^1\sin n\pi x\,dx
=-\frac{4}{ n^3\pi^3}\cos n\pi x\lims01=
-\frac{4}{ n^3\pi^2}(\cos n\pi-1)\\
&=\begin{cases}
\dst\frac{8}{(2m-1)^3\pi^3},&\text{if $n=2m-1$},\\
0,&\text{if  $n=2m$};
\end{cases}
\end{align*}
$S(x)=\dst\frac{8}{\pi^3}\sum_{n=1}^\infty\frac{1}{(2n-1)^3}
\sin\frac{(2n-1)\pi x}{ L}$.
From Definition~12.1.1,
\[
u(x,t)=\dst\frac{8}{\pi^3}\sum_{n=1}^\infty
\frac{1}{(2n-1)^3}
e^{-(2n-1)^2\pi^2t}\sin(2n-1)\pi x.
\]

\item
\begin{align*}
\alpha_1&=\frac{2}{\pi}\int_0^\pi x\sin^2 x\,dx=
\frac{1}{ \pi}\int_0^\pi x(1-\cos2x)\,dx
=\frac{x^2}{2\pi}\lims0\pi
-\frac{1}{\pi}\int_0^\pi x\cos2x\,dx\\
&=\frac{\pi}{2}-\frac{1}{2\pi}\left[x\sin2x\lims0\pi-\int_0^\pi
\cos2x\,dx\right]
=\frac{\pi}{2}+\frac{\sin2x}{4\pi}\lims0\pi=\frac{\pi}{2};
\end{align*}
if $n\ge2$, then
\begin{align*}
\alpha_n&=
\frac{2}{\pi}\int_0^\pi x\sin x\sin nx\,dx
=\frac{1}{\pi}\int_0^\pi x[\cos(n-1)x-\cos(n+1)x]\,dx\\
&=\frac{1}{\pi}\left[x\left[\frac{\sin(n-1)x}{ n-1}-
\frac{\sin(n+1)x}{ n+1}\right]\lims0\pi-\int_0^\pi
\left[\frac{\sin(n-1)x}{ n-1}-\frac{\sin(n+1)x}{ n+1}\right]\,dx\right]\\
&=\frac{1}{\pi}\left[\frac{\cos(n-1)x }{(n-1)^2}
-\frac{\cos(n+1)x}{(n+1)^2}\right]\lims0\pi
=\frac{1}{\pi}\left[\frac{1}{(n-1)^2}-\frac{1}{(n+1)^2}\right]
\left[(-1)^{n+1}-1\right]\\
&=\frac{4n}{(n^2-1)^2\pi}\left[(-1)^{n+1}-1\right]=
\begin{cases}
0& \text{ if }n=2m-1,\\
-\dst\frac{16m}{(4m^2-1)\pi}&\text{ if }n=2m;
\end{cases}
\end{align*}
$S(x)=\dst\frac{\pi}{2}\sin
x-\frac{16}{\pi}\sum_{n=1}^\infty\frac{n}{(4n^2-1)^2}\sin2nx$.
From Definition~12.1.1,
\[
u(x,t)=\dst\frac{\pi}{2}e^{-3t}\sin
x-\frac{16}{\pi}\sum_{n=1}^\infty\frac{n}{(4n^2-1)^2}e^{-12n^2t}\sin2nx.
\]



\item
Since $f(0)=f(L)=0$ and $f''(x)=-6x$, Theorem~11.3.5\part{b}
implies that
\begin{align*}
\alpha_n&=
\frac{36}{ n^2\pi^2}\int_0^3x\sin\frac{n\pi x}{ 3}\,dx
=-\frac{108}{ n^3\pi^3}\left[x\cos\frac{n\pi x}{ 3}\lims03
-\int_0^3\cos\frac{n\pi x}{ 3}\,dx\right]\\
&=(-1)^{n+1}\frac{108}{ n^3\pi^3}+
\frac{108}{ n^4\pi^4}\sin\frac{n\pi x}{ 3}\lims03=
(-1)^{n+1}\frac{324}{ n^3\pi^3};
\end{align*}
$S(x)=\dst-\frac{324}{\pi^3}\sum_{n=1}^\infty\frac{(-1)^n}{ n^3}
\sin\frac{n\pi x}{ 3}$.
From Definition~12.1.1,
$u(x,t)=\dst-\frac{324}{\pi^3}\sum_{n=1}^\infty\frac{(-1)^n}{
n^3}
e^{-4n^2\pi^2t/9}\sin\frac{n\pi x}{ 3}$.


\item
Since $f(0)=f(1)=f''(0)=f''(L)=0$ and $f^{(4)}=360x$,
Theorem~11.3.5\part{b} and Exercise~35\part{b} of
Section~11.3  imply that
\begin{align*}
\alpha_n&=
\frac{720}{ n^4\pi^4}\int_0^1x\sin n\pi x\,dx
=-\frac{720}{ n^5\pi^5}\left[x\cos n\pi x
\lims01-\int_0^1\cos n\pi x
\,dx\right]\\
&=(-1)^{n+1}\frac{720}{ n^5\pi^5}+\frac{720}{ n^6\pi^6}\sin\frac{n\pi
x}{ L}\lims01=
(-1)^{n+1}\frac{720}{ n^5\pi^5};
\end{align*}
$S(x)=\dst-\frac{720}{\pi^5}\sum_{n=1}^\infty\frac{(-1)^n}{ n^5}\sin n\pi
x$.
From Definition~12.1.1,
$u(x,t)=\dst-\frac{720}{\pi^5}\sum_{n=1}^\infty\frac{(-1)^n}{
n^5}e^{-7n^2\pi^2t}\sin n\pi x$.


\item
Since $f(0)=f(1)=f''(0)=f''(L)=0$ and $f^{(4)}=120(3x-1)$,
Theorem~11.3.5\part{b} and Exercise~35\part{b} of
Section~11.3 imply that
\begin{align*}
\alpha_n&=
\frac{240}{ n^4\pi^4}\int_0^1(3x-1)\sin n\pi x\,dx
=-\frac{240}{ n^5\pi^5}\left[(3x-1)\cos n\pi x
\lims01-3\int_0^1\cos n\pi x\,dx\right]\\
&=-\frac{240}{ n^5\pi^5}\left[(-1)^n2+1\right]
+\frac{720}{ n^6\pi^6}\sin n\pi x\lims01
=-\frac{240}{ n^5\pi^5}\left[1+(-1)^n2\right];
\end{align*}
$S(x)=\dst-\frac{240}{\pi^5}\sum_{n=1}^\infty\frac{1+(-1)^n2}{ n^5}\sin
n\pi x$.
From Definition~12.1.1,
\[
u(x,t)=\dst-\frac{240}{\pi^5}\sum_{n=1}^\infty\frac{1+(-1)^n2}{
n^5}e^{-2n^2\pi^2t}\sin n\pi x.
\]



\item
$\alpha_0=\dst\frac{1}{ 2}\int_0^2(x^2-4x)\,dx=\frac{1}{
L}\left(\frac{x^3}{3}-2x^2
\right)\lims02=-\frac{8}{3}$;  if $n\ge1$,
\begin{align*}
\alpha_n&=\int_0^2(x^2-4x)\cos\frac{n\pi x}{ 2}\,dx
=\frac{2}{ n\pi}\left[(x^2-4x)\sin\frac{n\pi x}{ 2}\lims02-
2\int_0^2
(x-2)\sin\frac{n\pi x}{ 2}\,dx\right]\\
&=\frac{8}{ n^2\pi^2}\left[(x-2)\cos\frac{n\pi x}{ 2}\lims02-
\int_0^2\cos\frac{n\pi x}{ 2}\,dx\right]=
\frac{16}{ n^2\pi^2}-\frac{32}{ n^3\pi^3}\sin\frac{n\pi x}{ 2}\lims02
=\frac{16}{ n^2\pi^2};
\end{align*}
$C(x)=\dst-\frac{8}{3}+\frac{16}{\pi^2}\sum_{n=1}^\infty\frac{1}{
n^2}\cos\frac{n\pi x}{ 2}$.
From Definition~12.1.3,
$u(x,t)=\dst-\frac{8}{3}+\frac{16}{\pi^2}\sum_{n=1}^\infty\frac{1}{
n^2}e^{-n^2\pi^2t}\cos\frac{n\pi x}{ 2}$.


\item
From Example~11.3.5,
$C(x)=\dst4-\frac{384}{\pi^4}\sum_{n=1}^\infty\frac{1}{
(2n-1)^4}\cos\frac{(2n-1)\pi x}{ 2}$.
From Definition~12.1.3,
$u(x,t)=\dst4-\frac{384}{\pi^4}\sum_{n=1}^\infty\frac{1}{
(2n-1)^4}e^{-3(2n-1)^2\pi^2t/4}\cos\frac{(2n-1)\pi x}{ 2}$.



\item
$\dst \alpha_0=\frac{1}{ L}\int_0^1(3x^4-4Lx^3)\,dx=\frac{1}{
L}\left(\frac{3x^5}{5}-x^4
\right)\lims01=-\frac{2}{5}$.
Since $f'(0)=f'(1)=0$ and  $f'''(x)=24(3x-1)$, Theorem~11.3.5\part{a}
implies that
\begin{align*}
\alpha_n&=
\frac{48}{ n^3\pi^3}\int_0^1(3x-1)\sin n\pi x\,dx
=-\frac{48}{ n^4\pi^4}\left[(3x-1)\cos n\pi x\lims01
-3\int_0^1\cos n\pi x\,dx\right]\\
&=-\frac{48}{ n^4\pi^4}\left[(-1)^n2+1\right]
+\frac{144}{ n^5\pi^5}\sin n\pi x\lims01=
-\frac{48}{ n^4\pi^4}\left[1+(-1)^n2\right],\quad n\ge1;
\end{align*}
$C(x)=\dst-\frac{2}{5}-\frac{48}{\pi^4}\sum_{n=1}^\infty\frac{1+(-1)^n2}{
n^4}\cos n\pi x$.
From Definition~12.1.3,
\[
u(x,t)=\dst-\frac{2}{5}-\frac{48}{\pi^4}\sum_{n=1}^\infty\frac{1+(-1)^n2}{
n^4}e^{-3n^2\pi^2t}\cos n\pi x.
\]


\item
$\dst \alpha_0=\frac{1}{ \pi}\int_0^\pi(x^4-2\pi x^3+\pi^2x^2)\,dx=\frac{1}{
\pi}\left(\frac{x^5}{5}-\frac{\pi x^4}{2}+\frac{\pi^2x^3}{3}
\right)\lims0\pi=\frac{\pi^4}{30}$.
Since  $f'(0)=f'(\pi)=0$ and  $f'''(x)=12(2x-\pi)$,
Theorem~11.3.5\part{a} implies that
\begin{align*}
\alpha_n&=
\frac{24}{ n^3\pi}\int_0^\pi(2x-\pi)\sin nx\,dx
=-\frac{24}{ n^4\pi}\left[(2x-\pi)\cos nx\lims0\pi
-2\int_0^\pi\cos nx\,dx\right]\\
&=-\frac{24}{ n^4\pi}\left[(-1)^n\pi+\pi\right]+
\frac{48}{ n^5\pi}\sin nx\lims0\pi=
-\frac{24}{ n^4}\left[1+(-1)^n\right]\\
&=\begin{cases}
0&\text{ if }n=2m-1,\\
\dst-\frac{3}{ m^4}&\text{ if }n=2m,
\end{cases}
n\ge1;
\end{align*}
$\dst C(x)=\frac{\pi^4}{30}-3\sum_{n=1}^\infty\frac{1}{
n^4}\cos2nx$.
From Definition~12.1.3,
$\dst u(x,t)=\frac{\pi^4}{30}-3\sum_{n=1}^\infty\frac{1}{
n^4}e^{-4n^2t}\cos2nx$.




\item
\begin{align*}
\alpha_n&=\frac{2}{
\pi}\int_0^\pi(\pi x-x^2)\sin\frac{(2n-1) x}{2}\,dx\\
&=-\frac{4}{(2n-1)\pi}\left[(\pi x-x^2)\cos\frac{(2n-1) x}{2}\lims0\pi-
\int_0^\pi (\pi-2x)\cos\frac{(2n-1) x}{2}\,dx\right]\\
&=\frac{8}{(2n-1)^2\pi}\left[(\pi-2x)\sin\frac{(2n-1) x}{2}\lims0\pi+2
\int_0^\pi\sin\frac{(2n-1) x}{2}\,dx\right]\\
&=(-1)^n\frac{8}{(2n-1)^2}-
\frac{32}{(2n-1)^3\pi}\cos\frac{(2n-1) x}{2}\lims0\pi\\
&=(-1)^n\frac{8}{(2n-1)^2} +\frac{32}{(2n-1)^3\pi};
\end{align*}
$S_M(x)=\dst8\sum_{n=1}^\infty\frac{1}{(2n-1)^2}\left[
(-1)^n+\frac{4}{(2n-1)\pi}\right]\sin\frac{(2n-1) x}{2}$. From
Definition~12.1.4, \[
u(x,t)=\dst8\sum_{n=1}^\infty\frac{1}{(2n-1)^2}\left[
(-1)^n+\frac{4}{(2n-1)\pi}\right]e^{-3(2n-1)^2t/4}\sin\frac{(2n-1) x}{ 2}. \]





\item
Since $f(0)=f'(1)=0$, and $f''(x)=6(1-2x)$,
Theorem~11.3.5\part{d} implies that
\begin{align*}
\alpha_n&=-\dst\frac{48}{(2n-1)^2\pi^2}\int_0^1(1-2x)
\sin\frac{(2n-1)\pi x}{2}\,dx\\
&=\dst\frac{96}{(2n-1)^3\pi^3}\left[(1-2x)\cos\frac{(2n-1)\pi
x}{2}\lims01+2\int_0^1\cos\frac{(2n-1)\pi x}{2}\,dx\right]\\
&=\dst\frac{96}{(2n-1)^3\pi^3}\left[-1+\frac{4}{(2n-1)\pi}
\sin\frac{(2n-1)\pi x}{2}\lims01\right]\\
&=\dst-\frac{96}{(2n-1)^3\pi^3}\left[1+(-1)^n\frac{4}{(2n-1)\pi}\right];
\end{align*}
$S_M(x)
=\dst-\frac{96}{\pi^3}\sum_{n=1}^\infty
\frac{1}{(2n-1)^3}\left[1+(-1)^n\frac{4}{(2n-1)\pi}\right]\sin\frac{(2n-1)\pi x
}{2}$.
From Definition~12.1.4,
\[
u(x,t)
=\dst-\frac{96}{\pi^3}\sum_{n=1}^\infty
\frac{1}{(2n-1)^3}\left[1+(-1)^n\frac{4}{(2n-1)\pi}\right]
e^{-(2n-1)^2\pi^2t}\sin\frac{(2n-1)\pi x}{2}.
\]



\item
Since $f(0)=f'(L)=f''(0)=0$ and $f'''(x)=6$,
Theorem~11.3.5\part{d} and Exercise~11.3.50\part{b}
imply that
\begin{align*}
\alpha_n&=-\frac{96}{(2n-1)^3\pi^3}\int_0^1
\cos\frac{(2n-1)\pi x}{2}\,dx\\
&=-\frac{192}{(2n-1)^4\pi^4}
\sin\frac{(2n-1)\pi x}{2}\lims01=
(-1)^n\frac{192}{(2n-1)^4\pi^4};
\end{align*}
$S_M(x)=\dst\frac{192}{\pi^4}\sum_{n=1}^\infty\frac{(-1)^n}{(2n-1)^4}
\sin\frac{(2n-1)\pi x}{2}$.
From Definition~12.1.4,
\[
u(x,t)=\dst\frac{192}{\pi^4}\sum_{n=1}^\infty\frac{(-1)^n}{(2n-1)^4}
e^{-(2n-1)^2\pi^2t}\sin\frac{(2n-1)\pi x}{2}.
\]




\item
Since $f(0)=f'(1)=f''(0)=0$ and $f'''(x)=12(2x-1)$,
Theorem~11.3.5\part{d} and Exercise~11.3.50\part{b}
imply that
\begin{align*}
\alpha_n&=-\frac{192}{(2n-1)^3\pi^3}\int_0^1(2x-1)
\cos\frac{(2n-1)\pi x}{2}\,dx\\
&=-\frac{384}{(2n-1)^4\pi^4}
\left[(2x-1)\sin\frac{(2n-1)\pi x}{2}\lims01
-2\int_0^1\sin\frac{(2n-1)\pi x}{2}\,dx\right]\\
&=-\frac{384}{(2n-1)^4\pi^4}
\left[(-1)^{n+1}1+\frac{4}{(2n-1)\pi}
\cos\frac{(2n-1)\pi x}{2}\lims01\right]\\
&=-\frac{384}{(2n-1)^4\pi^4}
\left[(-1)^{n+1}1-\frac{4}{(2n-1)\pi}\right]
=\frac{384}{(2n-1)^4\pi^4}
\left[(-1)^n+\frac{4}{(2n-1)\pi}\right];
\end{align*}
$S_M(x)=\dst\frac{384}{\pi^4}\sum_{n=1}^\infty\frac{1}{(2n-1)^4}
\left[(-1)^n+\frac{4}{(2n-1)\pi}\right]\sin\frac{(2n-1)\pi x}{2}$.
From Definition~12.1.4,
\[
u(x,t)=\dst\frac{384}{\pi^4}\sum_{n=1}^\infty\frac{1}{(2n-1)^4}
\left[(-1)^n+\frac{4}{(2n-1)\pi}\right]
e^{-(2n-1)^2\pi^2t}\sin\frac{(2n-1)\pi x}{2}.
\]


\addtocounter{enumi}{2}
\item%12.1.36
From Example~11.3.3,
$C_M(x)=
\dst\frac{8}{\pi^2}\sum_{n=1}^\infty\frac{1}{(2n-1)^2}
\cos\frac{(2n-1)\pi x}{2}$.
From Definition~12.1.5,
\[
u(x,t)=
\dst\frac{8}{\pi^2}\sum_{n=1}^\infty\frac{1}{(2n-1)^2}
e^{-3(2n-1)^2\pi^2t/4}
\cos\frac{(2n-1)\pi x}{2}.
\]




\item
Since $f'(0)=f(\pi)=0$ and $f''(x)=-2$,
Theorem~11.3.5\part{c}  implies that
\[
\alpha_n=
\frac{16}{(2n-1)^2\pi}
\int_0^\pi\cos\frac{(2n-1) x}{2}\,dx
=\frac{32}{(2n-1)^3\pi}\sin\frac{(2n-1) x}{2}\lims0\pi
=(-1)^{n+1}\frac{32}{(2n-1)^3\pi};
\]
$C_M(x)=-\dst
\frac{32}{\pi}\sum_{n=1}^\infty\frac{(-1)^n}{
(2n-1)^3}\cos\frac{(2n-1) x}{2}$.
From Definition~12.1.5,
\[
u(x,t)=-\dst
\frac{32}{\pi}\sum_{n=1}^\infty\frac{(-1)^n}{
(2n-1)^3}e^{-7(2n-1)^2t/4}\cos\frac{(2n-1) x}{2}.
\]



\item
Since $f'(0)=f(1)=0$ and $f''(x)=6(2x+1)$,
Theorem~11.3.5\part{c}  implies that
\begin{align*}
\alpha_n&=
-\frac{48L}{(2n-1)^2\pi^2}
\int_0^1(2x+1)\cos\frac{(2n-1)\pi x}{2}\,dx\\
&=-\frac{96}{(2n-1)^3\pi^3}\left[(2x+1)\sin\frac{(2n-1)\pi x}{2}
\lims01-2\int_0^1\sin\frac{(2n-1)\pi x}{2}\right]\,dx\\
&=-\frac{96}{(2n-1)^3\pi^3}
\left[(-1)^{n+1}3-\frac{4}{(2n-1)\pi}\cos\frac{(2n-1)\pi x}{2}
\lims01\right] \\
&=\frac{96}{(2n-1)^3\pi^3}
\left[(-1)^n3+\frac{4}{(2n-1)\pi}\right];
\end{align*}
$C_M(x)=\dst\frac{96}{\pi^3}\sum_{n=1}^\infty\frac{1}{(2n-1)^3}\left[
(-1)^n3+\frac{4}{(2n-1)\pi}\right]\cos\frac{(2n-1)\pi
x}{2}$.
From Definition~12.1.5,
\[
u(x,t)=\dst\frac{96}{\pi^3}\sum_{n=1}^\infty\frac{1}{(2n-1)^3}\left[
(-1)^n3+\frac{4}{(2n-1)\pi}\right]e^{-(2n-1)^2\pi^2t/4}\cos\frac{(2n-1)\pi
x}{2}.
\]


\item
Theorem~11.3.5\part{c}  and Exercise~11.3.42\part{b}
imply that
Since $f'(0)=f(1)=f''(1)=0$ and $f'''(x)=12(2x-1)$,
\begin{align*}
\alpha_n&=
\frac{192}{(2n-1)^3\pi^3}
\int_0^1(2x-1)\sin\frac{(2n-1)\pi x}{2}\,dx\\
&=-\frac{384}{(2n-1)^4\pi^4}\left[(2x-1)\cos\frac{(2n-1)\pi x}{2}
\lims01-2\int_0^1\cos\frac{(2n-1)\pi x}{2}\right]\,dx\\
&=-\frac{384}{(2n-1)^4\pi^4}
\left[1-\frac{4}{(2n-1)\pi}\sin\frac{(2n-1)\pi x}{2}
\lims01\right]
=-\frac{384}{(2n-1)^4\pi^4}
\left[1+\frac{(-1)^n4}{(2n-1)\pi}\right];
\end{align*}
$C_M(x)=-\dst\frac{384}{\pi^4}\sum_{n=1}^\infty\frac{1}{(2n-1)^4}\left[
1+\frac{(-1)^n4}{(2n-1)\pi}\right]\cos\frac{(2n-1)\pi
x}{2}$.
From Definition~12.1.5,
\[
u(x,t)=-\dst\frac{384}{\pi^4}\sum_{n=1}^\infty\frac{1}{(2n-1)^4}\left[
1+\frac{(-1)^n4}{(2n-1)\pi}\right]e^{-(2n-1)^2\pi^2t/4}\cos\frac{(2n-1)\pi
x}{2}.
\]


\item
$\alpha_n=\dst\frac{2}{ L}\int_0^{L/2}\sin\frac{n\pi x}{ L}\,dx=
-\frac{2}{ n\pi}\cos\frac{n\pi x}{ L}\lims0{L/2}=
\frac{2}{ n\pi}\left[1-\cos\frac{n\pi}{2}\right]$;\\
$S(x)=\dst\frac{2}{\pi}\sum_{n=1}^\infty\frac{1}{ n}
\left[1-\cos\frac{n\pi}{2}\right]\sin\frac{n\pi x}{
L}$.
From Definition~12.1.1,
\[
u(x,t)=\dst\frac{2}{\pi}\sum_{n=1}^\infty\frac{1}{ n}
\left[1-\cos\frac{n\pi}{2}\right]e^{-n^2\pi^2t^2/L^2}\sin\frac{n\pi x}{
L}.
\]

\item
\begin{align*}
\alpha_n&=\frac{2}{ L}\int_0^{L/2}\sin\frac{(2n-1)\pi x}{
2L}\,dx
=\dst-\frac{4}{(2n-1)\pi}\cos\frac{(2n-1)\pi x}{ 2L}\lims0{L/2}\\
&=\frac{4}{(2n-1)\pi}\left[1-\cos\frac{(2n-1)\pi)}{4}\right];
\end{align*}
$S_M(x)=\dst\frac{4}{\pi}\sum_{n=1}^\infty\frac{1}{2n-1}
\left[1-\cos\frac{(2n-1)\pi)}{4}\right]
\sin\frac{(2n-1)\pi x}{ 2L}$.
From Definition~12.1.4,
\[
u(x,t)=\dst\frac{4}{\pi}\sum_{n=1}^\infty\frac{1}{2n-1}
\left[1-\cos\frac{(2n-1)\pi)}{4}\right]
e^{-(2n-1)^2\pi^2a^2t/4L^2}\sin\frac{(2n-1)\pi x}{ 2L}.
\]


\item
Let $u(x,t)=v(x,t)+q(x)$; then $u_t=v_t$ and $u_{xx}=v_{xx}+q''$,
so
\[
\begin{gathered}
v_t=9v_{xx}+9q''-54x,\quad 0<x<4,\quad t>0,\\
v(0,t)=1-q(0),\quad v(4,t)=61-q(4),\quad t>0,\\
v(x,0)=2-x+x^3-q(x),\quad 0\le x\le4.
\end{gathered}
\tag{A}
\]
We want $q''(x)=6x,\ q(0)=1,\ q(4)=61$; $q(x)=x^3+a_1+a_2x$;
$q(0)=1\Rightarrow a_1=1$; $q(x)=x^3+1+a_2x$; $q(4)=61\Rightarrow
a_2=-1$; $q(x)=x^3+1-x$.  Now (A) reduces to
\[
\begin{gathered}
v_t=9v_{xx},\quad 0<x<4,\quad t>0,\\
v(0,t)=0,\quad v(4,t)=0,\quad t>0,\\
v(x,0)=1,\quad 0\le x\le4,
\end{gathered}
\]
which we solve by separation of variables.
\begin{align*}
 \alpha_n&=\frac{1}{ 2}\int_0^L\sin\frac{n\pi x}{ 4}\,dx=
-\dst\frac{2}{ n\pi}\cos\frac{n\pi x}{ 4}\lims04\\
&=
\frac{2}{ n\pi}[1-(-1)^n]=
\begin{cases}
\dst\frac{4}{(2m-1)\pi}&\text{ if }n=2m-1,\\
0&\text{ if }n=2m;
\end{cases}
\end{align*}
$S(x)=\dst\frac{4}{\pi}\sum_{n=1}^\infty\frac{1}{(2n-1)}
\sin\frac{(2n-1)\pi x}{ 4}$.
From Definition~12.1.1,
\[
v(x,t)=\dst\frac{4}{\pi}\sum_{n=1}^\infty\frac{1}{(2n-1)}
e^{-9\pi^2n^2t/16}\sin\frac{(2n-1)\pi x}{ 4}.
\]
Therefore,
\[
u(x,t)=1-x+x^3+\dst\frac{4}{\pi}\sum_{n=1}^\infty\frac{e^{-9\pi^2n^2t/16}}{(2n-1)}
\sin\frac{(2n-1)\pi x}{ 4}.
\]



\item
Let $u(x,t)=v(x,t)+q(x)$; then $u_t=v_t$ and $u_{xx}=v_{xx}+q''$,
so
\[
\begin{gathered}
v_t=3u_{xx}+3q''-18x,\quad 0<x<1,\quad t>0,\\
v_x(0,t)=-1-q'(0),\quad v(1,t)=-1-q(1),\quad t>0,\\
v(x,0)=x^3-2x-q(x),\quad 0\le x\le1.
\end{gathered}
\tag{A}
\]
We want $q''(x)=6x,\ q'(0)=-1,\ q(1)=-1$; $q'(x)=3x^2+a_2$;
$q'(0)=-1\Rightarrow a_2=-1$; $q'(x)=3x^2-1$;
 $q(x)=x^3-x+a_1x$; $q(1)=-1\Rightarrow
a_1=-1$; $q(x)=x^3-x-1$.  Now (A) reduces to
\begin{gather*}
v_t=3u_{xx},\quad 0<x<1,\quad t>0,\\
v_x(0,t)=0,\quad v(1,t)=0,\quad t>0,\\
v(x,0)=1-x,\quad 0\le x\le1.
\end{gather*}
From Example~11.3.3,
$C_M(x)=
\dst\frac{8}{\pi^2}\sum_{n=1}^\infty\frac{1}{(2n-1)^2}
\cos\frac{(2n-1)\pi x}{2}$.
From Definition~12.1.5,
$v(x,t)=
\dst\frac{8}{\pi^2}\sum_{n=1}^\infty\frac{1}{(2n-1)^2}
e^{-3(2n-1)^2\pi^2t/4}
\cos\frac{(2n-1)\pi x}{2}$.
Therefore,
\[
u(x,t)=-1-x+x^3+
\dst\frac{8}{\pi^2}\sum_{n=1}^\infty\frac{1}{(2n-1)^2}
e^{-3(2n-1)^2\pi^2t/4}
\cos\frac{(2n-1)\pi x}{2}.
\]



\item
Let $u(x,t)=v(x,t)+q(x)$; then $u_t=v_t$ and $u_{xx}=v_{xx}+q''$,
so
\[
\begin{gathered}
v_t=v_{xx}+q''+\pi^2\sin\pi x,\quad 0<x<1,\quad t>0,\\
v(0,t)=-q(0),\quad v_x(1,t)=-\pi-q'(1),\quad t>0,\\
v(x,0)=2\sin\pi x-q(x),\quad 0\le x\le1.
\end{gathered}
\tag{A}
\]
We want $q''(x)=-\pi^2\sin\pi x,\ q(0)=0,\ q'(1)=-\pi$;
$q'(x)=\pi\cos\pi x+a_2$;
$q'(1)=-\pi\Rightarrow a_2=0$; $q'(x)=\pi\cos\pi x$;
 $q(x)=\sin\pi x+a_1$; $q(0)=0\Rightarrow
a_1=0$; $q(x)=\sin\pi x$.  Now (A) reduces to
\begin{gather*}
v_t=v_{xx},\quad 0<x<1,\quad t>0,\\
v(0,t)=0,\quad v_x(1,t)=0,\quad t>0,\\
v(x,0)=\sin\pi x,\quad 0\le x\le1.
\end{gather*}
\begin{align*}
\alpha_n&=2\dst\int_0^1\sin\pi x\sin\frac{(2n-1)\pi x}{2}\,dx=
\int_0^1\left[\frac{\cos(2n-3)\pi x}{ 2}-
\frac{\cos(2n+1)\pi x}{2}\right]\,dx\\
&=\frac{2}{\pi}\left[\frac{\sin(2n-3)\pi x/2}{(2n-3)}-\frac{\sin(2n+1)\pi
x/2}{(2n+1)}\right]\lims01\\
&=(-1)^n\frac{2}{\pi}\left[\frac{1}{2n-3}-
\frac{1}{2n+1}\right]=(-1)^n\frac{8}{\pi}\frac{1}{(2n+1)(2n-3)};
\end{align*}
$S_M(x)=\dst{\frac{8}{\pi} \sum^\infty_{n=1}
\frac{(-1)^n}{(2n+1)(2n-3)} \sin\frac{(2n-1)\pi x}{
2}}$.
From Definition~12.1.4,
\[
v(x,t)=\dst{\frac{8}{\pi} \sum^\infty_{n=1}
\frac{(-1)^n}{(2n+1)(2n-3)} e^{-(2n-1)^2\pi^2 t/4}\sin\frac{(2n-1)\pi x}{
2}}.
\]
Therefore,
$\dst u(x,t)=\sin\pi x+\dst{\frac{8}{\pi} \sum^\infty_{n=1}
\frac{(-1)^n}{(2n+1)(2n-3)} e^{-(2n-1)^2\pi^2 t/4}\sin\frac{(2n-1)\pi x}{
2}}$.

\item
\begin{alist}
\item Since $f$ is piecewise smooth of $[0,L]$, there is a constant
$K$ such that $|f(x)|\le K$, $0\le x\le L$. Therefore,
$|\alpha_n|=\dst\frac{2}{ L}\left|\int_0^Lf(x)\sin\frac{n\pi x}{
L}\,dx\right|\le\frac{2}{ L}\int_0^L|f(x)|\,dx=2K$. Hence,
$\dst|\alpha_ne^{-n^2\pi^2a^2t/L}|\le2Ke^{-n^2\pi^2a^2t/L}$,
so $u(x,t)$ converges for all $x$ if $t>0$, by the comparison test.

\item
Let $t$ be a fixed positive number.
Apply Theorem~12.1.2 with $z=x$ and
$w_n(x)=\dst\alpha_ne^{-n^2\pi^2t/L^2}\sin\frac{n\pi x}{ L}$.
 Then
$w_n'(x)=\dst\frac{\pi}{ L}n\alpha_ne^{-n^2\pi^2t/L^2}\cos\frac{n\pi x}{
L}$,  so
$|w_n'(x)|\dst\le\frac{2K\pi}{ L}ne^{-n^2\pi^2t/L^2}$,
$-\infty<x<\infty$.
Since $\dst\sum_{n=1}^\infty ne^{-n^2\pi^2a^2t/L^2}$  converges if
$t>0$,   Theorem~12.1.1 (with
$z_1=x_1$ and
$z_2=x_2$  arbitrary) implies the conclusion.

\item
Since $\dst\sum_{n=1}^\infty n^2e^{-n^2\pi^2a^2t/L^2}$ also converges
if
$t>0$, an argument like that in \part{b} with
$w_n(x)=\dst n\alpha_ne^{-n^2\pi^2t/L^2}\cos\frac{n\pi x}{ L}$
yields the conclusion.

\item
Let $x$ be arbitrary, but fixed.
Apply Theorem~12.1.2 with $z=t$ and\\
$\dst w_n(t)=\alpha_ne^{-n^2\pi^2a^2t/L^2}\sin\frac{n\pi x}{ L}$.
Then
$\dst w_n'(t)=-\dst\frac{\pi^2a^2}{
L^2}n^2\alpha_ne^{-n^2\pi^2a^2t/L^2}\sin\frac{n\pi x}{ L}$,
so
$\dst|w_n'(t)|\le\frac{2K\pi^2a^2}{
L}n^2e^{-n^2\pi^2a^2t_0/L^2}$ if $t>t_0$.
Since $\dst\sum_{n=1}^\infty n^2e^{-n^2\pi^2a^2t_0/L^2}$  converges,
 Theorem~12.1.1 (with $z_1=t_0>0$ and $z_2=t_1$ arbitrary
implies the conclusion for $t\ge t_0$.
However, since $t_0$ is an arbitrary positive number, this holds
for $t>0$.
\end{alist}
\end{sectionSolutions}


\section{The Wave Equation}

\begin{sectionSolutions}
\addtocounter{enumi}{-1}
\item%12.2.1
$\beta_n=2\left[\int_0^{1/2}x\sin n\pi x
+\int_{1/2}^1(1-x)\sin n\pi x\,dx\right]$;
\begin{align*}
\int_0^{1/2}x\sin n\pi x\,dx
&=-\frac{1}{ n\pi}\left[x\cos n\pi x\lims0{1/2}-
\int_0^{1/2}\cos n\pi x\,dx\right]\\
&=-\frac{1}{2n\pi}\cos\frac{n\pi}{2}+\frac{1}{ n^2\pi^2}\sin n\pi x
\lims0{1/2}
=-\frac{1}{2n\pi}\cos\frac{n\pi}{2}+\frac{1}{ n^2\pi^2}\sin\frac{n\pi}{2};\\
\int_0^{1/2}(1-x)\sin n\pi x\,dx
&=-\frac{1}{n\pi}\left[(1-x)\cos n\pi x\lims{1/2}1+
\int_0^{1/2}\cos n\pi x\,dx\right]\\
&=\frac{1}{2n\pi}\cos\frac{n\pi}{2}-\frac{1}{ n^2\pi^2}\sin n\pi x
\lims{1/2}1
=\frac{1}{
2n\pi}\cos\frac{n\pi}{2}+\frac{1}{ n^2\pi^2}\sin\frac{n\pi}{2};\\
b_n&=\dst\frac{4}{ n^2\pi^2}\sin\frac{n\pi}{2}=
\begin{cases}
(-1)^{m+1}\dst\frac{4}{(2m-1)^2\pi^2}&\text{if }n=2m-1\\
0&\text{if }n=2m;
\end{cases}
\end{align*}
$S_g(x)=\dst\frac{4}{\pi^2}\sum_{n=1}^\infty
\frac{(-1)^{n+1}}{(2n-1)^2}\sin(2n-1)\pi x$.
From Definition~12.1.1,
\[
u(x,t)=\dst\frac{4}{3\pi^3}\sum_{n=1}^\infty
\frac{(-1)^{n+1}}{(2n-1)^3}\sin3(2n-1)\pi t\sin(2n-1)\pi x.
\]


\addtocounter{enumi}{-1}
\item%12.2.2
Since $f(0)=f(1)=0$ and $f''(x)=-2$,
Theorem~11.3.5\part{b} implies that
\begin{align*}
\alpha_n&=
\frac{4}{ n^2\pi^2}
\int_0^1\sin n\pi x\,dx
=-\frac{4}{ n^3\pi^3}\cos n\pi x\lims01=
-\frac{4}{ n^3\pi^2}(\cos n\pi-1)\\
&=\begin{cases}
\dst\frac{8}{(2m-1)^3\pi^3},&\text{if $n=2m-1$},\\
0,&\text{if  $n=2m$};\\
\end{cases}
\end{align*}
$S_f(x)=\dst\frac{8}{\pi^3}\sum_{n=1}^\infty\frac{1}{(2n-1)^3}
\sin(2n-1)\pi x$.
From Definition~12.1.1,
\[
u(x,t)=\dst\frac{8}{\pi^3}\sum_{n=1}^\infty\frac{1}{(2n-1)^3}
\cos3(2n-1)\pi t\sin(2n-1)\pi x.
\]




\item
Since $g(0)=g(1)=0$ and $g''(x)=-2$,
Theorem~11.3.5\part{b} implies that
\begin{align*}
\beta_n&=
\frac{4}{ n^2\pi^2}
\int_0^1\sin n\pi x\,dx
=-\frac{4}{ n^3\pi^3}\cos n\pi x\lims01=
-\frac{4}{ n^3\pi^2}(\cos n\pi-1)\\
&=\begin{cases}
\dst\frac{8}{(2m-1)^3\pi^3},&\text{if $n=2m-1$},\\
0,&\text{if  $n=2m$};\\
\end{cases}
\end{align*}
$S_g(x)=\dst\frac{8}{\pi^3}\sum_{n=1}^\infty\frac{1}{(2n-1)^3}
\sin(2n-1)\pi x$.
From Definition~12.1.1,
\[
u(x,t)=\dst\frac{8}{3\pi^4}\sum_{n=1}^\infty\frac{1}{(2n-1)^4}
\sin3(2n-1)\pi t\sin(2n-1)\pi x.
\]


\item
From Example~11.2.6,
$S_f(x)=\dst\frac{324}{\pi^3}\sum_{n=1}^\infty\frac{(-1)^n}{
n^3}\sin\frac{n\pi x}{3}$.
From Definition~12.1.1,
\[
u(x,t)=\dst\frac{324}{\pi^3}\sum_{n=1}^\infty\frac{(-1)^n}{
n^3}\cos\frac{8\pi nt}{3}\sin\frac{n\pi x}{3}.
\]




\item
From Example~11.2.6
$S_g(x)=\dst\frac{324}{\pi^3}\sum_{n=1}^\infty\frac{(-1)^n}{
n^3}\sin\frac{n\pi x}{3}$.
From Definition~12.1.1,
\[
u(x,t)=\dst\frac{81}{2\pi^4}\sum_{n=1}^\infty\frac{(-1)^n}{
n^4}\sin\frac{8\pi nt}{3}\sin\frac{n\pi x}{3}.
\]



\item
\begin{align*}
\alpha_1&=\frac{2}{\pi}\int_0^\pi x\sin^2 x\,dx=
\frac{1}{ \pi}\int_0^\pi x(1-\cos2x)\,dx
=\frac{x^2}{2\pi}\lims0\pi
-\frac{1}{\pi}\int_0^\pi x\cos2x\,dx\\
&=\frac{\pi}{2}-\frac{1}{2\pi}\left[x\sin2x\lims0\pi-\int_0^\pi
\cos2x\,dx\right]
=\frac{\pi}{2}+\frac{\sin2x}{4\pi}\lims0\pi=\frac{\pi}{2};
\end{align*}
if $n\ge2$, then
\begin{align*}
\alpha_n&=
\frac{2}{\pi}\int_0^\pi x\sin x\sin nx\,dx
=\frac{1}{\pi}\int_0^\pi x[\cos(n-1)x-\cos(n+1)x]\,dx\\
&=\frac{1}{\pi}\left[x\left[\frac{\sin(n-1)x}{ n-1}-
\frac{\sin(n+1)x}{ n+1}\right]\lims0\pi-\int_0^\pi
\left[\frac{\sin(n-1)x}{ n-1}-\frac{\sin(n+1)x}{ n+1}\right]\,dx\right]\\
&=\frac{1}{\pi}\left[\frac{\cos(n-1)x }{(n-1)^2}
-\frac{\cos(n+1)x}{(n+1)^2}\right]\lims0\pi
=\frac{1}{\pi}\left[\frac{1}{(n-1)^2}-\frac{1}{(n+1)^2}\right]
\left[(-1)^{n+1}-1\right]\\
&=\frac{4n}{(n^2-1)^2\pi}\left[(-1)^{n+1}-1\right]=
\begin{cases}
0& \text{ if }n=2m-1,\\
-\dst\frac{16m}{(4m^2-1)\pi}&\text{ if }n=2m;
\end{cases}
\end{align*}
$S_f(x)=\dst\frac{\pi}{2}\sin x-\frac{16}{\pi}
\sum_{n=1}^\infty\frac{n}{(4n^2-1)^2}\sin 2nx$.
From Definition~12.1.1,
\[
u(x,t)=\dst\frac{\pi}{2}\cos\sqrt{5}\,t\sin x-\frac{16}{\pi}
\sum_{n=1}^\infty\frac{n}{(4n^2-1)^2}\cos2n\sqrt5\,t\sin 2nx.
\]


\item
\begin{align*}
\beta_1&=\frac{2}{\pi}\int_0^\pi x\sin^2 x\,dx=
\frac{1}{ \pi}\int_0^\pi x(1-\cos2x)\,dx
=\frac{x^2}{2\pi}\lims0\pi
-\frac{1}{\pi}\int_0^\pi x\cos2x\,dx\\
&=\frac{\pi}{2}-\frac{1}{2\pi}\left[x\sin2x\lims0\pi-\int_0^\pi
\cos2x\,dx\right]
=\frac{\pi}{2}+\frac{\sin2x}{4\pi}\lims0\pi=\frac{\pi}{2};
\end{align*}
if $n\ge2$ then
\begin{align*}
\beta_n&=
\frac{2}{\pi}\int_0^\pi x\sin x\sin nx\,dx
=\frac{1}{\pi}\int_0^\pi x[\cos(n-1)x-\cos(n+1)x]\,dx\\
&=\frac{1}{\pi}\left[x\left[\frac{\sin(n-1)x}{ n-1}-
\frac{\sin(n+1)x}{ n+1}\right]\lims0\pi-\int_0^\pi
\left[\frac{\sin(n-1)x}{ n-1}-\frac{\sin(n+1)x}{ n+1}\right]\,dx\right]\\
&=\frac{1}{\pi}\left[\frac{\cos(n-1)x }{(n-1)^2}
-\frac{\cos(n+1)x}{(n+1)^2}\right]\lims0\pi
=\frac{1}{\pi}\left[\frac{1}{(n-1)^2}-\frac{1}{(n+1)^2}\right]
\left[(-1)^{n+1}-1\right]\\
&=\frac{4n}{(n^2-1)^2\pi}\left[(-1)^{n+1}-1\right]=
\begin{cases}
0& \text{ if }n=2m-1,\\
-\dst\frac{16m}{(4m^2-1)\pi}&\text{ if }n=2m;
\end{cases}
\end{align*}
$S_g(x)=\dst\frac{\pi}{2}\sin x-\frac{16}{\pi}
\sum_{n=1}^\infty\frac{n}{(4n^2-1)^2}\sin 2nx$.
From Definition~12.1.1,
\[
u(x,t)=\dst\frac{\pi}{2\sqrt5}\sin\sqrt{5}\,t\sin x-\frac{8}{\pi\sqrt5}
\sum_{n=1}^\infty\frac{1}{(4n^2-1)^2}\sin2n\sqrt5\,t\sin 2nx.
\]


\item
Since $f(0)=f(1)=f''(0)=f''(L)=0$ and $f^{(4)}=360x$,
Theorem~11.3.5\part{b} and Exercise~35\part{b} of
Section~11.3  imply that
\begin{align*}
\alpha_n&=
\frac{720}{ n^4\pi^4}\int_0^1x\sin n\pi x\,dx
=-\frac{720}{ n^5\pi^5}\left[x\cos n\pi x
\lims01-\int_0^1\cos n\pi x
\,dx\right]\\
&=(-1)^{n+1}\frac{720}{ n^5\pi^5}+\frac{720}{ n^6\pi^6}\sin\frac{n\pi
x}{ L}\lims01=
(-1)^{n+1}\frac{720}{ n^5\pi^5};
\end{align*}
$S_f(x)=\dst-\frac{720}{\pi^5}\sum_{n=1}^\infty\frac{(-1)^n}{ n^5}\sin
n\pi x$.
From Definition~12.1.1,
$u(x,t)=\dst-\frac{720}{\pi^5}\sum_{n=1}^\infty\frac{(-1)^n}{ n^5}
\cos3n\pi t\sin n\pi x$.


\item
\begin{alist}
\item $t$ must be in some interval of the form $[mL/a,(m+1)L/a]$.
If $\dst\frac{mL}{ a}\le t\le\left(m+\frac{1}{2}\right)\frac{L}{ a}$, then
\part{i}
holds with $0\le\tau\le L/2a$. If
$\dst\left(m+\frac{1}{2}\right)\frac{L}{ a}\le
t\le\frac{(m+1)L}{ a}$, then \part{ii} holds with $0\le\tau\le L/2a$.


\item Suppose that \part{i} holds. Since
\[
\cos\frac{(2n-1)\pi a}{ L}\left(\tau+\frac{mL}{ a}\right)
=\cos\frac{(2n-1)\pi a\tau}{ L}\cos(2n-1)m\pi=
(-1)^m\cos\frac{(2n-1)\pi a\tau}{ L},
\]
(A) implies that $u(x,t)=(-1)^mu(x,\tau)$.

 Suppose that \part{ii} holds. Since
\begin{align*}
\cos\frac{(2n-1)\pi a}{ L}\left(-\tau+\frac{(m+1)L}{ a}\right)
&=\cos\frac{(2n-1)\pi a\tau}{ L}\cos(2n-1)(m+1)\pi\\
&=(-1)^{m+1}\cos\frac{(2n-1)\pi a\tau}{ L},
\end{align*}
(B) implies that  that $u(x,t)=(-1)^{m+1}u(x,\tau)$.
\end{alist}

\item
Since $f'(0)=f(2)=0$ and $f''(x)=-2$,
Theorem~11.3.5(c) implies that
\begin{align*}
\alpha_n&=
\frac{32}{(2n-1)^2\pi^2}
\int_0^2\cos\frac{(2n-1)\pi x}{4}\,dx \\
&=\frac{128}{(2n-1)^3\pi^3}\sin\frac{(2n-1)\pi x}{4}\lims04
=(-1)^{n+1}\frac{128}{(2n-1)^3\pi^3};
\end{align*}
$C_{M\!f}(x)=-\dst
\frac{128}{\pi^3}\sum_{n=1}^\infty\frac{(-1)^n}{
(2n-1)^3}\cos\frac{(2n-1)\pi x}{4}$.
From Exercise~12.2.17,
\[
u(x,t)=-\dst
\frac{128}{\pi^3}\sum_{n=1}^\infty\frac{(-1)^n}{
(2n-1)^3}\cos\frac{3(2n-1)\pi t}{ 4}\,\cos\frac{(2n-1)\pi x}{4}.
\]


\item
Since $g'(0)=g(2)=0$ and $g''(x)=-2$,
Theorem~11.3.5(c) implies that
\begin{align*}
\beta_n&=
\frac{32}{(2n-1)^2\pi^2}
\int_0^2\cos\frac{(2n-1)\pi x}{4}\,dx \\
&=\frac{128}{(2n-1)^3\pi^3}\sin\frac{(2n-1)\pi x}{4}\lims04
=(-1)^{n+1}\frac{128}{(2n-1)^3\pi^3};
\end{align*}
$C_{M\!f}(x)=-\dst
\frac{128}{\pi^3}\sum_{n=1}^\infty\frac{(-1)^n}{
(2n-1)^3}\cos\frac{(2n-1)\pi x}{4}$.
From Exercise~12.2.17,
\[
u(x,t)=-\dst
\frac{512}{3\pi^4}\sum_{n=1}^\infty\frac{(-1)^n}{
(2n-1)^4}\sin\frac{3(2n-1)\pi t}{ 4}\,\cos\frac{(2n-1)\pi x}{4}.
\]

\item
Since $f'(0)=f(1)=0$ and $f''(x)=6(2x+1)$,
Theorem~11.3.5\part{c}  implies that
\begin{align*}
\alpha_n&=
-\frac{48L}{(2n-1)^2\pi^2}
\int_0^1(2x+1)\cos\frac{(2n-1)\pi x}{2}\,dx\\
&=-\frac{96}{(2n-1)^3\pi^3}\left[(2x+1)\sin\frac{(2n-1)\pi x}{2}
\lims01-2\int_0^1\sin\frac{(2n-1)\pi x}{2}\right]\,dx\\
&=-\frac{96}{(2n-1)^3\pi^3}
\left[(-1)^{n+1}3-\frac{4}{(2n-1)\pi}\cos\frac{(2n-1)\pi x}{2}
\lims01\right] \\
&=\frac{96}{(2n-1)^3\pi^3}
\left[(-1)^n3+\frac{4}{(2n-1)\pi}\right];\\
C_{M\!f}(x)&=\dst\frac{96}{\pi^3}\sum_{n=1}^\infty\frac{1}{(2n-1)^3}\left[
(-1)^n3+\frac{4}{(2n-1)\pi}\right]\cos\frac{(2n-1)\pi
x}{2}.
\end{align*}
From Exercise~12.2.17,
\[
u(x,t)=\dst\frac{96}{\pi^3}\sum_{n=1}^\infty\frac{1}{(2n-1)^3}\left[
(-1)^n3+\frac{4}{(2n-1)\pi}\right]\cos\frac{(2n-1)\sqrt5\,\pi
t}{2}\cos\frac{(2n-1)\pi x}{2}.
\]


\item
Since $g'(0)=g(1)=0$ and $g''(x)=6(2x+1)$,
Theorem~11.3.5\part{c}  implies that
\begin{align*}
\beta_n&=
-\frac{48L}{(2n-1)^2\pi^2}
\int_0^1(2x+1)\cos\frac{(2n-1)\pi x}{2}\,dx\\
&=-\frac{96}{(2n-1)^3\pi^3}\left[(2x+1)\sin\frac{(2n-1)\pi x}{2}
\lims01-2\int_0^1\sin\frac{(2n-1)\pi x}{2}\right]\,dx\\
&=-\frac{96}{(2n-1)^3\pi^3}
\left[(-1)^{n+1}3-\frac{4}{(2n-1)\pi}\cos\frac{(2n-1)\pi x}{2}
\lims01\right] \\
&=\frac{96}{(2n-1)^3\pi^3}
\left[(-1)^n3+\frac{4}{(2n-1)\pi}\right];
\end{align*}
$\dst
C_{M\!g}(x)=\dst\frac{96}{\pi^3}\sum_{n=1}^\infty\frac{1}{(2n-1)^3}\left[
(-1)^n3+\frac{4}{(2n-1)\pi}\right]\cos\frac{(2n-1)\pi
x}{2}$.
From Exercise~12.2.17,
\[
u(x,t)=\dst\frac{192}{\pi^4\sqrt5}\sum_{n=1}^\infty\frac{1}{(2n-1)^4}\left[
(-1)^n3+\frac{4}{(2n-1)\pi}\right]\sin\frac{(2n-1)\sqrt5\,\pi
t}{2}\cos\frac{(2n-1)\pi x}{2}.
\]


\item
Since $f'(0)=f(1)=f''(1)=0$ and $f'''(x)=24(x-1)$,
Theorem~11.3.5\part{c}  and Exercise~42\part{b} of
Section~11.3 imply that
\begin{align*}
\alpha_n&=
\frac{384}{(2n-1)^3\pi^3}
\int_0^1(x-1)\sin\frac{(2n-1)\pi x}{2}\,dx\\
&=-\frac{768}{(2n-1)^4\pi^4}\left[(x-1)\cos\frac{(2n-1)\pi x}{2}
\lims01-\int_0^1\cos\frac{(2n-1)\pi x}{2}\right]\,dx\\
&=-\frac{768}{(2n-1)^4\pi^4}
\left[1-\frac{2}{(2n-1)\pi}\sin\frac{(2n-1)\pi x}{2}
\lims01\right] \\
&=-\frac{768}{(2n-1)^4\pi^4}
\left[1+\frac{(-1)^n2}{(2n-1)\pi}\right];
\end{align*}
$\dst
C_{M\!f}(x)=-\dst\frac{384}{\pi^4}\sum_{n=1}^\infty\frac{1}{(2n-1)^4}\left[
1+\frac{(-1)^n4}{(2n-1)\pi}\right]\cos\frac{(2n-1)\pi
x}{2}$.
From Exercise~12.2.17,
\[
u(x,t)=-\dst\frac{384}{\pi^4}\sum_{n=1}^\infty\frac{1}{(2n-1)^4}\left[
1+\frac{(-1)^n4}{(2n-1)\pi}\right]\cos\frac{3(2n-1)\pi t}{2}\cos\frac{(2n-1)\pi
x}{2}.
\]


\item
Since $g'(0)=g(1)=g''(1)=0$ and $g'''(x)=24(x-1)$,
Theorem~11.3.5\part{c}  and Exercise~11.2.42\part{b}
 imply that
\begin{align*}
\beta_n&=
\frac{384}{(2n-1)^3\pi^3}
\int_0^1(x-1)\sin\frac{(2n-1)\pi x}{2}\,dx\\
&=-\frac{768}{(2n-1)^4\pi^4}\left[(x-1)\cos\frac{(2n-1)\pi x}{2}
\lims01-\int_0^1\cos\frac{(2n-1)\pi x}{2}\right]\,dx\\
&=-\frac{768}{(2n-1)^4\pi^4}
\left[1-\frac{2}{(2n-1)\pi}\sin\frac{(2n-1)\pi x}{2}
\lims01\right] \\
&=-\frac{768}{(2n-1)^4\pi^4}
\left[1+\frac{(-1)^n2}{(2n-1)\pi}\right];\\
C_{M\!g}(x)&=-\dst\frac{384}{\pi^4}\sum_{n=1}^\infty\frac{1}{(2n-1)^4}\left[
1+\frac{(-1)^n4}{(2n-1)\pi}\right]\cos\frac{(2n-1)\pi
x}{2}.
\end{align*}
From Exercise~12.2.17,
\[
u(x,t)=-\dst\frac{768}{3\pi^5}\sum_{n=1}^\infty\frac{1}{(2n-1)^5}\left[
1+\frac{(-1)^n4}{(2n-1)\pi}\right]\sin\frac{3(2n-1)\pi t}{2}\cos\frac{(2n-1)\pi
x}{2}.
\]


\item
Since $f'(0)=f(1)=f''(1)=0$ and $f'''(x)=24(x-1)$,
Theorem~11.3.5\part{c}  and Exercise~42\part{b} of
Section~11.3 imply that
\begin{align*}
\alpha_n&=
\frac{384}{(2n-1)^3\pi^3}
\int_0^1(x-1)\sin\frac{(2n-1)\pi x}{2}\,dx\\
&=-\frac{768}{(2n-1)^4\pi^4}\left[(x-1)\cos\frac{(2n-1)\pi x}{2}
\lims01-\int_0^1\cos\frac{(2n-1)\pi x}{2}\right]\,dx\\
&=-\frac{768}{(2n-1)^4\pi^4}
\left[1-\frac{2}{(2n-1)\pi}\sin\frac{(2n-1)\pi x}{2}
\lims01\right] \\
&=-\frac{768}{(2n-1)^4\pi^4}
\left[1+\frac{(-1)^n2}{(2n-1)\pi}\right];
\end{align*}
$\dst
C_{M\!f}(x)=-\dst\frac{768}{\pi^4}\sum_{n=1}^\infty\frac{1}{(2n-1)^4}\left[
1+\frac{(-1)^n2}{(2n-1)\pi}\right]\cos\frac{(2n-1)\pi
x}{2}$.
From Exercise~12.2.17,
\[
u(x,t)=-\dst\frac{768}{\pi^4}\sum_{n=1}^\infty\frac{1}{(2n-1)^4}\left[
1+\frac{(-1)^n2}{(2n-1)\pi}\right]\cos\frac{(2n-1)\pi t}{2}\cos\frac{(2n-1)\pi
x}{2}.
\]


\item
Setting $A=(2n-1)\pi x/2L$ and $B=(2n-1)\pi at/2L$ in the
 identities\\
$\cos A\cos B=\dst\frac{1}{2}[\cos(A+B)+\cos(A-B)]$
and
$\cos A\sin B=\dst\frac{1}{2}[\sin(A+B)-\sin(A-B)]$
yields
\[
\cos\frac{(2n-1)\pi at}{ 2L}\cos\frac{(2n-1)\pi x}{ 2L}
=\frac{1}{2}\left[\cos\frac{(2n-1)\pi(x+at)}{
2L}+\cos\frac{(2n-1)\pi(x-at)}{ 2L}\right]
\tag{A}
\]
and
\[
\begin{aligned}
\dst \sin\frac{(2n-1)\pi at}{ 2L}\cos\frac{(2n-1)\pi x}{ 2L}
&=\dst\frac{1}{2}\left[\sin\frac{(2n-1)\pi(x+at)}{2L}
-\sin\frac{(2n-1)\pi(x-at)}{2L}\right]\\
&=\dst\frac{(2n-1)\pi}{4L}\int_{x-at}^{x+at}
\cos\frac{(2n-1)\pi\tau}{2L}\,d\tau.
\end{aligned}
\tag{B}
\]
Since $\dst C_{M\!f}(x)=\sum_{n=1}^\infty\alpha_n\cos\frac{(2n-1)\pi
x}{ 2L}$,
(A) implies that
\[
\dst\sum_{n=1}^\infty\alpha_n \cos\frac{(2n-1)\pi at}{ 2L}\cos\frac{(2n-1)\pi
x}{ 2L}=
\dst\frac{1}{2}[C_{M\!f}(x+at)+C_{M\!f}(x-at)].
\tag{C}
\]
Since it can be shown that a mixed Fourier cosine series can be
integrated term by term between any two limits,
(B) implies that
\begin{align*}
\dst\sum_{n=1}^\infty\frac{2L\beta_n}{(2n-1)\pi a}
\sin\frac{(2n-1)\pi at}{ 2L}\cos\frac{(2n-1)\pi x}{2L}
&=\dst\frac{1}{2a}\sum_{n=1}^\infty \beta_n\int_{x-at}^{x+at}
\cos\frac{(2n-1)\pi\tau}{ 2L}\,d\tau\\
&=\dst\frac{1}{2a}\int_{x-at}^{x+at}\left(\sum_{n=1}^\infty \beta_n
\cos\frac{(2n-1)\pi\tau}{ 2L}\right)\,d\tau\\
&=\dst\frac{1}{2a}\int_{x-at}^{x+at}C_{M\!g}(\tau)\,d\tau.
\end{align*}
This and  (C) imply that
\[
u(x,t)=\frac{1}{2}[C_{M\!f}(x+at)+C_{M\!f}(x-at)]+
\frac{1}{2a}\int_{x-at}^{x+at}C_{M\!g}(\tau)\,d\tau.
\]



\item
 We begin by looking for functions of the form
$v(x,t)=X(x)T(t)$
that are not identically zero and satisfy
$v_{tt}=a^2v_{xx}$,
$v(0,t)=0$, $v_x(L,t)=0$
for all $(x,t)$. As shown in the text, $X$  and $T$
must satisfy
$X''+\lambda X=0$
and
(B) $T''+a^2\lambda T=0$
for the same value of $\lambda$.
Since $v(0,t)=X(0)T(t)$ and $v_x(L,t)=X'(L)T(t)$ and we don't want
$T$ to be  identically
zero,  $X(0)=0$ and $X'(L)=0$. Therefore,$\lambda$
must be an eigenvalue of
(C) $X''+\lambda X=0$, $X(0)=0$, $X'(L)=0$,
and $X$ must be a $\lambda$-eigenfunction.
 From Theorem~11.1.4, the eigenvalues of (C)
are $\lambda_n=\dst\frac{(2n-1)^2}{\pi^2/4L^2}$,
integer), with associated eigenfunctions
$X_n=\dst\sin\frac{(2n-1)\pi x}{2L}$, $n=1$, 2, 3,\dots.
Substituting $\lambda=\dst\frac{(2n-1)^2\pi^2}{ 4L^2}$  into
(B) yields
$T''+((2n-1)^2\pi^2a^2/4L^2)T=0$,
which has the general solution
\[
T_n=\alpha_n\cos\frac{(2n-1)\pi at}{2L}+\frac{2\beta_nL}{(2n-1)\pi a}\sin\frac{
(2n-1)\pi at}{ 2L},
\]
where $\alpha_n$  and $\beta_n$ are constants. Now let
\[
v_n(x,t)=X_n(x)T_n(t)=\left(\alpha_n\cos\frac{(2n-1)\pi at}{2L}+
\frac{2\beta_nL}{ (2n-1)\pi a}\sin\frac{(2n-1)\pi at}{2L}\right)
\sin\frac{(2n-1)\pi x}{2L}.
\]
Then
\[
\frac{\partial v_n}{\partial t}(x,t)=
\left(-\frac{(2n-1)\pi a }{2L}\alpha_n\sin\frac{(2n-1)\pi at}{2L}
+\beta_n\cos\frac{(2n-1)\pi at}{2L}\right)\sin\frac{(2n-1)\pi x}{2L},
\]
so
\[
v_n(x,0)=\alpha_n\sin\frac{(2n-1)\pi x}{2L}
\quad\text{and}\quad
\frac{\partial v_n}{\partial t}(x,0)=\beta_n
\sin\frac{(2n-1)\pi x}{2L}.
\]
Therefore,$v_n$ satisfies (A) with
$f(x)=\dst\alpha_n \sin\frac{(2n-1)\pi x}{ 2L}$ and
$g(x)=\dst\beta_n\sin
\frac{(2n-1)\pi x}{ 2L}$. More generally, if
$\alpha_1,\alpha_2,\dots,\alpha_m$
and $\beta_1,\beta_2,\dots,\beta_m$ are constants and
\[
u_m(x,t)=\sum_{n=1}^m\left(\alpha_n\cos\frac{(2n-1)\pi at}{
2L}+\frac{2\beta_nL}{ (2n-1)\pi a}\sin\frac{(2n-1)\pi at}{2L}\right)
\sin\frac{(2n-1)\pi x}{2L},
\]
then $u_m$ satisfies (A) with
\[
f(x)=\sum_{n=1}^m \alpha_n\sin\frac{(2n-1)\pi x}{2L}
\quad\text{and}\quad
g(x)=\sum_{n=1}^m \beta_n\sin\frac{(2n-1)\pi x}{2L}.
\]
 This motivates the definition.

\item
Since $f(0)=f'(1)=0$, and $f''(x)=6(1-2x)$,
Theorem~11.3.5\part{d} implies that
\begin{align*}
\alpha_n&=-\dst\frac{48}{(2n-1)^2\pi^2}\int_0^1(1-2x)
\sin\frac{(2n-1)\pi x}{2}\,dx\\
&=\dst\frac{96}{(2n-1)^3\pi^3}\left[(1-2x)\cos\frac{(2n-1)\pi
x}{2}\lims01+2\int_0^1\cos\frac{(2n-1)\pi x}{2}\,dx\right]\\
&=\dst\frac{96}{(2n-1)^3\pi^3}\left[-1+\frac{4}{(2n-1)\pi}
\sin\frac{(2n-1)\pi x}{2}\lims01\right]\\
&=\dst-\frac{96}{(2n-1)^3\pi^3}\left[1+(-1)^n\frac{4}{(2n-1)\pi}\right];
\end{align*}
$S_{M\!f}(x)
=\dst-\frac{96}{\pi^3}\sum_{n=1}^\infty
\frac{1}{(2n-1)^3}\left[1+(-1)^n\frac{4}{(2n-1)\pi}\right]\sin\frac{(2n-1)\pi x
}{2}$.
From Exercise~12.2.34,
\[
u(x,t)
=\dst-\frac{96}{\pi^3}\sum_{n=1}^\infty
\frac{1}{(2n-1)^3}\left[1+(-1)^n\frac{4}{(2n-1)\pi}\right]
\cos\frac{3(2n-1)\pi t}{2}\sin\frac{(2n-1)\pi x}{2}.
\]


\item
Since $g(0)=g'(1)=0$, and $g''(x)=6(1-2x)$,
Theorem~11.3.5\part{d} implies that
\begin{align*}
\beta_n&=-\dst\frac{48}{(2n-1)^2\pi^2}\int_0^1(1-2x)
\sin\frac{(2n-1)\pi x}{2}\,dx\\
&=\dst\frac{96}{(2n-1)^3\pi^3}\left[(1-2x)\cos\frac{(2n-1)\pi
x}{2}\lims01+2\int_0^1\cos\frac{(2n-1)\pi x}{2}\,dx\right]\\
&=\dst\frac{96}{(2n-1)^3\pi^3}\left[-1+\frac{4}{(2n-1)\pi}
\sin\frac{(2n-1)\pi x}{2}\lims01\right]\\
&=\dst-\frac{96}{(2n-1)^3\pi^3}\left[1+(-1)^n\frac{4}{(2n-1)\pi}\right];
\end{align*}
$S_{M\!g}(x)
=\dst-\frac{96}{\pi^3}\sum_{n=1}^\infty
\frac{1}{(2n-1)^3}\left[1+(-1)^n\frac{4}{(2n-1)\pi}\right]\sin\frac{(2n-1)\pi x
}{2}$.
From Exercise~12.2.34,
\[
u(x,t)
=\dst-\frac{64}{\pi^4}\sum_{n=1}^\infty
\frac{1}{(2n-1)^4}\left[1+(-1)^n\frac{4}{(2n-1)\pi}\right]
\sin\frac{3(2n-1)\pi t}{2}\sin\frac{(2n-1)\pi x}{2}.
\]

\item
Since $f(0)=f'(\pi)=f''(0)=0$ and $f'''(x)=6$,
Theorem~11.3.5\part{d} and Exercise~11.3.50\part{b} imply that
\begin{align*}
\alpha_n&
=-\frac{96}{(2n-1)^3\pi}\int_0^\pi\cos\frac{(2n-1) x}{2}\,dx
=-\frac{192}{(2n-1)^4\pi}\sin\frac{(2n-1)x}{2}\lims0\pi
\\&
=(-1)^n\frac{192}{(2n-1)^4\pi};\\
S_{M\!f}(x)&=\dst\frac{192}{\pi}\sum_{n=1}^\infty\frac{(-1)^n}{(2n-1)^4}
\sin\frac{(2n-1) x}{2}.\qquad
\text{From Exercise~12.2.34,}\\
u(x,t)&=\dst\frac{192}{\pi}\sum_{n=1}^\infty\frac{(-1)^n}{(2n-1)^4}
\cos\frac{(2n-1)\sqrt3\,t}{2}\sin\frac{(2n-1) x}{2}.
\end{align*}



\item
Since $g(0)=g'(\pi)=g''(0)=0$ and $g'''(x)=6$,
Theorem~11.3.5\part{d} and Exercise~50\part{b} imply that
\begin{align*}
\beta_n&
=-\frac{96}{(2n-1)^3\pi}\int_0^\pi\cos\frac{(2n-1) x}{2}\,dx
=-\frac{192}{(2n-1)^4\pi}\sin\frac{(2n-1)x}{2}\lims0\pi
\\&
=(-1)^n\frac{192}{(2n-1)^4\pi};\\
S_{M\!g}(x)&=\dst\frac{192}{\pi}\sum_{n=1}^\infty\frac{(-1)^n}{(2n-1)^4}
\sin\frac{(2n-1) x}{2}.\qquad\text{From Exercise~12.2.34,}\\
u(x,t)&=\dst\frac{384}{\sqrt3\,\pi}\sum_{n=1}^\infty\frac{(-1)^n}{(2n-1)^5}
\sin\frac{(2n-1)\sqrt3\,t}{2}\sin\frac{(2n-1) x}{2}.
\end{align*}

\item
Since $f(0)=f'(1)=f''(0)=0$ and $f'''(x)=12(2x-1)$,
Theorem~11.3.5\part{d} and Exercise~11.3.50\part{b}
imply that
\begin{align*}
\alpha_n&=-\frac{192}{(2n-1)^3\pi^3}\int_0^1(2x-1)
\cos\frac{(2n-1)\pi x}{2}\,dx\\
&=-\frac{384}{(2n-1)^4\pi^4}
\left[(2x-1)\sin\frac{(2n-1)\pi x}{2}\lims01
-2\int_0^1\sin\frac{(2n-1)\pi x}{2}\,dx\right]\\
&=-\frac{384}{(2n-1)^4\pi^4}
\left[(-1)^{n+1}+\frac{4}{(2n-1)\pi}
\cos\frac{(2n-1)\pi x}{2}\lims01\right]\\
&=-\frac{384}{(2n-1)^4\pi^4}
\left[(-1)^{n+1}-\frac{4}{(2n-1)\pi}\right]
=\frac{384}{(2n-1)^4\pi^4}
\left[(-1)^n+\frac{4}{(2n-1)\pi}\right];
\end{align*}
$\dst S_{M\!f}(x)=\dst\frac{384}{\pi^4}\sum_{n=1}^\infty\frac{1}{(2n-1)^4}
\left[(-1)^n+\frac{4}{(2n-1)\pi}\right]\sin\frac{(2n-1)\pi x}{2}$.
From Exercise~12.2.34,
\[
u(x,t)=\dst\frac{384}{\pi^4}\sum_{n=1}^\infty\frac{1}{(2n-1)^4}
\left[(-1)^n+\frac{4}{(2n-1)\pi}\right]\cos(2n-1)\pi
t\sin\frac{(2n-1)\pi x}{2}.
\]


\item
Since $g(0)=g'(1)=g''(0)=0$ and $g'''(x)=12(2x-1)$,
Theorem~11.3.5\part{d} and Exercise~11.3.50\part{b}
imply that
\begin{align*}
\beta_n&=-\frac{192}{(2n-1)^3\pi^3}\int_0^1(2x-1)
\cos\frac{(2n-1)\pi x}{2}\,dx\\
&=-\frac{384}{(2n-1)^4\pi^4}
\left[(2x-1)\sin\frac{(2n-1)\pi x}{2}\lims01
-2\int_0^1\sin\frac{(2n-1)\pi x}{2}\,dx\right]\\
&=-\frac{384}{(2n-1)^4\pi^4}
\left[(-1)^{n+1}+\frac{4}{(2n-1)\pi}
\cos\frac{(2n-1)\pi x}{2}\lims01\right]\\
&=-\frac{384}{(2n-1)^4\pi^4}
\left[(-1)^{n+1}-\frac{4}{(2n-1)\pi}\right]
=\frac{384}{(2n-1)^4\pi^4}
\left[(-1)^n+\frac{4}{(2n-1)\pi}\right];
\end{align*}
$\dst S_{M\!g}(x)=\dst\frac{384}{\pi^4}\sum_{n=1}^\infty\frac{1}{(2n-1)^4}
\left[(-1)^n+\frac{4}{(2n-1)\pi}\right]\sin\frac{(2n-1)\pi x}{2}$.
From Exercise~12.2.34,
\[
u(x,t)=\dst\frac{384}{\pi^5}\sum_{n=1}^\infty\frac{1}{(2n-1)^5}
\left[(-1)^n+\frac{4}{(2n-1)\pi}\right]\sin(2n-1)\pi
t\sin\frac{(2n-1)\pi x}{2}.
\]

\item
Since $f$ is continuous on $[0,L]$ and $f'(L)=0$,
Theorem~11.3.4 implies that $S_{M\!f}(x)=f(x)$, $0\le x\le L$. From
 Exercise~11.3.58,
 $S_{M\!f}$  is the odd periodic
extension (with period $2L$) of the function
$r(x)=
\begin{cases}
f(x),&0\le x\le L,\\f(2L-x),&L< x\le 2L,
\end{cases}$
which is continuous on $[0,2L]$.
Since $r(0)=r(2L)=f(0)=0$, $S_{M\!f}$  is continuous on
$(-\infty,\infty)$. Moreover,
$r'(x)=
\begin{cases}
f'(x),&0< x< L,\\-f'(2L-x),&L< x< 2L,
\end{cases}$
 $r_+'(0)=f'_+(0)$, $r_-'(2L)=-f_+'(0)$, and, since $f'_-(L)=0$,
$r'(L)=0$. Hence, $r$ is differentiable on $[0,2L]$. Since
$r(0)=r(2L)=f(0)=0$,
 Theorem~12.2.3\part{a} with $h=r$, $p=S_{M\!f}$,
and
$L$ replaced by $2L$  implies that $S_{M\!f}$  is differentiable
on $(-\infty,\infty)$. Similarly, $S_{M\!g}$ is differentiable
on $(-\infty,\infty)$.

Now we note that
$r''(x)=
\begin{cases}
f''(x),&0< x< L,\\f''(2L-x),&L< x< 2L,
\end{cases}$
$r''(L)=f''_-(L)$, and   $r_+''(0)=
r_-''(2L)=f_+''(0)=0$.
 Since $S_{M\!f}'$ is the even periodic
extension of $r'$, Theorem~12.2.3\part{b} with $h=r'$,
$q=S_{M\!f}'$,  and
$L$ replaced by $2L$ implies that  $S_{M\!f}'$  is differentiable
on $(-\infty,\infty)$.  Now follow the argument used to complete the
proof of Theorem~12.2.4.

\item
From Example~11.3.5,
$C_f(x)=\dst4-\frac{768}{\pi^4}\sum_{n=1}^\infty\frac{1}{
(2n-1)^4}\cos\frac{(2n-1)\pi x}{ 2}$.
From Exercise~12.2.49,
\[
u(x,t)=\dst4-\frac{768}{\pi^4}\sum_{n=1}^\infty\frac{1}{
(2n-1)^4}\cos\frac{\sqrt5(2n-1)\pi t}{2}\,\cos\frac{(2n-1)\pi x}{ 2}.
\]





\item
$\dst \alpha_0=\int_0^\pi(3x^4-4Lx^3)\,dx=
\left(\frac{3x^5}{5}-\pi x^4
\right)\lims0\pi=-\frac{2\pi^4}{5}$.
Since $f'(0)=f'(\pi)=0$ and  $f'''(x)=24(3x-\pi)$,
Theorem~11.3.5\part{a} implies that
\begin{align*}
\alpha_n&=
\frac{48}{ n^3\pi}\int_0^\pi(3x-\pi)\sin nx\,dx
=-\frac{48}{ n^4\pi}\left[(3x-\pi)\cos nx\lims0\pi
-3\int_0^\pi\cos nx\,dx\right]\\
&=-\frac{48}{ n^4\pi}\left[(-1)^n2\pi+\pi\right]
+\frac{144}{ n^5\pi}\sin nx\lims0\pi=
-\frac{48}{ n^4}\left[1+(-1)^n2\right],\quad n\ge1;
\end{align*}
$C_f(x)=\dst-\frac{2\pi^4}{5}-48\sum_{n=1}^\infty\frac{1+(-1)^n2}{
n^4}\cos n x$.
From Exercise~12.2.49,
\[
u(x,t)=\dst-\frac{2\pi^4}{5}-48\sum_{n=1}^\infty\frac{1+(-1)^n2}{
n^4}\cos2nt\cos nx.
\]



\item
$\dst \beta_0=\int_0^\pi(3x^4-4Lx^3)\,dx=\frac{1}{
\pi}\left(\frac{3x^5}{5}-\pi x^4
\right)\lims0\pi=-\frac{2\pi^4}{5}$.
Since $g'(0)=g'(\pi)=0$ and  $g'''(x)=24(3x-\pi)$,
Theorem~11.3.5\part{a} implies that
\begin{align*}
\beta_n&=
\frac{48}{ n^3\pi}\int_0^\pi(3x-\pi)\sin nx\,dx
=-\frac{48}{ n^4\pi}\left[(3x-\pi)\cos nx\lims0\pi
-3\int_0^\pi\cos nx\,dx\right]\\
&=-\frac{48}{ n^4\pi}\left[(-1)^n2\pi+\pi\right]
+\frac{144}{ n^5\pi}\sin nx\lims0\pi=
-\frac{48}{ n^4}\left[1+(-1)^n2\right],\quad n\ge1;
\end{align*}
$\dst C_g(x)=\dst-\frac{2\pi^4}{5}-48\sum_{n=1}^\infty\frac{1+(-1)^n2}{
n^4}\cos nx$.
From Exercise~12.2.49,
\[
u(x,t)=\dst-\frac{2\pi^4t}{5}-24\sum_{n=1}^\infty\frac{1+(-1)^n2}{
n^5}\sin2nt\cos nx.
\]


\item
$\dst \alpha_0=\frac{1}{ \pi}\int_0^\pi(x^4-2\pi x^3+\pi^2x^2)\,dx=\frac{1}{
\pi}\left(\frac{x^5}{5}-\frac{\pi x^4}{2}+\frac{\pi^2x^3}{3}
\right)\lims0\pi=\frac{\pi^4}{30}$.
Since  $f'(0)=f'(\pi)=0$ and  $f'''(x)=12(2x-\pi)$,
Theorem~11.3.5\part{a} implies that
\begin{align*}
\alpha_n&=
\frac{24}{ n^3\pi}\int_0^\pi(2x-\pi)\sin nx\,dx
=-\frac{24}{ n^4\pi}\left[(2x-\pi)\cos nx\lims0\pi
-2\int_0^\pi\cos nx\,dx\right]\\
&=-\frac{24}{ n^4\pi}\left[(-1)^n\pi+\pi\right]+
\frac{48}{ n^5\pi}\sin nx\lims0\pi=
-\frac{24}{ n^4}\left[1+(-1)^n\right]\\
&=\begin{cases}
0&\text{ if }n=2m-1,\\
\dst-\frac{3}{ m^4}&\text{ if }n=2m,
\end{cases}
n\ge1;
\end{align*}
$\dst C_f(x)=\frac{\pi^4}{30}-3\sum_{n=1}^\infty\frac{1}{
n^4}\cos2nx$.
From Exercise~12.2.49,
$u(x,t)=\dst\frac{\pi^4}{30}-3\sum_{n=1}^\infty\frac{1}{
n^4}\cos8nt\cos2nx$.


\item
$\dst \beta_0=\frac{1}{ \pi}\int_0^\pi(x^4-2\pi x^3+\pi^2x^2)\,dx=\frac{1}{
\pi}\left(\frac{x^5}{5}-\frac{\pi x^4}{2}+\frac{\pi^2x^3}{3}
\right)\lims0\pi=\frac{\pi^4}{30}$.
Since  $g'(0)=g'(\pi)=0$ and  $g'''(x)=12(2x-\pi)$,
Theorem~11.3.5\part{a} implies that
\begin{align*}
\beta_n&=
\frac{24}{ n^3\pi}\int_0^\pi(2x-\pi)\sin nx\,dx
=-\frac{24}{ n^4\pi}\left[(2x-\pi)\cos nx\lims0\pi
-2\int_0^\pi\cos nx\,dx\right]\\
&=-\frac{24}{ n^4\pi}\left[(-1)^n\pi+\pi\right]+
\frac{48}{ n^5\pi}\sin nx\lims0\pi=
-\frac{24}{ n^4}\left[1+(-1)^n\right]\\
&=\begin{cases}
0&\text{ if }n=2m-1,\\
\dst-\frac{3}{ m^4}&\text{ if }n=2m,
\end{cases}
\ge1;
\end{align*}
$\dst C_g(x)=\frac{\pi^4}{30}-3\sum_{n=1}^\infty\frac{1}{
n^4}\cos2nx$.
From Exercise~12.2.49,
$u(x,t)=\dst\frac{\pi^4t}{30}-\frac{3}{8}\sum_{n=1}^\infty\frac{1}{
n^5}\sin8nt\cos2nx$.


\item
Setting $A=n\pi x/L$ and $B=n\pi at/L$ in the
identities
$\cos A\cos B=\dst\frac{1}{2}[\cos(A+B)+\cos(A-B)]$
and
$\cos A\sin B=\dst\frac{1}{2}[\sin(A+B)-\sin(A-B)]$
yields
\[
\cos\frac{n\pi at}{ L}\cos\frac{n\pi x}{ L}
=\frac{1}{2}\left[\cos\frac{n\pi(x+at)}{
L}+\cos\frac{n\pi(x-at)}{ L}\right]
\tag{A}
\]
and
\[
\begin{aligned}
\dst \sin\frac{n\pi at}{ L}\cos\frac{n\pi x}{ L}
&=\dst\frac{1}{2}\left[\sin\frac{n\pi(x+at)}{L}
-\sin\frac{n\pi(x-at)}{L}\right]\\
&=\dst\frac{n\pi}{2L}\int_{x-at}^{x+at}
\cos\frac{n\pi\tau}{ L}\,d\tau.
\end{aligned}
\tag{B}
\]
Since $C_f(x)=\dst\alpha_0+\sum_{n=1}^\infty \alpha_n\sin\frac{n\pi x}{L}$,
(A) implies that
\[
\alpha_0+\dst\sum_{n=1}^\infty\alpha_n \cos\frac{n\pi at}{ L}\cos\frac{n\pi
x}{ L}
=\dst\frac{1}{2}[C_f(x+at)+C_f(x-at)].
\tag{C}
\]
Since it can be shown that a Fourier sine series can be integrated
term by term between any two limits,
(B) implies that
\begin{align*}
\beta_0t+\dst\sum_{n=1}^\infty\frac{\beta_nL}{ n\pi a}
\sin\frac{n\pi at}{ L}\cos\frac{n\pi x}{ L}
&=\beta_0t+\dst\frac{1}{2a}\sum_{n=1}^\infty \beta_n\int_{x-at}^{x+at}
\cos\frac{n\pi\tau}{ L}\,d\tau\\
&=\dst\frac{1}{2a}\int_{x-at}^{x+at}\left(\beta_0+\sum_{n=1}^\infty
\beta_n\cos\frac{n\pi\tau}{ L}\right)\,d\tau\\
&=\dst\frac{1}{2a}\int_{x-at}^{x+at}C_g(\tau)\,d\tau.
\end{align*}
This and  (C) imply that
\[
u(x,t)=\frac{1}{2}[C_f(x+at)+C_f(x-at)]+
\frac{1}{2a}\int_{x-at}^{x+at}C_g(\tau)\,d\tau.
\]

\item
\begin{alist}
\item
Since $|p_n(x)|\le1$ and $|q_n(t)|\le1$
for all $t$, $|k_np_n(x)q_n(t)|\le|k_n|$ for all $(x,t)$, and the
comparison test implies the conclusion.

\item If $t$ is fixed but arbitrary, then
 $|k_np_n'(x)q_n(t)|\le |\lambda|n|k_n|$, so Theorem~12.1.2
with $z=x$ and $w_n(x)=k_np_n(x)q_n(t)$ justifies term by term
differentiation with respect to $x$ on $(-\infty,\infty)$.
 If $x$ is fixed but arbitrary, then
 $|k_np_n(x)q_n'(t)|\le |\mu|n|k_n|$, so Theorem~12.1.2
with $z=t$ and $w_n(t)=k_np_n(x)q_n(t)$ justifies term by term
differentiation with respect to $t$ on $(-\infty,\infty)$.

\item The argument is similar to argument use in \part{b}.

\item Apply \part{b} and \part{c}  to the series
$\dst\sum_{n=1}^\infty\alpha_n\cos\frac{n\pi at}{
L}\sin\frac{n\pi x}{ L}$ and
$\dst\sum_{n=1}^\infty
\frac{\beta_nL}{ n\pi a}\sin\frac{n\pi at}{ L}
\sin\frac{n\pi x}{ L}$, recalling that the individual terms in
the series satisfy $u_{tt}=a^2u_{xx}$ for all $(x,t)$.

\item
$u(x,t)=\dst\frac{f(x+ct)+f(x-ct)}{2}+
\frac{1}{2a}\int_{x-at}^{x+at}g(u)\,du$.
\end{alist}

\item
\begin{align*}
 u(x,t)&=\frac{(x+at)+(x-at)}{2}+\frac{1}{2a}\int_{x-at}^{x+at}4au\,du
=x+2\int_{x-at}^{x+at}u\,du=x+u^2\lims{x-at}{x+at}\\
&=x+(x+at)^2-(x-at)^2=
x(1+4at).
\end{align*}


\item
\begin{align*}
\dst u(x,t)&=\frac{\sin(x+at)+\sin(x-at)}{2}+\frac{1}{2a}
\int_{x-at}^{x+at}a\cos u\,du\\
&=\frac{\sin(x+at)+\sin(x-at)}{2}
+\frac{\sin(x+at)-\sin(x-at)}{2}
=\sin(x+at).
\end{align*}




\item
\begin{align*}
u(x,t)&=\frac{(x+at)\sin(x+at)+(x-at)\sin(x-at)}{2}
+\frac{1}{2a}\int_{x-at}^{x+at}\sin u\,du\\
&=\frac{x[\sin(x+at)+\sin(x-at)]}{2}+\frac{at[\sin(x+at)-\sin(x-at)]}{2}\\
&+\frac{\cos(x-at)-\cos(x+at)}{2a}\\
&=x\sin x\cos at+at\cos x\sin at+\frac{\sin x\sin at}{ a}.
\end{align*}
\end{sectionSolutions}

\section{Laplace's Equation in Rectangular Coordinates}

\begin{sectionSolutions}
\item
Since $f(0)=f(1)=0$ and $f''(x)=2-6x$, Theorem~11.3.5\part{b}
implies that
\begin{align*}
\alpha_n&=
-\frac{4}{ n^2\pi^2}\int_0^1(4-6x)\sin\frac{n\pi x}{ 2}\,dx
=\frac{8}{ n^3\pi^3}\left[(4-6x)\cos\frac{n\pi x}{ 2}\lims02
+6\int_0^2\cos\frac{n\pi x}{ 2}\,dx\right]\\
&=-\frac{32}{ n^3\pi^3}\left(1+(-1)^n2\right)+\frac{96}{
n^4\pi^4}\sin\frac{n\pi x}{ 2}\lims02
=-\frac{32}{ n^3\pi^3}\left[1+(-1)^n2\right];
\end{align*}
$S(x)=\dst-\frac{32}{\pi^3}\sum_{n=1}^\infty\frac{[1+(-1)^n2]}{ n^3}
\sin\frac{n\pi x}{ 2}$.
From Example~12.3.1,
\[
u(x,y)=\dst-\frac{32}{\pi^3}\sum_{n=1}^\infty\frac{[1+(-1)^n2]\sinh
n\pi(3-y)/2}{ n^3\sinh3n\pi/2}
\sin\frac{n\pi x}{ 2}.
\]



\item
\begin{align*}
\alpha_1&=\frac{2}{\pi}\int_0^\pi x\sin^2 x\,dx=
\frac{1}{ \pi}\int_0^\pi x(1-\cos2x)\,dx
=\frac{x^2}{2\pi}\lims0\pi
-\frac{1}{\pi}\int_0^\pi x\cos2x\,dx\\
&=\frac{\pi}{2}-\frac{1}{2\pi}\left[x\sin2x\lims0\pi-\int_0^\pi
\cos2x\,dx\right]
=\frac{\pi}{2}+\frac{\sin2x}{4\pi}\lims0\pi=\frac{\pi}{2};
\end{align*}
if $n\ge2$, then
\begin{align*}
\alpha_n&=
\frac{2}{\pi}\int_0^\pi x\sin x\sin nx\,dx
=\frac{1}{\pi}\int_0^\pi x[\cos(n-1)x-\cos(n+1)x]\,dx\\
&=\frac{1}{\pi}\left[x\left[\frac{\sin(n-1)x}{ n-1}-
\frac{\sin(n+1)x}{ n+1}\right]\lims0\pi-\int_0^\pi
\left[\frac{\sin(n-1)x}{ n-1}-\frac{\sin(n+1)x}{ n+1}\right]\,dx\right]\\
&=\frac{1}{\pi}\left[\frac{\cos(n-1)x }{(n-1)^2}
-\frac{\cos(n+1)x}{(n+1)^2}\right]\lims0\pi
=\frac{1}{\pi}\left[\frac{1}{(n-1)^2}-\frac{1}{(n+1)^2}\right]
\left[(-1)^{n+1}-1\right]\\
&=\frac{4n}{(n^2-1)^2\pi}\left[(-1)^{n+1}-1\right]=
\begin{cases}
0& \text{ if }n=2m-1,\\
-\dst\frac{16m}{(4m^2-1)\pi}&\text{ if }n=2m;
\end{cases}
\end{align*}
$S(x)=\dst\frac{\pi}{2}\sin
x-\frac{16}{\pi}\sum_{n=1}^\infty\frac{n}{(4n^2-1)^2}\sin2nx$.
From Example~12.3.1,
\[
u(x,y)=\dst\frac{\pi}{2}\frac{\sinh(1-y)}{\sinh1}\sin
x-\frac{16}{\pi}\sum_{n=1}^\infty\frac{n\sinh2n(1-y)}{(4n^2-1)^2\sinh2n}\sin2nx.
\]



\item
  $\alpha_0=\dst\int_0^1(1-x)\,dx=-\frac{(1-x)^2}{2}\lims01=\frac{1}{2}$;
if $n\ge1$,
\begin{align*}
\alpha_n&=2\int_0^1(1-x)\cos n\pi x\,dx=
\frac{2}{ n\pi}\left[(1-x)\sin n\pi x\lims01+
\int_0^1\sin n\pi x\,dx\right]\\
&=-\frac{2}{ n^2\pi^2}\cos n\pi x\lims01
=\frac{2}{ n^2\pi^2}[1-(-1)^n]=
\begin{cases}
\dst\frac{4}{(2m-1)^2\pi^2} &\text{ if }n=2m-1,\\
0 &\text{ if }n=2m;
\end{cases}
\end{align*}
$C(x)=\dst\frac{1}{2}+\frac{4}{\pi^2}\sum_{n=1}^\infty\frac{1}{(2n-1)^2}
\cos(2n-1)\pi x$.
From Example~12.3.3,
\[
u(x,y)=\dst\frac{y}{2}+\frac{4}{\pi^3}\sum_{n=1}^\infty\frac{
\sinh(2n-1)\pi y}{(2n-1)^3\cosh2(2n-1)\pi}
\cos(2n-1)\pi x.
\]


\item
$\alpha_0=\dst\int_0^1(x-1)^2\,dx=\frac{(x-1)^3}{3}\lims01=\frac{1}{3}$;
if
$n\ge1$, then
\begin{align*}
\alpha_n&=2\int_0^1(x-1)^2\cos n\pi x\,dx=
\frac{2}{ n\pi}\left[(x-1)^2\sin n\pi x\lims01-2
\int_0^1(x-1)\sin n\pi x\,dx\right]\\
&=\frac{4}{ n^2\pi^2}\left[(x-1)\cos n\pi x\lims01
-\int_0^1\cos n\pi x\,dx\right]
=\frac{4}{ n^2\pi^2}-\frac{4}{ n^3\pi^3}\sin n\pi x\lims01=
\frac{4}{ n^2\pi^2};
\end{align*}
$C(x)=\dst\frac{1}{3}+\frac{4}{\pi^2}\sum_{n=1}^\infty\frac{1}{ n^2}\cos
n\pi x$.
From Example~12.3.3,
$u(x,y)=\dst\frac{y}{3}+\frac{4}{\pi^3}\sum_{n=1}^\infty\frac{\sinh n\pi y}{
n^3\cosh n\pi}\cos n\pi x$.



\item
Since $g(0)=g'(1)=0$, and $g''(y)=6(1-2y)$,
Theorem~11.3.5\part{d} implies that
\begin{align*}
\alpha_n&=-\dst\frac{48}{(2n-1)^2\pi^2}\int_0^1(1-2y)
\sin\frac{(2n-1)\pi y}{2}\,dy\\
&=\dst\frac{96}{(2n-1)^3\pi^3}\left[(1-2y)\cos\frac{(2n-1)\pi
y}{2}\lims01+2\int_0^1\cos\frac{(2n-1)\pi y}{2}\,dy\right]\\
&=\dst\frac{96}{(2n-1)^3\pi^3}\left[-1+\frac{4}{(2n-1)\pi}
\sin\frac{(2n-1)\pi y}{2}\lims01\right]\\
&=\dst-\frac{96}{(2n-1)^3\pi^3}\left[1+(-1)^n\frac{4}{(2n-1)\pi}\right];
\end{align*}
$S_M(y)
=\dst-\frac{96}{\pi^3}\sum_{n=1}^\infty
\frac{1}{(2n-1)^3}\left[1+(-1)^n\frac{4}{(2n-1)\pi}\right]\sin\frac{(2n-1)\pi
y}{2}$.
From Example~12.3.5,
\[
u(x,y)
=\dst-\frac{96}{\pi^3}\sum_{n=1}^\infty
\left[1+(-1)^n\frac{4}{(2n-1)\pi}\right]
\frac{\cosh(2n-1)\pi(x-2)/2}{(2n-1)^3\cosh2(2n-1)\pi/2}
\sin\frac{(2n-1)\pi y}{2}.
\]



\item
From Example~11.3.8.3,
\[
S_M(y)
=\dst\frac{96}{\pi^3}\sum_{n=1}^\infty
\frac{1}{(2n-1)^3}\left[3+(-1)^n\frac{4}{(2n-1)\pi}\right]
\sin\frac{(2n-1)\pi y}{2}.
\]
From Example~12.3.5,
\[
u(x,y)
=\dst\frac{96}{\pi^3}\sum_{n=1}^\infty
\left[3+(-1)^n\frac{4}{(2n-1)\pi}\right]
\frac{\cosh(2n-1)\pi(x-3)/2}{(2n-1)^3\cosh3(2n-1)\pi/2}
\sin\frac{(2n-1)\pi y}{2}.
\]



\item
\begin{align*}
c_n&=\frac{2}{3}\int_0^3(3y-y^2)\cos\frac{(2n-1)\pi y}{6}\,dy\\
&=\frac{4}{(2n-1)\pi}\left[(3y-y^2)\sin\frac{(2n-1)\pi
y}{6}\lims03-\int_0^3
(3-2y)\sin\frac{(2n-1)\pi y}{6}\,dy\right]\\
&=\frac{24}{(2n-1)^2\pi^2}\left[(3-2y)\cos\frac{(2n-1)\pi
y}{6}\lims03+2
\int_0^3\cos\frac{(2n-1)\pi y}{6}\,dy\right]\\
&=-\frac{72}{(2n-1)^2\pi^2}+
\frac{288}{(2n-1)^3\pi^3}\sin\frac{(2n-1)\pi y}{6}\lims03
=\frac{288}{(2n-1)^3\pi^3}\sin\frac{(2n-1)\pi y}{2}\\
&=-\frac{72}{(2n-1)^2\pi^2}
+(-1)^{n-1}\frac{288}{(2n-1)^3\pi^3};
\end{align*}
$C_M(y)=-\dst\frac{72}{\pi^2}\sum_{n=1}^\infty\frac{1}{(2n-1)^2}\left[
1+\frac{4(-1)^n}{(2n-1)\pi}\right]\cos\frac{(2n-1)\pi
y}{6}$.
From Example~12.3.7,
\[
u(x,y)=-\dst\frac{432}{\pi^3}\sum_{n=1}^\infty
\left[1+\frac{4(-1)^n}{(2n-1)\pi}\right]
\frac{\cosh(2n-1)\pi x/6}{(2n-1)^3\sinh(2n-1)\pi/3}
\cos\frac{(2n-1)\pi
y}{6}.
\]


\item
Since $g'(0)=g(1)=0$ and $g''(y)=-6y$,
Theorem~11.3.5\part{c}  implies that
\begin{align*}
\alpha_n&=
\frac{48}{(2n-1)^2\pi^2}
\int_0^1y\cos\frac{(2n-1)\pi y}{2}\,dy\\
&=\frac{96}{(2n-1)^3\pi^3}\left[y\sin\frac{(2n-1)\pi y}{2}
\lims01-\int_0^1\sin\frac{(2n-1)\pi y}{2}\,dy\right]\\
&=\frac{96}{(2n-1)^3\pi^3}\left[(-1)^{n+1}
+\frac{2}{(2n-1)\pi}\cos\frac{(2n-1)\pi y}{2}\lims01\right]\\
&=-\frac{96}{(2n-1)^3\pi^3}\left[(-1)^n
+\frac{2}{(2n-1)\pi}\right]\,dy;
\end{align*}
$C_M(y)=-\dst\frac{96}{\pi^3}\sum_{n=1}^\infty\frac{1}{(2n-1)^3}\left[
(-1)^n+\frac{2}{(2n-1)\pi}\right]\cos\frac{(2n-1)\pi
y}{2}$.
From Example~12.3.7,
\[
u(x,y)=-\dst\frac{192}{\pi^4}\sum_{n=1}^\infty
\frac{\cosh(2n-1)\pi x/2}{(2n-1)^4\sinh(2n-1)\pi/2}
\left[(-1)^n+\frac{2}{(2n-1)\pi}\right]
\cos\frac{(2n-1)\pi y}{2}.
\]



\item
The boundary conditions require products $v(x,y)=X(x)Y(y)$
such that
(A)  $X''+\lambda X=0$, $X'(0)=0$, $X'(a)=0$, and
(B)  $Y''-\lambda Y=0$, $Y(0)=1$, $Y(b)=0$.
From Theorem~11.1.3, the eigenvalues of (A) are
$\lambda=0$, with associated eigenfunction $X_0=1$, and
$\lambda_n=\dst\frac{n^2\pi^2}{ a^2}$, with associated eigenfunctions
$Y_n=\dst\cos\frac{n\pi x}{ a}$, $n=1$, $2$, $3$,\dots.
Substituting $\lambda=0$ into (B) yields $Y_0''=0$,
$Y_0(0)=1$, $Y_0(b)=0$,
so $Y_0(y)=\dst1-\frac{y}{ b}$.
Substituting $\lambda=\dst\frac{n^2\pi^2}{ a^2}$  into (B) yields
$Y_n''-({n^2\pi^2/a^2})Y_n=0$, $Y_n(0)=1$,  $Y_n(b)=0$, so
$Y_n=\dst\frac{\sinh n\pi(b-y)/a}{ \sinh n\pi b/a}$. Then
$v_n(x,y)=X_n(x)Y_n(y)=
\dst\frac{\sinh n\pi(b-y)/a}{\sinh n\pi
b/a}\cos\frac{n\pi x}{ a}$, so
$v_n(x,0)=\dst\cos\frac{n\pi x}{ a}$.
Therefore,$v_n$ is solution of the given problem  with
$f(x)=\dst\cos\frac{n\pi x}{ a}$. More generally,
 if $\alpha_0,\dots,\alpha_m$ are arbitrary constants,
then
$\dst
u_m(x,y)=\alpha_0\left(1-\frac{y}{ b}\right)+\sum_{n=1}^m\alpha_n\frac{\sinh
n\pi(b-y)/ a}{\sinh n\pi b/a}\cos
\frac{n\pi x}{ a}$
 is a solution of the given problem with
$\dst f(x)=\alpha_0+\sum_{n=1}^m\alpha_n\cos\frac{n\pi x}{ a}$.
Therefore, if $f$ is an arbitrary piecewise smooth function on
$[0,a]$  we define the formal solution of the given problem  to be
$\dst
u(x,y)=\alpha_0\left(1-\frac{y}{
b}\right)+\sum_{n=1}^\infty\alpha_n\frac{\sinh
n\pi(b-y)/ a}{\sinh n\pi b/a}\cos
\frac{n\pi x}{ a}$,
where
$\dst C(x)=\alpha_0+\sum_{n=1}^\infty \alpha_n\cos\frac{n\pi x}{ a}$
is the Fourier  cosine series of $f$ on $[0,a]$; that is,
$\dst\alpha_0=\frac{1}{ a}\int_0^a f(x)\,dx$ and
$\alpha_n=\dst\frac{2}{ a}\int_0^a f(x)\cos
\frac{n\pi x}{ a}\,dx$, $n\ge1$.


Now consider the special case.
$\dst \alpha_0=\frac{1}{ 2}\int_0^2(x^4-4x^3+4x^2)\,dx=\frac{1}{
2}\left(\frac{x^5}{5}-x^4+\frac{4x^3}{3}
\right)\lims02=\frac{6}{5}$.
Since  $f'(0)=f'(2)=0$ and  $f'''(x)=12(2x-2)$,
Theorem~11.3.5\part{a} implies that
\begin{align*}
\alpha_n&=
\frac{96}{ n^3\pi^3}\int_0^2(2x-2)\sin\frac{n\pi x}{ 2}\,dx
=-\frac{192}{ n^4\pi^4}\left[(2x-2)\cos\frac{n\pi x}{ 2}\lims02
-2\int_0^2\cos\frac{n\pi x}{ 2}\,dx\right]\\
&=-\frac{192}{ n^4\pi^4}\left[(-1)^n2+2\right]+
\frac{768}{ n^5\pi^5}\sin\frac{n\pi x}{ 2}\lims02=
-\frac{384}{ n^4\pi^4}\left[1+(-1)^n\right]\\
&=\begin{cases}
0&\text{ if }n=2m-1,\\
\dst-\frac{48}{ m^4\pi^4}&\text{ if }n=2m,
\end{cases}
n\ge1.
\end{align*}
$\dst C(x)=\frac{8}{15}-\frac{48}{\pi^4}\sum_{n=1}^\infty\frac{1}{
n^4}\cos n\pi x$;
$\dst u(x,y)=\frac{8(1-y)}{15}-\frac{48}{\pi^4}\sum_{n=1}^\infty\frac{1}{
n^4}\frac{\sinh n\pi(1-y)}{\sinh n\pi}\cos n\pi x$.


\item
The boundary conditions require products $v(x,y)=X(x)Y(y)$
such that
(A)  $X''+\lambda X=0$, $X(0)=0$, $X'(a)=0$, and
(B)  $Y''-\lambda Y=0$, $Y(0)=1$, $Y(b)=0$.
From Theorem~11.1.4, the eigenvalues of (A) are
$\lambda_n=\dst\frac{(2n-1)^2\pi^2}{4a^2}$, with associated
eigenfunctions
$Y_n=\dst\sin\frac{(2n-1)\pi x}{2a}$, $n=1$, $2$, $3$,\dots.
Substituting $\lambda=\dst\frac{(2n-1)^2\pi^2}{4a^2}$  into (B) yields
$Y_n''-((2n-1)^2\pi^2/4a^2)Y_n=0$, $Y_n(0)=1$, $Y_n(b)=0$, so
$Y_n=\dst\frac{\sinh (2n-1)\pi(b-y)/2a}{\sinh(2n-1) \pi b/2a}$. Then
$v_n(x,y)=X_n(x)Y_n(y)=\dst\frac{\sinh(2n-1)\pi(b-y)/2a}{\sinh(2n-1)\pi
b/2a}\sin\frac{(2n-1)\pi x}{2a}$, so
$v_n(x,0)=\dst\sin\frac{(2n-1)\pi x}{2a}$.
Therefore,$v_n$ is solution of the given problem  with
$f(x)=\dst\sin\frac{(2n-1)\pi x}{2a}$. More generally,
 if $\alpha_1,\dots,\alpha_m$ are arbitrary constants,
then
$\dst
u_m(x,y)=\sum_{n=1}^m\alpha_n\frac{\sinh(2n-1)\pi(b-y)/2a
}{\sinh(2n-1)\pi b/2a}\cos
\frac{(2n-1)\pi x}{2a}$
 is a solution of the given problem with
$\dst f(x)=\sum_{n=1}^m\alpha_n\sin\frac{(2n-1)\pi x}{2a}$.
Therefore, if $f$ is an arbitrary piecewise smooth function on
$[0,a]$  we define the formal solution of the given problem  to be
$\dst
u(x,y)=\sum_{n=1}^\infty\alpha_n\frac{\sinh(2n-1)\pi(b-y)/2a
}{\sinh(2n-1)\pi b/2a}\sin
\frac{(2n-1)\pi x}{2a}$, where
$\dst S_m(x)=\sum_{n=1}^\infty \alpha_n\sin\frac{(2n-1)\pi x}{2a}$
is the mixed Fourier  sine series of $f$ on $[0,a]$; that is,
$\alpha_n=\dst\frac{2}{ a}\int_0^a f(x)\sin\frac{(2n-1)\pi x}{2a}$.


Now consider the special case.
Since $f(0)=f'(L)=0$ and $f''(x)=-2$,
Theorem~11.3.5\part{d} implies that
\begin{gather*}
\alpha_n=
\frac{48}{(2n-1)^2\pi^2}\int_0^3\sin\frac{(2n-1)\pi x}{6}\,dx=
-\frac{288}{(2n-1)^3\pi^3}\cos\frac{(2n-1)\pi x}{6}\lims03
=\frac{288}{(2n-1)^3\pi^3};
\\
S_M(x)=\dst
\frac{288}{\pi^3}\sum_{n=1}^\infty\frac{1}{
(2n-1)^3}\sin\frac{(2n-1)\pi x}{6};
\\
u(x,y)=\dst
\frac{288}{\pi^3}\sum_{n=1}^\infty\frac{\sinh(2n-1)\pi(2-y)/6}{
(2n-1)^3\sinh(2n-1)\pi/3}\sin\frac{(2n-1)\pi x}{6}.
\end{gather*}

\item
The boundary conditions require products $v(x,y)=X(x)Y(y)$
such that
(A)  $X''+\lambda X=0$, $X'(0)=0$, $X'(a)=0$, and
(B)  $Y''-\lambda Y=0$, $Y'(0)=0$, $Y(b)=1$.
From Theorem~11.1.3, the eigenvalues of (A) are
$\lambda=0$, with associated eigenfunction $X_0=1$, and
$\lambda_n=\dst\frac{n^2\pi^2}{ a^2}$, with associated eigenfunctions
$Y_n=\dst\cos\frac{n\pi x}{ a}$, $n=1$, $2$, $3$,\dots.
Substituting $\lambda=0$ into (B) yields $Y_0''=0$,
$Y_0'(0)=0$, $Y_0(b)=1$,
so $Y_0=1$.
Substituting $\lambda=\dst\frac{n^2\pi^2}{ a^2}$  into (B) yields
$Y_n''-({n^2\pi^2/a^2})Y_n=0$, $Y_n'(0)=0$,  $Y_n(b)=1$, so
$Y_n=\dst\frac{\cosh n\pi y/a}{ \cosh n\pi b/a}$. Then
$v_n(x,y)=X_n(x)Y_n(y)=
\dst\frac{\sinh n\pi y/a}{\cosh n\pi
b/a}\cos\frac{n\pi x}{ a}$, so
$v_n(x,b)=\dst\cos\frac{n\pi x}{ a}$.
Therefore,$v_n$ is solution of the given problem  with
$f(x)=\dst\cos\frac{n\pi x}{ a}$. More generally,
 if $\alpha_0,\dots,\alpha_m$ are arbitrary constants,
then
$\dst u_m(x,y)=\alpha_0+\sum_{n=1}^m\alpha_n\frac{\cosh n\pi y/a
}{\cosh n\pi b/a}\cos
\frac{n\pi x}{ a}$
 is a solution of the given problem with
$\dst f(x)=\alpha_0+\sum_{n=1}^m\alpha_n\cos\frac{n\pi x}{ a}$.
Therefore, if $f$ is an arbitrary piecewise smooth function on
$[0,a]$  we define the formal solution of the given problem  to be
$\dst u(x,y)=\alpha_0+\sum_{n=1}^\infty\alpha_n\frac{\cosh n\pi y/a
}{\cosh n\pi b/a}\cos
\frac{n\pi x}{ a}$,
where
$\dst C(x)=\alpha_0+\sum_{n=1}^\infty \alpha_n\cos\frac{n\pi x}{ a}$
is the Fourier  cosine series of $f$ on $[0,a]$; that is,
$\dst\alpha_0=\frac{1}{ a}\int_0^a f(x)\,dx$ and
$\alpha_n=\dst\frac{2}{ a}\int_0^a f(x)\cos
\frac{n\pi x}{ a}\,dx$, $n\ge1$.

Now consider the special case.
\[
\dst \alpha_0=\frac{1}{ \pi}\int_0^\pi(x^4-2\pi x^3+\pi^2x^2)\,dx=\frac{1}{
\pi}\left(\frac{x^5}{5}-\frac{\pi x^4}{2}+\frac{\pi^2x^3}{3}
\right)\lims0\pi=\frac{\pi^4}{30}.
\]
Since  $f'(0)=f'(\pi)=0$ and  $f'''(x)=12(2x-\pi)$,
Theorem~11.3.5\part{a} implies that
\begin{align*}
\alpha_n&=
\frac{24}{ n^3\pi}\int_0^\pi(2x-\pi)\sin nx\,dx
=-\frac{24}{ n^4\pi}\left[(2x-\pi)\cos nx\lims0\pi
-2\int_0^\pi\cos nx\,dx\right]\\
&=-\frac{24}{ n^4\pi}\left[(-1)^n\pi+\pi\right]+
\frac{48}{ n^5\pi}\sin nx\lims0\pi=
-\frac{24}{ n^4}\left[1+(-1)^n\right]\\
&=\begin{cases}
0&\text{ if }n=2m-1,\\
\dst-\frac{3}{ m^4}&\text{ if }n=2m,
\end{cases}
n\ge1;
\end{align*}
$\dst C(x)=\frac{\pi^4}{30}-3\sum_{n=1}^\infty\frac{1}{
n^4}\cos 2nx$;
$\dst u(x,y)=\frac{\pi^4}{30}-3\sum_{n=1}^\infty\frac{1}{
n^4}\frac{\cosh2ny}{\cos2n}\cos 2nx$


\item
The boundary conditions require products $v(x,y)=X(x)Y(y)$
such that
(A)  $X''-\lambda X=0$, $X'(0)=0$, $X(a)=1$, and
(B)  $Y''+\lambda Y=0$, $Y(0)=0$, $Y(b)=0$.
From Theorem~11.1.2, the eigenvalues of (B) are
$\lambda_n=\dst\frac{n^2\pi^2}{ b^2}$, with associated eigenfunctions
$Y_n=\dst\sin\frac{n\pi y}{ b}$, $n=1$, $2$, $3$,\dots.
Substituting $\lambda=\dst\frac{n^2\pi^2}{ b^2}$  into (A) yields
$X_n''-({n^2\pi^2/b^2})X_n=0$, $X_n'(0)=0$, $X_n(a)=1$, so
$X_n=\dst\frac{\cosh n\pi x/b}{\cosh n\pi a/b}$. Then
$v_n(x,y)=X_n(x)Y_n(y)=\dst\frac{\cosh n\pi x/b}{\cosh n\pi
a/b}\sin\frac{n\pi y}{ b}$, so
$v_n(a,y)=\dst\sin\frac{n\pi y}{ b}$.
Therefore,$v_n$ is solution of the given problem  with
$g(y)=\dst\sin\frac{n\pi y}{ b}$. More generally,
 if $\alpha_1,\dots,\alpha_m$ are arbitrary constants,
then
$\dst u_m(x,y)=\sum_{n=1}^m\alpha_n\frac{\cosh n\pi x/b}{
\cosh n\pi a/b}
\sin\frac{n\pi y}{ b}$
 is a solution of the given problem with
$\dst g(y)=\sum_{n=1}^m\alpha_n\sin\frac{n\pi y}{ b}$.
Therefore, if $g$ is an arbitrary piecewise smooth function on
$[0,b]$  we define the formal solution of the given problem  to be
$\dst u(x,y)=\sum_{n=1}^\infty\alpha_n\frac{\cosh n\pi x/b}{
\cosh n\pi a/b}
\sin\frac{n\pi y}{ b}$, where
$\dst S(y)=\sum_{n=1}^\infty \alpha_n\sin\frac{n\pi y}{ b}$
is the Fourier  sine series of $g$ on $[0,b]$; that is,
$\dst\alpha_n=\frac{2}{ b}\int_0^b g(y)\sin
\frac{n\pi y}{ b}\,dy$.


Now consider the special case.
Since $g(0)=g(1)=g''(0)=g''(L)=0$ and $f^{(4)}(y)=24$,
Theorem~11.3.5\part{b} and Exercise~35\part{b} of
Section~11.3 imply that
\begin{align*}
\alpha_n&=
\frac{48}{ n^4\pi^4}\int_0^1\sin n\pi y\,dy
=-\frac{48}{ n^5\pi^5}\cos n\pi y\lims01\\
&=-\frac{48}{ n^5\pi^5}[(-1)^n-1]=
\begin{cases}
\dst\frac{96}{(2m-1)^5\pi^5}&\text{ if }n=2m-1\\
0&\text{ if } n=2m;
\end{cases}
\end{align*}
$S(y)=\dst\frac{96}{\pi^5}\sum_{n=1}^\infty\frac{1}{(2n-1)^5}\sin(2n-1)\pi
y$;
$u(x,y)=\dst\frac{96}{\pi^5}\sum_{n=1}^\infty\frac{\cosh(2n-1)\pi x
}{(2n-1)^5\cosh(2n-1)\pi}\sin(2n-1)\pi y$.


\item
The boundary conditions require products $v(x,y)=X(x)Y(y)$
such that
(A)  $X''-\lambda X=0$, $X'(0)=0$, $X'(a)=1$, and
(B)  $Y''+\lambda Y=0$, $Y(0)=0$, $Y(b)=0$.
From Theorem~11.1.2, the eigenvalues of (B) are
$\lambda_n=\dst\frac{n^2\pi^2}{ b^2}$, with associated eigenfunctions
$Y_n=\dst\sin\frac{n\pi y}{ b}$, $n=1$, $2$, $3$,\dots.
Substituting $\lambda=\dst\frac{n^2\pi^2}{ b^2}$  into (A) yields
$X_n''-({n^2\pi^2/b^2})X_n=0$, $X_n'(0)=0$, $X_n'(a)=1$, so
$X_n=\dst\frac{b}{ n\pi}\frac{\cosh n\pi x/b}{\sinh n\pi a/b}$. Then
$v_n(x,y)=X_n(x)Y_n(y)=\dst\frac{b}{ n\pi}\frac{\cosh n\pi x/b}{\sinh n\pi
a/b}\sin\frac{n\pi y}{ b}$, so
$\dst\frac{\partial v_n}{\partial x}(a,y)=\dst\sin\frac{n\pi y}{ b}$.
Therefore,$v_n$ is solution of the given problem  with
$g(y)=\dst\sin\frac{n\pi y}{ b}$. More generally,
 if $\alpha_1,\dots,\alpha_m$ are arbitrary constants,
then
$\dst u_m(x,y)=\frac{b}{ \pi}\sum_{n=1}^m\alpha_n
\frac{\cosh n\pi x/b}{ n \sinh n\pi a/b}
\sin\frac{n\pi y}{ b}$
 is a solution of the given problem with
$\dst g(y)=\sum_{n=1}^m\alpha_n\sin\frac{n\pi y}{ b}$.
Therefore, if $g$ is an arbitrary piecewise smooth function on
$[0,b]$  we define the formal solution of the given problem  to be
$\dst u(x,y)=\frac{b}{ \pi}\sum_{n=1}^\infty\alpha_n
\frac{\cosh n\pi x/b}{ n \sinh n\pi a/b}
\sin\frac{n\pi y}{ b}$, where
$\dst S(y)=\sum_{n=1}^\infty \alpha_n\sin\frac{n\pi y}{ b}$
is the Fourier  sine series of $g$ on $[0,b]$; that is,
$\dst\alpha_n=\frac{2}{ b}\int_0^b g(y)\sin
\frac{n\pi y}{ b}\,dy$.


Now consider the special case.
$\alpha_n=\dst\frac{1}{ 2}\left[\int_0^2y\sin\frac{n\pi y}{ 4}
+\int_2^4(4-y)\sin\frac{n\pi y}{ 4}\,dy\right]$;
\begin{align*}
\int_0^2y\sin\frac{n\pi y}{ 4}\,dy
&=-\frac{4}{ n\pi}\left[y\cos\frac{n\pi y}{ 4}\lims02-
\int_0^2\cos\frac{n\pi y}{ 4}\,dy\right]\\
&=-\frac{2}{ n\pi}\cos\frac{n\pi}{2}+\frac{4}{ n^2\pi^2}\sin\frac{n\pi y
}{ 4}\lims02
=-\frac{2}{
n\pi}\cos\frac{n\pi}{2}+\frac{4}{ n^2\pi^2}\sin\frac{n\pi}{2};\\
\int_0^2(4-y)\sin\frac{n\pi y}{ 4}\,dy&=-\frac{2}{
n\pi}\left[(4-y)\cos\frac{n\pi y}{ 4}\lims24+
\int_0^2\cos\frac{n\pi y}{ 4}\,dy\right]\\
&=\frac{2}{ n\pi}\cos\frac{n\pi}{2}-\frac{4}{ n^2\pi^2}\sin\frac{n\pi y
}{ 4}\lims24
=\frac{2}{
n\pi}\cos\frac{n\pi}{2}+\frac{4}{ n^2\pi^2}\sin\frac{n\pi}{2};\\
\alpha_n&=\dst\frac{16}{ n^2\pi^2}\sin\frac{n\pi}{2}=
\begin{cases}
(-1)^{m+1}\dst\frac{16}{(2m-1)^2\pi^2}&\text{if }n=2m-1\\
0&\text{if }n=2m;
\end{cases}
\\
S(y)&=\dst\frac{16}{\pi^2}\sum_{n=1}^\infty\frac{(-1)^{n+1}}{(2n-1)^2}
\sin\frac{(2n-1)\pi y}{ 4};
\\
u(x,y)&=\dst\frac{64}{\pi^3}\sum_{n=1}^\infty(-1)^{n+1}
\frac{\cosh(2n-1)\pi x/4}{(2n-1)^3\sinh(2n-1)\pi/4}
\sin\frac{(2n-1)\pi y}{ 4}.
\end{align*}


\item
The boundary conditions require products $v(x,y)=X(x)Y(y)$
such that
(A)  $X''-\lambda X=0$, $X'(0)=1$, $X(a)=0$, and
(B)  $Y''+\lambda Y=0$, $Y'(0)=0$, $Y'(b)=0$.
From Theorem~11.1.3, the eigenvalues of (B) are
$\lambda_0=0$, with associated eigenfunction $Y_0=1$, and
$\lambda_n=\dst\frac{n^2\pi^2}{ b^2}$, with associated eigenfunctions
$Y_n=\dst\cos\frac{n\pi y}{ b}$, $n=1$, $2$, $3$,\dots.
Substituting $\lambda_0=0$ into (A) yields $X_0''=0$, $X_0'(0)=1$
$X(a)=0$, so $X_0=x-a$.
Substituting $\lambda=\dst\frac{n^2\pi^2}{ b^2}$  into (A) yields
$X_n''-({n^2\pi^2/b^2})X_n=0$, $X_n'(0)=1$, $X_n(a)=0$, so
$X_n=\dst\frac{b}{ n\pi}\frac{\sinh n\pi(x-a)/b}{\cosh n\pi a/b}$.
Then
$v_n(x,y)=X_n(x)Y_n(y)=\dst\frac{b}{ n\pi}\frac{\sinh n\pi(x-a)/b}{\cosh
n\pi a/b}\cos\frac{n\pi y}{ b}$, so
$\dst\frac{\partial v_n}{\partial x}(0,y)=\dst\cos\frac{n\pi y}{ b}$.
Therefore,$v_n$ is solution of the given problem  with
$g(y)=\dst\cos\frac{n\pi y}{ b}$. More generally,
 if $\alpha_0,\dots,\alpha_m$ are arbitrary constants,
then
$\dst
u_m(x,y)=\alpha_0(x-a)+\frac{b}{\pi}\sum_{n=1}^m\alpha_n\frac{\sinh
n\pi(x-a)/b
}{ n\cosh n\pi a/b}
\cos\frac{n\pi y}{ b}$
 is a solution of the given problem with
$\dst g(y)=\sum_{n=1}^m\alpha_n\cos\frac{n\pi y}{ b}$.
Therefore, if $g$ is an arbitrary piecewise smooth function on
$[0,b]$  we define the formal solution of the given problem  to be
$\dst
u(x,y)=\alpha_0(x-a)+\frac{b}{\pi}\sum_{n=1}^\infty\alpha_n\frac{\sinh
n\pi(x-a)/b
}{ n\cosh n\pi a/b}
\cos\frac{n\pi y}{ b}$ where
$\dst C(y)=\alpha_0+\sum_{n=1}^\infty \alpha_n\cos\frac{n\pi y}{ b}$
is the Fourier cosine series of $g$ on $[0,b]$; that is,
$\dst\alpha_0=\frac{1}{ b}\int_0^bg(y)\cos\frac{n\pi y}{ b}\,dy$,
$\dst\alpha_n=\frac{2}{ b}\int_0^b g(y)\cos
\frac{n\pi y}{ b}\,dy$, $n\ge1$.

Now consider the special case.
From Example~11.3.1,
\begin{gather*}
C(y)=\dst\frac{\pi}{2}-\frac{4}{\pi}\sum_{n=1}^\infty\frac{1}{(2n-1)^2}
\cos(2n-1)y;
\\
u(x,y)=\dst\frac{\pi(x-2)}{2}-\frac{4}{\pi}\sum_{n=1}^\infty
\frac{\sinh(2n-1)(x-2)}{(2n-1)^3\cosh2(2n-1)}
\cos(2n-1)y.
\end{gather*}


\item
The boundary conditions require products $v(x,y)=X(x)Y(y)$
such that
(A)  $X''+\lambda X=0$, $X'(0)=0$, $X(a)=0$, and
(B)  $Y''-\lambda Y=0$, $Y(0)=1$, and $Y$ is bounded.
From Theorem~11.1.5, the eigenvalues of (A) are
$\lambda_n=\dst\frac{(2n-1)^2\pi^2}{4a^2}$, with associated
eigenfunctions
$Y_n=\dst\cos\frac{(2n-1)\pi x}{2a}$, $n=1$, $2$, $3$,\dots.
Substituting $\lambda=\dst\frac{(2n-1)^2\pi^2}{4a^2}$  into (B) yields
$Y_n''-((2n-1)^2\pi^2/4a^2)Y_n=0$, $Y_n(0)=1$, so
$Y_n=\dst e^{-(2n-1)\pi y/2a}$. Then
$v_n(x,y)=X_n(x)Y_n(y)=
\dst e^{-(2n-1)\pi y/2a}\cos\frac{(2n-1)\pi x}{2a}$, so
$v_n(x,0)=\dst\cos\frac{(2n-1)\pi x}{2a}$.
Therefore,$v_n$ is solution of the given problem  with
$f(x)=\dst\cos\frac{(2n-1)\pi x}{2a}$. More generally,
 if $\alpha_1,\dots,\alpha_m$ are arbitrary constants,
then
$u_m(x,y)=\dst\sum_{n=1}^m\alpha_n e^{-(2n-1)\pi y/2a}\cos
\frac{(2n-1)\pi x}{2a}$
 is a solution of the given problem with
$\dst f(x)=\sum_{n=1}^m\alpha_n\cos\frac{(2n-1)\pi x}{2a}$.
Therefore, if $f$ is an arbitrary piecewise smooth function on
$[0,a]$  we define the formal solution of the given problem  to be
$u(x,y)=\dst\sum_{n=1}^\infty\alpha_n e^{-(2n-1)\pi y/2a}\cos
\frac{(2n-1)\pi x}{2a}$, where
$\dst C_m(x)=\sum_{n=1}^\infty \alpha_n\cos\frac{(2n-1)\pi x}{2a}$
is the mixed Fourier  cosine series of $f$ on $[0,a]$; that is,
$\alpha_n=\dst\frac{2}{ a}\int_0^a f(x)\cos\frac{(2n-1)\pi x}{2a}$.




Now consider the special case.
Since $f'(0)=f(L)=0$ and $f''(x)=-2$,
Theorem~11.3.5\part{d} implies that
\begin{align*}
\alpha_n&=
\frac{16L}{(2n-1)^2\pi^2}
\int_0^3\cos\frac{(2n-1)\pi x}{6}\,dx \\
&=\frac{288}{(2n-1)^3\pi^3}\sin\frac{(2n-1)\pi x}{6}\lims03
=(-1)^{n+1}\frac{288}{(2n-1)^3\pi^3};\\
C_M(x)&=-\dst
\frac{288}{\pi^3}\sum_{n=1}^\infty\frac{(-1)^n}{
(2n-1)^3}\cos\frac{(2n-1)\pi x}{6};
\\
u(x,y)&=-\dst
\frac{288}{\pi^3}\sum_{n=1}^\infty\frac{(-1)^n}{
(2n-1)^3}e^{-(2n-1)\pi y/6}\cos\frac{(2n-1)\pi x}{6}.
\end{align*}


\item
The boundary conditions require products $v(x,y)=X(x)Y(y)$
such that
(A)  $X''+\lambda X=0$, $X(0)=0$, $X(a)=0$, and
(B)  $Y''-\lambda Y=0$, $Y'(0)=1$, and $Y$ is bounded.
From Theorem~11.1.2, the eigenvalues of (A) are
$\lambda_n=\dst\frac{n^2\pi^2}{ a^2}$, with associated eigenfunctions
$Y_n=\dst\sin\frac{n\pi x}{ a}$, $n=1$, $2$, $3$,\dots.
Substituting $\lambda=\dst\frac{n^2\pi^2}{ a^2}$  into (B) yields
$Y_n''-({n^2\pi^2/a^2})Y_n=0$, $Y_n'(0)=1$, so
$Y_n=-\dst \frac{a}{ n\pi}e^{-n\pi y/a}$. Then
$v_n(x,y)=X_n(x)Y_n(y)=-\dst \frac{a}{ n\pi}e^{-n\pi y/a}
\sin\frac{n\pi x}{ a}$, so
$\dst\frac{\partial v_n}{\partial y}(x,0)=\dst\sin\frac{n\pi x}{ a}$.
Therefore,$v_n$ is solution of the given problem  with
$f(x)=\dst\sin\frac{n\pi x}{ a}$. More generally,
 if $\alpha_1,\dots,\alpha_m$ are arbitrary constants,
then
$u_m(x,y)=-\dst\frac{a}{ \pi}\sum_{n=1}^m\frac{\alpha_n}{ n} e^{-n\pi
y/a}\sin \frac{n\pi x}{ a}$
 is a solution of the given problem with
$\dst f(x)=\sum_{n=1}^m\alpha_n\sin\frac{n\pi x}{ a}$.
Therefore, if $f$ is an arbitrary piecewise smooth function on
$[0,a]$  we define the formal solution of the given problem  to be
$u(x,y)=-\dst\frac{a}{ \pi}\sum_{n=1}^\infty\frac{\alpha_n}{ n} e^{-n\pi
y/a}\sin \frac{n\pi x}{ a}$, where
$\dst C(x)=\sum_{n=1}^\infty \alpha_n\sin\frac{n\pi x}{ a}$
is the Fourier  sine series of $f$ on $[0,a]$; that is,
$\alpha_n=\dst\frac{2}{ a}\int_0^a f(x)\sin
\frac{n\pi x}{ a}\,dx$.

Now consider the special case.
Since $f(0)=f(\pi)=0$ and $f''(x)=2\pi-6x$,
Theorem~11.3.5\part{b} implies that
\begin{align*}
\alpha_n&=
-\frac{2}{ n^2\pi}\int_0^\pi(2\pi-6x)\sin nx\,dx
=\frac{2}{ n^3\pi}\left[(2\pi-6x)\cos nx\lims0\pi
+6\int_0^\pi\cos nx\,dx\right]\\
&=-\frac{4}{ n^3}\left[1+(-1)^n2\right]+\frac{12}{
n^4\pi}\sin nx\lims0\pi
=-\frac{4}{ n^3}\left[1+(-1)^n2\right];
\end{align*}
$S(x)=\dst-4\sum_{n=1}^\infty\frac{[1+(-1)^n2]}{ n^3}
\sin nx$;
$u(x)=\dst4\sum_{n=1}^\infty\frac{[1+(-1)^n2]}{ n^4}
e^{-ny}\sin nx$;



\item
The boundary conditions require products $v(x,y)=X(x)Y(y)$
such that
(A)  $X''+\lambda X=0$, $X(0)=0$, $X'(a)=0$, and
(B)  $Y''-\lambda Y=0$, $Y'(0)=1$, and $Y$ is bounded.
From Theorem~11.1.4, the eigenvalues of (A) are
$\lambda_n=\dst\frac{(2n-1)^2\pi^2}{4a^2}$, with associated
eigenfunctions
$Y_n=\dst\sin\frac{(2n-1)\pi x}{2a}$, $n=1$, $2$, $3$,\dots.
Substituting $\lambda=\dst\frac{(2n-1)^2\pi^2}{4a^2}$  into (B) yields
$Y_n''-((2n-1)^2\pi^2/4a^2)Y_n=0$, $Y_n'(0)=1$, so
$Y_n=-\dst \frac{2a}{(2n-1)\pi}e^{-(2n-1)\pi y/2a}$. Then
$v_n(x,y)=X_n(x)Y_n(y)=-\dst \frac{2a}{(2n-1)\pi}e^{-(2n-1)\pi y/2a}
\sin\frac{(2n-1)\pi x}{2a}$, so
$\dst\frac{\partial v_n}{\partial y}(x,0)=\dst\sin\frac{(2n-1)\pi x}{2a}$.
Therefore,$v_n$ is solution of the given problem  with
$f(x)=\dst\sin\frac{(2n-1)\pi x}{2a}$. More generally,
 if $\alpha_1,\dots,\alpha_m$ are arbitrary constants,
then
$u_m(x,y)=\dst-\frac{2a}{\pi}\sum_{n=1}^m\frac{\alpha_n}{2n-1}e^{-(2n-1)\pi
y/2a}\sin
\frac{(2n-1)\pi x}{2a}$
 is a solution of the given problem with
$\dst f(x)=\sum_{n=1}^m\alpha_n\sin\frac{(2n-1)\pi x}{2a}$.
Therefore, if $f$ is an arbitrary piecewise smooth function on
$[0,a]$  we define the formal solution of the given problem  to be
$u(x,y)=\dst-\frac{2a}{\pi}\sum_{n=1}^\infty\frac{\alpha_n}{2n-1}e^{-(2n-1)\pi
y/2a}\sin
\frac{(2n-1)\pi x}{2a}$,where
$\dst S_m(x)=\sum_{n=1}^\infty \alpha_n\sin\frac{(2n-1)\pi x}{2a}$
is the mixed Fourier  sine series of $f$ on $[0,a]$; that is,
$\alpha_n=\dst\frac{2}{ a}\int_0^a f(x)\sin\frac{(2n-1)\pi x}{2a}$.

Now consider the special case.
\begin{align*}
\alpha_n&=\frac{2}{ 5}\int_0^5(5x-x^2)\sin\frac{(2n-1)\pi x}{10}\,dx\\
&=-\frac{4}{(2n-1)\pi}\left[(5x-x^2)\cos\frac{(2n-1)\pi
x}{10}\lims05-
\int_0^5
(5-2x)\cos\frac{(2n-1)\pi x}{10}\,dx\right]\\
&=\frac{40}{(2n-1)^2\pi^2}\left[(5-2x)\sin\frac{(2n-1)\pi
x}{10}\lims05+2
\int_0^5\sin\frac{(2n-1)\pi x}{10}\,dx\right]\\
&=(-1)^n\frac{200}{(2n-1)^2\pi^2}-
\frac{800}{(2n-1)^3\pi^3}\cos\frac{(2n-1)\pi x}{10}\lims05\\
&=(-1)^n\frac{200}{(2n-1)^2\pi^2}
+\frac{800}{(2n-1)^3\pi^3};\\
S_M(x)&=\dst\frac{200}{\pi^2}\sum_{n=1}^\infty\frac{1}{(2n-1)^2}\left[
(-1)^n+\frac{4}{(2n-1)\pi}\right]\sin\frac{(2n-1)\pi
x}{10};
\\
u(x,y)&=-\dst\frac{2000}{\pi^3}\sum_{n=1}^\infty\frac{1}{(2n-1)^3}\left[
(-1)^n+\frac{4}{(2n-1)\pi}\right]e^{-(2n-1)\pi y/10}\sin\frac{(2n-1)\pi
x}{10}.
\end{align*}


\item
Solving  BVP$(1,1,1,1)(f_0,0,0,0)$ requires
products $X(x)Y(y)$ such that
\[
X''+\lambda X=0,\quad X'(0)=0,\quad X'(a)=0;
\quad Y''-\lambda Y=0,\quad  Y'(b)=0,\quad Y'(0)=1.
\]
Hence, $X_n=\dst\cos\frac{n\pi x}{ a}$,
$Y_n=-\dst\frac{a\cosh n\pi(y-b)/a}{ n\pi\sinh n\pi b/a}$,
 and
$\dst c_1-\frac{a}{\pi}\sum_{n=1}^\infty A_n\frac{\cosh n\pi(y-b)/a}{
n\pi\sinh n\pi b/a}
\cos\frac{n\pi x}{ a}$
is a formal solution of  BVP$(1,1,1,1)(f_0,0,0,0)$ if $c_1$
is any constant  and $\dst\sum_{n=1}^\infty A_n\cos\frac{n\pi x}{ a}$
is the Fourier cosine expansion of $f_0$ on $[0,a]$, which is possible
if and only if $\dst\int_0^af_0(x)\,dx=0$.

Similarly,
$\dst c_2+\frac{a}{\pi}\sum_{n=1}^\infty B_n\frac{\cosh n\pi y/a}{
n\pi\sinh n\pi b/a}
\cos\frac{n\pi x}{ a}$
is a formal solution of  BVP$(1,1,1,1)(0,f_1,0,0)$ if $c_2$
is any constant  and $\dst\sum_{n=1}^\infty B_n\cos\frac{n\pi x}{ a}$
is the Fourier cosine expansion of $f_1$ on $[0,a]$, which is possible
if and only if $\dst\int_0^af_1(x)\,dx=0$.

 Interchanging $x$ and $y$ and $a$ and $b$ shows that
$\dst c_3-\frac{b}{\pi}\sum_{n=1}^\infty C_n\frac{\cosh n\pi(x-a)/b}{
n\pi\sinh n\pi a/b}
\cos\frac{n\pi y}{ b}$
is a formal solution of  BVP$(1,1,1,1)(0,0,g_0,0)$ if $c_3$
is any constant  and $\dst\sum_{n=1}^\infty C_n\cos\frac{n\pi y}{ b}$
is the Fourier cosine expansion of $g_0$ on $[0,b]$, which is
possible if and only if $\dst\int_0^bg_0(x)\,dx=0$, and
$\dst c_4+\frac{b}{\pi}\sum_{n=1}^\infty D_n\frac{\cosh n\pi x/b}{
n\pi\sinh n\pi a/b}
\cos\frac{n\pi y}{ b}$
is a formal solution of  BVP$(1,1,1,1)(0,0,0,g_1)$ if $c_4$
is any constant  and $\dst\sum_{n=1}^\infty D_n\cos\frac{n\pi y}{ b}$
is the Fourier cosine expansion of $g_1$ on $[0,b]$, which is
possible if and only if $\dst\int_0^bg_1(x)\,dx=0$.

Adding the four solutions yields
\begin{align*}
u(x,y)=C&{}+
\frac{a}{\pi}\sum_{n=1}^\infty\frac{B_n\cosh n\pi y/a-A_n\cosh
n\pi(y-b)/a}{ n\sinh n\pi b/a}\cos\frac{n\pi x}{ a}\\
&{}+\frac{b}{\pi}\sum_{n=1}^\infty\frac{D_n\cosh n\pi
x/b-C_n\cosh n\pi(x-a)/b}{ n\sinh n\pi a/b}\cos\frac{n\pi y}{ b},
\end{align*}
where $C$  is an arbitrary constant.
\end{sectionSolutions}

\section{Laplace's Equation in Polar Coordinates}

\begin{sectionSolutions}
\item
 $v(r,\theta)=R(r)\Theta(\theta)$ where
(A) $r^2R''+rR'-\lambda R=0$ and
$\Theta''+\lambda\Theta=0$, $\Theta(0)=0$, $\Theta(\gamma)=0$.
From Theorem~11.1.2,
$\lambda_n=\dst\frac{n^2\pi^2}{\gamma^2}$,
 $\Theta_n=\dst\sin\frac
{n\pi\theta}{\gamma}$,
$n=1$,  $2$, $3$,\dots.
Substituting $\lambda=\dst\frac{n^2\pi^2}{\gamma^2}$ into (A) yields
the Euler equation
$r^2R_n''+rR_n'-\dst\frac{n^2\pi^2}{\gamma^2} R_n=0$
for $R_n$.
The indicial polynomial  is
$\dst
\left(s-\frac{n\pi}{\gamma}\right)
\left(s+\frac{n\pi}{\gamma}\right)$,
so
$\dst R_n=c_1r^{n\pi/\gamma}+c_2r^{-n\pi/\gamma}$,
by Theorem~7.4.3. We want $R_n(\rho)=1$ and $R_n(\rho_0)=0$, so
$R_n(r)=\dst\frac{\rho_0^{-n\pi/\gamma}r^{n\pi/\gamma}-\rho_0^{n\pi/\gamma}
r^{-n\pi/\gamma}}{\rho_0^{-n\pi/\gamma}\rho^{n\pi/\gamma}-\rho_0^{n\pi/\gamma}
\rho^{-n\pi/\gamma}}$;
\[
\dst v_n(r,\theta)=
\dst\frac{\rho_0^{-n\pi/\gamma}r^{n\pi/\gamma}-\rho_0^{n\pi/\gamma}
r^{-n\pi/\gamma}}{\rho_0^{-n\pi/\gamma}\rho^{n\pi/\gamma}-\rho_0^{n\pi/\gamma}
\rho^{-n\pi/\gamma}}
\sin\frac{n\pi\theta}{\gamma};
\]
$\dst u(r,\theta)=\sum_{n=1}^\infty\alpha_n
\dst\frac{\rho_0^{-n\pi/\gamma}r^{n\pi/\gamma}-\rho_0^{n\pi/\gamma}
r^{-n\pi/\gamma}}{\rho_0^{-n\pi/\gamma}\rho^{n\pi/\gamma}-\rho_0^{n\pi/\gamma}
\rho^{-n\pi/\gamma}}
\sin\frac{n\pi\theta}{\gamma}$,
where
$S(\theta)=\dst\sum_{n=1}^\infty\gamma_n\sin\frac{n\pi\theta}{\gamma}$
is the Fourier sine series if $f$ on $[0,\gamma]$; that is,
$\alpha_n=\dst\frac{1}{\gamma}\int_0^\gamma
f(\theta)\sin\frac{n\pi\theta}{\gamma}\,d\theta$, $n=1$, $2$,
$3$,\dots.


\item
 $v(r,\theta)=R(r)\Theta(\theta)$ where
(A) $r^2R''+rR'-\lambda R=0$ and
$\Theta''+\lambda\Theta=0$, $\Theta'(0)=0$, $\Theta(\gamma)=0$.
From Theorem~11.1.5,
$\lambda_n=\dst\frac{(2n-1)^2\pi^2}{4\gamma^2}$,
 $\Theta_n=\dst\cos\frac
{(2n-1)\pi\theta}{2\gamma}$,
$n=1$,  $2$, $3$,\dots.
Substituting $\lambda=\dst\frac{(2n-1)^2\pi^2}{4\gamma^2}$ into (A)
yields the Euler equation
$r^2R_n''+rR_n'-\dst\frac{(2n-1)^2\pi^2}{4\gamma^2} R_n=0$
for $R_n$.
The indicial polynomial  is
$\dst
\left(s-\frac{(2n-1)\pi}{2\gamma}\right)
\left(s+\frac{(2n-1)\pi}{2\gamma}\right)$,
so
$\dst R_n=c_1r^{(2n-1)\pi/2\gamma}+c_2r^{-(2n-1)\pi/2\gamma}$,
by Theorem~7.4.3. We want $R_n$  to be bounded as $r\to0+$ and
$R_n(\rho)=1$, so we take
$R_n(r)=\dst\frac{r^{(2n-1)\pi/2\gamma}}{\rho^{(2n-1)\pi/2\gamma}}$;
$v_n(r,\theta)=
\dst\frac{r^{(2n-1)\pi/2\gamma}}{\rho^{(2n-1)\pi/2\gamma}}
\cos\frac
{(2n-1)\pi\theta}{2\gamma}$;
$u(r,\theta)=\dst\sum_{n=1}^\infty\alpha_n
\dst\frac{r^{(2n-1)\pi/2\gamma}}{\rho^{(2n-1)\pi/2\gamma}}
\cos\frac
{(2n-1)\pi\theta}{2\gamma}$,
where
$C_M(\theta)=\dst\sum_{n=1}^\infty\alpha_n\cos\frac{(2n-1)\pi\theta}{2\gamma}$
is the mixed Fourier cosine series of $f$ on $[0,\gamma]$; that is,
$\alpha_n=\dst\frac{2}{\gamma}\int_0^\gamma
f(\theta)\cos\frac{(2n-1)\pi\theta}{2\gamma}\,d\theta$,\quad
 $n=1$, $2$, $3$,\dots.


\item
 $v(r,\theta)=R(r)\Theta(\theta)$ where
(A) $r^2R''+rR'-\lambda R=0$
and
$\Theta''+\lambda\Theta=0$, $\Theta'(0)=0$, $\Theta'(\gamma)=0$.
From Theorem~11.1.3,  $\lambda_0=0$, $\Theta_0=1$;
$\lambda_n=\dst\frac{n^2\pi^2}{\gamma^2}$,
 $\Theta_n=\dst\cos\frac
{n\pi\theta}{\gamma}$,
$n=1$,  $2$, $3$,\dots.
Substituting $\lambda=0$ into (A) yields the
equation $r^2R_0''+rR_0'=0$ for $R_0$;
$R_0=c_1+c_2\ln r$. Since we want $R_0$ to be bounded as $r\to0+$
and  $R_0(\rho)=1$, $R_0(r)=1$; therefore $v_0(r,\theta)=1$.

Substituting $\lambda=\dst\frac{n^2\pi^2}{\gamma^2}$ into (A) yields
the Euler equation
$r^2R_n''+rR_n'-\dst\frac{n^2\pi^2}{\gamma^2} R_n=0$
for $R_n$.
The indicial polynomial  is
$\dst
\left(s-\frac{n\pi}{\gamma}\right)
\left(s+\frac{n\pi}{\gamma}\right)$,
so
$\dst R_n=c_1r^{n\pi/\gamma}+c_2r^{-n\pi/\gamma}$,
by Theorem~7.4.3.
Since we want  $R_n$ to be bounded as $r\to0+$ and $R_n(\rho)=1$,
$R_n(r)=\dst\frac{r^{n\pi/\gamma}}{\rho^{n\pi/\gamma}}$;
$v_n(r,\theta)=
\dst\frac{r^{n\pi/\gamma}}{\rho^{n\pi/\gamma}}\cos\frac{n\pi\theta}{\gamma}$,
$n=1$, $2$, $3$,\dots;
$\dst u(r,\theta)=\alpha_0+\sum_{n=1}^\infty\alpha_n
\dst\frac{r^{n\pi/\gamma}}{\rho^{n\pi/\gamma}}\cos\frac{n\pi\theta}{\gamma}$,
where
$F(\theta)=\dst\alpha_0+\sum_{n=1}^\infty\alpha_n\cos\frac{n\pi\theta}{\gamma}$
is the Fourier cosine series of $f$ on $[0,\gamma]$; that is,
$\alpha_0=\dst\frac{1}{\gamma}\int_0^\gamma f(\theta)\,d\theta$
and
$\alpha_n=\dst\frac{2}{\gamma}\int_0^\gamma
f(\theta)\cos\frac{n\pi\theta}{\gamma}\,d\theta$, $n=1$, $2$,
$3$,\dots.
\end{sectionSolutions}

\chapter[running={Boundary Value Problems for 2nd Order ODEs}]{Boundary Value Problems for Second Order Ordinary Differential Equations}

\section{Boundary Value Problems}

\begin{sectionSolutions}
\item By  inspection, $y_{p}=-x$;
$y=-x+c_{1}e^{x}+c_{2}e^{-x}$;
$y(0)=-2 \implies c_{1}+c_{2}=-2$;
$y(1)=1 \implies -1+c_{1}e+c_{2}/e=1$;
$
\begin{bmatrix}
1&1 \\e&1/e
\end{bmatrix}
\begin{bmatrix}
c_{1}\\c_{2}
\end{bmatrix}=
\begin{bmatrix}
-2\\2
\end{bmatrix}$;
$c_{1}=\dst{\frac{2}{e-1}}$;
$c_{2}=\dst{\frac{2e}{1-e}}$;
$y=-x+ \dst{\frac{2\left(e^{x}-e^{-(x-1)}\right)}{e-1}}$

\item
By inspection, $y_{p}=-x$; $y=-x+c_{1}e^{x}+c_{2}e^{-x}$;
$y'=-1+c_{1}e^{x}-c_{2}e^{-x}$;
 $y(0)+y'(0)=3 \implies -1+2c_{1}=3$; $c_{1}=2$;
$y(1)-y'(1)=2\implies\dst{\frac{2c_{2}}{e}}=2$; $c_{2}=e$;
$y=-x+2e^{x}+e^{-(x-1)}$


\item
$y_{p}=Ax^{2}e^{x}$; $y_{p}'=A(x^{2}+2x e^{x})$;
$y_{p}''=A(x^{2} e^{x}+4x e^{x}+2 e^{x})$;
\noindent
$y_{p}''-2y_{p}'+y_{p}=2A e^{x}= e^{x}$ if $A=1$;
$y_{p}=x^{2} e^{x}$;
$y=(x^{2}+c_{1}+c_{2}x) e^{x}$;
$y'=(x^{2}+2x+c_{1}+c_{2}+c_{2}x) e^{x}$;
$B_{1}(y)=3$ and $B_{2}(y)=6e\implies$ $-c_{1}-2c_{2}=3$,
$2c_{1}+3c_{2}=24$; $c_{1}=13$; $c_{2}=-8$;
$y=(x^{2}-8x+13)e^{x}$.


\item
$B_{1}(y)=y(0)$; $B_{2}(y)=y(1)-y'(1)$. Let $y_{1}=x$, $y_{2}=1$;
$B_{1}(y_{1})=B_{2}(y_{1})=0$. By variation of parameters, if
$y_{p}=u_{1}x+u_{2}$ where $u_{1}'x+u_{2}'=0$
and $u_{1}'=F$, then $y_{p}''=F(x)$. Let $u_{1}'=F$,
$u_{2}'=-xF$; $u_{1}=-\dst{\int_{x}^{1}}F(t)\,dt$,
$u_{2}=-\dst{\int_{0}^{x}}tF(t)\,dt$;
$y_{p}=-x\dst{\int_{x}^{1}}F(t)\,dt-\dst{\int_{0}^{x}}tF(t)\,dt$;
$y_{p}'=-\dst{\int_{x}^{1}}F(t)\,dt$;
$y=y_{p}+c_{1}x+c_{2}$. Since
$B_{1}(y_{p})=0$, $B_{1}(x)=0$ and $B_{1}(1)=1$, $B_{1}(y)=0\implies
c_{2}=0$; hence, $y=y_{p}+c_{1}x$. Since
$B_{2}(y_{p})=-\dst{\int_{0}^{1}}tF(t)\,dt$ and $B_{2}(x)=0$,
$B_{2}(y)=0\implies \dst{\int_{0}^{1}}tF(t)\,dt=0.$
There is no solution if this conditions does not hold.
If it does hold, then the solutions are $y=y_{p}+c_{1}x$, with   $c_{1}$
arbitrary.


\item
\begin{alist}
\item
The condition is  $b-a\ne (k+1/2)\pi$ ($k=$ integer).
Let $y_{1}=\sin(x-a)$ and $y_{2}=\cos(x-b)$. Then $y_{1}(a)=y_{2}'(b)=0$
and
$\{y_{1},y_{2}\}$ is linearly independent if $b-a\ne(k+1/2)\pi$
($k=$ integer), since
 \[
\begin{vmatrix}
\sin(x-a)&\cos(x-b)\\\cos(x-a)&-\sin(x-b)
\end{vmatrix}=-\cos(b-a)\ne0.
\]
Now Theorem 13.1.2  implies that (A) has a unique solution for any
continuous
$F$ and constants $k_{1}$ and $k_{2}$.
If $y=u_{1}\sin(x-a)+u_{2}\cos(x-b)$ where
\begin{align*}
u_{1}'\sin(x-a)+u_{2}'\cos(x-b)&=0 \\
u_{1}'\cos(x-a)-u_{2}'\sin(x-b)&=F,
\end{align*}
then $y''+y=F$.
\begin{gather*}
u_{1}'=F(x)\frac{\cos(x-b)}{\cos(b-a)}, \quad
u_{2}'= -F(x)\frac{\sin(x-a)}{\cos(b-a)},
\\
u_{1}=-\frac{1}{\cos(b-a)}\int_{x}^{b}F(t)\cos(t-b)\,dt,\quad
u_{2}=-\frac{1}{\cos(b-a)}\int_{a}^{x}F(t)\sin(t-b)\,dt;
\\
y=-\frac{\sin(x-a)}{\cos(b-a)}\int_{x}^{b}F(t)\cos(t-b)\,dt
-\frac{\cos(x-b)}{\cos(b-a)}\int_{a}^{x}F(t)\sin(t-a)\,dt.
\end{gather*}

\item If $b-a=(k+1/2)\pi$  ($k=$ integer), then
$y_{1}=\sin(x-a)$ satisfies both boundary conditions $y(a)=0$
and $y'(b)=0$. Let $y_{2}=\cos(x-a)$. If
$y_{p}=u_{1}\sin(x-a)+u_{2}\cos(x)$   where
\begin{align*}
u_{1}'\sin(x-a)+u_{2}'\cos(x-a)&=0,\\
u_{1}'\cos(x-a)-u_{2}'\sin(x-a)&=F,
\end{align*}
then $y_{p}''+y_{p}=F$; $u_{1}'=F\cos(x-a)$; $u_{2}'=-F\sin(x-a)$;
\begin{gather*}
u_{1}=-\int_{x}^{b}F(t)\cos(t-a)\,dt, \quad
u_{2}=-\int_{a}^{x}F(t)\sin(t-a)\,dt;
\\
y_{p}=-\sin(x-a)\int_{x}^{b}F(t)\cos(t-a)\,dt
-\cos(x-a)\int_{a}^{x} F(t)\sin(t-a)\,dt;
\\
y_{p}'=-\cos(x-a)\int_{x}^{b}F(t)\cos(t-a)\,dt
+\sin(x-a)\int_{a}^{x}F(t)\sin(t-a)\,dt.
\end{gather*}
The general solution of $y''+y=F$  is
$y=y_{p}+c_{1}\sin(x-a)+c_{2}\cos(x-a)$.  Since $b-a=(k+1/2)\pi$,
$y'(b)=0\implies y_{p}'(b)=(-1)^{k}\dst{\int_{a}^{b}}F(t)\sin(t-a)\,dt=0$;
therefore, $F$ must satisfy $\dst{\int_{a}^{b}}F(t)\sin(t-a)\,dt=0$.
In this case, the solutions of the boundary value problem are
$y=y_{p}+c_{1}\sin(x-a)$, with $c_{1}$ arbitrary.
\end{alist}

\item
Let $y_{1}=\sinh(x-a)$  and $y_{2}=\sinh(x-b)$. Then
$y_{1}(a)=0$, $y_{2}(b)=0$, and
\[
W(x)=\begin{vmatrix}
\sinh(x-a)&\sinh(x-b)\\\cosh(x-a)&\cosh(x-b)
\end{vmatrix}=\sinh(b-a)\ne0.
\]
(Since $W$ is constant (Theorem~5.1.4), evaluate it by
setting $x=b$.) From Theorem 13.1.2,  (A) has a unique solution
for any continuous $F$  and constants $k_{1}$ and $k_{2}$.
If $y=u_{1}\sinh(x-a)+u_{2}\sinh(x-b)$ where
\begin{align*}
u_{1}'\sinh(x-a)+u_{2}'\sinh(x-b)&=0 \\
u_{1}'\cosh(x-a)+u_{2}'\cosh(x-b)&=F,
\end{align*}
then $y''-y=F$.
\begin{gather*}
u_{1}'=-F(x)\frac{\sinh(x-b)}{\sinh(b-a)}, \quad
u_{2}'=F(x)\frac{\sinh(x-a)}{\sinh(b-a)},
\\
u_{1}=\frac{1}{\sinh(b-a)}\int_{x}^{b}F(t)\sinh(t-b)\,dt,\quad
u_{2}=\frac{1}{\sinh(b-a)}\int_{a}^{x}F(t)\sinh(t-a)\,dt;
\\
y=\frac{\sinh(x-a)}{\sinh(b-a)}\int_{x}^{b}F(t)\sinh(t-b)\,dt+
\frac{\sinh(x-b)}{\sinh(b-a)}\int_{a}^{x}F(t)\sinh(t-a)\,dt.
\end{gather*}


\item
Let $y_{1}=\cosh(x-a)$  and $y_{2}=\cosh(x-b)$. Then
$y_{1}'(a)=y_{2}'(b)=0$  and
\[
W(x)=
\begin{vmatrix}
\cosh(x-a)&\cosh(x-b)\\\sinh(x-a)&\sinh(x-b)
\end{vmatrix}=-\sinh(b-a)\ne0.
\]
(Since $W$ is constant  (Theorem~5.1.40, evaluate it by
setting $x=b$.)
If \\$y=u_{1}\cosh(x-a)+u_{2}\cosh(x-b)$ where
\begin{align*}
u_{1}'\cosh(x-a)+u_{2}'\cosh(x-b)&=0 \\
u_{1}'\sinh(x-a)+u_{2}'\sinh(x-b)&=F,
\end{align*}
then $y''-y=F$.
\begin{gather*}
u_{1}'=F(x)\frac{\cosh(x-b)}{\sinh(b-a)}, \quad
u_{2}'=-F(x)\frac{\cosh(x-a)}{\sinh(b-a)},
\\
u_{1}=-\frac{1}{\sinh(b-a)}\int_{x}^{b}F(t)\cosh(t-b)\,dt,\quad
u_{2}=-\frac{1}{\sinh(b-a)}\int_{a}^{x}F(t)\cosh(t-a)\,dt;
\\
y=-\frac{\cosh(x-a)}{\sinh(b-a)}\int_{x}^{b}F(t)\cosh(t-b)\,dt
-\frac{\cosh(x-b)}{\sinh(b-a)}\int_{a}^{x}F(t)\cosh(t-a)\,dt.
\end{gather*}


\item Let $y_{1}=\sin\omega x$, $y_{2}\sin\omega(x-\pi)$; then
$y_{1}(0)=0$, $y_{2}(\pi)=0$,
\[
W(x)=
\begin{vmatrix}
\sin\omega x&\sin\omega(x-\pi)\\
\omega\cos\omega x &\omega\cos\omega(x-\pi)
\end{vmatrix}=\omega\sin\omega\pi\ne0
\]
if  and only if $\omega$ is not an integer. If this is so, then
$y=u_{1}\sin\omega x+u_{2}\sin\omega(x-\pi)$ if
\begin{align*}
u_{1}'\sin\omega x+u_{2}'\sin\omega(x-\pi)&=0\\
\omega(u_{1}'\cos\omega x+u_{2}'\cos\omega(x-\pi))&=F;
\end{align*}
\begin{align*}
u_{1}'&=-F\frac{\sin\omega(x-\pi)}{\omega\sin\omega\pi}, \quad
u_{2}'=F\frac{\sin\omega x}{\omega\sin\omega\pi};
\\
u_{1}&=\frac{1}{\omega\sin \omega \pi}\int_{x}^{\pi}
F(t)\sin\omega(t-\pi)\,dt;\quad
u_{2}=\frac{1}{\omega\sin \omega \pi}\int_{0}^{x}
F(t)\sin\omega t\,dt;
\\
y&=\frac{1}{\omega\sin \omega \pi}\left(
\sin\omega x\int_{x}^{\pi}F(t)\sin\omega(t-\pi)\,dt +
\sin\omega(x-\pi) \int_{0}^{\pi}
F(t)\sin\omega t\,dt\right).
\end{align*}

If $\omega= n$ (positive integer), then $y_{1}=\sin nx$  is a nontrivial
solution of $y''+y=0$, $y(0)=0$, $y(\pi)=0$. Let $y_{2}=\cos nx$; then
$W(x)=
\begin{vmatrix}
\sin  nx &\cos nx \\ n\cos nx&-n\sin nx
\end{vmatrix}=-n$, and $y_{p}=u_{1}\sin nx+u_{2}\cos nx$
satisfies $y_{p}''+n^{2}y_{p}=0$  if
\begin{align*}
u_{1}'\sin nx +u_{2}'\cos nx&=0\\
nu_{1}'\cos nx -nu_{2}'\sin nx&=F;\\
\end{align*}
$u_{1}'=\dst{\frac{1}{n}}F\cos nx$,
$u_{2}'=-\dst{\frac{1}{n}}F\sin nx$;
$u_{1}=-\dst{\frac{1}{n}}\dst{\int_{x}^{\pi}}F(t)\cos nt\,dt$;
$u_{2}=-\dst{\frac{1}{n}}\dst{\int_{0}^{x}}F(t)\sin nt\,dt$;
\[
y_{p}=-\frac{1}{n}\left(\sin nx\int_{x}^{\pi}F(t)\cos nt\,dt
+\cos nx \int_{0}^{x}F(t)\sin nt\,dt\right);
\]
$y=y_{p}+c_{1}\sin nx+c_{2}\cos nx$. Since $y_{p}=0$, $y(0)=0$, so
$c_{2}=0$; $y=y_{p}+c_{1}\sin nx$. Since $y(\pi)=0$,
$\dst{\int_{0}^{\pi}}F(t)\sin nt\,dt=0$ is necessary for existence of a
solution. If this hold, then the solutions are $y=y_{p}+c_{1}\sin nx$,
with $c_{1}$ arbitrary.


\item
Let $y_{1}=\cos\omega x$; $y_{2}=\sin\omega(x-\pi)$; then
$y_{1}'(0)=y_{2}(\pi)=0$, and
\[
W(x)=
\begin{vmatrix}
\cos\omega x& \sin\omega(x-\pi)\\
-\omega\sin\omega x&\omega\cos\omega(x-\pi)
\end{vmatrix}=\omega\cos\omega\pi\ne0
\]
if and only if $\omega\ne n+1/2$ ($n=$ integer).
If this is so, then $y=u_{1}\cos \omega x+u_{2}\sin\omega(x-\pi)$
satisfies $y''+\omega^{2}y=F(x)$ if
\begin{align*}
u_{1}'\cos\omega x+u_{2}'\sin\omega(x-\pi)&=0\\
\omega(-u_{1}'\sin\omega x+u_{2}'\cos\omega(x-\pi))\omega&=F;
\end{align*}
then
\begin{gather*}
u_{1}'=-\frac{F\sin\omega(x-\pi)}{\omega\cos\omega\pi},\quad
u_{2}'=\frac{F\cos\omega x}{\omega\cos\omega\pi}, \quad
\\
u_{1}=\frac{1}{\omega\cos\omega\pi}
\int_{x}^{\pi}F(t) \sin\omega(t-\pi),dt, \quad
u_{2}=\frac{1}{\omega\cos\omega\pi}
\int_{0}^{x} F(t)\cos\omega t\,dt,\quad
\\
y=\frac{1}{\omega\cos\omega\pi}
\left(\sin\omega x \int_{x}^{\pi}F(t) \sin\omega(t-\pi),dt
+\sin\omega(x-\pi)\int_{0}^{x}F(t)\cos\omega t\,dt\right).
\end{gather*}

If  $\omega=n+1/2$ ($n=$ integer), then $y_{1}=\cos(n+1/2)x$ is a
nontrivial solution of $y''+y=0$, $y'(0)=y(\pi)=0$. Let
$y_{2}=\sin(n+1/2)x$; then
\[
W(x)=
\begin{vmatrix}
\cos(n+1/2)x&\sin(n+1/2)x\\
-(n+1/2)\sin(n+1/2)x&(n+1/2)\cos(n+1/2)x\\
\end{vmatrix}=n+1/2,
\]
so $y_{p}=u_{1}\cos(n+1/2)x+u_{2}\sin(n+1/2)x$  satisfies
$y_{p}''+(n+1/2)^{2}y_{p}=F$ if
\begin{align*}
u_{1}'\cos(n+1/2)x+u_{2}'\sin(n+1/2)x&=0 \\
-(n+1/2)u_{1}'\sin(n+1/2)x+(n+1/2)u_{2}'\cos(n+1/2)x&=F.
\end{align*}
\begin{align*}
u_{1}'&=-\frac{F\sin(n+1/2)x}{n+1/2};\quad
u_{2}'=\frac{F\cos(n+1/2)x}{n+1/2};
\\
u_{1}&=\int_{x}^{\pi} \frac{F(t)\sin(n+1/2)t}{n+1/2}\,dt; \quad
u_{2}=\int_{0}^{x} \frac{F(t)\cos(n+1/2)t}{n+1/2}\,dt; \quad
\\
y_{p}&=\frac{1}{n+\frac12}\left(
\cos(n+\frac12)x\int_{x}^{\pi} F(t)\sin(n+\frac12)t\,dt
+\sin(n+\frac12)x
\int_{0}^{x} F(t)\cos(n+\frac12)t\,dt\right);
\\
y_{p}'&=-
\sin(n+1/2)x\int_{x}^{\pi} F(t)\sin(n+1/2)t\,dt
+\cos(n+1/2)x
\int_{0}^{x} F(t)\cos(n+1/2)t\,dt;
\\
y&=y_{p}+c_{1}\cos(n+1/2)x+c_{2}\sin(n+1/2)x;
\\
y'&=y_{p}'+(n+1/2)(-c_{1}\sin(n+1/2)x+c_{2}\cos(n+1/2))x.
\end{align*}
Since $y_{p}'(0)=0$, $y'(0)=0\implies c_{2}=0$;
$y=y_{p}+c_{1}\cos(n+1/2)x$;  \\
$y'=y_{p}'+(n+1/2)c_{1}(n+1/2)x$; $y(\pi)=0\implies y_{p}(\pi)=0$.
Hence, $\dst{\int_{0}^{\pi}}F(y)\cos(n+1/2)t\,dt =0$
is a necessary condition for existence of a solution. If this holds,
then the solutions are $y=y_{p}+c_{1}\cos(n+1/2)x$  with $c_{1}$
arbitrary.



\item
Suppose $y=c_{1}z_{1}+c_{2}z_{2}$ is a nontrivial solution of
the homogeneous boundary value problem. Then
$B_{1}(y)=c_{1}B_{1}(z_{1})+c_{2}B_{1}(z_{2})=0$.
From Theorem~13.1.1 we may assume without loss of
generality that
$B_{1}(z_{2})\ne0$. Then
$c_{2}=-\dst{\frac{B_{1}(z_{1})}{B_{1}(z_{2})}}c_{1}$.
Therefore,   $y$ is constant multiple of
$y_{0}=B_{1}(z_{2})z_{1}-B_{1}(z_{1})z_{2}\ne0$. To that check $y$
satisfies the boundary conditions, note  that
$B_{1}(y_{0})=B_{1}(z_{2})B_{1}(z_{1})-B_{1}(z_{1})B_{1}(z_{2})=0$
$B_{2}(y_{0})=B_{1}(z_{2})B_{2}(z_{1})-B_{1}(z_{1})B_{2}(z_{2})=0$,
by Theorem~13.1.2.




\item
$y_{1}=a_{1}+a_{2}x$; $y_{1}(0)-2y_{1}'(0)=a_{1}-2a_{2}=0$  if $a_{1}=2$,
$a_{2}=1$;
$y_{1}=2+x$. $y_{2}=b_{1}+b_{2}x$; $y_{2}(1)=2y_{1}'(1)=b_{1}+3b_{2}=0$
if $b_{1}=3$, $b_{2}=-1$;  $y_{2}=3-x$.
\begin{gather*}
W(x)=
\begin{vmatrix}
2+x&3-x\\1&-1
\end{vmatrix}=-5;\quad
G(x,t)=
\begin{cases}
-\dst{\frac{(2+t)(3-x)}{5}}, &\quad 0\le t\le x,\\[.1in]
-\dst{\frac{(2+x)(3-t)}{5}}, &\quad x\le t\le 1.
\end{cases}\\
y=-\frac{1}{5}\left[(2+x)\int_{x}^{1}(3-t)F(t)\,dt+(3-x)\int_{0}^{x}(2+t)F(t)\,dt\right].\tag{B}
\end{gather*}

\begin{alist}
\item
With $F(x)=1$, (B)  becomes
\begin{align*}
y&=-\frac{1}{5}\left[(2+x)\int_{x}^{1}(3-t)\,dt+(3-x)\int_{0}^{x}(2+t)\,dt\right]\\
&=-\frac{1}{5}\left[(2+x)\left(\frac{x^{2}-6x+5}{2}\right)
+(3-x)\left(\frac{x^{2}+4x}{2}\right)\right] \\
&=\frac{x^{2}-x-2}{2}.
\end{align*}

\item
With $F(x)=x$, (B)  becomes
\begin{align*}
y&=-\frac{1}{5}\left[(2+x)\int_{x}^{1}(3t-t^{2})\,dt+(3-x)\int_{0}^{x}(2t+t^{2})\,dt\right]\\
&=-\frac{1}{5}\left[(2+x)\left(\frac{2x^{3}-9x^{2}+7}{6}\right)
+(3-x)\left(\frac{x^{3}+3x^{2}}{3}\right)\right] \\
&=\frac{5x^{3}-7x-14}{30}.
\end{align*}

\item
With $F(x)=x^{2}$, (B)  becomes
\begin{align*}
y&=-\frac{1}{5}\left[(2+x)\int_{x}^{1}(3t^{2}-t^{3})\,dt+(3-x)\int_{0}^{x}(2t^{2}+t^{3})\,dt\right]\\
&=-\frac{1}{5}\left[(2+x)\left(\frac{x^{4}-4x^{2}+3}{4}\right)
+(3-x)\left(\frac{3x^{4}+8x^{3}}{12}\right)\right] \\
&=\frac{5x^{4}-9x-18}{60}.
\end{align*}
\end{alist}

\item $y_{1}=x^{2}-x$, $y_{2}=x^{2}-2x$; then $y_{1}(1)=0$,
$y_{2}(2)=0$;
$W(x)=
\begin{vmatrix}
x^{2}-x &x^{2}-2x\\2x-1&2x-2
\end{vmatrix}=x^{2}$.
Since $P_{0}(x)=x^{2}$,
$G(x,t)=
\begin{cases}
\dst{\frac{(t-1)x(x-2)}{t^{3}}}, &\quad 1\le t\le x,\\[.1in]
\dst{\frac{x(x-1)(t-2)}{t^{3}}}, &\quad x\le t\le 2.
\end{cases}.$
\[
y=x(x-1)\int_{x}^{2}\frac{t-2}{t^{3}}F(t)\,dt+x(x-2)\int_{1}^{x}F(t)\,dt.
\tag{B}
\]

\begin{alist}
\item With $F(x)=2x^{3}$, (B) becomes
\begin{align*}
y&=2x(x-1)\int_{x}^{2}(t-2)\,dt +2x(x-2)\int_{1}^{x}(t-1)\,dt\\
&=-x(x-1)(x-2)^{2}+x(x-2)(x-1)^{2}=x(x-1)(x-2).
\end{align*}

\item With $F(x)=6x^{4}$, (B) becomes
\begin{align*}
y&=6x(x-1)\int_{x}^{2}(t-2)t\,dt +6x(x-2)\int_{1}^{x}(t-1)t\,dt\\
&=-2x(x-1)(x+1)(x-2)^{2}+x(x-2)(x-1)^{2}(2x+1)=x(x-1)(x-2)(x+3).
\end{align*}
\end{alist}

\item $y_{1}=a_{1}+a_{2}x$; $y_{1}'=a_{2}$;
$B_{1}(y_{1})=\alpha a_{1}+\beta a_{2}=0$ if $a_{1}=\beta$,
$a_{2}=-\alpha$; $y_{1}=\beta-\alpha  x$.
$y_{2}=b_{1}+b_{2}x$;
$y_{2}'=b_{2}$;
$B_{2}(y_{2})=\rho b_{1}+(\rho+\delta)b_{2}=0$ if $b_{1}=\rho+\delta$,
$b_{2}=-\rho$; $y_{2}=\rho+\delta-\rho x$;
$W(x)=
\begin{vmatrix}
\beta-\alpha x&\rho+\delta-\rho x\\-\alpha&-\rho
\end{vmatrix}
=\alpha(\rho+\delta)-\beta\rho$. From Theorem~13.1.2, (A) has a unique
solution if and only if $\alpha(\rho+\delta)-\beta\rho\ne0$. Then
\[
G(x,t)=
\begin{cases}
\dst{\frac{(\beta-\alpha t)(\rho+\delta-\rho x)}
{\alpha(\rho+\delta)-\beta\rho}}, \quad 0\le t\le x,\\[.1in]
\dst{\frac{(\beta-\alpha x)(\rho+\delta-\rho t)}
{\alpha(\rho+\delta)-\beta\rho}}, \quad x\le t\le 1.
\end{cases}
\]


\item
 $y_{1}=a_{1}\cos x+a_{2}\sin x$;
$y_{1}'=-a_{1}\sin  x+a_{2}\cos x$;
$B_{1}(y_{1})=\alpha a_{1}+\beta a_{2}=0$ if $a_{1}=\beta$,
$a_{2}=-\alpha$.
$y_{1}=\beta\cos x-\alpha\sin x$.
$y_{2}=b_{1}\cos x+b_{2}\sin x$;
$y_{2}'=-b_{1}\sin x+b_{2}\cos x$;
$B_{2}(y_{2})=\rho b_{2}-\delta b_{1}=0$ if $b_{1}=\rho$,
$b_{2}=\delta$;
 $y_{2}=\rho\cos x+\delta\sin x$;
$W(x)=
\begin{vmatrix}
\beta\cos x-\alpha\sin x&\rho\cos x+\delta\sin x\\
-\beta\sin x-\alpha\cos x&-\rho\sin x+\delta\cos x\\
\end{vmatrix}$
Since $W$ is constant, we can evaluate it with $x=0$:
$W=
\begin{vmatrix}
\beta&\rho\\-\alpha&\delta
\end{vmatrix}=\alpha\rho+\beta\delta$.
From Theorem~13.1.2, (A) has a unique
solution if and only if $\alpha\rho+\beta\delta\ne0$. Then
\[
G(x,t)=
\begin{cases}
\dst{\frac{(\beta\cos t-\alpha\sin  t)(\rho\cos x+\delta\sin  x)}
{\alpha\rho+\beta\delta}}, \quad 0\le t\le x,\\[.1in]
\dst{\frac{(\beta\cos x-\alpha\sin  x)(\rho\cos t+\delta\sin t))}
{\alpha\rho+\beta\delta}}, \quad x\le t\le \pi.
\end{cases}
\]


\item
 $y_{1}=e^{x}(a_{1}\cos x+a_{2}\sin x)$;
$y_{1}'=e^{x}[a_{1}(\cos x-\sin x)+a_{2}(\sin x+\cos x)]$;

$B_{1}(y_{1})=(\alpha+\beta)a_{1}+\beta a_{2}=0$ if $a_{1}=\beta$,
$a_{2}=-(\alpha+\beta)$.

$y_{1}=e^{x}(\beta\cos x-(\alpha+\beta)\sin x$.
 $y_{2}=e^{x}(b_{1}\cos x+b_{2}\sin x)$;

$y_{2}'=e^{x}[(b_{1}(\cos x-\sin x))+b_{2}(\sin x+\cos x)]$;

$B_{2}(y_{2})=-e^{\pi/2}[(\rho+\delta)b_{2}-\delta b_{1}]=0$  if
$b_{1}=\delta$, $b_{2}=(\rho+\delta)$;

$y_{2}=e^{x}[(\rho+\delta)\cos x+\delta\sin x)$;

To evaluate $W(x)$, we write $y_{1}=e^{x}v_{1}$ and
$y_{2}=e^{x}v_{2}$, where

$v_{1}=\beta\cos x-(\alpha+\beta)\sin x$ and
$v_{2}=(\rho+\delta)\cos x+\delta\sin x$.

Then $y_{1}'=y_{1}+e^{x}v_{1}'$
and $y_{2}=y_{2}+e^{x}y_{2}'$,
\begin{align*}
W(x)&=
\begin{vmatrix}
y_{1}&y_{2}\\ y_{1}+e^{x}v_{1}'&y_{2}+e^{x}v_{2}'
\end{vmatrix}=
\begin{vmatrix}
y_{1}&y_{2}\\ e^{x}v_{1}'&{x}v_{2}'
\end{vmatrix}= e^{2x}
\begin{vmatrix}
v_{1}&v_{2}\\ v_{1}'&{x}v_{2}'
\end{vmatrix}  \\
&=
\begin{vmatrix}
(\beta\cos x-(\alpha+\beta)\sin x&
(\rho+\delta)\cos x+\delta\sin x \\
(-\beta\sin x-(\alpha+\beta)\cos x&
(-(\rho+\delta)\sin x+\delta\cos x
\end{vmatrix}.
\end{align*}
Since $v_{i}''+v_{i}=0$, $i=1$, $2$,
Theorem~5.1.4 implies that $W(x)=Ke^{2x}$, where
is a constant that can be determined by setting $x=0$ in the determinant:
\[
W(x)e^{2x}
\begin{vmatrix}
\beta&\rho+\delta\\-\alpha-\beta&\delta
\end{vmatrix}
=[\beta\delta+(\alpha+\beta)(\rho+\delta)].
\]
From Theorem~13.1.2, the boundary value problem has a unique
solution if and only if $\beta\delta+(\alpha+\beta)(\rho+\delta)\ne0$.
In this case the Green's function is
\[
G(x,t)=
\begin{cases}
e^{x-t}\dst{\frac{[\beta\cos t-(\alpha+\beta)\sin t]
[\rho+\delta)\cos x+\delta\sin x]}
{\beta\delta+(\alpha+\beta)(\rho+\delta)}},&a\le t\le x \\
e^{x-t}\dst{\frac{[\beta\cos x-(\alpha+\beta)\sin x]
[\rho+\delta)\cos t+\delta\sin t]}
{\beta\delta+(\alpha+\beta)(\rho+\delta)}},&x\le t \le \pi/2.
\end{cases}
\]


\item Let $y_{p}=\dst{\int_{a}^{b}}G(x,t)F(t)\,dt$. From
Theorem~13.1.3, $Ly_{p}=F$, $B_{1}(y_{p})=0$, and $B_{2}(y_{p})=0$.
The solution of $Ly=F$, $B_{1}(y)=k_{1}$, and $B_{2}(y)=k_{2}$
is of the form $y=y_{p}+c_{1}y_{1}+c_{2}y_{2}$. Since
$B_{1}(y_{p})=0$ and $B_{1}(y_{1})=0$,
$B_{1}(y)=k_{1}\implies k_{1}=c_{2}B_{1}(y_{2})\implies
c_{2}=\dst{\frac{k_{1}}{B_{1}(y_{2})}}$.
Since $B_{2}(y_{p})=0$ and $B_{2}(y_{2})=0$,
$B_{2}(y)=k_{2}\implies k_{2}=c_{1}B_{2}(y_{1})\implies
c_{1}=\dst{\frac{k_{2}}{B_{2}(y_{1})}}$.
\end{sectionSolutions}

\section{Sturm-Liouville Problems}

\begin{sectionSolutions}
\item
$y''+\dst{\frac{1}{x}}y'+\left(1-\dst{\frac{\nu^{2}}{x^{2}}}\right)y=0$;
$\dst{\frac{p'}{p}}=\dst{\frac{1}{x}}$; $\ln|p|=\ln|x|$;
$p=x$; $xy''+y'=\dst{\left(x-\frac{\nu^{2}}{x}\right)}y=0$;
$(xy')'+\left(x-\dst{\frac{\nu^{2}}{x}}\right)y=0$.



\item
$y''+\dst{\frac{b}{x}}y'+\dst{\frac{c}{x^{2}}}y=0$;
$\dst{\frac{p'}{p}}=\dst{\frac{b}{x}}$; $\ln|p|=b\ln |x|$; $p=x^{b}$;

$x^{b}y''+bx^{b-1}y'+cx^{b-2}y=0$; $(x^{b}y')'+cx^{b-2}y=0$.




\item
$xy''+(1-x)y'+\alpha y=0$;
$y''+\left(\dst{\frac{1}{x}}-1\right)y'+\dst{\frac{\alpha}{x}}y=0$;
$\dst{\frac{p'}{p}}=\dst{\frac{1}{x}}-1$; $\ln|p|=\ln|x|-x$; $p=xe^{-x}$;
$xe^{-x}y''+(1-x)y'+\alpha e^{-x}y=0$; $(xe^{-x}y')'+\alpha e^{-x}y=0$.




\item If $\lambda$ is an eigenvalues of (A) and $y$ is a $\lambda$-eigenfunction,
 multiplying the differential equation in (B) by $y$ yields
$(xy')y+\dst{\frac{\lambda}{x}}y^{2}=0$;
\begin{gather*}
\lambda\int_{1}^{2}\frac{y^{2}(x)}{x^{2}}\,dx=
-\int_{1}^{2}(xy'(x))'y(x)\,dx=
-xy'(x)y(x)\bigg|_{1}^{2}+\int_{1}^{2}x(y'(x))^{2}\,dx;
\\
y(1)=y_(2)=0\implies xy'(x)y(x)\bigg|_{1}^{2}=0;\quad
\lambda\int_{1}^{2}\frac{y^{2}(x)}{x}\,dx=\int_{1}^{2}x(y'(x))^{2}\,dx.
\end{gather*}
Therefore $\lambda\ge 0.$ We must still show that $\lambda=0$  is not an
eigenvalue. To this end, suppose that $(xy')'=0$; then
$xy'=c_{1}$;   $y'=\dst{\frac{c_{1}}{x}}$; $y=c_{1}\ln|x|+c_{2}$;
$y(1)=0\implies c_{2}=0$; $y=c_{1}\ln|x|$; $y(2)=0\implies c_{1}=0$;
$y\equiv0$; therefore $\lambda=0$ is not an eigenvalue.



\item
Characteristic equation: $r^{2}+2r+1+\lambda=0$; $r=-1\pm\sqrt{-\lambda}$.

$\lambda=0$: $y=e^{-x}(c_{1}+c_{2}x)$; $y'=-e^{-x}(c_{1}-c_{2}+c_{2}x)$;
$y'(0)=0\implies c_{1}=c_{2}$; $y'=-c_{2}xe^{-x}$; $y'(1)=0\implies
-c_{2}/e=0\implies c_{2}=0$; $\lambda=0$ is not an eigenvalue.

 $\lambda=-k^{2}$, $k>0$: $r=-1\pm k$;
$y=e^{-x}(c_{1}\cosh kx+c_{2}\sinh kx)$;
$y'=-c_{1}e^{-x}(-\cosh kx+k\sinh kx)+c_{2}e^{-x}(-\sinh
kx+k\cosh kx)$.
 The boundary conditions require that\\
$c_{1}+c_{2}k=0$ \; and  \;
$(-\cosh k+k\sinh k)c_{1}+(-\sinh k+k\cosh k)c_{2}=0$.
This system has a nontrivial solution if and only if  $(1-k^{2})\sinh
k=0$. Let
$k=1$ and $c_{1}=c_{2}=1$; then $\lambda=-1$ is the only negative
eigenvalue, with associated eigenfunction $y=1$.





 $\lambda=k^{2}$, $k>0$:  $r=-1\pm ik$;
 $y=e^{-x}(c_{1}\cos kx+c_{2}\sin kx)$;  \\
$y'=c_{1}e^{-x}(-\cos kx-k\sin kx)+c_{2}e^{-x}(-\sin x+k\cos kx)$.
The boundary conditions require that  \\
 \centerline{$-c_{1}+c_{2}k=0$ \; and \;
$(-\cos k-k\sin k)c_{1}+(-\sin k+k\cos k)c_{2}=0$.}
\noindent
This system has a nontrivial solution if and only if  $(1+k^{2})\sin k=0$.
Let
$k=n\pi$ ($k$ a positive integer)and $c_{1}=n\pi$, $c_{2}=1$; then
$\lambda_{n}=n^{2}\pi^{2}$ is an eigenvalue,
 with associated eigenfunction $y_{n}=e^{-x}(n\pi\cos n\pi x+\sin n\pi
x)$.



\item
Characteristic equation: $r^{2}+\lambda=0$.

$\lambda=0:$
$y=c_{1}+c_{2}x$. $y(0)=0\implies c_{1}=0$, so $y=c_{2}x$.
Now $y(1)-2y'(1)=0\implies c_{2}=0$. Therefore $\lambda=0$
is not an eigenvalue.


$\lambda=-k^{2}$, $k>0$:
$y=c_{1}\cosh kx+c_{2}\sinh kx$; $y'=k(c_{1}\sinh kx+c_{2}\cosh
kx)$. $y'(0)\implies c_{2}=0$, so $y=c_{1}\cosh kx$. Now
$y(1)-2y'(1)=0\implies  c_{1}(\cosh k-2k \sinh k)=0$, which is
possible
with $c_{1}\ne0$ if and only if   $\tanh k=\dst{\frac{1}{2k}}$. Graphing
both sides of this equation on the same axes show that it has one positive
solution
$k_{0}$;
$y_{0}=\cosh k_{0}x$ is a $-k_{0}^{2}$-eigenfunction.



$\lambda=k^{2}$, $k>0$:
$y=c_{1}\cos kx+c_{2}\sin kx$; $y'=k(-c_{1}\sin kx+c_{2}\cos
kx)$. $y'(0)\implies c_{2}=0$, so $y=c_{1}\cos kx$. Now
$y(1)-2y'(1)=0\implies   c_{1}(\cos k+2k \sin k)=0$, which is
possible
with $c_{1}\ne0$ if and only if   $\tan k=-\dst{\frac{1}{2k}}$. Graphing
both sides of this equation on the same axes shows that it has
a solution $k_{n}$ in $((2n-1)\pi/2,n\pi)$, $n=1$, $2$, $3$, \dots;
$y_{n}=\cos k_{n}x$ is a
$k_{n}^{2}$-eigenfunction.





\item
Characteristic equation: $r^{2}+\lambda=0$.

$\lambda=0:$
 $y=c_{1}+c_{2}x$. The boundary conditions require  that
$c_{1}+2c_{2}=0$ and $c_{1}+\pi c_{2}=0$, which imply that
$c_{1}=c_{2}=0$, so $\lambda=0$ is not an eigenvalue.

$\lambda=-k^{2}$, $k>0$:
$y=c_{1}\cosh kx+c_{2}\sinh kx$; $y'=k(c_{1}\sinh kx+c_{2}\cosh
kx)$.
 The boundary conditions require that  \\
\centerline{$c_{1}+2kc_{2}=0$  \;  and \;  $c_{1}\cosh
k\pi+c_{2}\sinh 2k\pi=0$.}
\noindent
This system
has a nontrivial solution if and only if $\tanh k\pi=2k$.
 Graphing
both sides of this equation on the same axes shows that it has
a solution $k_{0}$ in $(0,\pi)$;
$y_{0}=2k_{0}\cosh k_{0}x-\sinh k_{0}x$ is a
$-k_{0}^{2}$-eigenfunction.



$\lambda=k^{2}$, $k>0$:
 $y=c_{1}\cos kx+c_{2}\sin kx$; $y'=k(-c_{1}\sin kx+c_{2}\cos
kx)$.
 The boundary conditions require that    \\
\centerline{$c_{1}+2kc_{2}=0$\; and \;$c_{1}\cos k\pi+c_{2}\sin k\pi=0$.}
\noindent
This system has
 a nontrivial solution if and only if $\tan k\pi=2k$.
Graphing both sides of this equation on the same axes shows that it
has a solution $k_{n}$ in $(n, n+1/2)$, $n=1$, $2$,$3$, \dots;
$y_{n}=2k_{n} \cos k_{n}x-\sin k_{n}x$ is a $k_{n}^{2}$-eigenfunction.





\item
Characteristic equation: $r^{2}+\lambda=0$.

$\lambda=0:$
$y=c_{1}+c_{2}x$. The  boundary conditions require that $c_{1}+c_{2}=0$
and
$c_{1}+4c_{2}=0$, so $c_{1}=c_{2}=0$.
Therefore  $\lambda=0$
is not an eigenvalue.


$\lambda=-k^{2}$, $k>0$:
 $y=c_{1}\cosh kx+c_{2}\sinh kx$; $y'=k(c_{1}\sinh kx+c_{2}\cosh
kx)$.  The boundary conditions require that\\
\centerline{$c_{1}+kc_{2}=0$ \; and  \;
$(\cosh 2k+2k\sinh 2k)c_{1}+(\sinh 2k+2k\cosh 2k)c_{2}=0$.}
\noindent
This system has
a  nontrivial solution if and only if
 $\tanh 2k =-\dst{\frac{k}{1-2k^{2}}}$.
 Graphing
both sides of this equation on the same axes shows that it has a
solution $k_{0}$ in $(1/\sqrt2)$; $y_{0}=k_{0}\cosh k_{0}x-\sinh k_{0}x$ is
a  $-k_{0}^{2}$- eigenfunction.



$\lambda=k^{2}$, $k>0$:
 $y=c_{1}\cos kx+c_{2}\sin kx$; $y'=k(-c_{1}\sin kx+c_{2}\cos
kx)$.  The boundary conditions require that  \\
\centerline{$c_{1}+kc_{2}=0$ \; and \;
$(\cos 2k-2k\sin 2k)c_{1}+(\sin 2k+2k\cos 2k)c_{2}=0$.}
\noindent
This system has  a
nontrivial solution if and only if
 $\tan 2k =-\dst{\frac{k}{1+2k^{2}}}$.
 Graphing
both sides of this equation on the same axes shows that it has
a solution $k_{n}$ in $((2n-1)\pi/4,n\pi/2)$, $n=1$, $2$,$3$, \dots;
$y_{n}=k_{n} \cos k_{n}x-\sin k_{n}x$ is a $k_{n}^{2}$-eigenfunction.





\item
Characteristic equation: $r^{2}+\lambda=0$.

$\lambda=0:$
$y=c_{1}+c_{2}x$. The  boundary conditions require that $3c_{1}+2c_{2}=0$
and
$3c_{1}+4c_{2}=0$, so $c_{1}=c_{2}=0$.
Therefore  $\lambda=0$
is not an eigenvalue.


$\lambda=-k^{2}$, $k>0$:
 $y=c_{1}\cosh kx+c_{2}\sinh kx$;  $y'=k(c_{1}\sinh kx+c_{2}\cosh
kx)$.  The boundary conditions require that        \\
\centerline{$3c_{1}+kc_{2}=0$ \; and  \;
$(3\cosh 2k-2k\sinh 2k)c_{1}+(3\sinh 3k-2k\cosh 2k)c_{2}=0$.}
\noindent
This system has a nontrivial solution if and only if
 $\tanh 2k =\dst{\frac{9k}{9+2k^{2}}}$.
 Graphing
both sides of this equation on the same axes shows that it has solutions
$y_{1}$ in $(1,2)$ and $y_{2}$ in $(5/2,7/2)$;
$y_{n}=k_{n}\cosh k_{n}x-3\sinh k_{n}x$ is   a
  $-k_{n}^{2}$-eigenfunction, $k=1$, $2$.



$\lambda=k^{2}$, $k>0$:
 $y=c_{1}\cos kx+c_{2}\sin kx$; $y'=k(-c_{1}\sin kx+c_{2}\cos
kx)$.  The boundary conditions require that \\
\centerline{$3c_{1}+kc_{2}=0$ \; and \;
$(3\cos 2k+2k\sin 2k)c_{1}+(3\sin 2k-2k\cos 2k)c_{2}=0$.}
\noindent
This system has  a   nontrivial solution if and only if
 $\tan 2k =-\dst{\frac{9k}{9-2k^{2}}}$.
 Graphing
both sides of this equation on the same axes shows that it has
 solutions $k_{0}$ in $(3/\sqrt{2},\pi)$ and
$k_{n}$ in $((2n+3)\pi/4,(n+2)\pi/3)$,
$n=1$, $2$,$3$, \dots;
$y_{n}=k_{n} \cos k_{n}x-3\sin k_{n}x$ is a $k_{n}^{2}$-eigenfunction.





\item
Characteristic equation: $r^{2}+\lambda=0$.

$\lambda=0:$
 $y=c_{1}+c_{2}x$. The  boundary conditions require that
$5c_{1}+2c_{2}=0$   and
$5c_1+3c_{2}$, so $c_{1}=c_{2}=0$.
Therefore $\lambda=0$ is not an eigenvalue.



$\lambda=-k^{2}$, $k>0$:
 $y=c_{1}\cosh kx+c_{2}\sinh kx$; $y'=k(c_{1}\sinh kx+c_{2}\cosh
kx)$.  The boundary conditions require that       \\
\centerline{$5c_{1}+2kc_{2}=0$ \; and \;
$(5\cosh k-2k\sinh k)c_{1}+(5\sinh k-2k\cosh k)c_{2}=0$.}
\noindent
This system has a nontrivial solution if and only if
 $\tanh k =\dst{\frac{20k}{25+4k^{2}}}$.
 Graphing
both sides of this equation on the same axes shows that it has
solutions
$k_{1}$ in $(1,2)$ and $k_{2}$ in $(5/2,7/2)$;
$y_{n}=2k_{n}\cosh k_{n}x-\sinh k_{n}x$ is
  $-k_{n}$-eigenfunction, $n=1$, $2$.



$\lambda=k^{2}$, $k>0$:
 $y=c_{1}\cos kx+c_{2}\sin kx$; $y'=k(-c_{1}\sin kx+c_{2}\cos
kx)$.  The boundary conditions require that \\
\centerline{$5c_{1}+2kc_{2}=0$ \; and  \;
$(5\cos k+2k\sin k)c_{1}+(5\sin k-2k\cos k)c_{2}=0$.}
\noindent
This system has  a   nontrivial solution if and only if
 $\tan k =-\dst{\frac{20k}{25-4k^{2}}}$.
 Graphing
both sides of this equation on the same axes shows that it has
a solution $k_{n}$ in\\ $((2n+1)\pi/2,(n+1)\pi)$, $n=1$, $2$,$3$,
\dots;
$y_{n}=2k_{n} \cos k_{n}x-3\sin k_{n}x$ is a $k_{n}^{2}$-eigenfunction.



\item
$\lambda=0$: $x^{2}y''-2xy'+2y=0$ is an Euler equation
with indicial equation
$r(r-1)-2r+2=(r-1)(r-2)=0$. $y=x(c_{1}+c_{2}x)$; $y(1)=y(2)=0\implies
c_{1}+c_{2}=c_{1}+2c_{2}=0\implies c_{1}=c_{2}=0$, so $\lambda=0$
is not an eigenvalue.

$\lambda=-k^{2}$; $k>0$:
$y=x(c_{1}\cosh k(x-1)+c_{2}\sinh k(x-1))$; $y(1)=0\implies c_{1}=0$;
$y=c_{2}x\sinh k(x-1)$; $y(2)=0\implies 2c_{2}\sinh k=0\implies c_{2}=0$;
$\lambda$ is not an eigenvalue.


$\lambda=k^{2}$; $k>0$:
$y=x(c_{1}\cos k(x-1)+c_{2}\sin k(x-1))$; $y(1)=0\implies c_{1}=0$;
$y=c_{2}x\sin k(x-1)$; $y(2)=0$ with $c_{2}\ne0$ if $k=n\pi$
($n$ a positive integer); $\lambda_{n}=n^{2}\pi^{2}$;
 $y_{n}=x\sin n\pi(x-1)$ is a $k_{n}^{2}$-eigenfunction.



\item
$\lambda=0$: $x^{2}y''-2xy'+2y=0$ is an Euler equation
with indicial equation
$r(r-1)-2r+2=(r-1)(r-2)=0$. $y=x(c_{1}+c_{2}x)$; $y'=c_{1}+2c_{2}x$;
$y(1)=y'(2)=0\implies c_{1}+c_{2}=c_{1}+4c_{2}=0\implies c_{1}=c_{2}=0$,
 so $\lambda=0$  is  not an eigenvalue.


$\lambda=-k^{2}$; $k>0$:
$y=x(c_{1}\cosh k(x-1)+c_{2}\sinh k(x-1))$; $y(1)=0\implies c_{1}=0$;
$y=c_{2}x\sinh k(x-1)$; $y'=c_{2}
(\sinh k(x-1)+kx\cosh k(x-1))$;
 $y'(2)=0\implies c_{2}(\sinh k+k\cosh k)\implies c_{2}=0$;
$\lambda$ is not an eigenvalue.


$\lambda=k^{2}$; $k>0$:
$y=x(c_{1}\cos k(x-1)+c_{2}\sin k(x-1))$; $y(1)=0\implies c_{1}=0$;
$y=c_{2}x\sin k(x-1)$;
$y'=c_{2}(\sin k(x-1)+kx\cos k(x-1))$;
 $y'(2)=0$ with $c_{2}\ne0$ if and only if $\sin k+2k\cos k=0$ or,
equivalently,
$\tan k=-2k$.
 Graphing
both sides of this equation on the same axes shows that it has
a  solution $k_{n}$ in $((2n-1)\pi/2,n\pi)$, $n=1$, $2$,$3$,
\dots;
 $y_{n}=x\sin k_{n}(x-1)$ is a $k_{n}^{2}$-eigenfunction.



\item
$\lambda=0$:
 $y=c_{1}+c_{2}x$. The boundary   conditions require that
$c_{1}+\alpha c_{2}=0$ and $c_{1}+(\pi+\alpha)c_{2}=0$, so $c_{1}=c_{2}=0$.
Therefore $\lambda=0$ is not an eigenvalue of (A).


$\lambda=-k^{2}$, $k>0$:
 $y=c_{1}\cosh kx+c_{2}\sinh kx$;
$y'=k(c_{1}\sinh kx+c_{2}\cosh kx)$. The boundary conditions  require
that
\[
\tag{D}
\begin{gathered}
c_{1}+\alpha kc_{2}=0\\
(\cosh k\pi+\alpha k\sinh k\pi)c_{1}+(\sinh k\pi+\alpha k\cosh
k\pi)c_{2}=0.
\end{gathered}
\]
This
 system has  a nontrivial solution if and only if
$(1-k^{2}\alpha^{2})\sinh k\pi=0$, which holds with $k>0$ if and only if
$k^{2}=\pm 1/\alpha$. Therefore $\lambda=-1/\alpha^{2}$ is the only
negative eigenvalue. We can choose $k=\pm1/\alpha$. Either way, the first
equation in (D) implies that $e^{-x/\alpha}$ is an associated
eigenfunction.


$\lambda=k^{2}$, $k>0$:
 $y=c_{1}\cos kx+c_{2}\sin kx$;
  $y'=k(-c_{1}\sin kx+c_{2}\cos kx)$. The boundary conditions  require
that
\[
\tag{E}
\begin{gathered}
c_{1}+\alpha kc_{2}=0\\
(\cos k\pi-\alpha k\sin k\pi)c_{1}+(\sin k\pi+\alpha k\cos k
\pi)c_{2}=0.
\end{gathered}
\]
This
 system has  a nontrivial solution if and only if
$(1+k^{2}\alpha^{2})\sin k\pi=0$. Choosing $k=n$  produces eigenvalues
$\lambda_{n}=n^{2}\pi^{2}$. Setting  $k=n$ in the first equation in
(E) yields $c_{1}+\alpha nc_{2}=0$, so
$y_{n}=n\alpha\cos nx-\sin nx$.



\item
 $y=c_{1}+c_{2}x$. The boundary conditions require that\\
$\alpha c_{1}+\beta c_{2}=0$ \;
and \; $\rho c_{1}+(\rho L+\delta)c_{2}=0$.
 This system  has a nontrivial solution if and only if $\alpha(\rho L+\delta)-\beta\rho=0$.




\item
\begin{alist}
\item
 $y=c_{1}\cos kx+c_{2}\sin kx$;
 $y'=k(-c_{1}\sin kx+c_{2}\cos kx)$. The boundary conditions  require
that  \\
\centerline{$\alpha c_{1}+\beta kc_{2}=0$ \; and \;
$(\rho\cos kL-\delta k\sin kL)c_{1}+(\rho\sin kL+\delta k\cos kL)c_{2}=0$.}
\noindent
This system has a nontrivial solution if and only if its determinant
is zero. This implies the conclusion.

\item
If $\alpha\delta-\beta\rho=0$, (A) reduces to
\[\tag{B}
(\alpha\rho+k^{2}\beta\delta)\sin kL=0.
\]
From the solution of Exercise~13.2.29\part{b},
$\alpha\rho+k^{2}\beta\delta>0$
for all $k>0$. Therefore the positive zeros of (B) are
$k_{n}=n\pi/L$, $n=1$, $2$, $3$, \dots, so the positive
eigenvalues (SL) are   $\lambda_{n}=n^{2}\pi^{2}/L^{2}$, $n=1$, $2$, $3$,
\dots.
\end{alist}


\item
Suppose $\lambda$ is an eigenvalue and $y$ is an associated eigenfunction.
From the solution of Exercise~13.2.31,
\[\tag{A}
\lambda\int_{a}^{b}r(x)y^{2}(x)\,dx
=p(a)y(a)y'(a)-p(b)y(b)y'(b)+
\int_{a}^{b}p(x)(y'(x))^{2}\,dx.
\]
If $\alpha\beta=0$ then either $y(a)=0$ or $y'(a)=0$, so $y(a)y'(a)=0$.
If $\alpha\beta<0$ then $y(a)=-\dst{\frac{\beta}{\alpha}}y'(a)$, so
\[\tag{B}
y(a)y'(a)=-\dst{\frac{\beta}{\alpha}}(y'(a))^{2}.
\]
Moreover, $y'(a)\ne0$ because if $y'(a)=0$ then $y(a)=0$, from (B), and
$y\equiv0$, a contradiction. Since $-\dst{\frac{\beta}{\alpha}}>0$
if $\alpha\beta<0$, we conclude that if $\alpha\beta\le 0$, then
\[\tag{C}
p(a)y(a)y'(a)\ge 0,
\]
with equality if and only if $\rho\delta=0$.
A similar argument shows that if $\rho\delta\ge 0$, then
\[\tag{D}
p(b)y(b)y'(b)\ge 0,
\]
with equality if and only if $\alpha\beta=0$.
Since
 $(\alpha\beta)^{2}+(\rho\delta)^{2}>0$, the inequality must hold in at
least one of (C) and (D). Now (A) implies that $\lambda>0$.
\end{sectionSolutions}

\end{document}
